\section{Analytic foundations and the transcendental Nakamaye theorem}

\begin{convention}[Geometric and pluripotential notation]
\label{ana:geometric-setting}
\uses{psh:local,psh:local-integrability,dist:psh-regular,dist:psh-germ-levi,dist:positive-measure,geom:flat-closed,dist:psh-representative-existence,fpgeom:hermitian-kahler,fpgeom:exterior-compatibility,lc:assembled,sing:local-potential,sing:multiplier-local,sing:multiplier-type-invariance,cp:inverse-coefficients,cp:inverse-current,fpglobal:theta-geometric-invariance,fpglobal:chart-independence,fpglobal:potential-independence,fpglobal:smooth-theta-criterion,as:analytic-potentials,as:two-matrices,pp:noncharging,pp:haar-noncharging,gc:smooth-global,fpenv:v-psh,fpenv:v-minimal,fpenv:p-psh,fpma:bt-existence,fpma:truncation-boundary,mpp:qe,mpp:noncharging-ae,mpp:qpsh-finite,fprooftop:definition,fprooftop:qe-invariance,gcp:global-properties,gcp:theta-family,bttr:step,bttr:consistency,bttr:mass-boundary,bc:quotient,bc:class-module,npl:measure,npl:restriction,npl:finite-mass,npl:limit-noncharging,npl:open-restriction,npl:bounded-agreement,fpmma:canonical,fpmma:bt-diagonal,fpmma:nonpluripolar-boundary,lpot:existence-global,lpot:semipositive-enlargement,posclass:quotient-predicates,posclass:class-nef-cone,poscon:background-big,poscon:big-psef,mnpp:limit,mnpp:restriction,mnpp:finite-mass,mnpp:noncharging,mnpp:open-locality,mnpp:bounded-agreement,mnpp:diagonal,mnpp:empty,gnpp:global-measure,gnpp:global-characterization,gnpp:finite-global-mass,gnpp:noncharging}
\source{transcendental-nakamaye-copy4-20260911:sec:prelim,transcendental-nakamaye-copy4-20260911:sec:qpsh,transcendental-nakamaye-copy4-20260911:sec:npp,transcendental-nakamaye-copy4-20260911:sec:nK}
Let \(X\) be a nonempty connected compact complex manifold of complex
dimension \(n>0\), with a K\"ahler form \(\omega\), and put
\(d d^c=(2\pi)^{-1}\mathrm i\partial\bar\partial\). A real
\((1,1)\)-class is pseudoeffective if it contains a closed positive
current, and is big if it contains a current dominating
\(\delta\omega\) for some \(\delta>0\). Such a current is called a
K\"ahler current. For a smooth closed real \((1,1)\)-form \(\theta\),
write \(\operatorname{PSH}(X,\theta)\) for the upper semicontinuous
quasi-plurisubharmonic functions in their canonical representatives,
\(u\not\equiv-\infty\), with
\(\theta_u:=\theta+d d^c u\ge0\), and put
\[
 V_\theta=\sup\{u\in\operatorname{PSH}(X,\theta):u\le0\}.
\]
The notation \(u\preceq v\) means \(u\le v+C\) for a constant \(C\).
The multiplier ideal \(\mathcal I(\lambda u)\) consists of holomorphic
germs \(f\) such that \(|f|^2e^{-\lambda u}\) is locally integrable.
For a coherent ideal \(\mathfrak a=(f_1,\ldots,f_N)\), put
\(\psi_{\mathfrak a}=\log\sum_j|f_j|^2\), considered up to locally
bounded functions. The restriction of an ideal to a closed complex
submanifold \(V\) is the image ideal
\(\operatorname{Res}_V\mathfrak a=\operatorname{im}(\mathfrak a\to
\mathcal O_V)\). A function has analytic singularities of type
\((c,\mathfrak a)\) if locally it equals
\(c\psi_{\mathfrak a}+O(1)\), where \(c\in\mathbb Q_{>0}\).
The phrase \emph{smooth analytic singularities} requires a smooth
remainder and allows a coefficient \(c\in\mathbb R_{>0}\). The phrase
\emph{analytic singularities with a remainder smooth off the poles} asks instead
only that the remainder be bounded near the point and \(C^\infty\) where the local
generators do not all vanish (Definition~\ref{as:smooth-off-pole-type}); this is the
convention the restricted volume is taken in, and it is the one Bergman-kernel
regularization produces.

For positive currents, angle brackets denote their non-pluripolar
product, and
\(\int_X\langle\theta_u^n\rangle\) is its total mass. Define
\[
 \begin{aligned}
 P_\theta[u]&=\sup{}^*\{v\in\operatorname{PSH}(X,\theta):
 v\le0,\ v\preceq u\},\\
 P_\theta(f)&=\sup{}^*\{v\in\operatorname{PSH}(X,\theta):
 v\le f\text{ outside a pluripolar set}\}.
 \end{aligned}
\]
Here \(\sup^*\) is upper semicontinuous regularization. If \(\alpha\)
is big, its non-K\"ahler locus is
\[
 E_{\mathrm{nK}}(\alpha)=
 \bigcap_{T\in\alpha,\ T\text{ K\"ahler current}}
 \{x:\nu(T,x)>0\}.
\]
For a pseudoeffective non-big class set \(E_{\mathrm{nK}}(\alpha)=X\).
For a pseudoeffective class with smooth representative \(\theta\), set
\[
 E_{\mathrm{nn}}(\alpha)=\bigcup_{\epsilon>0}
 E_{\mathrm{nK}}(\alpha+\epsilon\{\omega\}).
\]
A modification is a finite composition of blowups with smooth centers.
For an effective rational divisor \(D\) with a smooth Hermitian metric,
write \(\psi_D=\log|s_D|^2\), and \(\Theta_D\) for the curvature form of the
metric.
\end{convention}

\begin{proposition}[The non-nef locus of a big class through \(V_\theta\)]
\label{ana:nonnef-lelong-locus}
\uses{ana:geometric-setting,ana:nonkahler-locus}
\source{transcendental-nakamaye-copy4-20260911:sec:nK}
\lean{TranscendentalNakamaye.Pending.NonNefLocusLelong.nonNefLocus_eq_setOf_lelongNumber_pos,TranscendentalNakamaye.Pending.NonNefLocusLelong.nonNefLocus_subset_lelongNumber_pos,TranscendentalNakamaye.Pending.NonNefLocusLelong.kahlerCurrentAddSmul,TranscendentalNakamaye.Pending.NonNefLocusLelong.lelongNumber_eq_top_of_dim_zero,TranscendentalNakamaye.Pending.NonNefLocusLelong.nonNefLocus_eq_univ_of_dim_zero}
\leanok
Let \(\alpha\) be a pseudoeffective class with smooth representative \(\theta\), and let \(E_{\mathrm{nn}}(\alpha)\) be its non-nef locus
(Convention~\ref{ana:geometric-setting}). For a big class this is also \(\{x:\nu(V_\theta,x)>0\}\).
\end{proposition}
\begin{proof}
\uses{ana:minimal-singularity-lelong-vanishing,ana:nonkahler-locus}
\leanok
The inclusion \(\{x:\nu(V_\theta,x)>0\}\subseteq E_{\mathrm{nn}}(\alpha)\) is the last assertion of
Proposition~\ref{ana:minimal-singularity-lelong-vanishing}. For the reverse inclusion let \(u\) be any \(\theta\)-PSH potential and
\(\varepsilon>0\). Then \(\theta+\varepsilon\omega+dd^cu\ge\varepsilon\omega\), so \(\theta+\varepsilon\omega+dd^cu\) is a K\"ahler current in
\(\alpha+\varepsilon\{\omega\}\) whose Lelong number at \(x\) is \(\nu(u,x)\). A point of \(E_{\mathrm{nK}}(\alpha+\varepsilon\{\omega\})\) therefore
has \(\nu(u,x)>0\), and so does every point of \(E_{\mathrm{nn}}(\alpha)\). For a big class \(V_\theta\) is a \(\theta\)-PSH potential
(Proposition~\ref{ana:minimal-singularity-lelong-vanishing}), which gives the reverse inclusion.
\end{proof}
The proof of the reverse inclusion shows more: \(E_{\mathrm{nn}}(\alpha)\subseteq\{x:\nu(u,x)>0\}\) for every \(\theta\)-PSH potential \(u\),
for any class.

\begin{lemma}[Lelong--Poincar\'e for a metered divisor]
\label{ana:divisor-lelong-poincare}
\uses{ana:geometric-setting}
\source{transcendental-nakamaye-copy4-20260911:sec:qpsh}
\lean{TranscendentalNakamaye.Pending.LelongPoincareDivisor.lelongPoincare,TranscendentalNakamaye.Pending.LelongPoincareDivisor.ddCWeight_eq,TranscendentalNakamaye.Pending.LelongPoincareDivisor.integrationCurrent}
\leanok
For an effective rational divisor \(D\) with a smooth Hermitian metric, with \(\psi_D=\log|s_D|^2\) and curvature form \(\Theta_D\)
(Convention~\ref{ana:geometric-setting}),
\(d d^c\psi_D=[D]-\Theta_D\).
\end{lemma}
\begin{proof}
\leanok
\uses{bou:divisor-class}
This is the first assertion of Lemma~\ref{bou:divisor-class} for \(E=D\).
\end{proof}
Off \(\operatorname{Supp}D\) the identity reads \(d d^c\psi_D=-\Theta_D\), which is formalized
(Lemma~\ref{ana:divisor-weight-smooth-potential}).

\begin{lemma}[Lelong--Poincar\'e in one variable]
\label{ana:lelong-poincare-one-variable}
\source{demailly-cadg:III.2.15}
\lean{TranscendentalNakamaye.Pending.LelongPoincareOneVariable.integral_log_norm_mul_laplacian,TranscendentalNakamaye.Pending.LelongPoincareOneVariable.integral_log_norm_sub_mul_laplacian,TranscendentalNakamaye.Pending.LelongPoincareOneVariable.integral_log_norm_mul_laplacian_eq_zero,TranscendentalNakamaye.Pending.LelongPoincareOneVariable.integral_mul_laplacian_eq,TranscendentalNakamaye.Pending.LelongPoincareOneVariable.locallyIntegrable_log_norm,TranscendentalNakamaye.Pending.LelongPoincareOneVariable.integral_kernel,TranscendentalNakamaye.Pending.LelongPoincareOneVariable.laplacian_regLog}
\leanok
Let \(U\subset\mathbf C\) be open and let \(f\) be meromorphic on \(U\), vanishing identically near no point of \(U\), so that its order
\(\operatorname{ord}_u(f)\in\mathbf Z\) is defined at every \(u\in U\). Then for every \(\chi\in C^\infty_c(U,\mathbf R)\) (extended by zero to \(\mathbf C\)),
\[
 \int_{\mathbf C}\log|f|\,\Delta\chi\,d\lambda=2\pi\sum_{u\in U}\operatorname{ord}_u(f)\,\chi(u),
\]
where \(\Delta=\partial_x^2+\partial_y^2\), \(\lambda\) is Lebesgue measure, and the sum has finitely many nonzero terms. Since
\(dd^c\varphi=(4\pi)^{-1}\Delta\varphi\,d\lambda\) on functions for \(dd^c=(2\pi)^{-1}\mathrm i\partial\bar\partial\), this says that
\(dd^c\log|f|^2=\sum_{u}\operatorname{ord}_u(f)\,\delta_u\) as currents on \(U\): the Lelong--Poincar\'e equation in dimension one.
\end{lemma}
\begin{proof}
First, \(\int F\,\Delta\chi=\int\Delta F\,\chi\) whenever \(F\) is smooth on a neighbourhood of \(\operatorname{supp}\chi\) (multiply \(F\) by a cutoff and
integrate by parts twice). Hence \(\int\log|g|\,\Delta\chi=0\) when \(g\) is holomorphic without zeros near \(\operatorname{supp}\chi\), since
\(\log|g|\) is harmonic there.

The model is \(\int\log|z-u|\,\Delta\chi(z)\,d\lambda(z)=2\pi\chi(u)\); by translation take \(u=0\). The function \(\log|z|\) is locally integrable
(in polar coordinates \(r\log r\) is continuous). Put \(L_\varepsilon(z)=\tfrac12\log(|z|^2+\varepsilon^2)\), which is smooth with
\[
 \Delta L_\varepsilon(z)=\frac{2\varepsilon^2}{(|z|^2+\varepsilon^2)^2}=\varepsilon^{-2}K(z/\varepsilon),\qquad K(w)=\frac{2}{(|w|^2+1)^2},\qquad
 \int_{\mathbf C}K\,d\lambda=2\pi\int_0^\infty\frac{2r\,dr}{(r^2+1)^2}=2\pi .
\]
Then \(\int L_\varepsilon\,\Delta\chi=\int\Delta L_\varepsilon\,\chi=\int K(w)\,\chi(\varepsilon w)\,d\lambda(w)\), which tends to \(2\pi\chi(0)\) as
\(\varepsilon\to0\) by dominated convergence (bound \(K\sup|\chi|\)). On the other side, for \(0<\varepsilon\le1\) and \(z\ne0\),
\(\log|z|\le L_\varepsilon(z)\le L_1(z)\), so dominated convergence gives \(\int L_\varepsilon\,\Delta\chi\to\int\log|z|\,\Delta\chi\).

In general, choose an open \(W\) with \(\operatorname{supp}\chi\subset W\) and \(\overline W\subset U\) compact. The divisor of \(f\) on \(W\) has finite
support, and off a discrete subset of \(W\), hence \(\lambda\)-almost everywhere on \(W\),
\(\log|f|=\sum_u\operatorname{ord}_u(f)\log|z-u|+\log|g|\) with \(g\) holomorphic without zeros on \(W\). Integrate against \(\Delta\chi\), which
vanishes off \(\operatorname{supp}\chi\), and apply the two previous paragraphs; for \(u\notin W\) both \(\operatorname{ord}_u(f)\) (on \(W\)) and
\(\chi(u)\) vanish.
\end{proof}
This is the first step towards Lemma~\ref{ana:divisor-lelong-poincare}; the local equation \(dd^c\log|g|^2=[\operatorname{div}g]\) in several
variables is Theorem~\ref{ana:divisor-integration-current}, and the gluing is Theorem~\ref{ana:metered-divisor-lelong-poincare}.

\begin{lemma}[Lelong--Poincar\'e in slice form]
\label{ana:lelong-poincare-slices}
\uses{ana:lelong-poincare-one-variable}
\source{demailly-cadg:III.2.15}
\lean{TranscendentalNakamaye.Pending.LelongPoincareSlices.integral_log_norm_mul_sliceLaplacian,TranscendentalNakamaye.Pending.LelongPoincareSlices.integral_log_norm_mul_sliceLaplacian_pow_mul,TranscendentalNakamaye.Pending.LelongPoincareSlices.integral_log_norm_mul_sliceLaplacian_pow_mul',TranscendentalNakamaye.Pending.LelongPoincareSlices.locallyIntegrableOn_log_norm_pow_mul,TranscendentalNakamaye.Pending.LelongPoincareSlices.locallyIntegrable_log_norm_fst,TranscendentalNakamaye.Pending.LelongPoincareSlices.sliceLaplacian,TranscendentalNakamaye.Pending.LelongPoincareSlices.sliceLaplacian_eq,TranscendentalNakamaye.Pending.LelongPoincareSlices.sliceSet,TranscendentalNakamaye.Pending.LelongPoincareSlices.divisor_pow_mul,TranscendentalNakamaye.Pending.LelongPoincareSlices.meromorphicOrderAt_pow_mul}
\leanok
Let \(E\) be a real normed space with an \(s\)-finite Borel measure \(\mu\), let \(U\subset\mathbf C\times E\) be open, and put
\(U_w=\{s\in\mathbf C:(s,w)\in U\}\). Let \(G\) be a function on \(\mathbf C\times E\) such that \(\log|G|\) is locally integrable on \(U\) for
\(\lambda\otimes\mu\) and, for \(\mu\)-almost every \(w\), the slice \(G(\cdot,w)\) is meromorphic on \(U_w\) and vanishes identically near no point
of \(U_w\). Write \(\Delta_s\chi(s,w)=\Delta\bigl(\chi(\cdot,w)\bigr)(s)\) for the Laplacian in the first variable. Then for every
\(\chi\in C^\infty_c(U,\mathbf R)\)
\[
 \int_{\mathbf C\times E}\log|G|\,\Delta_s\chi\;d(\lambda\otimes\mu)
 =2\pi\int_E\sum_{s\in U_w}\operatorname{ord}_s\bigl(G(\cdot,w)\bigr)\,\chi(s,w)\;d\mu(w).
\]
In particular, let \(\mu\) be locally finite, \(E\) second countable, \(m\in\mathbf N\), and \(G(s,w)=s^mv(s,w)\) with \(v\) continuous and without
zeros on \(U\) and holomorphic in \(s\) on each \(U_w\). Then \(\log|G|\) is locally integrable on \(U\) and
\[
 \int_{\mathbf C\times E}\log|s^mv|\,\Delta_s\chi\;d(\lambda\otimes\mu)=2\pi m\int_E\chi(0,w)\,d\mu(w).
\]
\end{lemma}
\begin{proof}
\uses{ana:lelong-poincare-one-variable}
\leanok
The function \(\Delta_s\chi\) is \(D^2\chi\) evaluated on the pairs of vectors \(((1,0),(1,0))\) and \(((i,0),(i,0))\), so it is continuous and
vanishes off \(\operatorname{supp}\chi\). Hence \(\log|G|\,\Delta_s\chi\) is integrable (\(\log|G|\) is integrable on the compact set
\(\operatorname{supp}\chi\subset U\)), and Fubini's theorem writes the left-hand side as
\(\int_E\bigl(\int_{\mathbf C}\log|G(s,w)|\,\Delta(\chi(\cdot,w))(s)\,d\lambda(s)\bigr)\,d\mu(w)\). Each slice \(\chi(\cdot,w)\) is smooth, and its
support is compact and contained in \(U_w\). So for almost every \(w\), Lemma~\ref{ana:lelong-poincare-one-variable} on \(U_w\) gives
\(2\pi\sum_s\operatorname{ord}_s(G(\cdot,w))\,\chi(s,w)\) for the inner integral.

For the model, each slice \(s\mapsto s^mv(s,w)\) is holomorphic on \(U_w\). It has order \(m\) at \(0\) when \(0\in U_w\), and order \(0\) at every
other point of \(U_w\), because \(v\) has no zeros. When \(0\notin U_w\), \((0,w)\notin\operatorname{supp}\chi\), so \(\chi(0,w)=0\). In both
cases the sum is \(m\,\chi(0,w)\). Off the \((\lambda\otimes\mu)\)-null set \(\{s=0\}\) we have \(\log|G|=m\log|s|+\log|v|\).
\begin{itemize}
\item The first term is locally integrable: near a point it is integrable on the product of a closed disc with a set of finite \(\mu\)-measure,
because \(\log|s|\) is locally integrable on \(\mathbf C\).
\item The second term is continuous on \(U\).
\end{itemize}
\end{proof}
For \(E=\mathbf C^{n-1}\) with Lebesgue measure, \(dd^c\chi\wedge dV(w)=(4\pi)^{-1}\Delta_s\chi\,d\lambda\) for the \((n-1,n-1)\)-form
\(dV(w)\) of the last \(n-1\) coordinates. So the lemma says that \(dd^c\log|G|^2\) paired with \(\chi\,dV(w)\) is the slice measure
\(\int\sum_s\operatorname{ord}_s(G(\cdot,w))\,\chi(s,w)\,dV(w)\) of the divisor of \(G\). No analysis of the singular set of the divisor is needed.
For the model, this is the current of the hyperplane \(\{s=0\}\) with multiplicity \(m\). The integration current \([\operatorname{div}g]\) and its
identification with these slice measures in every unitary frame are Lemmas~\ref{ana:divisor-singular-projection} and~\ref{ana:divisor-coarea} and
Theorem~\ref{ana:divisor-integration-current}; the gluing of the local presentations is Theorem~\ref{ana:metered-divisor-lelong-poincare}.

\begin{lemma}[The Levi coefficient of \(dd^c\log|G|^2\) is the slice measure]
\label{ana:lelong-poincare-levi}
\uses{ana:lelong-poincare-slices,lc:ddc,dist:regular}
\source{demailly-cadg:III.2.15}
\lean{TranscendentalNakamaye.Pending.LelongPoincareLevi.ddC_log_normSq_apply_first,TranscendentalNakamaye.Pending.LelongPoincareLevi.ddC_log_normSq_pow_mul_apply_first,TranscendentalNakamaye.Pending.LelongPoincareLevi.leviForm_first,TranscendentalNakamaye.Pending.LelongPoincareLevi.locallyIntegrableOn_log_norm_of_log_normSq}
\leanok
Let \(E\) be a second countable complex normed space with an \(s\)-finite Borel measure, let \(\Omega\subset\mathbf C\times E\) be open, and let
\(G\) be as in Lemma~\ref{ana:lelong-poincare-slices} (on \(\Omega\)), with \(\log|G|^2\) locally integrable on \(\Omega\). Regard
\(dd^c\log|G|^2\) as the distribution-valued two-form of Definition~\ref{lc:ddc}. Then for every \(\varphi\in\mathcal D(\Omega)\)
\[
 dd^c\log|G|^2\bigl((1,0),\,i(1,0)\bigr)(\varphi)=\int_E\sum_{s\in\Omega_w}\operatorname{ord}_s\bigl(G(\cdot,w)\bigr)\,\varphi(s,w)\,dw .
\]
For \(G(s,w)=s^mv(s,w)\) as in Lemma~\ref{ana:lelong-poincare-slices} (and a locally finite measure), the right-hand side is
\(m\int_E\varphi(0,w)\,dw\).
\end{lemma}
\begin{proof}
\uses{ana:lelong-poincare-slices,lc:levi,dist:test-levi-pairing}
\leanok
By Proposition~\ref{lc:levi}, \(dd^cT(v,iv)=\pi^{-1}\mathcal L_vT\). For the regular distribution of a locally integrable \(u\),
\(\mathcal L_vT_u(\varphi)=\int u\,L_\varphi(v)\) (Lemma~\ref{dist:test-levi-pairing}), where \(L_\varphi(v)=\frac14(D^2\varphi(v,v)+D^2\varphi(iv,iv))\). For \(v=(1,0)\) this is
\(\frac14\Delta_s\varphi\). With \(\log|G|^2=2\log|G|\), the left-hand side is \(\frac1{2\pi}\int\log|G|\,\Delta_s\varphi\). This equals the right-hand side by
Lemma~\ref{ana:lelong-poincare-slices}, applied to \(\chi=\varphi\); the local integrability of \(\log|G|=\frac12\log|G|^2\) it needs follows
from the hypothesis here.
\end{proof}

\begin{lemma}[The area formula for \(C^1\) embeddings]
\label{ana:hausdorff-area-formula}
\lean{TranscendentalNakamaye.Pending.HausdorffAreaFormula.euclideanHausdorffMeasure_image_eq_lintegral,TranscendentalNakamaye.Pending.HausdorffAreaFormula.restrict_euclideanHausdorffMeasure_image,TranscendentalNakamaye.Pending.HausdorffAreaFormula.lintegral_image_eq_lintegral,TranscendentalNakamaye.Pending.HausdorffAreaFormula.integral_image_eq_integral,TranscendentalNakamaye.Pending.HausdorffAreaFormula.image_le_of_rel_close,TranscendentalNakamaye.Pending.HausdorffAreaFormula.le_image_of_rel_close,TranscendentalNakamaye.Pending.HausdorffAreaFormula.exists_local_bounds,TranscendentalNakamaye.Pending.HausdorffAreaFormula.continuous_normDet}
\leanok
Let \(E\) be a real inner product space of finite dimension \(k\) with Lebesgue measure \(\lambda\), let \(F\) be a real inner product space, and
write \(\mathcal H^k\) for the \(k\)-dimensional Hausdorff measure of \(F\), normalised to agree with Lebesgue measure on \(k\)-dimensional
spaces. For a linear map \(L\colon E\to F\) put \(J(L)=\sqrt{\det(L^*L)}\). Let \(f\) be of class \(C^1\) on an open set \(U\subset E\), and let
\(s\subset U\) be measurable, with \(f\) injective on \(s\) and \(Df(x)\) injective for every \(x\in s\). Then
\[
 \mathcal H^k(f(s))=\int_sJ(Df(x))\,d\lambda(x),\qquad\int_{f(s)}\varphi\,d\mathcal H^k=\int_sJ(Df(x))\,\varphi(f(x))\,d\lambda(x)
\]
for every nonnegative Borel function \(\varphi\) on \(F\), and for every real \(\varphi\) that is \(\mathcal H^k\)-measurable on \(f(s)\) (both sides
then being Bochner integrals).
\end{lemma}
\begin{proof}
\leanok
For a linear map \(A\) the formula \(\mathcal H^k(A(S))=J(A)\lambda(S)\) is classical. Suppose that on a set \(S\)
\(\|f(x)-f(y)-A(x-y)\|\le\varepsilon\|A(x-y)\|\) with \(\varepsilon<1\). Then \(f(x)\) depends only on \(Ax\), and \(f(S)\) is the image of \(A(S)\) under a
map that is \((1+\varepsilon)\)-Lipschitz; conversely \(A(S)\) is the image of \(f(S)\) under a \((1-\varepsilon)^{-1}\)-Lipschitz map. A
\(K\)-Lipschitz map multiplies \(\mathcal H^k\) by at most \(K^k\), so
\[
 \mathcal H^k(f(S))\le(1+\varepsilon)^kJ(A)\lambda(S),\qquad J(A)\lambda(S)\le(1-\varepsilon)^{-k}\mathcal H^k(f(S)).
\]
Let \(x_0\in s\) and \(A=Df(x_0)\), injective, so \(\|v\|\le K\|Av\|\) for some \(K\). By continuity of \(Df\) there is a ball \(B\) around \(x_0\)
inside \(U\) on which \(\|Df(x)-A\|\le\varepsilon/K\). The mean value theorem on the convex set \(B\) gives the relative closeness above on every
\(S\subset B\), and the same inequality for the linear maps \(Df(x)\), \(x\in B\) (applied to a ball, of positive finite measure), gives
\(J(Df(x))\le(1+\varepsilon)^kJ(A)\) and \(J(A)\le(1-\varepsilon)^{-k}J(Df(x))\).

Cover \(s\) by countably many such balls, make them disjoint, and intersect with \(s\). The images of the pieces are measurable (an injective
continuous image of a Borel set in a Polish space is Borel) and pairwise disjoint (\(f\) is injective on \(s\)), so both sides of the formula are
sums over the pieces, and on each piece they agree up to the factor \(c(\varepsilon)=(1+\varepsilon)^k(1-\varepsilon)^{-k}\). Hence each side
is at most \(c(\varepsilon)\) times the other, and \(\varepsilon\to0\) gives equality. Applying the set formula to \(s\cap f^{-1}(B)\) for Borel
\(B\) identifies \(\mathcal H^k\) on \(f(s)\) with the push-forward of \(J(Df)\lambda\) on \(s\), and the integral forms follow; \(J(Df)\) is
measurable since \(J\) is continuous (\(J(L)^2=\det(L^*L)\)).
\end{proof}

\begin{definition}[Regular points of a divisor, multiplicity and transverse form]
\label{ana:divisor-regular-points}
\lean{TranscendentalNakamaye.Pending.DivisorRegularLocus.RegularPresentation,TranscendentalNakamaye.Pending.DivisorRegularLocus.regularLocus,TranscendentalNakamaye.Pending.DivisorRegularLocus.divisorMult,TranscendentalNakamaye.Pending.DivisorRegularLocus.normalFun,TranscendentalNakamaye.Pending.DivisorRegularLocus.lineForm,TranscendentalNakamaye.Pending.DivisorRegularLocus.transverseForm,TranscendentalNakamaye.Pending.DivisorRegularLocus.RegularPresentation.analyticOrderAt_line,TranscendentalNakamaye.Pending.DivisorRegularLocus.RegularPresentation.le_analyticOrderAt_line,TranscendentalNakamaye.Pending.DivisorRegularLocus.RegularPresentation.m_eq,TranscendentalNakamaye.Pending.DivisorRegularLocus.RegularPresentation.exists_fderiv_eq_smul,TranscendentalNakamaye.Pending.DivisorRegularLocus.RegularPresentation.divisorMult_eq,TranscendentalNakamaye.Pending.DivisorRegularLocus.RegularPresentation.transverseForm_eq,TranscendentalNakamaye.Pending.DivisorRegularLocus.lineForm_apply_I,TranscendentalNakamaye.Pending.DivisorRegularLocus.lineForm_swap,TranscendentalNakamaye.Pending.DivisorRegularLocus.lineForm_I_I,TranscendentalNakamaye.Pending.DivisorRegularLocus.transverseLocus,TranscendentalNakamaye.Pending.DivisorRegularLocus.exists_isOpen_regularLocus,TranscendentalNakamaye.Pending.DivisorRegularLocus.exists_isOpen_transverseLocus}
\leanok
Let \(V\) be a complex inner product space and \(g\colon V\to\mathbf C\). A point \(x\) is a \emph{regular point of the divisor of \(g\)} if on
an open neighbourhood \(U\) of \(x\) one has \(g=h^mv\) with \(h,v\) holomorphic on \(U\), \(h(x)=0\), \(dh\neq0\) and \(v\neq0\) on \(U\), and
\(m\ge1\). If \(u\in V\) and \(dh(x)u\neq0\), then \(s\mapsto g(x+su)\) vanishes to order exactly \(m\) at \(0\); if \(dh(x)u=0\) it vanishes to order at least
\(2m\). Consequently two such presentations at \(x\) have the same exponent \(m\) and proportional differentials \(dh(x)\). The \emph{regular locus}
\(Z_{\mathrm{reg}}(g)\), the \emph{multiplicity} \(m_g(x)\) and the \emph{transverse form}
\[
 \tau_x(v,w)=\frac{\operatorname{Im}\bigl(\overline{\ell(v)}\,\ell(w)\bigr)}{\|\ell\|^2},\qquad \ell=dh(x),
\]
are therefore well defined. The form \(\tau_x\) is real, alternating and of type \((1,1)\); it is the area form of the complex line orthogonal to the
tangent hyperplane \(\ker\ell\), and \(\tau_x(v,iv)=|\ell(v)|^2/\|\ell\|^2\ge0\). For a direction \(u\) the \emph{transverse locus} is the set of
regular points with \(\tau_x(u,iu)\neq0\), equivalently \(dh(x)u\neq0\). Both loci are the zero set of \(g\) intersected with an open set.
\end{definition}

\begin{lemma}[Local graphs of the divisor and their area factor]
\label{ana:divisor-local-graph}
\uses{ana:divisor-regular-points}
\lean{TranscendentalNakamaye.Pending.DivisorRegularLocus.Splitting,TranscendentalNakamaye.Pending.DivisorRegularLocus.exists_splitting,TranscendentalNakamaye.Pending.DivisorRegularLocus.RegularPresentation.exists_graph,TranscendentalNakamaye.Pending.DivisorRegularLocus.RegularPresentation.exists_graph_slices,TranscendentalNakamaye.Pending.DivisorRegularLocus.normDet_graphDeriv,TranscendentalNakamaye.Pending.DivisorRegularLocus.lineForm_mul_graph,TranscendentalNakamaye.Pending.DivisorRegularLocus.RegularPresentation.normDet_mul_transverseForm,TranscendentalNakamaye.Pending.DivisorRegularLocus.det_one_add_comp_comm}
\leanok
Let \(V\) be finite-dimensional, \(u\in V\) a unit vector and \(V=\mathbf Cu\oplus\iota(E)\) an orthogonal splitting, \(\iota\) an isometric embedding with
projection \(\pi\). Let \(g=h^mv\) be a regular presentation at \(x\) with \(dh(x)u\neq0\). Then there are open sets \(N\ni x\) in \(V\) and \(W\ni\pi(x)\)
in \(E\) and a holomorphic \(k\) on \(W\) such that the zeros of \(g\) in \(N\) are exactly the points \(k(w)u+\iota w\), \(w\in W\); a point
\(su+\iota w\in N\) is a zero of \(g\) if and only if \(w\in W\) and \(s=k(w)\), and there \(s\mapsto g(\iota w+su)\) vanishes to order \(m\). If
\(V=\mathbf C^n\) and \(E=\mathbf C^{n-1}\), the graph map \(G(w)=k(w)u+\iota w\) satisfies
\[
 J(DG(w))\;\tau_{G(w)}(u,iu)=1\qquad(w\in W),
\]
where \(J\) is the volume factor of Lemma~\ref{ana:hausdorff-area-formula}.
\end{lemma}
\begin{proof}
\leanok
The map \(y\mapsto(h(y),\pi(y))\) has invertible differential at \(x\) (its inverse is \((c,w)\mapsto\frac{c-\ell(\iota w)}{\ell(u)}u+\iota w\),
\(\ell=dh(x)\)); the inverse function theorem gives \(N\), \(W\) and \(k(w)=\langle u,F^{-1}(0,w)\rangle\), holomorphic. A zero \(su+\iota w\) in \(N\)
satisfies \(h=0\), hence equals \(F^{-1}(0,w)\), so \(s=k(w)\); the order along the line is \(m\) by Definition~\ref{ana:divisor-regular-points},
since \(dh\,u\neq0\) on \(N\). For the area factor write \(\kappa=dk(w)=\langle b,\cdot\rangle\). Then \(DG(w)e=\kappa(e)u+\iota e\) and
\(\langle DG\,e,DG\,e'\rangle=\langle e+\kappa(e)b,e'\rangle\), so \(J(DG)^2=\det_{\mathbf R}(1+b\,b^*)=\det_{\mathbf R}\bigl((1+\|b\|^2)\,\mathrm{id}_{\mathbf C}\bigr)
=(1+\|\kappa\|^2)^2\) by \(\det(1+AB)=\det(1+BA)\). On the other hand \(h\circ G=0\) gives \(\ell\circ DG=0\) for \(\ell=dh(G(w))\); writing
\(\ell=\langle a,\cdot\rangle\), the component of \(a\) in \(\iota(E)\) is \(-\overline{\ell(u)}\,\iota b\), so \(\|\ell\|^2=|\ell(u)|^2(1+\|\kappa\|^2)\) and
\(\tau(u,iu)=|\ell(u)|^2/\|\ell\|^2=(1+\|\kappa\|^2)^{-1}\).
\end{proof}

\begin{lemma}[Almost every line misses the non-transverse points of a divisor]
\label{ana:divisor-singular-projection}
\uses{ana:divisor-regular-points,ana:divisor-local-graph}
\lean{TranscendentalNakamaye.Pending.DivisorSingularProjection.measure_image_bad_eq_zero,TranscendentalNakamaye.Pending.DivisorSingularProjection.exists_nhds_null_top,TranscendentalNakamaye.Pending.DivisorSingularProjection.exists_nhds_null_nat,TranscendentalNakamaye.Pending.DivisorSingularProjection.lineOrder_eq_iff,TranscendentalNakamaye.Pending.DivisorSingularProjection.iteratedDeriv_line,TranscendentalNakamaye.Pending.DivisorSingularProjection.dirDeriv,TranscendentalNakamaye.Pending.DivisorSingularProjection.infiniteOrderSet,TranscendentalNakamaye.Pending.DivisorSingularProjection.badStratum,TranscendentalNakamaye.Pending.DivisorSingularProjection.measure_image_null_of_local}
\leanok
Let \(g\) be holomorphic on an open set \(\Omega\) of a finite-dimensional complex inner product space \(V\), vanishing identically near no point of
\(\Omega\). Let \(u\) be a unit vector, \(V=\mathbf Cu\oplus\iota(E)\) an orthogonal splitting with projection \(\pi\), and \(\mu\) a Haar measure on
\(E\). Then the zeros of \(g\) in \(\Omega\) that are not in the transverse locus of Definition~\ref{ana:divisor-regular-points} (singular points of the
divisor, and regular points where it is tangent to \(u\)) have \(\mu\)-null image under \(\pi\). Equivalently, for \(\mu\)-almost every \(w\) every zero
of \(g\) on the line \(\{su+\iota w\}\cap\Omega\) is a transverse regular point.
\end{lemma}
\begin{proof}
\leanok
Stratify the zeros \(y\) by the order \(j\in\{1,2,\dots,\infty\}\) of \(s\mapsto g(y+su)\) at \(0\); it is the least \(i\) with \(\partial_u^ig(y)\neq0\).
It suffices to find, around every zero of order \(j\), a neighbourhood whose zeros of order \(j\) that are not transverse project to a null set: a
countable subcover finishes.

Order \(\infty\). Near \(y_0\), take a product neighbourhood \(\{su+\iota w:|s-s_0|<\varepsilon,\ \|w-w_0\|<\varepsilon\}\subset\Omega\). If \(y=s_1u+\iota w\)
has infinite order, then \(s\mapsto g(su+\iota w)\) vanishes near \(s_1\), hence on the whole disc, by the identity theorem. So the projection lies in
the set of \(w\) over which \(g\) vanishes on the whole disc. If some \(g(su+\iota w)\neq0\), this set lies in the zero set of the holomorphic
function \(w\mapsto g(su+\iota w)\), which is null; otherwise \(g\) vanishes on an open set, which is excluded.

Order \(j<\infty\). Put \(D=\partial_u^{j-1}g\). Then \(D(y_0)=0\), \(\partial_uD(y_0)\neq0\), and every zero of \(g\) of order \(j\) is a zero of \(D\).
By Lemma~\ref{ana:divisor-local-graph} applied to \(D=D^1\cdot1\), the zeros of \(D\) near \(y_0\) form a graph \(\{k(w)u+\iota w\}\). In the
coordinates \(F(z,t)=g\bigl((k(w_0+z)+t)u+\iota(w_0+z)\bigr)\), Taylor division in \(t\) gives \(F=t^aQ(z,t)\), where \(a\le j\) is the least index of a
Taylor coefficient \(c_a(z)\) of \(F(z,\cdot)\) that is not identically zero (some \(c_i\), \(i\le j\), is not, because \(F(0,\cdot)=g(y_0+\cdot\,u)\) has
order \(j\)); \(Q(z,0)=c_a(z)\). Let \(y\) be a zero of order \(j\) near \(y_0\); it lies on the graph, over \(w=w_0+z\). If \(c_a(z)\neq0\): for
\(a=0\) this contradicts \(F(z,0)=g(y)=0\); for \(a\ge1\), \(g=h^av\) near \(y\) with \(h=\langle u,\cdot\rangle-k\circ\pi\), \(v=Q(\pi-w_0,h)\neq0\) and
\(dh(u)=1\), so \(y\) is a transverse regular point. Hence the non-transverse zeros of order \(j\) project into \(w_0+\{c_a=0\}\), which is null.
\end{proof}

\begin{lemma}[Coarea formula for the transverse part of a divisor]
\label{ana:divisor-coarea}
\uses{ana:hausdorff-area-formula,ana:divisor-local-graph}
\lean{TranscendentalNakamaye.Pending.DivisorCoarea.lintegral_transverse_eq,TranscendentalNakamaye.Pending.DivisorCoarea.GraphPiece,TranscendentalNakamaye.Pending.DivisorCoarea.exists_graphPiece,TranscendentalNakamaye.Pending.DivisorCoarea.GraphPiece.graph_contDiffOn,TranscendentalNakamaye.Pending.DivisorCoarea.GraphPiece.graph_injOn,TranscendentalNakamaye.Pending.DivisorCoarea.GraphPiece.graph_fderiv_injective,TranscendentalNakamaye.Pending.DivisorCoarea.indicator_iUnion_eq_tsum,TranscendentalNakamaye.Pending.DivisorRegularLocus.RegularPresentation.restrictOpen}
\leanok
Let \(g\) be continuous on an open set \(\Omega\subset\mathbf C^n\), \(u\) a unit vector, \(\mathbf C^n=\mathbf Cu\oplus\iota(\mathbf C^m)\) an orthogonal
splitting, and \(Z_{\mathrm{tr}}\) the transverse locus of the divisor of \(g\) in the direction \(u\) (Definition~\ref{ana:divisor-regular-points}).
For every nonnegative Borel function \(F\),
\[
 \int_{Z_{\mathrm{tr}}\cap\Omega}\tau_x(u,iu)\,m_g(x)\,F(x)\,d\mathcal H^{2m}(x)
 =\int_{\mathbf C^m}\ \sum_{s\in\mathbf C,\ su+\iota w\in Z_{\mathrm{tr}}\cap\Omega}m_g(su+\iota w)\,F(su+\iota w)\;dw .
\]
\end{lemma}
\begin{proof}
\leanok
Every point of \(Z_{\mathrm{tr}}\cap\Omega\) has a neighbourhood \(N\subset\Omega\) as in Lemma~\ref{ana:divisor-local-graph}: the zeros in \(N\) form the
graph of \(G(w)=k(w)u+\iota w\), \(w\in W\), the multiplicity is constant on them, and \(J(DG)\,\tau(u,iu)=1\). Cover \(Z_{\mathrm{tr}}\cap\Omega\) by
countably many such \(N_i\) and make the traces \(Q_i\) on \(Z_{\mathrm{tr}}\cap\Omega\) disjoint. With \(P_i=G_i^{-1}(Q_i)\cap W_i\) we have
\(Q_i=G_i(P_i)\), and Lemma~\ref{ana:hausdorff-area-formula} gives
\(\int_{Q_i}\tau\,m\,F\,d\mathcal H^{2m}=\int_{P_i}J(DG_i)\,\tau\,m_i\,F\circ G_i=\int_{P_i}m_i\,F\circ G_i\). On the other side, a line
\(\{su+\iota w\}\) meets \(Q_i\) only over \(w\in P_i\), at \(s=k_i(w)\). Summing over \(i\) (monotone convergence) gives the formula.
\end{proof}

\begin{theorem}[Lelong--Poincar\'e equation on open sets of \(\mathbf C^n\)]
\label{ana:divisor-integration-current}
\uses{ana:divisor-regular-points,ana:divisor-singular-projection,ana:divisor-coarea,ana:lelong-poincare-slices,as:logarithmic-psh,psh:local-integrability,lc:closed-type,lc:positive-bridge,dist:smooth-domination}
\lean{TranscendentalNakamaye.Pending.DivisorIntegrationCurrent.ddC_log_normSq_eq_divisorCurrent,TranscendentalNakamaye.Pending.DivisorIntegrationCurrent.ddC_log_normSq_apply,TranscendentalNakamaye.Pending.DivisorIntegrationCurrent.divisorCurrent,TranscendentalNakamaye.Pending.DivisorIntegrationCurrent.divisorCurrent_apply,TranscendentalNakamaye.Pending.DivisorIntegrationCurrent.divisorDensity,TranscendentalNakamaye.Pending.DivisorIntegrationCurrent.isOneOne_divisorCurrent,TranscendentalNakamaye.Pending.DivisorIntegrationCurrent.locallyIntegrableOn_log_normSq,TranscendentalNakamaye.Pending.DivisorIntegrationCurrent.locallyIntegrableOn_divisorDensity,TranscendentalNakamaye.Pending.DivisorIntegrationCurrent.lintegral_indicator_divisorMult_lt_top,TranscendentalNakamaye.Pending.DivisorIntegrationCurrent.lintegral_transverse_eq_ofReal_ddC,TranscendentalNakamaye.Pending.DivisorIntegrationCurrent.ddC_apply_I_eq_integral_sliceSum,TranscendentalNakamaye.Pending.DivisorIntegrationCurrent.ofReal_sliceSum,TranscendentalNakamaye.Pending.DivisorIntegrationCurrent.measurePreserving_splitMap,TranscendentalNakamaye.Pending.DivisorIntegrationCurrent.eq_of_isOneOne,TranscendentalNakamaye.Pending.DivisorIntegrationCurrent.eq_of_forall_nonneg,TranscendentalNakamaye.Pending.DivisorIntegrationCurrent.sum_lineForm_basis}
\leanok
Let \(g\) be holomorphic on an open set \(\Omega\subset\mathbf C^n\), vanishing identically near no point of \(\Omega\), and let \(Z_{\mathrm{reg}}\) be the
regular part of its zero set, with multiplicity \(m_g\) and transverse form \(\tau\) (Definition~\ref{ana:divisor-regular-points}). Then
\(\log|g|^2\in L^1_{\mathrm{loc}}(\Omega)\); \(m_g\,\mathcal H^{2n-2}\) restricted to \(Z_{\mathrm{reg}}\) is locally finite in \(\Omega\); and
\[
 d d^c\log|g|^2=[\operatorname{div}g],\qquad
 [\operatorname{div}g](v,w)(\varphi)=\int_{Z_{\mathrm{reg}}\cap\Omega}\varphi\,m_g\,\tau(v,w)\,d\mathcal H^{2n-2},
\]
as distribution-valued real \(2\)-forms on \(\Omega\): for all \(v,w\in\mathbf C^n\) and \(\varphi\in C^\infty_c(\Omega,\mathbf R)\), with
\(d d^c=(2\pi)^{-1}\mathrm i\partial\bar\partial\) and \(\mathcal H^{2n-2}\) the Euclidean Hausdorff measure.
\end{theorem}
\begin{proof}
\leanok
\(\log|g|^2\) is plurisubharmonic (Lemma~\ref{as:logarithmic-psh}) and finite somewhere on every ball, hence locally integrable
(Theorem~\ref{psh:local-integrability}). Both sides are of type \((1,1)\) (Proposition~\ref{lc:closed-type}; \(\tau(iv,iw)=\tau(v,w)\)), so by
polarization (\(A(v,iw)=\tfrac12(Q(v+w)-Q(v)-Q(w))\) with \(Q(z)=A(z,iz)\), using Proposition~\ref{lc:positive-bridge}) it suffices to compare
them on \((u,iu)\) for unit vectors \(u\), and, writing \(\varphi=(\varphi+\chi)-\chi\) with \(\chi\ge|\varphi|\) smooth
(Lemma~\ref{dist:smooth-domination}), on tests \(\varphi\ge0\).

Fix \(u\) and an orthogonal splitting \(\mathbf C^n=\mathbf Cu\oplus\iota(\mathbf C^{n-1})\); \((s,w)\mapsto su+\iota w\) preserves Lebesgue measure.
By Lemma~\ref{ana:lelong-poincare-slices}, \(d d^c\log|g|^2(u,iu)(\varphi)=\int_w\sum_s\operatorname{ord}_s g(\cdot\,u+\iota w)\,\varphi(su+\iota w)\,dw\),
the inner sums being integrable. For almost every \(w\) every zero on the line is a transverse regular point
(Lemma~\ref{ana:divisor-singular-projection}), where the order along the line is \(m_g\); so the inner sum is
\(\sum_{s:\ su+\iota w\in Z_{\mathrm{tr}}}m_g\varphi\ge0\). By Lemma~\ref{ana:divisor-coarea} its integral is
\(\int_{Z_{\mathrm{tr}}\cap\Omega}\tau(u,iu)\,m_g\,\varphi\,d\mathcal H^{2n-2}\), and \(\tau(u,iu)=0\) on \(Z_{\mathrm{reg}}\setminus Z_{\mathrm{tr}}\).

Local finiteness: for an orthonormal basis \((e_j)\), \(\sum_j\tau(e_j,ie_j)=\sum_j|\ell(e_j)|^2/\|\ell\|^2=1\) on \(Z_{\mathrm{reg}}\), so for a compact
\(K\subset\Omega\) and \(\chi\ge\mathbf 1_K\) smooth, \(\int_{K\cap Z_{\mathrm{reg}}}m_g\,d\mathcal H^{2n-2}\le\sum_j d d^c\log|g|^2(e_j,ie_j)(\chi)<\infty\)
by the identity just proved for nonnegative tests (which does not use the finiteness). Measurability of the density holds piecewise on a
countable cover of \(Z_{\mathrm{reg}}\) by domains of local presentations.
\end{proof}

\begin{theorem}[Lelong--Poincar\'e for a metered effective \(\mathbf Q\)-divisor]
\label{ana:metered-divisor-lelong-poincare}
\uses{ana:divisor-integration-current,gcp:chart-ddc,gcp:global-ddc,gcp:properties,gc:smooth-global,dist:restriction,lc:closed-type,lc:positive-bridge}
\lean{TranscendentalNakamaye.Pending.LelongPoincareDivisor.lelongPoincare,TranscendentalNakamaye.Pending.LelongPoincareDivisor.chartLocallyIntegrable_weight,TranscendentalNakamaye.Pending.LelongPoincareDivisor.integrationCurrent,TranscendentalNakamaye.Pending.LelongPoincareDivisor.ddCWeight,TranscendentalNakamaye.Pending.LelongPoincareDivisor.ddCWeight_eq,TranscendentalNakamaye.Pending.LelongPoincareDivisor.isClosed_integrationCurrent,TranscendentalNakamaye.Pending.LelongPoincareDivisor.isTypeOneOne_integrationCurrent,TranscendentalNakamaye.Pending.LelongPoincareDivisor.isSquarePositive_integrationCurrent,TranscendentalNakamaye.Pending.LelongPoincareDivisor.constantDirectionPositive_integrationCurrent,TranscendentalNakamaye.Pending.LelongPoincareDivisor.restrictForm_coefficients_eq_sum,TranscendentalNakamaye.Pending.LelongPoincareDivisor.restrictForm_coefficients_eq_ddC,TranscendentalNakamaye.Pending.LelongPoincareDivisor.LocalPresentation,TranscendentalNakamaye.Pending.LelongPoincareDivisor.exists_localPresentation,TranscendentalNakamaye.Pending.LelongPoincareDivisor.haarFactor,TranscendentalNakamaye.Pending.LelongPoincareDivisor.restrictForm,TranscendentalNakamaye.Pending.LelongPoincareDivisor.restrictForm_ddC,TranscendentalNakamaye.Pending.LelongPoincareDivisor.restrictForm_exteriorDerivative,TranscendentalNakamaye.Pending.LelongPoincareDivisor.eq_of_forall_restrict_eq,TranscendentalNakamaye.Pending.LelongPoincareDivisor.positiveDistribution_of_forall_restrict,TranscendentalNakamaye.Pending.LelongPoincareDivisor.form_eq_of_forall_restrictForm_eq,TranscendentalNakamaye.Pending.LelongPoincareDivisor.ddC_regularDistribution_apply_of_contDiffOn,TranscendentalNakamaye.Pending.LelongPoincareDivisor.ddC_eq_sum_divisorCurrent_add,TranscendentalNakamaye.Pending.LelongPoincareDivisor.LocalPresentation.ae_chartFunction_eq}
\leanok
Let \(Y\) be a complex manifold of dimension \(n\ge1\), \(\psi_E\colon Y\to[-\infty,\infty)\) a function and \(\Theta_E\) a smooth closed real
\((1,1)\)-form such that every point of \(Y\) has an open neighbourhood \(U\) on which \(\psi_E=\sum_ia_i\log|g_i|^2+h\) and \(\Theta_E=-dd^ch\),
with finitely many \(a_i\in\mathbf Q_{\ge0}\), \(g_i\) holomorphic on \(U\) and vanishing identically on no open subset, and \(h\) smooth. Then
\(\psi_E\in L^1_{\mathrm{loc}}(Y)\), and the current \([E]:=dd^c\psi_E+\Theta_E\) is closed, of type \((1,1)\) and positive, and on every such
\(U\) it equals \(\sum_ia_i[\operatorname{div}g_i]\), the integration currents of Theorem~\ref{ana:divisor-integration-current} read in charts. In
particular \(dd^c\psi_E=[E]-\Theta_E\).
\end{theorem}
\begin{proof}
\leanok
Fix a chart and a presentation on \(U\); write \(G_i\), \(H\) for \(g_i\), \(h\) in the chart. The zero set of each \(G_i\) is null, so
\(\psi_E=\sum_ia_i\log|G_i|^2+H\) almost everywhere on the chart image \(W\) of \(U\), which is locally integrable
(Theorem~\ref{ana:divisor-integration-current}); this gives \(\psi_E\in L^1_{\mathrm{loc}}\). On \(W\), Theorem~\ref{ana:divisor-integration-current}
gives \(dd^c\sum_ia_i\log|G_i|^2=\sum_ia_i[\operatorname{div}G_i]\), and \(dd^cH=-\Theta_E\): for a test function \(\varphi\) on \(W\), replace \(H\)
by \(\chi H\) with \(\chi\) smooth, equal to \(1\) near \(\operatorname{supp}\varphi\) and supported in \(W\), which is smooth on the whole chart.
Restriction to open subsets commutes with \(dd^c\) and with the exterior derivative, and a distribution is determined by, and is positive if,
its restrictions to an open cover are (smooth partition of unity). On \(W\), \([E]=dd^c(\psi_E-h)\) is closed and of type \((1,1)\)
(Proposition~\ref{lc:closed-type}), and \([E]=\sum_ia_i[\operatorname{div}G_i]\) is positive on constant directions, the densities
\(m\,\tau(v,iv)\) being nonnegative; hence \([E]\) is closed, of type \((1,1)\) and positive on constant directions everywhere, and positive on
vector squares by Proposition~\ref{lc:positive-bridge}.

In the formalization the coefficients of currents in a chart are densities against the Haar measure normalized by a chosen basis \(b\); the
local identity then carries the constant factor \(c_b\), the density of that measure with respect to Lebesgue measure (\(c_b=1\) for an
orthonormal basis).
\end{proof}

\begin{definition}[First-pass restricted-volume data]
\label{ana:restricted-volume-scaffold}
\uses{fpgeom:complex-analytic-subvariety,as:analytic-potentials,as:smooth-off-pole-type,bc:quotient,posclass:quotient-predicates}
\lean{TranscendentalNakamaye.FirstPass.NumericalRestrictedVolumeScaffold.poleSet,TranscendentalNakamaye.FirstPass.NumericalRestrictedVolumeScaffold.mem_poleSet,TranscendentalNakamaye.FirstPass.NumericalRestrictedVolumeScaffold.HasManifoldSmoothAnalyticSingularities,TranscendentalNakamaye.FirstPass.NumericalRestrictedVolumeScaffold.PositiveError,TranscendentalNakamaye.FirstPass.NumericalRestrictedVolumeScaffold.AdmissibleRestrictedCurrentPresentation,TranscendentalNakamaye.FirstPass.NumericalRestrictedVolumeScaffold.NonNefLocusData,TranscendentalNakamaye.FirstPass.NumericalRestrictedVolumeScaffold.RestrictedACTopMassData,TranscendentalNakamaye.FirstPass.NumericalRestrictedVolumeScaffold.restrictedMassEnvelope,TranscendentalNakamaye.FirstPass.NumericalRestrictedVolumeScaffold.numericalRestrictedVolumeInf,TranscendentalNakamaye.FirstPass.NumericalRestrictedVolumeScaffold.numericalRestrictedVolumeInf_eq_zero_of_subset_nonNef,TranscendentalNakamaye.FirstPass.NumericalRestrictedVolumeScaffold.numericalRestrictedVolumeInf_eq_iInf_of_not_subset_nonNef,TranscendentalNakamaye.FirstPass.NumericalRestrictedVolumeScaffold.nullLocusInf,TranscendentalNakamaye.FirstPass.NumericalRestrictedVolumeScaffold.mem_nullLocusInf,TranscendentalNakamaye.FirstPass.NumericalRestrictedVolumeScaffold.carrier_subset_nullLocusInf_of_volume_eq_zero,TranscendentalNakamaye.FirstPass.NumericalRestrictedVolumeScaffold.carrier_subset_nullLocusInf_of_subset_nonNef}
\leanok
Fix a K\"ahler background, a Bott--Chern class $\alpha$, a closed
irreducible analytic subvariety $V$, and $\epsilon>0$.  An admissible
potential presentation consists of a closed smooth representative $\theta$ of
$\alpha$ and a potential $u$ with analytic singularities such that
$\theta+\epsilon\omega+d d^c u\geq0$ and $V$ is not contained in the pole set
of $u$.

Here ``analytic singularities'' is the convention of
Definition~\ref{as:smooth-off-pole-type}: in every chart
$u=c\log\sum_j|f_j|^2+b$ with $c>0$ real, the $f_j$ holomorphic, and $b$ bounded
near the point and $C^\infty$ at every nearby point where the $f_j$ do not all
vanish. Nothing is asked of $b$ \emph{at} a pole beyond boundedness. This is the
class in which the supremum below is taken, so it is also what every upper bound on
the mass envelope has to hold for; it is larger than the class with $b$ smooth
across the poles, and it is the one Bergman-kernel regularization produces.

Until the geometric constructions are available, let
$N(\alpha)\subset X$ and the absolutely-continuous top mass of every
admissible presentation be explicitly supplied data.  Their candidate mass
envelope is the supremum over admissible presentations.  One may form the
total auxiliary quantity
\[
 \operatorname{RV}^{\inf}_{X\mid V}(\alpha)=
 \begin{cases}
 0,&V\subset N(\alpha),\\
 \inf_{\epsilon>0}\sup_T m_{V,\epsilon}(T),&\text{otherwise},
 \end{cases}
\]
and the union of positive-dimensional $V$ where this value is zero.  The
supplied locus and mass remain data, and this infimum is not yet identified
with the manuscript's right limit.
\end{definition}

\begin{definition}[Presented non-K\"ahler and non-nef loci]
\label{ana:presented-nonnef-locus}
\uses{ana:restricted-volume-scaffold,ana:geometric-setting}
\lean{TranscendentalNakamaye.FirstPass.PresentedNonNefLocus.PresentedKahlerCurrent,TranscendentalNakamaye.FirstPass.PresentedNonNefLocus.presentedNonKahlerLocus,TranscendentalNakamaye.FirstPass.PresentedNonNefLocus.mem_presentedNonKahlerLocus,TranscendentalNakamaye.FirstPass.PresentedNonNefLocus.presentedNonKahlerLocus_eq_univ_of_not_nonempty,TranscendentalNakamaye.FirstPass.PresentedNonNefLocus.presentedNonNefLocus,TranscendentalNakamaye.FirstPass.PresentedNonNefLocus.mem_presentedNonNefLocus,TranscendentalNakamaye.FirstPass.PresentedNonNefLocus.presentedNonNefLocusData,TranscendentalNakamaye.FirstPass.PresentedNonNefLocus.presentedNonNefLocusData_locus}
\leanok
Relative to a fixed K\"ahler background, a presented K\"ahler current in a
Bott--Chern class consists of a closed smooth representative, a global
potential with smooth analytic singularities, and a positive lower bound
expressed by the existing theta-psh predicate.  Its presented non-K\"ahler
locus is the common pole set of all such presentations, and its presented
non-nef locus is the union of these loci over all positive K\"ahler
perturbations.  This gives an actual class-indexed locus for the first-pass
restricted-volume interface.

The word ``presented'' is essential: identifying this set with the
manuscript's $E_{\mathrm{nK}}$ and $E_{\mathrm{nn}}$ still requires the
current--potential equivalence, Demailly regularization, and the
pole-set--positive-Lelong theorem.
\end{definition}

\begin{lemma}[Presented currents give admissible competitors]
\label{ana:presented-admissible-bridge}
\uses{ana:presented-nonnef-locus,ana:restricted-volume-scaffold}
\lean{TranscendentalNakamaye.FirstPass.PresentedAdmissibleBridge.admissibleRestrictedCurrentPresentationOfPresented,TranscendentalNakamaye.FirstPass.PresentedAdmissibleBridge.nonempty_admissible_of_not_subset_presentedNonKahlerLocus}
\leanok
Let $\epsilon>0$.  A presented K\"ahler current in
$\alpha+\epsilon\{\omega\}$ whose pole set does not contain $V$ gives an
admissible restricted-current presentation for $\alpha$ at error
$\epsilon$: subtract $\epsilon\omega$ from its smooth representative, and
forget the remaining strict positive margin.  Consequently, if $V$ is not
contained in the presented non-K\"ahler locus of the perturbed class, the
admissible family is nonempty.
\end{lemma}

\begin{definition}[Lelong numbers and the non-K\"ahler locus]
\label{ana:nonkahler-locus}
\uses{ana:geometric-setting,ana:restricted-volume-scaffold,fpgeom:complex-analytic-subvariety,bc:quotient}
\source{transcendental-nakamaye-copy4-20260911:sec:nK}
\lean{TranscendentalNakamaye.FirstPass.NonKahlerLocus.lelongNumber,TranscendentalNakamaye.FirstPass.NonKahlerLocus.KahlerCurrent,TranscendentalNakamaye.FirstPass.NonKahlerLocus.IsBigClass,TranscendentalNakamaye.FirstPass.NonKahlerLocus.nonKahlerLocus,TranscendentalNakamaye.FirstPass.NonKahlerLocus.mem_nonKahlerLocus,TranscendentalNakamaye.FirstPass.NonKahlerLocus.nonKahlerLocus_eq_univ_of_not_isBigClass,TranscendentalNakamaye.FirstPass.NonKahlerLocus.nonNefLocus,TranscendentalNakamaye.FirstPass.NonKahlerLocus.mem_nonNefLocus,TranscendentalNakamaye.FirstPass.NonKahlerLocus.nonNefLocusData,TranscendentalNakamaye.FirstPass.NonKahlerLocus.IsSmoothNear,TranscendentalNakamaye.FirstPass.NonKahlerLocus.IsIrreducibleComponent,TranscendentalNakamaye.FirstPass.NonKahlerLocus.singularLocus}
\leanok
Let \(M\) be a complex manifold. The Lelong number of a function
\(u\colon M\to[-\infty,\infty)\) at \(x\) is computed in a holomorphic chart
centred at \(x\) as
\[
 \nu(u,x)=\sup\{\gamma\ge0:\ u(y)\le\gamma\log|y-x|+O(1)
 \text{ for }y\ne x\text{ near }x\}.
\]
For quasi-plurisubharmonic \(u\) this is the classical Lelong number, and for
a smooth closed real \((1,1)\)-form \(\theta\) it is the Lelong number of the
current \(\theta+dd^cu\). A K\"ahler current in a real Bott--Chern class
\(\alpha\) is written \(T=\theta+dd^cu\), with \(\theta\in\alpha\) smooth and
closed, such that \(T\ge\omega'\) for a K\"ahler form \(\omega'\); the class is
big if such a current exists. The non-K\"ahler locus is
\[
 E_{\mathrm{nK}}(\alpha)=\{x:\ \nu(T,x)>0\text{ for every K\"ahler current }T\in\alpha\},
\]
which is \(M\) when \(\alpha\) is not big, and for a K\"ahler form \(\omega\)
the non-nef locus is
\(E_{\mathrm{nn}}(\alpha)=\bigcup_{\epsilon>0}E_{\mathrm{nK}}(\alpha+\epsilon\{\omega\})\).
A potential is smooth near \(p\) if it agrees near \(p\) with a smooth real
function. A closed irreducible analytic subvariety \(V\) is an irreducible
component of a set \(S\) if \(V\subset S\) and every closed irreducible analytic
subvariety \(Z\) with \(V\subset Z\subset S\) equals \(V\). The singular locus
\(V_{\mathrm{sing}}\) is the complement in \(V\) of its regular points of
dimension \(\dim V\).
\end{definition}

\begin{definition}[Ideal sheaves, divisors and multiplier ideals]
\label{ana:ideal-sheaves}
\uses{ana:geometric-setting,ana:nonkahler-locus,sing:multiplier-local,bc:quotient,ana:metered-divisor-lelong-poincare}
\source{transcendental-nakamaye-copy4-20260911:sec:prelim,transcendental-nakamaye-copy4-20260911:lem:relative-principalization,transcendental-nakamaye-copy4-20260911:lem:exceptional-metric}
\lean{TranscendentalNakamaye.FirstPass.HolomorphicIdealSheaves.HolomorphicAt,TranscendentalNakamaye.FirstPass.HolomorphicIdealSheaves.IsInSpanAt,TranscendentalNakamaye.FirstPass.HolomorphicIdealSheaves.IsInSpanOnAt,TranscendentalNakamaye.FirstPass.HolomorphicIdealSheaves.GermIdeal,TranscendentalNakamaye.FirstPass.HolomorphicIdealSheaves.GermIdeal.Le,TranscendentalNakamaye.FirstPass.HolomorphicIdealSheaves.GermIdeal.IsCoherent,TranscendentalNakamaye.FirstPass.HolomorphicIdealSheaves.GermIdeal.IsInvertible,TranscendentalNakamaye.FirstPass.HolomorphicIdealSheaves.GermIdeal.mul,TranscendentalNakamaye.FirstPass.HolomorphicIdealSheaves.GermIdeal.pullback,TranscendentalNakamaye.FirstPass.HolomorphicIdealSheaves.GermIdeal.zeroLocus,TranscendentalNakamaye.FirstPass.HolomorphicIdealSheaves.vanishingIdeal,TranscendentalNakamaye.FirstPass.HolomorphicIdealSheaves.multiplierIdeal,TranscendentalNakamaye.FirstPass.HolomorphicIdealSheaves.EffectiveDivisor,TranscendentalNakamaye.FirstPass.HolomorphicIdealSheaves.EffectiveDivisor.support,TranscendentalNakamaye.FirstPass.HolomorphicIdealSheaves.MeteredEffectiveQDivisor,TranscendentalNakamaye.FirstPass.HolomorphicIdealSheaves.MeteredEffectiveQDivisor.support,TranscendentalNakamaye.FirstPass.GlobalAnalyticSingularities.HasAnalyticSingularitiesOfType,TranscendentalNakamaye.FirstPass.GlobalAnalyticSingularities.HasAnalyticSingularities}
\leanok
An ideal sheaf \(\mathfrak a\subset\mathcal O_M\) is described by its stalks of
holomorphic germs. It is coherent if every point has a neighbourhood on which
finitely many holomorphic functions generate every stalk; by Oka's theorem this
is coherence of an ideal subsheaf of \(\mathcal O_M\). Products
\(\mathfrak a\mathfrak b\), pullback ideals \(\mathfrak a\mathcal O_Y\) along
holomorphic maps \(f\colon Y\to M\) (generated by the germs \(a\circ f\)), the
ideal \(\mathcal I_A\) of germs vanishing on a set \(A\), and zero loci are the
usual ones. For a closed submanifold \(W\), an ideal of \(\mathcal O_W\) is
determined by the ambient germs whose restrictions lie in it, since every
holomorphic germ on \(W\) extends to the ambient manifold. An effective divisor
\(D\) is given by its invertible ideal \(\mathcal O_M(-D)\), locally generated
by one holomorphic function vanishing identically on no open set; its support
is its zero locus. The multiplier ideal \(\mathcal I(c\varphi)\) consists of
germs \(g\) with \(|g|^2e^{-c\varphi}\) locally integrable. An effective
\(\mathbf Q\)-divisor \(E\) with a smooth Hermitian metric on
\(\mathcal O_M(E)\) is given by \(\psi_E=\log|s_E|^2\) and the curvature form
\(\Theta_E\): locally \(\psi_E=\sum_ia_i\log|g_i|^2+h\) with \(a_i\in\mathbf
Q_{\ge0}\), local equations \(g_i\) vanishing identically on no open set,
\(h\) smooth, and \(\Theta_E=-dd^ch\), so that \(dd^c\psi_E=[E]-\Theta_E\);
the support of \(E\) is the polar set of \(\psi_E\). A function \(\varphi\) has
analytic singularities of type \((c,\mathfrak a)\), where
\(c\in\mathbf Q_{>0}\) and \(\mathfrak a\) is coherent, if near every point
\(\varphi=c\log\sum_i|g_i|^2+R\) for local generators \(g_i\) of
\(\mathfrak a\) and a bounded function \(R\).
\end{definition}

\begin{definition}[Blowups, modifications and strict transforms]
\label{ana:blowups}
\uses{ana:ideal-sheaves,fpgeom:complex-analytic-subvariety}
\source{transcendental-nakamaye-copy4-20260911:sec:prelim,transcendental-nakamaye-copy4-20260911:lem:relative-principalization,transcendental-nakamaye-copy4-20260911:lem:resolution}
\lean{TranscendentalNakamaye.FirstPass.BlowupSequences.IsClosedSubmanifold,TranscendentalNakamaye.FirstPass.BlowupSequences.IsClosedConnectedSubmanifold,TranscendentalNakamaye.FirstPass.BlowupSequences.exceptionalLocus,TranscendentalNakamaye.FirstPass.BlowupSequences.strictTransform,TranscendentalNakamaye.FirstPass.BlowupSequences.IsIsomorphismOver,TranscendentalNakamaye.FirstPass.BlowupSequences.IsBlowupAlong,TranscendentalNakamaye.FirstPass.BlowupSequences.IsBlowupAlong.contMDiff,TranscendentalNakamaye.FirstPass.BlowupSequences.IsBlowupSequence,TranscendentalNakamaye.FirstPass.BlowupSequences.IsAdaptedBlowupSequence,TranscendentalNakamaye.FirstPass.BlowupSequences.IsBlowupSequence.contMDiff,TranscendentalNakamaye.FirstPass.BlowupSequences.IsAdaptedBlowupSequence.isBlowupSequence,TranscendentalNakamaye.FirstPass.BlowupSequences.IsAdaptedBlowupSequence.contMDiff}
\leanok
A closed complex submanifold is a closed set which near each of its points is
biholomorphically a coordinate plane. For a holomorphic map
\(f\colon Y\to X\) between manifolds of the same dimension, the exceptional
locus \(\operatorname{Exc}(f)\) is the set where \(df\) is not bijective; the
strict transform of \(A\subset X\) is
\(\overline{f^{-1}(A)\setminus\operatorname{Exc}(f)}\); and \(f\) is an
isomorphism over \(S\subset X\) if it maps \(f^{-1}(S)\) bijectively onto
\(S\) with invertible differential at every point of \(f^{-1}(S)\). The blowup
of \(X\) along a nowhere dense closed submanifold \(C\) is a holomorphic map
\(\pi\colon Y\to X\) such that \(\mathcal I_C\mathcal O_Y\) is invertible and
every holomorphic map \(g\colon Z\to X\) from a complex manifold with
\(\mathcal I_C\mathcal O_Z\) invertible factors uniquely and holomorphically
through \(\pi\). A modification is a finite composition of blowups with smooth
centers. Such a composition is adapted to \(W\subset X\) if each center is
contained in a proper analytic subset of the current strict transform of \(W\)
or is disjoint from it.
\end{definition}

\begin{definition}[Gluing conditions]
\label{ana:gluing-conditions}
\uses{ana:nonkahler-locus}
\source{transcendental-nakamaye-copy4-20260911:prop:gluing,transcendental-nakamaye-copy4-20260911:prop:principalization}
\lean{TranscendentalNakamaye.FirstPass.AnalyticGluingData.ThetaPlurisubharmonicOn,TranscendentalNakamaye.FirstPass.AnalyticGluingData.IsThetaPshPotentialOn,TranscendentalNakamaye.FirstPass.AnalyticGluingData.IsKahlerCurrentPotential,TranscendentalNakamaye.FirstPass.AnalyticGluingData.DominatesKahlerFormOn,TranscendentalNakamaye.FirstPass.AnalyticGluingData.RelCompactSubset,TranscendentalNakamaye.FirstPass.AnalyticGluingData.GluingHypotheses,TranscendentalNakamaye.FirstPass.AnalyticGluingData.PrincipalizationGluingConclusions}
\leanok
Let \(Y\) be a complex manifold and \(\theta\) a smooth closed real
\((1,1)\)-form. A function \(u\) is \(\theta\)-psh on a set \(\Omega\) if near
every point of \(\Omega\) it is plurisubharmonic after adding a smooth local
potential of \(\theta\), and is not identically \(-\infty\) near any point of
\(\Omega\). The current \(\theta+dd^cu\) dominates a K\"ahler form on
\(\Omega\) if \(u\) is \((\theta-\omega')\)-psh on \(\Omega\) for a K\"ahler
form \(\omega'\) on \(Y\); on a compact K\"ahler manifold, after shrinking a
neighbourhood of a compact set, this is equivalent to dominating a K\"ahler form
defined only on the neighbourhood. For open sets \(U_1\Subset U_2\), a point
\(p\), and functions \(\varphi_G,\varphi_U\), the gluing hypotheses are
hypotheses (1)--(3) of Proposition~\ref{ana:gluing}. For a closed set \(W\), a
point \(y\), and an open set \(U\supset W\), the principalization conclusions
are conclusions (1)--(3) of Proposition~\ref{ana:principalization-data}.
\end{definition}

\begin{proposition}[Lebesgue decomposition of a signed coefficient measure]
\label{ana:coefficient-lebesgue-decomposition}
\lean{TranscendentalNakamaye.FirstPass.CurrentCoefficientLebesgue.density,TranscendentalNakamaye.FirstPass.CurrentCoefficientLebesgue.singularPart,TranscendentalNakamaye.FirstPass.CurrentCoefficientLebesgue.absolutelyContinuousPart,TranscendentalNakamaye.FirstPass.CurrentCoefficientLebesgue.measurable_density,TranscendentalNakamaye.FirstPass.CurrentCoefficientLebesgue.integrable_density,TranscendentalNakamaye.FirstPass.CurrentCoefficientLebesgue.singularPart_mutuallySingular,TranscendentalNakamaye.FirstPass.CurrentCoefficientLebesgue.absolutelyContinuousPart_absolutelyContinuous,TranscendentalNakamaye.FirstPass.CurrentCoefficientLebesgue.singularPart_add_absolutelyContinuousPart,TranscendentalNakamaye.FirstPass.CurrentCoefficientLebesgue.absolutelyContinuousPart_eq_self,TranscendentalNakamaye.FirstPass.CurrentCoefficientLebesgue.absolutelyContinuous_iff_absolutelyContinuousPart_eq_self,TranscendentalNakamaye.FirstPass.CurrentCoefficientLebesgue.singularPart_eq_zero_of_absolutelyContinuous,TranscendentalNakamaye.FirstPass.CurrentCoefficientLebesgue.decomposition_unique,TranscendentalNakamaye.FirstPass.CurrentCoefficientLebesgue.withDensity_indicator_eq_restrict,TranscendentalNakamaye.FirstPass.CurrentCoefficientLebesgue.mutuallySingular_restrict_left,TranscendentalNakamaye.FirstPass.CurrentCoefficientLebesgue.restrict_decomposition,TranscendentalNakamaye.FirstPass.CurrentCoefficientLebesgue.singularPart_restrict,TranscendentalNakamaye.FirstPass.CurrentCoefficientLebesgue.density_restrict,TranscendentalNakamaye.FirstPass.CurrentCoefficientLebesgue.pushforwardDensity,TranscendentalNakamaye.FirstPass.CurrentCoefficientLebesgue.measurable_pushforwardDensity,TranscendentalNakamaye.FirstPass.CurrentCoefficientLebesgue.pushforwardDensity_comp,TranscendentalNakamaye.FirstPass.CurrentCoefficientLebesgue.integrable_pushforwardDensity,TranscendentalNakamaye.FirstPass.CurrentCoefficientLebesgue.withDensity_pushforwardDensity,TranscendentalNakamaye.FirstPass.CurrentCoefficientLebesgue.toENNRealVectorMeasure_map,TranscendentalNakamaye.FirstPass.CurrentCoefficientLebesgue.mutuallySingular_map,TranscendentalNakamaye.FirstPass.CurrentCoefficientLebesgue.map_decomposition,TranscendentalNakamaye.FirstPass.CurrentCoefficientLebesgue.singularPart_map,TranscendentalNakamaye.FirstPass.CurrentCoefficientLebesgue.density_map_comp}
\leanok
For a signed measure coefficient $s$ and a positive reference measure $\mu$,
the Radon--Nikodym derivative and singular part give the canonical Lebesgue
decomposition
\[
 s=s_{\mathrm{sing}}+\mu\cdot f,
 \qquad s_{\mathrm{sing}}\perp\mu,
 \qquad f\in L^1(\mu).
\]
The two pieces are unique.  Restriction to measurable subsets multiplies the
density by the corresponding indicator and restricts the singular part.
Pushforward along a measurable embedding transports both pieces and pulls the
density back to the original density almost everywhere.

This proposition starts with a signed measure.  It does not yet prove that a
distributional coefficient of a restricted current is an order-zero signed
measure, nor does it construct the restriction of a current to
$V_{\mathrm{reg}}$.
\end{proposition}

\begin{definition}[Numerical restricted volume and null locus]
\label{ana:restricted-volume}
\uses{ana:nonkahler-locus,ana:geometric-setting,fpgeom:complex-analytic-subvariety,ana:restricted-volume-scaffold,ana:presented-nonnef-locus,ana:presented-admissible-bridge,ana:coefficient-lebesgue-decomposition}
\source{transcendental-nakamaye-copy4-20260911:eq:def-restricted-volume,transcendental-nakamaye-copy4-20260911:eq:null-locus}
\lean{TranscendentalNakamaye.FirstPass.NumericalRestrictedVolume.planeInclusion,TranscendentalNakamaye.FirstPass.NumericalRestrictedVolume.RegularPatch,TranscendentalNakamaye.FirstPass.NumericalRestrictedVolume.RegularPatch.parameterDomain,TranscendentalNakamaye.FirstPass.NumericalRestrictedVolume.RegularPatch.coordinateParametrization,TranscendentalNakamaye.FirstPass.NumericalRestrictedVolume.RegularPatch.toManifold,TranscendentalNakamaye.FirstPass.NumericalRestrictedVolume.perturbedRepresentative,TranscendentalNakamaye.FirstPass.NumericalRestrictedVolume.coordinateCurrentForm,TranscendentalNakamaye.FirstPass.NumericalRestrictedVolume.patchDensity,TranscendentalNakamaye.FirstPass.NumericalRestrictedVolume.RegularPieceFamily,TranscendentalNakamaye.FirstPass.NumericalRestrictedVolume.acTopMass,TranscendentalNakamaye.FirstPass.NumericalRestrictedVolume.acTopMassData,TranscendentalNakamaye.FirstPass.NumericalRestrictedVolume.numericalRestrictedVolume,TranscendentalNakamaye.FirstPass.NumericalRestrictedVolume.numericalRestrictedVolume_eq_zero_of_subset_nonNefLocus,TranscendentalNakamaye.FirstPass.NumericalRestrictedVolume.nullLocus,TranscendentalNakamaye.FirstPass.NumericalRestrictedVolume.mem_nullLocus}
\leanok
Let \(\alpha\) be pseudoeffective and let \(V\subset X\) be a closed
irreducible analytic subset of dimension \(k>0\). For \(\epsilon>0\),
let \(\mathcal A_\epsilon\) be the currents \(T\in\alpha\) with smooth
analytic singularities, with \(T\ge-\epsilon\omega\), and with polar
set not containing \(V\). If \(V\subset E_{\mathrm{nn}}(\alpha)\), define
\(\operatorname{RV}_{X\mid V}(\alpha)=0\). Otherwise define the
numerical restricted volume by the absolutely continuous integral
\[
 \operatorname{RV}_{X\mid V}(\alpha)=
 \lim_{\epsilon\downarrow0}\sup_{T\in\mathcal A_\epsilon}
 \int_{V_{\mathrm{reg}}}
 \left(\bigl((T+\epsilon\omega)|_{V_{\mathrm{reg}}}\bigr)_{
 \mathrm{ac}}\right)^k.
\]
The subscript \(\mathrm{ac}\) denotes the absolutely continuous part
of the coefficients of the restricted current.
Off the polar set \(P\) of \(T\) the current \(T+\epsilon\omega\) is a smooth
semipositive form, and \(P\cap V_{\mathrm{reg}}\) is Lebesgue-null in
\(V_{\mathrm{reg}}\); hence the integral equals the integral of the top power of
the smooth form \((T+\epsilon\omega)|_{V_{\mathrm{reg}}\setminus P}\), which is
the supremum, over finitely many disjoint measurable pieces of
\(V_{\mathrm{reg}}\setminus P\) lying in local holomorphic parametrizations of
\(V\), of the sums of the coordinate integrals of its density.
Finally,
\[
 \operatorname{Null}(\alpha)=
 \bigcup_{\substack{V\subset X\text{ closed irreducible analytic}\\
 \dim V>0,\ \operatorname{RV}_{X\mid V}(\alpha)=0}}V.
\]
\end{definition}

\begin{theorem}[Monotonicity of non-pluripolar mass]
\label{ana:mass-monotonicity}
\uses{mmass:monotonicity,cap:identity-principle}
\source{transcendental-nakamaye-copy4-20260911:it:monotone,ddnl-mixed-monotonicity:2.4,xia-sgpt:page-0073,xia-sgpt:page-0075}
\lean{TranscendentalNakamaye.Pending.GlobalMixedMassMonotone.globalMixedMass_mono_of_connected,TranscendentalNakamaye.Pending.GlobalMixedMassMonotone.globalMixedMass_mono_of_isThetaPshPotential,TranscendentalNakamaye.Pending.GlobalMixedMassMonotone.exists_ne_bot_of_moreSingularOn}
\leanok
Let \(Z\) be a connected compact K\"ahler manifold of dimension \(d\).
Let \(\theta_j\) be smooth closed real \((1,1)\)-forms, and let
\(u_j,v_j\in\operatorname{PSH}(Z,\theta_j)\), for \(1\le j\le d\).
If \(u_j\preceq v_j\) for every \(j\), then
\[
 \int_Z\langle\theta_{1,u_1}\wedge\cdots\wedge\theta_{d,u_d}\rangle
 \le
 \int_Z\langle\theta_{1,v_1}\wedge\cdots\wedge\theta_{d,v_d}\rangle.
\]
\end{theorem}

\begin{proof}
\uses{mmass:monotonicity,cap:identity-principle}
\leanok
The products are taken with \(d d^c=(2\pi)^{-1}\mathrm i\partial\bar\partial\) in both the convention and
Theorem~\ref{mmass:monotonicity}, so no renormalization enters. By Lemma~\ref{cap:identity-principle}, a theta-PSH
function on the connected manifold \(Z\) which is finite at one point is a locally nondegenerate theta-PSH potential;
this applies to \(u_j\), and to \(v_j\), which is finite wherever \(u_j\) is, because \(u_j\le v_j+C_j\).
If \(d>0\), Theorem~\ref{mmass:monotonicity} gives the inequality. If \(d=0\), the two empty products are the same
measure, so the inequality is an equality. (The argument of Theorem~\ref{mmass:monotonicity} uses \(d>0\) nowhere.)
\end{proof}

\begin{proposition}[Mass calculus]
\label{ana:mass-calculus}
\uses{ana:geometric-setting,gnpp:global-measure,ana:restricted-volume}
\source{transcendental-nakamaye-copy4-20260911:sec:npp}
\lean{TranscendentalNakamaye.Pending.GlobalNonPluripolar.globalMixedMeasure,TranscendentalNakamaye.Pending.MixedMassProductCalculus.globalMixedMeasure_sumPot_expand,TranscendentalNakamaye.Pending.MixedMassProductCalculus.globalMixedMeasure_update_zero,TranscendentalNakamaye.Pending.GlobalNonPluripolar.globalMixedMeasure_doesNotChargeManifoldPluripolar,TranscendentalNakamaye.Pending.RestrictedMassEquality.acTopMass_eq_diagonalPotentialMass_restrict',TranscendentalNakamaye.Pending.RestrictedMassEquality.acTopMass_eq_diagonalPotentialMass_restrict,TranscendentalNakamaye.Pending.RestrictedMassEquality.diagonalPotentialMass_restrict_le_acTopMass}
\leanok
Mixed non-pluripolar products of positive currents are positive and
multilinear with respect to nonnegative real combinations, and assign zero
mass to pluripolar sets. If an almost positive current has smooth analytic
singularities and restricts to a closed complex submanifold, its positive
perturbation has equal absolutely continuous and non-pluripolar top-degree
integrals on that submanifold.
\end{proposition}

\begin{proof}
\uses{gnpp:global-measure,mmass:product-calculus-np,cap:identity-principle,gnpp:noncharging,ana:restricted-volume,ana:submanifold-complex-structure,pc:f-restricted-volume-le,ana:fullspace-density-localization}
\leanok
The mixed products are the global non-pluripolar measures of Definition~\ref{gnpp:global-measure}; they are positive
Borel measures. A function of \(\operatorname{PSH}(X,\theta)\) is finite at some point, so on the connected manifold \(X\)
it is a locally nondegenerate \(\theta\)-PSH potential (Lemma~\ref{cap:identity-principle}). For such potentials
Lemma~\ref{mmass:product-calculus-np} expands the product of the nonnegative combinations
\(\theta_i=\sum_ac_{ia}\beta_a\), \(u_i=\sum_ac_{ia}q_a\) multilinearly, zero coefficients included. Pluripolar sets carry no
mass by Theorem~\ref{gnpp:noncharging}.

For the second sentence, let \(T=\theta+dd^cu\) have smooth analytic singularities, with \(T\ge-\epsilon\omega\), and let
\(V\) be the closed submanifold, of dimension \(k\), not contained in the polar set \(P\) of \(u\). As in
Definition~\ref{ana:restricted-volume}, where the absolutely continuous integral along \(V\) is defined, \(V\) is
irreducible, hence connected; with the atlas of Lemma~\ref{ana:submanifold-complex-structure} it is a compact complex
manifold. The restriction \(u|_V\) is \((\theta+\epsilon\omega)|_V\)-psh and finite at a point of \(V\setminus P\), hence a
potential on the connected \(V\), and \(P\cap V\) is its polar set.

The absolutely continuous integral is at most \(\int_V\langle((\theta+\epsilon\omega)|_V+dd^c(u|_V))^k\rangle\): this is the
bound in the proof of Lemma~\ref{pc:f-restricted-volume-le}. Conversely, every \(y\in V\) is the centre of a regular
parametrization \(\gamma_y\) of \(V\) (Definition~\ref{ana:restricted-volume}). The points of \(V\) whose flattened coordinate
in the chart of \(\gamma_y\) lies in a small ball \(B_y\) form an open neighbourhood \(N_y\) of \(y\) in \(V\), on which
\(\gamma_y\), read in the complex manifold \(V\), is a regular parametrization of the whole of \(V\). Finitely many
\(N_{y_1},\dots,N_{y_m}\) cover the compact \(V\). Put
\[
 B_i=\Bigl(N_{y_i}\setminus\bigcup_{j<i}N_{y_j}\Bigr)\setminus P,\qquad D_i=B_{y_i}\cap\gamma_{y_i}^{-1}(B_i).
\]
The sets \(B_i\) are Borel (\(P\) is the polar set of an upper semicontinuous function) and pairwise disjoint, and
\(\gamma_{y_i}(D_i)=B_i\), so the \(D_i\) form a finite family of disjoint regular pieces avoiding \(P\). On \(D_i\) the density
of \((T+\epsilon\omega)|_V\) in the parametrization \(\gamma_{y_i}\) of \(V\subseteq X\) equals that of the restricted current in
the same parametrization of the manifold \(V\) (chain rule, off the poles), and by Lemma~\ref{ana:fullspace-density-localization}
on \(V\) its integral over \(D_i\) is the non-pluripolar measure of \(B_i\). Summing over \(i\), the absolutely continuous
integral is at least the non-pluripolar mass of \(V\setminus P\), which is the mass of \(V\), because the polar set of a
potential is pluripolar and carries no mass (Theorem~\ref{gnpp:noncharging}).
\end{proof}

\begin{theorem}[Minimal-singularity formula for restricted volume]
\label{ana:ct-minimal-mass}
\uses{ana:restricted-volume,ana:actual-minimal-restricted-mass-formula,ana:actual-non-nef-smooth-kahler-perturbation,trace:actual-connected-submanifold-irreducible-carrier,ana:minimal-singularity-lelong-vanishing,fpenv:v-universal}
\source{ct22-restricted-volumes:sect2,ct22-restricted-volumes:nonp}
\lean{TranscendentalNakamaye.Pending.ActualMinimalRestrictedMassFormula.exists_minimal_restriction_mass_inf_and_tendsto,TranscendentalNakamaye.Pending.ActualNonNefSmoothKahlerPerturbation.exists_kahlerCurrent_smoothNear_of_not_subset_nonNefLocus}
\leanok
Suppose that $\alpha$ is pseudoeffective and that the positive-dimensional
closed connected complex submanifold $V$ is not contained in
$E_{\mathrm{nn}}(\alpha)$. For every $\epsilon>0$, the class
$\alpha+\epsilon\{\omega\}$ contains a K\"ahler current smooth near
some point of $V$, and a positive current
$T_{\min,\epsilon}$ with minimal singularities in that class restricts to $V$.
Its restricted mass is finite and
\[
 \operatorname{RV}_{X\mid V}(\alpha)
 =\lim_{\epsilon\downarrow0}\int_V
   \left\langle(T_{\min,\epsilon}|_V)^k\right\rangle,
 \qquad k=\dim V.
\]
Here the volume on the left uses positive real singularity coefficients and
remainders smooth off the poles (Definition~\ref{as:smooth-off-pole-type}): the
supremum defining it ranges over that class of presentations.
\end{theorem}
\begin{proof}
Theorem~\ref{trace:actual-connected-submanifold-irreducible-carrier}
realizes the given $V$ as an irreducible analytic carrier without changing
either its points or its dimension. Choose a smooth closed representative
$\theta\in\alpha$, and put
$U_\epsilon=V_{\theta+\epsilon\omega}$.
Theorem~\ref{ana:actual-non-nef-smooth-kahler-perturbation} supplies, for
every $\epsilon>0$, a K\"ahler current in the perturbed class smooth near
a point of $V$. In particular that class is big.
Proposition~\ref{ana:minimal-singularity-lelong-vanishing} makes
$U_\epsilon$ genuinely $(\theta+\epsilon\omega)$-PSH. The envelope
property~\ref{fpenv:v-universal}, after normalizing any competing potential
by an additive constant, shows that
\[
 T_{\min,\epsilon}=\theta+\epsilon\omega+dd^cU_\epsilon
\]
has minimal singularities.

Theorem~\ref{ana:actual-minimal-restricted-mass-formula} gives genuine
restrictions for this entire family and finite masses
\[
 M(\epsilon)=\int_V
 \left\langle(T_{\min,\epsilon}|_V)^k\right\rangle.
\]
To see the volume comparison explicitly, let $E(\epsilon)$ be the
admissible AC mass supremum defining restricted volume, with positive
real singularity coefficients and bounded remainders smooth off the
poles. The upper and lower comparisons give
\[
 M(\epsilon/2)\le E(\epsilon)\le M(\epsilon).
\]
Consequently
$\operatorname{RV}_{X\mid V}(\alpha)=\inf_{\epsilon>0}M(\epsilon)$.
If $0<\delta<\epsilon$, the normalized envelope $U_\delta$ is a
competitor for $U_\epsilon$, so $U_\delta\le U_\epsilon$.
Semipositive background enlargement and the ordered-potential mass
comparison imply $M(\delta)\le M(\epsilon)$.
Thus the right limit of $M$ at zero is its infimum, proving the stated
formula, including when this infimum is zero.
\end{proof}

\begin{lemma}[Monotonicity of the admissible mass envelope]
\label{ana:restricted-volume-monotone}
\uses{ana:restricted-volume,ana:restricted-volume-scaffold}
\source{ct22-restricted-volumes:voldef,transcendental-nakamaye-copy4-20260911:eq:def-restricted-volume}
\lean{TranscendentalNakamaye.Pending.RestrictedVolumeMonotonicity.restrictedMassEnvelope_mono,TranscendentalNakamaye.Pending.RestrictedVolumeMonotonicity.numericalRestrictedVolume_eq_numericalRestrictedVolumeInf}
\leanok
For \(0<\epsilon\le\epsilon'\),
\[
 \sup_{T\in\mathcal A_\epsilon}\int_{V_{\mathrm{reg}}}
 \bigl((T+\epsilon\omega)|_{V_{\mathrm{reg}}}\bigr)_{\mathrm{ac}}^k
 \le
 \sup_{T\in\mathcal A_{\epsilon'}}\int_{V_{\mathrm{reg}}}
 \bigl((T+\epsilon'\omega)|_{V_{\mathrm{reg}}}\bigr)_{\mathrm{ac}}^k .
\]
Consequently, if \(V\not\subset E_{\mathrm{nn}}(\alpha)\), the limit defining
\(\operatorname{RV}_{X\mid V}(\alpha)\) exists in \([0,\infty]\) and equals the
infimum over \(\epsilon>0\) of these suprema.
\end{lemma}
\begin{proof}
Let \(T\in\mathcal A_\epsilon\). Since \(T+\epsilon'\omega=(T+\epsilon\omega)
+(\epsilon'-\epsilon)\omega\ge0\), with the same potential and polar set,
\(T\in\mathcal A_{\epsilon'}\). Off the polar set \(P\) of \(T\), the restrictions
to \(V_{\mathrm{reg}}\setminus P\) satisfy
\((T+\epsilon'\omega)|\ge(T+\epsilon\omega)|\ge0\) as smooth \((1,1)\)-forms, and
the top power of a semipositive \((1,1)\)-form increases when a semipositive form
is added (the determinant of Hermitian matrices is monotone on the positive
cone). Integrating over \(V_{\mathrm{reg}}\setminus P\), which carries the whole
absolutely continuous integral, gives the inequality for each \(T\), hence for
the suprema. A nondecreasing function of \(\epsilon\in(0,\infty)\) with values in
\([0,\infty]\) has a right limit at \(0\), equal to its infimum.
\end{proof}

\begin{lemma}[Elementary comparisons of the loci]
\label{ana:locus-properties-elementary}
\uses{ana:nonkahler-locus}
\source{ct22-restricted-volumes:nonnef,transcendental-nakamaye-copy4-20260911:sec:nK}
\lean{TranscendentalNakamaye.Pending.LocusProperties.nonKahlerLocus_ne_univ,TranscendentalNakamaye.Pending.LocusProperties.nonNefLocus_subset_nonKahlerLocus}
\leanok
Let \(X\) be a complex manifold of dimension \(n\), let \(\omega\) be a K\"ahler
form and \(\alpha\) a real \((1,1)\)-class. Then
\(E_{\mathrm{nn}}(\alpha)\subset E_{\mathrm{nK}}(\alpha)\). If \(n>0\), \(X\) is
nonempty and \(\alpha\) is big, then \(E_{\mathrm{nK}}(\alpha)\ne X\).
\end{lemma}
\begin{proof}
\leanok
A K\"ahler current \(T\) in \(\alpha\) gives the K\"ahler current
\(T+\epsilon\omega\) in \(\alpha+\epsilon\{\omega\}\) with the same potential, hence
the same Lelong numbers; so a point with positive Lelong number for every
K\"ahler current in \(\alpha+\epsilon\{\omega\}\) has positive Lelong number for
every K\"ahler current in \(\alpha\). If \(\alpha\) is big, it carries a K\"ahler
current \(T=\theta+dd^cu\), and \(u\) is finite at some point \(y\). A positive
Lelong number of \(u\) at \(y\) would force \(u(y)=-\infty\) by the circle-mean
inequality, so \(y\notin E_{\mathrm{nK}}(\alpha)\).
\end{proof}

\begin{lemma}[Stokes' theorem and finite volume on a subvariety]
\label{ana:subvariety-stokes-mass}
\source{demailly-cadg:III.2.7}
\uses{ana:restricted-volume}
\lean{TranscendentalNakamaye.Pending.SubvarietyStokesMass.acTopMass_eq_of_smooth,TranscendentalNakamaye.Pending.SubvarietyStokesMass.acTopMass_ne_top_of_zero_potential}
\leanok
Let \((X,\beta)\) be a compact complex manifold with a K\"ahler form and let
\(V\subset X\) be an irreducible analytic subvariety of dimension \(k\). Then
\(\int_{V_{\mathrm{reg}}}\beta^k<\infty\), and for every smooth function \(v\) on
\(X\) with \(\beta+dd^cv\ge0\),
\[
 \int_{V_{\mathrm{reg}}}\bigl((\beta+dd^cv)|_{V_{\mathrm{reg}}}\bigr)^{k}
 =\int_{V_{\mathrm{reg}}}\beta^{k}.
\]
\end{lemma}
\begin{proof}
\uses{ana:background-mass-finite,ana:background-mass-stokes}
\leanok
By Lelong's theorem, \(V_{\mathrm{reg}}\) has locally finite \(2k\)-dimensional
volume near every point of \(X\), and the current of integration \([V]\),
\(\langle[V],\eta\rangle=\int_{V_{\mathrm{reg}}}\eta\), is a closed positive current
of bidimension \((k,k)\). Since \(X\) is compact and \(\beta^k/k!\) restricts to the
volume form of the metric \(\beta\) on \(V_{\mathrm{reg}}\) (Wirtinger), the first
integral is finite. For the second, \((\beta+dd^cv)^k-\beta^k=dd^cv\wedge S\) with
\(S=\sum_{i=0}^{k-1}(\beta+dd^cv)^i\wedge\beta^{k-1-i}\) a smooth closed form, so
\((\beta+dd^cv)^k-\beta^k=d(d^cv\wedge S)\) is exact with a smooth global primitive,
and \(\langle[V],d\eta\rangle=0\) because \([V]\) is closed.
\end{proof}

\begin{definition}[Background mass of a subvariety]
\label{ana:background-mass}
\uses{ana:restricted-volume}
\lean{TranscendentalNakamaye.Pending.SubvarietyStokesMassInterfaces.backgroundMass}
\leanok
Let \((X,\beta)\) be a compact complex manifold with a K\"ahler form and let
\(V\subset X\) be an irreducible analytic subvariety of dimension \(k\). The
\emph{background mass} of \(V\) is \(\int_{V_{\mathrm{reg}}}\beta^{k}\), computed as the
supremum, over all finite families of disjoint measurable pieces \(D_i\) of regular
parametrizations of \(V\), of \(\sum_i\int_{D_i}\) of the parameter density of \(\beta\).
It depends only on \((X,\beta)\) and \(V\): no class, current, potential or polar set
enters.
\end{definition}

\begin{lemma}[Change of regular patch]
\label{ana:patch-transition}
\uses{ana:background-mass}
\lean{TranscendentalNakamaye.Pending.RegularPatchTransition.transition,
TranscendentalNakamaye.Pending.RegularPatchTransition.toManifold_transition,
TranscendentalNakamaye.Pending.RegularPatchTransition.differentiableAt_transition,
TranscendentalNakamaye.Pending.RegularPatchTransition.injOn_transition,
TranscendentalNakamaye.Pending.RegularPatchTransition.actualTopPowerDensity_patchCoordinateForm_transition,
TranscendentalNakamaye.Pending.RegularPatchTransition.lintegral_actualTopPowerDensity_transition}
\leanok
Let \(P,Q\) be two regular parametrizations of \(V\) of dimension \(k\). On the
open set of \(Q\)-parameters whose image lies in the chart and the flattening
domain of \(P\) there is a holomorphic, injective change of parameters \(\tau\)
with \(P\circ\tau=Q\), and for every smooth real \((1,1)\)-form \(\theta\) the
parameter densities satisfy
\[
 \rho_{\theta,Q}(w)=\det{}_{\mathbf R}\bigl(d\tau_w\bigr)\,\rho_{\theta,P}(\tau w).
\]
Consequently, for every measurable piece \(S\) of that open set,
\(\int_{\tau(S)}\rho_{\theta,P}=\int_S\rho_{\theta,Q}\): the two patches assign
the same mass to the part of \(V\) they both parametrize.
\end{lemma}
\begin{proof}
\uses{ana:background-mass}
\leanok
Let \(\varphi_P\) be the preferred chart at the centre of \(P\) and \(\Phi_P\) its
flattening, which sends \(V\) into the coordinate \(k\)-plane, and let \(\mathrm{pr}\)
be the projection onto that plane. Put
\(\tau(w)=\mathrm{pr}\bigl(\Phi_P(\varphi_P(Q(w)))\bigr)\). Because \(Q(w)\in V\),
the point \(\Phi_P(\varphi_P(Q(w)))\) lies in the coordinate plane, so \(\mathrm{pr}\)
loses nothing and \(\tau(w)\) is a \(P\)-parameter; applying \(\Phi_P^{-1}\) and
\(\varphi_P^{-1}\) gives \(P(\tau w)=Q(w)\). No inverse of \(Q\) is used, and
\(\tau\) is holomorphic, being a composite of \(Q\), a chart transition, \(\Phi_P\)
and a linear projection. It is injective because \(Q\) is.

Write \(T=\varphi_P\circ\varphi_Q^{-1}\) for the chart transition. Then
\(\gamma_P\circ\tau=T\circ\gamma_Q\) on the whole open set, so by the chain rule
\(dT\circ d\gamma_Q=d\gamma_P\circ d\tau\) there. The coefficient of \(\theta\) in
the chart of \(Q\) is its coefficient in the chart of \(P\) precomposed with
\(dT\), so the pullback of \(\theta\) along \(\gamma_Q\) is its pullback along
\(\gamma_P\) at \(\tau w\), precomposed with \(d\tau_w\); the top exterior power of
a precomposition multiplies by the real determinant, which gives the density
identity. Finally \(\tau\) is holomorphic, so \(\det_{\mathbf R}(d\tau_w)\) is the
squared modulus of the complex determinant and in particular nonnegative; the
change-of-variables formula for an injective differentiable map then turns the
density identity into the equality of the two parameter integrals.
\end{proof}

\begin{lemma}[The regular-piece mass depends only on the region covered]
\label{ana:piece-mass-comparison}
\uses{ana:background-mass,ana:patch-transition}
\lean{TranscendentalNakamaye.Pending.RegularPieceMassComparison.inter_preimage_image_toManifold,
TranscendentalNakamaye.Pending.RegularPieceMassComparison.backgroundPieceMass_mono_of_image_subset,
TranscendentalNakamaye.Pending.RegularPieceMassComparison.backgroundPieceMass_congr_of_iUnion_image_eq,
TranscendentalNakamaye.Pending.RegularPieceMassComparison.unionFamily,
TranscendentalNakamaye.Pending.RegularPieceMassComparison.directed_backgroundPieceMass}
\leanok
Let \(V\) be an irreducible analytic subvariety of dimension \(k\) of the compact
K\"ahler manifold \((X,\beta)\). A \emph{regular-piece family} is a finite family
\(F=(D_i,\gamma_i)_{i<N}\) of measurable sets \(D_i\) of parameters of regular
parametrizations \(\gamma_i\) of \(V\) whose images \(\gamma_i(D_i)\subset
V_{\mathrm{reg}}\) are pairwise disjoint; its background parameter mass is
\[
 \mathbf m(F)=\sum_{i<N}\int_{D_i}\rho_{\beta,\gamma_i}.
\]
If every image of \(F\) is contained in \(\bigcup_j\gamma'_j(D'_j)\) for a second
regular-piece family \(G=(D'_j,\gamma'_j)_{j<N'}\), then \(\mathbf m(F)\le\mathbf
m(G)\). Consequently two families with the same union of images have the same
mass, and any two families are dominated by a single family whose images cover
both; the supremum defining the background mass of \(V\) is therefore directed.
Only the background density \(\rho_\beta\) enters: no class, current, potential
or polar set.
\end{lemma}
\begin{proof}
\uses{ana:background-mass,ana:patch-transition}
\leanok
Write \(P_j=\gamma'_j\) and \(Q_i=\gamma_i\), and let \(\Delta_{ji}\) and
\(\tau_{ji}\) be the transition domain and the change of parameters from \(Q_i\)
to \(P_j\) of Lemma~\ref{ana:patch-transition}.

\emph{The overlap is a set of parameters.} Lemma~\ref{ana:patch-transition} says
that whenever a \(Q_i\)-parameter and a \(P_j\)-parameter have the same image,
the first lies in \(\Delta_{ji}\) and \(\tau_{ji}\) carries it to the second.
Hence, for measurable \(D_i\) and \(D'_j\) inside the respective parameter
domains,
\[
 S_{ij}:=D_i\cap Q_i^{-1}\bigl(P_j(D'_j)\bigr)
 =D_i\cap\Delta_{ji}\cap\tau_{ji}^{-1}(D'_j).
\]
The right-hand side lives entirely in the parameter space: \(\Delta_{ji}\) is
open and \(\tau_{ji}\) is holomorphic on it, so \(S_{ij}\) is measurable, and
\(\tau_{ji}(S_{ij})\subset D'_j\) is measurable because \(\tau_{ji}\) is
injective and differentiable there. Nothing is asserted about the measurability
of \(P_j(D'_j)\) inside \(X\).

\emph{Slicing one piece of \(F\).} For fixed \(i\), the sets \(S_{ij}\),
\(j<N'\), are pairwise disjoint: they are the \(Q_i\)-preimages of the pairwise
disjoint sets \(P_j(D'_j)\). They exhaust \(D_i\), because by hypothesis every
\(Q_i(w)\) with \(w\in D_i\) lies in some \(P_j(D'_j)\). Therefore
\[
 \int_{D_i}\rho_{\beta,Q_i}=\sum_{j<N'}\int_{S_{ij}}\rho_{\beta,Q_i}
 =\sum_{j<N'}\int_{\tau_{ji}(S_{ij})}\rho_{\beta,P_j},
\]
the second equality being the change-of-parameters identity of
Lemma~\ref{ana:patch-transition} applied on \(S_{ij}\subset\Delta_{ji}\).

\emph{Reassembling along \(G\).} Fix \(j\) and let \(i\) vary. The sets
\(\tau_{ji}(S_{ij})\) are pairwise disjoint: \(P_j\) is injective and
\(P_j(\tau_{ji}(S_{ij}))=Q_i(S_{ij})\subset\gamma_i(D_i)\), and the images of
\(F\) are pairwise disjoint. All of them lie in \(D'_j\), so
\[
 \sum_{i<N}\int_{\tau_{ji}(S_{ij})}\rho_{\beta,P_j}
 =\int_{\bigsqcup_{i<N}\tau_{ji}(S_{ij})}\rho_{\beta,P_j}
 \le\int_{D'_j}\rho_{\beta,P_j}.
\]
Exchanging the two finite sums gives
\(\mathbf m(F)\le\mathbf m(G)\). Applying this in both directions gives equality
when the two unions of images coincide.

\emph{Directedness.} Given \(F\) and \(G\), let \(G'\) consist, for each \(j\),
of the \(G\)-parameters in \(D'_j\) whose image no piece of \(F\) already sees.
Writing \(\Delta'_{ij}\) and \(\tau'_{ij}\) for the transition domain and the change
of parameters from \(P_j\) to \(Q_i\), the overlap identity with the roles of the
two patches exchanged gives
\[
 D'_j\cap P_j^{-1}\bigl(Q_i(D_i)\bigr)
 =D'_j\cap\Delta'_{ij}\cap(\tau'_{ij})^{-1}(D_i),
\]
so that set of fresh parameters is
\(D'_j\setminus\bigcup_{i<N}\bigl(D'_j\cap\Delta'_{ij}\cap(\tau'_{ij})^{-1}(D_i)\bigr)\),
a measurable subset of \(D'_j\).
The family consisting of the pieces of \(F\) together with those of \(G'\) is
again a regular-piece family: the images of its \(G'\)-pieces are disjoint from
all images of \(F\) by construction and from each other because \(G\) is a
family. Its union of images is the union of those of \(F\) and \(G\), so
monotonicity bounds both \(\mathbf m(F)\) and \(\mathbf m(G)\) by its mass.
\end{proof}

\begin{lemma}[$dd^cv$ of a smooth function is a global smooth two-form]
\label{ana:ddc-global-form}
\uses{hgeom:dc-naturality,hgeom:exterior-derivative}
\lean{TranscendentalNakamaye.Pending.SmoothDDCForm.dcForm,
TranscendentalNakamaye.Pending.SmoothDDCForm.ddcForm}
\leanok
Let \(X\) be a complex manifold and let \(v\) be a real function on \(X\) that is
smooth in every holomorphic chart. Reading \(v\) in each chart and applying the
normalized operators there produces chartwise coefficient fields that obey the
tensor transformation law. They are therefore a global smooth real one-form
\(d^cv\) and a global smooth real two-form \(dd^cv\) on \(X\).
\end{lemma}
\begin{proof}
\uses{hgeom:dc-naturality,hgeom:exterior-derivative}
\leanok
Let \(e,f\) be holomorphic charts and let \(T=f\circ e^{-1}\) be the transition,
a biholomorphism between open subsets of \(\mathbf C^n\). On the overlap,
\(v\circ e^{-1}=(v\circ f^{-1})\circ T\), so Lemma~\ref{hgeom:dc-naturality},
applied with \(F=T\) and \(\rho=v\circ f^{-1}\), gives
\[
 d^c(v\circ e^{-1})_z=\bigl(d^c(v\circ f^{-1})\bigr)_{T(z)}\circ DT_z
\]
at every point \(z\) of the overlap. That is exactly the transformation law for a
global one-form, so the family \((d^c(v\circ e^{-1}))_e\) is a smooth real
one-form \(d^cv\) on \(X\). Its global exterior derivative
(Definition~\ref{hgeom:exterior-derivative}) is a smooth real two-form whose
coefficient in the chart \(e\) is \(dd^c(v\circ e^{-1})\); this is \(dd^cv\).
\end{proof}

\begin{lemma}[The comparison holds for every smooth form, and for a smooth potential]
\label{ana:piece-mass-comparison-smooth}
\uses{ana:piece-mass-comparison,ana:patch-transition,ana:ddc-global-form,ana:restricted-volume}
\lean{TranscendentalNakamaye.Pending.PatchFormMass.patchFormDensity,
TranscendentalNakamaye.Pending.PatchFormMass.formPieceMass,
TranscendentalNakamaye.Pending.PatchFormMass.formPieceMass_mono_of_image_subset,
TranscendentalNakamaye.Pending.PatchFormMass.formPieceMass_congr_of_iUnion_image_eq,
TranscendentalNakamaye.Pending.PatchFormMass.directed_formPieceMass,
TranscendentalNakamaye.Pending.ACTopMassComparison.currentForm,
TranscendentalNakamaye.Pending.ACTopMassComparison.coordinateCurrentForm_eq,
TranscendentalNakamaye.Pending.ACTopMassComparison.acPieceMass_mono_of_image_subset,
TranscendentalNakamaye.Pending.ACTopMassComparison.directed_acPieceMass,
TranscendentalNakamaye.Pending.ACTopMassComparison.acTopMass_eq_iSup_formPieceMass}
\leanok
Let \(a\) be any smooth real two-form on \(X\) and write
\(\mathbf m_a(F)=\sum_{i<N}\int_{D_i}\rho_{a,\gamma_i}\) for the parameter mass a
regular-piece family \(F\) assigns to \(a\). Then
Lemma~\ref{ana:piece-mass-comparison} holds verbatim for \(\mathbf m_a\): a
family whose images are covered by those of a second family has no larger mass,
two families with the same union of images have the same mass, and the masses of
the families avoiding a fixed subset of \(X\) form a directed set.

Moreover, let \(\theta\) be a smooth closed real \((1,1)\)-form, let
\(\varepsilon>0\), and let \(v\) be a real function on \(X\) that is smooth in
every chart. Then the coordinate forms of the current
\(\theta+\varepsilon\beta+dd^cv\) are the chart coefficients of the single smooth
two-form \(a_v=\theta+\varepsilon\beta+dd^cv\), so the restricted mass along
\(V\) of an admissible presentation with representative \(\theta\), error
\(\varepsilon\) and potential \(v\) is \(\sup_F\mathbf m_{a_v}(F)\), the supremum
being over the same regular-piece families as for the background mass: a real
potential has empty polar set.
\end{lemma}
\begin{proof}
\uses{ana:piece-mass-comparison,ana:patch-transition,ana:ddc-global-form}
\leanok
The change-of-parameters identity of Lemma~\ref{ana:patch-transition} is proved
for an arbitrary smooth real two-form, and the proof of
Lemma~\ref{ana:piece-mass-comparison} uses nothing about the family beyond the
measurability of its parameter pieces, their inclusion in the parameter domains,
and the pairwise disjointness of their images; the set a family is required to
avoid never enters. Both assertions of the first paragraph therefore follow by
repeating that proof with \(\rho_\beta\) replaced by \(\rho_a\), and with the two
families allowed to avoid different sets.

For the second paragraph, \(dd^cv\) is a global smooth two-form by
Lemma~\ref{ana:ddc-global-form}, so \(a_v=\theta+\varepsilon\beta+dd^cv\) is one
as well, and its coefficient in the chart \(e\) is the coefficient of
\(\theta+\varepsilon\beta\) there plus \(dd^c(v\circ e^{-1})\) — which is by
definition the coordinate form of the current of the presentation in that chart.
The parameter densities of the presentation are therefore the parameter densities
\(\rho_{a_v}\), and the two masses are suprema of the same quantity. Finally a
real-valued potential takes the value \(-\infty\) nowhere, so its polar set is
empty and the families indexing the two suprema are the same.
\end{proof}

\begin{lemma}[Every compact part of the regular locus is covered by one family]
\label{ana:piece-family-compact-cover}
\uses{ana:restricted-volume,ana:piece-mass-comparison}
\lean{TranscendentalNakamaye.Pending.RegularPieceCompactCover.emptyFamily,
TranscendentalNakamaye.Pending.RegularPieceCompactCover.singletonFamily,
TranscendentalNakamaye.Pending.RegularPieceCompactCover.exists_family_iUnion_image_superset,
TranscendentalNakamaye.Pending.RegularPieceCompactCover.exists_family_covering_of_isCompact}
\leanok
Let \(V\) be a closed irreducible analytic subvariety of dimension \(k\) of the
complex manifold \(X\), let \(A\subset X\) be closed, and let \(K\subset
V_{\mathrm{reg}}\) be compact with \(K\cap A=\emptyset\). Then there is a
regular-piece family \(F\) whose images avoid \(A\) and satisfy
\(K\subset\bigcup_i\gamma_i(D_i)\).

More primitively: finitely many measurable parameter sets \(D_1,\dots,D_N\) of
finitely many regular parametrizations \(\gamma_1,\dots,\gamma_N\), whose images
avoid \(A\), are all contained in the union of the images of a SINGLE
regular-piece family avoiding \(A\); no disjointness of the \(\gamma_i(D_i)\) is
assumed.
\end{lemma}
\begin{proof}
\uses{ana:restricted-volume,ana:piece-mass-comparison}
\leanok
\emph{The primitive statement.} Induct on \(N\). For \(N=0\) take the empty
family. For the step, let \(F\) be a family covering the first \(N\) images and
let \(G\) be the one-piece family \((D_{N+1},\gamma_{N+1})\) --- legitimate
because every point parametrized by a regular parametrization is a regular point
of the certified dimension, and its single image avoids \(A\) by hypothesis. The
union family of Lemma~\ref{ana:piece-mass-comparison} has images covering those
of \(F\) and of \(G\), hence all \(N+1\) of the given images. This is exactly
where the overlaps are removed: the parameters of \(G\) whose image \(F\)
already sees form the overlap set of that lemma, which lives in the parameter
space and is therefore measurable.

\emph{The compact statement.} Let \(x\in K\). Since \(x\) is a regular point of
dimension \(k\), a single regular parametrization \(\gamma_x\) is defined near
the parameter \(0\) with \(\gamma_x(0)=x\); choose \(r_x>0\) with
\(B(0,r_x)\) inside its parameter domain. The flattening defining \(\gamma_x\)
is a biholomorphism of an open set of the chart carrying \(V\) onto a
coordinate plane, so \(\gamma_x(B(0,r_x))\) is relatively open in \(V\): there
is an open \(U_x\subset X\) with \(x\in U_x\) and \(V\cap U_x\subset
\gamma_x(B(0,r_x))\). Put \(W_x=U_x\setminus A\), open and containing \(x\)
because \(A\) is closed and \(x\notin A\).

Compactness of \(K\) leaves finitely many \(W_{x_1},\dots,W_{x_N}\) covering
\(K\). Put \(D_i=B(0,r_{x_i})\cap\gamma_{x_i}^{-1}(W_{x_i})\), an open --- in
particular measurable --- subset of the parameter domain whose image lies in
\(W_{x_i}\), hence avoids \(A\). These images cover \(K\): if \(y\in K\) then
\(y\in W_{x_i}\) for some \(i\), and \(y\in V\), so \(y=\gamma_{x_i}(w)\) for
some \(w\in B(0,r_{x_i})\), and then \(w\in D_i\). The primitive statement now
supplies a single family whose images contain all the \(\gamma_i(D_i)\), hence
\(K\).
\end{proof}

\begin{lemma}[Continuity of the parameter density]
\label{ana:patch-density-continuity}
\uses{ana:restricted-volume,ana:patch-transition}
\lean{TranscendentalNakamaye.Pending.RegularPatchTransition.continuousOn_actualTopPowerDensity_patchCoordinateForm}
\leanok
Let \(a\) be a smooth real two-form on \(X\) and \(\gamma\) a regular
parametrization of \(V\). Then \(w\mapsto\mathcal D_a(\gamma,w)\), the density of
the top power of \(\gamma^*a\) with respect to Lebesgue measure in the
parameters, is continuous on the parameter domain. In particular its integral
over a compact parameter set inside that domain is finite.
\end{lemma}
\begin{proof}
\leanok
The density of the top power of a real two-form \(A\) on \(\mathbf C^k\) is, up to
the factor \(2^{-k}\), the alternating sum of products of the values of \(A\) on
the standard real frame, hence a polynomial in those values and so a continuous
function of \(A\). The parametrization \(\gamma\) is the restriction to a
coordinate plane of the inverse of a biholomorphic flattening, so it is smooth on
the parameter domain and \(w\mapsto D\gamma(w)\) is continuous there; the
coefficient matrix of \(a\) in the chart is smooth, hence \(w\mapsto a_{\gamma(w)}\)
is continuous. Composition of a two-form with a continuously varying linear map is
continuous, so \(\mathcal D_a(\gamma,\cdot)\) is continuous. A continuous function
on a compact set is bounded, and the compact set has finite Lebesgue measure, so
the integral is finite.
\end{proof}

\begin{lemma}[Inner regularity of a parameter integral]
\label{ana:inner-regular-parameter-integral}
\lean{TranscendentalNakamaye.Pending.RegularPieceInnerRegular.exists_isCompact_subset_lt_setLIntegral}
\leanok
Let \(\mu\) be a locally finite Borel measure on a proper metric space \(E\)
which is inner regular with respect to compact sets on sets of finite measure,
let \(S\subset E\) be measurable and let \(f\ge0\) be \(\mu|_S\)-measurable.
Then
\[
 \int_S f\,d\mu=\sup\Bigl\{\int_K f\,d\mu\;:\;K\subset S\ \text{compact}\Bigr\}.
\]
\end{lemma}
\begin{proof}
\leanok
Only ``\(\ge\)'' needs proof. Fix \(c<\int_Sf\,d\mu\). Put
\(A_m=S\cap \overline B(x_0,m)\) and \(f_m=\min(f,m)\mathbf 1_{A_m}\) for a fixed
\(x_0\in E\). The sequence \(f_m\) increases pointwise to \(f\mathbf 1_S\), since
every point lies in \(\overline B(x_0,m)\) for large \(m\) and
\(\min(a,m)\uparrow a\) for every \(a\in[0,\infty]\). Monotone convergence gives
\(\int_S f\,d\mu=\sup_m\int_{A_m}\min(f,m)\,d\mu\), so there is \(m\) with
\(c<t:=\int_{A_m}\min(f,m)\,d\mu\). The set \(A_m\) has finite measure, being
contained in a compact set, and \(t\le m\,\mu(A_m)<\infty\); choose \(r>0\) with
\(c+r<t\) and set \(\delta=r/(m+1)>0\).

Inner regularity of \(\mu\) supplies a compact \(K\subset A_m\) with
\(\mu(A_m\setminus K)<\delta\). Since \(A_m\subset K\cup(A_m\setminus K)\),
\[
 t\le\int_K\min(f,m)\,d\mu+\int_{A_m\setminus K}\min(f,m)\,d\mu
  \le\int_Kf\,d\mu+m\,\mu(A_m\setminus K)\le\int_Kf\,d\mu+r .
\]
Hence \(c+r<\int_Kf\,d\mu+r\) and, \(r\) being finite, \(c<\int_Kf\,d\mu\).
\end{proof}

\begin{lemma}[Compact regular pieces carry all the mass]
\label{ana:piece-family-compact-exhaustion}
\uses{ana:restricted-volume,ana:piece-mass-comparison-smooth,ana:inner-regular-parameter-integral,ana:patch-density-continuity}
\lean{TranscendentalNakamaye.Pending.RegularPieceInnerRegular.formPieceMass_eq_iSup_trimFamily,
TranscendentalNakamaye.Pending.RegularPieceInnerRegular.exists_compactPieces_lt_formPieceMass,
TranscendentalNakamaye.Pending.RegularPieceInnerRegular.iSup_formPieceMass_eq_iSup_compactPieces}
\leanok
Let \(a\) be a smooth real two-form on \(X\), let \(V\) be a closed irreducible
analytic subvariety, let \(A\subset X\) and let \(F=(\gamma_i,D_i)_{i<N}\) be a
regular-piece family avoiding \(A\). Then
\[
 \mathcal M_a(F)=\sup\bigl\{\mathcal M_a(F')\bigr\},
\]
the supremum running over the families \(F'=(\gamma_i,D_i')_{i<N}\) with the same
parametrizations whose pieces \(D_i'\subset D_i\) are compact. Consequently
\(\sup_F\mathcal M_a(F)\) is unchanged if only families with compact pieces are
allowed.
\end{lemma}
\begin{proof}
\uses{ana:inner-regular-parameter-integral,ana:patch-density-continuity}
\leanok
A family with pieces \(D_i'\subset D_i\) is again a regular-piece family avoiding
\(A\): the defining conditions --- measurability, inclusion in the parameter
domain, image inside \(V_{\mathrm{reg}}\), disjointness from \(A\) and pairwise
disjointness of the images --- are all inherited by subsets, and a compact subset
of a metric space is measurable. So ``\(\ge\)'' holds termwise.

For ``\(\le\)'', the density \(w\mapsto\mathcal D_a(\gamma_i,w)\) is continuous on
the parameter domain, hence measurable on \(D_i\), so
Lemma~\ref{ana:inner-regular-parameter-integral} applies to each of the \(N\)
parameter integrals separately and gives
\[
 \int_{D_i}\mathcal D_a(\gamma_i,\cdot)
 =\sup_{K\subset D_i\ \mathrm{compact}}\int_K\mathcal D_a(\gamma_i,\cdot).
\]
The \(N\) suprema may be exchanged with the finite sum because the family of
choices is directed: given two choices \((K_i)_i\) and \((K_i')_i\), the choice
\((K_i\cup K_i')_i\) is again admissible --- a union of two compact subsets of
\(D_i\) is a compact subset of \(D_i\) --- and dominates both termwise. Hence
\(\sum_i\sup_{K_i}\int_{K_i}=\sup_{(K_i)_i}\sum_i\int_{K_i}\), which is the
asserted supremum.
\end{proof}

\begin{lemma}[Exhaustion of \(V_{\mathrm{reg}}\) by compact regular pieces]
\label{ana:piece-family-exhaustion}
\uses{ana:piece-family-compact-exhaustion,ana:piece-family-compact-cover,ana:piece-mass-comparison-smooth}
\lean{TranscendentalNakamaye.Pending.RegularPieceInnerRegular.exists_isCompact_carrying_mass}
\leanok
Let \(a\) be a smooth real two-form on \(X\), let \(A\subset X\) be closed and let
\(F\) be a regular-piece family of \(V\) avoiding \(A\). For every
\(c<\mathcal M_a(F)\) there is a compact set \(K\subset V_{\mathrm{reg}}\) with
\(K\cap A=\emptyset\) such that some regular-piece family avoiding \(A\) covers
\(K\), and every regular-piece family \(G\) avoiding \(A\) whose images cover
\(K\) satisfies \(c<\mathcal M_a(G)\).
\end{lemma}
\begin{proof}
\uses{ana:piece-family-compact-exhaustion,ana:piece-family-compact-cover,ana:piece-mass-comparison-smooth}
\leanok
By Lemma~\ref{ana:piece-family-compact-exhaustion} there is a family
\(F'=(\gamma_i,D_i')_{i<N}\) with compact pieces \(D_i'\subset D_i\) and
\(c<\mathcal M_a(F')\). Put \(K=\bigcup_i\gamma_i(D_i')\). Each
\(\gamma_i\) is continuous on its parameter domain, so \(\gamma_i(D_i')\) is
compact and \(K\), a finite union of compact sets, is compact; it lies in
\(V_{\mathrm{reg}}\) and misses \(A\) because \(F'\) is a regular-piece family
avoiding \(A\).

A family covering \(K\) exists by
Lemma~\ref{ana:piece-family-compact-cover}. If \(G\) avoids \(A\) and its images
cover \(K\), then in particular every \(\gamma_i(D_i')\) is contained in the
union of the images of \(G\), so
\(\mathcal M_a(F')\le\mathcal M_a(G)\) by
Lemma~\ref{ana:piece-mass-comparison-smooth}, whence \(c<\mathcal M_a(G)\).
\end{proof}

\begin{lemma}[A parameter ball of regular points, trapped by an ambient open set]
\label{ana:regular-ball-ambient-trace}
\uses{ana:restricted-volume,ana:patch-transition}
\lean{TranscendentalNakamaye.Pending.RegularPieceExhaustion.exists_localPatchData}
\leanok
Let \(V\) be a closed irreducible analytic subvariety of dimension \(k\) of the
complex manifold \(X\), let \(A\subset X\) be closed, and let \(x\in
V_{\mathrm{reg}}\setminus A\). Then there are a regular parametrization
\(\gamma\) of \(V\) centred at \(x\), an open parameter ball \(\Omega\) inside its
parameter domain with \(\gamma(\Omega)\subset V_{\mathrm{reg}}\setminus A\), and
an OPEN set \(U\subset X\) with
\[
 x\in U,\qquad \gamma(\Omega)\subset U,\qquad V\cap U\subset\gamma(\Omega).
\]
\end{lemma}
\begin{proof}
\uses{ana:restricted-volume,ana:patch-transition}
\leanok
Regularity of \(V\) at \(x\) of dimension \(k\) is by definition the data of a
local biholomorphism flattening \(V\) near \(x\); it is exactly the data of a
regular parametrization \(\gamma\) centred at \(x\), so take that one, and then
\(\gamma(0)=x\). The parameter domain of \(\gamma\) is open and contains \(0\),
and every point it parametrizes is a regular point of \(V\) of dimension \(k\);
so some ball \(B(0,r_0)\) lies in the parameter domain and has
\(\gamma(B(0,r_0))\subset V_{\mathrm{reg}}\). The complement of \(A\) is open and
contains \(\gamma(0)=x\), and \(\gamma\) is continuous at \(0\), so
\(\gamma(B(0,r_1))\cap A=\emptyset\) for some \(r_1>0\). Put
\(\Omega=B(0,\min(r_0,r_1))\).

For \(U\): a regular parametrization maps an open set of parameters onto a
relatively open part of \(V\), i.e.\ there is an open \(U\subset X\) with
\(\gamma(\Omega)\subset U\) and \(V\cap U\subset\gamma(\Omega)\). Finally
\(x=\gamma(0)\in\gamma(\Omega)\subset U\) because \(0\in\Omega\).
\end{proof}

\begin{lemma}[One monotone sequence of families exhausts \(V_{\mathrm{reg}}\)]
\label{ana:piece-family-monotone-exhaustion}
\uses{ana:regular-ball-ambient-trace,ana:piece-family-compact-cover,ana:piece-mass-comparison-smooth}
\lean{TranscendentalNakamaye.Pending.RegularPieceExhaustion.coveredRegion,
TranscendentalNakamaye.Pending.RegularPieceExhaustion.runningUnion,
TranscendentalNakamaye.Pending.RegularPieceExhaustion.exists_exhaustion}
\leanok
Let \(V\) be a closed irreducible analytic subvariety of the complex manifold
\(X\) and let \(A\subset X\) be closed. There is a sequence
\((H_m)_{m\in\mathbf N}\) of regular-piece families of \(V\) avoiding \(A\) whose
covered regions \(R_m=\bigcup_i\gamma^{H_m}_i(D^{H_m}_i)\) increase,
\(R_m\subset R_{m+1}\), and which absorb every compact part of
\(V_{\mathrm{reg}}\setminus A\): for every compact \(K\subset
V_{\mathrm{reg}}\) with \(K\cap A=\emptyset\) there is an \(m\) with \(K\subset
R_m\).

The sequence depends on \(V\) and \(A\) only --- on no form, no class and no
current.
\end{lemma}
\begin{proof}
\uses{ana:regular-ball-ambient-trace,ana:piece-family-compact-cover,ana:piece-mass-comparison-smooth}
\leanok
Write \(S=V_{\mathrm{reg}}\setminus A\). If \(S=\emptyset\) take every \(H_m\)
empty: a compact \(K\subset V_{\mathrm{reg}}\) missing \(A\) lies in \(S\), hence
is empty. So assume \(S\neq\emptyset\).

By Lemma~\ref{ana:regular-ball-ambient-trace} each \(x\in S\) has a
parametrization \(\gamma_x\), a parameter ball \(\Omega_x\) with
\(\gamma_x(\Omega_x)\subset S\), and an open \(U_x\subset X\) with \(x\in U_x\)
and \(V\cap U_x\subset\gamma_x(\Omega_x)\). The \(U_x\) are open in \(X\) and
cover \(S\). A complex manifold is second countable, so a family of open sets has
the same union as a countable subfamily; \(S\) being nonempty that subfamily is
nonempty, and we may enumerate it as \(U_{x_0},U_{x_1},\dots\). Write
\(\gamma_m=\gamma_{x_m}\) and \(\Omega_m=\Omega_{x_m}\).

Each \(\Omega_m\) is open, hence measurable, lies in the parameter domain of
\(\gamma_m\) and has image missing \(A\); so by the primitive form of
Lemma~\ref{ana:piece-family-compact-cover} the first \(N+1\) of them are all
contained in the covered region of a single regular-piece family \(G_N\) avoiding
\(A\). Put \(H_0=G_0\) and \(H_{N+1}=H_N\sqcup G_{N+1}\), the union of families of
Lemma~\ref{ana:piece-mass-comparison-smooth}, whose covered region contains those
of both. Then \(R_N\subset R_{N+1}\), and \(R_N\) contains the covered region of
\(G_N\), hence \(\gamma_m(\Omega_m)\) for every \(m\le N\).

Let \(K\subset V_{\mathrm{reg}}\) be compact with \(K\cap A=\emptyset\); then
\(K\subset S\subset\bigcup_m U_{x_m}\), so finitely many \(U_{x_m}\), say with
indices \(m\le N\), already cover \(K\). For \(y\in K\) pick such an \(m\): since
\(y\in V\) and \(y\in U_{x_m}\), we get \(y\in V\cap U_{x_m}\subset
\gamma_m(\Omega_m)\subset R_N\). Hence \(K\subset R_N\).
\end{proof}

\begin{lemma}[The exhaustion computes the regular-piece supremum of every form]
\label{ana:piece-family-exhaustion-supremum}
\uses{ana:piece-family-monotone-exhaustion,ana:piece-family-exhaustion}
\lean{TranscendentalNakamaye.Pending.RegularPieceExhaustion.iSup_formPieceMass_eq_iSup_of_exhaustion,
TranscendentalNakamaye.Pending.RegularPieceExhaustion.exists_exhaustion_iSup_formPieceMass,
TranscendentalNakamaye.Pending.RegularPieceExhaustion.monotone_formPieceMass,
TranscendentalNakamaye.Pending.RegularPieceExhaustion.tendsto_formPieceMass_of_exhaustion}
\leanok
Let \(A\subset X\) be closed and let \((H_m)_m\) be a sequence of regular-piece
families of \(V\) avoiding \(A\) that absorbs every compact part of
\(V_{\mathrm{reg}}\setminus A\), as in
Lemma~\ref{ana:piece-family-monotone-exhaustion}. Then for EVERY smooth real
two-form \(a\) on \(X\),
\[
 \sup_m\mathcal M_a(H_m)=\sup_F\mathcal M_a(F),
\]
the right-hand supremum running over all regular-piece families of \(V\) avoiding
\(A\).

If moreover the covered regions increase, \(R_m\subset R_{m+1}\), then
\(m\mapsto\mathcal M_a(H_m)\) is nondecreasing and
\(\mathcal M_a(H_m)\to\sup_F\mathcal M_a(F)\).

In particular there is one monotone sequence \((H_m)_m\), chosen before any form
is named, which computes the regular-piece supremum of every smooth real two-form
at once --- and therefore of two different forms simultaneously, as the limits of
two convergent sequences indexed by the same \(m\).
\end{lemma}
\begin{proof}
\uses{ana:piece-family-monotone-exhaustion,ana:piece-family-exhaustion}
\leanok
``\(\le\)'' is immediate, each \(H_m\) being one of the families \(F\). For
``\(\ge\)'' it suffices to show \(c<\sup_m\mathcal M_a(H_m)\) for every family
\(F\) and every \(c<\mathcal M_a(F)\). By
Lemma~\ref{ana:piece-family-exhaustion} there is a compact \(K\subset
V_{\mathrm{reg}}\) with \(K\cap A=\emptyset\) such that every regular-piece family
avoiding \(A\) whose covered region contains \(K\) has mass \(>c\). By hypothesis
\(K\subset R_m\) for some \(m\), so \(c<\mathcal M_a(H_m)\le\sup_m\mathcal
M_a(H_m)\).

If the covered regions increase then every image of \(H_m\) lies in \(R_m\subset
R_{m+1}\), so \(\mathcal M_a(H_m)\le\mathcal M_a(H_{m+1})\) by
Lemma~\ref{ana:piece-mass-comparison-smooth}; a nondecreasing sequence in
\([0,\infty]\) converges to its supremum, which is \(\sup_F\mathcal M_a(F)\) by
the above.

The last assertion is this statement applied to the sequence of
Lemma~\ref{ana:piece-family-monotone-exhaustion}, which is produced without
reference to \(a\).
\end{proof}

\begin{definition}[A local piece property]
\label{ana:local-piece-property}
\uses{ana:background-mass}
\lean{TranscendentalNakamaye.Pending.RegularPieceGoodExhaustion.LocalPieceProperty,
TranscendentalNakamaye.Pending.RegularPieceGoodExhaustion.IsGoodFamily}
\leanok
A \emph{local piece property} of \(V\) assigns to every regular patch
\(P=(\gamma,\Omega_P)\) of \(V\) a collection \(\mathcal G(P)\) of subsets of
\(\Omega_P\) such that
\begin{enumerate}
 \item[(i)] \(\mathcal G(P)\) is closed under passing to subsets: if
   \(S\in\mathcal G(P)\) and \(S'\subseteq S\) then \(S'\in\mathcal G(P)\);
 \item[(ii)] every parameter \(w\in\Omega_P\) has an open neighbourhood
   \(\Omega\subseteq\Omega_P\) with \(\Omega\in\mathcal G(P)\).
\end{enumerate}
A regular-piece family \(F\) is \emph{good} for \(\mathcal G\) when every piece
\(D_i\) of \(F\) satisfies \(D_i\in\mathcal G(P_i)\), \(P_i\) being the patch of
that piece.
\end{definition}

\begin{lemma}[The exhaustion may be taken good for any local piece property]
\label{ana:piece-family-good-exhaustion}
\uses{ana:local-piece-property,ana:regular-ball-ambient-trace,ana:piece-family-compact-cover,ana:piece-mass-comparison-smooth,ana:piece-family-monotone-exhaustion}
\lean{TranscendentalNakamaye.Pending.RegularPieceGoodExhaustion.exists_exhaustion_good,
TranscendentalNakamaye.Pending.RegularPieceGoodExhaustion.exists_goodFamily_iUnion_image_superset,
TranscendentalNakamaye.Pending.RegularPieceGoodExhaustion.exists_localPatchData_good}
\leanok
Let \(A\subset X\) be closed and let \(\mathcal G\) be a local piece property of
\(V\). Then the sequence of Lemma~\ref{ana:piece-family-monotone-exhaustion} can be
chosen good for \(\mathcal G\): there is a sequence \((H_m)_m\) of regular-piece
families of \(V\) avoiding \(A\) whose covered regions increase, which absorbs every
compact part of \(V_{\mathrm{reg}}\setminus A\), and every piece of every \(H_m\) of
which lies in \(\mathcal G\) of its patch.

Consequently, if two smooth real two-forms \(a,b\) on \(X\) satisfy
\(\sup_m\mathcal M_a(H_m)=\sup_m\mathcal M_b(H_m)\) for one such sequence, then
\(\sup_F\mathcal M_a(F)=\sup_F\mathcal M_b(F)\), the suprema running over all
regular-piece families of \(V\) avoiding \(A\).
\end{lemma}
\begin{proof}
\uses{ana:local-piece-property,ana:regular-ball-ambient-trace,ana:piece-family-compact-cover,ana:piece-mass-comparison-smooth,ana:piece-family-exhaustion-supremum}
\leanok
First refine the local step. Let \(x\in V_{\mathrm{reg}}\setminus A\) and take
\(\gamma_x\), \(\Omega_x\) and \(U_x\) as in
Lemma~\ref{ana:regular-ball-ambient-trace}. Since \(x\in V\cap
U_x\subset\gamma_x(\Omega_x)\) there is a parameter \(w_x\in\Omega_x\) with
\(\gamma_x(w_x)=x\). By clause (ii) of Definition~\ref{ana:local-piece-property}
there is an open \(\Omega'\subseteq\Omega_{\gamma_x}\) containing \(w_x\) with
\(\Omega'\in\mathcal G(\gamma_x)\); replace \(\Omega_x\) by
\(\Omega_x\cap\Omega'\), which is again open, still contains \(w_x\), still has
image in \(V_{\mathrm{reg}}\setminus A\), and lies in \(\mathcal G(\gamma_x)\) by
clause (i). Applying the second half of Lemma~\ref{ana:regular-ball-ambient-trace}
to the smaller ball produces a new open \(U_x'\supseteq\gamma_x(\Omega_x\cap
\Omega')\) with \(V\cap U_x'\subset\gamma_x(\Omega_x\cap\Omega')\), and
\(x=\gamma_x(w_x)\in U_x'\). So the local step may be carried out with a good
parameter set.

Goodness now survives the two constructions by which the families of
Lemma~\ref{ana:piece-family-monotone-exhaustion} are assembled. A single
measurable parameter set, read as a one-piece family, keeps that set as its only
piece, so it is good as soon as the set is. The union \(F\sqcup G\) of
Lemma~\ref{ana:piece-mass-comparison-smooth} keeps every piece of \(F\) and
replaces every piece \(D_j\) of \(G\) by the subset of \(D_j\) whose image no piece
of \(F\) already sees; both are then in \(\mathcal G\) of their (unchanged) patch,
the first by hypothesis and the second by clause (i). By induction the finite
unions \(G_N\) of the good parameter balls, and the running unions
\(H_N\), are good.

The proof of Lemma~\ref{ana:piece-family-monotone-exhaustion} is otherwise
unchanged: the countable subfamily, the finite subcover of a compact \(K\), and the
monotonicity of the covered regions use nothing about the parameter sets beyond
openness. This gives the sequence.

For the consequence, Lemma~\ref{ana:piece-family-exhaustion-supremum} applied to
that one sequence with \(a\) and with \(b\) rewrites both suprema over all families
as suprema along \((H_m)_m\), and the hypothesis is the rewritten identity.
\end{proof}

\begin{lemma}[Locally finite volume at a regular point]
\label{ana:background-mass-local-regular}
\uses{ana:background-mass,ana:patch-transition}
\lean{TranscendentalNakamaye.Pending.SubvarietyStokesMassInterfaces.exists_localBackgroundBound_of_mem_regularLocus}
\leanok
Let \(x\) be a regular point of \(V\) of dimension \(k\). Then there are a
neighbourhood \(U\subset X\) of \(x\) and a constant \(C<\infty\) such that every
finite family of disjoint measurable pieces of regular parametrizations of \(V\)
whose images lie in \(U\) has total background parameter mass at most \(C\).
\end{lemma}
\begin{proof}
\uses{ana:background-mass,ana:patch-transition}
\leanok
Let \(P\) be the regular parametrization of \(V\) at \(x\) given by the
regularity data, \(B\) a closed parameter ball inside its parameter domain, and
\(B^\circ\) its interior. Since \(P\) maps open sets of parameters onto
relatively open parts of \(V\), there is an open \(U\subset X\) with
\(P(B^\circ)\subset U\) and \(V\cap U\subset P(B^\circ)\); in particular
\(x=P(0)\in U\).

Let \((D_i,Q_i)\) be a finite family of disjoint measurable pieces with images in
\(U\). Each image point lies in \(V\cap U\subset P(B^\circ)\), so by
Lemma~\ref{ana:patch-transition} every \(D_i\) lies in the transition domain from
\(Q_i\) to \(P\), and \(\tau_i(D_i)\subset B^\circ\). The sets \(\tau_i(D_i)\) are
measurable, being images of measurable sets under injective maps differentiable
there, and pairwise disjoint, because \(P\circ\tau_i=Q_i\) and the original images
are disjoint. Lemma~\ref{ana:patch-transition} also gives
\(\int_{D_i}\rho_{\beta,Q_i}=\int_{\tau_i(D_i)}\rho_{\beta,P}\), so the total mass
is \(\int_{\bigsqcup_i\tau_i(D_i)}\rho_{\beta,P}\le\int_{B}\rho_{\beta,P}=:C\).
The density \(\rho_{\beta,P}\) is continuous on the parameter domain, hence bounded
on the compact ball \(B\), so \(C<\infty\); and \(C\) depends only on \(x\), not on
the family.
\end{proof}

\begin{lemma}[Upper comparison for the genuine top-power density]
\label{ana:top-power-upper-comparison}
\uses{fpma:top-power-determinant}
\lean{TranscendentalNakamaye.Pending.TopPowerDensityUpperBound.actualTopPowerDensity_le_pow_mul,
TranscendentalNakamaye.Pending.TopPowerDensityUpperBound.apply_I_smul_self_le_norm_mul,
TranscendentalNakamaye.Pending.TopPowerDensityUpperBound.actualTopPowerDensity_comp_le_pow_mul}
\leanok
Let \(A,B\) be real \((1,1)\)-forms on \(\mathbf C^{k}\) with \(A\) semipositive,
and suppose \(A(v,iv)\le c\,B(v,iv)\) for every \(v\). Then
\(A^{\wedge k}\le c^{k}B^{\wedge k}\) on the standard real frame. Moreover every
real two-form satisfies \(A(v,iv)\le\|A\|\,\|v\|^{2}\); and if
\(L:\mathbf C^{k}\to\mathbf C^{n}\) is complex linear and \(A\) is a semipositive
\((1,1)\)-form on \(\mathbf C^{n}\) with \(A(z,iz)\le c\|z\|^{2}\), then
\((L^{\star}A)^{\wedge k}\le c^{k}(L^{\star}\omega)^{\wedge k}\), where \(\omega\)
is the standard Hermitian form of \(\mathbf C^{n}\).
\end{lemma}
\begin{proof}
\uses{fpma:top-power-determinant}
\leanok
The difference \(D=cB-A\) is again a real \((1,1)\)-form and is semipositive by
hypothesis, so monotonicity of the top exterior power on semipositive
\((1,1)\)-forms gives \(A^{\wedge k}\le(A+D)^{\wedge k}=(cB)^{\wedge k}
=c^{k}B^{\wedge k}\), the last step because the \(k\)-fold wedge of \(cB\)
scales by \(c^{k}\).

For the norm bound, \(A(v,iv)\le|A(v,iv)|\le\|A\|\,\|v\|\,\|iv\|=\|A\|\,\|v\|^{2}\).

For the pullback, \(L\) commutes with the complex structure, so \(L^{\star}A\)
and \(L^{\star}\omega\) are again real \((1,1)\)-forms, \(L^{\star}A\) is
semipositive, and for \(v\in\mathbf C^{k}\)
\[
 (L^{\star}A)(v,iv)=A(Lv,iLv)\le c\|Lv\|^{2}=c\,(L^{\star}\omega)(v,iv).
\]
The first assertion applies.
\end{proof}

\begin{lemma}[The parameter density in one fixed chart, and the flat density]
\label{ana:flat-patch-density}
\uses{ana:patch-transition,ana:top-power-upper-comparison}
\lean{TranscendentalNakamaye.Pending.FlatPatchDensityComparison.chartDomain,
TranscendentalNakamaye.Pending.FlatPatchDensityComparison.chartPullbackDerivative,
TranscendentalNakamaye.Pending.FlatPatchDensityComparison.patchCoordinateForm_eq_chart,
TranscendentalNakamaye.Pending.FlatPatchDensityComparison.flatPatchDensity,
TranscendentalNakamaye.Pending.FlatPatchDensityComparison.actualTopPowerDensity_patchCoordinateForm_le,
TranscendentalNakamaye.Pending.FlatPatchDensityComparison.actualTopPowerDensity_patchCoordinateForm_le_norm}
\leanok
Fix a chart \(e\) of \(X\) and a regular parametrization \(\gamma\) of \(V\) of
dimension \(k\). At a parameter \(w\) whose image lies in the source of \(e\),
write \(L_{e,\gamma}(w)=D(e\circ\gamma)_w:\mathbf C^{k}\to\mathbf C^{n}\).

For every smooth real two-form \(a\) on \(X\), the pullback of \(a\) along
\(\gamma\) read in parameters equals \(L_{e,\gamma}(w)^{\star}\) of the
coefficient of \(a\) in the chart \(e\) at \(e(\gamma(w))\): the chart at the
centre of \(\gamma\) does not appear. Writing
\(\rho^{\flat}_{e,\gamma}(w)=(L_{e,\gamma}(w)^{\star}\omega)^{\wedge k}\) for the
flat parameter density of \(e\), if moreover \(a\) is a semipositive
\((1,1)\)-form whose coefficient at \(e(\gamma(w))\) satisfies
\(a(z,iz)\le c\|z\|^{2}\), then \(\rho_{a,\gamma}(w)\le
c^{k}\rho^{\flat}_{e,\gamma}(w)\); in particular this holds with
\(c=\|a_{e}(e(\gamma(w)))\|\).
\end{lemma}
\begin{proof}
\uses{ana:patch-transition,ana:top-power-upper-comparison}
\leanok
Let \(e_\gamma\) be the chart at the centre of \(\gamma\) and \(T=e\circ
e_\gamma^{-1}\) the transition, defined and holomorphic at the \(e_\gamma\)
coordinates of \(\gamma(w)\). By the chain rule
\(L_{e,\gamma}(w)=DT\circ D(e_\gamma\circ\gamma)_w\), and the tensor
transformation law of a smooth real two-form says exactly that its coefficient in
\(e_\gamma\) is the coefficient in \(e\) precomposed with \(DT\). Composing the
two precompositions gives the first assertion. Both \(DT\) and
\(D(e_\gamma\circ\gamma)_w\) are complex linear, hence so is
\(L_{e,\gamma}(w)\), and the comparison is
Lemma~\ref{ana:top-power-upper-comparison} applied to
\(L=L_{e,\gamma}(w)\). The last clause is the same lemma's norm bound.
\end{proof}

\begin{lemma}[The flat density is a Gram determinant]
\label{ana:flat-density-gram}
\uses{fpma:top-power-determinant,ana:flat-patch-density}
\lean{TranscendentalNakamaye.Pending.FlatDensityGramDeterminant.rectGram,
TranscendentalNakamaye.Pending.FlatDensityGramDeterminant.normalizedHermitianMatrix_standard_comp,
TranscendentalNakamaye.Pending.FlatDensityGramDeterminant.actualTopPowerDensity_standard_comp_eq_gram,
TranscendentalNakamaye.Pending.FlatDensityGramDeterminant.rectGram_eq,
TranscendentalNakamaye.Pending.FlatDensityGramDeterminant.flatPatchDensity_eq_gram}
\leanok
Let \(L:\mathbf C^{k}\to\mathbf C^{n}\) be complex linear and let \(\omega\) be the
standard Hermitian form of \(\mathbf C^{n}\). Then the genuine top exterior power
of \(L^{\star}\omega\), evaluated on the standard real frame of \(\mathbf C^{k}\),
is
\[
 (L^{\star}\omega)^{\wedge k}=k!\,\operatorname{Re}\det G(L),\qquad
 G(L)_{ij}=\langle Le_i,Le_j\rangle,
\]
and \(G(L)=L^{\mathsf H}L\) as matrices. In particular the flat parameter density
\(\rho^{\flat}_{e,\gamma}(w)\) of Lemma~\ref{ana:flat-patch-density} is
\(k!\,\operatorname{Re}\det\bigl(L^{\mathsf H}L\bigr)\) for
\(L=D(e\circ\gamma)_w\).
\end{lemma}
\begin{proof}
\uses{fpma:top-power-determinant,ana:flat-patch-density}
\leanok
Because \(L\) commutes with the complex structure, \(L^{\star}\omega\) is a real
\((1,1)\)-form, and it is semipositive since
\((L^{\star}\omega)(v,iv)=\omega(Lv,iLv)=\|Lv\|^{2}\ge0\). So
Lemma~\ref{fpma:top-power-determinant} applies:
\((L^{\star}\omega)^{\wedge k}=k!\pi^{-k}\det H\) with
\(H_{ij}=\pi\bigl((L^{\star}\omega)(e_i,ie_j)-i\,(L^{\star}\omega)(e_i,e_j)\bigr)\).

Now \(\omega(u,v)=\operatorname{Im}\langle u,v\rangle\), so, writing
\(z=\langle Le_i,Le_j\rangle\),
\[
 (L^{\star}\omega)(e_i,ie_j)=\operatorname{Im}\langle Le_i,iLe_j\rangle
 =\operatorname{Im}(iz)=\operatorname{Re}z,\qquad
 (L^{\star}\omega)(e_i,e_j)=\operatorname{Im}z,
\]
whence \(H_{ij}=\pi(\operatorname{Re}z-i\operatorname{Im}z)=\pi\overline z\), that
is \(H=\pi\,\overline{G(L)}\). Therefore
\(\det H=\pi^{k}\overline{\det G(L)}\) and
\(k!\pi^{-k}\det H=k!\,\overline{\det G(L)}\), whose real part is
\(k!\operatorname{Re}\det G(L)\).

Finally, expanding the Hermitian inner product of \(\mathbf C^{n}\) in
coordinates, \(G(L)_{ij}=\sum_{p}\overline{L_{pi}}L_{pj}=(L^{\mathsf H}L)_{ij}\).
\end{proof}

\begin{lemma}[Cauchy--Binet]
\label{ana:cauchy-binet}
\lean{TranscendentalNakamaye.Pending.CauchyBinet.det_mul_rect,
TranscendentalNakamaye.Pending.CauchyBinet.det_conjTranspose_mul_self,
TranscendentalNakamaye.Pending.CauchyBinet.re_det_conjTranspose_mul_self}
\leanok
Let \(R\) be a commutative ring, \(A\) a \(k\times n\) matrix and \(B\) an
\(n\times k\) matrix over \(R\). For a \(k\)-element subset
\(I\subset\{1,\dots,n\}\) write \(A_{\bullet I}\) for the \(k\times k\) matrix of
the columns of \(A\) indexed by \(I\), and \(B_{I\bullet}\) for the \(k\times k\)
matrix of the rows of \(B\) indexed by \(I\), both listed in increasing order.
Then
\[
 \det(AB)=\sum_{|I|=k}\det(A_{\bullet I})\,\det(B_{I\bullet}).
\]
In particular, for a complex \(n\times k\) matrix \(A\),
\(\det(A^{\mathsf H}A)=\sum_{|I|=k}\lvert\det A_{I\bullet}\rvert^{2}\); it is a
nonnegative real number.
\end{lemma}
\begin{proof}
\leanok
By the Leibniz formula and bilinearity,
\[
 \det(AB)=\sum_{\sigma\in\mathfrak S_k}\varepsilon(\sigma)
 \prod_{i=1}^{k}\sum_{j=1}^{n}A_{\sigma(i)j}B_{ji}
 =\sum_{p:\{1,\dots,k\}\to\{1,\dots,n\}}
 \Bigl(\sum_{\sigma}\varepsilon(\sigma)\prod_i A_{\sigma(i)p(i)}\Bigr)
 \prod_i B_{p(i)i},
\]
the middle step expanding the product of the \(k\) sums over all maps \(p\). The
inner sum is \(\det(A_{\bullet p})\), the determinant of the \(k\times k\) matrix
whose \(i\)-th column is the \(p(i)\)-th column of \(A\). If \(p\) is not
injective that matrix has two equal columns, so the term vanishes.

An injective \(p\) factors uniquely as \(p=\iota_I\circ\tau\) with
\(I=p(\{1,\dots,k\})\), \(\iota_I\) its increasing enumeration and
\(\tau\in\mathfrak S_k\); conversely every such pair gives an injective \(p\).
Permuting columns, \(\det(A_{\bullet p})=\varepsilon(\tau)\det(A_{\bullet I})\),
while \(\prod_i B_{p(i)i}=\prod_i (B_{I\bullet})_{\tau(i)i}\). Summing over
\(\tau\) at fixed \(I\),
\(\sum_{\tau}\varepsilon(\tau)\prod_i (B_{I\bullet})_{\tau(i)i}
=\det(B_{I\bullet})\), which gives the formula.

For the Hermitian case apply it with \(A^{\mathsf H}\) in place of \(A\) and
\(A\) in place of \(B\): the columns of \(A^{\mathsf H}\) indexed by \(I\) form
\((A_{I\bullet})^{\mathsf H}\), whose determinant is
\(\overline{\det A_{I\bullet}}\), so the \(I\)-term is
\(\overline{\det A_{I\bullet}}\det A_{I\bullet}
=\lvert\det A_{I\bullet}\rvert^{2}\).
\end{proof}

\begin{lemma}[The flat density is the sum of the squared coordinate Jacobians]
\label{ana:flat-density-minors}
\uses{ana:flat-density-gram,ana:cauchy-binet}
\lean{TranscendentalNakamaye.Pending.FlatDensityMinorExpansion.actualTopPowerDensity_standard_comp_eq_sum_minors,
TranscendentalNakamaye.Pending.FlatDensityMinorExpansion.flatPatchDensity_eq_sum_minors}
\leanok
\source{demailly-cadg:III.2.6}
For a complex linear \(L:\mathbf C^{k}\to\mathbf C^{n}\),
\[
 (L^{\star}\omega)^{\wedge k}=k!\sum_{|I|=k}\lvert\det L_{I}\rvert^{2},
\]
\(L_I\) being the \(k\times k\) minor of the matrix of \(L\) on the rows \(I\).
In particular the flat parameter density of a chart \(e\) along a regular
parametrization \(\gamma\) is
\(\rho^{\flat}_{e,\gamma}(w)=k!\sum_{|I|=k}\lvert\det L_I(w)\rvert^{2}\) with
\(L(w)=D(e\circ\gamma)_w\).
\end{lemma}
\begin{proof}
\uses{ana:flat-density-gram,ana:cauchy-binet}
\leanok
Lemma~\ref{ana:flat-density-gram} gives
\((L^{\star}\omega)^{\wedge k}=k!\operatorname{Re}\det(L^{\mathsf H}L)\), and
Lemma~\ref{ana:cauchy-binet} evaluates \(\det(L^{\mathsf H}L)\) as the real number
\(\sum_{|I|=k}\lvert\det L_I\rvert^{2}\).
\end{proof}

\begin{lemma}[Area formula with multiplicity]
\label{ana:multiplicity-area}
\lean{TranscendentalNakamaye.Pending.MultiplicityAreaBound.exists_disjointed_injOn,
TranscendentalNakamaye.Pending.MultiplicityAreaBound.lintegral_abs_det_eq_tsum_image,
TranscendentalNakamaye.Pending.MultiplicityAreaBound.sum_lintegral_abs_det_le}
\leanok
Let \(E\) be a finite-dimensional real normed space with a Haar measure \(\mu\),
let \(S_1,\dots,S_m\subset E\) be measurable, and let \(f_i\) be strictly
differentiable at every point of \(S_i\), with derivative \(Df_i(x)\). Let
\(\Delta\subset E\) be measurable with \(f_i(S_i)\subset\Delta\) for every \(i\),
and let \(N\in\mathbf N\) satisfy
\(\sum_i\#\bigl(S_i\cap f_i^{-1}(y)\bigr)\le N\) for every \(y\in\Delta\). Then
\[
 \sum_i\int_{S_i}\lvert\det Df_i\rvert\,d\mu\le N\,\mu(\Delta).
\]
\end{lemma}
\begin{proof}
\leanok
Fix \(i\). At a point \(x\in S_i\) with \(\det Df_i(x)\neq0\) the derivative is
invertible, so by the inverse function theorem \(f_i\) is injective on an open
neighbourhood \(U_x\) of \(x\). The open sets \(U_x\), \(x\) ranging over
\(\{x\in S_i:\det Df_i(x)\neq0\}\), cover that set; \(E\) is second countable, so
countably many of them, \(U_1,U_2,\dots\), have the same union. Put
\(D_{i,j}=S_i\cap\bigl(U_j\setminus\bigcup_{l<j}U_l\bigr)\). These are measurable,
pairwise disjoint in \(j\), contained in \(S_i\), \(f_i\) is injective on each,
and their union contains every point of \(S_i\) at which \(\det Df_i\neq0\).

The integrand vanishes at the remaining points of \(S_i\), so
\[
 \int_{S_i}\lvert\det Df_i\rvert\,d\mu
 =\sum_j\int_{D_{i,j}}\lvert\det Df_i\rvert\,d\mu
 =\sum_j\mu\bigl(f_i(D_{i,j})\bigr),
\]
the last equality being the change-of-variables formula for a differentiable map
injective on a measurable set; the image sets are measurable because a continuous
injective image of a Borel set is Borel.

Each \(f_i(D_{i,j})\) lies in \(\Delta\). For \(y\in\Delta\) the pairs \((i,j)\)
with \(y\in f_i(D_{i,j})\) are at most \(N\) in number: choosing a preimage in
each \(D_{i,j}\) gives, for each \(i\), distinct points of \(S_i\cap f_i^{-1}(y)\),
because the \(D_{i,j}\) are disjoint in \(j\). Integrating the resulting
inequality \(\sum_{i,j}\mathbf 1_{f_i(D_{i,j})}\le N\,\mathbf 1_{\Delta}\) gives
\(\sum_{i,j}\mu(f_i(D_{i,j}))\le N\mu(\Delta)\).
\end{proof}

\begin{lemma}[The projected parametrization and its Jacobian]
\label{ana:projected-parametrization}
\uses{ana:flat-patch-density,ana:flat-density-minors}
\lean{TranscendentalNakamaye.Pending.ProjectedPatchJacobian.hasStrictFDerivAt_projectedParametrization,
TranscendentalNakamaye.Pending.ProjectedPatchJacobian.det_projectedDerivative,
TranscendentalNakamaye.Pending.ProjectedPatchJacobian.chartPullbackDerivative_eq}
\leanok
Let \(e\) be a chart of \(X\), \(\gamma\) a regular parametrization of \(V\) of
dimension \(k\), and \(I\) a \(k\)-element set of coordinates of \(\mathbf C^{n}\)
with projection \(\pi_I:\mathbf C^{n}\to\mathbf C^{k}\). On the open set of
parameters \(w\) whose image \(\gamma(w)\) lies in the source of \(e\), the map
\(f_I=\pi_I\circ e\circ\gamma:\mathbf C^{k}\to\mathbf C^{k}\) is strictly
differentiable, and its real Jacobian determinant is
\[
 \det{}_{\mathbf R}\,D(f_I)_w=\lvert\det L_I(w)\rvert^{2},\qquad
 L(w)=D(e\circ\gamma)_w .
\]
\end{lemma}
\begin{proof}
\uses{ana:flat-patch-density,ana:flat-density-minors}
\leanok
The set in question is open, being the intersection of the open parameter domain
with the preimage of an open set under a continuous map. On it \(e\circ\gamma\) is
a composition of the parametrization and a chart transition, both holomorphic,
so \(f_I\) is holomorphic; a holomorphic map of several variables is \(C^{\infty}\)
as a map of real variables, hence strictly differentiable, and its real derivative
is the restriction of scalars of its complex derivative,
\(D(f_I)_w=\pi_I\circ L(w)\).

The real determinant of a complex linear endomorphism of \(\mathbf C^{k}\) is the
squared modulus of its complex determinant. The matrix of \(\pi_I\circ L(w)\) in
the standard bases has \((p,j)\) entry the \(I(p)\)-th coordinate of
\(L(w)e_j\), i.e. it is \(L_I(w)\); so
\(\det_{\mathbf R}D(f_I)_w=\lvert\det L_I(w)\rvert^{2}\).
\end{proof}

\begin{proposition}[Local flat volume from the sheet numbers]
\label{ana:flat-volume-fibre-bound}
\uses{ana:flat-density-minors,ana:multiplicity-area,ana:projected-parametrization}
\lean{TranscendentalNakamaye.Pending.FlatVolumeFibreBound.sum_lintegral_flatPatchDensity_le}
\leanok
\source{demailly-cadg:III.2.6}
Let \(e\) be a chart of \(X\) and let \(D_1,\dots,D_m\) be measurable pieces of
regular parametrizations \(\gamma_1,\dots,\gamma_m\) of \(V\), all inside the set
of parameters that \(e\) sees. Let \(\Delta\subset\mathbf C^{k}\) be measurable
and \(N\in\mathbf N\), and suppose that for every \(k\)-element set \(I\) of
coordinates the maps \(f_{I,i}=\pi_I\circ e\circ\gamma_i\) send \(D_i\) into
\(\Delta\) and satisfy
\(\sum_i\#\bigl(D_i\cap f_{I,i}^{-1}(y)\bigr)\le N\) for almost every \(y\in\Delta\).
Then
\[
 \sum_i\int_{D_i}\rho^{\flat}_{e,\gamma_i}
 \le k!\binom{n}{k}N\operatorname{vol}(\Delta).
\]
\end{proposition}
\begin{proof}
\uses{ana:flat-density-minors,ana:multiplicity-area,ana:projected-parametrization}
\leanok
By Lemma~\ref{ana:flat-density-minors} the integrand is
\(k!\sum_{|I|=k}\lvert\det L_{I,i}\rvert^{2}\), and by
Lemma~\ref{ana:projected-parametrization} the \(I\)-summand is the real Jacobian
determinant of \(f_{I,i}\), which is nonnegative. Interchanging the finite sums
over \(i\) and over \(I\),
\[
 \sum_i\int_{D_i}\rho^{\flat}_{e,\gamma_i}
 =k!\sum_{|I|=k}\ \sum_i\int_{D_i}\lvert\det D(f_{I,i})\rvert .
\]
Lemma~\ref{ana:projected-parametrization} also makes each \(f_{I,i}\) strictly
differentiable on \(D_i\), so Lemma~\ref{ana:multiplicity-area} bounds the inner
double sum by \(N\operatorname{vol}(\Delta)\) for each \(I\). There are
\(\binom{n}{k}\) sets \(I\).
\end{proof}

\begin{proposition}[Almost-everywhere fibre count of an adapted projection]
\label{ana:adapted-projection-ae-fibre}
\uses{lag:parametrization-sets,ana:holomorphic-zero-set-null}
\source{demailly-cadg:II.4.19}
\lean{TranscendentalNakamaye.Pending.ParametrizationFibreBound.exists_fibreBound}
\leanok
This is the ADAPTED projection of the parametrization data, not an arbitrary coordinate
projection of an arbitrary chart (for which generic rank \(d\) does not give a ramified covering; see Remark~\ref{ana:lelong-basis-form-route}).
Let \(A\) be an analytic germ at \(0\in\mathbf C^{n}\) with parametrization data of
dimension \(d\) and degree \(q\), and let \(\Delta=\Delta'_{\rho'}\times\Delta''_{\rho''}\)
be a parametrizing polydisc for it, with adapted projection \(\pi(z)=z'\). Shrink the base
radius to \(\rho=\min(\rho',\varrho)\), where \(\varrho\) is the radius of the ball on which
the discriminant \(\delta\) is holomorphic, and let
\(\Delta_\rho=\Delta'_\rho\times\Delta''_{\rho''}\). Then \(\Delta_\rho\) is an open
neighbourhood of \(0\), the ball \(\Delta'_\rho\) is measurable of finite volume and
contains \(\pi(A\cap\Delta_\rho)\), and for almost every \(z'\in\mathbf C^{d}\)
\[
  \#\bigl(A\cap\Delta_\rho\cap\pi^{-1}(z')\bigr)\le q .
\]
\end{proposition}
\begin{proof}
\uses{lag:parametrization-sets,ana:holomorphic-zero-set-null}
\leanok
The set of \(z'\) at which the bound can fail is contained in
\(\{z'\;:\;\lvert z'\rvert<\rho\ \text{and}\ \delta(z')=0\}\), and that set is Lebesgue-null.
Indeed \(\delta\) is holomorphic on the ball of radius \(\varrho\) and does not vanish
identically there --- otherwise it would vanish on a neighbourhood of \(0\), contrary to the
parametrization data --- so Lemma~\ref{ana:holomorphic-zero-set-null} applies, and
\(\rho\le\varrho\).

It remains to see the inclusion. Let \(z'\) satisfy the two negations. If
\(\lvert z'\rvert\ge\rho\) then no point of \(A\cap\Delta_\rho\) projects to \(z'\), because
every point of \(\Delta_\rho\) has \(\lvert\pi(z)\rvert<\rho\); so the fibre is empty and
the bound holds. If \(\lvert z'\rvert<\rho\) and \(\delta(z')\ne0\) then, since
\(\rho\le\rho'\), the fibre \(A\cap\Delta_\rho\cap\pi^{-1}(z')\) is contained in
\(A\cap\Delta\cap\pi^{-1}(z')\), which by Proposition~\ref{lag:parametrization-sets}(2) is
finite with exactly \(q\) points.
\end{proof}

\begin{remark}[A sheet bound with a free direction is empty]
\label{ana:free-direction-empty}
\uses{ana:restricted-volume}
Let \(A\) be an analytic germ at \(0\in\mathbf C^{n}\). The assertion
\begin{quote}
there are \(d\), a \(\mathbf C\)-linear \(B:\mathbf C^{n}\to\mathbf C^{d}\), a neighbourhood
\(U\) of \(0\), a measurable \(\Delta\) of finite volume containing \(B(A\cap U)\) and an
integer \(N\) with \(\#(A\cap U\cap B^{-1}(y))\le N\) for almost every \(y\)
\end{quote}
holds for EVERY \(A\), with no hypothesis at all. Take \(d=n\) and \(B\) the identity: every
fibre is a single point. Take instead any \(d\ge1\) and \(B=0\): the image of \(A\cap U\) is
\(\{0\}\), a Lebesgue-null set, so the fibre over almost every \(y\) is empty.

So all the content of the local parametrization theorem sits in the DIRECTION, and a
statement of this shape is informative only when the direction is exhibited --- through the
data that produces it, as in Proposition~\ref{ana:germ-sheet-bound} --- or prescribed up to
\(\delta\), as in Lemma~\ref{ana:adapted-projection-germ}, where a degenerate \(B\) is
excluded because the degenerate directions form a proper analytic subset of
\(\operatorname{Hom}(\mathbf C^{n},\mathbf C^{d})\), whose points are at positive distance
from a direction off it.
\end{remark}

\begin{proposition}[Sheet bound of an irreducible germ, in coordinates it chooses]
\label{ana:germ-sheet-bound}
\uses{lag:noether-normalization,lag:parametrization-sets,ana:adapted-projection-ae-fibre}
\lean{TranscendentalNakamaye.Pending.IrreducibleGermSheetBound.exists_germSheetBound}
\leanok
\source{demailly-cadg:II.4.19}
Let \(A\) be an irreducible analytic germ at \(0\in\mathbf C^{n}\). There are a linear change
of coordinates \(L\) of \(\mathbf C^{n}\) putting the vanishing ideal of \(A\) in Noether
position over \(\mathcal O_{d}\), with parametrization data for \(L^{-1}A\) of dimension
\(d\), an open neighbourhood \(U\) of \(0\), a measurable \(\Delta\subset\mathbf C^{d}\) of
finite volume and an integer \(N\), such that the transported adapted projection
\(B=\pi\circ L^{-1}\) satisfies \(B(A\cap U)\subset\Delta\) and
\[
  \#\bigl(A\cap U\cap B^{-1}(y)\bigr)\le N\qquad\text{for almost every }y\in\mathbf C^{d}.
\]
The parametrization data is part of the conclusion, not an auxiliary object:
Remark~\ref{ana:free-direction-empty} is why.
\end{proposition}
\begin{proof}
\uses{lag:noether-normalization,lag:parametrization-sets,ana:adapted-projection-ae-fibre}
\leanok
The vanishing ideal of \(A\) is prime, so by Proposition~\ref{lag:noether-normalization} a
linear change of coordinates \(L\) puts it in Noether position over \(\mathcal O_d\), and the
transformed ideal is the vanishing ideal of \(L^{-1}A\); Proposition~\ref{lag:parametrization-sets}
then gives parametrization data for \(L^{-1}A\) in those coordinates, and
Proposition~\ref{ana:adapted-projection-ae-fibre} an almost-everywhere fibre bound for the
adapted projection \(\pi\), on a polydisc \(U_0\) around \(0\) with base set \(\Delta\) and
sheet number \(N\).

Take \(B=\pi\circ L^{-1}\) and \(U=L(U_0)\). A point \(z\) lies in \(A\cap U\cap B^{-1}(y)\)
exactly when \(L^{-1}z\) lies in \(L^{-1}A\cap U_0\cap\pi^{-1}(y)\), so the two fibres
correspond bijectively under the injective map \(L\) and have the same cardinality; the same
correspondence sends \(A\cap U\) to \(L^{-1}A\cap U_0\), so \(B(A\cap U)\subset\Delta\). Only
the source of the projection has moved: the target \(\mathbf C^{d}\), the set \(\Delta\) and
the null set of exceptional \(y\) are unchanged.
\end{proof}

\begin{lemma}[Genericity of the normalizing direction, for a finite family]
\label{ana:dense-noether-normalization}
\uses{lag:noether-normalization,lag:parametrization-sets}
\lean{TranscendentalNakamaye.Pending.IrreducibleGermSheetBound.exists_parametrizationData_near}
\leanok
\source{demailly-cadg:II.4.19}
Let \(A_1,\dots,A_r\) be irreducible analytic germs at \(0\in\mathbf C^{n}\), let \(d\le n\)
be a dimension such that each \(A_i\) is put in Noether position over \(\mathcal O_d\), with
parametrization data of dimension \(d\), by SOME linear change of coordinates, let
\(B_0:\mathbf C^{n}\to\mathbf C^{d}\) be \(\mathbf C\)-linear and let \(\delta>0\). Then a
SINGLE linear change of coordinates \(L\) with \(\lVert\pi\circ L^{-1}-B_0\rVert<\delta\),
\(\pi\) the projection on the first \(d\) coordinates, does so for all of them at once, with
parametrization data of dimension \(d\) for every \(L^{-1}A_i\).

Proposition~\ref{lag:noether-normalization} produces one normalizing \(L\) per germ and
asserts nothing about it; the finite family, and not a single germ, is the form the consumer
needs, because the germ of a subvariety at a point is reducible and one direction must serve
all its components at once.

The bound \(d\le n\) is part of the statement and cannot be read off the family: the maps
\(\pi\circ L^{-1}\) all vanish on the coordinates of index \(>n\), so for \(d>n\) none of them
approximates a \(B_0\) with a nonzero such coordinate. For \(r\ge1\) the bound follows from
the parametrization data of any one germ, but the empty family carries no such evidence. For
\(d=0\), and for \(n=0\), the target is the zero space, \(B_0\) and every \(\pi\circ L^{-1}\)
are the zero map, and any of the given changes of coordinates serves.
\end{lemma}
\begin{proof}
\uses{lag:noether-normalization-near-identity,lag:transported-projections-dense,ana:parametrization-dim-intrinsic,lag:parametrization-sets}
\leanok
By Lemma~\ref{lag:transported-projections-dense} there is a linear automorphism \(K\) of
\(\mathbf C^{n}\) with \(\lVert\pi\circ K^{-1}-B_0\rVert<\delta/2\). Let
\[
  C=(\lVert\pi\rVert+1)(\lVert K^{-1}\rVert+1),\qquad \varepsilon=\frac{\delta}{2C}.
\]
Each \(A_i\) is irreducible, so its vanishing ideal is prime, and so is its image \(Q_i\)
under the automorphism of \(\mathcal O_n\) induced by \(K\). Apply
Proposition~\ref{lag:noether-normalization-near-identity} to \(Q_1,\dots,Q_r\) with tolerance
\(\varepsilon\): one automorphism \(U\) of \(\mathbf C^{n}\) with
\(\lVert U^{-1}-\mathrm{id}\rVert\le\varepsilon\) puts every \(Q_i\) in Noether position. Put
\(L=K\circ U\). Composing the two induced automorphisms of \(\mathcal O_n\) shows that the
image of the vanishing ideal of \(A_i\) under the one induced by \(L\) is in Noether position;
that image is the vanishing ideal of \(L^{-1}A_i\), so
Proposition~\ref{lag:parametrization-sets} supplies parametrization data for each
\(L^{-1}A_i\).

Its dimension is \(d\). Indeed by hypothesis \(A_i\) already carries parametrization data of
dimension \(d\), in other coordinates, and by
Lemma~\ref{ana:parametrization-dim-intrinsic} the dimension of parametrization data of an
irreducible germ does not depend on the data.

Finally, \(L^{-1}=U^{-1}\circ K^{-1}\), so
\[
  \pi\circ L^{-1}-B_0=\pi\circ(U^{-1}-\mathrm{id})\circ K^{-1}+\bigl(\pi\circ K^{-1}-B_0\bigr),
\]
whose first term has norm at most
\(\lVert\pi\rVert\,\varepsilon\,\lVert K^{-1}\rVert\le\varepsilon C=\delta/2\) and whose
second has norm \(<\delta/2\).
\end{proof}

\begin{lemma}[The dimension of parametrization data is intrinsic]
\label{ana:parametrization-dim-intrinsic}
\uses{lag:parametrization-sets,lag:local-statement}
\lean{TranscendentalNakamaye.Pending.GermParametrizationDimUnique.dim_eq_dim}
\leanok
\source{demailly-cadg:II.5.3}
Let \(A\) be an irreducible analytic germ at \(0\in\mathbf C^{n}\) and let \(L\), \(L'\) be
linear automorphisms of \(\mathbf C^{n}\) putting \(A\) in Noether position, with
parametrization data \(P\) for \(L^{-1}A\) and \(P'\) for \(L'^{-1}A\). Then
\(\dim P=\dim P'\).
\end{lemma}
\begin{proof}
\uses{ana:parametrization-regular-accumulate,lag:local-statement}
\leanok
By Lemma~\ref{ana:parametrization-regular-accumulate} regular points of \(A\) of dimension
\(\dim P\) occur in every neighbourhood of \(0\). By Theorem~\ref{lag:local-statement} there
is a neighbourhood \(\Omega\) of \(0\) on which \(A\) is pure of some dimension, i.e. on which
every regular point of \(A\cap\Omega'\), \(\Omega'\subset\Omega\) open, has one and the same
dimension; the previous sentence identifies that dimension with \(\dim P\). Applying
Lemma~\ref{ana:parametrization-regular-accumulate} to \(P'\) produces a regular point of
\(A\) of dimension \(\dim P'\) inside \(\Omega\), whence \(\dim P'=\dim P\).
\end{proof}

\begin{lemma}[Regular points of a parametrized germ accumulate at the origin]
\label{ana:parametrization-regular-accumulate}
\uses{lag:parametrization-sets,ana:restricted-volume}
\lean{TranscendentalNakamaye.Pending.GermComponentDimension.exists_regularAtDim_near_zero}
\leanok
\source{demailly-cadg:II.5.3}
Let \(A\) be an analytic germ at \(0\in\mathbf C^{n}\), let \(L\) be a linear automorphism of
\(\mathbf C^{n}\) and let \(P\) be parametrization data for \(L^{-1}A\), of dimension
\(k\). Then every neighbourhood of \(0\) contains a point \(y\) and an open set
\(\Delta\ni y\) such that \(y\) is a regular point of \(A\cap\Delta\) of dimension \(k\).
\end{lemma}
\begin{proof}
\leanok
Work first with \(L^{-1}A\). The discriminant of \(P\) is not identically zero, so base points
\(x\in\mathbf C^{k}\) with \(\mathrm{disc}(x)\ne0\) occur arbitrarily close to \(0\); over
such a base point the fibre of the adapted projection is nonempty, having exactly \(\deg P\ge1\)
points, and every point of it is a regular point of \(L^{-1}A\cap\Delta_0\) of dimension
\(k\) for the polydisc \(\Delta_0\) of the data. Moreover the fibre over a base point of norm
\(<\eta\) lies, for \(\eta\) small, in a prescribed small polydisc, hence in a prescribed
small ball. Thus regular points of dimension \(k\) of \(L^{-1}A\cap\Delta_0\) occur arbitrarily
near \(0\). Regularity in \(\mathbf C^{n}\) is invariant under a linear automorphism, so
applying \(L\) transports such a point \(y'\) to a point \(y=Ly'\) of norm at most
\(\lVert L\rVert\lVert y'\rVert\), regular of dimension \(k\) for
\(A\cap L(\Delta_0)\).
\end{proof}

\begin{lemma}[Sheet bound of a finite family of germs in a prescribed direction]
\label{ana:germ-sheet-bound-near}
\uses{ana:germ-sheet-bound,ana:dense-noether-normalization,ana:germ-sheet-union}
\lean{TranscendentalNakamaye.Pending.IrreducibleGermSheetBound.exists_germSheetBound_near}
\leanok
\source{demailly-cadg:III.2.6}
In the situation of Lemma~\ref{ana:dense-noether-normalization}, the conclusion of
Proposition~\ref{ana:germ-sheet-bound} holds for the union \(A_1\cup\dots\cup A_r\) with a
single \(B\) satisfying \(\lVert B-B_0\rVert<\delta\).
\end{lemma}
\begin{proof}
\uses{ana:dense-noether-normalization,ana:germ-sheet-union}
\leanok
Lemma~\ref{ana:dense-noether-normalization} supplies one normalizing change of coordinates
\(L\) with \(\lVert\pi\circ L^{-1}-B_0\rVert<\delta\) together with parametrization data of
dimension \(d\) for every \(L^{-1}A_i\). The transport in the proof of
Proposition~\ref{ana:germ-sheet-bound} gives a sheet bound for each \(A_i\) in the SAME
direction \(\pi\circ L^{-1}\), and Lemma~\ref{ana:germ-sheet-union} assembles them.
\end{proof}

\begin{lemma}[Finite unions of sheet bounds in one direction]
\label{ana:germ-sheet-union}
\uses{ana:restricted-volume}
\lean{TranscendentalNakamaye.Pending.ChartGermSheetBound.exists_sheetBound_iUnion}
\leanok
Let \(A_1,\dots,A_r\) be subsets of \(\mathbf C^{n}\) and let
\(B:\mathbf C^{n}\to\mathbf C^{d}\) be \(\mathbf C\)-linear. If each \(A_i\) has a sheet bound
in the direction \(B\), so has \(A_1\cup\dots\cup A_r\).
\end{lemma}
\begin{proof}
\leanok
Intersect the neighbourhoods, union the base sets --- a finite union of measurable sets of
finite volume --- and add the sheet numbers. The fibre of the union over \(y\) is contained in
the union of the fibres, whose cardinalities add up; and a finite intersection of sets of full
measure has full measure.
\end{proof}

\begin{lemma}[From the chart germ to the subvariety]
\label{ana:chart-germ-to-subvariety}
\uses{ana:restricted-volume}
\lean{TranscendentalNakamaye.Pending.ChartGermSheetBound.exists_adaptedProjection_of_chartGermSheetBound}
\leanok
Let \(V\) be a closed irreducible analytic subvariety of \(X\) of dimension \(k\), let
\(x\in V\), let \(e\) be the chart at \(x\), and let \(B:\mathbf C^{n}\to\mathbf C^{k}\) be
\(\mathbf C\)-linear. If the chart germ of \(V\) at \(x\) --- the set \(e(V)\) translated so
that \(e(x)\) sits at the origin --- has a sheet bound in the direction \(B\), then the
conclusion of Lemma~\ref{ana:adapted-projection-germ} holds at \(x\) with the SAME \(B\).
\end{lemma}
\begin{proof}
\leanok
Let \(U_0,\Delta_0,N\) be the germ data and put \(c=B(e(x))\). Take
\(U=\{z\in\operatorname{dom}e:e(z)-e(x)\in U_0\}\), which is open and contains \(x\) because
\(e\) is continuous on its domain and \(0\in U_0\), and \(\Delta=\Delta_0+c\), which is
measurable of the same finite volume because translation preserves Lebesgue measure.

For \(z\in U\cap V\) the point \(w=e(z)-e(x)\) lies in the chart germ and in \(U_0\), and
\(B(e(z))-c=B(w)\in\Delta_0\), so \(B(e(z))\in\Delta\). Fix \(y\). The map \(z\mapsto e(z)-e(x)\)
is injective on \(U\), because \(e\) is injective on its domain, and it sends
\(\{z\in U\cap V:B(e(z))=y\}\) into the germ fibre over \(y-c\); so the first set has at most
as many points as the second. Finally the set of \(y\) for which the germ fibre over \(y-c\)
exceeds \(N\) is the translate by \(c\) of a null set, hence null.
\end{proof}

\begin{lemma}[Degenerate projections: null image, hence empty fibres]
\label{ana:projection-null-image-fibre}
\uses{ana:restricted-volume}
\lean{TranscendentalNakamaye.Pending.SingleProjectionSheetBound.ae_encard_eq_zero_of_measure_image_eq_zero,TranscendentalNakamaye.Pending.SingleProjectionSheetBound.exists_singleProjectionSheetBound_of_measure_image_eq_zero}
\leanok
Let \(V\) be a closed irreducible analytic subvariety of \(X\) of dimension \(k\), let
\(x\in V\), let \(e\) be the chart at \(x\), let \(I\) be a \(k\)-element set of coordinates,
and suppose that for SOME neighbourhood \(U_0\) of \(x\) the image \(\pi_I(e(U_0\cap V))\) is
Lebesgue-null in \(\mathbf C^{k}\). Then there are a neighbourhood \(U\) of \(x\) inside the source of \(e\) and a measurable
\(\Delta\subset\mathbf C^{k}\) of finite volume with \(\pi_I(e(V\cap U))\subset\Delta\) such that for almost every \(y\in\mathbf C^{k}\) the fibre
\(\{z\in V\cap U:\pi_I(e(z))=y\}\) is empty.
\end{lemma}
\begin{proof}
\leanok
For any map \(f\) and any set \(S\), the set of values \(y\) whose fibre
\(\{z\in S:f(z)=y\}\) is nonempty is contained in \(f(S)\); so if \(f(S)\) is null then
almost every fibre is empty, i.e. has cardinality at most \(0\). Apply this to
\(f=\pi_I\circ e\) and to \(S=U\cap V\) with \(U=U_0\cap U_1\), where \(U_1\) is the part of
the source of \(e\) whose chart image lies in the unit ball around \(e(x)\); the image over
\(S\) is contained in the image over \(U_0\cap V\), hence null. For the target set take
\(\Delta=\overline B(\pi_I(e(x)),\|\pi_I\|)\), which contains \(\pi_I(e(U\cap V))\) because
\(\|\pi_I(e(z))-\pi_I(e(x))\|\le\|\pi_I\|\,\|e(z)-e(x)\|\le\|\pi_I\|\) there, and which is
measurable of finite volume.
\end{proof}

\begin{lemma}[The projected parametrization for an arbitrary linear projection]
\label{ana:projected-parametrization-linear}
\uses{ana:flat-patch-density,ana:projected-parametrization}
\lean{TranscendentalNakamaye.Pending.LinearProjectionJacobian.hasStrictFDerivAt_linearProjectedParametrization,
TranscendentalNakamaye.Pending.LinearProjectionJacobian.det_comp_linear}
\leanok
Let \(e\) be a chart of \(X\), \(\gamma\) a regular parametrization of \(V\) of
dimension \(k\) and \(B:\mathbf C^{n}\to\mathbf C^{k}\) a \(\mathbf C\)-linear map. Where the
chart sees the patch, \(f_B=B\circ e\circ\gamma\) is strictly differentiable, and its real
Jacobian determinant at \(w\) is \(\lvert\det(BL(w))\rvert^{2}\), where \(L(w)\) is the
complex derivative of \(e\circ\gamma\) at \(w\).
\end{lemma}
\begin{proof}
\uses{ana:projected-parametrization}
\leanok
The map \(f_B\) is the composition of the holomorphic map \(e\circ\gamma\) with a
\(\mathbf C\)-linear map, hence holomorphic, hence strictly differentiable as a map of real
variables, with real derivative the restriction of scalars of \(B\circ L(w)\). The real
determinant of a complex linear endomorphism of \(\mathbf C^{k}\) is the squared modulus of
its complex determinant, and the matrix of \(B\circ L(w)\) in the standard bases is the
product of the matrix of \(B\) with that of \(L(w)\). For \(B=\pi_I\) this is
Lemma~\ref{ana:projected-parametrization}.
\end{proof}

\begin{lemma}[A spanning family of decomposable forms controls the flat density]
\label{ana:plucker-spanning-gram}
\uses{ana:flat-density-minors,ana:cauchy-binet}
\lean{TranscendentalNakamaye.Pending.PluckerSpanningBound.exists_gram_le}
\leanok
\source{demailly-cadg:III.2.6}
Let \(S\) be finite and let \(B_s:\mathbf C^{n}\to\mathbf C^{k}\) be \(\mathbf C\)-linear maps
whose Pl\"ucker vectors \(\bigl(\det (B_s)_{\cdot,I}\bigr)_{\lvert I\rvert=k}\) separate the
points of \(\Lambda^{k}\mathbf C^{n}\). There is a constant \(C=C(B)\ge0\) such that
\[
  \operatorname{Re}\det(A^{H}A)\le C\sum_{s\in S}\lvert\det(B_sA)\rvert^{2}
\]
for every \(n\times k\) complex matrix \(A\).

This is the first step of Demailly's proof of Lemma~\ref{ana:local-flat-volume}: writing the
monomials \(dz_K\wedge d\bar z_L\) in a basis of decomposable \((k,k)\)-forms
\(\beta_s=i\beta_{s,1}\wedge\overline{\beta_{s,1}}\wedge\dots\wedge
i\beta_{s,k}\wedge\overline{\beta_{s,k}}\) expresses the coefficient measures of a positive
current as bounded linear combinations of the nonnegative measures \(T\wedge\beta_s\). The
decomposable form \(\beta_s\) is the same datum as the map \(B_s\) whose rows are the
\(\beta_{s,j}\), and its pairing with the \(k\)-plane spanned by the columns of \(A\) is
\(\det(B_sA)\).
\end{lemma}
\begin{proof}
\uses{ana:flat-density-minors,ana:cauchy-binet}
\leanok
By Cauchy--Binet (Lemma~\ref{ana:cauchy-binet}),
\(\det(B_sA)=\sum_{\lvert I\rvert=k}\det (B_s)_{\cdot,I}\,\xi_I\), where
\(\xi_I=\det A_{I,\cdot}\) is the Pl\"ucker vector of \(A\). So \(s\mapsto\det(B_sA)\) is the
image of \(\xi\) under the linear map \(P:\Lambda^{k}\mathbf C^{n}\to\mathbf C^{S}\) whose
rows are the Pl\"ucker vectors of the \(B_s\), and the hypothesis says exactly that \(P\) is
injective. On the finite-dimensional space \(\Lambda^{k}\mathbf C^{n}\) an injective linear
map is bounded below: there is \(K\) with \(\lVert\xi\rVert\le K\lVert P\xi\rVert\) for all
\(\xi\), the norms being Euclidean. Finally
\(\operatorname{Re}\det(A^{H}A)=\sum_{\lvert I\rvert=k}\lvert\xi_I\rvert^{2}
=\lVert\xi\rVert^{2}\) and \(\lVert P\xi\rVert^{2}=\sum_s\lvert\det(B_sA)\rvert^{2}\), so
the claim holds with \(C=K^{2}\).
\end{proof}

\begin{proposition}[Flat volume from one spanning family of projections]
\label{ana:decomposable-basis-flat-volume}
\uses{ana:plucker-spanning-gram,ana:multiplicity-area,ana:projected-parametrization-linear,ana:flat-patch-density}
\lean{TranscendentalNakamaye.Pending.DecomposableBasisFlatVolume.exists_flatVolumeConst}
\leanok
\source{demailly-cadg:III.2.6}
Let \(V\) have dimension \(k\) and let \(B_s:\mathbf C^{n}\to\mathbf C^{k}\), \(s\in S\)
finite, be a family whose Pl\"ucker vectors separate the points of
\(\Lambda^{k}\mathbf C^{n}\). There is a constant \(C=C(B)<\infty\) with the following
property. For every chart \(e\), every measurable \(\Delta\subset\mathbf C^{k}\), every
\(N\in\mathbf N\) and every finite family of disjoint measurable pieces \(D_i\) of regular
parametrizations \(\gamma_i\) of \(V\) inside the set of parameters that \(e\) sees, if each
\(f_{s,i}=B_s\circ e\circ\gamma_i\) maps \(D_i\) into \(\Delta\) and
\(\sum_i\#\bigl(D_i\cap f_{s,i}^{-1}(y)\bigr)\le N\) for almost every \(y\in\Delta\) and
every \(s\), then
\[
  \sum_i\int_{D_i}\rho^{\flat}_{e,\gamma_i}\le C\,N\operatorname{vol}(\Delta).
\]
Unlike Proposition~\ref{ana:flat-volume-fibre-bound}, the family \(B\) is arbitrary: the
\(\binom nk\) coordinate projections of \(e\) are one admissible choice among many.
\end{proposition}
\begin{proof}
\uses{ana:plucker-spanning-gram,ana:multiplicity-area,ana:projected-parametrization-linear}
\leanok
By Lemma~\ref{ana:flat-density-minors} the flat density at \(w\) is
\(k!\operatorname{Re}\det(L_i(w)^{H}L_i(w))\), which by
Lemma~\ref{ana:plucker-spanning-gram} is at most
\(k!\,C_0\sum_{s}\lvert\det(B_sL_i(w))\rvert^{2}\); by
Lemma~\ref{ana:projected-parametrization-linear} the \(s\)-summand is the real Jacobian
determinant of \(f_{s,i}\), and \(f_{s,i}\) is strictly differentiable on \(D_i\). Summing
over \(i\) and interchanging with the finite sum over \(s\),
Lemma~\ref{ana:multiplicity-area} bounds
\(\sum_i\int_{D_i}\lvert\det Df_{s,i}\rvert\) by \(N\operatorname{vol}(\Delta)\) for each
\(s\). Take \(C=k!\,C_0\,\#S\).
\end{proof}

\begin{lemma}[Spanning is an open condition]
\label{ana:plucker-spanning-open}
\uses{ana:plucker-spanning-gram}
\lean{TranscendentalNakamaye.Pending.PluckerSpanningOpen.isOpen_setOf_spansPlucker,
TranscendentalNakamaye.Pending.PluckerSpanningOpen.exists_dist_lt_spansPlucker}
\leanok
\source{demailly-cadg:III.2.6}
For \(S\) finite, the families \(B\in\operatorname{Hom}(\mathbf C^{n},\mathbf C^{k})^{S}\)
whose Pl\"ucker vectors separate the points of \(\Lambda^{k}\mathbf C^{n}\) form an open set.
Equivalently, such a family admits a \(\delta>0\) with the property that every family \(B'\)
with \(\lVert B'_s-B_s\rVert<\delta\) for all \(s\) again separates points.

This is the clause ``\((\beta_s)\) remains a basis of \(\Lambda^{p,p}\)'' of Demailly's
perturbation step, and it is what makes the perturbation of a single \(\beta_s\) free: each
\(\beta_{s,j}\) may be moved on its own, the others staying put.
\end{lemma}
\begin{proof}
\uses{ana:plucker-spanning-gram}
\leanok
Separation of points is the injectivity of the linear map
\(P_B:\Lambda^{k}\mathbf C^{n}\to\mathbf C^{S}\) whose rows are the Pl\"ucker vectors of the
\(B_s\), as in the proof of Lemma~\ref{ana:plucker-spanning-gram}. The assignment
\(\Pi\mapsto P_\Pi\) is LINEAR in the array \(\Pi\) of Pl\"ucker coordinates, hence continuous
because that array ranges in a finite-dimensional space; and \(B\mapsto\Pi(B)\) is continuous
because each \(\Pi(B)_{s,I}=\det (B_s)_{\cdot,I}\) is a polynomial in the entries of \(B_s\).
So \(B\mapsto P_B\) is continuous, and the injective linear maps out of the
finite-dimensional space \(\Lambda^{k}\mathbf C^{n}\) form an open subset of
\(\operatorname{Hom}(\Lambda^{k}\mathbf C^{n},\mathbf C^{S})\): if \(P\) is injective then
\(\lVert\xi\rVert\le K\lVert P\xi\rVert\) for some \(K\), and every \(Q\) with
\(\lVert Q-P\rVert<K^{-1}\) satisfies \(\lVert Q\xi\rVert\ge(K^{-1}-\lVert Q-P\rVert)
\lVert\xi\rVert>0\) for \(\xi\ne0\). The preimage of an open set under a continuous map is
open.
\end{proof}

\begin{lemma}[Pure dimension of the germ of an irreducible subvariety]
\label{ana:pure-dimension-germ}
\uses{lag:germ-decomposition,lag:noether-normalization,ana:restricted-volume}
\lean{TranscendentalNakamaye.Pending.AdaptedProjectionAssembly.exists_pureDimensionParametrizationData}
\leanok
Let \(V\) be a closed irreducible analytic subvariety of \(X\) of dimension \(k\), let
\(x\in V\), and let \(A_1,\dots,A_r\) be an irredundant family of irreducible analytic germs
whose union is, near \(0\), the chart germ of \(V\) at \(x\). Then each \(A_i\) is put in
Noether position over \(\mathcal O_{k}\) by some linear change of coordinates, with
parametrization data of dimension exactly \(k\).

An irreducible analytic set is of pure dimension, and the dimension of \(V\) is certified
once and for all by a single regular point; what is asserted is that the local irreducible
components at EVERY point of \(V\), and not only near that one, have that same dimension.

The irredundancy is not decoration. Adjoining the germ \(\{0\}\) to a decomposition changes
neither the union nor any sheet bound, and \(\{0\}\) has Noether dimension \(0\); so without
it no common dimension exists. For \(r=0\) nothing is asserted, which happens only when the
chart germ is empty near \(0\), and it is not, since \(x\in V\). For \(k=0\) the point \(x\)
is isolated in \(V\) and every component is \(\{0\}\).
\end{lemma}
\begin{proof}
\uses{lag:noether-normalization,lag:parametrization-sets,ana:parametrization-regular-accumulate,ana:pure-dimension,lag:thin-subgerm}
\leanok
Fix \(i\). Proposition~\ref{lag:noether-normalization} and
Proposition~\ref{lag:parametrization-sets} put \(A_i\) in Noether position by some \(L\), with
parametrization data \(P\) for \(L^{-1}A_i\); what has to be shown is \(\dim P=k\).

By Lemma~\ref{ana:parametrization-regular-accumulate} the regular points of \(A_i\) of
dimension \(\dim P\) accumulate at \(0\), and they are dense in \(A_i\) near \(0\). The
pairwise intersections \(A_i\cap A_l\), \(l\ne i\), are thin in \(A_i\)
(Proposition~\ref{lag:thin-subgerm}), so by irredundancy there is a nonempty open set \(O\)
meeting \(A_i\) and disjoint from every \(A_l\), \(l\ne i\), and inside it a regular point
\(y\) of \(A_i\cap O\) of dimension \(\dim P\). On \(O\) the chart germ of \(V\) coincides with
\(A_i\), so \(y\) is a regular point of the chart germ, hence — after undoing the chart and the
translation, regularity being a local and chart-invariant notion — a regular point of \(V\) of
dimension \(\dim P\). By Theorem~\ref{ana:pure-dimension} every regular point of the
irreducible \(V\) has dimension \(k\), so \(\dim P=k\).
\end{proof}

\begin{lemma}[One direction, prescribed up to \(\delta\), adapted to the germ]
\label{ana:adapted-projection-germ}
\uses{lag:parametrization-sets,ana:adapted-projection-ae-fibre,ana:restricted-volume,ana:germ-sheet-bound-near,ana:chart-germ-to-subvariety,ana:pure-dimension-germ,lag:germ-decomposition}
\lean{TranscendentalNakamaye.Pending.AdaptedProjectionGerm.exists_adaptedProjection}
\leanok
\source{demailly-cadg:III.2.6}
Let \(V\) be a closed irreducible analytic subvariety of \(X\) of dimension \(k\), let
\(x\in V\), let \(e\) be the chart at \(x\), let \(B_0:\mathbf C^{n}\to\mathbf C^{k}\) be
\(\mathbf C\)-linear and let \(\delta>0\). There are a \(\mathbf C\)-linear
\(B:\mathbf C^{n}\to\mathbf C^{k}\) with \(\lVert B-B_0\rVert<\delta\), a neighbourhood
\(U\subset X\) of \(x\) inside the source of \(e\), a measurable
\(\Delta\subset\mathbf C^{k}\) of finite volume with \(B(e(V\cap U))\subset\Delta\), and an
integer \(N\), such that for almost every \(y\in\mathbf C^{k}\)
\[
  \#\{z\in V\cap U\;:\;B(e(z))=y\}\le N .
\]

This is the clause ``each \(\pi_s\) defines a ramified covering of
\(A\cap(\Delta'\times\Delta'')\)'' of Demailly's perturbation step: ONE germ, ONE direction,
and the direction prescribed only up to \(\delta\), which is what makes it compatible with
Lemma~\ref{ana:plucker-spanning-open}. It is the local parametrization theorem
(Proposition~\ref{lag:parametrization-sets}) together with the genericity of the projection
direction: the directions for which the germ is not a ramified covering of its image form a
proper analytic subset of \(\operatorname{Hom}(\mathbf C^{n},\mathbf C^{k})\), whose
complement is therefore dense, so \(B_0\) may be prescribed and only approximated. The
proof is an assembly. The chart germ of \(V\) at \(x\) --- the set \(e(V)\) translated so
that \(e(x)\) sits at the origin --- is an analytic germ, hence near \(0\) an irredundant
finite union of irreducible germs \(A_1,\dots,A_m\) (Theorem~\ref{lag:germ-decomposition}).
Each of them has dimension \(k\), because an irreducible analytic set is of pure dimension:
that is Lemma~\ref{ana:pure-dimension-germ}. Lemma~\ref{ana:germ-sheet-bound-near} then
supplies ONE direction \(B\) within \(\delta\) of \(B_0\) which is a ramified covering of
every \(A_j\), hence of their union; a sheet bound depends only on the germ, so it holds for
the chart germ itself; and Lemma~\ref{ana:chart-germ-to-subvariety} carries it to \(V\) with
the same \(B\).

For \(k=0\) the target is the zero space, \(B=B_0=0\) serves, and the fibre clause says that
a small enough \(U\) meets \(V\) in finitely many points, which holds because a
\(0\)-dimensional analytic set is discrete. The almost-everywhere quantifier cannot be
dropped: the fibre of a BAD direction through \(x\) may be positive-dimensional.
\end{lemma}

\begin{proposition}[Demailly's perturbed decomposable basis]
\label{ana:decomposable-basis-sheet}
\uses{ana:plucker-spanning-open,ana:adapted-projection-germ,ana:restricted-volume}
\lean{TranscendentalNakamaye.Pending.DecomposableBasisSheetBound.exists_decomposableBasisSheetBound}
\leanok
\source{demailly-cadg:III.2.6}
Let \(V\) be a closed irreducible analytic subvariety of \(X\) of dimension \(k\), let
\(x\in V\) and let \(e\) be the chart at \(x\). There are a finite family
\(B_s:\mathbf C^{n}\to\mathbf C^{k}\), \(s\in S\), whose Pl\"ucker vectors separate the
points of \(\Lambda^{k}\mathbf C^{n}\), a neighbourhood \(U\subset X\) of \(x\) inside the
source of \(e\), a measurable \(\Delta\subset\mathbf C^{k}\) of finite volume and an integer
\(N\), such that for every \(s\) the map \(B_s\circ e\) sends \(V\cap U\) into \(\Delta\) and
\[
  \#\{z\in V\cap U\;:\;B_s(e(z))=y\}\le N\qquad\text{for almost every }y\in\mathbf C^{k}.
\]
\end{proposition}
\begin{proof}
\uses{ana:plucker-spanning-open,ana:adapted-projection-germ}
\leanok
Take for \(S\) the set of \(k\)-element sets \(I\) of coordinates and for the initial family
the coordinate projections \(\pi_I\) of \(e\). Their Pl\"ucker matrix is the identity, so they
separate points. Let \(\delta>0\) be the radius Lemma~\ref{ana:plucker-spanning-open}
attaches to that family, and let \(B_I,U_I,\Delta_I,N_I\) be the data
Lemma~\ref{ana:adapted-projection-germ} supplies for the direction \(\pi_I\) and this
\(\delta\). Then \(\lVert B_I-\pi_I\rVert<\delta\) for every \(I\), so the perturbed family
\((B_I)_I\) again separates points.

Put \(U=\operatorname{source}(e)\cap\bigcap_IU_I\), \(\Delta=\bigcup_I\Delta_I\) and
\(N=\max_IN_I\). Then \(U\) is open and contains
\(x\), being a finite intersection of such sets, and lies in the source of \(e\) by
construction; \(\Delta\) is measurable with
\(\operatorname{vol}\Delta\le\sum_I\operatorname{vol}\Delta_I<\infty\); and for each \(I\),
\(U\subset U_I\) gives \(B_I(e(V\cap U))\subset\Delta_I\subset\Delta\) and, for almost every
\(y\),
\[
  \#\{z\in V\cap U:B_I(e(z))=y\}\le\#\{z\in V\cap U_I:B_I(e(z))=y\}\le N_I\le N .
\]
\end{proof}

\begin{remark}[Two routes to the local flat volume, and what each costs]
\label{ana:lelong-basis-form-route}
\uses{ana:decomposable-basis-sheet}
\source{demailly-cadg:III.2.6}
Lemma~\ref{ana:local-flat-volume} below is proved through
Proposition~\ref{ana:decomposable-basis-sheet}, in which the projections are perturbed one at
a time. The alternative would be to fix the chart and run over its \(\binom nk\) coordinate
projections, which is what the Cauchy--Binet expansion of the flat density
(Lemma~\ref{ana:flat-density-minors}) produces directly. A fixed coordinate projection need not
be a ramified covering of the germ (the fibre over \(0\) of the first two coordinates on the image
of \((t,s)\mapsto(t,ts,s)\) is a line), and bounding its fibres is the finiteness of the local degree
of a generically finite projection, which is not proved here; only its degenerate case is
(Lemma~\ref{ana:projection-null-image-fibre}). Perturbing instead costs nothing on the algebraic
side, because separation of points is open (Lemma~\ref{ana:plucker-spanning-open}), and leaves only
the local parametrization theorem in coordinates adapted to the germ
(Lemma~\ref{ana:adapted-projection-germ}).
\end{remark}

\begin{lemma}[Lelong: locally finite FLAT volume in one chart]
\label{ana:local-flat-volume}
\uses{ana:restricted-volume,ana:flat-patch-density,ana:decomposable-basis-flat-volume,ana:decomposable-basis-sheet}
\source{demailly-cadg:III.2.6}
\lean{TranscendentalNakamaye.Pending.LocalFlatVolumeBound.exists_localFlatVolumeBound}
\leanok
Let \(V\) be a closed irreducible analytic subvariety of dimension \(k\) of \(X\)
and \(x\in V\). There are a neighbourhood \(U\subset X\) of \(x\) inside the
source of the chart \(e\) at \(x\), and a constant \(N<\infty\), such that every
finite family of disjoint measurable pieces of regular parametrizations of \(V\)
whose images lie in \(U\) has total FLAT parameter mass
\(\sum_i\int_{D_i}\rho^{\flat}_{e,\gamma_i}\le N\).

No K\"ahler form, class, presentation, potential or polar set occurs: one chart,
one subvariety, and the flat metric of that chart.
\end{lemma}
\begin{proof}
\uses{ana:restricted-volume,ana:decomposable-basis-flat-volume,ana:decomposable-basis-sheet}
\leanok
Work in the coordinates of the chart \(e\) at \(x\). Let \(B,U,\Delta,N\) be the data
Proposition~\ref{ana:decomposable-basis-sheet} attaches to \(x\): a family of projections
whose Pl\"ucker vectors separate points, a neighbourhood of \(x\) inside the source of \(e\),
a measurable target of finite volume and a sheet number. Let \(C=C(B)\) be the constant
Proposition~\ref{ana:decomposable-basis-flat-volume} attaches to that family.

Let \(D_1,\dots,D_m\) be disjoint measurable pieces of regular parametrizations of \(V\)
whose images lie in \(U\). For each \(s\), \(f_{s,i}=B_s\circ e\circ\gamma_i\) maps \(D_i\)
into \(\Delta\), and the fibre count of the family over a point is at most the fibre count of
\(V\cap U\) — the pieces are disjoint and the parametrizations are injective — hence at most
\(N\) for almost every \(y\). So
Proposition~\ref{ana:decomposable-basis-flat-volume} gives
\(\sum_i\int_{D_i}\rho^{\flat}_{e,\gamma_i}\le C\,N\operatorname{vol}(\Delta)\), a bound
independent of the family.
\end{proof}

\begin{lemma}[Lelong: locally finite volume at a singular point]
\label{ana:background-mass-local-singular}
\uses{ana:background-mass,ana:flat-patch-density,ana:local-flat-volume}
\source{demailly-cadg:III.2.6}
\lean{TranscendentalNakamaye.Pending.SubvarietyStokesMassInterfaces.exists_localBackgroundBound_of_notMem_regularLocus}
\leanok
The same conclusion at a point \(x\in V\) that is not a regular point of
dimension \(k\).
\end{lemma}
\begin{proof}
\uses{ana:background-mass,ana:flat-patch-density,ana:local-flat-volume}
\leanok
No property of the point beyond \(x\in V\) is used. Let \(e\) be the chart at
\(x\). The coefficient of \(\beta\) in \(e\) is smooth, hence continuous, on the
target of \(e\), so it attains a maximum norm \(C_0\) on a closed coordinate ball
\(\overline B(e(x),\varrho)\) contained in that target; let \(U_1\) be the
preimage of the open ball \(B(e(x),\varrho)\). For a regular parametrization
\(\gamma\) and a parameter \(w\) with \(\gamma(w)\in U_1\),
Lemma~\ref{ana:flat-patch-density} gives
\(\rho_{\beta,\gamma}(w)\le C_0^{k}\rho^{\flat}_{e,\gamma}(w)\), because
\(\beta\) is a semipositive \((1,1)\)-form and its coefficient at
\(e(\gamma(w))\) has norm at most \(C_0\).

Let \(U_2\) and \(N\) be the neighbourhood and the constant of
Lemma~\ref{ana:local-flat-volume} at \(x\), and put \(U=U_1\cap U_2\). A finite
family of disjoint pieces with images in \(U\) has total background mass at most
\(C_0^{k}\) times its total flat mass, which is at most \(N\); so
\(C=C_0^{k}N<\infty\) works.
\end{proof}

\begin{lemma}[Lelong: locally finite volume of the regular locus]
\label{ana:background-mass-local}
\uses{ana:background-mass-local-regular,ana:background-mass-local-singular}
\lean{TranscendentalNakamaye.Pending.SubvarietyStokesMassInterfaces.exists_localBackgroundBound}
\leanok
Every point \(x\in V\) has a neighbourhood \(U\subset X\) and a constant
\(C<\infty\) such that every finite family of disjoint measurable pieces of
regular parametrizations of \(V\) whose images lie in \(U\) has total background
parameter mass at most \(C\).
\end{lemma}
\begin{proof}
\uses{ana:background-mass-local-regular,ana:background-mass-local-singular}
\leanok
Either \(x\) is a regular point of \(V\) of dimension \(k\), and
Lemma~\ref{ana:background-mass-local-regular} applies, or it is not, and
Lemma~\ref{ana:background-mass-local-singular} applies.
\end{proof}

\begin{lemma}[Lelong: the regular locus has finite volume]
\label{ana:background-mass-finite}
\uses{ana:background-mass,ana:background-mass-local}
\source{demailly-cadg:III.2.7}
\lean{TranscendentalNakamaye.Pending.SubvarietyStokesMassInterfaces.backgroundMass_ne_top,TranscendentalNakamaye.Pending.SubvarietyStokesMass.acTopMass_ne_top_of_zero_potential}
\leanok
The background mass of \(V\) is finite. Equivalently, the absolutely continuous top mass of
the presentation of the zero class by the zero form and the zero potential is finite.
\end{lemma}
\begin{proof}
\uses{ana:background-mass,ana:background-mass-local}
\leanok
\(V\) is closed in the compact manifold \(X\), hence compact, so finitely many of
the neighbourhoods of the previous lemma, say \(U_1,\dots,U_N\) with constants
\(C_1,\dots,C_N\), cover \(V\). Let \((D_i)_{i}\) be a finite family of disjoint
measurable pieces of regular parametrizations of \(V\). Each \(D_i\) maps into
\(V_{\mathrm{reg}}\subset V\), so it is the union of the finitely many measurable
subsets \(D_i\cap\gamma_i^{-1}(U_j)\), \(1\le j\le N\), and therefore
\(\int_{D_i}\le\sum_{j}\int_{D_i\cap\gamma_i^{-1}(U_j)}\). For each fixed \(j\)
the pieces \(D_i\cap\gamma_i^{-1}(U_j)\) are again disjoint pieces of regular
parametrizations, now with all images in \(U_j\), so their masses sum to at most
\(C_j\). Exchanging the two finite sums gives
\(\sum_i\int_{D_i}\le\sum_{j}C_j<\infty\), a bound independent of the family, so
the supremum defining the background mass is finite.
\end{proof}

\begin{lemma}[Stokes along an exhaustion is enough]
\label{ana:background-mass-stokes-exhaustion}
\uses{ana:background-mass,ana:piece-family-exhaustion-supremum,ana:piece-mass-comparison-smooth}
\lean{TranscendentalNakamaye.Pending.SubvarietyStokesMassInterfaces.iSup_formPieceMass_currentForm_eq_of_exhaustion}
\leanok
Let \(v\) be a smooth real function on \(X\) and let \(a_v=\beta+dd^cv\). If
\[
 \sup_m\mathcal M_{a_v}(H_m)=\sup_m\mathcal M_{\beta}(H_m)
\]
for one monotone sequence \((H_m)_m\) of regular-piece families of \(V\) avoiding
\(\emptyset\) that absorbs every compact part of \(V_{\mathrm{reg}}\), then
\[
 \sup_F\mathcal M_{a_v}(F)=\sup_F\mathcal M_{\beta}(F),
\]
the suprema running over all regular-piece families of \(V\).
\end{lemma}
\begin{proof}
\uses{ana:background-mass,ana:piece-family-exhaustion-supremum}
\leanok
Lemma~\ref{ana:piece-family-monotone-exhaustion} produces such a sequence
\((H_m)_m\), and Lemma~\ref{ana:piece-family-exhaustion-supremum} applies to it
with \(a=a_v\) and with \(a=\beta\), rewriting both sides of the conclusion as
suprema along that one sequence. The hypothesis is then exactly the rewritten
identity.

The force of the statement is that the SAME sequence serves both forms: the
suprema over all families are compared through a single exhaustion, whereas for a
fixed family the two masses genuinely differ.
\end{proof}

\begin{lemma}[The Stokes defect is a density on the regular locus]
\label{ana:stokes-defect-density}
\uses{ana:background-mass,ana:patch-transition,fpmma:mixed-telescope,fpmma:mixed-pullback}
\lean{TranscendentalNakamaye.Pending.PatchDensityTelescope.patchFormDensity_sub,TranscendentalNakamaye.Pending.PatchDensityTelescope.patchFormDensity_currentForm_sub,TranscendentalNakamaye.Pending.PatchDensityTelescope.sum_actualMixedTopDensity_telescopeFamily_transition}
\leanok
Let \(V\subset X\) be an irreducible analytic subvariety of dimension \(k\), let
\(a\) and \(b\) be smooth real two-forms on \(X\) and let \(c=a+b\). For a regular
patch \(P=(\gamma,\Omega)\) of \(V\) write \(a_P(w)\) for the coordinate two-form
of \(a\) pulled back to the parameter \(w\in\Omega\), and likewise for \(b\) and
\(c\); then \(c_P(w)=a_P(w)+b_P(w)\) and, with \(T_j\) as in
Lemma~\ref{fpmma:mixed-telescope},
\[
 \delta_P(w):=\rho_c(P,w)-\rho_a(P,w)=\sum_{j=0}^{k-1}W\bigl(T_j(a_P(w),b_P(w))\bigr),
\]
where \(\rho\) denotes the parameter density of a form on the patch. Moreover, for
two patches \(P,Q\) and a parameter \(w\) in their transition domain,
\[
 \delta_Q(w)=\det{}_{\mathbf R}\bigl(D\varphi_{PQ}(w)\bigr)\,\delta_P(\varphi_{PQ}(w)),
\]
\(\varphi_{PQ}\) being the transition of Lemma~\ref{ana:patch-transition}: the
defect obeys the same change-of-parameters law as the densities themselves.

With \(a=\beta\) and \(b=dd^cv\) for a smooth real \(v\) on \(X\), \(\delta\) is the
pointwise Stokes defect of \(\beta+dd^cv\) against \(\beta\) along \(V\).
\end{lemma}
\begin{proof}
\uses{ana:patch-transition,fpmma:mixed-telescope,fpmma:mixed-pullback}
\leanok
The coordinate two-form of a form on a patch is its coefficient in the preferred
chart at the centre of the patch, evaluated at \(\gamma(w)\) and precomposed with
\(D\gamma(w)\). Both taking the chart coefficient and precomposing with a linear
map are additive, so \(c_P(w)=a_P(w)+b_P(w)\), and the first display is
Lemma~\ref{fpmma:mixed-telescope} applied to \(A=a_P(w)\) and \(B=b_P(w)\).

For the second display, Lemma~\ref{ana:patch-transition} gives
\(a_Q(w)=a_P(\varphi_{PQ}(w))\circ D\varphi_{PQ}(w)\) and the same for \(b\). Each
slot of each family \(T_j\) is built from its two arguments using addition only, and
precomposition with a linear map is additive, so every slot of
\(T_j(a_Q(w),b_Q(w))\) is the corresponding slot of
\(T_j(a_P(\varphi_{PQ}(w)),b_P(\varphi_{PQ}(w)))\) precomposed with
\(D\varphi_{PQ}(w)\). The determinant clause of Lemma~\ref{fpmma:mixed-pullback}
multiplies the mixed density of each term by
\(\det_{\mathbf R}(D\varphi_{PQ}(w))\), and summing over \(j\) gives the claim.
\end{proof}

\begin{lemma}[The potential slot of the defect is the \(dd^c\) of the restriction]
\label{ana:stokes-defect-ddc}
\uses{ana:stokes-defect-density,ana:ddc-global-form}
\lean{TranscendentalNakamaye.Pending.PatchDensityTelescope.patchCoordinateForm_ddcForm,TranscendentalNakamaye.Pending.PatchDensityTelescope.patchFormDensity_currentForm_sub_ddc}
\leanok
Let \(v\) be a smooth real function on \(X\), let \(P=(\gamma,\Omega)\) be a regular
patch of \(V\) and let \(w\in\Omega\). Then the coordinate two-form of \(dd^cv\)
pulled back to the parameter \(w\) is the \(dd^c\), taken in the \(k=\dim V\) complex
parameters, of the restriction of \(v\) to the branch:
\[
 (dd^cv)_P(w)=dd^c\bigl(v\circ\gamma_{\mathrm{man}}\bigr)(w),
\]
\(\gamma_{\mathrm{man}}\) being the parametrization read as a map into \(X\).
Consequently the Stokes defect of Lemma~\ref{ana:stokes-defect-density} with
\(a=\beta\) and \(b=dd^cv\) is, in patch parameters, the sum of the \(k\) mixed
densities in which \(dd^c(v\circ\gamma_{\mathrm{man}})\) occupies exactly one slot and
every other slot is the pullback of \(\beta\).
\end{lemma}
\begin{proof}
\uses{ana:stokes-defect-density}
\leanok
The parametrization \(\gamma\) is holomorphic on \(\Omega\), hence \(C^\infty\) there
(complex differentiability on an open subset of \(\mathbf C^k\) implies smoothness), and
it maps \(\Omega\) into the target of the preferred chart at the centre of \(P\). The
chain rule for the normalized \(dd^c\) under a holomorphic map therefore applies to
\(\rho\circ\gamma\), where \(\rho\) is \(v\) read in that chart:
\[
 dd^c(\rho\circ\gamma)(w)=\bigl(dd^c\rho\bigr)(\gamma(w))\circ D\gamma(w),
\]
whose right-hand side is by definition the coordinate two-form of \(dd^cv\) pulled back
to \(w\), and whose left-hand side is \(dd^c(v\circ\gamma_{\mathrm{man}})(w)\) because
\(\gamma_{\mathrm{man}}\) is the chart inverse composed with \(\gamma\). The second
assertion is this identity substituted into Lemma~\ref{ana:stokes-defect-density}.
\end{proof}

\begin{lemma}[The Stokes defect is locally a mixed Monge--Amp\`ere density of the parameters]
\label{ana:stokes-defect-mixed}
\uses{ana:stokes-defect-ddc,lpot:existence-global,fpmma:geometric-density}
\lean{TranscendentalNakamaye.Pending.PatchDensityLocalPotential.exists_localPotential_patchCoordinateForm,TranscendentalNakamaye.Pending.PatchDensityLocalPotential.exists_localPotential_patchFormDensity_currentForm_sub}
\leanok
Let \(v\) be a smooth real function on \(X\), let \(P=(\gamma,\Omega)\) be a regular patch
of \(V\), \(k=\dim V\), and let \(w\in\Omega\). Then there are an open
\(\Omega'\subseteq\Omega\) containing \(w\) and a smooth real function \(\varphi\) on
\(\Omega'\) such that the pullback of \(\beta\) to the parameters is \(dd^c\varphi\) on
\(\Omega'\), and for every \(w'\in\Omega'\)
\[
 \rho_{\beta+dd^cv}(P,w')-\rho_{\beta}(P,w')
  =\sum_{j=0}^{k-1}D_{u_{j,1},\ldots,u_{j,k}}(w'),
\]
where \(u_j\) is the telescope pattern of Lemma~\ref{fpmma:mixed-telescope} applied to the
pair \(\bigl(\varphi,\;v\circ\gamma_{\mathrm{man}}\bigr)\) and \(D\) is the mixed
determinant density of complex Hessians.

In particular the defect is, locally in the parameters, the mixed Monge--Amp\`ere density
of \(k\) smooth real functions of \(k\) complex variables: no ambient object occurs in it.
\end{lemma}
\begin{proof}
\uses{ana:stokes-defect-ddc,lpot:existence-global,fpmma:geometric-density}
\leanok
By Corollary~\ref{lpot:existence-global} the closed real \((1,1)\)-form \(\beta\) has a
smooth local potential \(\rho\) on an open \(W\) of the target of the chart at the centre
of \(P\) containing \(\gamma(w)\). Put \(\Omega'=\Omega\cap\gamma^{-1}(W)\), open because
\(\gamma\) is continuous, and \(\varphi=\rho\circ\gamma\). The parametrization is
holomorphic, so the chain rule used in Lemma~\ref{ana:stokes-defect-ddc} gives, for
\(w'\in\Omega'\),
\[
 dd^c\varphi(w')=\bigl(dd^c\rho\bigr)(\gamma(w'))\circ D\gamma(w')
  =\beta_P(w'),
\]
the last equality because \(dd^c\rho\) is the coefficient of \(\beta\) on \(W\).

Together with Lemma~\ref{ana:stokes-defect-ddc} both arguments of the telescope are now
\(dd^c\)s of smooth real functions on \(\Omega'\), namely of \(\varphi\) and of
\(v\circ\gamma_{\mathrm{man}}\). Since \(dd^c\) is additive on smooth functions, every slot
of the \(j\)-th telescope family is the \(dd^c\) of the corresponding member of the
telescope pattern of the two functions. Lemma~\ref{fpmma:geometric-density} identifies the
mixed exterior density of a family of \(dd^c\)s with the mixed determinant density of the
corresponding complex Hessians, which is the displayed right-hand side.
\end{proof}

\begin{lemma}[Both unsigned slots of the Stokes telescope are plurisubharmonic]
\label{ana:stokes-defect-psh}
\uses{ana:stokes-defect-mixed,stokesred:density-nonneg,fpma:hermitian-hessian,levi:real-hessian}
\lean{TranscendentalNakamaye.Pending.TelescopeSlotPositivity.coefficientSemipositive_patchCoordinateForm_kahler,
TranscendentalNakamaye.Pending.TelescopeSlotPositivity.coefficientSemipositive_patchCoordinateForm_currentForm,
TranscendentalNakamaye.Pending.TelescopeSlotPositivity.exists_localPotential_telescope_posSemidef,
TranscendentalNakamaye.Pending.TelescopeSlotPositivity.posSemidef_complexHessian_telescopeFunctions_of_ne,
TranscendentalNakamaye.Pending.SmoothDDCHessianPositivity.complexHessian_posSemidef_of_leviForm_nonneg,
TranscendentalNakamaye.Pending.SmoothDDCHessianPositivity.complexHessian_posSemidef_of_coefficientSemipositive,
TranscendentalNakamaye.Pending.SmoothDDCHessianPositivity.complexHessian_posSemidef_of_contDiffOn}
\leanok
Keep the notation of Lemma~\ref{ana:stokes-defect-mixed}, and assume moreover that the
presentation is admissible \emph{at error} \(\epsilon=1\), with representative \(0\) and
potential the smooth function \(v\), so that its pole set is empty. Then the open set
\(\Omega'\) and the local potential \(\varphi\) may be chosen so that, in addition to the
density identity,
\[
 H_{\varphi}(w')\succeq 0\quad\text{and}\quad
 H_{\varphi+v\circ\gamma_{\mathrm{man}}}(w')\succeq 0
 \qquad\text{for every }w'\in\Omega' .
\]
Consequently every slot \(i\neq j\) of the telescope family \(u_j\) has a positive
semidefinite complex Hessian on \(\Omega'\).
\end{lemma}
\begin{proof}
\uses{ana:stokes-defect-mixed,stokesred:density-nonneg,fpma:hermitian-hessian,levi:real-hessian}
\leanok
Take \(\Omega'\) and \(\varphi\) from Lemma~\ref{ana:stokes-defect-mixed}, so that the
pullback of \(\beta\) to the parameters is \(dd^c\varphi\) and that of
\(\beta+dd^cv\) is \(dd^c(\varphi+v\circ\gamma_{\mathrm{man}})\) on \(\Omega'\).

The K\"ahler background \(\beta\) is semipositive at every point of \(X\); its pullback
along the holomorphic parametrization of a regular patch is therefore semipositive at
every parameter, since pulling a semipositive \((1,1)\)-form back along a holomorphic map
preserves semipositivity. Admissibility at error \(\epsilon\) gives semipositivity of
\(\theta+\epsilon\beta+dd^cv\) off the pole set of the potential; with representative
\(\theta=0\) and \(\epsilon=1\) that object is exactly the current form
\(\beta+dd^cv\), and the pole set is empty here, so it is semipositive at the image of
every parameter and the same pullback argument applies. The restriction
\(\epsilon=1\) is not cosmetic: at \(\epsilon\neq1\) admissibility controls
\(\epsilon\beta+dd^cv\) and says nothing about \(\beta+dd^cv\), while the density
identity of Lemma~\ref{ana:stokes-defect-mixed} compares against the unscaled
\(\beta\).

It remains to pass from semipositivity of a \(dd^c\) to semidefiniteness of a complex
Hessian. For a function of class \(C^2\) at a point, the complex diagonal of \(dd^c\)
evaluated on a vector is \(\pi^{-1}\) times the Levi form in that direction, and the
complex Hessian is Hermitian with
\(\langle x,H_u(z)x\rangle=\mathcal L u(z)(x)\); so a nonnegative Levi form in every
direction is exactly a positive semidefinite complex Hessian. Both \(\varphi\) and
\(\varphi+v\circ\gamma_{\mathrm{man}}\) are \(C^\infty\) on the open set \(\Omega'\),
hence \(C^2\) at each of its points.

Finally, the slot \(i\) of \(u_j\) carries \(\varphi+v\circ\gamma_{\mathrm{man}}\) when
\(i<j\) and \(\varphi\) when \(i>j\); so every \(i\neq j\) is one of the two functions
just treated.
\end{proof}

\begin{lemma}[An exhaustion whose every piece carries a local potential]
\label{ana:stokes-defect-good-exhaustion}
\uses{ana:stokes-defect-mixed,ana:local-piece-property,ana:piece-family-good-exhaustion}
\lean{TranscendentalNakamaye.Pending.RegularPieceGoodExhaustion.HasLocalTelescopePotential,
TranscendentalNakamaye.Pending.RegularPieceGoodExhaustion.localTelescopePotential,
TranscendentalNakamaye.Pending.RegularPieceGoodExhaustion.exists_exhaustion_localTelescopePotential}
\leanok
Let \(v\) be a smooth real function on \(X\) and \(k=\dim V\). There is one monotone
sequence \((H_m)_m\) of regular-piece families of \(V\), absorbing every compact part
of \(V_{\mathrm{reg}}\), with the following property. For every \(m\) and every piece
\(D\) of \(H_m\), with patch \(P=(\gamma,\Omega_P)\), there are an open
\(\Omega\subseteq\Omega_P\) with \(D\subseteq\Omega\) and a smooth real function
\(\varphi\) on \(\Omega\) such that \(v\circ\gamma_{\mathrm{man}}\) is smooth on
\(\Omega\) and, for every \(w'\in\Omega\),
\[
 \rho_{\beta+dd^cv}(P,w')-\rho_{\beta}(P,w')
  =\sum_{j=0}^{k-1}D_{u_{j,1},\ldots,u_{j,k}}(w'),
\]
\(u_j\) being the telescope pattern of \(\bigl(\varphi,\,v\circ\gamma_{\mathrm{man}}\bigr)\).

So the defect may be integrated over each piece as the mixed Monge--Amp\`ere density of
\(k\) smooth real functions of \(k\) complex variables, with no local choice left to make
inside a piece.
\end{lemma}
\begin{proof}
\uses{ana:stokes-defect-mixed,ana:piece-family-good-exhaustion}
\leanok
Let \(\mathcal G(P)\) consist of those \(S\subseteq\Omega_P\) that lie in some open
\(\Omega\subseteq\Omega_P\) carrying a \(\varphi\) as in the statement. Clause (i) of
Definition~\ref{ana:local-piece-property} is immediate: a subset of \(S\) lies in the
same \(\Omega\). Clause (ii) is exactly Lemma~\ref{ana:stokes-defect-mixed}, which
produces, for every parameter \(w\in\Omega_P\), an open \(\Omega'\ni w\) inside
\(\Omega_P\) and a \(\varphi\) on it with the displayed identity, and on which
\(v\circ\gamma_{\mathrm{man}}\) is smooth because \(v\) is smooth on \(X\) and
\(\gamma\) is holomorphic; taking \(S=\Omega'\) gives \(\Omega'\in\mathcal G(P)\).
So \(\mathcal G\) is a local piece property, and
Lemma~\ref{ana:piece-family-good-exhaustion} with \(A=\emptyset\) produces the
sequence.
\end{proof}

\begin{lemma}[An exhaustion whose local potentials are plurisubharmonic]
\label{ana:stokes-defect-good-exhaustion-psh}
\uses{ana:stokes-defect-psh,ana:local-piece-property,ana:piece-family-good-exhaustion,ana:stokes-defect-good-exhaustion}
\lean{TranscendentalNakamaye.Pending.PositiveTelescopeExhaustion.HasPositiveLocalTelescopePotential,
TranscendentalNakamaye.Pending.PositiveTelescopeExhaustion.hasLocalTelescopePotential_of_positive,
TranscendentalNakamaye.Pending.PositiveTelescopeExhaustion.positiveLocalTelescopePotential,
TranscendentalNakamaye.Pending.PositiveTelescopeExhaustion.exists_exhaustion_positiveLocalTelescopePotential}
\leanok
In the situation of Lemma~\ref{ana:stokes-defect-good-exhaustion}, assume in addition that
the presentation is admissible at error \(\epsilon=1\) with representative \(0\) and
potential the smooth function \(v\). Then the sequence \((H_m)_m\) may be chosen so that
the local data attached to every piece satisfy, besides the density identity,
\[
 H_{\varphi}(w')\succeq 0\quad\text{and}\quad H_{\varphi+v\circ\gamma_{\mathrm{man}}}(w')\succeq 0
 \qquad\text{for every }w'\in\Omega .
\]
Forgetting the two certificates gives back the property of
Lemma~\ref{ana:stokes-defect-good-exhaustion}, so the strengthened sequence is in
particular one of the sequences produced there.
\end{lemma}
\begin{proof}
\uses{ana:stokes-defect-psh,ana:local-piece-property,ana:piece-family-good-exhaustion}
\leanok
Enlarge the family \(\mathcal G(P)\) of the proof of
Lemma~\ref{ana:stokes-defect-good-exhaustion} by requiring of the witnessing pair
\((\Omega,\varphi)\) the two semidefiniteness statements on all of \(\Omega\). Clause (i)
of Definition~\ref{ana:local-piece-property} is unaffected --- a subset of \(S\) is
witnessed by the same pair. Clause (ii) is Lemma~\ref{ana:stokes-defect-psh}, which
produces, near every parameter, ONE pair \((\Omega',\varphi)\) carrying the smoothness, both
semidefiniteness statements and the density identity together. So the enlarged
\(\mathcal G\) is again a local piece property and
Lemma~\ref{ana:piece-family-good-exhaustion} applies to it verbatim.

This is not a cosmetic strengthening. Lemma~\ref{ana:stokes-defect-good-exhaustion}
records only smoothness of \(\varphi\) and the density identity, and the identity alone says
nothing about \(H_\varphi\): at \(k=1\) it reads \((a+b)-a=b\). So a potential selected by
that lemma carries no plurisubharmonicity, while the fixed slots of the estimate of
Theorem~\ref{cln:telescope-cutoff} must be plurisubharmonic. The two certificates are
exactly what closes that gap.
\end{proof}

\begin{lemma}[The cutoff bound on one piece of the exhaustion]
\label{ana:stokes-defect-piece-cutoff}
\uses{ana:stokes-defect-good-exhaustion-psh,cln:telescope-flat-slot,cln:telescope-defect-cutoff,cln:telescope-cutoff}
\lean{TranscendentalNakamaye.Pending.PositiveTelescopeExhaustion.FormPiecePositiveTelescopeData,
TranscendentalNakamaye.Pending.PositiveTelescopeExhaustion.FormPiecePositiveTelescopeData.of_positiveLocalTelescopePotential,
TranscendentalNakamaye.Pending.PositiveTelescopeExhaustion.abs_integral_mul_mixedDensity_telescope_piece_le,
TranscendentalNakamaye.Pending.PositiveTelescopeExhaustion.abs_integral_mul_telescopeIntegrand_piece_le}
\leanok
Let \(H\) be a finite regular-piece family every piece of which carries a
plurisubharmonic local telescope potential, and fix for each piece \(i\) a witnessing pair
\((\Omega_i,\varphi_i)\). Let \(\chi\) be smooth on \(\Omega_i\) with compact support
contained in \(\Omega_i\), let \(A\ge0\) satisfy
\(-A\,\mathrm{Id}\preceq H_\chi\preceq A\,\mathrm{Id}\) on \(\operatorname{supp}\chi\), and
let \(C\) bound \(|v\circ\gamma_{\mathrm{man}}|\) there. Then for every \(j<k\)
\[
 \Bigl|\int\chi\,D_{u_j}\,d\lambda\Bigr|
  \;\le\;C\,A\int_{\operatorname{supp}\chi}D_{u_j^{\flat}}\,d\lambda ,
\]
\(u_j\) being the telescope family of \((\varphi_i,\,v\circ\gamma_{\mathrm{man}})\) with
signed index \(j\) and \(u_j^{\flat}\) its flat slot replacement. Summing over \(j\), the
same bound holds for the cutoff against the piece's whole Stokes defect:
\[
 \Bigl|\int\chi\bigl(\rho_{\beta+dd^cv}(P,\cdot)-\rho_{\beta}(P,\cdot)\bigr)d\lambda\Bigr|
  \;\le\;C\,A\sum_{j<k}\int_{\operatorname{supp}\chi}D_{u_j^{\flat}}\,d\lambda .
\]
\end{lemma}
\begin{proof}
\uses{ana:stokes-defect-good-exhaustion-psh,cln:telescope-flat-slot,cln:telescope-defect-cutoff}
\leanok
The first display is Lemma~\ref{cln:telescope-flat-slot} on the open set \(\Omega_i\), whose
two plurisubharmonicity hypotheses are the certificates carried by the witnessing pair of
the piece and whose remaining hypotheses are those on \(\chi\); the second is
Theorem~\ref{cln:telescope-defect-cutoff} on the same data, together with the density
identity of the witnessing pair, which on the piece identifies the sum of the \(k\)
telescope densities with the signed density defect.
\end{proof}

\begin{lemma}[The collar obligation does not depend on the chosen potential]
\label{ana:stokes-obligation-potential-free}
\uses{ana:stokes-defect-good-exhaustion,ana:stokes-defect-good-exhaustion-psh}
\lean{TranscendentalNakamaye.Pending.PositiveTelescopeExhaustion.telescopeIntegrand_congr,
TranscendentalNakamaye.Pending.PositiveTelescopeExhaustion.pairingData_transport,
TranscendentalNakamaye.Pending.PositiveTelescopeExhaustion.remainderData_transport,
TranscendentalNakamaye.Pending.PositiveTelescopeExhaustion.finitePairingRemainderData_of_positive,
TranscendentalNakamaye.Pending.CertifiedClosedCurrentStokes.closedCurrentFinitePairingRemainderData_of_certified}
\leanok
Fix a finite regular-piece family \(H\) and two choices of local telescope data for it,
i.e.\ two families of witnessing pairs \((\Omega_i,\varphi_i)\) and
\((\Omega_i',\varphi_i')\). Then on each piece the two telescope integrands
\(\sum_{j<k}D_{u_j}\) and \(\sum_{j<k}D_{u_j'}\) coincide. Consequently local
weighted-pushforward measures, test functions and a real remainder bounding their finite
pairing sum for one choice do the same for the other, with the SAME remainder.
\end{lemma}
\begin{proof}
\uses{ana:stokes-defect-good-exhaustion}
\leanok
On the piece itself both telescope integrands equal the signed density defect
\(\rho_{\beta+dd^cv}(P,\cdot)-\rho_{\beta}(P,\cdot)\), by the density identity of the
respective witnessing pair; the two identities hold on neighbourhoods of the piece and are
compared on the piece, where both apply. Since the pairing datum consists of measures and
test functions together with the a.e.\ identification of the telescope integrand on the
piece, and since the bound it carries mentions only the measures and the test functions,
transporting it changes nothing but the integrand it is matched against.
\end{proof}

\begin{lemma}[Cartan generators cut out a closed analytic set]
\label{ana:stokes-cartan-generators}
\uses{ana:cartan-coherence,ana:ideal-sheaves,fpgeom:complex-analytic-subvariety}
\source{demailly-cadg:II.4.29}
\lean{TranscendentalNakamaye.Pending.StokesSingularLocus.exists_cartanGenerators}
\leanok
Let \(A\) be a closed analytic subset of a complex manifold \(M\). Every point has an open neighbourhood \(U\) with holomorphic \(s_1,\dots,s_N\) on \(U\)
that generate the ideal \(\mathcal I_A\) at every point of \(U\), and \(A\cap U=\{y\in U:s_1(y)=\dots=s_N(y)=0\}\).
\end{lemma}
\begin{proof}
\leanok
The generators are those of Cartan's theorem (Theorem~\ref{ana:cartan-coherence}). Each \(s_i\) generates its own germ, so it lies in \(\mathcal I_A\) and
vanishes on \(A\). If \(y\in U\setminus A\), a neighbourhood of \(y\) misses the closed set \(A\), so the constant \(1\) lies in \(\mathcal I_{A,y}\); writing
\(1=\sum_ja_jt_j\) near \(y\) with each \(t_j\) among the \(s_i\), some \(s_i(y)\ne0\).
\end{proof}

\begin{lemma}[The singular locus of \(V\) is analytic, cut out by Jacobian minors]
\label{ana:stokes-singular-locus-analytic}
\uses{ana:stokes-cartan-generators,ana:cartan-coherence,lag:regular-locus-dense,lag:irreducible-germ-pure,lag:thin-subgerm,lag:regular-point-basics,fpgeom:complex-analytic-subvariety,ana:stokes-cutoffs}
\source{demailly-cadg:II.4.31,demailly-cadg:II.4.29,demailly-cadg:II.4.26}
\lean{TranscendentalNakamaye.Pending.StokesSingularLocusAnalytic.evalBasis,TranscendentalNakamaye.Pending.StokesSingularLocusAnalytic.evalBasis_injective,TranscendentalNakamaye.Pending.StokesSingularLocusAnalytic.exists_cols_det_ne_zero,TranscendentalNakamaye.Pending.StokesSingularLocusAnalytic.exists_minor_ne_zero_of_le_finrank,TranscendentalNakamaye.Pending.StokesSingularLocusAnalytic.injective_fderiv_toFun,TranscendentalNakamaye.Pending.StokesSingularLocusAnalytic.linearIndependent_normal_comp,TranscendentalNakamaye.Pending.StokesSingularLocusAnalytic.extendFirst,TranscendentalNakamaye.Pending.StokesSingularLocusAnalytic.proj_extendFirst,TranscendentalNakamaye.Pending.StokesSingularLocusAnalytic.extendFirst_extendFirst,TranscendentalNakamaye.Pending.StokesSingularLocusAnalytic.extendFirst_of_eq,TranscendentalNakamaye.Pending.StokesSingularLocusAnalytic.extendFirst_apply_of_lt,TranscendentalNakamaye.Pending.StokesSingularLocusAnalytic.extendFirst_apply_of_le,TranscendentalNakamaye.Pending.StokesSingularLocusAnalytic.extendFirst_mem_plane_iff,TranscendentalNakamaye.Pending.StokesSingularLocusAnalytic.differentiableOn_extendFirst,TranscendentalNakamaye.Pending.StokesSingularLocusAnalytic.extendLHE,TranscendentalNakamaye.Pending.StokesSingularLocusAnalytic.exists_chartFlattening,TranscendentalNakamaye.Pending.StokesSingularLocusAnalytic.differentiableAt_chart,TranscendentalNakamaye.Pending.StokesSingularLocusAnalytic.fderiv_mem_span_of_isInSpanAt,TranscendentalNakamaye.Pending.StokesSingularLocusAnalytic.holomorphicAt_coord_chart,TranscendentalNakamaye.Pending.StokesSingularLocusAnalytic.le_finrank_span_fderiv_of_regularAtDim,TranscendentalNakamaye.Pending.StokesSingularLocusAnalytic.regularAtDim_of_minor_ne_zero,TranscendentalNakamaye.Pending.StokesSingularLocusAnalytic.singularLocus,TranscendentalNakamaye.Pending.StokesSingularLocusAnalytic.mem_singularLocus,TranscendentalNakamaye.Pending.StokesSingularLocusAnalytic.singularLocus_subset,TranscendentalNakamaye.Pending.StokesSingularLocusAnalytic.disjoint_singularLocus_regularLocus,TranscendentalNakamaye.Pending.StokesSingularLocusAnalytic.isClosed_singularLocus,TranscendentalNakamaye.Pending.StokesSingularLocusAnalytic.isComplexAnalyticSubset_singularLocus,TranscendentalNakamaye.Pending.StokesSingularLocusAnalytic.exists_stokesCutoffs_singularLocus}
\leanok
Let \(V\subset M\) be an irreducible analytic subvariety of dimension \(k\) of a complex manifold of dimension \(n\), and
\(V_{\mathrm{sing}}=V\setminus V_{\mathrm{reg}}\), \(V_{\mathrm{reg}}\) the points where \(V\) is regular of dimension \(k\). Near every point of \(M\) there are
holomorphic functions \(g_1,\dots,g_N\) generating the ideal \(\mathcal I_V\) at every point of a neighbourhood \(U\), and on \(U\), in the preferred chart,
\[
 V_{\mathrm{sing}}\cap U=\Bigl\{z\in U:\ g_1(z)=\dots=g_N(z)=0,\ \det\Bigl(\frac{\partial g_j}{\partial z_l}(z)\Bigr)_{j\in J,\,l\in L}=0\ \text{for all }|J|=|L|=n-k\Bigr\}.
\]
In particular \(V_{\mathrm{sing}}\) is a closed analytic subset of \(M\), disjoint from \(V_{\mathrm{reg}}\); and if \(M\) is compact with a K\"ahler form \(\omega\),
the cutoffs of Lemma~\ref{ana:stokes-cutoffs} exist for \(A=V_{\mathrm{sing}}\): they vanish near \(V_{\mathrm{sing}}\) and tend to \(1\) at every point of
\(V_{\mathrm{reg}}\).
\end{lemma}
\begin{proof}
\leanok
The generators exist by Cartan's theorem (Lemma~\ref{ana:stokes-cartan-generators}), and \(V\cap U\) is their common zero set. Let \(y\in V\cap U\).

If \(y\) is a regular point of dimension \(k\), \(V\) is near \(y\) the zero set of the \(n-k\) normal coordinates \(u_1,\dots,u_{n-k}\) of a flattening
(read in the chart); they vanish on \(V\), so \(u_i=\sum_jh_{ij}g_j\) near \(y\) with \(h_{ij}\) holomorphic, and at \(y\), where every \(g_j\) vanishes,
\(du_i(y)=\sum_jh_{ij}(y)\,dg_j(y)\). The \(du_i(y)\) are independent (the flattening is biholomorphic), so the \(dg_j(y)\) span a space of dimension
\(\ge n-k\); choosing \(n-k\) independent ones among them and then \(n-k\) columns on which they are independent gives a nonzero \((n-k)\)-minor.

Conversely, let a minor with rows \(J\) be nonzero at \(y\). The implicit function theorem flattens \(Z=\{g_j=0,\ j\in J\}\), a manifold of dimension
\(k\) containing \(V\), near \(y\): \(\Phi(Z)\) is the coordinate plane \(\mathbf C^k\times0\). On a ball \(D\) of \(\mathbf C^k\) around \(0=\Phi(y)\), the set
\(D_V=\{w\in D:\Phi^{-1}(w,0)\in V\}\) is the common zero set of the holomorphic \(G_i(w)=g_i(\Phi^{-1}(w,0))\), and contains \(0\). By the density of regular
points (Theorem~\ref{lag:regular-locus-dense}, in its coordinate form) \(D_V\) has a regular point \(w_1\) of some dimension \(j\le k\), flattened by a
local biholomorphism \(\Phi_1\) of \(\mathbf C^k\); then \((\Phi_1\times\mathrm{id})\circ\Phi\) flattens \(V\) at \(\Phi^{-1}(w_1,0)\) onto a coordinate
plane of dimension \(j\) (near that point, \(V=V\cap Z\) because \(V\subseteq Z\)). As \(V\) is irreducible it is of pure dimension \(k\)
(Proposition~\ref{lag:irreducible-germ-pure}), so \(j=k\), i.e.\ \(D_V\) contains a neighbourhood of \(w_1\). Each \(G_i\) vanishes there, hence on the
whole ball by the identity theorem (Proposition~\ref{lag:thin-subgerm}); so \(V=Z\) near \(y\), and \(y\) is regular of dimension \(k\)
(Lemma~\ref{lag:regular-point-basics}). Thus the displayed set is exactly \(V_{\mathrm{sing}}\cap U\).

Regularity is an open condition along \(V\), so \(V_{\mathrm{sing}}\) is closed; with the presentation it is analytic, and Lemma~\ref{ana:stokes-cutoffs}
applies to it. This is Demailly's proof of II.4.31 for a pure-dimensional set.
\end{proof}

\begin{lemma}[A global quasi-plurisubharmonic function with poles exactly along an analytic set]
\label{ana:stokes-pole-potential}
\uses{ana:cartan-coherence,bou:concentration-potentials,posclass:background}
\source{dp04-kahler-cone:2.1,demailly-cadg:III.2.2}
\lean{TranscendentalNakamaye.Pending.StokesPolePotential.isCoherent_vanishingIdeal,TranscendentalNakamaye.Pending.StokesPolePotential.leviForm_congr_of_eventuallyEq,TranscendentalNakamaye.Pending.StokesPolePotential.continuousAt_chart_of_chartSmooth,TranscendentalNakamaye.Pending.StokesPolePotential.exists_polePotential}
\leanok
Let \((M,\omega)\) be a compact K\"ahler manifold and \(A\subset M\) a closed analytic subset. There are a smooth function \(W\colon M\to[0,1)\)
with \(W=0\) exactly on \(A\), and \(C>0\), such that at every point where \(W>0\), in every chart, the Levi form of \(\log W\) satisfies
\(L_{\log W}(v)\ge-C\,\omega(v,Jv)\). Consequently \(u=\frac12\log W\colon M\to[-\infty,0)\) is smooth on \(M\setminus A\), \(u=-\infty\) exactly on \(A\), and
\(dd^cu\ge-\frac C{2\pi}\,\omega\) on \(M\setminus A\).
\end{lemma}
\begin{proof}
\leanok
By Cartan's theorem (Theorem~\ref{ana:cartan-coherence}) the ideal sheaf \(\mathcal I_A\) is coherent, so Lemma~\ref{bou:concentration-potentials} gives a
weight \(U=\sum_a\theta_a^2\sum_i|g_{a,i}|^2\) built from local generators of \(\mathcal I_A\): it is smooth, \(U=0\) on \(A\), \(U>0\) off the closed set \(A\), and
\(dd^c\log U+A'\omega\ge0\) where \(U>0\) (part (3), the case \(s=0\)); in a chart the coefficient of \(dd^c\) on a line is the Levi form divided by \(\pi\), so
\(L_{\log U}(v)\ge-\pi A'\omega(v,Jv)\). \(U\) is continuous on the compact \(M\), hence \(U<K\) for a constant \(K\); put \(W=U/K\) and \(C=\pi A'\). Near a point
where \(W>0\), \(\log W=\log U-\log K\), which has the same Levi form.
\end{proof}

\begin{lemma}[The cutoffs at a point: the convex glue and the Levi bound]
\label{ana:stokes-cutoff-pointwise}
\uses{psh:local}
\source{demailly-cadg:III.2.2}
\lean{TranscendentalNakamaye.Pending.StokesCutoff.realHessian_comp,TranscendentalNakamaye.Pending.StokesCutoff.leviForm_comp,TranscendentalNakamaye.Pending.StokesCutoff.mul_leviForm_le_leviForm_comp,TranscendentalNakamaye.Pending.StokesCutoff.exp_mul_leviForm_le_leviForm_exp,TranscendentalNakamaye.Pending.StokesCutoff.glue,TranscendentalNakamaye.Pending.StokesCutoff.contDiff_glue,TranscendentalNakamaye.Pending.StokesCutoff.glue_eq_zero,TranscendentalNakamaye.Pending.StokesCutoff.glue_nonneg,TranscendentalNakamaye.Pending.StokesCutoff.glue_exp_nonneg,TranscendentalNakamaye.Pending.StokesCutoff.monotone_glue,TranscendentalNakamaye.Pending.StokesCutoff.glue_one,TranscendentalNakamaye.Pending.StokesCutoff.deriv_glue_nonneg,TranscendentalNakamaye.Pending.StokesCutoff.deriv_deriv_glue_nonneg,TranscendentalNakamaye.Pending.StokesCutoff.deriv_glue_one,TranscendentalNakamaye.Pending.StokesCutoff.deriv_glue_mul_le,TranscendentalNakamaye.Pending.StokesCutoff.glue_le_one,TranscendentalNakamaye.Pending.StokesCutoff.glue_exp_eq_zero,TranscendentalNakamaye.Pending.StokesCutoff.glue_exp_le_one,TranscendentalNakamaye.Pending.StokesCutoff.tendsto_glue_exp,TranscendentalNakamaye.Pending.StokesCutoff.neg_le_leviForm_glue_exp}
\leanok
Let \(\theta(t)=e^2\exp(-1/(t-\frac12))\) for \(t>\frac12\) and \(\theta(t)=0\) otherwise. Then \(\theta\) is smooth, \(\theta=0\) on \((-\infty,\frac12]\),
\(\theta(1)=1\), \(\theta'\ge0\), \(\theta''\ge0\) on \((-\infty,1]\), \(\theta'(1)=4\), and \(0\le\theta(e^{t/m})\le1\) for \(t\le0\), \(m>0\);
\(\theta(e^{t/m})=0\) for \(t\le-m\log2\), and \(\theta(e^{t/m})\to1\) as \(m\to\infty\) for every real \(t\). Moreover, for a function \(u\) on
\(\mathbf C^n\) which is \(\mathcal C^2\) near \(z\) with \(u(z)\le0\), \(m>0\), \(c\ge0\) and a vector \(v\) with \(L_u(z)(v)\ge-c\),
\[
 L_{\theta(e^{u/m})}(z)(v)\ \ge\ -\frac{4c}{m},
\]
where \(L\) is the Levi form \((D^2u(v,v)+D^2u(iv,iv))/4\).
\end{lemma}
\begin{proof}
\leanok
The second-order chain rule gives \(L_{\theta\circ f}=\theta'(f)L_f+\theta''(f)|\partial f|^2\ge\theta'(f)L_f\) where \(\theta''(f)\ge0\); with
\(\theta=\exp\) it gives \(L_{e^{u/m}}\ge e^{u/m}L_u/m\). Writing \(g(s)=\exp(-1/s)\) for \(s>0\), \(g'(s)=s^{-2}g(s)\) and
\(g''(s)=s^{-4}(1-2s)g(s)\ge0\) for \(0<s\le\frac12\), which gives the sign of \(\theta''\) on \((-\infty,1]\); \(\theta'\) is then nondecreasing there, so
\(0\le\theta'(s)s\le\theta'(1)=4\) for \(0\le s\le1\). At \(z\), with \(s=e^{u(z)/m}\in(0,1]\):
\(L_{\theta(e^{u/m})}\ge\theta'(s)\,s\,L_u/m\ge-\theta'(s)\,s\,c/m\ge-4c/m\).
\end{proof}

\begin{lemma}[Smooth cutoffs with almost positive Hessian]
\label{ana:stokes-cutoffs}
\uses{ana:stokes-pole-potential,ana:stokes-cutoff-pointwise}
\source{demailly-cadg:III.2.2}
\lean{TranscendentalNakamaye.Pending.StokesPolePotential.cutoff,TranscendentalNakamaye.Pending.StokesPolePotential.threshold,TranscendentalNakamaye.Pending.StokesPolePotential.threshold_pos,TranscendentalNakamaye.Pending.StokesPolePotential.cutoff_eq_zero_of_lt,TranscendentalNakamaye.Pending.StokesPolePotential.cutoff_nonneg,TranscendentalNakamaye.Pending.StokesPolePotential.cutoff_le_one,TranscendentalNakamaye.Pending.StokesPolePotential.tendsto_cutoff,TranscendentalNakamaye.Pending.StokesPolePotential.cutoff_chart_eventuallyEq,TranscendentalNakamaye.Pending.StokesPolePotential.cutoff_chart_eventuallyEq_zero,TranscendentalNakamaye.Pending.StokesPolePotential.contDiffAt_log_chart,TranscendentalNakamaye.Pending.StokesPolePotential.chartSmooth_cutoff,TranscendentalNakamaye.Pending.StokesPolePotential.neg_le_leviForm_cutoff,TranscendentalNakamaye.Pending.StokesPolePotential.exists_stokesCutoffs}
\leanok
In the situation of Lemma~\ref{ana:stokes-pole-potential}, let \(\theta\) be the function of Lemma~\ref{ana:stokes-cutoff-pointwise}
(smooth, convex and nondecreasing on \((-\infty,1]\), \(\theta=0\) on \((-\infty,\frac12]\), \(\theta(1)=1\)). For \(m>0\) put \(\chi_m=\theta(e^{u/m})\) where
\(W>0\) and \(\chi_m=0\) on \(A\). Then \(\chi_m\) is smooth on \(M\), \(0\le\chi_m\le1\), \(\chi_m=0\) on the open neighbourhood \(\{W<4^{-m}\}\) of \(A\),
\(\chi_m\to1\) pointwise on \(M\setminus A\) as \(m\to\infty\) through the integers, and, in every chart, \(L_{\chi_m}(v)\ge-\frac{2C}m\,\omega(v,Jv)\); that is,
\[
 dd^c\chi_m+\frac{2C}{\pi m}\,\omega\ \ge\ 0 .
\]
\end{lemma}
\begin{proof}
\leanok
Where \(0<W<4^{-m}\), \(u/m\le-\log2\), so \(e^{u/m}\le\frac12\) and \(\chi_m=0\); hence \(\chi_m=0\) on the open set \(\{W<4^{-m}\}\), which contains \(A=\{W=0\}\).
On the open set \(\{W>0\}\), \(\chi_m\) is a composition of smooth functions. These two open sets cover \(M\), so \(\chi_m\) is smooth. As \(u<0\),
\(e^{u/m}\in(0,1)\) and \(0\le\chi_m\le1\); at a point off \(A\), \(e^{u/m}\to1\) and \(\chi_m\to\theta(1)=1\). For the Levi form at a point where \(W>0\): \(u\le0\)
and \(L_u=\frac12L_{\log W}\ge-c\) with \(c=\frac C2\omega(v,Jv)\ge0\), so Lemma~\ref{ana:stokes-cutoff-pointwise} gives \(L_{\chi_m}(v)\ge-4c/m=-\frac{2C}m\omega(v,Jv)\);
near a point of \(A\), \(\chi_m=0\) and the bound is trivial. In a chart the coefficient of \(dd^c\chi_m\) on a line is \(L_{\chi_m}/\pi\).
\end{proof}

\begin{lemma}[A countable decomposition of \(V_{\mathrm{reg}}\); integrals over \(V_{\mathrm{reg}}\)]
\label{ana:stokes-regular-decomposition}
\uses{ana:regular-ball-ambient-trace,ana:piece-family-exhaustion-supremum,ana:patch-transition,ana:background-mass}
\lean{TranscendentalNakamaye.Pending.StokesRegularDecomposition.RegularDecomposition,TranscendentalNakamaye.Pending.StokesRegularDecomposition.RegularDecomposition.image_subset_regularLocus,TranscendentalNakamaye.Pending.StokesRegularDecomposition.RegularDecomposition.exists_mem_image,TranscendentalNakamaye.Pending.StokesRegularDecomposition.RegularDecomposition.partialFamily,TranscendentalNakamaye.Pending.StokesRegularDecomposition.RegularDecomposition.coveredRegion_partialFamily,TranscendentalNakamaye.Pending.StokesRegularDecomposition.RegularDecomposition.coveredRegion_partialFamily_mono,TranscendentalNakamaye.Pending.StokesRegularDecomposition.RegularDecomposition.exhausts_partialFamily,TranscendentalNakamaye.Pending.StokesRegularDecomposition.RegularDecomposition.formPieceMass_partialFamily,TranscendentalNakamaye.Pending.StokesRegularDecomposition.RegularDecomposition.iSup_formPieceMass_eq_tsum,TranscendentalNakamaye.Pending.StokesRegularDecomposition.nonempty_regularDecomposition,TranscendentalNakamaye.Pending.StokesRegularDecomposition.enorm_eq_ofReal_add_ofReal_neg,TranscendentalNakamaye.Pending.StokesRegularDecomposition.IsDensityField,TranscendentalNakamaye.Pending.StokesRegularDecomposition.IsDensityField.neg,TranscendentalNakamaye.Pending.StokesRegularDecomposition.IsDensityField.add,TranscendentalNakamaye.Pending.StokesRegularDecomposition.IsDensityField.sum,TranscendentalNakamaye.Pending.StokesRegularDecomposition.IsDensityField.mul_comp,TranscendentalNakamaye.Pending.StokesRegularDecomposition.RegularDecomposition.setLIntegral_piece_eq,TranscendentalNakamaye.Pending.StokesRegularDecomposition.RegularDecomposition.tsum_setLIntegral_eq,TranscendentalNakamaye.Pending.StokesRegularDecomposition.RegularDecomposition.integrableOn_piece_and_hasSum,TranscendentalNakamaye.Pending.StokesRegularDecomposition.RegularDecomposition.measure,TranscendentalNakamaye.Pending.StokesRegularDecomposition.RegularDecomposition.aestronglyMeasurable_measure_iff,TranscendentalNakamaye.Pending.StokesRegularDecomposition.RegularDecomposition.lintegral_measure,TranscendentalNakamaye.Pending.StokesRegularDecomposition.RegularDecomposition.integrable_measure_iff,TranscendentalNakamaye.Pending.StokesRegularDecomposition.RegularDecomposition.hasSum_integral_measure,TranscendentalNakamaye.Pending.StokesRegularDecomposition.RegularDecomposition.integrable_measure_and_integral_eq,TranscendentalNakamaye.Pending.StokesRegularDecomposition.patchMixedDensity,TranscendentalNakamaye.Pending.StokesRegularDecomposition.patchMixedDensity_const,TranscendentalNakamaye.Pending.StokesRegularDecomposition.isDensityField_patchMixedDensity,TranscendentalNakamaye.Pending.StokesRegularDecomposition.continuous_actualMixedTopDensity,TranscendentalNakamaye.Pending.StokesRegularDecomposition.continuousOn_patchCoordinateForm,TranscendentalNakamaye.Pending.StokesRegularDecomposition.continuousOn_patchMixedDensity}
\leanok
Let \(V\) be an irreducible analytic subvariety of dimension \(k\) of a complex manifold \(X\). For smooth real two-forms \(a_1,\dots,a_k\) on \(X\) and a
regular patch \(\gamma\colon\Omega\to V\) write \(D_\gamma(a_1,\dots,a_k)\) for the density, in the parameters, of \(\gamma^*a_1\wedge\dots\wedge\gamma^*a_k\)
(the mixed density of the pulled-back coordinate forms), and \(\rho_\gamma(a)=D_\gamma(a,\dots,a)\).
\begin{enumerate}
\item There are regular patches \(\gamma_j\colon\Omega_j\to V\) (\(j\in\mathbf N\)) and measurable \(E_j\subseteq\Omega_j\) whose images \(\gamma_j(E_j)\) are pairwise
      disjoint with union \(V_{\mathrm{reg}}\), such that every compact subset of \(V_{\mathrm{reg}}\) lies in \(\gamma_0(E_0)\cup\dots\cup\gamma_{N-1}(E_{N-1})\) for
      some \(N\). The first \(N\) pieces form a regular-piece family \(G_N\), and \((G_N)_N\) is an exhaustion of \(V_{\mathrm{reg}}\).
\item For every smooth real two-form \(a\) with \(\rho_{\gamma_j}(a)\ge0\) on \(E_j\) for all \(j\),
      \(\sup_F\mathcal M_a(F)=\sum_j\int_{E_j}\rho_{\gamma_j}(a)\).
\item (Density fields) Call \(Y\) a \emph{density field} if it assigns to every regular patch \(\gamma'\) a function \(Y_{\gamma'}\) on its
      parameter domain with \(Y_{\gamma'}(w)=\det_{\mathbf R}D\tau(w)\,Y_{\gamma''}(\tau(w))\) for the transition \(\tau\) from \(\gamma'\) to \(\gamma''\).
      Mixed densities \(D_{\gamma'}(a_1,\dots,a_k)\) of smooth two-forms are density fields, continuous on the parameter domain; sums,
      negatives and products \((g\circ\gamma')\,Y_{\gamma'}\) with a function \(g\) on \(X\) of density fields are density fields.
\item (Transfer) Let \(\mu\) be the measure on \(\mathbf N\times\mathbf C^k\) with \(\int F\,d\mu=\sum_j\int_{E_j}F(j,\cdot)\). Let \(Y\) be a density field,
      \(\gamma\colon\Omega\to V\) a regular patch and \(\Omega'\subseteq\Omega\) measurable, such that \(Y_{\gamma_j}(w)=0\) whenever \(w\in E_j\) and
      \(\gamma_j(w)\notin\gamma(\Omega')\), \(Y_{\gamma_j}\) is measurable on \(E_j\) for every \(j\), and \(Y_\gamma\) is integrable on \(\Omega'\). Then
      \((j,w)\mapsto Y_{\gamma_j}(w)\) is \(\mu\)-integrable (so \(\sum_j\int_{E_j}|Y_{\gamma_j}|<\infty\)) and
      \(\sum_j\int_{E_j}Y_{\gamma_j}=\int Y\,d\mu=\int_{\Omega'}Y_\gamma\).
\end{enumerate}
\end{lemma}
\begin{proof}
\leanok
(1) Every point \(x\in V_{\mathrm{reg}}\) has a regular patch \(\gamma_x\colon\Omega_x\to V_{\mathrm{reg}}\) and an open \(U_x\ni x\) of \(X\) with
\(V\cap U_x\subseteq\gamma_x(\Omega_x)\) (Lemma~\ref{ana:regular-ball-ambient-trace}). \(X\) is second countable, so countably many \(U_j=U_{x_j}\) cover
\(V_{\mathrm{reg}}\). Put \(E_j=\{w\in\Omega_j:\gamma_j(w)\in U_j\setminus\bigcup_{i<j}U_i\}\), measurable since \(\gamma_j\) is continuous on the open
\(\Omega_j\). Then \(\gamma_j(E_j)=(V\cap U_j)\setminus\bigcup_{i<j}U_i\): these sets are pairwise disjoint, their first \(N\) have union
\(V\cap\bigcup_{i<N}U_i\), and their union is \(V_{\mathrm{reg}}\). A compact subset of \(V_{\mathrm{reg}}\) is covered by finitely many \(U_i\).
(2) \(\mathcal M_a(G_N)=\sum_{j<N}\int_{E_j}\rho_{\gamma_j}(a)\) increases to \(\sup_F\mathcal M_a(F)\) along the exhaustion
(Lemma~\ref{ana:piece-family-exhaustion-supremum}).
(3) The mixed density transforms under a real-linear change of variables by its determinant (as the top-power density does,
Lemma~\ref{ana:patch-transition}), and the pulled-back coordinate forms transform by the derivative of the transition; continuity is that of a polynomial in
the coefficients and the derivative of the parametrization. (4) For each \(j\), \(Y_{\gamma_j}\) vanishes off the set \(S_j\) of parameters \(w\in E_j\) with \(\gamma_j(w)\in\gamma(\Omega')\), on which the transition
\(\tau_j\) to \(\gamma\) is defined, holomorphic and injective, with \(\gamma\circ\tau_j=\gamma_j\) (Lemma~\ref{ana:patch-transition}). By the change of
variables formula and the transformation law (\(\det_{\mathbf R}D\tau_j=|\det_{\mathbf C}D\tau_j|^2\ge0\)), \(\int_{S_j}|Y_{\gamma_j}|=\int_{\tau_j(S_j)}|Y_\gamma|\) and
\(\int_{S_j}Y_{\gamma_j}=\int_{\tau_j(S_j)}Y_\gamma\). The sets \(\tau_j(S_j)\) are measurable, pairwise disjoint (their images \(\gamma(\tau_j(S_j))\subseteq\gamma_j(E_j)\)
are, and \(\gamma\) is injective) and cover \(\Omega'\) (every point of \(\gamma(\Omega')\subseteq V_{\mathrm{reg}}\) lies in some \(\gamma_j(E_j)\)). Summing over
\(j\) gives both claims, first for the nonnegative and nonpositive parts of \(Y\) (themselves density fields, \(\det_{\mathbf R}D\tau_j\ge0\)), then for \(Y\).
\end{proof}

\begin{lemma}[Integration by parts on \(V_{\mathrm{reg}}\) against a cutoff vanishing near \(V_{\mathrm{sing}}\)]
\label{ana:stokes-patch-ibp}
\uses{ana:stokes-regular-decomposition,bts:slot-ibp,ana:stokes-defect-mixed,ana:stokes-defect-ddc,fpmma:geometric-density,hgeom:chart-partition,ana:regular-ball-ambient-trace}
\source{demailly-cadg:III.2.7,demailly-cadg:III.3}
\lean{TranscendentalNakamaye.Pending.StokesPatchIBP.HasPatchPotentials,TranscendentalNakamaye.Pending.StokesPatchIBP.isOpen_parameterDomain,TranscendentalNakamaye.Pending.StokesPatchIBP.contDiffOn_comp_toManifold,TranscendentalNakamaye.Pending.StokesPatchIBP.hasPatchPotentials_closed,TranscendentalNakamaye.Pending.StokesPatchIBP.hasPatchPotentials_ddcForm,TranscendentalNakamaye.Pending.StokesPatchIBP.HasPatchPotentials.add,TranscendentalNakamaye.Pending.StokesPatchIBP.patchCoordinateForm_update,TranscendentalNakamaye.Pending.StokesPatchIBP.patchMixedDensity_update_ddcForm_eq,TranscendentalNakamaye.Pending.StokesPatchIBP.mixedDensity_update_eq_zero,TranscendentalNakamaye.Pending.StokesPatchIBP.setIntegral_patch_ibp,TranscendentalNakamaye.Pending.StokesPatchIBP.smoothDDC_finsetSum,TranscendentalNakamaye.Pending.StokesPatchIBP.patchMixedDensity_update_ddcForm_eq_zero,TranscendentalNakamaye.Pending.StokesPatchIBP.patchMixedDensity_update_ddcForm_sum,TranscendentalNakamaye.Pending.StokesPatchIBP.continuousOn_mul_patchMixedDensity,TranscendentalNakamaye.Pending.StokesPatchIBP.exists_ibpChart,TranscendentalNakamaye.Pending.StokesPatchIBP.ae_mem_piece,TranscendentalNakamaye.Pending.StokesPatchIBP.toManifold_injOn,TranscendentalNakamaye.Pending.StokesPatchIBP.integral_measure_ibp}
\leanok
Let \(V\subset X\) be an irreducible analytic subvariety of dimension \(k\) of a compact complex manifold, with the decomposition
\((\gamma_j,E_j)_j\) and the measure \(\mu\) of Lemma~\ref{ana:stokes-regular-decomposition}. Let \(\chi\) and \(w\) be smooth functions on \(X\) with
\(\operatorname{supp}\chi\cap V\subseteq V_{\mathrm{reg}}\) (for instance \(\chi=0\) on a neighbourhood of \(V_{\mathrm{sing}}\)), \(i_0\) a slot, and \(a_i\) (\(i\ne i_0\)) smooth
two-forms on \(X\) whose pull-backs have smooth local potentials in the parameters of every regular patch (closed real \((1,1)\)-forms,
\(dd^c\) of smooth functions, and their sums). Then \((j,u)\mapsto\chi(\gamma_j(u))\,D_{\gamma_j}(a[i_0\mapsto dd^cw])(u)\) and
\((j,u)\mapsto w(\gamma_j(u))\,D_{\gamma_j}(a[i_0\mapsto dd^c\chi])(u)\) are \(\mu\)-integrable and
\[
 \sum_j\int_{E_j}(\chi\circ\gamma_j)\,D_{\gamma_j}(a[i_0\mapsto dd^cw])\ =\ \sum_j\int_{E_j}(w\circ\gamma_j)\,D_{\gamma_j}(a[i_0\mapsto dd^c\chi]).
\]
That is, \(\int_{V_{\mathrm{reg}}}\chi\,dd^cw\wedge\eta=\int_{V_{\mathrm{reg}}}w\,dd^c\chi\wedge\eta\) for \(\eta=\bigwedge_{i\ne i_0}a_i\).
\end{lemma}
\begin{proof}
\leanok
\(K=\operatorname{supp}\chi\cap V\) is compact and contained in \(V_{\mathrm{reg}}\). At \(x\in K\) take the regular patch \(\gamma_x\) centred at \(x\), an open
\(\Omega_x\ni0\) of its parameters on which the slots \(a_i\), \(i\ne i_0\), have smooth potentials \(\varphi_{x,i}\), a ball \(B_x=B(0,r_x)\) with
\(\overline B_x\subseteq\Omega_x\), and an open \(U_x\ni x\) of \(X\) with \(V\cap U_x\subseteq\gamma_x(B_x)\) (Lemma~\ref{ana:regular-ball-ambient-trace}). The sets
\(U_x\) and \(X\setminus K\) cover \(X\); Lemma~\ref{hgeom:chart-partition} gives smooth \(\rho_b\ge0\), finitely many, with \(\sum_b\rho_b=1\) and each
\(\operatorname{supp}\rho_b\) inside one member of the cover. Put \(g_b=\rho_b\chi\), so \(\sum_bg_b=\chi\) and, by multilinearity in the slot \(i_0\) and linearity of
\(dd^c\), \(\sum_bD(a[i_0\mapsto dd^cg_b])=D(a[i_0\mapsto dd^c\chi])\). If \(\operatorname{supp}\rho_b\subseteq X\setminus K\), \(g_b\) vanishes near every point of \(V\), and
both terms vanish. If \(\operatorname{supp}\rho_b\subseteq U_x\), then a point \(\gamma_x(u)\) with \(u\) in the parameter domain and \(\gamma_x(u)\in\operatorname{supp}\rho_b\)
has \(u\in B_x\) (injectivity of \(\gamma_x\)), so \(g_b\circ\gamma_x\) vanishes on \(\Omega_x\) off the compact \(\overline B_x\); in the parameters of \(\gamma_x\) the two
densities are the mixed densities of the complex Hessians of \(\varphi_x[i_0\mapsto w\circ\gamma_x]\) and \(\varphi_x[i_0\mapsto g_b\circ\gamma_x]\)
(Lemma~\ref{fpmma:geometric-density}; \(dd^c\) commutes with the holomorphic \(\gamma_x\), Lemma~\ref{ana:stokes-defect-ddc}), and the slot integration by parts
(Lemma~\ref{bts:slot-ibp}) gives \(\int_{\Omega_x}(g_b\circ\gamma_x)\,D(a[i_0\mapsto dd^cw])=\int_{\Omega_x}(w\circ\gamma_x)\,D(a[i_0\mapsto dd^cg_b])\). Both integrands are
density fields (Lemma~\ref{ana:stokes-regular-decomposition}(3)) vanishing at the parameters of every patch whose image lies off \(\gamma_x(\Omega_x)\) (such points of
\(V\) lie off \(\operatorname{supp}\rho_b\)), continuous, and integrable on \(\Omega_x\) (continuous on \(\overline B_x\), zero on \(\Omega_x\setminus\overline B_x\)), so by
Lemma~\ref{ana:stokes-regular-decomposition}(4) each side is the corresponding \(\mu\)-integral. Summing over \(b\) gives the identity.
\end{proof}

\begin{lemma}[The limit step of the one-sided probe]
\label{ana:stokes-probe-limit}
\source{demailly-cadg:III.2.3}
\lean{TranscendentalNakamaye.Pending.StokesProbeLimit.abs_integral_mul_le,TranscendentalNakamaye.Pending.StokesProbeLimit.integral_eq_zero_of_probe}
\leanok
Let \(\mu\) be a measure, \(f\) and \(b\ge0\) integrable, \(v\) measurable with \(lo\le v\le hi\), and for \(m\ge0\) let \(\chi_m\) be measurable with
\(0\le\chi_m\le1\) and \(\chi_m\to1\) almost everywhere, \(q_m\) integrable and \(\epsilon_m\ge0\) with \(\epsilon_m\to0\) and \(q_m+\epsilon_mb\ge0\). If
\(\int\chi_mf\,d\mu=\int(v-c)\,q_m\,d\mu\) for every \(m\) and every constant \(c\), then \(|\int\chi_mf\,d\mu|\le\epsilon_m(hi-lo)\int b\,d\mu\) for every
\(m\), and \(\int f\,d\mu=0\).
\end{lemma}
\begin{proof}
\leanok
Write \(\int(v-c)q_m=\int(v-c)(q_m+\epsilon_mb)-\epsilon_m\int(v-c)b\). For \(c=hi\) the first integral is \(\le0\) and
\(-\int(v-hi)b\le(hi-lo)\int b\); for \(c=lo\) the first integral is \(\ge0\) and \(\int(v-lo)b\le(hi-lo)\int b\). This gives the bound, which tends to
\(0\); and \(\int\chi_mf\to\int f\) by dominated convergence, \(|\chi_mf|\le|f|\).
\end{proof}

\begin{theorem}[Stokes on the regular locus: a smooth potential does not change the mass of \(V_{\mathrm{reg}}\)]
\label{ana:stokes-regular-locus-identity}
\uses{ana:stokes-probe-limit,ana:stokes-singular-locus-analytic,ana:stokes-pole-potential,ana:stokes-cutoffs,ana:stokes-patch-ibp,ana:stokes-regular-decomposition,ana:background-mass-finite,ana:background-mass,gnpp:kahler-domination}
\source{demailly-cadg:III.2.7,demailly-cadg:III.2.2,demailly-cadg:III.2.3}
\lean{TranscendentalNakamaye.Pending.StokesRegularLocusIdentity.teleFamily,TranscendentalNakamaye.Pending.StokesRegularLocusIdentity.isSemipositiveOneOne_teleFamily,TranscendentalNakamaye.Pending.StokesRegularLocusIdentity.actualMixedTopDensity_mono,TranscendentalNakamaye.Pending.StokesRegularLocusIdentity.actualMixedTopDensity_update_add_smul,TranscendentalNakamaye.Pending.StokesRegularLocusIdentity.isSemipositiveOneOne_smul_sub,TranscendentalNakamaye.Pending.StokesRegularLocusIdentity.sum_teleFamily_nonneg,TranscendentalNakamaye.Pending.StokesRegularLocusIdentity.sum_teleFamily_le,TranscendentalNakamaye.Pending.StokesRegularLocusIdentity.actualTopPowerDensity_add_sub_eq_sum_teleFamily,TranscendentalNakamaye.Pending.StokesRegularLocusIdentity.coordinateParametrization_mem_target,TranscendentalNakamaye.Pending.StokesRegularLocusIdentity.isSemipositiveOneOne_patchCoordinateForm,TranscendentalNakamaye.Pending.StokesRegularLocusIdentity.patchCoordinateForm_apply_line,TranscendentalNakamaye.Pending.StokesRegularLocusIdentity.exists_smul_sub_isSemipositiveOneOne,TranscendentalNakamaye.Pending.StokesRegularLocusIdentity.patchCoordinateForm_add_smul,TranscendentalNakamaye.Pending.StokesRegularLocusIdentity.patchCoordinateForm_ddcForm_sub_const,TranscendentalNakamaye.Pending.StokesRegularLocusIdentity.slots,TranscendentalNakamaye.Pending.StokesRegularLocusIdentity.patchMixedDensity_slots,TranscendentalNakamaye.Pending.StokesRegularLocusIdentity.teleLin,TranscendentalNakamaye.Pending.StokesRegularLocusIdentity.teleLin_apply,TranscendentalNakamaye.Pending.StokesRegularLocusIdentity.integral_indicator_patchFormDensity,TranscendentalNakamaye.Pending.StokesRegularLocusIdentity.iSup_formPieceMass_currentForm_eq}
\leanok
Let \((X,\beta)\) be a compact K\"ahler manifold, \(V\subset X\) an irreducible analytic subvariety of dimension \(k\ge1\) and \(v\) a smooth
function with \(\beta+dd^cv\ge0\). Then
\(\int_{V_{\mathrm{reg}}}(\beta+dd^cv)^k=\int_{V_{\mathrm{reg}}}\beta^k\), both sides finite.
(In Lean \(v\) is the potential of an admissible presentation of the zero class with representative \(0\), the two sides are the
regular-piece suprema of Lemma~\ref{ana:background-mass}, and \(k\ge1\) is not needed; the right side is finite by
Lemma~\ref{ana:background-mass-finite}.)
\end{theorem}
\begin{proof}
\leanok
Put \(S=\sum_{i<k}(\beta+dd^cv)^i\wedge\beta^{k-1-i}\), so that \((\beta+dd^cv)^k-\beta^k=dd^cv\wedge S\); \(S\) is a sum of products of
smooth closed semipositive forms, and \(\beta+dd^cv\le C_v\beta\) on \(X\), so \(0\le\beta\wedge S\le kC_v^{k-1}\beta^k\) on
\(V_{\mathrm{reg}}\) and \(\int_{V_{\mathrm{reg}}}\beta\wedge S<\infty\) by Lemma~\ref{ana:background-mass-finite}; likewise \(|dd^cv\wedge S|\) is
dominated by a multiple of \(\beta^k\).

\(V_{\mathrm{sing}}\) is a closed analytic subset (Lemma~\ref{ana:stokes-singular-locus-analytic}); let \(u\), \(C\) and \(\chi_m\) be as in
Lemmas~\ref{ana:stokes-pole-potential} and~\ref{ana:stokes-cutoffs} for \(A=V_{\mathrm{sing}}\) and \(\omega=\beta\), and put
\(\epsilon_m=2C/(\pi m)\), so that \(dd^c\chi_m+\epsilon_m\beta\ge0\). Lemma~\ref{ana:stokes-patch-ibp}, applied to \(\chi=\chi_m\) and \(w=v-c\) for a constant \(c\) (the \(\gamma_i\)
running through the factors of each term of \(S\)), gives
\(\int_{V_{\mathrm{reg}}}\chi_m\,dd^cv\wedge S=\int_{V_{\mathrm{reg}}}(v-c)\,dd^c\chi_m\wedge S\). On \(V_{\mathrm{reg}}\) the form
\((dd^c\chi_m+\epsilon_m\beta)\wedge S\) is a nonnegative top form. With \(c=\sup v\), \(v-c\le0\), so
\(\int(v-c)\,dd^c\chi_m\wedge S\le-\epsilon_m\int(v-c)\beta\wedge S\le\epsilon_m\operatorname{osc}(v)\int\beta\wedge S\); with \(c=\inf v\) the reverse
inequality holds. Hence
\[
 \Bigl|\int_{V_{\mathrm{reg}}}\chi_m\,dd^cv\wedge S\Bigr|\ \le\ \epsilon_m\,\operatorname{osc}(v)\int_{V_{\mathrm{reg}}}\beta\wedge S\ \longrightarrow\ 0 .
\]
As \(V_{\mathrm{sing}}\cap V_{\mathrm{reg}}=\emptyset\), \(\chi_m\to1\) at every point of \(V_{\mathrm{reg}}\), and \(|\chi_m\,dd^cv\wedge S|\) is dominated by an
integrable form; by dominated convergence \(\int_{V_{\mathrm{reg}}}dd^cv\wedge S=0\). (This last paragraph is Lemma~\ref{ana:stokes-probe-limit}
for the measure \(\mu\) on \(\mathbf N\times\mathbf C^k\) with \(\int F\,d\mu=\sum_j\int_{E_j}F(j,\cdot)\), \((\gamma_j,E_j)_j\) the decomposition of
Lemma~\ref{ana:stokes-regular-decomposition}, and \(f(j,\cdot)\), \(b(j,\cdot)\), \(q_m(j,\cdot)\) the densities \(D_{\gamma_j}\) of \(dd^cv\wedge S\), \(\beta\wedge S\),
\(dd^c\chi_m\wedge S\), \(\chi_m(j,\cdot)=\chi_m\circ\gamma_j\), \(v(j,\cdot)=v\circ\gamma_j\), \(lo=\inf v\), \(hi=\sup v\); by
Lemma~\ref{ana:stokes-regular-decomposition}(2), \(\int f\,d\mu=\int_{V_{\mathrm{reg}}}(\beta+dd^cv)^k-\int_{V_{\mathrm{reg}}}\beta^k\).)
\end{proof}
This is the argument of Remark~\ref{ana:stokes-one-sided-probe}: its two prerequisites are Lemma~\ref{ana:stokes-patch-ibp} (Stokes for forms
vanishing near \(V_{\mathrm{sing}}\), here in the Hessian form the project already has in the parameters) and the cutoffs of
Lemma~\ref{ana:stokes-cutoffs}, whose poles lie exactly on the analytic set \(V_{\mathrm{sing}}\) of Lemma~\ref{ana:stokes-singular-locus-analytic}.
No current on \(X\), no Skoda--El Mir extension and no Cauchy--Schwarz inequality is used.

\begin{lemma}[Local weighted-pushforward data exist for every finite family]
\label{ana:stokes-weighted-pairing-data}
\uses{ana:stokes-defect-good-exhaustion-psh,ana:background-mass-finite}
\lean{TranscendentalNakamaye.Pending.StokesWeightedPairing.exhaustFun,TranscendentalNakamaye.Pending.StokesWeightedPairing.norm_le_exhaustFun,TranscendentalNakamaye.Pending.StokesWeightedPairing.continuous_infDist_inv_on,TranscendentalNakamaye.Pending.StokesWeightedPairing.continuousOn_exhaustFun,TranscendentalNakamaye.Pending.StokesWeightedPairing.measurable_exhaustFun,TranscendentalNakamaye.Pending.StokesWeightedPairing.isCompact_exhaustFun_le,TranscendentalNakamaye.Pending.StokesWeightedPairing.ramp,TranscendentalNakamaye.Pending.StokesWeightedPairing.continuous_ramp,TranscendentalNakamaye.Pending.StokesWeightedPairing.ramp_nonneg,TranscendentalNakamaye.Pending.StokesWeightedPairing.ramp_le_one,TranscendentalNakamaye.Pending.StokesWeightedPairing.ramp_eq_zero,TranscendentalNakamaye.Pending.StokesWeightedPairing.ramp_eq_one,TranscendentalNakamaye.Pending.StokesWeightedPairing.weightProfile,TranscendentalNakamaye.Pending.StokesWeightedPairing.weightTerm_eq_zero,TranscendentalNakamaye.Pending.StokesWeightedPairing.weightTerm_nonneg,TranscendentalNakamaye.Pending.StokesWeightedPairing.weightTerm_eq_zero_of_ge,TranscendentalNakamaye.Pending.StokesWeightedPairing.summable_weightTerm,TranscendentalNakamaye.Pending.StokesWeightedPairing.one_le_weightProfile,TranscendentalNakamaye.Pending.StokesWeightedPairing.continuous_weightProfile,TranscendentalNakamaye.Pending.StokesWeightedPairing.le_weightProfile,TranscendentalNakamaye.Pending.StokesWeightedPairing.ofReal_weightProfile_le,TranscendentalNakamaye.Pending.StokesWeightedPairing.exists_properWeight,TranscendentalNakamaye.Pending.StokesWeightedPairing.parameterNbhd,TranscendentalNakamaye.Pending.StokesWeightedPairing.parameterOf,TranscendentalNakamaye.Pending.StokesWeightedPairing.isOpen_parameterNbhd,TranscendentalNakamaye.Pending.StokesWeightedPairing.continuousOn_parameterOf,TranscendentalNakamaye.Pending.StokesWeightedPairing.toManifold_mem_parameterNbhd,TranscendentalNakamaye.Pending.StokesWeightedPairing.parameterOf_toManifold,TranscendentalNakamaye.Pending.StokesWeightedPairing.toManifold_parameterOf,TranscendentalNakamaye.Pending.StokesWeightedPairing.exists_continuousMap_toManifold_eq,TranscendentalNakamaye.Pending.StokesWeightedPairing.exists_localWeightedPatchMeasure_test,TranscendentalNakamaye.Pending.StokesWeightedPairing.continuousOn_telescopeIntegrand,TranscendentalNakamaye.Pending.StokesWeightedPairing.nonempty_finiteTelescopeLocalWeightedPairingData,TranscendentalNakamaye.Pending.StokesWeightedPairing.finiteTelescopePairingRemainderData_of_integrableOn,TranscendentalNakamaye.Pending.StokesWeightedPairing.finiteTelescopePairingRemainderData_of_integrableOn_remainder,TranscendentalNakamaye.Pending.StokesWeightedPairing.integrableOn_telescopeIntegrand,TranscendentalNakamaye.Pending.StokesWeightedPairing.finiteTelescopePairingRemainderData_of_finite,TranscendentalNakamaye.Pending.StokesWeightedPairing.finiteTelescopePairingRemainderData_of_finite_remainder}
\leanok
Let \(V\) be an irreducible analytic subvariety of a compact complex manifold \(X\), \(G\) a finite regular-piece family of \(V\) with pieces
\(D_i\subseteq\Omega_i\), where \(\Omega_i\) is an open subset of the parameter domain of the patch \(\gamma_i\) on which the telescope integrand \(f_i\) of
Lemma~\ref{ana:stokes-defect-good-exhaustion-psh} is defined (its potentials are smooth there), and suppose \(f_i\) is integrable on \(D_i\) --- for
instance because both masses \(\mathcal M_{\beta+dd^cv}(G)\) and \(\mathcal M_\beta(G)\) are finite, the two densities being nonnegative on the pieces.
Then for every \(i\) there are a measurable density \(r_i\colon D_i\to[0,\infty)\) with \(\int_{D_i}r_i<\infty\) and a continuous function \(t_i\) on \(X\) with
\(f_i=r_i\,(t_i\circ\gamma_i)\) on \(D_i\); so the finite weighted push-forwards \(\gamma_{i*}(r_i\,d\lambda|_{D_i})\) paired with \(t_i\) reproduce
\(\int_{D_i}f_i\), and \(R=\bigl|\sum_i\int_{D_i}f_i\bigr|=|\mathcal M_{\beta+dd^cv}(G)-\mathcal M_\beta(G)|\) bounds the sum of the pairings.
\end{lemma}
\begin{proof}
\leanok
Fix \(i\) and drop it; \(f\) is continuous on \(\Omega\) (its potentials are smooth there) and integrable on \(D\).
(1) \emph{An exhaustion function.} \(g(w)=d(w,\mathbf C^k\setminus\Omega)^{-1}+|w|\) (the first term read as \(0\) if \(\Omega=\mathbf C^k\)) is continuous on
\(\Omega\), and \(\{w\in\Omega:g(w)\le c\}\) is closed --- on it the distance to the complement is at least \(c^{-1}\) --- and bounded, hence compact.
(2) \emph{A proper integrable weight.} Let \(u=(1+|w|)^{-r}+|f|\) with \(r>2k\), so that \(u>0\) is continuous on \(\Omega\) and integrable on \(D\), and
let \(\mu=u\,d\lambda|_D\), a finite measure. As \(\mu\{g>T\}\to0\) when \(T\to\infty\), there are thresholds \(t_l\ge l\), increasing, with
\(\mu\{g>t_l\}\le4^{-l}\). Put \(\Phi(x)=1+\sum_l2^l\,\mathrm{clamp}(x-t_l)\), \(\mathrm{clamp}(y)=\max(0,\min(1,y))\): locally a finite sum, so \(\Phi\) is
continuous, \(\Phi\ge1\), \(\Phi\ge2^l\) on \([t_l+1,\infty)\), and \(\int\Phi\circ g\,d\mu\le\mu(\mathbf C^k)+\sum_l2^l4^{-l}<\infty\). So \(\psi=\Phi\circ g\) is
continuous on \(\Omega\), \(\int_D\psi u<\infty\), and \(\psi\ge L\) off the compact set \(\{g\le t_l+1\}\subseteq\Omega\) once \(2^l\ge L\).
(3) \emph{The test function.} \(h=f/(\psi u)\) is continuous on \(\Omega\) and \(|h|\le1/\psi\), so \(h\to0\) off the compact subsets of \(\Omega\).
In the chart \(e\) of the patch, with \(\Phi_\gamma\) its flattening and \(\mathrm{pr}\) the projection onto the first \(k\) coordinates,
\(y\mapsto\mathrm{pr}\,\Phi_\gamma(e(y))\) is continuous on an open set \(N\) of \(X\) containing \(\gamma(\Omega)\) and inverts \(\gamma\) on \(V\cap N\). Let
\(t_0=h\circ\gamma^{-1}\) on \(\gamma(\Omega)\) and \(t_0=0\) on \(V\setminus\gamma(\Omega)\). It is continuous on \(V\): near a point of \(\gamma(\Omega)\) it is
\(h\circ\mathrm{pr}\circ\Phi_\gamma\circ e\); a point \(y\in V\setminus\gamma(\Omega)\) lies off the compact set \(\gamma(K)\), \(K\subseteq\Omega\) compact with
\(|h|\le\varepsilon\) off \(K\), and near \(y\) every value of \(t_0\) is \(0\) or a value of \(h\) off \(K\). \(V\) is closed in the compact (hence normal)
\(X\), so Tietze extends \(t_0\) to a continuous \(t\) on \(X\), with \(t\circ\gamma=h\) on \(\Omega\).
(4) With \(r=\psi u\): \(\int_Dr<\infty\), and \(r\,(t\circ\gamma)=\psi u\cdot f/(\psi u)=f\) on \(D\). Summing over \(i\), the sum of the pairings is
\(\sum_i\int_{D_i}f_i\), which is \(\mathcal M_{\beta+dd^cv}(G)-\mathcal M_\beta(G)\) when the two masses are finite.
\end{proof}

\begin{corollary}[The collar remainder tends to \(0\) along an exhaustion]
\label{ana:stokes-exhaustion-remainder}
\uses{ana:stokes-regular-locus-identity,ana:stokes-weighted-pairing-data,ana:piece-family-exhaustion-supremum,ana:background-mass-finite,stokesred:current-mass-finite,stokesred:density-nonneg,ana:stokes-defect-good-exhaustion-psh,ana:stokes-defect-good-exhaustion}
\lean{TranscendentalNakamaye.Pending.StokesRegularLocusIdentity.tendsto_abs_formPieceMass_sub_of_exhaustion,TranscendentalNakamaye.Pending.StokesRegularLocusIdentity.pairingRemainder_of_exhaustion,TranscendentalNakamaye.Pending.StokesRegularLocusIdentity.pairingRemainder_of_exhaustion_remainder,TranscendentalNakamaye.Pending.StokesRegularLocusIdentity.tendsto_pairingRemainder_of_exhaustion,TranscendentalNakamaye.Pending.StokesRegularLocusIdentity.pairingRemainderOfTelescope,TranscendentalNakamaye.Pending.StokesRegularLocusIdentity.pairingRemainderOfTelescope_remainder,TranscendentalNakamaye.Pending.StokesRegularLocusIdentity.tendsto_pairingRemainderOfTelescope}
\leanok
Let \(X\), \(\beta\), \(V\), \(v\) be as in Theorem~\ref{ana:stokes-regular-locus-identity}, and let \((H_m)_m\) be finite regular-piece families of \(V\)
whose pieces carry the certified local potentials of Lemma~\ref{ana:stokes-defect-good-exhaustion-psh}, with covered regions increasing in \(m\)
and every compact subset of \(V_{\mathrm{reg}}\) contained in the covered region of some \(H_m\). For every \(m\) let the local weighted-pushforward
data of \(H_m\) be those of Lemma~\ref{ana:stokes-weighted-pairing-data}, with remainder
\(R_m=|\mathcal M_{\beta+dd^cv}(H_m)-\mathcal M_\beta(H_m)|\). Then \(R_m\to0\). The same holds for any choice of the telescope data of the
families \(H_m\) (in particular for the uncertified data of Lemma~\ref{ana:stokes-defect-good-exhaustion}).
\end{corollary}
\begin{proof}
\leanok
Both masses of \(H_m\) are finite (Lemmas~\ref{stokesred:current-mass-finite} and~\ref{ana:background-mass-finite}) and both densities are
nonnegative on the pieces (Lemma~\ref{stokesred:density-nonneg}), so Lemma~\ref{ana:stokes-weighted-pairing-data} applies to every \(H_m\). Along
the exhaustion \(\mathcal M_a(H_m)\to\sup_F\mathcal M_a(F)\) for \(a=\beta+dd^cv\) and \(a=\beta\)
(Lemma~\ref{ana:piece-family-exhaustion-supremum}); the two suprema are equal by Theorem~\ref{ana:stokes-regular-locus-identity} and finite by
Lemma~\ref{ana:background-mass-finite}. So the two real sequences \(\mathcal M_a(H_m)\) converge to the same real number and \(R_m\to0\).
\end{proof}
This is the proof of Lemma~\ref{ana:stokes-certified-input} below.

\begin{lemma}[The closed-current collar remainder along the exhaustion]
\label{ana:stokes-certified-input}
\uses{ana:stokes-defect-good-exhaustion-psh,ana:stokes-obligation-potential-free,ana:background-mass-finite,ana:piece-family-exhaustion-supremum}
\source{demailly-cadg:III.2.7,demailly-cadg:III.2.3,demailly-cadg:III.2.2}
\lean{TranscendentalNakamaye.Pending.CertifiedClosedCurrentStokes.CertifiedFinitePairingRemainderData,
TranscendentalNakamaye.Pending.CertifiedClosedCurrentStokes.certifiedFinitePairingRemainderData_of_closed_integrationCurrent,
TranscendentalNakamaye.Pending.CertifiedClosedCurrentStokes.formPieceMassCollarError_of_certified,
TranscendentalNakamaye.Pending.CertifiedClosedCurrentStokes.formPieceMassCommonLimit_of_certified,
TranscendentalNakamaye.Pending.CertifiedClosedCurrentStokes.iSup_formPieceMass_global_eq_of_certified}
\leanok
Along the exhaustion of Lemma~\ref{ana:stokes-defect-good-exhaustion-psh}, the
positive-dimensional form-current argument supplies, for every \(m\) and every piece of
\(H_m\), a finite local weighted-pushforward measure and a test function reproducing that
piece's telescope integrand, together with a real remainder \(R_m\) bounding the finite sum
of the resulting pairings, with \(R_m\to0\). Consequently the finite-piece masses of \(\beta+dd^cv\) and of \(\beta\) have a common limit along the
exhaustion, and their suprema over all regular-piece families agree.
\end{lemma}
\begin{proof}
\uses{ana:stokes-exhaustion-remainder,ana:stokes-regular-locus-identity,ana:stokes-weighted-pairing-data,ana:piece-family-exhaustion-supremum,ana:background-mass-finite}
\leanok
This is Corollary~\ref{ana:stokes-exhaustion-remainder}; in detail: let \((H_m)_m\) be the exhaustion. For every \(m\) and every piece of \(H_m\), Lemma~\ref{ana:stokes-weighted-pairing-data} supplies the local
weighted-pushforward measure and the test function, and the sum of the resulting pairings is the sum over the pieces of the integrals of the
telescope integrands, that is \(\mathcal M_{\beta+dd^cv}(H_m)-\mathcal M_\beta(H_m)\), both finite (Lemma~\ref{ana:background-mass-finite}). Take
\(R_m=|\mathcal M_{\beta+dd^cv}(H_m)-\mathcal M_\beta(H_m)|\). Along the exhaustion \(\mathcal M_a(H_m)\) increases to
\(\sup_F\mathcal M_a(F)=\int_{V_{\mathrm{reg}}}a^k\) for \(a=\beta+dd^cv\) and for \(a=\beta\) (Lemma~\ref{ana:piece-family-exhaustion-supremum}), and the two
limits are finite and equal by Theorem~\ref{ana:stokes-regular-locus-identity}. Hence \(R_m\to0\).
\end{proof}
The data are asked along the exhaustion and only there: for one fixed family the two masses genuinely differ
(Lemma~\ref{ana:background-mass-stokes-exhaustion}), so no remainder tending to \(0\) exists along a sequence of families that does not exhaust
\(V_{\mathrm{reg}}\).

\begin{remark}[The one-sided bound is not cheaper, and for smooth potentials no Cauchy--Schwarz is needed]
\label{ana:stokes-one-sided-probe}
\uses{ana:stokes-certified-input,ana:background-mass-finite,ana:kahler-restricted-mass-bound}
\source{demailly-cadg:III.2.7,demailly-cadg:III.2.3,demailly-cadg:III.2.2}
Downstream, Lemma~\ref{ana:background-mass-stokes} is consumed only through the inequality
\(\int_{V_{\mathrm{reg}}}(\beta+dd^cv)^k\le C\) for one finite \(C\)
(Lemma~\ref{ana:kahler-restricted-mass-bound}). This remark records why that one-sided form
does not avoid the content of Lemma~\ref{ana:stokes-certified-input}, and what the cheapest
argument for either form consumes. Nothing here is formalized.

Write \(S=\sum_{i<k}(\beta+dd^cv)^i\wedge\beta^{k-1-i}\), a smooth closed form on \(X\)
which is positive on \(V_{\mathrm{reg}}\), so that
\((\beta+dd^cv)^k-\beta^k=dd^cv\wedge S\). Assume:
\begin{enumerate}
\item[(P1)] \emph{Stokes for compactly supported forms on \(V_{\mathrm{reg}}\):}
\(\int_{V_{\mathrm{reg}}}d\eta=0\) for every smooth \((2k-1)\)-form \(\eta\) on \(X\) vanishing
on a neighbourhood of \(V_{\mathrm{sing}}\);
\item[(P2)] \emph{a global quasi-plurisubharmonic cutoff family:} smooth
\(\chi_m:X\to[0,1]\), each vanishing on a neighbourhood of \(V_{\mathrm{sing}}\), with
\(\chi_m\to1\) almost everywhere on \(V_{\mathrm{reg}}\) and
\(dd^c\chi_m\ge-\epsilon_m\beta\), \(\epsilon_m\to0\).
\end{enumerate}
For (P2) it suffices to have a quasi-plurisubharmonic \(u\le0\) on \(X\), smooth off a proper
analytic subset \(A\not\supseteq V\) containing \(V_{\mathrm{sing}}\), with \(u\to-\infty\) at
\(A\) and \(dd^cu\ge-C\beta\): with \(\theta\) smooth, convex and increasing, \(\theta=0\) on
\([0,1/2]\) and \(\theta(1)=1\), the functions \(\chi_m=\theta(e^{u/m})\) vanish where
\(u<-m\log2\), tend to \(1\) off \(A\), and satisfy
\(dd^c\chi_m\ge\theta'(e^{u/m})\,e^{u/m}\,dd^cu/m\ge-(C\theta'(1)/m)\,\beta\).

Since \(\chi_m\) vanishes near \(V_{\mathrm{sing}}\), two applications of (P1) give
\(\int_{V_{\mathrm{reg}}}\chi_m\,dd^cv\wedge S=\int_{V_{\mathrm{reg}}}v\,dd^c\chi_m\wedge S\) and
\(\int_{V_{\mathrm{reg}}}dd^c\chi_m\wedge S=0\); so \(v\) may be replaced by \(v-c\) for any
constant \(c\). As \(S\) is positive, \(dd^c\chi_m\wedge S+\epsilon_m\beta\wedge S\ge0\) on
\(V_{\mathrm{reg}}\). Taking \(c=\sup v\), so that \(v-c\le0\), and then \(c=\inf v\), so that
\(v-c\ge0\), gives
\[
 \Bigl|\int_{V_{\mathrm{reg}}}\chi_m\,dd^cv\wedge S\Bigr|
 \le\epsilon_m\,\bigl(\sup v-\inf v\bigr)\int_{V_{\mathrm{reg}}}\beta\wedge S ,
\]
and the last integral is finite by Lemma~\ref{ana:background-mass-finite}, because
\(\beta+dd^cv\le C_v\beta\) on \(X\). Dominated convergence, against the same finite mass,
removes \(\chi_m\). Hence \(\int_{V_{\mathrm{reg}}}dd^cv\wedge S=0\).

Three consequences.
\emph{(a)} The upper bound alone uses the choice \(c=\sup v\) and nothing less: the same (P1),
the same (P2), the same finiteness. So the one-sided statement is \emph{not} materially easier.
The constant \(\epsilon_m(\sup v-\inf v)\) depends on \(v\), but it multiplies
\(\epsilon_m\to0\), so the non-uniformity in \(v\) that a Chern--Levine--Nirenberg bound
would carry does not survive the limit.
\emph{(b)} The \emph{sharp} one-sided bound \(\int_{V_{\mathrm{reg}}}(\beta+dd^cv)^k\le
\int_{V_{\mathrm{reg}}}\beta^k\), asked for every K\"ahler background, is equivalent to the
equality: for smooth \(v\) and \(|s|\) small, \(sv\) is admissible because \(\beta\) is
K\"ahler on the compact \(X\), so the polynomial \(s\mapsto\int(\beta+s\,dd^cv)^k\) has a
maximum at \(s=0\) and its derivative \(k\int dd^cv\wedge\beta^{k-1}\) vanishes; applied to the
K\"ahler backgrounds \(\beta+dd^cw\), \(w\) small, and polarized, this kills every mixed term.
Without (P1) no argument is visible at all: a local estimate controls the negative part of
\(dd^cv\) on \(V_{\mathrm{reg}}\) by \(\beta\), but not its positive part.
\emph{(c)} For the smooth potentials this node is about, the Cauchy--Schwarz step and the
\(i\,d'd''v_k^2\ge2\,i\,d'v_k\wedge d''v_k\) step of the Skoda--El Mir argument in
Lemma~\ref{ana:stokes-certified-input} are unnecessary: a \emph{global} quasi-plurisubharmonic
cutoff and a sign suffice. That argument needs the cutoffs only locally and the current only
closed off the set; here the current is smooth on \(X\) and the cutoff is global, which is
what makes the sign usable. The prerequisites are unchanged: (P1) is the manifold Stokes
theorem absent from the pinned Mathlib, and (P2) needs an analytic \(A\supseteq
V_{\mathrm{sing}}\) with \(A\not\supseteq V\) presented by local holomorphic generators,
which the project's \(V_{\mathrm{sing}}\) (a bare set difference) does not carry.
\end{remark}

\begin{lemma}[Stokes: a smooth potential does not change the mass]
\label{ana:background-mass-stokes}
\uses{ana:background-mass,ana:background-mass-stokes-exhaustion}
\source{demailly-cadg:III.2.7,demailly-cadg:III.2.3}
\lean{TranscendentalNakamaye.Pending.SubvarietyStokesMassInterfaces.acTopMass_eq_backgroundMass}
\leanok
For every smooth function \(v\) on \(X\) with \(\beta+dd^cv\ge0\),
\[
 \int_{V_{\mathrm{reg}}}\bigl((\beta+dd^cv)|_{V_{\mathrm{reg}}}\bigr)^{k}
 =\int_{V_{\mathrm{reg}}}\beta^{k},
\]
that is, the mass of \(\beta+dd^cv\) along \(V\) is the background mass of \(V\).
\end{lemma}
\begin{proof}
\uses{ana:background-mass,ana:background-mass-stokes-exhaustion,ana:stokes-defect-density,ana:stokes-defect-good-exhaustion,ana:stokes-defect-good-exhaustion-psh,ana:stokes-certified-input,ana:background-mass-finite,fpmma:hessian-ibp}
\leanok
By Lemma~\ref{ana:background-mass-stokes-exhaustion} it is enough to prove the
identity along one exhausting sequence \((H_m)_m\) of regular-piece families.

The integrands differ by the density of Lemma~\ref{ana:stokes-defect-density}:
pointwise on \(V_{\mathrm{reg}}\),
\((\beta+dd^cv)^k-\beta^k=dd^cv\wedge S\) with
\(S=\sum_{i=0}^{k-1}(\beta+dd^cv)^i\wedge\beta^{k-1-i}\), and that lemma also shows
the difference to be a genuine density on \(V_{\mathrm{reg}}\), so it may be
integrated patch by patch along \((H_m)_m\).

By Lemma~\ref{ana:stokes-defect-good-exhaustion} the sequence may moreover be chosen so
that on each piece the defect is the mixed Monge--Amp\`ere density
\(\sum_{j<k}D_{u_{j,1},\ldots,u_{j,k}}\) of \(k\) smooth real functions of the \(k\)
parameters, the functions being a local potential \(\varphi\) of \(\beta\) read in the
parameters together with \(v\circ\gamma_{\mathrm{man}}\). For such a density,
integration by parts against a compactly supported test function -- exchanging the
Hessians of two of the functions, Lemma~\ref{fpmma:hessian-ibp} -- is available in the
parameters alone.

Three steps then remain, and they are what the statement still rests on. The pieces of a
family are disjoint but do not cover \(V_{\mathrm{reg}}\), so the boundary terms of the
per-piece identities do not cancel and the sum must be organised by a partition of unity
subordinate to the pieces rather than by a single cut-off. The exchange of Hessians is
stated for functions smooth on all of \(\mathbf C^k\), so \(\varphi\) and
\(v\circ\gamma_{\mathrm{man}}\) must be extended from \(\Omega\) to globally smooth
functions agreeing near the support of the test function; only values there enter either
side. Finally the boundary terms must be shown to vanish in the limit. In classical language these
three steps are the cut-off argument that justifies pairing the closed current \([V]\) with
\(\eta=d^cv\wedge S\), \(S=\sum_{i<k}(\beta+dd^cv)^i\wedge\beta^{k-1-i}\), on the
non-compact manifold \(V_{\mathrm{reg}}\); that is, Lelong's theorem (Demailly, III.2.7),
which is the Skoda--El Mir trivial extension of \([V_{\mathrm{reg}}]\) across the complete
pluripolar set \(V_{\mathrm{sing}}\) (Demailly, III.2.3). The cut-offs are \(\theta\circ
v_k\) with \(v_k\) plurisubharmonic, and only one derivative lands on them; Cauchy--Schwarz
then splits the defect into a factor that need only stay bounded --- bounded by one
integration by parts onto a \emph{fixed} test function, using
\(i\,d'd''v_k^2\ge2\,i\,d'v_k\wedge d''v_k\) --- and a factor that vanishes by dominated
convergence against the finite mass of Lemma~\ref{ana:background-mass-finite}. In
particular no smallness \emph{rate} for the background mass of a shrinking collar is needed;
finiteness is what the argument consumes.

They are isolated as Lemma~\ref{ana:stokes-certified-input}, and the exhaustion they are
asked along is the strengthened one of Lemma~\ref{ana:stokes-defect-good-exhaustion-psh}:
by Lemma~\ref{ana:stokes-defect-piece-cutoff} the per-piece cutoff estimate is then
available, which it is not along the exhaustion of
Lemma~\ref{ana:stokes-defect-good-exhaustion} alone. Nothing is lost by asking for more:
Lemma~\ref{ana:stokes-obligation-potential-free} shows that the obligation along the
strengthened exhaustion implies the obligation along the weaker one, because the collar
datum does not depend on which local potential was selected.
\end{proof}

\begin{lemma}[Restricted masses in a K\"ahler class]
\label{ana:kahler-restricted-mass-bound}
\uses{ana:restricted-volume}
\source{ct22-restricted-volumes:finite}
\lean{TranscendentalNakamaye.Pending.KahlerRestrictedMassBound.exists_acTopMass_le_of_zero_representative}
\leanok
Let \((X,\beta)\) be a compact complex manifold with a K\"ahler form and let
\(V\subset X\) be an irreducible analytic subvariety of dimension \(k\). There is a
constant \(C<\infty\) such that for every \(\beta\)-plurisubharmonic function \(u\)
with analytic singularities in the sense of Definition~\ref{as:smooth-off-pole-type}
--- remainder bounded near each point and smooth off the poles --- whose polar set
\(P\) does not contain \(V\),
\[
 \int_{V_{\mathrm{reg}}\setminus P}\bigl((\beta+dd^cu)|_{V_{\mathrm{reg}}}\bigr)^{k}\le C.
\]
\end{lemma}
\begin{proof}
\uses{ana:subvariety-stokes-mass}
\leanok
Let \(P:\mathbb R\to\mathbb R\) be smooth and convex with \(P(s)=0\) for \(s\le-1\),
\(P(s)=s\) for \(s\ge1\) and \(0\le P'\le1\). For \(j\in\mathbb N\) put
\(u_j=-j+P(u+j)\). Near a point where \(u<-j-1\), in particular near a pole
(\(u\) is upper semicontinuous), \(u_j\) is the constant \(-j\); near a point where
\(u>-j-2\), \(u\) is a smooth function \(g\) and \(u_j=-j+P(g+j)\). So \(u_j\) is
smooth, and where \(g\) is defined
\[
 \beta+dd^cu_j=(1-P')\beta+P'(\beta+dd^cg)+P''\,i\partial g\wedge\bar\partial g\ge0 ,
\]
while \(\beta+dd^cu_j=\beta\ge0\) where \(u_j\) is constant. Moreover \(u_j=u\)
near every point where \(u>-j+1\).

On \(V_{\mathrm{reg}}\setminus P\) the sets \(\{u\ge-j+2\}\) increase to
\(V_{\mathrm{reg}}\setminus P\), and on them
\((\beta+dd^cu)^k=(\beta+dd^cu_j)^k\). By monotone convergence and
Lemma~\ref{ana:subvariety-stokes-mass},
\[
 \int_{V_{\mathrm{reg}}\setminus P}(\beta+dd^cu)^k
 =\lim_{j}\int_{(V_{\mathrm{reg}}\setminus P)\cap\{u\ge-j+2\}}(\beta+dd^cu_j)^k
 \le\sup_j\int_{V_{\mathrm{reg}}}(\beta+dd^cu_j)^k=\int_{V_{\mathrm{reg}}}\beta^k<\infty .
\]
Take \(C=\int_{V_{\mathrm{reg}}}\beta^k\).
\end{proof}

\begin{lemma}[Uniform bound for restricted absolutely continuous masses]
\label{ana:rv-mass-bound}
\uses{ana:restricted-volume}
\source{ct22-restricted-volumes:finite}
\lean{TranscendentalNakamaye.Pending.RestrictedMassFiniteness.exists_acTopMass_le}
\leanok
Let \((X,\omega)\) be a compact complex manifold with a K\"ahler form, let
\(\alpha\) be a real \((1,1)\)-class, let \(V\subset X\) be an irreducible analytic
subvariety of dimension \(k\), and let \(A>0\). There is a constant \(C<\infty\)
such that for every \(0<\epsilon\le A\) and every closed real \((1,1)\)-current
\(T\in\alpha\) with analytic singularities in the sense of
Definition~\ref{as:smooth-off-pole-type}, not
containing \(V\) in its polar set, and with \(T+\epsilon\omega\ge0\),
\[
 \int_{V_{\mathrm{reg}}}\bigl((T+\epsilon\omega)|_{V_{\mathrm{reg}}}\bigr)_{\mathrm{ac}}^{k}\le C.
\]
\end{lemma}
\begin{proof}
\uses{ana:kahler-restricted-mass-bound,gnpp:kahler-domination}
\leanok
Fix a smooth representative \(\theta_0\) of \(\alpha\) and \(c>0\) with
\(c\omega-\theta_0\ge0\), and put \(\beta=(c+A+1)\omega\), a K\"ahler form. Let
\(T=\theta'+dd^cu\) be as in the statement. Then \(\theta_0=\theta'+dd^ch\) for a
smooth function \(h\), and \(u-h\) is \((\theta_0+\epsilon\omega)\)-plurisubharmonic.
Since \(\beta-\theta_0-\epsilon\omega=(c\omega-\theta_0)+(A+1-\epsilon)\omega\ge0\),
\(u-h\) is also \(\beta\)-plurisubharmonic. It has analytic singularities in the
same convention and the same polar set \(P\) as \(u\): the generators are unchanged,
the new remainder \(b-h\) is bounded near each point because \(h\) is continuous, and
it is smooth wherever \(b\) is. Off \(P\),
\[
 \beta+dd^c(u-h)=(\theta'+\epsilon\omega+dd^cu)+(\beta-\theta_0-\epsilon\omega),
\]
the sum of two semipositive \((1,1)\)-forms. Restricted to \(V_{\mathrm{reg}}\), the top
power of a semipositive form does not decrease when a semipositive form is added.
Hence
\(\int_{V_{\mathrm{reg}}}((T+\epsilon\omega)|_{V_{\mathrm{reg}}})_{\mathrm{ac}}^k\le
\int_{V_{\mathrm{reg}}\setminus P}((\beta+dd^c(u-h))|_{V_{\mathrm{reg}}})^k\), which is at most the
constant of Lemma~\ref{ana:kahler-restricted-mass-bound}.
\end{proof}

\begin{lemma}[Demailly regularization of K\"ahler currents]
\label{ana:kahler-current-regularization}
\uses{ana:nonkahler-locus,ana:ideal-sheaves,ana:demailly-algebraic-approximation}
\source{ct22-restricted-volumes:sect2}
\lean{TranscendentalNakamaye.Pending.KahlerCurrentRegularization.exists_kahlerCurrent_poleSet_subset_lelongNumber_pos,TranscendentalNakamaye.Pending.KahlerCurrentRegularizationAlgebraicApproximation.exists_kahlerCurrent_poleSet_subset_lelongNumber_pos}
\leanok
Let \(X\) be a compact complex manifold and \(T\) a K\"ahler current in a real
\((1,1)\)-class \(\alpha\). There is a K\"ahler current \(S\in\alpha\) with
analytic singularities, both with a rational coefficient and bounded remainder
along a coherent ideal and with a smooth remainder in every chart, whose polar set
is contained in \(\{x:\nu(T,x)>0\}\).
\end{lemma}
\begin{proof}
\uses{ana:demailly-algebraic-approximation,ana:analytic-type-lelong,ana:kahler-current-max}
\leanok
This is the deduction in the proof of Boucksom's Theorem~3.17 (\emph{Divisorial Zariski
decompositions}). Write \(T=\theta+dd^c\varphi\ge\omega'\) with \(\omega'\) a K\"ahler form,
and apply Theorem~\ref{ana:demailly-algebraic-approximation} to the form \(\theta-\omega'\)
and the potential \(\varphi\) of the positive current \((\theta-\omega')+dd^c\varphi\), with
\(\omega'\) as the K\"ahler form. It gives \(\varphi_k\) with analytic singularities in both
conventions of clause \emph{(i)}, decreasing to \(\varphi\), and
\(\delta_k\to0\) with \((\theta-\omega')+\delta_k\omega'+dd^c\varphi_k\ge0\). Choose \(k\) with
\(\delta_k<1\): then \(\theta+dd^c\varphi_k\ge(1-\delta_k)\omega'\), so
\(S=\theta+dd^c\varphi_k\) is a K\"ahler current in the same class \(\alpha\), with the
analytic singularities of \(\varphi_k\). At a pole \(x\) of \(\varphi_k\),
Lemma~\ref{ana:analytic-type-lelong} gives \(\nu(\varphi_k,x)>0\); and
\(\varphi_k\ge\varphi\) gives \(\nu(\varphi_k,x)\le\nu(\varphi,x)\) (the first step of the
proof of Lemma~\ref{ana:kahler-current-max}). Hence every pole of \(S\) lies in
\(\{\nu(T,\cdot)>0\}\).
\end{proof}

\begin{lemma}[Noether property of analytic sets]
\label{ana:analytic-noether}
\uses{fpgeom:complex-analytic-subvariety}
\lean{TranscendentalNakamaye.Pending.AnalyticNoetherProperty.exists_minimal_of_isComplexAnalyticSubset}
\leanok
Every nonempty family of analytic subsets of a compact complex manifold has a
member that is minimal for inclusion.
\end{lemma}
\begin{proof}
\uses{lag:germ-noetherian,lag:germ-decomposition,lag:local-statement-connected,ana:regular-component}
\leanok
It suffices to show that every descending chain \(A_0\supseteq A_1\supseteq\cdots\)
of analytic subsets of a compact complex manifold \(X\) is stationary: if the family
had no minimal member, dependent choice would produce a strictly descending chain in it.

Fix \(p\in X\). In a chart centred at \(p\) the vanishing ideals of the germs of the
\(A_i\) at \(p\) form an ascending chain in the Noetherian ring of holomorphic germs,
so they are equal for \(i\ge N\). An analytic germ is determined by its vanishing
ideal, hence for every \(i\ge N\) there is a neighbourhood \(U_i\) of \(p\) with
\(A_N\cap U_i\subseteq A_i\).

We show that the neighbourhood can be taken independent of \(i\). If \(p\notin A_N\),
take \(U=X\setminus A_N\), on which every \(A_i\), \(i\ge N\), is empty. Otherwise
take local decomposition data for \(A_N\) at \(p\) with connected regular parts: a
neighbourhood \(W\) of \(p\), pieces \(P_1,\dots,P_r\) through \(p\) with
\(A_N\cap W=\bigcup_jP_j\), open neighbourhoods \(\Delta_j\) of \(p\), and connected
sets \(G_j\) of regular points of \(A_N\), dense in \(P_j\cap\Delta_j\). Put
\(U=W\cap\bigcap_j\Delta_j\) and let \(i\ge N\). Since \(p\in P_j\cap\Delta_j\subseteq
\overline{G_j}\), the set \(G_j\) meets \(U_i\) at a regular point \(g\). Let \(C\) be
the connected component of the regular locus of \(A_N\) through \(g\); it contains the
connected set \(G_j\). The analytic set \(A_i\) contains \(C\cap U_i\), a nonempty
relatively open part of \(C\), so by the identity theorem along \(C\) it contains \(C\).
Since \(A_i\) is closed, \(\overline{G_j}\subseteq A_i\) for every \(j\), and therefore
\(A_N\cap U\subseteq A_i\). Together with \(A_i\subseteq A_N\) this gives
\(A_i\cap U=A_N\cap U\) for all \(i\ge N\).

By compactness finitely many such neighbourhoods \(U_{p_1},\dots,U_{p_s}\) cover \(X\);
with \(N=\max_kN_{p_k}\) we get \(A_i=A_N\) for all \(i\ge N\).
\end{proof}

\begin{lemma}[Maximum of two K\"ahler currents]
\label{ana:kahler-current-max}
\uses{ana:nonkahler-locus}
\lean{TranscendentalNakamaye.Pending.KahlerCurrentMax.lelongNumber_le_of_eventually_le,TranscendentalNakamaye.Pending.KahlerCurrentMax.lelongNumber_sub_le,TranscendentalNakamaye.Pending.KahlerCurrentMax.exists_kahlerCurrent_lelongNumber_le_both}
\leanok
Let \(S\) and \(T\) be K\"ahler currents in the same real \((1,1)\)-class
\(\alpha\) on a compact complex manifold. There is a K\"ahler current
\(R\in\alpha\) with \(\nu(R,x)\le\min(\nu(S,x),\nu(T,x))\) for every \(x\).
\end{lemma}
\begin{proof}
\leanok
If \(u\le v\) near \(x\), every bound \(v\le\gamma\log|z-x|+C\) is a bound for
\(u\), so \(\nu(v,x)\le\nu(u,x)\). If \(h\) is real and continuous at \(x\), then
\(h\le h(x)+1\) near \(x\), and a bound \(u-h\le\gamma\log|z-x|+C\) gives
\(u\le\gamma\log|z-x|+C+h(x)+1\); hence \(\nu(u-h,x)\le\nu(u,x)\).

Write \(S=\theta_S+dd^cu_S\ge\omega_S\) and \(T=\theta_T+dd^cu_T\ge\omega_T\).
Since \(\theta_S\) and \(\theta_T\) represent the same class,
\(\theta_S=\theta_T+dd^ch\) for a smooth function \(h\). By compactness there is
a K\"ahler form \(\omega'\) with \(\omega_S-\omega'\ge0\) and \(\omega_T-\omega'\ge0\).
Then \(u_S\) and \(u_T-h\) are both \((\theta_S-\omega')\)-plurisubharmonic, hence
so is \(w=\max(u_S,u_T-h)\), and \(w\) is not identically \(-\infty\) near any
point because \(w\ge u_S\). Thus \(R=\theta_S+dd^cw\ge\omega'\) is a K\"ahler current
in \(\alpha\). From \(w\ge u_S\) and \(w\ge u_T-h\) the first paragraph gives
\(\nu(R,x)\le\nu(S,x)\) and \(\nu(R,x)\le\nu(u_T-h,x)\le\nu(T,x)\).
\end{proof}

\begin{lemma}[Strict transforms under compositions]
\label{ana:strict-transform-composition}
\uses{ana:blowups}
\lean{TranscendentalNakamaye.Pending.StrictTransformComposition.map_nhds_eq_of_not_mem_exceptionalLocus,TranscendentalNakamaye.Pending.StrictTransformComposition.isOpen_compl_exceptionalLocus,TranscendentalNakamaye.Pending.StrictTransformComposition.strictTransform_comp}
\leanok
Let \(f:Y\to X\) and \(g:Z\to Y\) be holomorphic maps between complex manifolds of
the same dimension. Then \(\operatorname{Exc}(g)\) is closed, \(g\) maps every
neighbourhood of a point \(z\notin\operatorname{Exc}(g)\) onto a neighbourhood of
\(g(z)\), and for every \(A\subset X\) the strict transforms satisfy
\(\widetilde A^{f\circ g}=\widetilde{(\widetilde A^{f})}^{g}\), where
\(\widetilde A^{f}=\overline{f^{-1}(A)\setminus\operatorname{Exc}(f)}\).
\end{lemma}
\begin{proof}
\leanok
Near \(z\), in charts \(e\) at \(z\) and \(e'\) at \(g(z)\), write
\(g=e'^{-1}\circ F\circ e\) with \(F\) smooth. For \(y\) near \(z\) the
differential of \(g\) at \(y\) is \(D(e'^{-1})\circ DF(e(y))\circ De(y)\) with
invertible outer factors, so it is bijective exactly when \(DF(e(y))\) is.
Invertible operators form an open set and \(DF\) is continuous, so the set of
points with bijective differential is open. At such a point \(DF(e(z))\) is
invertible, and the inverse function theorem gives
\(F(\text{nbhd of }e(z))=\text{nbhd of }F(e(z))\); composing with the charts,
\(g\) maps neighbourhoods of \(z\) onto neighbourhoods of \(g(z)\).

Since \(\operatorname{Exc}(f\circ g)=\operatorname{Exc}(g)\cup g^{-1}\operatorname{Exc}(f)\),
with \(S=f^{-1}(A)\setminus\operatorname{Exc}(f)\) we have
\((f\circ g)^{-1}(A)\setminus\operatorname{Exc}(f\circ g)=g^{-1}(S)\setminus\operatorname{Exc}(g)\).
Its closure is contained in \(\overline{g^{-1}(\overline S)\setminus\operatorname{Exc}(g)}\).
Conversely let \(z\notin\operatorname{Exc}(g)\) with \(g(z)\in\overline S\), and let
\(t\) be a neighbourhood of \(z\). Then \(t\setminus\operatorname{Exc}(g)\) is a
neighbourhood of \(z\), its image is a neighbourhood of \(g(z)\), hence meets \(S\):
some \(w\in t\setminus\operatorname{Exc}(g)\) has \(g(w)\in S\). So
\(z\in\overline{g^{-1}(S)\setminus\operatorname{Exc}(g)}\), and taking closures gives
the reverse inclusion.
\end{proof}

\begin{lemma}[Blowup sequences are surjective]
\label{ana:blowup-surjective}
\uses{ana:blowups,ana:strict-transform-composition}
\lean{TranscendentalNakamaye.Pending.StrictTransformComposition.surjective_of_isBlowupAlong,TranscendentalNakamaye.Pending.StrictTransformComposition.IsBlowupSequence.compactSpace_surjective,TranscendentalNakamaye.Pending.StrictTransformComposition.image_exceptionalLocus_subset_comp,TranscendentalNakamaye.Pending.StrictTransformComposition.not_strictTransform_subset_image_exceptionalLocus}
\leanok
Let \(X\) be a compact complex manifold and \(\mu:\widetilde X\to X\) a finite
composition of blowups with smooth centres. Then \(\widetilde X\) is compact and
\(\mu\) is surjective. If \(f:Y\to X\) and \(g:Z\to Y\) are surjective
holomorphic maps between manifolds of the same dimension and
\(V\not\subset(f\circ g)(\operatorname{Exc}(f\circ g))\), then
\(V\not\subset f(\operatorname{Exc}(f))\) and
\(\widetilde V^{f}\not\subset g(\operatorname{Exc}(g))\).
\end{lemma}
\begin{proof}
\leanok
A blowup \(g:Z\to Y\) along a nowhere dense centre \(C\) of a compact manifold is
an isomorphism over \(Y\setminus C\), and \(Z\) is compact. Its image is compact,
hence closed, and contains the dense set \(Y\setminus C\), so \(g\) is onto. By
induction along the sequence every intermediate manifold is compact and the
composition is onto.

If \(y=f(x)\) with \(x\in\operatorname{Exc}(f)\), write \(x=g(w)\); then
\(w\in g^{-1}\operatorname{Exc}(f)\subset\operatorname{Exc}(f\circ g)\), so
\(f(\operatorname{Exc} f)\subset(f\circ g)(\operatorname{Exc}(f\circ g))\), which
gives the first claim. For the second, suppose
\(\widetilde V^{f}\subset g(\operatorname{Exc} g)\) and let \(v\in V\), \(v=f(y)\).
If \(y\in\operatorname{Exc}(f)\), then \(v\in(f\circ g)(\operatorname{Exc}(f\circ g))\)
as above. Otherwise \(y\in f^{-1}(V)\setminus\operatorname{Exc}(f)\subset\widetilde V^{f}\),
so \(y=g(w)\) with \(w\in\operatorname{Exc}(g)\subset\operatorname{Exc}(f\circ g)\)
and \(v=(f\circ g)(w)\). Thus \(V\subset(f\circ g)(\operatorname{Exc}(f\circ g))\), a
contradiction.
\end{proof}

\begin{lemma}[Restricted volumes do not depend on the K\"ahler form]
\label{ana:rv-kahler-independence}
\uses{ana:restricted-volume,ana:restricted-volume-monotone}
\source{ct22-restricted-volumes:sect2}
\lean{TranscendentalNakamaye.Pending.RestrictedVolumeKahlerIndependence.changeKahler,TranscendentalNakamaye.Pending.RestrictedVolumeKahlerIndependence.acTopMass_le_changeKahler,TranscendentalNakamaye.Pending.RestrictedVolumeKahlerIndependence.numericalRestrictedVolume_kahler_independent}
\leanok
On a compact complex manifold, for K\"ahler forms \(\omega,\omega'\), every real
\((1,1)\)-class \(\alpha\) and every irreducible subvariety \(V\),
\(\operatorname{RV}^{\omega}_{X\mid V}(\alpha)=\operatorname{RV}^{\omega'}_{X\mid V}(\alpha)\).
\end{lemma}
\begin{proof}
\leanok
By compactness there is \(c>0\) with \(c\omega'-\omega\ge0\). If
\(T=\theta+dd^cu\) is admissible at error \(\epsilon\) for \(\omega\), then
\(\theta+c\epsilon\omega'+dd^cu=(\theta+\epsilon\omega+dd^cu)+\epsilon(c\omega'-\omega)\ge0\),
so \(T\) is admissible at error \(c\epsilon\) for \(\omega'\). Off the polar set, on
each regular patch, \((T+c\epsilon\omega')|_V=(T+\epsilon\omega)|_V+\epsilon(c\omega'-\omega)|_V\)
is a sum of semipositive \((1,1)\)-forms, and the top power density is monotone on
such sums, so the absolutely continuous restricted top mass does not decrease.
Hence the envelope for \(\omega\) at \(\epsilon\) is at most the envelope for
\(\omega'\) at \(c\epsilon\). The non-nef locus does not depend on the K\"ahler form,
so either both volumes vanish or both are infima of the envelopes
(Lemma~\ref{ana:restricted-volume-monotone}); reindexing \(\epsilon\mapsto c\epsilon\)
gives \(\operatorname{RV}^{\omega}\le\operatorname{RV}^{\omega'}\), and symmetry gives
equality.
\end{proof}

\begin{lemma}[Preimages of analytic sets]
\label{ana:analytic-preimage}
\uses{fpgeom:complex-analytic-subvariety,ana:blowups}
\lean{TranscendentalNakamaye.Pending.AnalyticSubsetPreimage.isComplexAnalyticSubset_preimage,TranscendentalNakamaye.Pending.AnalyticSubsetPreimage.isComplexAnalyticSubset_of_isClosedSubmanifold,TranscendentalNakamaye.Pending.AnalyticSubsetPreimage.isComplexAnalyticSubset_of_locally_eq}
\leanok
Let \(g:Z\to M\) be a holomorphic map of complex manifolds. If \(A\subset M\) is
analytic, so is \(g^{-1}(A)\). A closed complex submanifold is analytic. If \(A\) is
analytic, \(B\) is closed, and every point of \(B\) has a neighbourhood \(N\) with
\(A\cap N=B\cap N\), then \(B\) is analytic.
\end{lemma}
\begin{proof}
\leanok
Near a point outside a closed set, the constant equation \(1=0\) presents the set.
If \(A\) is the common zero set of \(f_1,\dots,f_r\) near \(g(z)\), then near \(z\)
the set \(g^{-1}(A)\) is the common zero set of the holomorphic functions \(f_i\circ g\).
Near a point of a submanifold flattened to \(\mathbf C^k\times0\) by \(\Phi\), the
coordinates \(\Phi_i\), \(i>k\), are equations. In the last statement, near a point of
\(B\) the equations of \(A\) restricted to \(N\) present \(B\).
\end{proof}

\begin{lemma}[Regular points across local biholomorphisms]
\label{ana:regular-transport}
\uses{fpgeom:complex-analytic-subvariety,ana:blowups}
\lean{TranscendentalNakamaye.Pending.RegularAtDimTransport.regularAtDim_of_preimage}
\leanok
Let \(g:Z\to Y\) be holomorphic between manifolds of the same dimension,
\(z\notin\operatorname{Exc}(g)\), and \(A'\subset Z\), \(A\subset Y\) with
\(A'\cap O=g^{-1}(A)\cap O\) for an open \(O\ni z\). If \(A\) is regular of dimension
\(d\) at \(g(z)\), then \(A'\) is regular of dimension \(d\) at \(z\).
\end{lemma}
\begin{proof}
\leanok
In charts at \(z\) and \(g(z)\), \(g\) is a holomorphic map \(F\) with invertible
differential at the centre, so by the holomorphic inverse function theorem it is a
biholomorphism \(\Psi\) near the centre. If \(\Phi\) flattens \(A\) to
\(\mathbf C^d\times0\) near \(g(z)\), then \(\Phi\circ\Psi\), restricted so that its
source lies in the chart image of \(O\), flattens \(A'\) near \(z\).
\end{proof}

\begin{lemma}[A regular component minus an analytic set]
\label{ana:component-complement-connected}
\uses{ana:regular-component,lag:hypersurface-complement-connected}
\lean{TranscendentalNakamaye.Pending.ComponentComplementConnected.isPreconnected_component_diff_and_subset_closure}
\leanok
Let \(A\) and \(T\) be analytic subsets of a complex manifold and \(K\) the connected
component through a regular point of the regular locus of \(A\). If \(K\not\subset T\),
then \(K\setminus T\) is connected and dense in \(K\).
\end{lemma}
\begin{proof}
\leanok
Let \(x\in K\) and take a flattening ball \(B\) of \(A\) at \(x\), on which \(A\) is a
ball in \(\mathbf C^j\) contained in \(K\), small enough that the local equations of
\(T\) at \(x\) are defined on it. If all of them vanished on \(A\cap B\), then \(T\)
would contain a nonempty relatively open part of \(K\), hence \(K\) by the identity
theorem along \(K\) (Lemma~\ref{ana:regular-component}). So one equation, read on the
ball of \(\mathbf C^j\), is not identically zero, hence not identically zero near the
centre, and by Lemma~\ref{lag:hypersurface-complement-connected} the ball minus its zero
set is connected and dense. The set \((A\cap B)\setminus T\) lies between this set and
its closure, so it is connected and dense in \(A\cap B\).

Density of \(K\setminus T\) in \(K\) follows. Suppose \(K\setminus T=U_1\sqcup U_2\) with
\(U_1,U_2\) relatively open. For each \(x\in K\) the connected set
\((A\cap B_x)\setminus T\) lies in \(U_1\) or in \(U_2\); let \(K_i\) be the set of
\(x\) for which it lies in \(U_i\). Two such neighbourhoods of the same point meet in a
nonempty part of \(K\setminus T\) by density, so \(K_1\cap K_2=\emptyset\); both are
relatively open and they cover \(K\). As \(K\) is connected, one of them, say \(K_2\),
is empty, and then \(U_2=\emptyset\).
\end{proof}

\begin{lemma}[Fibres of a blowup and the image of its exceptional locus]
\label{ana:blowup-exceptional-image}
\uses{ana:blowups,ana:strict-transform-composition,ana:analytic-preimage}
\lean{TranscendentalNakamaye.Pending.ExceptionalFibres.interior_preimage_eq_empty,TranscendentalNakamaye.Pending.ExceptionalFibres.eq_of_apply_eq_of_not_mem_exceptionalLocus,TranscendentalNakamaye.Pending.ExceptionalFibres.preimage_image_exceptionalLocus,TranscendentalNakamaye.Pending.ExceptionalFibres.isClosed_image_exceptionalLocus,TranscendentalNakamaye.Pending.ExceptionalImageCodimension.mem_exceptionalLocus_of_mem,TranscendentalNakamaye.Pending.ExceptionalImageCodimension.not_mem_exceptionalLocus,TranscendentalNakamaye.Pending.ExceptionalImageAnalytic.isComplexAnalyticSubset_image_exceptionalLocus}
\leanok
Let \(\pi:Z\to M\) be the blowup of a compact complex manifold along a nowhere dense
closed submanifold \(C\), and \(T=\pi(\operatorname{Exc}\pi)\). Then
\(\pi^{-1}(C)\) has empty interior; if \(z\notin\operatorname{Exc}\pi\), then
\(\pi^{-1}(\pi(z))=\{z\}\); \(\pi^{-1}(T)=\operatorname{Exc}\pi\); and \(T\) is the
union of the components of \(C\) of codimension at least two, a closed analytic subset of
\(M\) contained in \(C\).
\end{lemma}
\begin{proof}
\leanok
Near any point, \(\mathcal I_C\mathcal O_Z\) is generated by one function \(h\) vanishing
on no nonempty open set, and \(h\) is a combination of pullbacks of functions vanishing on
\(C\), so \(\pi^{-1}(C)\subset\{h=0\}\) locally; hence \(\pi^{-1}(C)\) has empty interior.

Let \(z\notin\operatorname{Exc}\pi\) and \(\pi(z')=\pi(z)\) with \(z'\ne z\). Choose
disjoint open \(U\ni z\), \(U'\ni z'\). By Lemma~\ref{ana:strict-transform-composition}
\(\pi(U)\) is a neighbourhood of \(\pi(z)\), so \(U'\cap\pi^{-1}(\operatorname{int}\pi(U))\)
is a nonempty open set, not contained in \(\pi^{-1}(C)\). A point \(u\) of it over
\(M\setminus C\) equals \(\pi(u_1)\) with \(u_1\in U\), and \(\pi\) is injective over
\(M\setminus C\), so \(u=u_1\in U\cap U'\), a contradiction. Consequently
\(\pi^{-1}(T)=\operatorname{Exc}\pi\). Since \(\operatorname{Exc}\pi\subset\pi^{-1}(C)\) is
closed and \(Z\) is compact, \(T\) is a closed subset of \(C\).

Let \(y\in C\), flatten \(C\) near \(y\) to \(\mathbf C^k\times0\) by \(\Phi\), and let
\(\sigma_j(w)_i=w_iw_j\) for \(i\ge k\), \(i\ne j\), and \(\sigma_j(w)_i=w_i\) otherwise
(\(j\ge k\)). Near every \(z\) over the flattening domain, \(\Phi\circ\pi=\sigma_j\circ W\)
for some \(j\) and a holomorphic \(W\), and \(W\) inverts the holomorphic lift \(\ell_j\) of
\(\Phi^{-1}\circ\sigma_j\) through \(\pi\). In particular \(k<n\). If \(k\le n-2\) and
\(\pi(z)\in C\), then \(W(z)_j=0\), so \(\sigma_j\) is constant along
\(W(z)+t e_l\) for some \(l\ge k\), \(l\ne j\); \(d\sigma_j(W(z))\) is not injective, and
since \(d\Phi\circ d\pi=d\sigma_j\circ dW\) the differential \(d\pi_z\) is not bijective:
\(z\in\operatorname{Exc}\pi\). If \(k=n-1\), then \(\sigma_j=\mathrm{id}\) and
\(\ell_j\circ\Phi\) is a holomorphic section of \(\pi\) through \(z\), so \(d\pi_z\) is
surjective and \(z\notin\operatorname{Exc}\pi\). As \(\pi\) is surjective, near a point of
\(T\) the centre has codimension at least two and \(T\) agrees with \(C\); near a point
outside \(T\) the closed set \(T\) is absent. By Lemma~\ref{ana:analytic-preimage}
\(T\) is analytic.
\end{proof}

\begin{lemma}[Strict transform under one blowup]
\label{ana:strict-transform-blowup}
\uses{ana:blowups,fpgeom:complex-analytic-subvariety,ana:blowup-exceptional-image,ana:analytic-preimage,ana:regular-transport,ana:component-complement-connected,ana:closure-difference-analytic,ana:regular-component,ana:pure-dimension,lag:regular-locus-dense}
\source{ct22-restricted-volumes:invar}
\lean{TranscendentalNakamaye.Pending.RestrictedVolumeOneBlowup.exists_strictTransform_subvariety_of_isBlowupAlong,TranscendentalNakamaye.Pending.StrictTransformSubvariety.eq_closure_component,TranscendentalNakamaye.Pending.StrictTransformSubvariety.exists_strictTransform_subvariety_of_isBlowupAlong}
\leanok
Let \(g:Z\to Y\) be the blowup of a compact complex manifold along a nowhere dense closed
submanifold and \(W\subset Y\) a closed irreducible analytic subset with
\(W\not\subset g(\operatorname{Exc}(g))\). Then the strict transform
\(\widetilde W=\overline{g^{-1}(W)\setminus\operatorname{Exc}(g)}\) is a closed
irreducible analytic subset of \(Z\) of dimension \(\dim W\).
\end{lemma}
\begin{proof}
\leanok
Let \(T=g(\operatorname{Exc}g)\), a closed analytic set with
\(g^{-1}(T)=\operatorname{Exc}g\) (Lemma~\ref{ana:blowup-exceptional-image}), and
\(\Omega=Z\setminus\operatorname{Exc}g\). Then
\(\widetilde W=\overline{g^{-1}(W)\setminus g^{-1}(T)}\) is analytic
(Lemmas~\ref{ana:analytic-preimage} and~\ref{ana:closure-difference-analytic}), and
\(\widetilde W\cap\Omega=g^{-1}(W)\cap\Omega\).

First, an irreducible analytic set \(A\) with a regular point \(c\) equals
\(\overline K\), \(K\) the component of its regular locus through \(c\): indeed
\(A=\overline K\cup\overline{A\setminus\overline K}\) is a union of two analytic sets
(Lemmas~\ref{ana:regular-component} and~\ref{ana:closure-difference-analytic}), and near
\(c\) the set \(A\) lies in \(K\), so \(c\notin\overline{A\setminus\overline K}\).

Since \(W\setminus T\) is a nonempty relatively open subset of \(W\), it contains a
regular point \(c\), of dimension \(\dim W\) by pure dimensionality. Let \(\tilde c\) be a
preimage (\(g\) is surjective); \(\tilde c\in\Omega\), and by
Lemma~\ref{ana:regular-transport} \(\widetilde W\) is regular of dimension \(\dim W\) at
every point \(u\in\Omega\) with \(g(u)\) a regular point of \(W\) of that dimension; in
particular at \(\tilde c\). Let \(K\) be the component through \(c\), so
\(W=\overline K\) and \(K\not\subset T\); by
Lemma~\ref{ana:component-complement-connected} \(K\setminus T\) is connected and dense
in \(K\). The restriction of \(g\) to \(\Omega\) is injective
(Lemma~\ref{ana:blowup-exceptional-image}) and open, and its image contains
\(K\setminus T\); hence \(S=\Omega\cap g^{-1}(K\setminus T)\) is connected. It consists
of regular points of \(\widetilde W\) and contains \(\tilde c\), so it lies in the
component \(\widetilde K\) of the regular locus of \(\widetilde W\) through \(\tilde c\).
Every \(u\in g^{-1}(W)\cap\Omega\) lies in \(\overline S\): each neighbourhood of \(u\)
in \(\Omega\) maps onto a neighbourhood of \(g(u)\in\overline{K\setminus T}\). Hence
\(\widetilde W=\overline{\widetilde K}\), which is irreducible
(Lemma~\ref{ana:regular-component}), and \(\widetilde W\) has dimension \(\dim W\).
\end{proof}

\begin{theorem}[K\"ahler currents with polar set the non-K\"ahler locus]
\label{ana:nonkahler-locus-analytic}
\uses{ana:nonkahler-locus,ana:ideal-sheaves}
\source{ct22-restricted-volumes:sect2}
\lean{TranscendentalNakamaye.Pending.NonKahlerLocusPoleSet.exists_kahlerCurrent_poleSet_eq_nonKahlerLocus}
\leanok
Let \(X\) be a compact complex manifold of dimension \(n>0\) and \(\alpha\) a big
real \((1,1)\)-class. There is a K\"ahler current in \(\alpha\) with analytic
singularities (with a rational coefficient and bounded remainder, and with a smooth
remainder in every chart) whose polar set is exactly \(E_{\mathrm{nK}}(\alpha)\). In
particular \(E_{\mathrm{nK}}(\alpha)\) is an analytic subset.
\end{theorem}
\begin{proof}
\uses{ana:kahler-current-regularization,ana:analytic-noether,ana:kahler-current-max}
\leanok
The polar sets of K\"ahler currents in \(\alpha\) with analytic singularities form a
nonempty family (Lemma~\ref{ana:kahler-current-regularization}) of analytic sets, so
it has a minimal member \(P(S)\) (Lemma~\ref{ana:analytic-noether}). A positive
Lelong number forces a pole, so \(E_{\mathrm{nK}}(\alpha)\subset P(S)\). Conversely,
for any K\"ahler current \(T\in\alpha\), let \(R\) be as in
Lemma~\ref{ana:kahler-current-max} and \(S'\) a regularization of \(R\); then
\(P(S')\subset\{\nu(R)>0\}\subset\{\nu(S)>0\}\subset P(S)\), so \(P(S')=P(S)\) by
minimality and \(\nu(T,x)\ge\nu(R,x)>0\) on \(P(S)\). Hence
\(P(S)=E_{\mathrm{nK}}(\alpha)\). The polar set of a potential with analytic
singularities is locally the common zero set of finitely many holomorphic functions,
hence analytic.
\end{proof}

\begin{lemma}[Poles of an analytic singularity type have positive Lelong number]
\label{ana:analytic-type-lelong}
\uses{ana:nonkahler-locus,as:rational-type}
\lean{TranscendentalNakamaye.Pending.QuasiEquisingularKahlerCurrent.exists_generatorLog_le,TranscendentalNakamaye.Pending.QuasiEquisingularKahlerCurrent.lelongNumber_pos_of_hasAnalyticSingularityTypeAt,TranscendentalNakamaye.Pending.AnalyticPoleLelongLowerBound.ofReal_two_mul_le_lelongNumber_of_logarithmicTypeAt}
\leanok
Let \(u\) be an extended-real function on a complex manifold and let \(x\) be a point
at which the chart germ of \(u\) has a logarithmic presentation
\(u=c\log\sum_{j}|f_j|^2+b\) with \(c>0\) rational, the \(f_j\) holomorphic near \(x\)
and \(b\) bounded near \(x\). If \(u(x)=-\infty\) then \(\nu(u,x)\ge 2c>0\).
\end{lemma}
\begin{proof}
\leanok
Since \(u(x)=-\infty\) and \(b\) is finite, \(f_j(x)=0\) for every \(j\). Each \(f_j\) is
complex differentiable at \(x\), so \(|f_j(w)|\le c_j\|w-x\|\) near \(x\); with
\(A=1+\sum_j c_j^2\) this gives \(\sum_j|f_j(w)|^2\le A\|w-x\|^2\) near \(x\), hence
\[
 \log\sum_j|f_j(w)|^2\le 2\log\|w-x\|+\log A
\]
(the inequality being trivial where the left side is \(-\infty\)). Multiplying by \(c\)
and adding the bound \(|b|\le C_b\) yields
\(u(w)\le 2c\log\|w-x\|+(c\log A+C_b)\) near \(x\), which is exactly the defining
condition for \(\nu(u,x)\ge2c\).
\end{proof}

\begin{lemma}[Polar sets of analytic singularity types are analytic]
\label{ana:analytic-type-pole-set}
\uses{as:rational-type,fpgeom:complex-analytic-subvariety}
\lean{TranscendentalNakamaye.Pending.NonKahlerLocusPoleSetBounded.isComplexAnalyticSubset_poleSet_of_analyticSingularityType}
\leanok
If a function \(u\) on a complex manifold has analytic singularity type in every chart,
that is, near every point \(u=c\log\sum_j|f_j|^2+b\) with \(c>0\) rational, the \(f_j\)
holomorphic and \(b\) bounded, then \(\{u=-\infty\}\) is an analytic subset.
\end{lemma}
\begin{proof}
\leanok
Near a given point the presentation holds on one open set, and there
\(c\log\sum_j|f_j|^2+b=-\infty\) exactly when every \(f_j\) vanishes, because \(c>0\)
and \(b\) is real-valued. So \(\{u=-\infty\}\) coincides on that set with the common
zero locus of the finitely many holomorphic functions \(f_j\).
\end{proof}

\begin{theorem}[Regularization of a K\"ahler current, bounded convention]
\label{ana:regularization-kahler-current-bounded}
\uses{ana:nonkahler-locus,trace:qe-definition}
\source{transcendental-nakamaye-copy4-20260911:sec:prelim,xia-sgpt:page-0049}
\lean{TranscendentalNakamaye.Pending.QuasiEquisingularKahlerCurrent.exists_kahlerCurrent_analyticSingularityType_poleSet_subset}
\leanok
Let \(X\) be a compact complex manifold and \(T=\theta+dd^cu\ge\omega_0\) a K\"ahler
current in a real \((1,1)\)-class \(\alpha\). There is a K\"ahler current \(S\in\alpha\)
whose potential has analytic singularity type in every chart and whose polar set is
contained in \(\{x:\nu(u,x)>0\}\).
\end{theorem}
\begin{proof}
\uses{trace:qe-existence,ana:analytic-type-lelong}
\leanok
Apply Theorem~\ref{trace:qe-existence} to \(u\in\operatorname{PSH}(X,\theta-\omega_0)\)
with background \(\omega_0\): there are \(\epsilon_k\downarrow0\) and \(u_k\) decreasing
pointwise to \(u\), each with analytic singularity type in every chart and with
\(u_k\in\operatorname{PSH}(X,\theta-\omega_0+\epsilon_k\omega_0)\). Fix \(k\) with
\(\epsilon_k<1\). Since
\(\theta-\omega_0+\epsilon_k\omega_0=\theta-(1-\epsilon_k)\omega_0\) and
\((1-\epsilon_k)\omega_0\) is K\"ahler, \(S:=\theta+dd^cu_k\) is a K\"ahler current in
\(\alpha\). If \(u_k(x)=-\infty\) then \(\nu(u_k,x)>0\) by
Lemma~\ref{ana:analytic-type-lelong}, and \(u\le u_k\) gives
\(\nu(u,x)\ge\nu(u_k,x)>0\).
\end{proof}

\begin{theorem}[K\"ahler currents with polar set the non-K\"ahler locus, bounded convention]
\label{ana:nonkahler-locus-analytic-bounded}
\uses{ana:nonkahler-locus,as:rational-type}
\source{transcendental-nakamaye-copy4-20260911:sec:nK}
\lean{TranscendentalNakamaye.Pending.NonKahlerLocusPoleSetBounded.exists_kahlerCurrent_analyticSingularityType_poleSet_eq_nonKahlerLocus,TranscendentalNakamaye.Pending.NonKahlerLocusPoleSetBounded.isComplexAnalyticSubset_nonKahlerLocus,TranscendentalNakamaye.Pending.NonKahlerLocusPoleSet.isComplexAnalyticSubset_nonKahlerLocus}
\leanok
Let \(X\) be a compact complex manifold of dimension \(n>0\) and \(\alpha\) a big real
\((1,1)\)-class. There is a K\"ahler current in \(\alpha\) whose potential has analytic
singularity type in every chart and whose polar set is exactly
\(E_{\mathrm{nK}}(\alpha)\). In particular \(E_{\mathrm{nK}}(\alpha)\) is an analytic
subset.
\end{theorem}
\begin{proof}
\uses{ana:regularization-kahler-current-bounded,ana:analytic-noether,ana:kahler-current-max,ana:analytic-type-pole-set}
\leanok
The polar sets of K\"ahler currents in \(\alpha\) whose potentials have analytic
singularity type form a family of analytic sets
(Lemma~\ref{ana:analytic-type-pole-set}), nonempty by
Theorem~\ref{ana:regularization-kahler-current-bounded} applied to any K\"ahler current
in \(\alpha\); let \(P(S)\) be a minimal member (Lemma~\ref{ana:analytic-noether}). A
positive Lelong number forces a pole, so \(E_{\mathrm{nK}}(\alpha)\subset P(S)\).
Conversely, let \(T\in\alpha\) be any K\"ahler current, let \(R\) be as in
Lemma~\ref{ana:kahler-current-max} and let \(S'\) be the current supplied by
Theorem~\ref{ana:regularization-kahler-current-bounded} applied to \(R\). Then
\(P(S')\subset\{\nu(R)>0\}\subset\{\nu(S)>0\}\subset P(S)\), so \(P(S')=P(S)\) by
minimality, and \(\nu(T,x)\ge\nu(R,x)>0\) for \(x\in P(S)\). Hence
\(P(S)=E_{\mathrm{nK}}(\alpha)\), and this set is analytic by
Lemma~\ref{ana:analytic-type-pole-set}.
\end{proof}

\begin{lemma}[Blowups along a centre of codimension one]
\label{ana:blowup-codim-one}
\uses{ana:blowups}
\lean{TranscendentalNakamaye.Pending.BlowupExistenceCodimOne.isBlowupAlong_id_of_codimOne,TranscendentalNakamaye.Pending.BlowupExistenceCodimOne.exists_isBlowupAlong_of_codimOne}
\leanok
Let \(C\) be a nowhere dense closed submanifold of a complex manifold \(X\) of dimension \(n\),
of dimension \(n-1\) at each of its points. Then the identity of \(X\) is a blowup of \(X\) along
\(C\).
\end{lemma}
\begin{proof}
\leanok
The ideal of a smooth hypersurface is already invertible, and a holomorphic map pulling the
ideal of \(C\) back to an invertible ideal factors uniquely through the identity.
\end{proof}

\begin{lemma}[Existence of the blowup along a smooth centre]
\label{ana:blowup-existence}
\uses{ana:blowups,ana:blowup-codim-one}
\source{transcendental-nakamaye-copy4-20260911:sec:prelim}
\lean{TranscendentalNakamaye.Pending.BlowupExistence.exists_isBlowupAlong}
\leanok
Let \(C\) be a nowhere dense closed submanifold of a complex manifold \(X\). There is a blowup
\(\pi:Z\to X\) of \(X\) along \(C\).
\end{lemma}
\begin{proof}
\uses{bex:blowup-exists}
\leanok
Take \(Z=\tilde X\) of Section~\ref{sec:blowup-existence} (Demailly's construction,
\((X\setminus C)\sqcup P(NC)\) with the monomial atlas); it is a blowup of \(X\) along \(C\) by
Theorem~\ref{bex:blowup-exists}.
\end{proof}
Every statement about blowups in this development takes a blowup as given; this lemma says that
one always exists, and nothing uses it. In codimension one the blowup is the identity
(Lemma~\ref{ana:blowup-codim-one}).

\begin{lemma}[Blowups of compact manifolds are compact]
\label{ana:blowup-compact}
\uses{ana:blowups}
\lean{TranscendentalNakamaye.Pending.BlowupCompact.compactSpace_of_isBlowupAlong}
\leanok
Let \(g:Z\to Y\) be a blowup of a compact complex manifold along a nowhere dense closed
submanifold. Then \(Z\) is compact.
\end{lemma}
\begin{proof}
\leanok
Every point of \(Y\) has a neighbourhood \(K\) with \(g^{-1}(K)\) compact: off the centre
\(g\) is a biholomorphism near the point, and near a point of the centre the universal
property identifies \(g^{-1}(K)\) with a closed subset of \(K\times\mathbf P^{r}\) for a
coordinate polydisc \(K\). Finitely many such \(K\) cover \(Y\).
\end{proof}

\begin{lemma}[Subadditivity and finiteness of Lelong numbers]
\label{ana:lelong-additivity}
\uses{ana:nonkahler-locus}
\lean{TranscendentalNakamaye.Pending.LelongNumberAdditivity.lelongNumber_add_le,TranscendentalNakamaye.Pending.LelongNumberAdditivity.lelongNumber_ne_top}
\leanok
For quasi-plurisubharmonic \(u,v\) on a complex manifold and any point \(x\),
\(\nu(u+v,x)\le\nu(u,x)+\nu(v,x)\).  If moreover \(n>0\) and \(u\) is not identically
\(-\infty\) near any point, then \(\nu(u,x)<\infty\).  Neither extra hypothesis can be
dropped: the identically \(-\infty\) germ is quasi-plurisubharmonic with infinite Lelong
number, and in dimension \(0\) every punctured neighbourhood filter is trivial.
\end{lemma}
\begin{proof}
\uses{ana:nonkahler-locus}
Both statements are local, and in a coordinate chart centred at \(x\) the Lelong number
is by definition
\[
 \nu(u,x)=\sup\Gamma(u),\qquad
 \Gamma(u)=\{\gamma\ge0:\ u(z)\le\gamma\log|z|+C\ \text{near }0\ \text{for some }C\}.
\]

\emph{Finiteness.} Let \(u\) be quasi-plurisubharmonic near \(0\) and not identically
\(-\infty\) there, and let \(n>0\). Write \(M_u(r)=\sup_{|z|=r}u\) for \(0<r<r_0\). For
a plurisubharmonic function \(M_u(e^{t})\) is a nondecreasing convex function of \(t\)
(Hadamard's three-circle inequality), and it is finite for every \(r<r_0\) because \(u\)
is not identically \(-\infty\). A finite convex nondecreasing function of
\(t=\log r\) on \((-\infty,\log r_0)\) has a slope at \(-\infty\), namely
\(L=\lim_{t\to-\infty}M_u(e^t)/t\in[0,\infty)\), and lies above every chord, so
\(M_u(e^t)\ge L\,t+C_0\) for \(t\) below any fixed \(t_1<\log r_0\) with
\(C_0=M_u(e^{t_1})-L\,t_1\). If \(\gamma\in\Gamma(u)\), then \(M_u(r)\le\gamma\log r+C\)
for small \(r\); dividing by \(\log r<0\) and letting \(r\downarrow0\) gives
\(L\ge\gamma\). Hence \(\Gamma(u)\subset[0,L]\) and \(\nu(u,x)\le L<\infty\). Both
hypotheses are used: the identically \(-\infty\) germ has \(M_u\equiv-\infty\) and every
\(\gamma\) admissible, and for \(n=0\) the punctured neighbourhood filter is trivial, so
again every \(\gamma\) is admissible.

\emph{Subadditivity.} The elementary implication goes the wrong way: if
\(\gamma\in\Gamma(u)\) and \(\gamma'\in\Gamma(v)\) then
\(u+v\le(\gamma+\gamma')\log|z|+(C+C')\), so \(\gamma+\gamma'\in\Gamma(u+v)\) and
\(\nu(u+v,x)\ge\nu(u,x)+\nu(v,x)\). The stated inequality is the other one, and it is not
available from the supremum description alone: it comes from the classical identification
of \(\sup\Gamma(u)\) with the mean-value Lelong number
\[
 \nu(u,x)=\lim_{r\downarrow0}\frac{\lambda(u,r)}{\log r},\qquad
 \lambda(u,r)=\fint_{|z|=r}u,
\]
which holds for every plurisubharmonic germ that is not identically \(-\infty\). Granting
it, \(\lambda(\cdot,r)\) is exactly additive, \(\lambda(u+v,r)=\lambda(u,r)+\lambda(v,r)\),
so the two Lelong numbers add and in particular \(\nu(u+v,x)\le\nu(u,x)+\nu(v,x)\). If
either \(u\) or \(v\) is identically \(-\infty\) near \(x\) the right-hand side is
\(+\infty\) and there is nothing to prove.
\end{proof}

\begin{lemma}[Minimal multiplicity on the non-nef locus]
\label{ana:boucksom-minimal-multiplicity}
\uses{ana:nonkahler-locus}
\lean{TranscendentalNakamaye.Pending.BoucksomMinimalMultiplicity.exists_uniform_lelongNumber_lowerBound_of_mem_nonNefLocus}
\leanok
Let \((Y,\omega)\) be compact K\"ahler, \(\alpha\) a real \((1,1)\)-class and
\(y\in E_{\mathrm{nn}}(\alpha)\). There are \(c>0\) and \(t_0>0\) such that every K\"ahler
current in \(\alpha+t\{\omega\}\), for \(0<t\le t_0\), has Lelong number at least \(c\) at
\(y\).
\end{lemma}
\begin{proof}
\uses{ana:nonkahler-locus,fpenv:v-psh,fpenv:v-minimal}
\leanok
By definition \(E_{\mathrm{nn}}(\alpha)=\bigcup_{\varepsilon>0}E_{\mathrm
{nK}}(\alpha+\varepsilon\{\omega\})\), so fix \(\varepsilon>0\) with
\(y\in E_{\mathrm{nK}}(\alpha+\varepsilon\{\omega\})\) and put
\(s=\tfrac23\varepsilon\), \(t_0=\tfrac13\varepsilon\).

If \(\alpha+s\{\omega\}\) carries no K\"ahler current, then neither does
\(\alpha+t\{\omega\}\) for \(t\le t_0<s\): adding \((s-t)\omega\) to a K\"ahler current of
\(\alpha+t\{\omega\}\) produces one of \(\alpha+s\{\omega\}\) with the same potential. The
assertion is then vacuous and any \(c>0\) will do.

Otherwise fix a K\"ahler current \(T_0\in\alpha+s\{\omega\}\), let \(\Theta\) be a smooth
representative of \(\alpha+s\{\omega\}\) and let
\(V=\sup\{u:\ u\ \Theta\text{-psh},\ u\le0\}\) be the canonical envelope. By
Theorem~\ref{fpenv:v-psh} the envelope is itself a \(\Theta\)-psh potential (its
upper-semicontinuous regularization changes nothing), and \(V\le0\). Since
\(\Theta+dd^cV\ge0\),
\[
 \Theta+(\varepsilon-s)\omega+dd^cV\ \ge\ (\varepsilon-s)\omega
 =\tfrac13\varepsilon\,\omega>0,
\]
so \(V\) is the potential of a K\"ahler current in \(\alpha+\varepsilon\{\omega\}\). As
\(y\in E_{\mathrm{nK}}(\alpha+\varepsilon\{\omega\})\), every K\"ahler current of that class
has positive Lelong number at \(y\); put \(c=\nu(V,y)>0\).

Let now \(T=\theta_T+dd^cu\) be any K\"ahler current in \(\alpha+t\{\omega\}\) with
\(0<t\le t_0\). Discarding its K\"ahler lower bound, \(u\) is \(\theta_T\)-psh, hence also
\((\theta_T+(s-t)\omega)\)-psh, and \(\theta_T+(s-t)\omega\) has Bott--Chern class
\(\alpha+s\{\omega\}=\{\Theta\}\). So \(\Theta=\theta_T+(s-t)\omega+dd^cf\) for some
chart-smooth real \(f\), and
\(\Theta+dd^c(u-f)=\theta_T+(s-t)\omega+dd^cu\ge0\), i.e. \(u-f\) is \(\Theta\)-psh.
Theorem~\ref{fpenv:v-minimal} gives a constant \(C\) with \(u-f\le V+C\).

Lelong numbers reverse inequalities of quasi-plurisubharmonic functions, are unchanged by
adding a constant, and are unchanged by adding a smooth function; therefore
\[
 c=\nu(V,y)=\nu(V+C,y)\le\nu(u-f,y)=\nu(u,y)=\nu(T,y).
\]
Both \(f\) and \(C\) depend on \(T\), but \(c\) and \(t_0\) do not.
\end{proof}

\begin{lemma}[Transformation of the non-nef locus under one blowup]
\label{ana:nonnef-transformation-blowup}
\uses{ana:nonkahler-locus,ana:blowups}
\source{ct22-restricted-volumes:nonnef}
\lean{TranscendentalNakamaye.Pending.ModificationInputs.nonNefLocus_pullback_of_isBlowupAlong}
\leanok
Let \(g:Z\to Y\) be a blowup of a compact complex manifold along a nowhere dense closed
submanifold, \(\omega\) and \(\omega_Z\) K\"ahler forms on \(Y\) and \(Z\), and
\(\alpha=\{\theta\}\) a big class. Then
\(E_{\mathrm{nn}}(g^*\alpha)=g^{-1}(E_{\mathrm{nn}}(\alpha))\).
\end{lemma}
\begin{proof}
\uses{ana:boucksom-minimal-multiplicity,ana:lelong-additivity,ana:nonkahler-locus-analytic}
\leanok
For \(\subset\), argue contrapositively: let \(z\) with \(g(z)\notin
E_{\mathrm{nn}}(\alpha)\), so some K\"ahler current \(T\in\alpha+\delta\{\omega\}\) has
\(\nu(T,g(z))=0\). Replace \(T\), by Theorem~\ref{ana:nonkahler-locus-analytic}, by a K\"ahler current \(S\) in the same
class whose polar set is exactly \(E_{\mathrm{nK}}(\alpha+\delta\{\omega\})\); then the
potential \(u_S\) is finite at \(g(z)\), so \(u_S\circ g\) is finite at \(z\). Adding the
K\"ahler perturbation \(\epsilon\omega_Z-\delta g^*\omega\), which has a smooth potential,
gives a K\"ahler current in \(g^*\alpha+\epsilon\{\omega_Z\}\) whose potential is
\(u_S\circ g\) up to a smooth term, hence of Lelong number \(0\) at \(z\) — at exceptional
points as well, because no logarithmic term along the exceptional divisor was introduced.
Therefore \(z\notin E_{\mathrm{nn}}(g^*\alpha)\).
For \(\supset\), a point \(y\in E_{\mathrm{nn}}(\alpha)\) has a uniform positive lower
bound for the Lelong numbers of all K\"ahler currents in the small perturbations
\(\alpha+t\{\omega\}\) (Lemma~\ref{ana:boucksom-minimal-multiplicity}); a K\"ahler current
upstairs pushes forward to one downstairs whose Lelong number at \(y\) is at least the
upstairs one, up to the finite correction of the exceptional potential
(Lemma~\ref{ana:lelong-additivity}), so every point of \(g^{-1}(y)\) lies in
\(E_{\mathrm{nn}}(g^*\alpha)\).
\end{proof}

\begin{lemma}[Irreducible subvarieties and relatively open subsets]
\label{ana:irreducible-open-density}
\uses{fpgeom:complex-analytic-subvariety}
\lean{TranscendentalNakamaye.Pending.StrictTransformLocus.subset_of_inter_isOpen_subset,TranscendentalNakamaye.Pending.StrictTransformLocus.exists_subset_of_inter_isOpen_subset_iUnion}
\leanok
Let \(V\) be a closed irreducible analytic subset of a complex manifold \(M\) and
\(U\subset M\) open with \(V\cap U\ne\emptyset\). If \(V\cap U\subset A\) for an
analytic subset \(A\), then \(V\subset A\). If \(V\cap U\subset\bigcup_{m\in\mathbf N}A_m\)
for analytic subsets \(A_m\), then \(V\subset A_m\) for some \(m\).
\end{lemma}
\begin{proof}
\uses{ana:pure-dimension,ana:strict-transform-blowup,ana:component-complement-connected}
\leanok
Suppose \(V\not\subset A\). By Theorem~\ref{ana:pure-dimension} there is a regular point
\(c\notin A\) of \(V\) of dimension \(\dim V\); let \(K\) be the component of the regular
locus through \(c\). Then \(V=\overline K\) (Lemma~\ref{ana:strict-transform-blowup})
and \(K\not\subset A\), so \(K\subset\overline{K\setminus A}\)
(Lemma~\ref{ana:component-complement-connected}). Hence
\(V\subset\overline{V\setminus A}\), and the open set \(U\), which meets \(V\), meets
\(V\setminus A\), contradicting \(V\cap U\subset A\).
For the countable version, \(O=V\cap U\) is a nonempty locally closed subset of a
locally compact Hausdorff space, hence a Baire space, covered by the relatively closed
sets \(O\cap A_m\). Some \(O\cap A_m\) has nonempty interior in \(O\), so there is an
open \(U'\) with \(\emptyset\ne V\cap U\cap U'\subset A_m\); apply the first part.
\end{proof}

\begin{lemma}[Polar sets of smooth analytic singularities]
\label{ana:smooth-pole-set-analytic}
\uses{ana:restricted-volume,fpgeom:complex-analytic-subvariety}
\lean{TranscendentalNakamaye.Pending.StrictTransformLocus.isComplexAnalyticSubset_poleSet_of_smoothOffPole,TranscendentalNakamaye.Pending.StrictTransformLocus.isComplexAnalyticSubset_poleSet}
\leanok
If \(u\) has analytic singularities in every chart --- in the convention of
Definition~\ref{as:smooth-off-pole-type}, so with a remainder only bounded near the
point, and a fortiori if the remainder is smooth --- its polar set \(\{u=-\infty\}\)
is an analytic subset.
\end{lemma}
\begin{proof}
\leanok
Near a point, in a chart, \(u=c\log\sum_j|f_j|^2+b\) with \(c>0\), \(f_j\) holomorphic
and \(b\) real-valued, so \(u=-\infty\) exactly where all \(f_j\) vanish. This is the
one place where the value of \(b\) at a pole is looked at, and all that is used of it
is that it is a real number rather than \(\pm\infty\); no regularity of \(b\)
enters.
\end{proof}

\begin{lemma}[A strict transform is not inside the preimage of a smaller analytic set]
\label{ana:strict-transform-not-in-preimage}
\uses{ana:blowups,ana:irreducible-open-density,ana:blowup-surjective,ana:blowup-exceptional-image}
\lean{TranscendentalNakamaye.Pending.StrictTransformLocus.not_strictTransform_subset_preimage}
\leanok
Let \(g:Z\to Y\) be a blowup of a compact complex manifold along a nowhere dense closed
submanifold and \(W\) a closed irreducible analytic subset with
\(W\not\subset g(\operatorname{Exc} g)\), with strict transform \(\widetilde W\). If \(A\) is
analytic and \(W\not\subset A\), then \(\widetilde W\not\subset g^{-1}(A)\).
\end{lemma}
\begin{proof}
\uses{ana:irreducible-open-density,ana:blowup-surjective,ana:blowup-exceptional-image}
\leanok
Let \(w\in W\setminus g(\operatorname{Exc} g)\). As \(g\) is onto
(Lemma~\ref{ana:blowup-surjective}), \(w=g(z)\) with \(z\notin\operatorname{Exc} g\), so
\(z\in g^{-1}(W)\setminus\operatorname{Exc} g\subset\widetilde W\). Thus
\(\widetilde W\subset g^{-1}(A)\) implies \(W\cap U\subset A\) for the open set
\(U=Y\setminus g(\operatorname{Exc} g)\) (the image is closed,
Lemma~\ref{ana:blowup-exceptional-image}), which meets \(W\), and
Lemma~\ref{ana:irreducible-open-density} gives \(W\subset A\).
\end{proof}

\begin{lemma}[Fibres of a sequence of blowups off its exceptional locus]
\label{ana:blowup-sequence-fibres}
\uses{ana:blowups,ana:blowup-exceptional-image,ana:blowup-surjective}
\lean{TranscendentalNakamaye.Pending.BlowupSequenceIsomorphismOff.eq_of_apply_eq_of_not_mem_exceptionalLocus,TranscendentalNakamaye.Pending.BlowupSequenceIsomorphismOff.isIsomorphismOver_compl_of_image_exceptionalLocus_subset}
\leanok
Let \(\mu:\widetilde X\to X\) be a finite composition of blowups with smooth centres. If
\(y\notin\operatorname{Exc}(\mu)\), then \(\mu^{-1}(\mu(y))=\{y\}\). If moreover \(X\) is
compact and \(T\supseteq\mu(\operatorname{Exc}\mu)\), then \(\mu\) is an isomorphism over
\(X\setminus T\).
\end{lemma}
\begin{proof}
\leanok
Induction on the number of blowups, the empty composition being the identity. Let
\(\mu=f\circ g\) with \(g\) one blowup. Since
\(\operatorname{Exc}(f\circ g)=\operatorname{Exc}(g)\cup g^{-1}\operatorname{Exc}(f)\), a point
\(y\notin\operatorname{Exc}(\mu)\) has \(y\notin\operatorname{Exc}g\) and
\(g(y)\notin\operatorname{Exc}f\). If \(\mu(y')=\mu(y)\), the induction hypothesis for \(f\)
gives \(g(y')=g(y)\), and then \(y'=y\) because the fibre of one blowup through a point off its
exceptional locus is that point (Lemma~\ref{ana:blowup-exceptional-image}).

For the second claim, \(\mu^{-1}(X\setminus T)\) is disjoint from \(\operatorname{Exc}\mu\), so
the differential of \(\mu\) is invertible at each of its points, and \(\mu\) is injective on it
by the first claim. It maps onto \(X\setminus T\) because \(\mu\) is surjective
(Lemma~\ref{ana:blowup-surjective}).
\end{proof}

\begin{lemma}[Strict transform along a sequence of blowups]
\label{ana:strict-transform-sequence}
\uses{ana:blowups,ana:strict-transform-blowup,ana:strict-transform-composition,ana:blowup-surjective,fpgeom:complex-analytic-subvariety}
\lean{TranscendentalNakamaye.Pending.StrictTransformSubvarietySequence.exists_strictTransform_subvariety_of_isBlowupSequence,TranscendentalNakamaye.Pending.StrictTransformSubvarietySequence.not_subset_image_exceptionalLocus_of_image_exceptionalLocus_subset_singular,TranscendentalNakamaye.Pending.StrictTransformSubvarietySequence.isConnected_strictTransform_of_isBlowupSequence}
\leanok
Let \(X\) be a compact complex manifold, \(\mu:\widetilde X\to X\) a finite composition of
blowups with smooth centres, and \(V\subset X\) a closed irreducible analytic subvariety with
\(V\not\subset\mu(\operatorname{Exc}\mu)\). Then the strict transform \(\widetilde V\) of \(V\)
is a closed irreducible analytic subvariety of \(\widetilde X\) of dimension \(\dim V\). If
\(\mu(\operatorname{Exc}\mu)\) is contained in the singular locus of \(V\), then
\(V\not\subset\mu(\operatorname{Exc}\mu)\) holds and \(\widetilde V\) is connected.
\end{lemma}
\begin{proof}
\leanok
Induction on the number of blowups; for the identity, \(\widetilde V=\overline V=V\). Let
\(\mu=f\circ g\) with \(g\) one blowup. By Lemma~\ref{ana:blowup-surjective},
\(V\not\subset f(\operatorname{Exc}f)\) and the strict transform \(\widetilde V^{f}\) is not
contained in \(g(\operatorname{Exc}g)\). By induction \(\widetilde V^{f}\) is a closed
irreducible analytic subvariety of dimension \(\dim V\), so
Lemma~\ref{ana:strict-transform-blowup} applies to it and \(g\); its strict transform under
\(g\) is \(\widetilde V^{f\circ g}\) (Lemma~\ref{ana:strict-transform-composition}).

If \(\mu(\operatorname{Exc}\mu)\subset V_{\mathrm{sing}}\), a regular point of \(V\) lies
outside \(\mu(\operatorname{Exc}\mu)\). Then \(\widetilde V\) is irreducible, hence equal to the
closure of the connected component of its regular locus through any of its regular points (as
in the proof of Lemma~\ref{ana:strict-transform-blowup}), and the closure of a connected set is
connected.
\end{proof}

\begin{lemma}[Strict transforms meet the regular part over every open set]
\label{ana:strict-transform-generic-point}
\uses{ana:blowups,ana:blowup-surjective,ana:blowup-sequence-fibres,ana:strict-transform-sequence,ana:strict-transform-composition,ana:blowup-exceptional-image,ana:irreducible-open-density}
\source{transcendental-nakamaye-copy4-20260911:prop:principalization}
\lean{TranscendentalNakamaye.Pending.StepSixGenericPoint.exists_mem_inter_notMem_image_exceptionalLocus,TranscendentalNakamaye.Pending.StepSixGenericPoint.not_subset_image_exceptionalLocus_of_strictTransform_nonempty,TranscendentalNakamaye.Pending.StepSixGenericPoint.exists_mem_strictTransform_notMem_exceptionalLocus}
\leanok
Let \(X\) be a compact complex manifold, \(f:Y\to X\) a finite composition of blowups with
smooth centres, and \(V\subset X\) a closed irreducible analytic subvariety with strict
transform \(W\).
\begin{enumerate}
\item If \(V\not\subset f(\operatorname{Exc}f)\), every open set meeting \(V\) meets
\(V\setminus f(\operatorname{Exc}f)\).
\item If \(W\ne\emptyset\), then \(V\not\subset f(\operatorname{Exc}f)\), and for every open
\(O\subset X\) meeting \(V\) there is \(y\in W\setminus\operatorname{Exc}f\) with
\(f(y)\in O\).
\end{enumerate}
\end{lemma}
\begin{proof}
\leanok
(i) Induction on the number of blowups; for the identity \(\operatorname{Exc}f=\emptyset\). Let
the sequence be \(f\circ g\) with \(g:Z\to Y\) one blowup, suppose
\(V\not\subset(f\circ g)(\operatorname{Exc}(f\circ g))\), and let \(O\) be open with
\(V\cap O\ne\emptyset\). By Lemma~\ref{ana:blowup-surjective}, \(Y\) is compact, \(f\) and \(g\)
are onto, \(V\not\subset f(\operatorname{Exc}f)\), and the strict transform \(W'\) of \(V\) under
\(f\) is not contained in \(T=g(\operatorname{Exc}g)\). By induction there is
\(x_1\in V\cap O\) with \(x_1\notin f(\operatorname{Exc}f)\); write \(x_1=f(y_1)\), so
\(y_1\notin\operatorname{Exc}f\) and \(y_1\in f^{-1}(V)\setminus\operatorname{Exc}f\subset W'\).
The set \(O_1=f^{-1}(O)\setminus\operatorname{Exc}f\) is open
(Lemma~\ref{ana:strict-transform-composition}) and contains \(y_1\). The strict transform
\(W'\) is a closed irreducible analytic subvariety (Lemma~\ref{ana:strict-transform-sequence})
and \(T\) is analytic (Lemma~\ref{ana:blowup-exceptional-image}); if \(W'\cap O_1\subset T\),
then \(W'\subset T\) by Lemma~\ref{ana:irreducible-open-density}, which is excluded. So there is
\(y_2\in W'\cap O_1\setminus T\). As \(W'\subset f^{-1}(V)\), the point \(x_2=f(y_2)\) lies in
\(V\cap O\). Suppose \(x_2=f(g(z))\) with
\(z\in\operatorname{Exc}(f\circ g)=\operatorname{Exc}g\cup g^{-1}\operatorname{Exc}f\). Since
\(y_2\notin\operatorname{Exc}f\), Lemma~\ref{ana:blowup-sequence-fibres} gives \(g(z)=y_2\);
then \(z\notin\operatorname{Exc}g\) because \(y_2\notin T\), and
\(g(z)=y_2\notin\operatorname{Exc}f\), a contradiction. So
\(x_2\in V\cap O\setminus(f\circ g)(\operatorname{Exc}(f\circ g))\).

(ii) If \(W\ne\emptyset\), its dense subset \(f^{-1}(V)\setminus\operatorname{Exc}f\) contains a
point \(y\). If \(f(y)=f(z)\) with \(z\in\operatorname{Exc}f\), then \(z=y\) by
Lemma~\ref{ana:blowup-sequence-fibres}, which is impossible; so
\(f(y)\in V\setminus f(\operatorname{Exc}f)\). Given \(O\), part (i) gives
\(x\in V\cap O\setminus f(\operatorname{Exc}f)\); write \(x=f(y)\) (\(f\) is onto). Then
\(y\notin\operatorname{Exc}f\), so \(y\in f^{-1}(V)\setminus\operatorname{Exc}f\subset W\).
\end{proof}
Copy-4 obtains (i) from the fact that \(f(\operatorname{Exc}f)\cap V\) is a proper analytic
subset of \(V\). Analyticity of \(f(\operatorname{Exc}f)\) for a sequence of blowups is an
instance of the proper mapping theorem; the induction above uses it for one blowup only, where
it is Lemma~\ref{ana:blowup-exceptional-image}.

\begin{lemma}[Strict transforms over the complement of an analytic set]
\label{ana:strict-transform-dense-over-complement}
\uses{ana:blowups,ana:blowup-surjective,ana:blowup-sequence-fibres,ana:strict-transform-generic-point,ana:irreducible-open-density}
\source{transcendental-nakamaye-copy4-20260911:prop:principalization}
\lean{TranscendentalNakamaye.Pending.StepFourNoComponentContainsW.exists_mem_notMem_exceptionalLocus_notMem_of_inter_strictTransform}
\leanok
Let \(X\) be a compact complex manifold, \(\pi:Y\to X\) a finite composition of blowups with
smooth centres, \(V\subset X\) a closed irreducible analytic subvariety with strict transform
\(W\), and \(Z\subset X\) an analytic subset with \(V\not\subset Z\). Every open set
\(O\subset Y\) meeting \(W\) contains a point \(p\notin\operatorname{Exc}\pi\) with
\(\pi(p)\in V\setminus Z\).
\end{lemma}
\begin{proof}
\leanok
Since \(W\) is the closure of \(\pi^{-1}(V)\setminus\operatorname{Exc}\pi\) and \(O\) is open,
\(O\) contains a point \(y_1\) of \(\pi^{-1}(V)\setminus\operatorname{Exc}\pi\). The manifold
\(Y\) is compact and \(\pi\) is onto (Lemma~\ref{ana:blowup-surjective}), so
\(U=X\setminus\pi(Y\setminus O)\) is open. It contains \(\pi(y_1)\): if
\(\pi(y')=\pi(y_1)\), then \(y'=y_1\in O\) by Lemma~\ref{ana:blowup-sequence-fibres}. Thus
\(U\) meets \(V\), and \(V\cap U\not\subset Z\) by Lemma~\ref{ana:irreducible-open-density}.
The set \(U\setminus Z\) is open (\(Z\) is closed) and meets \(V\), and
\(V\not\subset\pi(\operatorname{Exc}\pi)\) because \(W\ne\emptyset\); by
Lemma~\ref{ana:strict-transform-generic-point} there is
\(x\in V\cap(U\setminus Z)\) outside \(\pi(\operatorname{Exc}\pi)\). Let \(\pi(p)=x\). Then
\(p\notin\operatorname{Exc}\pi\), \(\pi(p)\in V\setminus Z\), and \(p\in O\) because
\(x\in U\).
\end{proof}

\begin{lemma}[Strict transforms leave analytic sets and non-nef loci]
\label{ana:strict-transform-loci}
\uses{ana:blowups,ana:nonkahler-locus,ana:strict-transform-not-in-preimage}
\lean{TranscendentalNakamaye.Pending.StrictTransformLocus.not_strictTransform_subset_nonNefLocus}
\leanok
Let \(g:Z\to Y\) be a blowup of a compact complex manifold along a nowhere dense closed
submanifold and \(W\) a closed irreducible analytic subset with
\(W\not\subset g(\operatorname{Exc} g)\), with strict transform \(\widetilde W\). If
\(A\) is analytic and \(W\not\subset A\), then \(\widetilde W\not\subset g^{-1}(A)\). If
\(n>0\), \(\alpha=\{\theta\}\) is big and \(W\not\subset E_{\mathrm{nn}}(\alpha)\), then
\(\widetilde W\not\subset E_{\mathrm{nn}}(g^*\alpha)\) for all K\"ahler forms.
\end{lemma}
\begin{proof}
\uses{ana:irreducible-open-density,ana:blowup-surjective,ana:blowup-exceptional-image,ana:nonnef-transformation-blowup,ana:nonkahler-locus-analytic}
\leanok
Let \(w\in W\setminus g(\operatorname{Exc} g)\). As \(g\) is onto
(Lemma~\ref{ana:blowup-surjective}), \(w=g(z)\) with \(z\notin\operatorname{Exc} g\), so
\(z\in g^{-1}(W)\setminus\operatorname{Exc} g\subset\widetilde W\). Thus
\(\widetilde W\subset g^{-1}(B)\) implies \(W\cap U\subset B\) for the open set
\(U=Y\setminus g(\operatorname{Exc} g)\) (the image is closed,
Lemma~\ref{ana:blowup-exceptional-image}), which meets \(W\).
For \(B=A\) analytic, Lemma~\ref{ana:irreducible-open-density} gives \(W\subset A\).
For the non-nef locus, \(E_{\mathrm{nn}}(g^*\alpha)=g^{-1}(E_{\mathrm{nn}}(\alpha))\)
(Lemma~\ref{ana:nonnef-transformation-blowup}). A K\"ahler current in
\(\alpha+\frac1{m+1}\{\omega\}\) gives one in \(\alpha+\epsilon\{\omega\}\) with the same
potential when \(\frac1{m+1}<\epsilon\), so
\(E_{\mathrm{nn}}(\alpha)=\bigcup_m E_{\mathrm{nK}}(\alpha+\frac1{m+1}\{\omega\})\), a
countable union of analytic sets (Theorem~\ref{ana:nonkahler-locus-analytic}; these
classes are big). By Lemma~\ref{ana:irreducible-open-density},
\(W\subset E_{\mathrm{nK}}(\alpha+\frac1{m+1}\{\omega\})\subset E_{\mathrm{nn}}(\alpha)\)
for some \(m\), a contradiction.
\end{proof}

\begin{lemma}[Pulling back admissible presentations]
\label{ana:presentation-pullback}
\uses{ana:restricted-volume,ana:blowups}
\lean{TranscendentalNakamaye.Pending.AdmissiblePresentationPullback.pullbackPresentation,TranscendentalNakamaye.Pending.AdmissiblePresentationPullback.coordinateCurrentForm_pullbackPresentation,TranscendentalNakamaye.Pending.AdmissiblePresentationPullback.hasManifoldAnalyticSingularitiesSmoothOffPole_comp,TranscendentalNakamaye.Pending.AdmissiblePresentationPullback.hasManifoldSmoothAnalyticSingularities_comp,TranscendentalNakamaye.Pending.AdmissiblePresentationPullback.ofRepresentative_pullback_congr}
\leanok
Let \(g:Z\to Y\) be a blowup along a nowhere dense closed submanifold, \(\omega,\omega_Z\)
K\"ahler forms with \(c\omega_Z-g^*\omega\ge0\) for some \(c>0\), and
\(T=\theta'+dd^c u\) a current in \(\alpha=\{\theta\}\) with smooth analytic
singularities and \(T+\epsilon\omega\ge0\). For every closed irreducible analytic
\(\widetilde W\not\subset g^{-1}(P(u))\), the current
\(g^*T=g^*\theta'+dd^c(u\circ g)\in g^*\alpha\) has smooth analytic singularities,
satisfies \(g^*T+c\epsilon\omega_Z=g^*(T+\epsilon\omega)+\epsilon(c\omega_Z-g^*\omega)\ge0\),
and is admissible along \(\widetilde W\) at error \(c\epsilon\).
\end{lemma}
\begin{proof}
\leanok
Writing \(\theta'=\theta+dd^c h\) with \(h\) smooth gives
\(g^*\theta'=g^*\theta+dd^c(h\circ g)\), since \(dd^c\) commutes with holomorphic maps; so
\(g^*\theta'\) represents \(g^*\alpha\). Locally \(u=c'\log\sum|f_j|^2+b\) in a chart of
\(Y\); composing with the holomorphic chart expression \(F\) of \(g\) gives
\(u\circ g=c'\log\sum|f_j\circ F|^2+b\circ F\). In the convention of
Definition~\ref{as:smooth-off-pole-type} this is again a presentation: the generators
\(f_j\circ F\) are holomorphic and \(\sum_j|f_j\circ F|^2\) vanishes exactly on
\(F^{-1}\) of the pole set, so \(b\circ F\) is smooth off the new pole set, and it is
bounded near the point because \(F\) is continuous there. The pullback of a
\((\theta'+\epsilon\omega)\)-psh function by a holomorphic map is
\(g^*(\theta'+\epsilon\omega)\)-psh, and adding the semipositive form
\(\epsilon(c\omega_Z-g^*\omega)\) gives positivity. Finally \(u\circ g\) is not
\(-\infty\) on any open set: an open subset of \(Z\) contains a point off the preimage of
the centre, near which \(g\) is open, and \(u\) is finite somewhere on every open set.
\end{proof}

\begin{lemma}[Transport of regular patches]
\label{ana:patch-transport}
\uses{ana:restricted-volume,ana:regular-transport}
\lean{TranscendentalNakamaye.Pending.RegularPatchPullback.exists_regularPatch_pullback}
\leanok
Let \(g:Z\to Y\) be holomorphic, \(W\subset Y\) and \(\widetilde W\subset Z\) closed
irreducible analytic subsets of the same dimension \(k\) with
\(\widetilde W\cap O=g^{-1}(W)\cap O\) for an open \(O\), \(\gamma\) a regular
parametrization of \(W\) defined at \(w_0\), and \(z\in O\setminus\operatorname{Exc} g\)
with \(g(z)=\gamma(w_0)\). There are a regular parametrization \(\tilde\gamma\) of
\(\widetilde W\) and \(r>0\) such that for \(|w-w_0|<r\):
\(\tilde\gamma(w-w_0)\in O\setminus\operatorname{Exc} g\),
\(g(\tilde\gamma(w-w_0))=\gamma(w)\), and for every closed real \((1,1)\)-form \(\beta\) and every
function \(v\) smooth near \(\gamma(w)\),
\((\tilde\gamma^*g^*\beta)(w-w_0)=(\gamma^*\beta)(w)\) and
\((\tilde\gamma^*dd^c(v\circ g))(w-w_0)=(\gamma^*dd^c v)(w)\).
\end{lemma}
\begin{proof}
\leanok
In the charts at \(z\) and at the centre of \(\gamma\), \(g\) is a holomorphic map \(F\)
whose differential at the chart image of \(z\) is invertible, so it has a local
holomorphic inverse. If \(\Phi\) flattens \(W\) with \(\gamma(w)=\Phi^{-1}(w,0)\), then
\((\Phi\circ F)-(w_0,0)\) flattens \(\widetilde W\) at \(z\) (translation by a point of
\(\mathbf C^k\times0\) preserves the plane), and its parametrization is
\(\tilde\gamma(v)=F^{-1}(\gamma(v+w_0))\). Hence \(\gamma=F\circ\tilde\gamma\circ(\cdot-w_0)\)
near \(w_0\), and the identities follow from the chain rule, the coordinate formula for
pullbacks of forms, and \(dd^c(v\circ F)=F^*dd^c v\).
\end{proof}

\begin{lemma}[Transport of regular parametrizations down a local biholomorphism]
\label{ana:patch-pushforward}
\uses{ana:restricted-volume,ana:regular-transport}
\source{ct22-restricted-volumes:invar}
\lean{TranscendentalNakamaye.Pending.RegularPatchPushforward.exists_regularPatch_pushforward}
\leanok
Let \(g:Z\to Y\) be holomorphic, \(W\subset Y\) and \(\widetilde W\subset Z\) closed
irreducible analytic subsets of the same dimension \(k\) with
\(\widetilde W\cap O=g^{-1}(W)\cap O\) for an open \(O\), \(\tilde\gamma\) a regular
parametrization of \(\widetilde W\), and \(v_0\) a parameter with
\(\tilde\gamma(v_0)\in O\setminus\operatorname{Exc} g\). There are a regular
parametrization \(\gamma\) of \(W\) centred at \(g(\tilde\gamma(v_0))\) and \(r>0\) such
that for \(|v-v_0|<r\): \(\tilde\gamma(v)\in O\setminus\operatorname{Exc} g\),
\(g(\tilde\gamma(v))=\gamma(v-v_0)\), and for every closed real \((1,1)\)-form \(\beta\)
and every function \(u\) smooth near \(\gamma(v-v_0)\),
\((\tilde\gamma^*g^*\beta)(v)=(\gamma^*\beta)(v-v_0)\) and
\((\tilde\gamma^*dd^c(u\circ g))(v)=(\gamma^*dd^c u)(v-v_0)\).
\end{lemma}
\begin{proof}
\leanok
Read \(g\) in the chart at the centre of \(\tilde\gamma\) and the chart at
\(g(\tilde\gamma(v_0))\); it becomes a holomorphic map \(F\), and its differential at the
chart image \(u_0\) of \(\tilde\gamma(v_0)\) is invertible because \(\tilde\gamma(v_0)\)
is not exceptional, so \(F\) has a local holomorphic inverse \(F^{-1}\) at \(F(u_0)\).
Let \(\widetilde\Phi\) flatten \(\widetilde W\) near \(u_0\), with
\(\tilde\gamma(v)=\widetilde\Phi^{-1}(v,0)\). Then
\((\widetilde\Phi\circ F^{-1})-(v_0,0)\) is defined near the chart image of
\(g(\tilde\gamma(v_0))\) and flattens \(W\) there. Indeed a nearby point \(y\) is
\(F(x)\) for the unique \(x=F^{-1}(y)\) close to \(u_0\); that \(x\) lies in \(O\), so
\(y\in W\) if and only if \(x\in\widetilde W\), if and only if \(\widetilde\Phi(x)\) lies
in \(\mathbf C^k\times0\); and subtracting the point \((v_0,0)\) of
\(\mathbf C^k\times0\) preserves that plane. It sends \(g(\tilde\gamma(v_0))\) to the
origin, again because \(\widetilde\Phi(u_0)=(v_0,0)\). Its parametrization is therefore
\(\gamma(v)=F(\tilde\gamma(v+v_0))\), so \(\gamma(v-v_0)=F(\tilde\gamma(v))\) and
\(g(\tilde\gamma(v))=\gamma(v-v_0)\). The two identities follow from the chain rule, the
coordinate formula for pullbacks of forms, and \(dd^c(u\circ F)=F^*dd^c u\).
\end{proof}

\begin{lemma}[The mass across the centre, with the downstairs weight abstract]
\label{ana:mass-across-centre-abstract}
\uses{ana:restricted-volume,ana:blowups,ana:strict-transform-blowup}
\source{ct22-restricted-volumes:invar}
\lean{TranscendentalNakamaye.Pending.RestrictedVolumePushforwardMass.acTopMass_le_of_patchDensity_le_general}
\leanok
Let \(g:Z\to Y\) be holomorphic, \(W\subset Y\) and \(\widetilde W\subset Z\) closed
irreducible analytic subsets of the same dimension, and \(A\subset Z\) a closed analytic
set containing \(\operatorname{Exc} g\), off which \(g\) is injective, with
\(\widetilde W\not\subset A\) and \(\widetilde W\setminus A=g^{-1}(W)\setminus A\). Let
\(T'\) be admissible along \(\widetilde W\). Let \(\mathcal A\subset Y\) be a set,
\(\Phi\ge0\) a weight on pairs (regular parametrization of \(W\), parameter) and
\(B\in[0,\infty]\), and assume
\begin{enumerate}
\item \(g(z)\notin\mathcal A\) whenever \(z\notin A\) and the potential of \(T'\) is
finite at \(z\);
\item whenever a regular parametrization \(\tilde\gamma\) of \(\widetilde W\) at a
parameter \(v'\) with \(\tilde\gamma(v')\notin A\) and a regular parametrization
\(\gamma\) of \(W\) at a parameter \(v\) satisfy \(g(\tilde\gamma(v'))=\gamma(v)\) and
the two identities of Lemma~\ref{ana:patch-pushforward}, and the potential of \(T'\) is
finite at \(\tilde\gamma(v')\), the restricted top-power density of \(T'\) along
\(\tilde\gamma\) at \(v'\) is at most \(\Phi(\gamma,v)\); and
\item \(\sum_k\int_{S_k}\Phi(\gamma_k,\cdot)\le B\) for every finite family
\((\gamma_k,S_k)\) of regular parametrizations of \(W\) with measurable parameter pieces
whose images are pairwise disjoint subsets of \(W_{\mathrm{reg}}\) avoiding
\(\mathcal A\).
\end{enumerate}
Then the restricted absolutely continuous top mass of \(T'\) along \(\widetilde W\) is at
most \(B\).
\end{lemma}
\begin{proof}
\uses{ana:patch-pushforward,ana:patch-parameter-null,ana:irreducible-open-density}
\leanok
The mass of \(T'\) along \(\widetilde W\) is the supremum, over finite families of
regular parametrizations of \(\widetilde W\) with measurable parameter pieces whose
images are pairwise disjoint subsets of \(\widetilde W_{\mathrm{reg}}\) avoiding the
polar set of \(T'\), of the sums of the integrals of the parameter densities. Fix such a
family and one of its pieces. By Lemma~\ref{ana:irreducible-open-density} the analytic
set \(A\) contains no nonempty relatively open part of \(\widetilde W\), so by
Lemma~\ref{ana:patch-parameter-null} the parameters of the piece whose point lies in
\(A\) form a Lebesgue-null set and may be discarded without changing the integral. Every
remaining parameter \(v_0\) has \(\tilde\gamma(v_0)\notin A\), hence
\(\tilde\gamma(v_0)\notin\operatorname{Exc} g\), so Lemma~\ref{ana:patch-pushforward}
supplies a radius \(r(v_0)>0\) and a regular parametrization of \(W\) over
\(B(v_0,r(v_0))\). Countably many of these balls already cover what is left (Lindel\"of),
and replacing them by their disjointification gives countably many disjoint measurable
sets \(D_j\) whose union carries the whole integral; hence that integral is the supremum
over \(N\) of the partial sums over \(j<N\).

Fix \(N\). Translate each \(D_j\) by \(-v_0^{(j)}\), a measure-preserving map of
\(\mathbf C^k\). The pushed parametrizations together with the translated pieces form a
finite family for \(W\): the pieces are measurable and land in the parameter domains, their
images consist of regular points of \(W\), they avoid \(\mathcal A\) by (i), and
they are pairwise disjoint, because \(g\) is injective off \(A\), a regular
parametrization is injective on its parameter domain, and the original images were
pairwise disjoint. By (ii) the integrand does not decrease under the translation, so the
stage-\(N\) sum is at most \(\sum_k\int_{S_k}\Phi(\gamma_k,\cdot)\) for that family,
hence at most \(B\) by (iii). Letting \(N\to\infty\) and then taking the supremum over
families gives the inequality.
\end{proof}

\begin{lemma}[The mass across the centre, from a pointwise density comparison]
\label{ana:mass-across-centre}
\uses{ana:mass-across-centre-abstract}
\source{ct22-restricted-volumes:invar}
\lean{TranscendentalNakamaye.Pending.RestrictedVolumePushforwardMass.acTopMass_le_of_patchDensity_le,TranscendentalNakamaye.Pending.RestrictedVolumePushforwardMass.acTopMass_le_of_patchDensity_le_of_exceptionalLocus,TranscendentalNakamaye.Pending.RestrictedVolumePushforwardMass.acTopMass_le_of_patchDensity_le_of_isBlowupAlong}
\leanok
Keep the geometric hypotheses of Lemma~\ref{ana:mass-across-centre-abstract} on \(g\),
\(W\), \(\widetilde W\) and \(A\). Let \(T'\) be admissible along \(\widetilde W\) and
\(T\) admissible along \(W\), and assume
\begin{enumerate}
\item the potential of \(T\) is finite at \(g(z)\) whenever \(z\notin A\) and the
potential of \(T'\) is finite at \(z\); and
\item whenever a regular parametrization \(\tilde\gamma\) of \(\widetilde W\) at a
parameter \(v'\) with \(\tilde\gamma(v')\notin A\) and a regular parametrization
\(\gamma\) of \(W\) at a parameter \(v\) satisfy \(g(\tilde\gamma(v'))=\gamma(v)\) and
the two identities of Lemma~\ref{ana:patch-pushforward}, and the potential of \(T'\) is
finite at \(\tilde\gamma(v')\), the restricted top-power density of \(T'\) along
\(\tilde\gamma\) at \(v'\) is at most that of \(T\) along \(\gamma\) at \(v\).
\end{enumerate}
Then the restricted absolutely continuous top mass of \(T'\) along \(\widetilde W\) is at
most that of \(T\) along \(W\). For the blowup \(g\) of a compact \(Y\) along a nowhere
dense closed submanifold \(C\), with \(\widetilde W\) the strict transform of a \(W\)
satisfying \(W\not\subset g(\operatorname{Exc} g)\), the set \(A=\operatorname{Exc} g\)
satisfies the geometric hypotheses.
\end{lemma}
\begin{proof}
\uses{ana:mass-across-centre-abstract,ana:blowup-exceptional-image,ana:analytic-preimage,ana:strict-transform-loci,ana:blowups}
\leanok
Apply Lemma~\ref{ana:mass-across-centre-abstract} with \(\mathcal A\) the polar set of
\(T\), with \(\Phi(\gamma,v)\) the restricted top-power density of \(T\) along \(\gamma\)
at \(v\), and with \(B\) the restricted absolutely continuous top mass of \(T\) along
\(W\). Hypotheses (i) and (ii) of that lemma are the present (i) and (ii), and its
hypothesis (iii) holds because the mass of \(T\) along \(W\) is by definition the
supremum of exactly those sums.

For the blowup, take \(A=\operatorname{Exc} g\). A fibre of \(g\) through a
non-exceptional point is a singleton, so \(g\) is injective off
\(\operatorname{Exc} g\), and the exceptional locus is the full preimage
\(g^{-1}(g(\operatorname{Exc} g))\) of its image; that image is analytic
(Lemma~\ref{ana:blowup-exceptional-image}), so \(\operatorname{Exc} g\) is analytic
(Lemma~\ref{ana:analytic-preimage}). Off \(\operatorname{Exc} g\) the strict transform of
\(W\) is \(g^{-1}(W)\). Finally \(\widetilde W\subset\operatorname{Exc} g
=g^{-1}(g(\operatorname{Exc} g))\) would force \(W\subset g(\operatorname{Exc} g)\)
(Lemma~\ref{ana:strict-transform-loci}), which is excluded.
\end{proof}

\begin{lemma}[Measurability of the parameter density of an admissible presentation]
\label{ana:patch-density-measurable}
\uses{ana:restricted-volume,ana:patch-density-continuity,ana:smooth-pole-set-analytic}
\lean{TranscendentalNakamaye.Pending.PatchDensityMeasurability.continuousAt_patchDensity,TranscendentalNakamaye.Pending.PatchDensityMeasurability.aemeasurable_ofReal_patchDensity}
\leanok
Let \(T=\theta+dd^c u\) be admissible along \(V\) at error \(\epsilon\), so that \(u\) has
analytic singularities in every chart in the sense of
Definition~\ref{as:smooth-off-pole-type}, and let \(\gamma\) be a
regular parametrization of \(V\). Then the parameter density
\(w\mapsto\mathcal D_{T+\epsilon\omega}(\gamma,w)\) is continuous at every parameter of
the domain of \(\gamma\) whose point does not lie in the polar set \(\{u=-\infty\}\).
Consequently it is measurable on every measurable parameter set contained in that domain
whose image avoids the polar set; in particular on every piece of a family computing the
restricted absolutely continuous top mass of \(T\).
\end{lemma}
\begin{proof}
\uses{ana:patch-density-continuity}
\leanok
Fix a parameter \(w_0\) with \(\gamma(w_0)\notin\{u=-\infty\}\) and work in the chart at
the centre of \(\gamma\). Near \(\gamma(w_0)\) one has \(u=c\log\sum_j|f_j|^2+b\) with
\(c>0\), the \(f_j\) holomorphic and \(b\) smooth, and not all \(f_j\) vanish at
\(\gamma(w_0)\); so on a neighbourhood \(U\) of that point \(u\) coincides with the smooth
real function \(c\log\sum_j|f_j|^2+b\). The operator \(dd^c\) depends only on the germ, so
on \(U\) it agrees with \(dd^c\) of that smooth function, which is a smooth two-form.
Hence the coordinate two-form of \(T+\epsilon\omega\), the sum of the smooth coefficient
of \(\theta+\epsilon\omega\) and of \(dd^c u\), is continuous on \(U\). The remaining
steps are those of Lemma~\ref{ana:patch-density-continuity}: \(\gamma\) is smooth and
\(w\mapsto D\gamma(w)\) is continuous on the parameter domain, composition of a two-form
with a continuously varying linear map is continuous, and the top-power density is a
polynomial in the values of the form on the standard real frame. A function continuous on
a set is measurable there.
\end{proof}

\begin{lemma}[The mass across the centre, from a liminf density comparison]
\label{ana:mass-across-centre-fatou}
\uses{ana:mass-across-centre-abstract,ana:patch-density-measurable}
\source{ct22-restricted-volumes:invar}
\lean{TranscendentalNakamaye.Pending.RestrictedVolumePushforwardFatou.acTopMass_le_liminf_of_patchDensity_liminf_le,TranscendentalNakamaye.Pending.RestrictedVolumePushforwardFatou.acTopMass_le_liminf_of_patchDensity_liminf_le_of_isBlowupAlong}
\leanok
Keep the geometric hypotheses of Lemma~\ref{ana:mass-across-centre} on \(g\), \(W\),
\(\widetilde W\) and \(A\). Let \(T'\) be admissible along \(\widetilde W\) and let
\((T_m)_{m\in\mathbf N}\) be a sequence of currents admissible along \(W\) such that
\begin{enumerate}
\item for every \(m\), the potential of \(T_m\) is finite at \(g(z)\) whenever
\(z\notin A\) and the potential of \(T'\) is finite at \(z\); and
\item whenever regular parametrizations \(\tilde\gamma\) of \(\widetilde W\) at \(v'\) and
\(\gamma\) of \(W\) at \(v\) are as in hypothesis (ii) of
Lemma~\ref{ana:mass-across-centre}, the restricted top-power density of \(T'\) along
\(\tilde\gamma\) at \(v'\) is at most
\(\liminf_{m}\mathcal D_{T_m}(\gamma,v)\).
\end{enumerate}
Then the restricted absolutely continuous top mass of \(T'\) along \(\widetilde W\) is at
most \(\liminf_m\) of the restricted absolutely continuous top masses of \(T_m\) along
\(W\).
\end{lemma}
\begin{proof}
\uses{ana:mass-across-centre-abstract,ana:patch-density-measurable}
\leanok
Apply Lemma~\ref{ana:mass-across-centre-abstract} with \(\mathcal A=\bigcup_m\{u_m=-\infty\}\), the union of the polar sets of the
\(T_m\), which hypothesis (i) makes admissible, and with
\(\Phi(\gamma,v)=\liminf_m\mathcal D_{T_m}(\gamma,v)\), which hypothesis (ii) makes
admissible. Let \((\gamma_k,S_k)_{k<K}\) be a finite family for \(W\) avoiding
\(\mathcal A\). Each \(S_k\) is a measurable parameter set inside the domain of
\(\gamma_k\) whose image avoids every \(\{u_m=-\infty\}\), so by
Lemma~\ref{ana:patch-density-measurable} each \(\mathcal D_{T_m}(\gamma_k,\cdot)\) is
measurable on \(S_k\), and Fatou's lemma gives
\[\int_{S_k}\liminf_m\mathcal D_{T_m}(\gamma_k,\cdot)\le
\liminf_m\int_{S_k}\mathcal D_{T_m}(\gamma_k,\cdot).\]
Summing over \(k<K\) and using that \(\liminf\) is superadditive along a finite sum,
\[\sum_{k<K}\int_{S_k}\Phi(\gamma_k,\cdot)\le
\liminf_m\sum_{k<K}\int_{S_k}\mathcal D_{T_m}(\gamma_k,\cdot).\]
For each fixed \(m\) the family \((\gamma_k,S_k)\) avoids \(\{u_m=-\infty\}\), so the inner
sum is at most the mass of \(T_m\) along \(W\); hence the right-hand side is at most
\(\liminf_m\) of those masses, which is the required uniform bound.
\end{proof}

\begin{lemma}[Zero sets of holomorphic functions are null]
\label{ana:holomorphic-zero-set-null}
\lean{TranscendentalNakamaye.Pending.HolomorphicZeroSetNull.measure_inter_preimage_zero_eq_zero}
\leanok
Let \(E\) be a finite-dimensional complex vector space with a Haar measure \(\mu\),
\(U\subset E\) convex and open, and \(f\) holomorphic on \(U\) with \(f(x_1)\ne0\) for some
\(x_1\in U\). Then \(\mu(U\cap f^{-1}(0))=0\).
\end{lemma}
\begin{proof}
\leanok
In one variable the zeros of a holomorphic function on a connected open set that is not
identically zero are discrete, hence null. A measurable set whose intersection with every
complex line \(y+\mathbf C v\) is null is null, by Fubini and translation invariance. Let
\(x_0\in U\) and \(\delta>0\) with \(f\ne0\) on \(B(x_1,\delta)\subset U\). For \(y\) close
to \(x_0\), the line \(y+\mathbf C(x_1-x_0)\) meets \(B(x_1,\delta)\), and its intersection
with the convex set \(U\) is connected; so \(U\cap f^{-1}(0)\) meets it in a null set.
Thus \(U\cap f^{-1}(0)\) is null near every point, and countably many such neighbourhoods
cover it.
\end{proof}

\begin{lemma}[Analytic sets are null in regular parameters]
\label{ana:patch-parameter-null}
\uses{ana:restricted-volume,ana:holomorphic-zero-set-null}
\lean{TranscendentalNakamaye.Pending.PatchParameterNull.volume_parameterDomain_inter_preimage_eq_zero}
\leanok
Let \(\gamma:D\to V\) be a regular parametrization of a closed irreducible analytic subset
\(V\) and \(A\) an analytic subset containing no nonempty relatively open subset of \(V\).
Then \(D\cap\gamma^{-1}(A)\) is Lebesgue-null.
\end{lemma}
\begin{proof}
\leanok
The parametrization maps open subsets of \(D\) onto relatively open subsets of \(V\), so
\(S=D\cap\gamma^{-1}(A)\) has empty interior. Near \(w\in S\), \(A\) is the common zero set
of holomorphic \(h_1,\dots,h_m\) in a chart, so on a small ball \(B\ni w\),
\(S\cap B=\bigcap_j (h_j\circ\gamma)^{-1}(0)\). Not every \(h_j\circ\gamma\) vanishes
identically on \(B\), since \(S\) has empty interior; so \(S\cap B\) is null by
Lemma~\ref{ana:holomorphic-zero-set-null}. Countably many such balls cover \(S\).
\end{proof}

\begin{lemma}[The restricted mass does not decrease under pullback]
\label{ana:mass-pullback}
\uses{ana:restricted-volume,ana:presentation-pullback}
\lean{TranscendentalNakamaye.Pending.RestrictedVolumePullbackLower.acTopMass_le_pullbackPresentation,TranscendentalNakamaye.Pending.RestrictedVolumePullbackLower.patchDensity_le_pullback}
\leanok
In the situation of Lemma~\ref{ana:presentation-pullback}, with \(Y\) compact,
\(W\not\subset g(\operatorname{Exc} g)\), \(W\not\subset P(u)\) and \(\widetilde W\) the
strict transform of \(W\), of dimension \(k=\dim W\),
\[
 \int_{W_{\mathrm{reg}}}\bigl((T+\epsilon\omega)|_{W_{\mathrm{reg}}}\bigr)_{\mathrm{ac}}^k
 \le\int_{\widetilde W_{\mathrm{reg}}}\bigl((g^*T+c\epsilon\omega_Z)|_{\widetilde W_{\mathrm{reg}}}\bigr)_{\mathrm{ac}}^k .
\]
\end{lemma}
\begin{proof}
\uses{ana:irreducible-open-density,ana:blowup-exceptional-image,ana:blowup-surjective,ana:patch-transport,ana:patch-parameter-null}
\leanok
Fix finitely many disjoint measurable pieces \(D_i\) in regular parametrizations
\(\gamma_i\) of \(W\) avoiding \(P(u)\). The set \(E=g(\operatorname{Exc} g)\) is a closed
analytic set (Lemma~\ref{ana:blowup-exceptional-image}) containing no relatively open part
of \(W\) (Lemma~\ref{ana:irreducible-open-density}), so the parameters over \(E\) form a
null set (Lemma~\ref{ana:patch-parameter-null}) and may be removed. Off
\(\operatorname{Exc} g\), \(\widetilde W=g^{-1}(W)\). Each remaining parameter \(w_0\) lies
over \(g(z)\) with \(z\notin\operatorname{Exc} g\) (\(g\) is onto,
Lemma~\ref{ana:blowup-surjective}), so Lemma~\ref{ana:patch-transport} gives a
parametrization of \(\widetilde W\) on a ball \(B(w_0,r)\). Countably many such balls
cover the good part of \(D_i\); disjointify them into \(D_{ij}\), with parametrizations
\(\tilde\gamma_{ij}\) centred over \(w_{ij}\). On \(D_{ij}\),
\[
 \tilde\gamma_{ij}^*(g^*T+c\epsilon\omega_Z)(w-w_{ij})
 =\gamma_i^*(T+\epsilon\omega)(w)+\tilde\gamma_{ij}^*\bigl(\epsilon(c\omega_Z-g^*\omega)\bigr)(w-w_{ij}),
\]
a sum of semipositive \((1,1)\)-forms, so the top-power density of the left side dominates
that of \(\gamma_i^*(T+\epsilon\omega)\). For each \(N\) the translated pieces
\(D_{ij}-w_{ij}\), \(j<N\), are pieces of regular parametrizations of \(\widetilde W\)
avoiding \(g^{-1}(P(u))\), with disjoint images because their images under \(g\) are
disjoint and parametrizations are injective. By translation invariance of Lebesgue measure
their total mass is at least \(\sum_i\sum_{j<N}\int_{D_{ij}}\) of the density of
\(\gamma_i^*(T+\epsilon\omega)\); letting \(N\to\infty\) (monotone convergence) and taking
the supremum over the pieces gives the claim.
\end{proof}

\begin{lemma}[Restricted volumes do not decrease under one blowup]
\label{ana:rv-le-pullback}
\uses{ana:restricted-volume,ana:blowups,ana:nonkahler-locus,ana:strict-transform-blowup}
\source{ct22-restricted-volumes:invar}
\lean{TranscendentalNakamaye.Pending.RestrictedVolumeOneBlowup.numericalRestrictedVolume_le_pullback}
\leanok
Let \(Y\) be a compact complex manifold, \(g:Z\to Y\) a blowup along a nowhere
dense closed submanifold, \(\alpha=\{\theta\}\) a big class on \(Y\), and \(W\) a
positive-dimensional closed irreducible analytic subset with
\(W\not\subset g(\operatorname{Exc}(g))\) and strict transform \(\widetilde W\) of
dimension \(\dim W\). For all K\"ahler forms on \(Y\) and \(Z\),
\(\operatorname{RV}_{Y\mid W}(\alpha)\le\operatorname{RV}_{Z\mid\widetilde W}(g^*\alpha)\).
\end{lemma}
\begin{proof}
\uses{ana:strict-transform-loci,ana:smooth-pole-set-analytic,ana:mass-pullback,ana:restricted-volume-monotone,ana:blowup-compact,gnpp:kahler-domination}
\leanok
If \(W\subset E_{\mathrm{nn}}(\alpha)\) the left side is \(0\). Otherwise
\(\widetilde W\not\subset E_{\mathrm{nn}}(g^*\alpha)\) (Lemma~\ref{ana:strict-transform-loci}),
and both sides are infima over \(\epsilon>0\) of the admissible mass envelopes
(Lemma~\ref{ana:restricted-volume-monotone}). Pick \(c>0\) with
\(c\omega_Z-g^*\omega\ge0\) (\(Z\) is compact, Lemma~\ref{ana:blowup-compact}, so
Lemma~\ref{gnpp:kahler-domination} applies). For an \(\epsilon\)-admissible \(T\) along
\(W\), its polar set is analytic (Lemma~\ref{ana:smooth-pole-set-analytic}) and does not
contain \(W\), so \(\widetilde W\not\subset g^{-1}(P(T))\)
(Lemma~\ref{ana:strict-transform-loci}); by Lemma~\ref{ana:mass-pullback} the envelope of
\(W\) at \(\epsilon\) is at most the envelope of \(\widetilde W\) at \(c\epsilon\). Taking
infima gives the inequality.
\end{proof}

\begin{lemma}[Change of background form in the admissible mass envelope]
\label{ana:restricted-mass-background-change}
\uses{ana:restricted-volume}
\source{ct22-restricted-volumes:sect2}
\lean{TranscendentalNakamaye.Pending.RestrictedMassBackgroundChange.restrictedMassEnvelope_le_of_isSemipositive}
\leanok
Let \(\omega,\omega'\) be K\"ahler forms on a complex manifold \(X\), let
\(\alpha\) be a Bott--Chern class, \(V\subset X\) a closed irreducible analytic
subvariety and \(\epsilon,\epsilon'>0\) errors such that
\(\epsilon'\omega'-\epsilon\omega\) is a semipositive \((1,1)\)-form. Then
\[
 \mathcal E_{\omega}(\alpha,V,\epsilon)\le\mathcal E_{\omega'}(\alpha,V,\epsilon')
\]
for the admissible mass envelopes.
\end{lemma}
\begin{proof}
\uses{ana:restricted-volume,ana:rv-kahler-independence}
\leanok
Put \(c=\epsilon'/\epsilon>0\). Then
\(c\omega'-\omega=\epsilon^{-1}(\epsilon'\omega'-\epsilon\omega)\) is semipositive
and \(c\epsilon=\epsilon'\), so this is the case treated in the proof of
Lemma~\ref{ana:rv-kahler-independence}: an \(\epsilon\)-admissible current
\(T=\theta+\epsilon\omega+dd^cu\) along \(V\) for \(\omega\) is, with the same
representative and the same potential, \(c\epsilon\)-admissible for \(\omega'\),
because
\(\theta+c\epsilon\omega'+dd^cu=(\theta+\epsilon\omega+dd^cu)+\epsilon(c\omega'-\omega)\);
and off the poles, on every regular patch, the restricted form grows by the
semipositive \((1,1)\)-form \(\epsilon(c\omega'-\omega)|_V\), so the absolutely
continuous top-power density does not decrease and neither does the restricted
mass. Taking the supremum over \(\epsilon\)-admissible currents for \(\omega\)
gives the claim.
\end{proof}

\begin{lemma}[\(\log|f|^2\) is pluriharmonic off the zeros of \(f\)]
\label{ana:log-norm-pluriharmonic}
\lean{TranscendentalNakamaye.Pending.PluriharmonicLogNorm.smoothDDC_re_comp_eq_zero,TranscendentalNakamaye.Pending.PluriharmonicLogNorm.smoothDDC_log_normSq_eq_zero}
\leanok
Let \(f\) be holomorphic on an open subset of \(\mathbf C^n\) and \(f(z)\neq0\). Then
\(dd^c\log|f|^2=0\) near \(z\).
\end{lemma}
\begin{proof}
\leanok
For \(g\) holomorphic, \(\operatorname{Re}g\) is the composition of the holomorphic \(g\)
with a real-linear map, and \(dd^c\) of a real-linear map vanishes, so
\(dd^c\operatorname{Re}g=0\) by naturality of \(dd^c\) under holomorphic maps. On a
neighbourhood of \(z\) where \(f/f(z)\) avoids the negative real axis, the principal
logarithm gives the holomorphic \(g=2\log(f/f(z))+\log|f(z)|^2\) with
\(\operatorname{Re}g=\log|f|^2\); hence \(dd^c\log|f|^2=0\) there.
\end{proof}

\begin{lemma}[Off its poles a function with smooth analytic singularities is smooth]
\label{ana:smooth-off-poles}
\uses{ana:smooth-pole-set-analytic}
\lean{TranscendentalNakamaye.FirstPass.AdmissibleCurrentFormPositivity.exists_smooth_eventuallyEq_of_hasSmoothAnalyticSingularitiesAt,TranscendentalNakamaye.FirstPass.AdmissibleCurrentFormPositivity.SmoothOffPoleSet.of_hasManifoldSmoothAnalyticSingularities,TranscendentalNakamaye.FirstPass.AdmissibleCurrentFormPositivity.isOneOne_and_semipositive_add_smoothDDC_of_smoothOffPole}
\leanok
Let \(u\) be an extended-real function with smooth analytic singularities on a complex
manifold and \(u(x)\neq-\infty\). Then in every chart around \(x\) the chart expression of
\(u\) agrees, on a neighbourhood, with a smooth real function.
\end{lemma}
\begin{proof}
\leanok
Near \(x\) one has \(u=c\log\sum_j|f_j|^2+b\) with \(c>0\), the \(f_j\) holomorphic and
\(b\) smooth and real. Since \(u(x)\neq-\infty\), some \(f_j(x)\neq0\), so
\(\sum_j|f_j|^2>0\) on a neighbourhood and \(\log\sum_j|f_j|^2\) is smooth there; hence so
is \(u\).
\end{proof}

\begin{lemma}[Local holomorphic inverse]
\label{ana:local-holomorphic-inverse}
\lean{TranscendentalNakamaye.FirstPass.exists_localHolomorphicEquivAt}
\leanok
Let \(F\) be holomorphic on an open \(N\subset\mathbf C^n\) and let the complex
differential of \(F\) at \(p\in N\) be injective. Then \(F\) restricts to a
biholomorphism between open neighbourhoods of \(p\) and of \(F(p)\).
\end{lemma}
\begin{proof}
\leanok
An injective endomorphism of a finite-dimensional space is invertible, so \(dF_p\) is a
linear isomorphism and the inverse function theorem gives open neighbourhoods on which
\(F\) is a diffeomorphism. The inverse of a bijective holomorphic map with holomorphic
local inverse is holomorphic, the differential being the inverse of \(dF\), which is
complex linear.
\end{proof}

\begin{lemma}[Pushforward of a \(\theta\)-plurisubharmonic function along an unramified map]
\label{ana:theta-psh-pushforward}
\uses{ana:nonkahler-locus,ana:local-holomorphic-inverse}
\lean{TranscendentalNakamaye.Pending.ThetaPshPushforward.thetaPshOff_comp_section}
\leanok
Let \(g:Z\to Y\) be holomorphic, \(A\subset Y\) closed with
\(g(z)\notin A\Rightarrow z\notin\operatorname{Exc}(g)\), and let \(s:Y\to Z\) satisfy
\(g\circ s=\operatorname{id}\) off \(A\) and \(s\circ g=\operatorname{id}\) over
\(Y\setminus A\). Let \(w\) be \(\theta_Z\)-plurisubharmonic on \(Z\) and let \(r\) be a
real function which is, near every point over \(Y\setminus A\), a smooth local potential of
\(\theta_Z-g^*\theta_Y\). Then \((w+r)\circ s\) is \(\theta_Y\)-plurisubharmonic on
\(Y\setminus A\).
\end{lemma}
\begin{proof}
\uses{ana:local-holomorphic-inverse}
Fix \(y\notin A\) and put \(z=s(y)\); then \(z\notin\operatorname{Exc}(g)\), so \(g\) has
a holomorphic local inverse \(\sigma\) at \(y\) with \(\sigma(y)=z\)
(Lemma~\ref{ana:local-holomorphic-inverse}), and \(\sigma=s\) near \(y\) off \(A\) by
uniqueness of the preimage. Let \(\rho_Y\) be a smooth local potential of \(\theta_Y\) near
\(y\). Then \(\rho_Y\circ g+r\) is a smooth local potential of \(\theta_Z\) near \(z\), so
\(w+\rho_Y\circ g+r\) is plurisubharmonic there; composing with the holomorphic
\(\sigma\) preserves plurisubharmonicity and gives \((w+r)\circ s+\rho_Y\)
plurisubharmonic near \(y\), which is exactly
\(\theta_Y\)-plurisubharmonicity of \((w+r)\circ s\) at \(y\).
\end{proof}

\begin{lemma}[Extension of a \(\theta\)-plurisubharmonic function across a codimension-two set]
\label{ana:theta-psh-codim-two-extension}
\uses{ana:nonkahler-locus}
\lean{TranscendentalNakamaye.Pending.ThetaPshCodimTwoExtension.thetaPlurisubharmonic_extensionAcross,TranscendentalNakamaye.Pending.ThetaPshCodimTwoExtension.extensionAcross_eq_of_notMem}
\leanok
Let \(K\) be a closed subset of a complex manifold which, around each of its points, is
carried by a local biholomorphism into a coordinate plane of codimension at least two, and
let \(\theta\) have smooth local potentials. If \(v\) is \(\theta\)-plurisubharmonic on the
complement of \(K\), then its upper-semicontinuous extension
\(\psi(x)=\limsup_{y\to x,\,y\notin K}v(y)\) is \(\theta\)-plurisubharmonic on the whole
manifold and equals \(v\) off \(K\).
\end{lemma}
\begin{proof}
\leanok
The equality off \(K\) is upper semicontinuity of \(v\) there. For the extension, work in a
chart in which \(K\) lies in a coordinate plane of codimension at least two, and add a
smooth local potential of \(\theta\) to reduce to an ordinary plurisubharmonic \(u\) on the
complement of that plane. Upper semicontinuity of the extension is immediate from its
definition as a \(\limsup\). For the sub-mean-value inequality on a complex line \(L\)
through a point of \(K\): the lines through a point that meet the coordinate plane in more
than that point form a set of measure zero in the space of directions, because the plane has
codimension at least two; so the inequality holds on a dense set of lines, and passes to the
limit by upper semicontinuity. The extension is therefore plurisubharmonic, hence \(u\) and
with it \(v\) extends \(\theta\)-plurisubharmonically across \(K\).
\end{proof}

\begin{lemma}[Zero loci of coherent ideals]
\label{ana:zero-locus-coherent}
\uses{ana:ideal-sheaves,fpgeom:complex-analytic-subvariety}
\source{transcendental-nakamaye-copy4-20260911:sec:qpsh}
\lean{TranscendentalNakamaye.Pending.StepFourNoComponentContainsW.mem_zeroLocus_iff_of_generators,TranscendentalNakamaye.Pending.StepFourNoComponentContainsW.poleSet_eq_zeroLocus,TranscendentalNakamaye.Pending.StepFourNoComponentContainsW.isComplexAnalyticSubset_zeroLocus,TranscendentalNakamaye.Pending.StepFourNoComponentContainsW.holomorphicAt_of_mem_of_isCoherent,TranscendentalNakamaye.Pending.StepFourNoComponentContainsW.zeroLocus_pullback}
\leanok
Let \(\mathfrak a\) be an ideal sheaf on a complex manifold \(M\), with zero locus
\(Z(\mathfrak a)\), the set of points at which some germ of \(\mathfrak a\) does not vanish
excluded.
\begin{enumerate}
\item On an open set on which holomorphic functions \(g_1,\dots,g_N\) generate every stalk of
\(\mathfrak a\), \(Z(\mathfrak a)=\{g_1=\dots=g_N=0\}\).
\item If \(\varphi\) has analytic singularities of type \((c,\mathfrak a)\), then
\(\{\varphi=-\infty\}=Z(\mathfrak a)\).
\item If \(\mathfrak a\) is coherent, \(Z(\mathfrak a)\) is a complex analytic subset of \(M\),
and every germ of \(\mathfrak a\) at a point is holomorphic there.
\item If \(\mathfrak a\) is coherent and \(f:Y\to M\) is holomorphic, then
\(Z(\mathfrak a\mathcal O_Y)=f^{-1}Z(\mathfrak a)\).
\end{enumerate}
\end{lemma}
\begin{proof}
\leanok
(i) Each \(g_i\) is a germ of \(\mathfrak a\), so a point of \(Z(\mathfrak a)\) is a common zero
of the \(g_i\); conversely every germ of \(\mathfrak a\) at a common zero is a holomorphic
combination of the \(g_i\), hence vanishes there.

(ii) Near a point, \(\varphi=c\log\sum_i|g_i|^2+R\) with \(c>0\), local generators \(g_i\) of
\(\mathfrak a\) and \(R\) bounded, so \(\varphi=-\infty\) exactly where all \(g_i\) vanish, which
is \(Z(\mathfrak a)\) there by (i).

(iii) By coherence every point has a neighbourhood with finitely many holomorphic generators,
and (i) describes \(Z(\mathfrak a)\) there as their common zero set. A germ of \(\mathfrak a\)
is a holomorphic combination of these generators, hence holomorphic.

(iv) The germs of \(\mathfrak a\mathcal O_Y\) at \(y\) are holomorphic combinations of germs
\(a\circ f\) with \(a\) a germ of \(\mathfrak a\) at \(f(y)\). If \(f(y)\in Z(\mathfrak a)\), all
such \(a\) vanish at \(f(y)\), so every germ of \(\mathfrak a\mathcal O_Y\) vanishes at \(y\).
If \(f(y)\notin Z(\mathfrak a)\), some germ \(a\) of \(\mathfrak a\) has \(a(f(y))\ne0\); it is
holomorphic by (iii), so \(a\circ f\) is a germ of \(\mathfrak a\mathcal O_Y\) not vanishing at
\(y\).
\end{proof}

\begin{lemma}[Smooth metrics on effective divisors]
\label{ana:divisor-metric}
\uses{ana:ideal-sheaves,ana:blanchard-form,ana:zero-locus-coherent}
\source{transcendental-nakamaye-copy4-20260911:prop:principalization,transcendental-nakamaye-copy4-20260911:sec:qpsh}
\lean{TranscendentalNakamaye.Pending.DivisorMetric.IsDivisorMetric,TranscendentalNakamaye.Pending.DivisorMetric.exists_divisorMetric,TranscendentalNakamaye.Pending.DivisorMetric.hasAnalyticSingularitiesOfType_of_isDivisorMetric,TranscendentalNakamaye.Pending.StepFourNoComponentContainsW.support_eq_of_isDivisorMetric}
\leanok
A smooth metric on an effective divisor \(D\) of a complex manifold \(M\) is an effective
\(\mathbf Q\)-divisor \(E\) with a smooth metric whose weight is, near every point,
\(\psi_E=\log|g|^2+h\) with \(g\) a local generator of \(\mathcal O_M(-D)\) and \(h\) smooth. If
\(M\) is compact, every effective divisor carries a smooth metric with the same support. For any
smooth metric \(E\) on \(D\), \(\psi_E\) has analytic singularities of type
\((1,\mathcal O_M(-D))\), and \(\operatorname{Supp}E=\operatorname{Supp}D\).
\end{lemma}
\begin{proof}
\leanok
Cover \(M\) by finitely many open sets \(S_j\) on which \(\mathcal O_M(-D)\) has a generator
\(g_j\), take a smooth partition of unity \((\chi_a)\) subordinate to the cover, with
\(\chi_a\) supported in \(S_{j(a)}\), and put \(U=\sum_a\chi_a|g_{j(a)}|^2\) and
\(\psi_E=\log U\). Two generators of one invertible ideal differ by a nowhere vanishing
holomorphic factor: the two divisibility relations give \(uvg=g\), and \(g\) vanishes
identically on no nonempty open set, so \(uv=1\). Hence near a point of \(S_{j_0}\),
\(U=|g_{j_0}|^2\sum_a\chi_a|u_a|^2\) with \(u_a=g_{j(a)}/g_{j_0}\) holomorphic and nowhere zero,
and \(\psi_E=\log|g_{j_0}|^2+h\) with \(h=\log\sum_a\chi_a|u_a|^2\) smooth. This is the local
form of the weight of a metered divisor (Definition~\ref{ana:ideal-sheaves}); the forms
\(-dd^ch\) of two such local factorizations agree on overlaps and glue to its curvature, as in
Lemma~\ref{ana:blanchard-form}~(ii). Its support is locally \(\{g_{j_0}=0\}\), that of \(D\).

For any smooth metric, \(h\) is bounded near each point (on a compact neighbourhood), so
\(\psi_E=\log|g|^2+O(1)\) with \(g\) generating \(\mathcal O_M(-D)\), which is coherent being
invertible: this is analytic singularities of type \((1,\mathcal O_M(-D))\). Then
\(\operatorname{Supp}E=\{\psi_E=-\infty\}=Z(\mathcal O_M(-D))=\operatorname{Supp}D\) by
Lemma~\ref{ana:zero-locus-coherent}.
\end{proof}

\begin{lemma}[Off its support the weight of a metered divisor is a potential of \(-\Theta_E\)]
\label{ana:divisor-weight-smooth-potential}
\uses{ana:ideal-sheaves,ana:exceptional-metric}
\source{ct22-restricted-volumes:invar}
\lean{TranscendentalNakamaye.Pending.MeteredDivisorSmoothOffSupport.exists_isSmoothLocalPotential_neg_curvature}
\leanok
Let \(E\) be an effective \(\mathbf Q\)-divisor with a smooth metric on a complex manifold
\(Z\), with weight \(\psi_E\) and curvature \(\Theta_E\), so that
\(dd^c\psi_E=[E]-\Theta_E\). At every point \(z\notin\operatorname{Supp}E\) and in every
chart around it, \(\psi_E\) is real valued and smooth on a neighbourhood, and
\(dd^c\psi_E=-\Theta_E\) there.
\end{lemma}
\begin{proof}
\uses{ana:ideal-sheaves,ana:log-norm-pluriharmonic}
Near \(z\) the divisor has a single-generator presentation
\(\psi_E=\tfrac1q\log|G|^2+h\) with \(G\) holomorphic, \(q\in\mathbf N_{>0}\) and \(h\)
smooth, and \(\Theta_E=-dd^ch\); such a presentation is obtained from the defining one
\(\psi_E=\sum_ia_i\log|g_i|^2+h\) by clearing the common denominator of the
\(a_i=p_i/q\) and setting \(G=\prod_ig_i^{p_i}\), which changes no term of the identity
but the generators. Since \(z\notin\operatorname{Supp}E=\{\psi_E=-\infty\}\) we have
\(G(z)\neq0\), so on a neighbourhood \(W\) of \(z\) the function \(\log|G|^2\) is smooth,
whence so is \(\psi_E\); and \(\log|G|^2\) is pluriharmonic there
(Lemma~\ref{ana:log-norm-pluriharmonic}), so
\(dd^c\psi_E=\tfrac1qdd^c\log|G|^2+dd^ch=-\Theta_E\) on \(W\).
\end{proof}

\begin{lemma}[Pushing the exceptionally corrected current across the centre]
\label{ana:pushforward-potential}
\uses{ana:restricted-volume,ana:blowups,ana:exceptional-metric,ana:divisor-weight-smooth-potential}
\source{ct22-restricted-volumes:invar}
\lean{TranscendentalNakamaye.Pending.ExceptionalPushforwardPotential.exists_thetaPshPotential_pushforward}
\leanok
Let \(g:Z\to Y\) be a blowup of a compact complex manifold along \(C\), let \(\omega\) be
a K\"ahler form on \(Y\), let \(E\) be an effective \(\mathbf Q\)-divisor with a smooth
metric supported on \(\operatorname{Exc}(g)\) and put
\(\widetilde\omega=g^*\omega-\Theta_E\). Let \(\theta\) be a closed real
\((1,1)\)-form on \(Y\), let \(\theta_Z\) represent \(g^*\{\theta\}\) and let \(u\) be a
\((\theta_Z+\epsilon_0\widetilde\omega)\)-plurisubharmonic potential on \(Z\). Then there
are a \((\theta+\epsilon_0\omega)\)-plurisubharmonic potential \(\psi\) on \(Y\) and a
smooth real \(h_Z\) with \(g^*\theta=\theta_Z+dd^ch_Z\) such that
\[
 \psi(g(z))=u(z)-h_Z(z)+\epsilon_0\psi_E(z)\qquad\text{for every }
 z\notin\operatorname{Exc}(g).
\]
\end{lemma}
\begin{proof}
\uses{ana:divisor-weight-smooth-potential,ana:theta-psh-pushforward,ana:theta-psh-codim-two-extension,ana:blowup-exceptional-image,ana:blowup-surjective}
Write \(w=u-h_Z\), which is \((g^*\theta+\epsilon_0\widetilde\omega)\)-plurisubharmonic.
The discrepancy between the two backgrounds is
\[
 (g^*\theta+\epsilon_0\widetilde\omega)-g^*(\theta+\epsilon_0\omega)
 =\epsilon_0(\widetilde\omega-g^*\omega)=-\epsilon_0\Theta_E ,
\]
and off \(\operatorname{Supp}E\) the function \(\epsilon_0\psi_E\) is a smooth local
potential of \(-\epsilon_0\Theta_E\) (Lemma~\ref{ana:divisor-weight-smooth-potential}).
Since \(\operatorname{Supp}E\subset\operatorname{Exc}(g)\) and the exceptional locus is the
full preimage of the exceptional centre \(C_2=g(\operatorname{Exc}(g))\), the function
\(w+\epsilon_0\psi_E\) is \(g^*(\theta+\epsilon_0\omega)\)-plurisubharmonic at every point
over \(Y\setminus C_2\). The map \(g\) is surjective (Lemma~\ref{ana:blowup-surjective})
and injective off \(\operatorname{Exc}(g)\) (Lemma~\ref{ana:blowup-exceptional-image}), so it has
a section \(s\) which is the inverse there, and
\(v=(w+\epsilon_0\psi_E)\circ s\) is \((\theta+\epsilon_0\omega)\)-plurisubharmonic on
\(Y\setminus C_2\) (Lemma~\ref{ana:theta-psh-pushforward}). The centre has codimension at
least two flattenings, so \(v\) extends across it
(Lemma~\ref{ana:theta-psh-codim-two-extension}) to a
\((\theta+\epsilon_0\omega)\)-plurisubharmonic \(\psi\) on \(Y\) with \(\psi=v\) off
\(C_2\); the displayed identity is that equality read at \(y=g(z)\), using
\(s(g(z))=z\). Finally \(\psi\) is not identically \(-\infty\) on any open set: an open
\(O\subset Y\) has \(g^{-1}O\not\subset\operatorname{Exc}(g)\) because the exceptional
locus has empty interior (Lemma~\ref{ana:blowup-exceptional-image}), and \(w\) is finite
somewhere on the open set \(g^{-1}O\setminus\operatorname{Exc}(g)\).
\end{proof}

\begin{lemma}[A local biholomorphism carries a smooth germ downwards]
\label{ana:pushforward-smooth-germ}
\uses{ana:blowups}
\source{ct22-restricted-volumes:invar}
\lean{TranscendentalNakamaye.Pending.ExceptionalPushforwardSmooth.eventually_contDiffAt_chartFunction_of_comp}
\leanok
Let \(g:Z\to Y\) be holomorphic and \(z\notin\operatorname{Exc}(g)\), and let \(\psi\) be
an extended-real function on \(Y\). If the chart expression of \(\psi\circ g\) is smooth
near the coordinate of \(z\) in some chart of \(Z\), then the chart expression of \(\psi\)
is smooth near the coordinate of \(g(z)\) in every chart of \(Y\) around \(g(z)\).
\end{lemma}
\begin{proof}
\uses{ana:patch-pushforward,ana:local-holomorphic-inverse}
The differential of \(g\) at \(z\) is bijective, so the chart expression
\(F=\varphi\circ g\circ\phi^{-1}\) has injective differential at \(\phi(z)\) and is a
local biholomorphism there (Lemma~\ref{ana:local-holomorphic-inverse}); let \(\Psi\) be a
local inverse of \(F\) near \(\varphi(g(z))\). For \(y\) near \(\varphi(g(z))\) one has
\(\varphi^{-1}(y)=g(\phi^{-1}(\Psi(y)))\), so the chart expression of \(\psi\) equals the
chart expression of \(\psi\circ g\) precomposed with the holomorphic \(\Psi\), and is
therefore smooth there.
\end{proof}

\begin{lemma}[The parameter density upstairs is the parameter density of the pushforward]
\label{ana:pushforward-density-identity}
\uses{ana:restricted-volume,ana:pushforward-potential,ana:patch-pushforward}
\source{ct22-restricted-volumes:invar}
\lean{TranscendentalNakamaye.Pending.ExceptionalPushforwardDensity.patchDensity_eq_potentialPatchDensity_pushforward}
\leanok
In the situation of Lemma~\ref{ana:pushforward-potential}, let \(T'\) be
\(\epsilon_0\)-admissible along \(\widetilde W\) with representative \(\theta_Z\) and
potential \(u\), and let \(\psi\) be the pushed-down potential. If regular
parametrizations \(\tilde\gamma\) of \(\widetilde W\) at \(v'\) and \(\gamma\) of \(W\) at
\(v\) are related as in hypothesis (ii) of
Lemma~\ref{ana:mass-across-centre-abstract} with \(A=\operatorname{Exc}(g)\), and
\(\tilde\gamma(v')\) lies off \(\operatorname{Exc}(g)\) and off the polar set of \(u\),
then the restricted top-power density of \(T'+\epsilon_0\widetilde\omega\) along
\(\tilde\gamma\) at \(v'\) equals that of \(\theta+\epsilon_0\omega+dd^c\psi\) along
\(\gamma\) at \(v\).
\end{lemma}
\begin{proof}
\uses{ana:pushforward-potential,ana:divisor-weight-smooth-potential,ana:pushforward-smooth-germ,ana:patch-transport,ana:patch-pushforward,ana:smooth-off-poles}
Write \(z=\tilde\gamma(v')\). Off the polar set of \(u\) its chart expression is smooth
(Lemma~\ref{ana:smooth-off-poles}), and off \(\operatorname{Exc}(g)\) the same holds for
\(h_Z\) and for \(\psi_E\) (Lemma~\ref{ana:divisor-weight-smooth-potential}); so near
\(z\) the identity of Lemma~\ref{ana:pushforward-potential} reads
\(u=\psi\circ g+h_Z-\epsilon_0\psi_E\) between smooth real functions, and
\[
 dd^cu=dd^c(\psi\circ g)+dd^ch_Z+\epsilon_0\Theta_E .
\]
Adding \(\theta_Z+\epsilon_0\widetilde\omega=g^*\theta-dd^ch_Z
+\epsilon_0(g^*\omega-\Theta_E)\), the two terms in \(h_Z\) and the two in \(\Theta_E\)
cancel and
\[
 \theta_Z+\epsilon_0\widetilde\omega+dd^cu
 =g^*(\theta+\epsilon_0\omega)+dd^c(\psi\circ g)
\]
as two-forms near \(z\). Restricting to \(\widetilde W\) along \(\tilde\gamma\), the first
term is carried to \((\theta+\epsilon_0\omega)\) along \(\gamma\) by the form identity of
Lemma~\ref{ana:patch-pushforward}, and the second to \(dd^c\psi\) along \(\gamma\) by its
\(dd^c\) identity, which applies because the chart expression of \(\psi\) is smooth near
\(\gamma(v)=g(z)\) (Lemma~\ref{ana:pushforward-smooth-germ}). The two top-power densities
of equal restricted two-forms agree.
\end{proof}

\begin{lemma}[Regularization of a quasi-plurisubharmonic potential along a subvariety]
\label{ana:potential-regularization}
\uses{ana:restricted-volume,ana:nonkahler-locus}
\source{ct22-restricted-volumes:invar}
\lean{TranscendentalNakamaye.Pending.PushforwardPotentialRegularization.exists_regularization_of_thetaPshPotential}
\leanok
Let \((Y,\omega)\) be compact K\"ahler, \(\alpha=\{\theta\}\) big, \(W\) a
positive-dimensional closed irreducible analytic subset, \(\epsilon\ge2\epsilon_0>0\), and
let \(\psi\) be a \((\theta+\epsilon_0\omega)\)-plurisubharmonic potential which is not
identically \(-\infty\) on \(W\). Then there is a sequence \((T_m)_{m\in\mathbf N}\) of
currents \(\epsilon\)-admissible along \(W\) for \(\alpha\) such that
\begin{enumerate}
\item the potential of every \(T_m\) is finite at every point where \(\psi\) is finite; and
\item at every regular parametrization \(\gamma\) of \(W\) and parameter \(v\) at which
\(\psi\) is finite and its chart expression is smooth, the restricted top-power density of
\(\theta+\epsilon_0\omega+dd^c\psi\) along \(\gamma\) at \(v\) is at most
\(\liminf_m\) of the restricted top-power densities of \(T_m+\epsilon\omega\) along
\(\gamma\) at \(v\).
\end{enumerate}
\end{lemma}

\begin{proof}
\uses{ana:potential-regularization-source}
\leanok
The hypotheses \(\epsilon\ge2\epsilon_0\) and \(\epsilon_0>0\) give
\(\epsilon_0<2\epsilon_0\le\epsilon\), which is the only relation between the two
errors used by Lemma~\ref{ana:potential-regularization-source}; that lemma is exactly
the present statement with the strict gap in place of the factor \(2\), and with the
bigness of \(\alpha\) and the positive-dimensionality of \(W\) omitted. Neither of
those two hypotheses enters the construction, which regularizes an almost-positive
current on \(Y\) and never restricts it to \(W\); they are retained here only because
the surrounding CT22 statement carries them.
\end{proof}

\begin{lemma}[The density lower limit reduces to a pointwise Hessian bound]
\label{ana:patch-density-liminf}
\uses{ana:piece-mass-comparison-smooth,ana:restricted-volume}
\lean{TranscendentalNakamaye.Pending.PotentialPatchDensityComparison.potentialCoordinateForm_congr,
TranscendentalNakamaye.Pending.PotentialPatchDensityComparison.potentialPatchDensity_congr,
TranscendentalNakamaye.Pending.PotentialPatchDensityComparison.potentialCoordinateForm_perturbedRepresentative,
TranscendentalNakamaye.Pending.PotentialPatchDensityComparison.ofReal_potentialPatchDensity_le_liminf_of_eventually,
TranscendentalNakamaye.Pending.PotentialPatchDensityComparison.ofReal_potentialPatchDensity_le_liminf_of_eventually_germ}
\leanok
Fix a regular patch \(P\) of \(V\) and a parameter \(v\) of it, and write \(D(a,\psi)\) for
the parameter density of \(a+dd^c\psi\) at \(v\), for a closed real \((1,1)\)-form \(a\) and
an extended-real potential \(\psi\).
\begin{enumerate}
\item \(D(a,\psi)\) depends on \(\psi\) only through the germ at \(v\) of the real chart
expression of \(\psi\).
\item Let \((T_m)_m\) be currents \(\epsilon\)-admissible along \(V\). If, for all large
\(m\), the coordinate \((1,1)\)-form of \(T_m+\epsilon\omega\) dominates that of
\(a+dd^c\psi\) on complex lines at the chart image of \(v\), then
\(D(a,\psi)\le\liminf_mD(T_m+\epsilon\omega)\).
\item In particular it is enough that, for all large \(m\), the germ of the potential of
\(T_m\) at \(v\) is that of \(\psi\) and the smooth representative of \(T_m+\epsilon\omega\)
dominates \(a\) on complex lines at that point: the lower limit is then attained by an
eventually constant sequence.
\end{enumerate}
\end{lemma}
\begin{proof}
\uses{ana:piece-mass-comparison-smooth}
\leanok
Clause 1 is the corresponding property of \(dd^c\) of a real function: the chartwise
\(dd^c\) at a point is determined by the germ of the function there. Clause 2 is the
determinant monotonicity \(\det A\le\det(A+B)\) for positive semidefinite \((1,1)\)-forms
\(A,B\), transported through the complex-linear differential of the parametrization
(Lemma~\ref{ana:piece-mass-comparison-smooth}), followed by the fact that a quantity
eventually below a sequence is below its lower limit. Clause 3 substitutes the germ equality
of clause 1 into clause 2, the two coordinate forms then differing only by the semipositive
difference of the smooth representatives.
\end{proof}

\begin{lemma}[A chart-smooth function is quasi-plurisubharmonic]
\label{ana:smooth-quasi-psh}
\uses{ana:ddc-global-form}
\lean{TranscendentalNakamaye.Pending.SmoothDDCForm.exists_pos_thetaPlurisubharmonic_smooth}
\leanok
Let \((X,\omega)\) be a compact K\"ahler manifold, let \(\theta_i\), \(i\) in a finite set,
be closed real \((1,1)\)-forms and let \(v_i\) be real functions on \(X\) that are smooth in
every holomorphic chart. Then there is \(C>0\) such that every \(\theta_i+C\omega\) is
semipositive and every \(v_i\) is \((\theta_i+C\omega)\)-plurisubharmonic. In particular
\(dd^cv_i\ge-C\omega\) for every \(i\).
\end{lemma}
\begin{proof}
\uses{ana:ddc-global-form}
\leanok
By Lemma~\ref{ana:ddc-global-form} each \(dd^cv_i\) is a global smooth real two-form. A
finite family of continuous Hermitian coefficient fields on a compact manifold is dominated
by a fixed multiple of a positive definite one, applied first to the backgrounds
\(\theta_i\) and then to the forms \(dd^cv_i\); take for \(C\) the larger of the two
constants.
\end{proof}

\begin{lemma}[A collar cut-off with a lower \(dd^c\) bound]
\label{ana:collar-cutoff}
\uses{ana:smooth-quasi-psh}
\lean{TranscendentalNakamaye.Pending.CompactCollarCutoff.exists_chartSmooth_collar_ddc_lower}
\leanok
Let \(X\) be a compact complex manifold with a K\"ahler form \(\omega\), and let
\(K\subset W\subset X\) with \(K\) compact and \(W\) open. Then there is a real function
\(\rho\) on \(X\), smooth in every holomorphic chart, with \(-1\le\rho\le0\), \(\rho=0\) on
a neighbourhood of \(K\) and \(\rho=-1\) off \(W\); and, by
Lemma~\ref{ana:smooth-quasi-psh}, a constant \(C>0\) with \(dd^c\rho\ge-C\omega\).
\end{lemma}
\begin{proof}
\uses{ana:smooth-quasi-psh}
\leanok
For each \(x\in K\) choose, by the chart bump construction, a chart-smooth \(h_x\) with
\(0\le h_x\le1\), compact support inside \(W\) and \(h_x=1\) near \(x\). Finitely many of
the open sets \(\operatorname{int}\{h_x=1\}\) cover \(K\); say those of
\(h_{x_1},\dots,h_{x_r}\). Put
\[
 \rho=-\prod_{i=1}^r(1-h_{x_i}).
\]
Each factor takes values in \([0,1]\), so \(-1\le\rho\le0\); the product vanishes wherever
some \(h_{x_i}\) equals \(1\), hence on the open set \(\bigcup_i\operatorname{int}\{h_{x_i}=1\}\supset K\);
and it equals \(1\) wherever every \(h_{x_i}\) vanishes, hence off \(W\). A finite product
of chart-smooth functions is chart-smooth, so \(\rho\) is. The \(dd^c\) bound is
Lemma~\ref{ana:smooth-quasi-psh} applied to \(\rho\).

Only the LOWER bound is asserted, because only it is used: in
Lemma~\ref{ana:potential-regularization-source} the local branch is \(\psi+s\rho\), whose
background is \(\theta-s\,dd^c\rho\), so what the positivity budget must pay for is
semipositivity of \(s\,dd^c\rho+(\Lambda/2)\omega\).
\end{proof}

\begin{lemma}[Exhaustion of an open set by compacta with collars]
\label{ana:collar-exhaustion}
\lean{IsOpen.exists_nested_collar_exhaustion}
\leanok
Let \(X\) be a compact Hausdorff second-countable space and \(U\subset X\) open. Then there
are compacta \(K_m\) and open sets \(W_m,V_m\) with \(K_m\subset K_{m+1}\),
\(\bigcup_mK_m=U\), \(K_m\subset W_m\), \(\overline{W_m}\subset V_m\) and
\(\overline{V_m}\subset U\) with \(\overline{V_m}\) compact.
\end{lemma}
\begin{proof}
\leanok
\(U\) is locally compact (an open subset of a compact Hausdorff space) and second countable,
hence \(\sigma\)-compact, so it carries a compact exhaustion; push it forward along the
inclusion to get the increasing \(K_m\subset U\) with union \(U\). Given \(K_m\), local
compactness and regularity give an open \(V_m\supset K_m\) with \(\overline{V_m}\subset U\)
compact, and applying the same to \(K_m\subset V_m\) gives \(W_m\).
\end{proof}

\begin{lemma}[A theta-PSH potential is finite where \(\operatorname{toReal}\) is continuous]
\label{ana:toreal-finiteness}
\lean{TranscendentalNakamaye.Pending.ThetaPshContinuityLocus.ne_bot_of_mem_continuityLocus}
\leanok
Let \(u\) be a \(\theta\)-plurisubharmonic potential on \(X\) and let \(x\) be a point near
which the real function \(\operatorname{toReal}u\) --- equal to \(u\) where \(u\) is finite
and to \(0\) at the poles --- is continuous. Then \(u(x)>-\infty\).
\end{lemma}
\begin{proof}
\leanok
Suppose \(u(x)=-\infty\), so \(\operatorname{toReal}u(x)=0\). Upper semicontinuity of \(u\)
gives a neighbourhood \(V\) of \(x\) on which \(u<-1\), and continuity of
\(\operatorname{toReal}u\) gives one on which \(|\operatorname{toReal}u|<1/2\); shrink to
their intersection. At a point of \(V\) where \(u\) is finite one has
\(\operatorname{toReal}u=u<-1\), contradicting \(|\operatorname{toReal}u|<1/2\). So
\(u\equiv-\infty\) on \(V\), an open set --- which the local nondegeneracy clause of
``\(\theta\)-PSH potential'' forbids.

This is the step that makes the hypothesis of clause \emph{(ii)} of
Lemma~\ref{ana:potential-regularization} usable. That hypothesis asserts smoothness of the
chart expression of \(\operatorname{toReal}\psi\) near one point together with
\(\psi>-\infty\) at that point, and \(\operatorname{toReal}\) is blind to poles; nothing in it
gives finiteness on a \emph{neighbourhood}, which the collar construction needs, because
\(\psi_k\ge\psi\) does not force \(\psi_k=-\infty\) where \(\psi=-\infty\), so the glued
potential cannot equal \(\psi\) at a pole.
\end{proof}

\begin{lemma}[The smooth locus is open, and chart-independent]
\label{ana:smooth-locus-open}
\lean{TranscendentalNakamaye.Pending.ThetaPshSmoothLocus.isOpen_smoothLocus}
\leanok
For \(u:X\to[-\infty,\infty)\) let \(U\) be the set of points near which \(u\) agrees with a
real function that is smooth read in \emph{every} holomorphic chart. Then \(U\) is open, and
\(u\) is finite on \(U\).
\end{lemma}
\begin{proof}
\leanok
Finiteness is immediate: on \(U\) the function \(u\) equals a real function. For openness it
suffices that agreement with a smooth real germ in ONE chart implies it in every chart at
every nearby point. If \(u=g\circ\varphi_e\) near \(x\) with \(g\) smooth on an open
\(V\subset\varphi_e(X)\), and \(f\) is another chart around a point \(y\) of
\(\varphi_e^{-1}(V)\), then near \(\varphi_f(y)\) one has
\(u\circ\varphi_f^{-1}=g\circ(\varphi_e\circ\varphi_f^{-1})\), and the chart transition is
\(C^\infty\); so \(g\circ(\varphi_e\circ\varphi_f^{-1})\) is the required smooth witness in
the chart \(f\).
\end{proof}

\begin{lemma}[The density clause is asserted only on the smooth locus]
\label{ana:smooth-locus-clause}
\uses{ana:toreal-finiteness,ana:smooth-locus-open}
\lean{TranscendentalNakamaye.Pending.ThetaPshSmoothLocus.mem_smoothLocus_of_eventually_contDiffAt}
\leanok
Let \(\psi\) be a \(\theta\)-plurisubharmonic potential, \(e\) a holomorphic chart and \(z\) a
point of its target such that the chart expression of \(\operatorname{toReal}\psi\) is
\(C^\infty\) near \(z\). Then \(\varphi_e^{-1}(z)\) lies in the smooth locus of \(\psi\).
\end{lemma}
\begin{proof}
\uses{ana:toreal-finiteness,ana:smooth-locus-open}
\leanok
Smoothness implies continuity, so by Lemma~\ref{ana:toreal-finiteness} the function \(\psi\)
is finite at every point of an open neighbourhood of \(\varphi_e^{-1}(z)\) --- the continuity
locus is open and \(\psi\) is finite on all of it. Where \(\psi\) is finite it equals the
real number \(\operatorname{toReal}\psi\), so near \(z\) the chart expression of \(\psi\) IS
the smooth function of the hypothesis, and Lemma~\ref{ana:smooth-locus-open} transports that
witness to every chart.
\end{proof}

\begin{lemma}[The collar glue keeps analytic singularities, with no alignment condition]
\label{ana:offpole-collar-glue}
\uses{rmax:definition,rmax:gap,rmax:max-append,ana:analytic-type-pole-set}
\lean{TranscendentalNakamaye.Pending.RegularizedMaxOffPoleCollarGluing.hasManifoldAnalyticSingularitiesSmoothOffPole_eregularizedMaxGlued}
\leanok
Let \(V'\Subset V\subset X\) be open, \(\delta>0\), and let \(f,g\) satisfy \(f+\delta\le g\)
on \(V\setminus V'\), with \(g<\infty\) everywhere. Suppose \(g\) has analytic singularities
and \(f\) is a smooth real germ at every point of \(V\). Then the glue
\(M_\delta(g,f)\) on \(V\), \(g\) off \(V\), has analytic singularities.
\end{lemma}
\begin{proof}
\uses{rmax:definition,rmax:gap,rmax:max-append,ana:analytic-type-pole-set}
\leanok
Off \(V\) the collar makes the glue equal to \(g\) as a germ
(Lemma~\ref{rmax:gap}), so the conclusion is inherited. Inside \(V\) the glue is
\(M_\delta(g,f)\), and Lemma~\ref{rmax:max-append} concatenates the two generator families
provided both branches carry the SAME positive coefficient. That is free here: a smooth real
germ is presented by the one-element constant family \((1)\), whose logarithm vanishes
identically, so the coefficient never appears and the local branch simply adopts the
coefficient of \(g\). No alignment of the two presentations at a common pole is needed,
because the local branch has no poles. The smoothness of the local branch is essential and is
what Lemma~\ref{ana:smooth-locus-clause} supplies: a merely continuous branch does not have
analytic singularities.
\end{proof}

\begin{theorem}[Demailly: approximation with analytic singularities]
\label{ana:demailly-algebraic-approximation}
\uses{ana:analytic-type-pole-set,ana:geometric-setting,ana:ideal-sheaves,as:smooth-off-pole-type}
\source{bou02-volume:2.6}
\lean{TranscendentalNakamaye.Pending.Demailly1992AlgebraicApproximation.demailly1992_exists_algebraicApproximation,TranscendentalNakamaye.Pending.Demailly1992AlgebraicApproximation.exists_algebraicApproximation_smoothOffPole,TranscendentalNakamaye.Pending.DemaillyAlgebraicApproximation.exists_algebraicApproximation,TranscendentalNakamaye.Pending.KahlerCurrentRegularizationAlgebraicApproximation.demailly1992_exists_algebraicApproximation_analyticSingularities,TranscendentalNakamaye.Pending.DemaillyApproximationSmoothOffPole.exists_algebraicApproximation,TranscendentalNakamaye.Pending.DemaillyApproximationSmoothOffPole.hasManifoldAnalyticSingularitiesSmoothOffPole_of_dim_zero}
\leanok
Let \((M,\omega)\) be a compact K\"ahler manifold, \(\alpha\) a smooth closed real
\((1,1)\)-form, and \(\psi\) an \(\alpha\)-plurisubharmonic potential (so
\(\psi\in\operatorname{PSH}(M,\alpha)\) is \(-\infty\) on no open set). There are functions
\(\psi_k\colon M\to[-\infty,\infty)\) and real numbers \(\delta_k\), \(k\in\mathbf N\), such that
\begin{enumerate}
\item each \(\psi_k\) has analytic singularities of type \((c_k,\mathfrak a_k)\) with
\(c_k\in\mathbf Q_{>0}\) and \(\mathfrak a_k\) coherent (Definition~\ref{ana:ideal-sheaves}), and in
every chart a presentation \(c\log\sum_i|f_i|^2+b\) with \(b\) bounded and \(C^\infty\) off the common
zeros of the \(f_i\) (Definition~\ref{as:smooth-off-pole-type});
\item \(k\mapsto\psi_k(x)\) is nonincreasing and \(\inf_k\psi_k(x)=\psi(x)\) at every \(x\in M\);
\item[(iv)] \(\delta_k\ge0\), \(\delta_k\to0\), and \(\psi_k\) is an
\((\alpha+\delta_k\omega)\)-plurisubharmonic potential.
\end{enumerate}
\end{theorem}
\begin{proof}
\uses{dem:algebraic-approximation-global,dem:decreasing-approximation,dem:smooth-off-poles}
\leanok
If \(M\) is empty there is nothing to prove. If \(\dim M=0\), the points of \(M\) are open,
\(\psi\) is finite at every point and locally constant, and \(\psi_k=\psi\), \(\delta_k=0\) work:
near \(x\), \(\psi=1\cdot\log|1|^2+\psi(x)\) with the unit ideal, a constant (hence bounded and
smooth) remainder, and no poles.

Let \(\dim M\ge1\). Lemma~\ref{dem:decreasing-approximation} gives the sequence
\(\psi_k=\varphi_{2^k}+E2^{-k}\) of glued chart Bergman weights and \(\delta_k=(C_B+1)C_\chi2^{-k}\) with
clauses \emph{(ii)} and \emph{(iv)}. Clause \emph{(i)} for the same \(\psi_k\), with \(c_k=2^{-k}\) and
\(\mathfrak a_k=\mathcal I(2^k\psi)\), is Lemma~\ref{dem:analytic-singularities} (the type) together with
Lemma~\ref{dem:smooth-off-poles} (the remainder smooth off the poles). Everything except the smoothness of the
remainder is Theorem~\ref{dem:algebraic-approximation-global}.

This is Demailly's proof: local Bergman kernels of the weights \(e^{-m\Psi}\), their
subadditivity \(K_{2m}\le CK_m^2\) (Lemma~\ref{berg:subadditivity}) for the monotonicity along
\(m=2^k\), the Ohsawa--Takegoshi lower bound and the sub-mean value upper bound for the convergence
(Lemma~\ref{berg:weight-convergence}), and Nadel's theorem for the singularities
(Lemma~\ref{dem:nadel}). The gluing differs from Demailly's: the local weights are glued with the
symmetric regularized maximum over a cover and with a constant \(C_B\) fixed before \(\alpha\)
(Copy-4's Step~3; Lemma~\ref{ana:regularization-gluing-rmax}), and clause \emph{(iii)} is not
proved.
\end{proof}
This is Demailly's regularization theorem (\emph{Regularization of closed positive currents and
intersection theory}, J.\ Algebraic Geom.\ \textbf{1} (1992)), in the form in which it enters the
development. Boucksom quotes it for a compact complex manifold \(X\) with a Hermitian form
\(\omega\) and a closed almost positive \((1,1)\)-current \(T=\alpha+dd^c\varphi\ge\gamma\), \(\gamma\)
a continuous real form: there are almost plurisubharmonic \(\varphi_k\) with \emph{(i)} algebraic
singularities, \emph{(ii)} decreasing to \(\varphi\) at every point, \emph{(iii)}
\(\nu(\varphi_k,x)\to\nu(\varphi,x)\) uniformly in \(x\), and \emph{(iv)}
\(\alpha+dd^c\varphi_k\ge\gamma-\epsilon_k\omega\) with \(\epsilon_k\downarrow0\). The statement
above is the case of a K\"ahler \(\omega\) and \(\gamma=0\), with clause \emph{(iii)} omitted, clause
\emph{(i)} asked in the two forms used below (both implied by algebraic singularities), clause
\emph{(ii)} as a nonincreasing sequence with infimum \(\psi\), and clause \emph{(iv)} as
\(0\le\delta_k\to0\). Each of these changes weakens the statement; the clauses keep the source's
numbering. The forms used by the collar construction of
Lemma~\ref{ana:potential-regularization-source} (chartwise clause \emph{(i)} only) and by
Lemma~\ref{ana:kahler-current-regularization} are obtained from it by dropping clauses.

\begin{lemma}[Demailly regularization of an almost-positive potential with a strict error gap]
\label{ana:potential-regularization-source}
\uses{ana:restricted-volume,ana:nonkahler-locus,ana:analytic-type-pole-set}
\source{ct22-restricted-volumes:invar,bou02-volume:2.6,bou02-volume:2.7}
\lean{TranscendentalNakamaye.Pending.PushforwardPotentialRegularization.ct22Lemma43_exists_regularization_of_thetaPshPotential}
\leanok
Let \((Y,\omega)\) be compact K\"ahler, \(\theta\) a smooth closed real
\((1,1)\)-form, \(W\) a closed irreducible analytic subset, \(0<\epsilon_0<\epsilon\),
and let \(\psi\) be a \((\theta+\epsilon_0\omega)\)-plurisubharmonic potential which is
not identically \(-\infty\) on \(W\). Then there is a sequence \((T_m)_{m\in\mathbf N}\)
of currents \(\epsilon\)-admissible along \(W\) for \(\{\theta\}\) satisfying clauses
\emph{(i)} and \emph{(ii)} of Lemma~\ref{ana:potential-regularization}.
\end{lemma}
\begin{proof}
\uses{ana:demailly-algebraic-approximation,ana:smooth-quasi-psh,ana:collar-cutoff,ana:collar-exhaustion,ana:smooth-locus-open,ana:smooth-locus-clause,ana:offpole-collar-glue,ana:patch-density-liminf,rmax:definition,rmax:gap,rmax:max-append,bts:dini,ana:analytic-type-pole-set}
\leanok
Write \(T=\theta+\epsilon_0\omega+dd^c\psi\); it is closed and positive because \(\psi\) is
\((\theta+\epsilon_0\omega)\)-plurisubharmonic. Apply
Theorem~\ref{ana:demailly-algebraic-approximation} with \(\alpha=\theta+\epsilon_0\omega\)
and the potential \(\psi\): there are almost plurisubharmonic \(\psi_k\) with
analytic singularities, decreasing to \(\psi\) at every point, and \(\delta_k\ge0\) with
\(\delta_k\to0\) and
\[
 \theta+\epsilon_0\omega+dd^c\psi_k\ge-\delta_k\omega .
\]
Two consequences are recorded once. Since \(\psi_k\ge\psi\), the pole set of \(\psi_k\) is
contained in that of \(\psi\). And a function with analytic singularities is continuous as a
map to \([-\infty,\infty)\), so each \(\psi_k\) is real valued and continuous on every set
where \(\psi\) is finite. Write \(\Lambda=\epsilon-\epsilon_0>0\).

\emph{The open set on which the density clause is asserted.} Let \(U\) be the
\emph{smooth locus} of \(\psi\): the set of points near which \(\psi\) agrees with a real
function that is smooth read in every holomorphic chart
(Lemma~\ref{ana:smooth-locus-open}). It is open by construction, \(\psi\) is finite on it,
and the local branch below will be smooth there --- which is what the gluing
Lemma~\ref{ana:offpole-collar-glue} requires, a merely continuous branch not having
analytic singularities. By Lemma~\ref{ana:smooth-locus-clause} the parameters over which
clause \emph{(ii)} quantifies all have their image in \(U\); this is the one step where the
hypothesis of clause \emph{(ii)} has to be improved before it can be used, since that
hypothesis controls only \(\operatorname{toReal}\psi\), which is blind to poles.

Fix compact \(K_m\) and open \(W_m,V_m\) with
\[
 K_m\subset W_m,\qquad \overline{W_m}\subset V_m,\qquad \overline{V_m}\subset U
 \text{ compact},\qquad K_m\subset K_{m+1},\qquad \bigcup_mK_m=U ,
\]
which exist because \(U\) is an open subset of a compact Hausdorff second-countable space,
hence locally compact and \(\sigma\)-compact (Lemma~\ref{ana:collar-exhaustion}). Note that
\(K_m\subset\operatorname{int}K_{m+1}\) is \emph{not} needed: the cut-off of
Lemma~\ref{ana:collar-cutoff} already vanishes on a neighbourhood of \(K_m\), so a monotone
family with union \(U\) suffices for every point of \(U\) to be in the core of all late
stages.

\emph{The two branches.} Fix \(m\). By Lemma~\ref{ana:collar-cutoff} choose a chart-smooth
\(\rho_m\) with \(-1\le\rho_m\le0\), \(\rho_m=0\) on a neighbourhood \(N_m\) of \(K_m\),
\(\rho_m=-1\) off \(W_m\), and \(C_m>0\) with \(dd^c\rho_m\ge-C_m\omega\). Choose
\(0<s_m\le\Lambda/(2C_m)\) and then \(k_m\) with
\[
 \delta_{k_m}\le\Lambda/2\qquad\text{and}\qquad\sup_{\overline{V_m}}(\psi_{k_m}-\psi)<s_m/4 .
\]
Both are possible: \(\delta_k\to0\), and on the compact set \(\overline{V_m}\subset U\)
the functions \(\psi_k\) are continuous and decrease to the continuous \(\psi\), so
Lemma~\ref{bts:dini} makes that supremum tend to \(0\). Put \(\eta_m=s_m/2\),
\(\delta=s_m/4\), and
\[
 w_m=\begin{cases}
 M_{\delta}\bigl(\psi+s_m\rho_m,\ \psi_{k_m}-\eta_m\bigr)&\text{on }V_m,\\
 \psi_{k_m}-\eta_m&\text{off }V_m,
 \end{cases}
\]
with \(M_\delta\) the regularized maximum of Definition~\ref{rmax:definition}. The
\emph{local} branch is \(\psi+s_m\rho_m\), defined on \(V_m\subset U\), where \(\psi\) is a
smooth real function; the
\emph{global} branch is the constant downward shift \(\psi_{k_m}-\eta_m\). The shift is what
makes the construction work: without it \(\psi_{k_m}\ge\psi\) everywhere, so no maximum
against \(\psi\) could ever select \(\psi\).

\emph{The global branch wins on the collar, so \(w_m\) is well defined and \(\epsilon\)-admissible.}
On \(V_m\setminus W_m\) one has \(\rho_m=-1\), hence
\[
 (\psi+s_m\rho_m)+\delta=\psi-s_m+\tfrac{s_m}4\le\psi_{k_m}-\tfrac{3s_m}4
 \le\psi_{k_m}-\eta_m ,
\]
so the gap hypothesis of Lemma~\ref{rmax:gap} holds on the whole of \(V_m\setminus W_m\) and
\(w_m=\psi_{k_m}-\eta_m\) there; in particular the two cases of the definition agree near
\(\partial V_m\) and \(w_m\) is a globally defined potential. The local branch is
\((\theta+\epsilon_0\omega-s_m\,dd^c\rho_m)\)-plurisubharmonic --- one subtracts the
\(dd^c\) that was added to the potential --- hence, by \(dd^c\rho_m\ge-C_m\omega\),
\((\theta+\epsilon_0\omega+s_mC_m\omega)\)-plurisubharmonic; and the global branch is
\((\theta+\epsilon_0\omega+\delta_{k_m}\omega)\)-plurisubharmonic. Both constants are at most
\(\Lambda/2\), so \(w_m\) is \((\theta+\epsilon_0\omega+(\Lambda/2)\omega)\)-plurisubharmonic
and
\[
 \theta+\epsilon\omega+dd^cw_m
 =\bigl(\theta+\epsilon_0\omega+\tfrac\Lambda2\omega+dd^cw_m\bigr)+\tfrac\Lambda2\omega\ge0 .
\]
Only the lower bound \(dd^c\rho_m\ge-C_m\omega\) is used, which is why
Lemma~\ref{ana:collar-cutoff} asserts no other.
Off \(V_m\) the potential \(w_m\) is \(\psi_{k_m}-\eta_m\), which has algebraic
singularities; on \(V_m\) the local branch is smooth, so by
Lemma~\ref{rmax:max-append} the regularized maximum again has analytic singularities, with
pole set contained in that of \(\psi_{k_m}\), hence in that of \(\psi\). Therefore
\(T_m=\theta+dd^cw_m\) is \(\epsilon\)-admissible along \(W\) for \(\{\theta\}\), and clause
\emph{(i)} holds for every \(m\) at once. That the gap is used exactly here --- once to absorb
\(\delta_{k_m}\) and once to pay for the collar --- is the whole role of
\(\epsilon_0<\epsilon\).

\emph{The local branch wins near \(K_m\), so \(w_m=\psi\) there.} On \(N_m\) one has
\(\rho_m=0\) and \(\psi_{k_m}-\psi<s_m/4\), hence
\[
 (\psi_{k_m}-\eta_m)+\delta<\psi+\tfrac{s_m}4-\tfrac{s_m}2+\tfrac{s_m}4=\psi
 =\psi+s_m\rho_m ,
\]
so Lemma~\ref{rmax:gap} in its other direction gives \(w_m=\psi\) on \(N_m\). The transition
region, where \(-1<\rho_m<0\) and neither gap holds, lies inside \(W_m\setminus N_m\), and no
information about \(w_m\) is needed there.

\emph{Clause \emph{(ii)}.} Let \(\gamma\) be a regular parametrization of \(W\) and \(v\) a
parameter with \(\gamma(v)\in U\). Since the \(K_m\) are increasing with union \(U\), there is
an \(m_0\) with \(\gamma(v)\in K_m\) for every \(m\ge m_0\), and \(K_m\subset N_m\), so for
those \(m\) the germ of \(w_m\) at \(\gamma(v)\) is the germ of \(\psi\). The restricted top-power density at a parameter
depends only on that germ and on the smooth representative, and it is monotone in a
semipositive summand, so for all large \(m\) the density of \(T_m+\epsilon\omega\) along
\(\gamma\) at \(v\) equals that of
\(\theta+\epsilon\omega+dd^c\psi=(\theta+\epsilon_0\omega+dd^c\psi)+\Lambda\omega\),
which is at least the density of \(\theta+\epsilon_0\omega+dd^c\psi\). A sequence eventually
bounded below by a number has lower limit at least that number; this is exactly the reduction
of Lemma~\ref{ana:patch-density-liminf}, and it gives clause \emph{(ii)} with an eventually
constant sequence, so no convergence theorem for absolutely continuous parts is needed.

Neither the bigness of \(\{\theta\}\) nor \(W\not\subset\operatorname{poleSet}\psi\) enters;
they are carried because the surrounding CT22 statement carries them.

The source's own route to clause \emph{(ii)} is different and weaker: the refinement of
Demailly's approximation which also controls the absolutely continuous parts
(\(\source{bou02-volume:2.7}\): the approximants may be chosen with
\((T_k)_{\mathrm{ac}}\to T_{\mathrm{ac}}\) almost everywhere) gives the densities almost
everywhere, which with Fatou's lemma is what Collins--Tosatti use at a generic point of
\(W\). The construction above reaches the pointwise form the Lean statement asserts, so that
refinement is not needed.
\end{proof}

\begin{lemma}[Regularization of the pushforward across the centre]
\label{ana:pushforward-regularization}
\uses{ana:restricted-volume,ana:blowups,ana:exceptional-metric,ana:potential-regularization,ana:pushforward-potential}
\source{ct22-restricted-volumes:invar}
\lean{TranscendentalNakamaye.Pending.ExceptionalPushforwardFatouReduction.exists_pushforwardRegularization}
\leanok
In the situation of Lemma~\ref{ana:rv-le-pullback}, let \(E\) be an effective
\(\mathbf Q\)-divisor on \(Z\) with a smooth metric, supported on
\(\operatorname{Exc}(g)\), such that \(\widetilde\omega=g^*\omega-\Theta_E\) is K\"ahler,
let \(\epsilon\ge2\epsilon_0>0\) and let \(T'\) be \(\epsilon_0\)-admissible along
\(\widetilde W\) for \(g^*\alpha\) with respect to \(\widetilde\omega\). Then there is a
sequence \((T_m)_{m\in\mathbf N}\) of currents \(\epsilon\)-admissible along \(W\) for
\(\alpha\) with respect to \(\omega\) such that
\begin{enumerate}
\item for every \(m\), the potential of \(T_m\) is finite at \(g(z)\) whenever
\(z\notin\operatorname{Exc}(g)\) and the potential of \(T'\) is finite at \(z\); and
\item whenever regular parametrizations \(\tilde\gamma\) of \(\widetilde W\) at \(v'\)
and \(\gamma\) of \(W\) at \(v\) are related as in hypothesis (ii) of
Lemma~\ref{ana:mass-across-centre-abstract} with \(A=\operatorname{Exc}(g)\), the
restricted top-power density of \(T'+\epsilon_0\widetilde\omega\) along \(\tilde\gamma\)
at \(v'\) is at most \(\liminf_m\) of the restricted top-power densities of
\(T_m+\epsilon\omega\) along \(\gamma\) at \(v\).
\end{enumerate}
\end{lemma}
\begin{proof}
\uses{ana:pushforward-potential,ana:pushforward-density-identity,ana:pushforward-smooth-germ,ana:potential-regularization,ana:smooth-pole-set-analytic,ana:blowup-exceptional-image,ana:strict-transform-loci,ana:mass-across-centre-abstract,fpgeom:complex-analytic-subvariety}
\leanok
\emph{Pushing forward.} From \(T'+\epsilon_0\widetilde\omega\ge0\) and
\(\widetilde\omega=g^*\omega-\Theta_E\) we get
\(S=T'-\epsilon_0\Theta_E\ge-\epsilon_0g^*\omega\); the divisor \(E\) is exceptional, so
the class of \(\Theta_E\) pushes forward to zero and \(T=g_*S\) lies in \(\alpha\) with
\(T\ge-\epsilon_0\omega\). Lemma~\ref{ana:pushforward-potential} realises this in
potentials: there is a \((\theta+\epsilon_0\omega)\)-plurisubharmonic \(\psi\) on \(Y\) with
\(u=\psi\circ g+h_Z-\epsilon_0\psi_E\) off \(\operatorname{Exc}(g)\), where \(u\) is the
potential of \(T'\) and \(g^*\theta=\theta_Z+dd^ch_Z\); so \(T=\theta+dd^c\psi\).

\emph{\(\psi\) is finite somewhere on \(W\).} The exceptional locus is analytic and so is
the polar set \(P(u)\) (Lemma~\ref{ana:smooth-pole-set-analytic}). The strict transform
\(\widetilde W\) is contained in neither: not in \(P(u)\) because \(T'\) is admissible
along \(\widetilde W\), and not in \(\operatorname{Exc}(g)\) because
\(W\not\subset g(\operatorname{Exc}(g))\) (Lemma~\ref{ana:strict-transform-loci}). Since
\(\widetilde W\) is irreducible it is not the union of the two proper analytic subsets
\(\widetilde W\cap P(u)\) and \(\widetilde W\cap\operatorname{Exc}(g)\), so some
\(z\in\widetilde W\) lies outside both. Then \(g(z)\in W\), and \(\psi(g(z))\) is finite
because \(u(z)\), \(h_Z(z)\) and \(\psi_E(z)\) all are.

\emph{Regularizing.} Apply Lemma~\ref{ana:potential-regularization} to \(\psi\): it
produces currents \(T_m\) that are \(\epsilon\)-admissible along \(W\) for \(\alpha\), whose
potentials are finite wherever \(\psi\) is, and whose densities \(\liminf\)-dominate the
density of \(\theta+\epsilon_0\omega+dd^c\psi\) at every parameter where \(\psi\) is finite
with smooth chart expression.

\emph{Clause (i).} For \(z\notin\operatorname{Exc}(g)\) at which \(u\) is finite, the
displayed identity makes \(\psi(g(z))\) finite, since \(\psi_E\) is finite off
\(\operatorname{Supp}E\subset\operatorname{Exc}(g)\); so no \(T_m\) has a pole at \(g(z)\).

\emph{Clause (ii).} Let \(\tilde\gamma\) at \(v'\) and \(\gamma\) at \(v\) be as stated.
By Lemma~\ref{ana:pushforward-density-identity} the density of
\(T'+\epsilon_0\widetilde\omega\) along \(\tilde\gamma\) at \(v'\) EQUALS the density of
\(\theta+\epsilon_0\omega+dd^c\psi\) along \(\gamma\) at \(v\): the two exceptional
corrections \(-\epsilon_0\Theta_E\) in \(\widetilde\omega\) and \(+\epsilon_0\Theta_E\) in
\(dd^c(\epsilon_0\psi_E)\) cancel exactly, so nothing is lost in \(\epsilon_0\). The chart
expression of \(\psi\) is smooth at \(\gamma(v)=g(\tilde\gamma(v'))\)
(Lemma~\ref{ana:pushforward-smooth-germ}), so clause (ii) of
Lemma~\ref{ana:potential-regularization} bounds that density by
\(\liminf_m\) of the densities of \(T_m+\epsilon\omega\).
\end{proof}

\begin{lemma}[Pushforward of one admissible current across the centre]
\label{ana:rv-pushforward-one-current}
\uses{ana:restricted-volume,ana:blowups,ana:strict-transform-blowup,ana:exceptional-metric}
\source{ct22-restricted-volumes:invar}
\lean{TranscendentalNakamaye.Pending.ExceptionalPushforwardMass.acTopMass_le_restrictedMassEnvelope_of_exceptionalCorrection,TranscendentalNakamaye.Pending.ExceptionalPushforwardFatouReduction.acTopMass_le_restrictedMassEnvelope_of_exceptionalCorrection}
\leanok
In the situation of Lemma~\ref{ana:rv-le-pullback}, let \(E\) be an effective
\(\mathbf Q\)-divisor on \(Z\) with a smooth metric, supported on
\(\operatorname{Exc}(g)\), such that \(\widetilde\omega=g^*\omega-\Theta_E\) is a
K\"ahler form on \(Z\). Let \(\epsilon_0>0\) and let \(T'\) be
\(\epsilon_0\)-admissible along \(\widetilde W\) for \(g^*\alpha\) with respect to
\(\widetilde\omega\). Then for every \(\epsilon\ge2\epsilon_0\),
\[
 \int_{\widetilde W_{\mathrm{reg}}}\bigl((T'+\epsilon_0\widetilde\omega)|\bigr)^{k}_{\mathrm{ac}}
 \le\mathcal E_{\omega}(\alpha,W,\epsilon).
\]
\end{lemma}
\begin{proof}
\uses{ana:pushforward-regularization,ana:mass-across-centre-fatou,ana:mass-across-centre}
\leanok
Lemma~\ref{ana:pushforward-regularization} supplies a sequence \((T_m)_m\) of currents
\(\epsilon\)-admissible along \(W\) for \(\alpha\) whose two clauses are exactly the two
hypotheses of Lemma~\ref{ana:mass-across-centre-fatou} at \(A=\operatorname{Exc}(g)\),
whose geometric hypotheses hold for the strict transform of a \(W\) not contained in
\(g(\operatorname{Exc} g)\) (Lemma~\ref{ana:mass-across-centre}). Hence
\[
 \int_{\widetilde W_{\mathrm{reg}}}\bigl((T'+\epsilon_0\widetilde\omega)|\bigr)^{k}_{\mathrm{ac}}
 \le\liminf_m\int_{W_{\mathrm{reg}}}\bigl((T_m+\epsilon\omega)|\bigr)^k_{\mathrm{ac}} .
\]
Each \(T_m\) is \(\epsilon\)-admissible along \(W\) for \(\alpha\), so each term of the
sequence is at most \(\mathcal E_{\omega}(\alpha,W,\epsilon)\), and therefore so is its
\(\liminf\).
\end{proof}

\begin{lemma}[Pushforward of an admissible current across the centre]
\label{ana:rv-pushforward-envelope}
\uses{ana:restricted-volume,ana:blowups,ana:strict-transform-blowup}
\source{ct22-restricted-volumes:invar}
\lean{TranscendentalNakamaye.Pending.RestrictedVolumePullbackUpper.exists_restrictedMassEnvelope_pullback_le_at_error,
TranscendentalNakamaye.Pending.RestrictedVolumePullbackEnvelope.exists_restrictedMassEnvelope_le_at_error}
\leanok
In the situation of Lemma~\ref{ana:rv-le-pullback}, for every \(\epsilon>0\) there is
\(\epsilon'>0\) such that the admissible mass envelope of \(\widetilde W\) at \(\epsilon'\)
for \(g^*\alpha\) is at most the admissible mass envelope of \(W\) at \(\epsilon\) for
\(\alpha\).
\end{lemma}
\begin{proof}
\uses{ana:blowup-compact,ana:exceptional-metric,ana:restricted-mass-background-change,ana:rv-pushforward-one-current,gnpp:kahler-domination}
\leanok
The manifold \(Z\) is compact (Lemma~\ref{ana:blowup-compact}), and
Lemma~\ref{ana:exceptional-metric} applied to the modification \(g\) gives an
effective \(\mathbf Q\)-divisor \(E\) with a smooth metric, supported on
\(\operatorname{Exc}(g)\), such that \(\widetilde\omega=g^*\omega-\Theta_E\) is a
K\"ahler form on \(Z\). By Lemma~\ref{gnpp:kahler-domination} there is \(A>0\)
with \(\omega_Z\le A\widetilde\omega\). Set \(\epsilon_0=\epsilon/2\) and
\(\epsilon'=\epsilon_0/A\).

Then \(\epsilon_0\widetilde\omega-\epsilon'\omega_Z=\epsilon'(A\widetilde\omega-\omega_Z)\)
is semipositive, so by Lemma~\ref{ana:restricted-mass-background-change}
\[
 \mathcal E_{\omega_Z}(g^*\alpha,\widetilde W,\epsilon')
 \le\mathcal E_{\widetilde\omega}(g^*\alpha,\widetilde W,\epsilon_0).
\]
Every \(\epsilon_0\)-admissible current along \(\widetilde W\) for
\(\widetilde\omega\) has restricted mass at most
\(\mathcal E_{\omega}(\alpha,W,\epsilon)\) by
Lemma~\ref{ana:rv-pushforward-one-current}, since \(2\epsilon_0=\epsilon\); taking
the supremum bounds the right-hand side above by
\(\mathcal E_{\omega}(\alpha,W,\epsilon)\). The witness
\(\epsilon'=\epsilon/(2A)\) depends only on \(\epsilon\), on \(E\) and on the two
K\"ahler forms.
\end{proof}

\begin{lemma}[Restricted volumes do not increase under one blowup]
\label{ana:rv-pullback-le}
\uses{ana:restricted-volume,ana:blowups,ana:nonkahler-locus,ana:strict-transform-blowup}
\source{ct22-restricted-volumes:invar}
\lean{TranscendentalNakamaye.Pending.RestrictedVolumeOneBlowup.numericalRestrictedVolume_pullback_le}
\leanok
In the situation of Lemma~\ref{ana:rv-le-pullback},
\(\operatorname{RV}_{Z\mid\widetilde W}(g^*\alpha)\le\operatorname{RV}_{Y\mid W}(\alpha)\).
\end{lemma}
\begin{proof}
\uses{ana:rv-pushforward-envelope,ana:strict-transform-loci,ana:restricted-volume-monotone}
\leanok
If \(W\subset E_{\mathrm{nn}}(\alpha)\) then, by the transformation rule
\(E_{\mathrm{nn}}(g^*\alpha)=g^{-1}E_{\mathrm{nn}}(\alpha)\) and closedness,
\(\widetilde W\subset E_{\mathrm{nn}}(g^*\alpha)\), so the left side is \(0\).
Otherwise \(\widetilde W\not\subset E_{\mathrm{nn}}(g^*\alpha)\)
(Lemma~\ref{ana:strict-transform-loci}) and both sides are infima over \(\epsilon>0\) of
the admissible mass envelopes (Lemma~\ref{ana:restricted-volume-monotone}); for each
downstairs \(\epsilon\), Lemma~\ref{ana:rv-pushforward-envelope} produces an upstairs
\(\epsilon'\) whose envelope is dominated, and taking infima gives the inequality.
\end{proof}

\begin{lemma}[Bimeromorphic invariance of restricted volumes for one blowup]
\label{ana:rv-birational-blowup}
\uses{ana:strict-transform-blowup,ana:rv-le-pullback,ana:rv-pullback-le}
\source{ct22-restricted-volumes:invar}
\lean{TranscendentalNakamaye.Pending.RestrictedVolumeOneBlowup.exists_strictTransform_numericalRestrictedVolume_eq_of_isBlowupAlong}
\leanok
Let \(Y\) be a compact complex manifold, \(g:Z\to Y\) the blowup along a nowhere
dense closed submanifold, \(\alpha=\{\theta\}\) a big class on \(Y\), and
\(W\subset Y\) a positive-dimensional closed irreducible analytic subset with
\(W\not\subset g(\operatorname{Exc}(g))\). Then the strict transform
\(\widetilde W\) is a closed irreducible analytic subset of dimension \(\dim W\) and
\(\operatorname{RV}_{Z\mid\widetilde W}(g^*\alpha)=\operatorname{RV}_{Y\mid W}(\alpha)\)
for all K\"ahler forms on \(Y\) and \(Z\).
\end{lemma}
\begin{proof}
\leanok
Lemma~\ref{ana:strict-transform-blowup} gives \(\widetilde W\), and
Lemmas~\ref{ana:rv-le-pullback} and~\ref{ana:rv-pullback-le} give the two
inequalities.
\end{proof}

\begin{lemma}[The regular-patch density is the determinant density]
\label{ana:patch-density-determinant}
\uses{ana:restricted-volume,fpma:top-power-determinant}
\source{ct22-restricted-volumes:sect2}
\lean{TranscendentalNakamaye.Pending.FullSpaceGeneralDensity.patchDensity_eq_normalizedTopPowerDensity}
\leanok
Let \(T=\theta+\epsilon\omega+dd^cu\) be an admissible current along an irreducible analytic
subvariety \(V\) of dimension \(k\), and let \(\gamma\) be a regular parametrization of
\(V\). At a parameter \(w\) whose image avoids the polar set of \(u\), the parameter
density of \(\gamma^*T\) equals \(k!\,\pi^{-k}\det H\), where \(H\) is the Hermitian
coefficient matrix of \(\gamma^*T\) at \(w\).
\end{lemma}
\begin{proof}
\uses{fpma:top-power-determinant}
\leanok
The differential of a regular parametrization is \(\mathbb C\)-linear, and off its polar set
an admissible current is a semipositive real \((1,1)\)-form in coordinates; hence the
pullback \(\gamma^*T(w)\) is a semipositive real \((1,1)\)-form and
Lemma~\ref{fpma:top-power-determinant} applies.
\end{proof}

\begin{lemma}[One regular closed ball carries the ambient measure]
\label{ana:fullspace-closed-ball-measure}
\uses{ana:patch-density-determinant,gnpp:global-measure}
\source{ct22-restricted-volumes:sect2}
\lean{TranscendentalNakamaye.Pending.FullSpaceGeneralClosedBall.exists_fullSpace_closedBall_lintegral_eq_globalMixedMeasure}
\leanok
Let \(V=X\) and let \(\gamma\) be a regular parametrization of \(X\) around a point off the
polar set of an admissible current \(T\). There is \(r>0\) such that the integral of the
parameter density of \(\gamma^*T\) over the closed ball \(\overline B(w,r)\) equals the
ambient non-pluripolar measure \(\operatorname{NP}(T)\) of \(\gamma(\overline B(w,r))\).
\end{lemma}
\begin{proof}
\uses{ana:patch-density-determinant}
\leanok
Off the poles the potential has a smooth local representative, so on a small coordinate
neighbourhood \(T=dd^c(g+\rho)\) with \(g+\rho\) smooth and plurisubharmonic. By
Lemma~\ref{ana:patch-density-determinant} and the holomorphic \(dd^c\) chain rule the
parameter density of \(\gamma^*T\) is the smooth Monge--Amp\`ere density of
\((g+\rho)\circ\gamma\), whose integral over \(\overline B(w,r)\) is the smooth measure of
the image. Shrinking the neighbourhood so that it lies in one strict truncation level of a
Bedford--Taylor patch identifies that smooth measure with the restriction of the ambient
non-pluripolar measure.
\end{proof}

\begin{lemma}[Localization of the closed-ball identity]
\label{ana:fullspace-density-localization}
\uses{ana:fullspace-closed-ball-measure}
\source{ct22-restricted-volumes:sect2}
\lean{TranscendentalNakamaye.Pending.FullSpaceAmbientComparison.lintegral_patchDensity_eq_globalMixedMeasure_image}
\leanok
Let \(V=X\) and let \(D\) be a measurable piece of a regular parametrization \(\gamma\) of
\(X\) whose image is regular and avoids the polar set of an admissible current \(T\). Then
\(\gamma(D)\) is measurable and the integral of the parameter density of \(\gamma^*T\) over
\(D\) equals \(\operatorname{NP}(T)(\gamma(D))\).
\end{lemma}
\begin{proof}
\uses{ana:fullspace-closed-ball-measure}
\leanok
The identity of Lemma~\ref{ana:fullspace-closed-ball-measure} holds, with the same local
data, for every measurable subset of the closed ball it produces: the only step of its proof
that used compactness of the piece was the measurability of \(\gamma(\,\cdot\,)\), and
\(\gamma\) is one leg of a biholomorphism between open sets, so \(\gamma(D)\) is the
preimage of \(D\) under the other leg intersected with an open set, hence measurable
whenever \(D\) is.

Now every parameter of \(D\) has such a closed ball, and the parameter domain is second
countable, so countably many of those balls cover \(D\). Intersect them with \(D\) and
disjointify the resulting countable family in parameter space; each member is a measurable
subset of one of the balls, so the identity of the previous paragraph applies to it. Both
sides are countably additive over that partition: the left side by countable additivity of
the integral, the right side because \(\gamma\) is injective on the parameter domain, so the
images are pairwise disjoint measurable sets. Summing gives the claim, and
\(\gamma(D)\) is the countable union of those measurable images.
\end{proof}

\begin{lemma}[Restricted mass on the whole space is at most the ambient mass]
\label{ana:fullspace-ambient-comparison}
\uses{ana:restricted-volume,ana:fullspace-density-localization}
\source{ct22-restricted-volumes:sect2}
\lean{TranscendentalNakamaye.Pending.FullSpaceAmbientComparison.acTopMass_le_perturbedDiagonalPotentialMass}
\leanok
For \(V=X\) and any admissible current \(T\) along \(X\), the absolutely continuous
restricted mass of \(T\) is at most the ambient non-pluripolar mass
\(\operatorname{NP}(T)(X)\).
\end{lemma}
\begin{proof}
\uses{ana:fullspace-density-localization}
\leanok
A finite family of disjoint measurable pieces \(D_i\) of regular parametrizations avoiding
the polar set has, by Lemma~\ref{ana:fullspace-density-localization}, total parameter mass
\(\sum_i\operatorname{NP}(T)(\gamma_i(D_i))\); the images are measurable and pairwise
disjoint, so the sum is \(\operatorname{NP}(T)\) of their union, hence at most
\(\operatorname{NP}(T)(X)\). Taking the supremum over the families gives the claim.
\end{proof}

\begin{lemma}[Raising the error does not decrease the ambient mass of a presentation]
\label{ana:ambient-mass-raise-error}
\uses{ana:restricted-volume,ana:restricted-volume-scaffold,mmass:mass}
\lean{TranscendentalNakamaye.Pending.PerturbedAmbientEnvelopeMonotone.perturbedDiagonalPotentialMass_le_raiseError}
\leanok
Let \(0<\epsilon\le\epsilon'\) and let \(T=\theta+dd^cu\) be an
\(\epsilon\)-admissible presentation of \(\alpha\) along \(V\). Then the same
representative and potential form an \(\epsilon'\)-admissible presentation, and
\[
 \operatorname{NP}\bigl(\theta+\epsilon\omega+dd^cu\bigr)(X)
 \ \le\ \operatorname{NP}\bigl(\theta+\epsilon'\omega+dd^cu\bigr)(X).
\]
\end{lemma}
\begin{proof}
\uses{mmass:diagonal-background-monotone,lpot:semipositive-enlargement,posclass:background}
\leanok
Admissibility at \(\epsilon'\) is the enlargement corollary
\ref{lpot:semipositive-enlargement} applied to the semipositive form
\((\epsilon'-\epsilon)\omega\): the representative, the potential, its analytic
singularity type and the non-degeneracy along \(V\) are all unchanged.

For the masses, if \(\epsilon=\epsilon'\) there is nothing to prove. Otherwise put
\(c=\epsilon'-\epsilon>0\) and \(\alpha_0=\theta+\epsilon\omega\), so that
\(\theta+\epsilon'\omega=\alpha_0+c\omega\) with \(\omega\) semipositive. The two
masses are the diagonal non-pluripolar total masses of the one potential \(u\) in
the two backgrounds \(\alpha_0\) and \(\alpha_0+c\omega\), so
Theorem~\ref{mmass:diagonal-background-monotone} gives the inequality.
\end{proof}

\begin{lemma}[The perturbed ambient volume envelope is nondecreasing in the error]
\label{ana:ambient-volume-envelope-monotone}
\uses{ana:restricted-volume,ana:restricted-volume-scaffold,mmass:mass}
\lean{TranscendentalNakamaye.Pending.PerturbedAmbientEnvelopeMonotone.perturbedDiagonalPotentialEnvelope_mono}
\leanok
For \(0<\epsilon\le\epsilon'\),
\[
 \sup_{T\in\mathcal A_\epsilon}\operatorname{NP}(T+\epsilon\omega)(X)
 \ \le\ \sup_{T\in\mathcal A_{\epsilon'}}\operatorname{NP}(T+\epsilon'\omega)(X),
\]
the suprema being over the \(\epsilon\)- and \(\epsilon'\)-admissible
presentations of \(\alpha\) along \(V\).
\end{lemma}
\begin{proof}
\uses{ana:ambient-mass-raise-error}
\leanok
Each \(\epsilon\)-admissible presentation is carried to an
\(\epsilon'\)-admissible one without decreasing its ambient mass
(Lemma~\ref{ana:ambient-mass-raise-error}); its mass is therefore at most the
supremum on the right. Take the supremum on the left.
\end{proof}

\begin{lemma}[A positive cluster of ambient volumes is a positive bound at every error]
\label{ana:ambient-volume-uniform-bound}
\uses{ana:ambient-volume-envelope-monotone}
\lean{TranscendentalNakamaye.Pending.PerturbedAmbientEnvelopeMonotone.exists_hasUniformAmbientMassLowerBound_of_hasPositiveAmbientMassCluster,TranscendentalNakamaye.Pending.PerturbedAmbientEnvelopeMonotone.isBigClass_of_hasPositiveAmbientMassCluster_of_uniform}
\leanok
Write \(E(\epsilon)=\sup_{T\in\mathcal A_\epsilon}\operatorname{NP}(T+\epsilon\omega)(X)\).
Let \((\epsilon_k)_k\) be positive errors with \(\epsilon_k\to0\), and suppose there
is \(\eta>0\) with \(\eta<E(\epsilon_k)\) for arbitrarily large \(k\). Then
\(\eta<E(\epsilon)\) for \emph{every} \(\epsilon>0\).

Consequently, to prove that a pseudoeffective \(\alpha\) with such a cluster is
big it suffices to prove the sequence-free implication: if \(E(\epsilon)>\eta\)
for every \(\epsilon>0\) and some \(\eta>0\), then \(\alpha\) is big.
\end{lemma}
\begin{proof}
\uses{ana:ambient-volume-envelope-monotone}
\leanok
Fix \(\epsilon>0\). Since \(\epsilon_k\to0\) there is \(N\) with
\(\epsilon_k<\epsilon\) for all \(k\ge N\), and by hypothesis some \(k\ge N\) has
\(\eta<E(\epsilon_k)\). Monotonicity
(Lemma~\ref{ana:ambient-volume-envelope-monotone}) gives
\(E(\epsilon_k)\le E(\epsilon)\), hence \(\eta<E(\epsilon)\).

The second assertion is immediate from the first: the cluster hypothesis supplies
some \(\eta>0\), and the first part upgrades it to the bound at every error, which
is the hypothesis of the implication.
\end{proof}

\begin{lemma}[Minkowski's determinant inequality]
\label{ana:minkowski-determinant}
\source{bou02-volume:3.5,bou02-volume:page-0014}
\lean{TranscendentalNakamaye.Pending.MinkowskiDeterminant.one_add_rpow_prod_le,TranscendentalNakamaye.Pending.MinkowskiDeterminant.det_eq_ofReal_prod_eigenvalues,TranscendentalNakamaye.Pending.MinkowskiDeterminant.re_det_nonneg,TranscendentalNakamaye.Pending.MinkowskiDeterminant.det_eq_ofReal_re,TranscendentalNakamaye.Pending.MinkowskiDeterminant.det_one_add_eq_prod,TranscendentalNakamaye.Pending.MinkowskiDeterminant.rpow_inv_card_re_det_add_le_of_ne_zero,TranscendentalNakamaye.Pending.MinkowskiDeterminant.rpow_inv_card_re_det_add_le,TranscendentalNakamaye.Pending.MinkowskiDeterminant.re_det_le_re_det_add,TranscendentalNakamaye.Pending.MinkowskiDeterminant.rpow_inv_card_re_det_smul,TranscendentalNakamaye.Pending.MinkowskiDeterminant.convex_setOf_posSemidef,TranscendentalNakamaye.Pending.MinkowskiDeterminant.concaveOn_rpow_inv_card_re_det}
\leanok
Let \(A,B\) be Hermitian nonnegative \(n\times n\) complex matrices, \(n\ge1\). Then
\[
 \det(A)^{1/n}+\det(B)^{1/n}\le\det(A+B)^{1/n};
\]
in particular \(\det A\le\det(A+B)\). The map \(A\mapsto\det(A)^{1/n}\) is homogeneous of degree one and concave on the convex cone of
Hermitian nonnegative matrices.
\end{lemma}
\begin{proof}
\leanok
For nonnegative reals \(\lambda_1,\dots,\lambda_n\), the arithmetic--geometric mean inequality applied to the numbers
\(1/(1+\lambda_i)\) and to the numbers \(\lambda_i/(1+\lambda_i)\), whose sums add up to \(n\), gives
\[
 \Bigl(\prod_i\frac1{1+\lambda_i}\Bigr)^{1/n}+\Bigl(\prod_i\frac{\lambda_i}{1+\lambda_i}\Bigr)^{1/n}\le1,\qquad\text{that is}\qquad
 1+\Bigl(\prod_i\lambda_i\Bigr)^{1/n}\le\Bigl(\prod_i(1+\lambda_i)\Bigr)^{1/n}.
\]
If \(\det A=\det B=0\) the inequality is clear; by symmetry let \(A\) be invertible and write \(A=y^*y\). Then \(X=(y^{-1})^*By^{-1}\) is
Hermitian nonnegative, with eigenvalues \(\lambda_i\ge0\), \(A+B=y^*(1+X)y\) and \(B=y^*Xy\), so \(\det(A+B)=\prod_i(1+\lambda_i)\det A\) and
\(\det B=\prod_i\lambda_i\det A\) (diagonalize \(X\) by a unitary matrix); the displayed inequality, multiplied by \(\det(A)^{1/n}\), is the
claim. Monotonicity follows, since \(\det(B)^{1/n}\ge0\). Homogeneity is \(\det(tA)=t^n\det A\), and concavity follows from homogeneity and
the inequality applied to \(tA\) and \((1-t)B\).
\end{proof}
This is the matrix inequality behind the semicontinuity of the top power of the absolutely continuous part of positive currents and the
log-concavity of the volume in Boucksom's proof of Lemma~\ref{ana:volume-semicontinuity}.

\begin{lemma}[The determinant as an infimum of traces]
\label{ana:determinant-trace-inf}
\uses{ana:minkowski-determinant}
\source{bou02-volume:3.5}
\lean{TranscendentalNakamaye.Pending.DeterminantTraceInf.re_trace_eq_sum_eigenvalues,TranscendentalNakamaye.Pending.DeterminantTraceInf.rpow_inv_card_re_det_le_re_trace_div,TranscendentalNakamaye.Pending.DeterminantTraceInf.re_trace_mul_nonneg,TranscendentalNakamaye.Pending.DeterminantTraceInf.rpow_inv_card_re_det_mul_le,TranscendentalNakamaye.Pending.DeterminantTraceInf.posDef_add_smul_one,TranscendentalNakamaye.Pending.DeterminantTraceInf.exists_posDef_det_eq_one_re_trace_div_le,TranscendentalNakamaye.Pending.DeterminantTraceInf.isGLB_re_trace_div}
\leanok
Let \(A\) be a Hermitian nonnegative \(n\times n\) complex matrix, \(n\ge1\). Then \(\det(X)^{1/n}\le\operatorname{tr}(X)/n\) for every Hermitian nonnegative
\(X\), \(\det(A)^{1/n}\det(B)^{1/n}\le\operatorname{tr}(BA)/n\) for every Hermitian nonnegative \(B\), and
\[
 \det(A)^{1/n}=\inf\Bigl\{\tfrac1n\operatorname{tr}(BA)\ :\ B\ \text{positive definite},\ \det B=1\Bigr\}.
\]
\end{lemma}
\begin{proof}
\leanok
The first inequality is the inequality of arithmetic and geometric means for the eigenvalues of \(X\), which are nonnegative, with product \(\det X\) and sum
\(\operatorname{tr}X\). For the second write \(B=y^*y\) and apply the first to \(X=yAy^*\): \(\det X=\det A\det B\) and \(\operatorname{tr}X=\operatorname{tr}(y^*yA)=\operatorname{tr}(BA)\).
Hence \(\det(A)^{1/n}\) is a lower bound of the set. For \(t>0\) the matrix \(A_t=A+t\) is positive definite; \(B_t=\det(A_t)^{1/n}A_t^{-1}\) is positive definite with
\(\det B_t=1\), and \(\operatorname{tr}(A_t^{-1}A)=\operatorname{tr}(1-tA_t^{-1})\le n\) because \(A_t^{-1}\) is nonnegative, so
\(\frac1n\operatorname{tr}(B_tA)\le\det(A_t)^{1/n}\), which tends to \(\det(A)^{1/n}\) as \(t\downarrow0\).
\end{proof}
For a positive \((1,1)\)-form \(\alpha\) with matrix \(A\) and a positive definite \((1,1)\)-form \(\beta\) with matrix \(B\), one has
\(\alpha\wedge\beta^{n-1}=(n-1)!\det(B)\operatorname{tr}(B^{-1}A)\,\omega_0^n/n!\); with this identity (Lemma~\ref{ana:mixed-trace-infimum} below) the lemma identifies the infimum \(D_C\) of
Lemma~\ref{ana:ac-density-liminf}, for \(q=n-1\) and \(C\) the forms \(\beta^{n-1}/n!\) with \(\beta^n=\omega_0^n\), with \(\det(T_{\mathrm{ac}})^{1/n}\).

\begin{lemma}[The mixed product \(\beta^{n-1}\wedge\alpha\) as a trace, and \(\det^{1/n}\) as an infimum over \(\beta\)]
\label{ana:mixed-trace-infimum}
\uses{ana:determinant-trace-inf,fpmma:coefficient,fpmma:geometric-density,fpma:top-power-determinant-general,bt:constant-form-currents}
\source{bou02-volume:3.5}
\lean{TranscendentalNakamaye.Pending.OneOneFormMixedTrace.sum_perm_apply,TranscendentalNakamaye.Pending.OneOneFormMixedTrace.det_updateRow_eq_sum_adjugate,TranscendentalNakamaye.Pending.OneOneFormMixedTrace.mixedCoefficient_update_const,TranscendentalNakamaye.Pending.OneOneFormMixedTrace.normalizedHermitianMatrix_snoc,TranscendentalNakamaye.Pending.OneOneFormMixedTrace.topDensity_wedge_formPower,TranscendentalNakamaye.Pending.OneOneFormMixedTrace.exists_isSemipositiveOneOne_normalizedHermitianMatrix_eq,TranscendentalNakamaye.Pending.OneOneFormMixedTrace.isGLB_topDensity_wedge_formPower}
\leanok
Let \(n\ge1\). For a real \((1,1)\)-form \(\gamma\) on \(\mathbb C^n\) write \(H_\gamma\) for its normalized Hermitian matrix,
\((H_\gamma)_{jk}=\pi\bigl(\gamma(e_j,ie_k)-i\gamma(e_j,e_k)\bigr)\), and \(\langle\eta\rangle\) for the value of a top-degree form \(\eta\) on the standard real
frame \(e_1,ie_1,\dots,e_n,ie_n\), so that \(\langle\gamma^n\rangle=n!\,\pi^{-n}\operatorname{Re}\det H_\gamma\)
(Lemma~\ref{fpma:top-power-determinant-general}).
\begin{enumerate}
\item For complex \(n\times n\) matrices \(A,B\), the mixed row coefficient of the family \((B,\dots,B,A,B,\dots,B)\) (one slot \(A\)) is
 \((n-1)!\operatorname{tr}(A\operatorname{adj}B)\).
\item For real \((1,1)\)-forms \(\alpha,\beta\),
 \[
  \langle\beta^{n-1}\wedge\alpha\rangle=\frac{(n-1)!}{\pi^n}\operatorname{Re}\operatorname{tr}\bigl(H_\alpha\operatorname{adj}H_\beta\bigr);
 \]
 in particular \(\langle\beta^{n-1}\wedge\alpha\rangle=\frac{(n-1)!}{\pi^n}\det(H_\beta)\operatorname{Re}\operatorname{tr}(H_\beta^{-1}H_\alpha)\) when \(H_\beta\) is invertible.
\item Every Hermitian nonnegative \(n\times n\) matrix is \(H_\beta\) for a semipositive \((1,1)\)-form \(\beta\).
\item For a semipositive \((1,1)\)-form \(\alpha\),
 \[
  \inf\Bigl\{\langle\beta^{n-1}\wedge\alpha\rangle\ :\ \beta\ \text{semipositive},\ H_\beta\ \text{positive definite},\ \det H_\beta=1\Bigr\}
  =\frac{n!}{\pi^n}\Bigl(\frac{\pi^n}{n!}\langle\alpha^n\rangle\Bigr)^{1/n}=\frac{n!}{\pi^n}\det(H_\alpha)^{1/n}.
 \]
\end{enumerate}
\end{lemma}
\begin{proof}
\leanok
(1) For a permutation \(\sigma\), the summand of the mixed coefficient (Definition~\ref{fpmma:coefficient}) is the determinant of \(B\) with its row
\(j=\sigma^{-1}(i)\) replaced by the row \(j\) of \(A\), where \(i\) is the slot of \(A\); by the Laplace expansion along that row it equals
\(\sum_kA_{jk}(\operatorname{adj}B)_{kj}\). As \(\sigma\) runs over \(S_n\), \(\sigma^{-1}(i)\) takes every value \(j\) exactly \((n-1)!\) times, so the sum is
\((n-1)!\sum_{j,k}A_{jk}(\operatorname{adj}B)_{kj}=(n-1)!\operatorname{tr}(A\operatorname{adj}B)\).

(2) \(\beta^{n-1}\wedge\alpha\) is the wedge product of the family \((\beta,\dots,\beta,\alpha)\) of real \((1,1)\)-forms, whose value on the standard frame is
\(\pi^{-n}\operatorname{Re}\) of the mixed row coefficient of \((H_\beta,\dots,H_\beta,H_\alpha)\) (Lemma~\ref{fpmma:geometric-density}); apply (1). If \(H_\beta\) is
invertible, \(\operatorname{adj}H_\beta=\det(H_\beta)H_\beta^{-1}\).

(3) Let \(M\) be Hermitian nonnegative. Then \(N=\pi^{-1}\overline M=\pi^{-1}M^{\mathsf T}\) is Hermitian nonnegative, so \(N=y^*y\) for a complex matrix \(y\). Let
\(S\) be the \(\mathbb C\)-linear map with matrix \(y\) and \(\beta=S^*\omega_0\) the pullback of the form \(\omega_0\) of Definition~\ref{bt:constant-form-currents}; it is semipositive, since
\(\beta(v,iv)=\omega_0(Sv,iSv)=|Sv|^2\ge0\) by \(\mathbb C\)-linearity of \(S\). The matrix of a pullback is \(H_{S^*\omega_0}=\pi\,\overline{y^*y}\), which is \(\pi\overline N=M\).

(4) Put \(A=H_\alpha\), Hermitian nonnegative. By (2), for \(\beta\) in the set, \(\langle\beta^{n-1}\wedge\alpha\rangle=\frac{n!}{\pi^n}\cdot\frac1n\operatorname{Re}\operatorname{tr}(H_\beta^{-1}A)\),
and \(H_\beta^{-1}\) is positive definite of determinant \(1\). Conversely, by (3) every positive definite \(B\) of determinant \(1\) is \(H_\beta^{-1}\) for a
semipositive \(\beta\) (take \(H_\beta=B^{-1}\)), which then lies in the set. Hence the set is \(\frac{n!}{\pi^n}\) times the set of
Lemma~\ref{ana:determinant-trace-inf}, whose infimum is \(\det(A)^{1/n}\), and \(\det A=\frac{\pi^n}{n!}\langle\alpha^n\rangle\).
\end{proof}

\begin{lemma}[Boucksom: semicontinuity of the absolutely continuous part, liminf form]
\label{ana:ac-density-liminf}
\uses{bt:lebesgue-decomposition,bt:order-zero,bt:constant-form-currents}
\source{bou02-volume:3.5}
\lean{TranscendentalNakamaye.Pending.LocalCurrentSemicontinuity.bump,TranscendentalNakamaye.Pending.LocalCurrentSemicontinuity.bump_rOut,TranscendentalNakamaye.Pending.LocalCurrentSemicontinuity.bump_rOut_le,TranscendentalNakamaye.Pending.LocalCurrentSemicontinuity.tendsto_bump_rOut,TranscendentalNakamaye.Pending.LocalCurrentSemicontinuity.eventually_bump_rOut_lt,TranscendentalNakamaye.Pending.LocalCurrentSemicontinuity.mol,TranscendentalNakamaye.Pending.LocalCurrentSemicontinuity.mol_nonneg,TranscendentalNakamaye.Pending.LocalCurrentSemicontinuity.contDiff_mol,TranscendentalNakamaye.Pending.LocalCurrentSemicontinuity.continuous_mol,TranscendentalNakamaye.Pending.LocalCurrentSemicontinuity.mol_eq_zero,TranscendentalNakamaye.Pending.LocalCurrentSemicontinuity.dist_lt_of_mol_ne_zero,TranscendentalNakamaye.Pending.LocalCurrentSemicontinuity.integral_mol,TranscendentalNakamaye.Pending.LocalCurrentSemicontinuity.integrable_mol,TranscendentalNakamaye.Pending.LocalCurrentSemicontinuity.hasCompactSupport_mol,TranscendentalNakamaye.Pending.LocalCurrentSemicontinuity.tsupport_mol_subset,TranscendentalNakamaye.Pending.LocalCurrentSemicontinuity.bumpTest,TranscendentalNakamaye.Pending.LocalCurrentSemicontinuity.bumpTest_apply,TranscendentalNakamaye.Pending.LocalCurrentSemicontinuity.eventually_closedBall_subset,TranscendentalNakamaye.Pending.LocalCurrentSemicontinuity.ae_of_forall_ball,TranscendentalNakamaye.Pending.LocalCurrentSemicontinuity.ae_tendsto_integral_mol_mul,TranscendentalNakamaye.Pending.LocalCurrentSemicontinuity.ae_tendsto_integral_mol_singular,TranscendentalNakamaye.Pending.LocalCurrentSemicontinuity.apply_bumpTest,TranscendentalNakamaye.Pending.LocalCurrentSemicontinuity.ae_tendsto_apply_bumpTest,TranscendentalNakamaye.Pending.LocalCurrentSemicontinuity.acInf,TranscendentalNakamaye.Pending.LocalCurrentSemicontinuity.acInf_le,TranscendentalNakamaye.Pending.LocalCurrentSemicontinuity.acInf_eq_of_subset_closure,TranscendentalNakamaye.Pending.LocalCurrentSemicontinuity.measurable_acInf,TranscendentalNakamaye.Pending.LocalCurrentSemicontinuity.lintegral_mol_mul_acInf_le,TranscendentalNakamaye.Pending.LocalCurrentSemicontinuity.ae_liminf_acInf_le}
\leanok
Let \(U\subset\mathbb C^n\) be open, \(q\ge0\), and let \(T_k\), \(T\) be positive currents of bidimension \((q,q)\) on \(U\) with
\(T_k\to T\) weakly, i.e.\ \(T_k(\gamma)(\varphi)\to T(\gamma)(\varphi)\) for every constant \((q,q)\)-form \(\gamma\in\Lambda_q\) and every test function \(\varphi\).
With the notation of Proposition~\ref{bt:lebesgue-decomposition}, write \(a_T(x)=h(x)\,L_x\), so that \(T_{\mathrm{ac}}(\gamma)=a_T(\gamma)\,\lambda\). For a set
\(C\) of strongly positive constant \((q,q)\)-forms put
\[
 D_C(T)(x)=\inf_{c\in C}a_T(x)(c)\in[0,\infty].
\]
Then for Lebesgue-almost every \(x\in U\),
\[
 D_C(T)(x)\ \ge\ \liminf_{k\to\infty}D_C(T_k)(x).
\]
Moreover, for every \(\gamma\in\Lambda_q\) and Lebesgue-almost every \(x\in U\), \(T(\gamma)(\rho_\varepsilon(x-\cdot))\to a_T(x)(\gamma)\) as
\(\varepsilon\to0\), where \(\rho_\varepsilon\) is a normalized smooth bump function supported in the ball of radius \(\varepsilon\).
\end{lemma}
\begin{proof}
\leanok
\emph{Mollified values.} Fix \(\gamma\) and write \(T(\gamma)=f_\gamma\sigma_T=a_T(\gamma)\lambda+L(\gamma)\sigma_T^{\mathrm{sg}}\). For the absolutely continuous
part, \(\int\rho_\varepsilon(x-y)a_T(y)(\gamma)\,dy\to a_T(x)(\gamma)\) almost everywhere by Lebesgue's differentiation theorem for mollifiers (applied to
\(a_T(\gamma)\) cut off to a closed ball inside \(U\), which is integrable). For the singular part, \(|L_y(\gamma)|\le M\) and \(\rho_\varepsilon\le
K/\lambda(B(x,\varepsilon))\) on \(B(x,\varepsilon)\), so its mollified value is at most \(MK\,\sigma^{\mathrm{sg}}_T(B(x,\varepsilon))/\lambda(B(x,\varepsilon))\), which tends to
\(d\sigma^{\mathrm{sg}}_T/d\lambda=0\) almost everywhere (Besicovitch's differentiation theorem, applied to the finite restrictions of \(\sigma^{\mathrm{sg}}_T\) to closed
balls inside \(U\)). A property that holds almost everywhere near every point of \(U\) holds almost everywhere on \(U\) (a countable subcover).

\emph{Semicontinuity.} If \(C=\emptyset\) both sides are \(+\infty\). Otherwise replace \(C\) by a countable dense subset: \(a_T(x)\) is continuous on the
finite-dimensional space \(\Lambda_q\), so the infimum does not change, and \(D_C(T)\) becomes measurable. Let \(G=\liminf_kD_C(T_k)\). For \(c\in C\) and a mollifier
\(\rho=\rho_\varepsilon(z-\cdot)\) whose ball lies in \(U\): the singular part of \(T_k\) is positive and \(\rho\ge0\), so
\(\int\rho\,D_C(T_k)\,d\lambda\le\int\rho\,a_{T_k}(c)\,d\lambda=T_{k,\mathrm{ac}}(c)(\rho)\le T_k(c)(\rho)\). By Fatou's lemma and the weak convergence,
\[
 \int\rho\,G\,d\lambda\le\liminf_k\int\rho\,D_C(T_k)\,d\lambda\le\lim_kT_k(c)(\rho)=T(c)(\rho).
\]
Since \(\rho\) is a positive constant on the ball of half the radius, \(G\) has finite integral on a closed ball around every point of \(U\); hence \(G\) is finite
almost everywhere on \(U\) and, by the first part applied to \(G\), \(\int\rho_\varepsilon(x-y)G(y)\,dy\to G(x)\) for almost every \(x\in U\). At such \(x\) which is
also a point of convergence of the mollified values of \(T(c)\) for the countably many \(c\) of \(C\), letting \(\varepsilon\to0\) in the last display gives
\(G(x)\le a_T(x)(c)\) for every \(c\in C\).
\end{proof}
This is Boucksom's semicontinuity lemma in the form its proof gives: with \(\liminf\) (the source states \(\limsup\), obtained afterwards by passing to subsequences
pointwise), and for every functional \(D_C\). Boucksom uses the concavity of \(A\mapsto\det(A)^{1/n}\) on Hermitian nonnegative matrices; here the infimum over \(C\)
plays that role, since an infimum of linear functionals passes through the integrals and limits directly. For \(q=n-1\) and \(C\) the forms \(\beta^{n-1}/n!\) of the
positive definite \((1,1)\)-forms \(\beta\) with \(\beta^n=\omega_0^n\), \(D_C(T)\) is \(\det(T_{\mathrm{ac}})^{1/n}\) relative to \(\omega_0\), by the inequality of
arithmetic and geometric means (Lemma~\ref{ana:minkowski-determinant} is the same matrix inequality); that identification is not part of this lemma.

\begin{lemma}[Duality between \(\Lambda_{n-1}\) and real \((1,1)\)-forms]
\label{ana:one-one-density-duality}
\uses{ana:mixed-trace-infimum,ana:determinant-trace-inf,bt:constant-form-currents,bt:form-products}
\source{bou02-volume:3.5}
\lean{TranscendentalNakamaye.Pending.LocalCurrentDeterminant.oneWedge_apply_cons,TranscendentalNakamaye.Pending.LocalCurrentDeterminant.oneWedge_oneWedge_apply_cons_cons,TranscendentalNakamaye.Pending.LocalCurrentDeterminant.eq_zero_of_forall_topDensity_oneWedge,TranscendentalNakamaye.Pending.LocalCurrentDeterminant.topDensity_coordinateWedge_oneForms,TranscendentalNakamaye.Pending.LocalCurrentDeterminant.pairing,TranscendentalNakamaye.Pending.LocalCurrentDeterminant.pairing_apply,TranscendentalNakamaye.Pending.LocalCurrentDeterminant.pairing_injective,TranscendentalNakamaye.Pending.LocalCurrentDeterminant.exists_eq_topDensity,TranscendentalNakamaye.Pending.LocalCurrentDeterminant.det_complexJ,TranscendentalNakamaye.Pending.LocalCurrentDeterminant.topDensity_compJ,TranscendentalNakamaye.Pending.LocalCurrentDeterminant.topDensity_coordinateWedge_pi11,TranscendentalNakamaye.Pending.LocalCurrentDeterminant.exists_isOneOne_eq_topDensity,TranscendentalNakamaye.Pending.LocalCurrentDeterminant.adjugate_eq_det_smul_inv,TranscendentalNakamaye.Pending.LocalCurrentDeterminant.posSemidef_of_nonneg,TranscendentalNakamaye.Pending.LocalCurrentDeterminant.isGLB_topDensity_wedge_formPower_of_posSemidef}
\leanok
Let \(n\ge1\), let \(\Lambda_{n-1}\) be the space of constant real \((n-1,n-1)\)-forms on \(\mathbb C^n\) of Definition~\ref{bt:constant-form-currents}, and write
\(\langle\eta\rangle\) and \(H_\gamma\) as in Lemma~\ref{ana:mixed-trace-infimum}.
\begin{enumerate}
\item A real alternating form \(\gamma\) of degree \(2n-2\) on \(\mathbb C^n\) with \(\langle\ell_1\wedge\ell_2\wedge\gamma\rangle=0\) for all real linear forms \(\ell_1,\ell_2\)
 is zero.
\item The pairing \((\gamma,\alpha)\mapsto\langle\gamma\wedge\alpha\rangle\) of \(\Lambda_{n-1}\) with the real two-forms is nondegenerate on the \(\Lambda_{n-1}\) side, and
 every linear functional \(L\) on \(\Lambda_{n-1}\) is \(L(\gamma)=\langle\gamma\wedge\alpha\rangle\) for a real \((1,1)\)-form \(\alpha\).
\item If moreover \(L\ge0\) on the strongly positive forms and \(\alpha\) is a real \((1,1)\)-form with \(L=\langle\cdot\wedge\alpha\rangle\), then \(H_\alpha\) is
 Hermitian nonnegative and
 \[
  \inf\Bigl\{L(\beta^{n-1})\ :\ \beta\ \text{semipositive},\ H_\beta\ \text{positive definite},\ \det H_\beta=1\Bigr\}=\frac{n!}{\pi^n}\det(H_\alpha)^{1/n}.
 \]
\end{enumerate}
\end{lemma}
\begin{proof}
\leanok
(1) Suppose \(\gamma(v_1,\dots,v_{2n-2})\ne0\). Then the \(v_i\) are linearly independent; extend them by \(x,y\) to a real basis of \(\mathbb C^n\) and let \(\ell_1,\ell_2\) be
the coordinate functionals of \(x\) and \(y\) in this basis. Expanding \((\ell\wedge A)(u_0,\dots,u_k)=\sum_i(-1)^i\ell(u_i)A(u_0,\dots,\widehat{u_i},\dots,u_k)\) twice
gives \((\ell_1\wedge\ell_2\wedge\gamma)(x,y,v_1,\dots,v_{2n-2})=\gamma(v)\ne0\). But a top-degree form is the multiple \(\langle\eta\rangle\det\) of the determinant
in the standard frame, so \(\langle\ell_1\wedge\ell_2\wedge\gamma\rangle=0\) forces \(\ell_1\wedge\ell_2\wedge\gamma=0\), a contradiction.

(2) Two-forms commute with the forms of \(\Lambda_{n-1}\) and the product is associative (Lemma~\ref{bt:form-products}), so
\(\langle\gamma\wedge(\ell_1\wedge\ell_2)\rangle=\langle\ell_1\wedge\ell_2\wedge\gamma\rangle\); if \(\gamma\in\Lambda_{n-1}\) pairs to zero with every two-form, (1) gives
\(\gamma=0\). The spaces are finite dimensional, so the transpose of the injective map \(\gamma\mapsto\langle\gamma\wedge\cdot\rangle\) is surjective, and with the bidual
identification \(L=\langle\cdot\wedge\alpha\rangle\) for a real two-form \(\alpha\). Replace \(\alpha\) by its \((1,1)\)-part \(\Pi_{11}\alpha=\frac12(\alpha+J^*\alpha)\),
\(J=i\): the real determinant of \(J\) is \(|\det_{\mathbb C}J|^2=|i^n|^2=1\), so \(\langle J^*\eta\rangle=\langle\eta\rangle\) for top-degree \(\eta\), and the forms of
\(\Lambda_{n-1}\) are \(J\)-invariant (Lemma~\ref{bt:form-products}); hence \(\langle\gamma\wedge J^*\alpha\rangle=\langle J^*(\gamma\wedge\alpha)\rangle=\langle\gamma\wedge\alpha\rangle\).

(3) Let \(C\) be positive definite. By Lemma~\ref{ana:mixed-trace-infimum}(3) there is a semipositive \(\beta\) with \(H_\beta=C^{-1}\); then \(\beta^{n-1}\) is strongly
positive and, by Lemma~\ref{ana:mixed-trace-infimum}(2) with \(\operatorname{adj}(C^{-1})=\det(C^{-1})\,C\),
\(0\le L(\beta^{n-1})=\frac{(n-1)!}{\pi^n}\det(C)^{-1}\operatorname{Re}\operatorname{tr}(H_\alpha C)\). So \(\operatorname{Re}\operatorname{tr}(H_\alpha C)\ge0\) for every
positive definite \(C\); for \(C=vv^*+t\) and \(t\downarrow0\) this gives \(v^*H_\alpha v\ge0\), and \(H_\alpha\) is Hermitian since \(\alpha\) is of type \((1,1)\). The
infimum is computed as in Lemma~\ref{ana:mixed-trace-infimum}(4), which uses only \(H_\alpha\ge0\): the set is \(\frac{n!}{\pi^n}\) times the set of
Lemma~\ref{ana:determinant-trace-inf} for \(A=H_\alpha\).
\end{proof}

\begin{lemma}[Boucksom: semicontinuity of \(\det(T_{\mathrm{ac}})^{1/n}\)]
\label{ana:ac-determinant-semicontinuity}
\uses{ana:one-one-density-duality,ana:ac-density-liminf,bt:lebesgue-decomposition}
\source{bou02-volume:3.5}
\lean{TranscendentalNakamaye.Pending.LocalCurrentDeterminant.detCone,TranscendentalNakamaye.Pending.LocalCurrentDeterminant.detCone_subset_posCone,TranscendentalNakamaye.Pending.LocalCurrentDeterminant.exists_isOneOne_acDensity_eq,TranscendentalNakamaye.Pending.LocalCurrentDeterminant.ae_acInf_detCone_eq,TranscendentalNakamaye.Pending.LocalCurrentDeterminant.ae_liminf_rpow_det_le,TranscendentalNakamaye.Pending.LocalCurrentDeterminant.acForm,TranscendentalNakamaye.Pending.LocalCurrentDeterminant.isOneOne_acForm,TranscendentalNakamaye.Pending.LocalCurrentDeterminant.acWeight_mul_densityForm_eq,TranscendentalNakamaye.Pending.LocalCurrentDeterminant.acMatrix,TranscendentalNakamaye.Pending.LocalCurrentDeterminant.acWeight_mul_densityForm_formPower,TranscendentalNakamaye.Pending.LocalCurrentDeterminant.ae_posSemidef_acMatrix,TranscendentalNakamaye.Pending.LocalCurrentDeterminant.ae_acInf_detCone_eq_acMatrix,TranscendentalNakamaye.Pending.LocalCurrentDeterminant.ae_liminf_rpow_det_acMatrix_le}
\leanok
Let \(U\subset\mathbb C^n\) be open, \(n\ge1\), and let \(T\) be a positive current of bidimension \((n-1,n-1)\) on \(U\), with \(a_T(x)=h(x)L_x\) as in
Lemma~\ref{ana:ac-density-liminf}. For every \(x\) there is a real \((1,1)\)-form \(\alpha_x\) with \(a_T(x)(\gamma)=\langle\gamma\wedge\alpha_x\rangle\) on
\(\Lambda_{n-1}\); put \(A_T(x)=H_{\alpha_x}\), the \emph{density matrix} of \(T_{\mathrm{ac}}\), so that
\(a_T(x)(\beta^{n-1})=\frac{(n-1)!}{\pi^n}\operatorname{Re}\operatorname{tr}\bigl(A_T(x)\operatorname{adj}H_\beta\bigr)\) for semipositive \(\beta\). Let
\(C=\{\beta^{n-1}:\beta\ \text{semipositive},\ H_\beta\ \text{positive definite},\ \det H_\beta=1\}\). Then for almost every \(x\in U\), \(A_T(x)\) is Hermitian
nonnegative and \(D_C(T)(x)=\frac{n!}{\pi^n}\det(A_T(x))^{1/n}\), for every choice of \(\alpha_x\). If positive currents \(T_k\) of bidimension \((n-1,n-1)\) converge
weakly to \(T\), then for almost every \(x\in U\)
\[
 \det\bigl(A_T(x)\bigr)^{1/n}\ \ge\ \liminf_{k\to\infty}\det\bigl(A_{T_k}(x)\bigr)^{1/n}.
\]
\end{lemma}
\begin{proof}
\leanok
The form \(\alpha_x\) exists by Lemma~\ref{ana:one-one-density-duality}(2) applied to the linear functional \(a_T(x)\), and the formula for \(a_T(x)(\beta^{n-1})\) is
Lemma~\ref{ana:mixed-trace-infimum}(2). For almost every \(x\), \(a_T(x)\ge0\) on the strongly positive forms (Proposition~\ref{bt:lebesgue-decomposition}), and then
Lemma~\ref{ana:one-one-density-duality}(3) gives \(A_T(x)\ge0\) and \(D_C(T)(x)=\frac{n!}{\pi^n}\det(A_T(x))^{1/n}\). The forms of \(C\) are strongly positive. Discarding the
countably many null sets for \(T\) and the \(T_k\), Lemma~\ref{ana:ac-density-liminf} gives
\(\frac{n!}{\pi^n}\det(A_T(x))^{1/n}=D_C(T)(x)\ge\liminf_kD_C(T_k)(x)=\frac{n!}{\pi^n}\liminf_k\det(A_{T_k}(x))^{1/n}\) for almost every \(x\in U\).
\end{proof}
This is Boucksom's Lemma~3.5 (arXiv numbering) in the \(\liminf\) form its proof gives, for the coefficients of a positive \((1,1)\)-current read in the
\((n-1,n-1)\)-bidimension convention of Definition~\ref{bt:constant-form-currents}; the density matrix is relative to the standard form, i.e.\(\det(A_T)\) is \(\det(T_{\mathrm{ac}})\) up to the factor fixed by the normalization of Lemma~\ref{ana:mixed-trace-infimum}.

\begin{lemma}[The density form of \(T_{\mathrm{ac}}\) is unique and measurable]
\label{ana:ac-density-form-measurable}
\uses{ana:ac-determinant-semicontinuity,ana:one-one-density-duality,bt:lebesgue-decomposition,bt:order-zero}
\source{bou02-volume:1.3}
\lean{TranscendentalNakamaye.Pending.LocalCurrentDensityMatrix.eq_sum_coordinateVector,TranscendentalNakamaye.Pending.LocalCurrentDensityMatrix.eq_zero_of_normalizedHermitianMatrix_eq_zero,TranscendentalNakamaye.Pending.LocalCurrentDensityMatrix.isOneOne_sub,TranscendentalNakamaye.Pending.LocalCurrentDensityMatrix.isOneOne_add,TranscendentalNakamaye.Pending.LocalCurrentDensityMatrix.isOneOne_smul,TranscendentalNakamaye.Pending.LocalCurrentDensityMatrix.isOneOne_zero,TranscendentalNakamaye.Pending.LocalCurrentDensityMatrix.isOneOne_sum,TranscendentalNakamaye.Pending.LocalCurrentDensityMatrix.wedgeDensity,TranscendentalNakamaye.Pending.LocalCurrentDensityMatrix.wedgeDensity_apply,TranscendentalNakamaye.Pending.LocalCurrentDensityMatrix.eq_of_forall_topDensity_eq,TranscendentalNakamaye.Pending.LocalCurrentDensityMatrix.acForm_eq_of_forall,TranscendentalNakamaye.Pending.LocalCurrentDensityMatrix.dualForm,TranscendentalNakamaye.Pending.LocalCurrentDensityMatrix.isOneOne_dualForm,TranscendentalNakamaye.Pending.LocalCurrentDensityMatrix.topDensity_dualForm,TranscendentalNakamaye.Pending.LocalCurrentDensityMatrix.acForm_eq_sum,TranscendentalNakamaye.Pending.LocalCurrentDensityMatrix.measurable_acForm_apply,TranscendentalNakamaye.Pending.LocalCurrentDensityMatrix.measurable_acMatrix_apply,TranscendentalNakamaye.Pending.LocalCurrentDensityMatrix.measurable_det_acMatrix,TranscendentalNakamaye.Pending.LocalCurrentDensityMatrix.measurable_rpow_det_acMatrix}
\leanok
Let \(n=m+1\ge1\) and write \(\langle\eta\rangle\), \(H_\gamma\) and \(\Lambda_{n-1}\) as in Lemma~\ref{ana:one-one-density-duality}.
\begin{enumerate}
\item A real \((1,1)\)-form \(\alpha\) on \(\mathbb C^n\) with \(H_\alpha=0\) is zero. If real \((1,1)\)-forms \(\alpha,\alpha'\) satisfy
 \(\langle\gamma\wedge\alpha\rangle=\langle\gamma\wedge\alpha'\rangle\) for every \(\gamma\in\Lambda_{n-1}\), then \(\alpha=\alpha'\).
\item Let \(T\) be a positive current of bidimension \((n-1,n-1)\) on an open set \(U\subset\mathbb C^n\), with \(a_T(x)=h(x)L_x\) as in
 Lemma~\ref{ana:ac-density-liminf}. For every \(x\) the real \((1,1)\)-form \(\alpha_x\) with \(a_T(x)=\langle\cdot\wedge\alpha_x\rangle\) on \(\Lambda_{n-1}\)
 (Lemma~\ref{ana:ac-determinant-semicontinuity}) is unique. If \(b_1,\dots,b_N\) is a basis of \(\Lambda_{n-1}\) and \(\alpha^{(i)}\) is the real \((1,1)\)-form
 with \(\langle\gamma\wedge\alpha^{(i)}\rangle\) the \(i\)-th coordinate of \(\gamma\), then \(\alpha_x=\sum_i a_T(x)(b_i)\,\alpha^{(i)}\).
\item Consequently \(x\mapsto\alpha_x(v,w)\) for fixed vectors \(v,w\), the entries of the density matrix \(A_T(x)=H_{\alpha_x}\), \(x\mapsto\det A_T(x)\) and
 \(x\mapsto(\operatorname{Re}\det A_T(x))^{1/n}\) are measurable functions on \(\mathbb C^n\).
\end{enumerate}
\end{lemma}
\begin{proof}
\leanok
(1) For a form \(\alpha\) of type \((1,1)\) the Hermitian entry \(G(v,w)=\alpha(v,iw)-i\,\alpha(v,w)\) is \(\mathbb C\)-linear in \(v\) and conjugate linear in \(w\), and
\((H_\alpha)_{jk}=\pi\,G(e_j,e_k)\). So \(H_\alpha=0\) gives \(G=0\) on the standard basis, hence everywhere, and \(\alpha(v,w)=-\operatorname{Im}G(v,w)=0\). For the
second claim, the zero functional on \(\Lambda_{n-1}\) is represented by the \((1,1)\)-forms \(\alpha-\alpha'\) and \(\alpha'-\alpha\); by
Lemma~\ref{ana:one-one-density-duality}(3) both \(H_{\alpha-\alpha'}\) and \(H_{\alpha'-\alpha}=-H_{\alpha-\alpha'}\) are Hermitian nonnegative, so
\(H_{\alpha-\alpha'}=0\) and \(\alpha=\alpha'\) by the first claim.

(2) Uniqueness is (1). The forms \(\alpha^{(i)}\) exist by Lemma~\ref{ana:one-one-density-duality}(2). The form \(\sum_ia_T(x)(b_i)\,\alpha^{(i)}\) is of type \((1,1)\),
and since \(\alpha\mapsto\langle\gamma\wedge\alpha\rangle\) is linear its pairing with \(\gamma=\sum_i\gamma_ib_i\) is \(\sum_ia_T(x)(b_i)\gamma_i=a_T(x)(\gamma)\); by
uniqueness it is \(\alpha_x\).

(3) The coefficients \(a_T(x)(b_i)=h(x)f_{b_i}(x)\) are measurable: \(h=d\sigma_T/d\lambda\) (Proposition~\ref{bt:lebesgue-decomposition}) and the coefficient densities
\(f_{b_i}\) (Proposition~\ref{bt:order-zero}) are measurable. By (2), \(\alpha_x(v,w)\) is a finite linear combination of them with constant coefficients; the entries of
\(H_{\alpha_x}\) are real linear combinations of such values, and the determinant is a polynomial in the entries.
\end{proof}

\begin{lemma}[Uniqueness of the absolutely continuous density]
\label{ana:ac-density-unique}
\uses{bt:lebesgue-decomposition,bt:order-zero,bt:recursive-product-local}
\source{bou02-volume:4.2}
\lean{TranscendentalNakamaye.Pending.LocalCurrentPullback.coeffMeasure,TranscendentalNakamaye.Pending.LocalCurrentPullback.measurable_ofReal_densityForm,TranscendentalNakamaye.Pending.LocalCurrentPullback.coeffMeasure_compl,TranscendentalNakamaye.Pending.LocalCurrentPullback.coeffMeasure_lt_top,TranscendentalNakamaye.Pending.LocalCurrentPullback.integral_coeffMeasure,TranscendentalNakamaye.Pending.LocalCurrentPullback.coeffMeasure_eq,TranscendentalNakamaye.Pending.LocalCurrentPullback.ae_eq_ofReal_acDensity}
\leanok
Let \(U\subset\mathbb C^n\) be open and \(T\) a positive current of bidimension \((q,q)\) on \(U\), with trace measure \(\sigma_T\), density form \(L_x\)
and \(a_T(x)=h(x)L_x\), \(h=d\sigma_T/d\lambda\), as in Proposition~\ref{bt:lebesgue-decomposition}. Let \(\gamma\) be a strongly positive constant
\((q,q)\)-form and \(\mu\) a measure on \(\mathbb C^n\) with \(\mu(\mathbb C^n\setminus U)=0\), finite on the compact subsets of \(U\), such that
\(T(\gamma)(\varphi)=\int\varphi\,d\mu\) for every test function \(\varphi\) on \(U\). If \(\mu=s+f\lambda\) with \(s\perp\lambda\) and \(f\colon\mathbb C^n\to[0,\infty]\)
measurable, then \(f(x)=a_T(x)(\gamma)\) for Lebesgue-almost every \(x\) (recall \(a_T(x)(\gamma)\ge0\) almost everywhere,
Proposition~\ref{bt:lebesgue-decomposition}).
\end{lemma}
\begin{proof}
\leanok
Since \(\gamma\) is strongly positive, \(L_x(\gamma)\ge0\) for \(\sigma_T\)-almost every \(x\) (Proposition~\ref{bt:lebesgue-decomposition}), so the measure
\(\nu=L(\gamma)\,\sigma_T\) represents \(T(\gamma)\) (Proposition~\ref{bt:order-zero}); it is carried by \(U\) and finite on the compact subsets of \(U\), because
\(L(\gamma)\) is bounded. Two such measures with the same integrals of all test functions on \(U\) are equal (the uniqueness step of
Lemma~\ref{bt:recursive-product-local}), so \(\mu=\nu\). Splitting \(\sigma_T=\sigma_T^{\mathrm{sing}}+h\lambda\) gives
\(\nu=L(\gamma)\sigma_T^{\mathrm{sing}}+h\,L(\gamma)\lambda\), whose first term is singular to \(\lambda\). By the uniqueness of the Lebesgue decomposition, \(f\) and
\(hL(\gamma)\) both agree almost everywhere with the Radon--Nikodym derivative \(d\mu/d\lambda\).
\end{proof}

\begin{lemma}[The pull-back of a positive current along a biholomorphism]
\label{ana:current-pullback}
\uses{bt:constant-form-currents,bt:order-zero,bt:lebesgue-decomposition}
\source{bou02-volume:4.2}
\lean{TranscendentalNakamaye.Pending.LocalCurrentPullback.compContinuousLinearMap_mem_formSpan,TranscendentalNakamaye.Pending.LocalCurrentPullback.pullSpan,TranscendentalNakamaye.Pending.LocalCurrentPullback.coe_pullSpan,TranscendentalNakamaye.Pending.LocalCurrentPullback.pullSpan_mem_posCone,TranscendentalNakamaye.Pending.LocalCurrentPullback.pullSpan_pullSpan,TranscendentalNakamaye.Pending.LocalCurrentPullback.pullSpan_id,TranscendentalNakamaye.Pending.LocalCurrentPullback.continuous_pullSpan_apply,TranscendentalNakamaye.Pending.LocalCurrentPullback.continuous_repr_pullSpan_apply,TranscendentalNakamaye.Pending.LocalCurrentPullback.measurable_densityForm_pullSpan,TranscendentalNakamaye.Pending.LocalCurrentPullback.onOpen,TranscendentalNakamaye.Pending.LocalCurrentPullback.onOpen_of_mem,TranscendentalNakamaye.Pending.LocalCurrentPullback.measurable_onOpen,TranscendentalNakamaye.Pending.LocalCurrentPullback.pbMeasure,TranscendentalNakamaye.Pending.LocalCurrentPullback.pbDensity,TranscendentalNakamaye.Pending.LocalCurrentPullback.measurable_pbDensity,TranscendentalNakamaye.Pending.LocalCurrentPullback.pbDensity_apply_of_mem,TranscendentalNakamaye.Pending.LocalCurrentPullback.pbDensity_comp,TranscendentalNakamaye.Pending.LocalCurrentPullback.pbDensity_add,TranscendentalNakamaye.Pending.LocalCurrentPullback.pbDensity_smul,TranscendentalNakamaye.Pending.LocalCurrentPullback.measurable_onOpen_inv,TranscendentalNakamaye.Pending.LocalCurrentPullback.pbMeasure_apply,TranscendentalNakamaye.Pending.LocalCurrentPullback.preimage_onOpen_inter_subset,TranscendentalNakamaye.Pending.LocalCurrentPullback.pbMeasure_compl,TranscendentalNakamaye.Pending.LocalCurrentPullback.pbMeasure_lt_top,TranscendentalNakamaye.Pending.LocalCurrentPullback.integral_pbMeasure,TranscendentalNakamaye.Pending.LocalCurrentPullback.integrable_pbMeasure_iff,TranscendentalNakamaye.Pending.LocalCurrentPullback.exists_abs_repr_pullSpan_le,TranscendentalNakamaye.Pending.LocalCurrentPullback.exists_abs_pbDensity_le,TranscendentalNakamaye.Pending.LocalCurrentPullback.locallyIntegrableOn_pbDensity,TranscendentalNakamaye.Pending.LocalCurrentPullback.integrable_mul_pbDensity,TranscendentalNakamaye.Pending.LocalCurrentPullback.pullback,TranscendentalNakamaye.Pending.LocalCurrentPullback.pullback_apply_eq,TranscendentalNakamaye.Pending.LocalCurrentPullback.pullback_apply,TranscendentalNakamaye.Pending.LocalCurrentPullback.integrableOn_pullback_integrand,TranscendentalNakamaye.Pending.LocalCurrentPullback.isPositive_pullback}
\leanok
Let \(\tau\colon A\to B\) be a biholomorphism between open subsets of \(\mathbb C^n\), with inverse \(\sigma\), and let \(T\) be a positive current of bidimension
\((q,q)\) on \(B\), with trace measure \(\sigma_T\) and density form \(L_y\) (Proposition~\ref{bt:lebesgue-decomposition}). For a \(\mathbb C\)-linear map \(S\) of
\(\mathbb C^n\), the pull-back \(S^*\gamma\) of a constant form \(\gamma\in\Lambda_q\) lies in \(\Lambda_q\), and \(S^*\) maps strongly positive forms to strongly
positive forms. The formula
\[
 (\tau^*T)(\gamma)(\varphi)=\int\varphi(x)\,L_{\tau(x)}\bigl(D\sigma(\tau(x))^*\gamma\bigr)\,d(\sigma_*\sigma_T|_B)(x)
 =\int_B\varphi(\sigma(y))\,L_y\bigl(D\sigma(y)^*\gamma\bigr)\,d\sigma_T(y)
\]
for \(\gamma\in\Lambda_q\) and test functions \(\varphi\) on \(A\) defines a current \(\tau^*T\) of bidimension \((q,q)\) on \(A\), the \emph{pull-back} of \(T\):
\(\langle\tau^*T,\varphi\gamma\rangle=\langle T,\sigma^*(\varphi\gamma)\rangle\), the current \(T\) being evaluated on the variable form
\(\sigma^*(\varphi\gamma)=\varphi\circ\sigma\cdot D\sigma^*\gamma\) through its order-zero representation. The pull-back \(\tau^*T\) is positive.
\end{lemma}
\begin{proof}
\leanok
A \(\mathbb C\)-linear \(S\) commutes with \(J=i\), so \(S^*\) maps products of semipositive \((1,1)\)-forms to such products, hence the strongly positive cone into
itself and its span \(\Lambda_q\) into itself. The measure \(\sigma_*\sigma_T|_B\) is carried by \(A\) and finite on compact subsets \(K\subset A\), since
\(\sigma^{-1}(K)\cap B=\tau(K)\) is compact in \(B\). The density \(x\mapsto L_{\tau(x)}(D\sigma(\tau(x))^*\gamma)\) is measurable and bounded on compact subsets
of \(A\): \(L_y\) is bounded in \(y\) on each basis form, and the coefficients of \(D\sigma(y)^*\gamma\) are continuous on \(B\). So the formula is integration of
test functions against a locally finite measure with a locally bounded density, a distribution, and it is linear in \(\gamma\); the second expression is the change
of variables \(x=\sigma(y)\). Positivity: if \(\sum_k\varphi_k(x)\gamma_k\) is strongly positive for every \(x\), then so is
\(D\sigma(y)^*\bigl(\sum_k\varphi_k(\sigma(y))\gamma_k\bigr)\), and \(L_y\) is nonnegative on strongly positive forms for \(\sigma_T\)-almost every \(y\); hence
\(\sum_k(\tau^*T)(\gamma_k)(\varphi_k)=\int_BL_y\bigl(D\sigma(y)^*\sum_k\varphi_k(\sigma(y))\gamma_k\bigr)\,d\sigma_T(y)\ge0\).
\end{proof}

\begin{lemma}[Boucksom: \((\tau^*T)_{\mathrm{ac}}=\tau^*(T_{\mathrm{ac}})\); the chart invariance of \(\det(A_T)/\det(H_\omega)\)]
\label{ana:current-pullback-ac}
\uses{ana:current-pullback,ana:ac-density-unique,ana:ac-density-form-measurable,ana:ac-determinant-semicontinuity,fpma:top-power-determinant-general,bou:integration-current}
\source{bou02-volume:4.2}
\lean{TranscendentalNakamaye.Pending.LocalCurrentPullback.differentiableOn_inv,TranscendentalNakamaye.Pending.LocalCurrentPullback.differentiableAt_of_mem,TranscendentalNakamaye.Pending.LocalCurrentPullback.differentiableAt_inv_of_mem,TranscendentalNakamaye.Pending.LocalCurrentPullback.ae_comp_of_ae,TranscendentalNakamaye.Pending.LocalCurrentPullback.preimage_onOpen_inter_eq,TranscendentalNakamaye.Pending.LocalCurrentPullback.map_withDensity_restrict,TranscendentalNakamaye.Pending.LocalCurrentPullback.mutuallySingular_map,TranscendentalNakamaye.Pending.LocalCurrentPullback.ae_pbDensity_nonneg,TranscendentalNakamaye.Pending.LocalCurrentPullback.measurable_ofReal_pbDensity,TranscendentalNakamaye.Pending.LocalCurrentPullback.pbCoeffMeasure,TranscendentalNakamaye.Pending.LocalCurrentPullback.integral_pbCoeffMeasure,TranscendentalNakamaye.Pending.LocalCurrentPullback.pbCoeffMeasure_compl,TranscendentalNakamaye.Pending.LocalCurrentPullback.pbCoeffMeasure_lt_top,TranscendentalNakamaye.Pending.LocalCurrentPullback.pbJacobianWeight,TranscendentalNakamaye.Pending.LocalCurrentPullback.measurable_pbJacobianWeight,TranscendentalNakamaye.Pending.LocalCurrentPullback.pbMeasure_eq,TranscendentalNakamaye.Pending.LocalCurrentPullback.pbCoeffMeasure_eq,TranscendentalNakamaye.Pending.LocalCurrentPullback.ae_acDensity_pullback_of_mem_posCone,TranscendentalNakamaye.Pending.LocalCurrentPullback.ae_acDensity_pullback,TranscendentalNakamaye.Pending.LocalCurrentPullback.fderiv_inv_comp,TranscendentalNakamaye.Pending.LocalCurrentPullback.fderiv_real_eq,TranscendentalNakamaye.Pending.LocalCurrentPullback.det_restrictScalars_eq_normSq,TranscendentalNakamaye.Pending.LocalCurrentPullback.det_fderiv_ne_zero,TranscendentalNakamaye.Pending.LocalCurrentPullback.isOneOne_comp_restrictScalars,TranscendentalNakamaye.Pending.LocalCurrentPullback.compContinuousLinearMap_restrictScalars_eq,TranscendentalNakamaye.Pending.LocalCurrentPullback.topDensity_coordinateWedge_comp,TranscendentalNakamaye.Pending.LocalCurrentPullback.ae_acForm_pullback,TranscendentalNakamaye.Pending.LocalCurrentPullback.det_transpose_mul_mul_conj,TranscendentalNakamaye.Pending.LocalCurrentPullback.ae_acMatrix_pullback,TranscendentalNakamaye.Pending.LocalCurrentPullback.ae_det_acMatrix_pullback,TranscendentalNakamaye.Pending.LocalCurrentPullback.ae_det_acMatrix_div_eq}
\leanok
Let \(\tau\colon A\to B\), \(\sigma\) and \(T\) be as in Lemma~\ref{ana:current-pullback}, and write \(a_T(y)=h_T(y)L_y\) as in Lemma~\ref{ana:ac-density-liminf}.
\begin{enumerate}
\item For Lebesgue-almost every \(x\in A\) and every \(\gamma\in\Lambda_q\),
 \(a_{\tau^*T}(x)(\gamma)=\lvert\det_{\mathbb R}D\tau(x)\rvert\,a_T(\tau(x))\bigl(D\sigma(\tau(x))^*\gamma\bigr)\).
\item Let \(q=n-1\) and let \(\alpha_x(T)\) be the density form of \(T_{\mathrm{ac}}\) (Lemma~\ref{ana:ac-density-form-measurable}). For almost every \(x\in A\),
 \(\alpha_x(\tau^*T)=D\tau(x)^*\alpha_{\tau(x)}(T)\); consequently \(A_{\tau^*T}(x)=J(x)^{\mathsf T}A_T(\tau(x))\overline{J(x)}\) and
 \(\det A_{\tau^*T}(x)=\lvert\det J(x)\rvert^2\det A_T(\tau(x))\), where \(J(x)\) is the complex Jacobian matrix of \(\tau\) at \(x\).
\item If \(\omega_A\), \(\omega_B\) assign real two-forms to the points of \(A\) and \(B\), with \(\omega_B\) of type \((1,1)\) and
 \(\omega_A(x)=D\tau(x)^*\omega_B(\tau(x))\) on \(A\), then for almost every \(x\in A\)
 \[
  \frac{\det A_{\tau^*T}(x)}{\det H_{\omega_A(x)}}=\frac{\det A_T(\tau(x))}{\det H_{\omega_B(\tau(x))}} .
 \]
\end{enumerate}
\end{lemma}
\begin{proof}
\leanok
(1) Two facts about \(\tau\). Since \(\sigma\) is \(C^1\) on \(B\), it maps Lebesgue-null subsets of \(B\) to null sets, so a property holding at almost every point of
\(B\) holds at \(\tau(x)\) for almost every \(x\in A\); and \(\sigma_*\) maps measures singular to \(\lambda\), restricted to \(B\), to measures singular to \(\lambda\). By the
change of variables formula, \(\sigma_*\bigl((r\lambda)|_B\bigr)=\lvert\det_{\mathbb R}D\tau\rvert\,(r\circ\tau)\,\lambda|_A\) for every measurable \(r\ge0\).
Fix a strongly positive \(\gamma\). The measure \(\mu=L_{\tau(\cdot)}\bigl(D\sigma(\tau(\cdot))^*\gamma\bigr)\,\sigma_*\sigma_T|_B\) has a nonnegative density
(\(D\sigma^*\) preserves strong positivity, Lemma~\ref{ana:current-pullback}), represents \((\tau^*T)(\gamma)\), is carried by \(A\) and is finite on its compact subsets.
Splitting \(\sigma_T=\sigma_T^{\mathrm{sing}}+h_T\lambda\) writes \(\mu=s+\lvert\det_{\mathbb R}D\tau\rvert\,h_T(\tau(\cdot))\,L_{\tau(\cdot)}(D\sigma(\tau(\cdot))^*\gamma)\,\lambda|_A\) with
\(s\perp\lambda\), and Lemma~\ref{ana:ac-density-unique} gives (1) for \(\gamma\) (using \(h_T(\tau(x))<\infty\) and \(a_T(\tau(x))\ge0\) on strongly positive forms for
almost every \(x\in A\), by the first fact). Every \(\gamma\in\Lambda_q\) is the difference of the strongly positive forms \(C\omega_0^q+\gamma\) and \(C\omega_0^q\) for some
\(C\ge0\) (Definition~\ref{bt:constant-form-currents}), and both sides of (1) are linear in \(\gamma\); applying this to a basis of \(\Lambda_q\) gives (1) for all \(\gamma\)
outside one null set.

(2) Put \(L=D\tau(x)\), \(M=D\sigma(\tau(x))\), so \(M\circ L=\mathrm{id}\). For \(\gamma\in\Lambda_{n-1}\) and \(\alpha=\alpha_{\tau(x)}(T)\),
\(\langle\gamma\wedge L^*\alpha\rangle=\langle L^*(M^*\gamma\wedge\alpha)\rangle=\det_{\mathbb R}L\,\langle M^*\gamma\wedge\alpha\rangle\) (Lemma~\ref{bou:integration-current}), and
\(\det_{\mathbb R}L=\lvert\det_{\mathbb C}L\rvert^2\ge0\). By (1) the \((1,1)\)-form \(L^*\alpha\) represents \(a_{\tau^*T}(x)\), so it is \(\alpha_x(\tau^*T)\) by the uniqueness
in Lemma~\ref{ana:ac-density-form-measurable}. The coefficient matrix of the pull-back of a \((1,1)\)-form along a \(\mathbb C\)-linear map with matrix \(J\) is
\(J^{\mathsf T}H\overline J\) (Lemma~\ref{fpma:top-power-determinant-general}), and \(\det(J^{\mathsf T}H\overline J)=\lvert\det J\rvert^2\det H\).

(3) The same transformation law holds for \(\omega_A(x)=L^*\omega_B(\tau(x))\), and \(\det J(x)\ne0\) because \(M\circ L=\mathrm{id}\); the factors \(\lvert\det J(x)\rvert^2\)
cancel.
\end{proof}
On a complex manifold \(X\) the chart readings of a positive \((1,1)\)-current are related by such pull-backs along the transition maps; together with the
measurability of Lemma~\ref{ana:ac-density-form-measurable}, (3) then makes \(\det(A_T)/\det(H_\omega)\), i.e.\ \(\det_\omega(T_{\mathrm{ac}})=T_{\mathrm{ac}}^n/\omega^n\), a function
on \(X\) defined almost everywhere. For the chart readings of \(\theta+dd^cu\) this needs \(\tau^*(dd^cw)=dd^c(w\circ\tau)\), which is not part of these lemmas.

\begin{lemma}[Naturality of \(dd^c\) under biholomorphisms: \(\tau^*(dd^cw)=dd^c(w\circ\tau)\)]
\label{ana:current-pullback-ddc}
\uses{ana:current-pullback,ana:current-pullback-ac,ana:ac-density-unique,bt:recursive-product-positive-closed,bt:recursive-product-smooth,bt:recursive-product-local,bt:collar-comparison,bt:continuous-convergence,hgeom:ddc-naturality,psh:holomorphic-pullback,psh:locally-bounded-decreasing-regularization}
\source{bou02-volume:4.2}
\lean{TranscendentalNakamaye.Pending.LocalCurrentPullbackDdc.ae_forall_eq_of_posCone,TranscendentalNakamaye.Pending.LocalCurrentPullbackDdc.formDensity,TranscendentalNakamaye.Pending.LocalCurrentPullbackDdc.formDensity_apply,TranscendentalNakamaye.Pending.LocalCurrentPullbackDdc.formDensity_nonneg,TranscendentalNakamaye.Pending.LocalCurrentPullbackDdc.ddcMul_oneCurrent_apply,TranscendentalNakamaye.Pending.LocalCurrentPullbackDdc.continuousOn_formDensity,TranscendentalNakamaye.Pending.LocalCurrentPullbackDdc.formWeight,TranscendentalNakamaye.Pending.LocalCurrentPullbackDdc.measurable_formWeight,TranscendentalNakamaye.Pending.LocalCurrentPullbackDdc.formWeight_compl,TranscendentalNakamaye.Pending.LocalCurrentPullbackDdc.formWeight_lt_top,TranscendentalNakamaye.Pending.LocalCurrentPullbackDdc.integral_formWeight,TranscendentalNakamaye.Pending.LocalCurrentPullbackDdc.pullback_apply_of_form,TranscendentalNakamaye.Pending.LocalCurrentPullbackDdc.isLocallyBoundedPshOn_of_isSmoothPshOn,TranscendentalNakamaye.Pending.LocalCurrentPullbackDdc.isPositive_ddcMul_oneCurrent,TranscendentalNakamaye.Pending.LocalCurrentPullbackDdc.pullback_ddcMul_of_isSmoothPshOn,TranscendentalNakamaye.Pending.LocalCurrentPullbackDdc.preOpen,TranscendentalNakamaye.Pending.LocalCurrentPullbackDdc.preOpen_le,TranscendentalNakamaye.Pending.LocalCurrentPullbackDdc.mem_preOpen,TranscendentalNakamaye.Pending.LocalCurrentPullbackDdc.isBiholomorphicPair_preOpen,TranscendentalNakamaye.Pending.LocalCurrentPullbackDdc.ae_densityForm_restrict_eq,TranscendentalNakamaye.Pending.LocalCurrentPullbackDdc.pullback_restrict,TranscendentalNakamaye.Pending.LocalCurrentPullbackDdc.exists_testFunction_eq_one,TranscendentalNakamaye.Pending.LocalCurrentPullbackDdc.pbCoeffFun,TranscendentalNakamaye.Pending.LocalCurrentPullbackDdc.continuousOn_pbCoeffFun,TranscendentalNakamaye.Pending.LocalCurrentPullbackDdc.measurable_pbCoeffFun,TranscendentalNakamaye.Pending.LocalCurrentPullbackDdc.isLocallyBoundedOn_pbCoeffFun,TranscendentalNakamaye.Pending.LocalCurrentPullbackDdc.pullback_apply_eq_sum_mulBorel,TranscendentalNakamaye.Pending.LocalCurrentPullbackDdc.tendstoLocallyUniformlyOn_const_seq,TranscendentalNakamaye.Pending.LocalCurrentPullbackDdc.weakTendsto_pullback,TranscendentalNakamaye.Pending.LocalCurrentPullbackDdc.restrict_oneCurrent,TranscendentalNakamaye.Pending.LocalCurrentPullbackDdc.weakTendsto_ddcMul_oneCurrent,TranscendentalNakamaye.Pending.LocalCurrentPullbackDdc.isLocallyBoundedPshOn_comp,TranscendentalNakamaye.Pending.LocalCurrentPullbackDdc.pullback_congr,TranscendentalNakamaye.Pending.LocalCurrentPullbackDdc.pullback_ddcMul}
\leanok
Let \(\tau\colon A\to B\) be a biholomorphism between open subsets of \(\mathbb C^n\), with inverse \(\sigma\), and write \(dd^cw\) for the positive closed current
\(dd^c(w[1])\) of bidimension \((n-1,n-1)\) of a locally bounded psh function \(w\) (Proposition~\ref{bt:recursive-product-positive-closed}).
\begin{enumerate}
\item For a function \(g\) smooth on an open set \(U\), \((dd^cg)(\gamma)(\varphi)=\int\varphi\,\langle\gamma\wedge dd^cg\rangle\,d\lambda\).
\item If \(T\) is a positive current on \(B\) with \(T(\gamma)(\varphi)=\int\varphi\,\langle\gamma\wedge\alpha\rangle\,d\lambda\) for a continuous form \(\alpha\) that is
 semipositive of type \((1,1)\) on \(B\), then \((\tau^*T)(\gamma)(\varphi)=\int\varphi(x)\,\langle\gamma\wedge D\tau(x)^*\alpha(\tau(x))\rangle\,dx\).
\item The pull-back of Lemma~\ref{ana:current-pullback} commutes with restriction to open subsets \(V\subset B\) (with \(A\cap\tau^{-1}(V)\) on the source side),
 and it is weakly continuous: if positive currents \(T_k\to T\) weakly on \(B\), then \(\tau^*T_k\to\tau^*T\) weakly on \(A\).
\item If locally bounded psh functions \(f_j\) decrease pointwise to a locally bounded psh \(w\) on \(U\), then \(dd^cf_j\to dd^cw\) weakly.
\item For every locally bounded psh function \(w\) on \(B\), \(\tau^*(dd^cw)=dd^c(w\circ\tau)\) on \(A\).
\end{enumerate}
\end{lemma}
\begin{proof}
\leanok
(1) For \(g\) smooth, \(dd^c(g[1])=dd^cg\wedge[1]\) (Proposition~\ref{bt:recursive-product-smooth}), i.e.
\((dd^cg)(\gamma)(\varphi)=-\frac1{4\pi}\sum_{a,b}\langle\gamma\wedge P_{ab}\rangle\int\varphi\,\partial_a\partial_bg\,d\lambda\), and
\(dd^cg=-\frac1{4\pi}\sum_{a,b}\partial_a\partial_bg\,P_{ab}\) on \(U\) (the coordinate formula used in Lemma~\ref{bt:collar-comparison}).

(2) For a strongly positive \(\gamma\) the density \(\langle\gamma\wedge\alpha\rangle\) is nonnegative. Taking \(\gamma=\omega_0^{n-1}\), the trace measure of \(T\) is
\(\langle\omega_0^{n-1}\wedge\alpha\rangle\lambda|_B\) (both measures are carried by \(B\), finite on its compact subsets and have the same integrals of test functions), so
\(\sigma_T\ll\lambda\). For strongly positive \(\gamma\), Lemma~\ref{ana:ac-density-unique} (with \(s=0\)) gives \(a_T(y)(\gamma)=\langle\gamma\wedge\alpha(y)\rangle\) for almost every
\(y\in B\); by linearity, writing every form of \(\Lambda_{n-1}\) as a difference of two strongly positive ones on a basis, this holds for all \(\gamma\) at once almost
everywhere. Then, by Lemma~\ref{ana:current-pullback}, \((\tau^*T)(\gamma)(\varphi)=\int_B\varphi(\sigma(y))\,\langle D\sigma(y)^*\gamma\wedge\alpha(y)\rangle\,dy\), and the change of
variables \(y=\tau(x)\) with the identity \(\lvert\det_{\mathbb R}D\tau\rvert\langle M^*\gamma\wedge\alpha\rangle=\langle\gamma\wedge L^*\alpha\rangle\) of the proof of
Lemma~\ref{ana:current-pullback-ac} gives the formula.

(3) Restriction: on \(V\) the trace measure and the density form of \(T|_V\) are those of \(T\) (Proposition~\ref{bt:recursive-product-local}), and the integrand
\(\varphi(\sigma(y))L_y(D\sigma(y)^*\gamma)\) vanishes for \(y\in B\setminus V\) when \(\varphi\) is supported in \(A\cap\tau^{-1}(V)\). Continuity: let \(\chi\) be a test function on \(B\)
equal to \(1\) on the compact set \(\tau(\operatorname{supp}\varphi)\), and \(a_i(y)=\varphi(\sigma(y))\,(D\sigma(y)^*\gamma)_i\), continuous on \(B\), the coordinates being taken in a basis
\((b_i)\) of \(\Lambda_q\). Then \(\langle\tau^*T,\varphi\gamma\rangle=\sum_i(a_iT)(b_i)(\chi)\), and \((a_iT_k)(b_i)(\chi)\to(a_iT)(b_i)(\chi)\) by the continuity of multiplication by a
continuous weight under weak convergence (Theorem~\ref{bt:continuous-convergence}).

(4) \(dd^c\) is weakly continuous (Theorem~\ref{bt:continuous-convergence}), and \((f_j[1])(\gamma)(\varphi)=\langle\gamma\rangle\int\varphi f_j\,d\lambda\to\langle\gamma\rangle\int\varphi w\,d\lambda\) by
dominated convergence: on \(\operatorname{supp}\varphi\), \(w\le f_j\le f_0\) and both bounds are bounded.

(5) Fix \(\gamma\) and \(\varphi\). Choose an open \(V\) with compact closure in \(B\) containing \(\tau(\operatorname{supp}\varphi)\), and smooth psh \(f_j\) on \(V\) decreasing to \(w\)
(Corollary~\ref{psh:locally-bounded-decreasing-regularization}). On \(A'=A\cap\tau^{-1}(V)\), which contains \(\operatorname{supp}\varphi\): by (2) with \(\alpha=dd^cf_j\) (semipositive, by (1)) and
the naturality \(dd^c(f_j\circ\tau)=D\tau^*(dd^cf_j\circ\tau)\) of the smooth \(dd^c\) (Lemma~\ref{hgeom:ddc-naturality}), \(\tau^*(dd^cf_j)=dd^c(f_j\circ\tau)\). By (4) and (3), the left
side tends to \(\tau^*(dd^cw|_V)\) and the right side to \(dd^c(w\circ\tau)\) on \(A'\) (\(f_j\circ\tau\) and \(w\circ\tau\) are locally bounded psh, Theorem~\ref{psh:holomorphic-pullback}).
Restriction (3) and the locality of \(dd^c(w[1])\) (Proposition~\ref{bt:recursive-product-local}) identify both sides on \(\varphi\gamma\) with the values on \(A\).
\end{proof}
With Lemma~\ref{ana:current-pullback-ac}, this shows that on a complex manifold the chart readings of \(\theta+dd^cu\), for \(u\) locally bounded and
\(\theta=dd^c\rho\) locally, transform by pull-back along the chart transitions, so that \(\det(A_T)/\det(H_\omega)\) is one function on the manifold, defined almost everywhere.
For an unbounded \(u\) the current \(dd^c(u[1])\) of Proposition~\ref{bt:recursive-product-positive-closed} is not defined; Lemma~\ref{ana:l1-psh-current} and
Proposition~\ref{ana:l1-psh-ddc} below give \(dd^c(u[1])\) for every psh \(u\in L^1_{\mathrm{loc}}\), with the same naturality.

\begin{lemma}[The current \(u[1]\) of a psh function \(u\in L^1_{\mathrm{loc}}\)]
\label{ana:l1-psh-current}
\uses{psh:local-integrability,bt:constant-form-currents,bt:ddc-currents,bt:recursive-product-local}
\source{demailly-cadg:I.4.17,demailly-cadg:III.1.18}
\lean{TranscendentalNakamaye.Pending.LocalCurrentL1Ddc.IsL1PshOn,TranscendentalNakamaye.Pending.LocalCurrentL1Ddc.IsL1PshOn.subset_pshIntegrabilityLocus,TranscendentalNakamaye.Pending.LocalCurrentL1Ddc.IsL1PshOn.locallyIntegrableOn,TranscendentalNakamaye.Pending.LocalCurrentL1Ddc.IsL1PshOn.ae_ne_bot,TranscendentalNakamaye.Pending.LocalCurrentL1Ddc.IsL1PshOn.mono,TranscendentalNakamaye.Pending.LocalCurrentL1Ddc.l1Current,TranscendentalNakamaye.Pending.LocalCurrentL1Ddc.l1Current_apply,TranscendentalNakamaye.Pending.LocalCurrentL1Ddc.l1Current_congr_ae,TranscendentalNakamaye.Pending.LocalCurrentL1Ddc.restrict_l1Current,TranscendentalNakamaye.Pending.LocalCurrentL1Ddc.integrable_testFunction_mul,TranscendentalNakamaye.Pending.LocalCurrentL1Ddc.weakTendsto_l1Current,TranscendentalNakamaye.Pending.LocalCurrentL1Ddc.ddcL1,TranscendentalNakamaye.Pending.LocalCurrentL1Ddc.restrict_ddcL1}
\leanok
Let \(U\subset\mathbb C^n\) be open and \(u\colon\mathbb C^n\to[-\infty,\infty]\) psh on \(U\) such that every point of \(U\) has points \(y\) with \(u(y)\ne-\infty\) arbitrarily
close to it (\(u\) is not identically \(-\infty\) near any point of \(U\)). Then \(u\in L^1_{\mathrm{loc}}(U)\) and \(u>-\infty\) almost everywhere on \(U\); the same holds on every open
\(V\subset U\). For such \(u\) put
\[
 (u[1])(\gamma)(\varphi)=\langle\gamma\rangle\int\varphi\,u\,d\lambda\qquad(\gamma\in\Lambda_n,\ \varphi\in\mathcal D(U)),
\]
a current of bidimension \((n,n)\) on \(U\), and \(dd^c(u[1])\), of bidimension \((n-1,n-1)\) (Definition~\ref{bt:ddc-currents}). Both depend only on the class of \(u\) almost
everywhere on \(U\), and their restrictions to an open \(V\subset U\) are the currents of \(u|_V\). If \(f_k\to g\) almost everywhere on \(U\) with \(|f_k|\le F\) for some \(F\in
L^1_{\mathrm{loc}}(U)\), then \(f_k[1]\to g[1]\) weakly.
\end{lemma}
\begin{proof}
\leanok
\uses{psh:local-integrability}
Let \(x\in U\) and \(B=B(x,r)\subset U\). By hypothesis \(B\) contains a point where \(u\ne-\infty\); \(B\) is connected, so by Theorem~\ref{psh:local-integrability} every point of \(B\),
in particular \(x\), has an open neighbourhood \(V_x\) on which \(u\) is finite almost everywhere and integrable. Hence \(u\in L^1_{\mathrm{loc}}(U)\), and the set
\(N=\{x\in U:u(x)=-\infty\}\) has \(\lambda(N\cap V_x)=0\) for every \(x\in N\); since \(\mathbb C^n\) is second countable, \(\lambda(N)=0\). The hypothesis passes to open subsets.
The current is the regular distribution of \(u\) times the top-degree density \(\langle\gamma\rangle\), linear in \(\gamma\). Changing \(u\) on a null subset of \(U\) does not change
\(\int\varphi u\), and a test function on \(V\) extended by zero to \(U\) has the same integral against \(u\); \(dd^c\) commutes with restriction (Proposition~\ref{bt:recursive-product-local}).
For the convergence: \(\varphi F\) is integrable (\(F\) is integrable on the compact support of \(\varphi\)), \(|\varphi f_k|\le|\varphi||F|\) and \(\varphi f_k\to\varphi g\) almost everywhere,
so \(\int\varphi f_k\to\int\varphi g\) by dominated convergence.
\end{proof}

\begin{proposition}[\(dd^c(u[1])\) for psh \(u\in L^1_{\mathrm{loc}}\): positivity, agreement with Bedford--Taylor, naturality]
\label{ana:l1-psh-ddc}
\uses{ana:l1-psh-current,bttr:truncation,bttr:bounded-psh,bt:recursive-product-positive-closed,bt:recursive-product-forms,bt:continuous-convergence,bt:ddc-currents,ana:current-pullback,ana:current-pullback-ddc,psh:holomorphic-pullback}
\source{demailly-cadg:I.5.8,demailly-cadg:III.1.18}
\lean{TranscendentalNakamaye.Pending.LocalCurrentL1Ddc.abs_truncationReal_le,TranscendentalNakamaye.Pending.LocalCurrentL1Ddc.tendsto_truncationReal,TranscendentalNakamaye.Pending.LocalCurrentL1Ddc.IsL1PshOn.weakTendsto_l1Current_truncation,TranscendentalNakamaye.Pending.LocalCurrentL1Ddc.l1Current_coe_eq_mulBorel,TranscendentalNakamaye.Pending.LocalCurrentL1Ddc.ddcL1_coe_eq_ddcMul,TranscendentalNakamaye.Pending.LocalCurrentL1Ddc.IsL1PshOn.weakTendsto_ddcMul_truncation,TranscendentalNakamaye.Pending.LocalCurrentL1Ddc.IsL1PshOn.isPositive_ddcL1,TranscendentalNakamaye.Pending.LocalCurrentL1Ddc.isClosedAny_ddcL1,TranscendentalNakamaye.Pending.LocalCurrentL1Ddc.IsL1PshOn.comp,TranscendentalNakamaye.Pending.LocalCurrentL1Ddc.pullback_ddcL1}
\leanok
Let \(u\) be as in Lemma~\ref{ana:l1-psh-current} on an open \(U\subset\mathbb C^n\).
\begin{enumerate}
\item The truncations \(u_k=\max\{u,-k\}\) (Definition~\ref{bttr:truncation}) satisfy \(u_k[1]\to u[1]\) and \(dd^c(u_k[1])\to dd^c(u[1])\) weakly.
\item If \(w\) is locally bounded psh on \(U\), then \(w[1]=w\cdot[1]\) is the product of Proposition~\ref{bt:recursive-product-positive-closed}, so \(dd^c(w[1])\) is the current of that
 proposition.
\item \(dd^c(u[1])\) is a closed positive current.
\item If \(\tau\colon A\to B\) is a biholomorphism between open subsets of \(\mathbb C^n\) and \(u\) is as above on \(B\), then \(u\circ\tau\) is as above on \(A\), and
 \(\tau^*\bigl(dd^c(u[1])\bigr)=dd^c\bigl((u\circ\tau)[1]\bigr)\), with the pull-back of Lemma~\ref{ana:current-pullback}.
\end{enumerate}
\end{proposition}
\begin{proof}
\leanok
(1) The \(u_k\) are locally bounded psh (Lemma~\ref{bttr:bounded-psh}), hence in \(L^1_{\mathrm{loc}}\). Where \(u\) is finite, \(u_k=u\) as soon as \(k\ge-u\), and
\(|u_k|\le|u|\) (if \(u<-k\) then \(|u_k|=k<|u|\)); both hold almost everywhere on \(U\) by Lemma~\ref{ana:l1-psh-current}, so the convergence clause there applies with
\(F=|u|\), and \(dd^c\) is weakly continuous (Theorem~\ref{bt:continuous-convergence}).

(2) The coefficient \([1](\gamma)\) is the constant density \(\langle\gamma\rangle\) against \(\lambda|_U\), and the product \(w\cdot[1]\) may be computed against any such
representation of the coefficients (Proposition~\ref{bt:recursive-product-forms}); so \((w\cdot[1])(\gamma)(\varphi)=\int_U\varphi\,w\,\langle\gamma\rangle\,d\lambda=(w[1])(\gamma)(\varphi)\).

(3) By (2), \(dd^c(u_k[1])=dd^c(u_k\cdot[1])\) is positive (Proposition~\ref{bt:recursive-product-positive-closed}), and a weak limit of positive currents is positive
(Theorem~\ref{bt:continuous-convergence}); \(dd^c\) of any current is closed (Definition~\ref{bt:ddc-currents}).

(4) \(u\circ\tau\) is psh (Theorem~\ref{psh:holomorphic-pullback}). Given \(x\in A\) and a neighbourhood \(V\) of \(x\), \(\sigma^{-1}(V\cap A)\cap B\) is a neighbourhood of \(\tau(x)\)
(\(\sigma=\tau^{-1}\) is continuous and \(\sigma(\tau(x))=x\)), so it contains a point \(y\) with \(u(y)\ne-\infty\); then \(\sigma(y)\in V\) and \(u(\tau(\sigma(y)))=u(y)\ne-\infty\).
The truncations of \(u\circ\tau\) are \(u_k\circ\tau\). By Lemma~\ref{ana:current-pullback-ddc}(5) and (2), \(\tau^*(dd^c(u_k[1]))=dd^c((u_k\circ\tau)[1])\). As \(k\to\infty\) the left
side tends to \(\tau^*(dd^c(u[1]))\) weakly (by (1) and the weak continuity of \(\tau^*\), Lemma~\ref{ana:current-pullback-ddc}(3)), and the right side to \(dd^c((u\circ\tau)[1])\) (by (1)
for \(u\circ\tau\)). Weak limits are unique.
\end{proof}

\begin{lemma}[Boucksom: upper semicontinuity of the volume]
\label{ana:volume-semicontinuity}
\uses{ana:coefficient-lebesgue-decomposition,ana:geometric-setting,ana:actual-original-current-volume-sequential-upper-semicontinuity,posclass:actual-original-current-pseudoeffective-iff,ana:actual-original-current-volume-semipositive-monotone,ana:actual-original-current-volume-kahler-right-limit}
\source{bou02-volume:page-0014}
\lean{TranscendentalNakamaye.Pending.PseudoeffectiveOriginalCurrentEquivalence.classPseudoeffectiveViaPotentials_iff_exists_original_current,TranscendentalNakamaye.Pending.OriginalCurrentVolumeUpperSemicontinuity.limsup_originalCurrentACVolume_le,TranscendentalNakamaye.Pending.OriginalCurrentVolumeRightLimit.originalCurrentACVolume_le_add_semipositive,TranscendentalNakamaye.Pending.OriginalCurrentVolumeRightLimit.tendsto_originalCurrentACVolume_add_smul_right}
\leanok
Let \((X,\omega)\) be a compact K\"ahler manifold of dimension \(n\). For a
pseudoeffective real \((1,1)\)-class \(\beta\) put
\[
 v(\beta)=\sup_{T}\int_X\bigl(T_{\mathrm{ac}}\bigr)^{n},
\]
the supremum over all closed positive currents \(T\in\beta\), where
\(T_{\mathrm{ac}}\) is the absolutely continuous part of the coefficients of
\(T\). If \(\beta_j\to\beta\) inside the pseudoeffective cone, then
\(v(\beta)\ge\limsup_j v(\beta_j)\). Since \(v\) is nondecreasing under adding a
semipositive class, it follows that
\(\lim_{\epsilon\downarrow0}v(\alpha+\epsilon\{\omega\})=v(\alpha)\) for every
pseudoeffective \(\alpha\).
\end{lemma}
\begin{proof}
The standing convention supplies a nonempty connected compact K\"ahler manifold of dimension $n>0$. Write $n=m+1$. The actual potential-pseudoeffective predicate is equivalent to containment of an ORIGINAL closed positive current in the same smooth Bott--Chern class. Thus it is exactly the pseudoeffective cone in this statement; neither analytic singularities nor an incoming potential is required. The original coefficientwise absolutely continuous part gives the canonical global AC top measure, and its mass on $X$ is precisely the term in the defining supremum $v(\beta)$.

Let $\beta_j\to\beta$ be the given sequence inside that cone. Apply the genuine original-current sequential upper-semicontinuity producer to the same classes and their actual quotient convergence. It internally chooses original currents of these classes whose whole AC measures are calibrated by their true volumes, refines the frequent set of any hypothetical strict limsup violation to a nontrivial ultrafilter, and produces a same-class positive weak limit using the full original tests. The genuine arbitrary-filter lower-calibration inequality, which passes fixed finite reference integrals to their scalar limits before differentiation rather than invoking arbitrary-filter Fatou, contradicts the defining volume supremum of that limit. Therefore $v(\beta)\geq\limsup_jv(\beta_j)$, including an initially infinite limsup and zero limiting volume.

For any smooth closed semipositive form $\eta$, its zero-potential ORIGINAL current $D$ is internally proved closed and fully positive with class $[\eta]$. For EVERY original admissible $T\in\alpha$, real linearity gives $[T+D]=\alpha+[\eta]$, and the actual coefficient AC part adds almost everywhere. On the common finite-basis conull set the genuine Hermitian density matrices are positive semidefinite and $A_{T+D}=A_T+A_D$. Determinant monotonicity consequently gives the ENTIRE measure inequality $\mu_T^{\rm AC}\leq\mu_{T+D}^{\rm AC}$. It preserves possibly infinite density integrals. Evaluation on $X$ and the original supremum bound for $T+D$, followed by the supremum over ALL original $T$, prove $v(\alpha)\leq v(\alpha+[\eta])$. The same constructed current sum and the current-cone equivalence show that semipositive addition preserves the genuine pseudoeffective cone.

In particular $f(\epsilon)=v(\alpha+\epsilon[\omega])$ is monotone: if $s\leq t$, apply the preceding inequality to $(t-s)\omega\geq0$. Choose $\epsilon_j=1/(j+1)>0$. A genuine smooth representative $\theta$ of $\alpha$ is selected internally from the quotient. Continuity of addition and real scalar multiplication in the actual smooth coefficient-jet topology gives $\theta+\epsilon_j\omega\to\theta$; continuity of the same quotient projection gives $\alpha+\epsilon_j[\omega]\to\alpha$. Each perturbed class is pseudoeffective by the preceding current sum. The already proved upper bound therefore gives $\limsup_j f(\epsilon_j)\leq v(\alpha)$. Since $v(\alpha)=f(0)\leq f(\epsilon_j)$, we also have $v(\alpha)\leq\liminf_j f(\epsilon_j)$. Completeness of the extended nonnegative order and its Hausdorff order topology force $f(\epsilon_j)\to v(\alpha)$.

Monotonicity gives the ENTIRE right-neighborhood limit of $f$ as $\inf f((0,\infty))$: values are bounded below by the infimum, and for every larger threshold a positive error with value below that threshold exists, after which all smaller positive errors satisfy the same upper bound. The reciprocal path tends to zero inside $(0,\infty)$, so uniqueness of limits identifies this infimum with $v(\alpha)$. Hence $\lim_{\epsilon\downarrow0}v(\alpha+\epsilon[\omega])=v(\alpha)$. No strictly positive limiting volume, incoming calibration, weak-limit class certificate or supplied volume limit has been used.
\end{proof}

\begin{theorem}[Grauert--Riemenschneider criterion, K\"ahler case]
\label{ana:volume-big-criterion}
\uses{ana:nonkahler-locus,ana:volume-semicontinuity,ana:geometric-setting,ana:actual-original-current-volume-big-iff,posclass:actual-original-current-pseudoeffective-iff,ana:actual-nef-original-current-volume,trace:model-potential-big,posclass:representative-iff}
\source{bou02-volume:4.8,dp04-kahler-cone:2.12}
\lean{TranscendentalNakamaye.Pending.OriginalCurrentVolumeBigness.isBigClass_iff_originalCurrentACVolume_pos,TranscendentalNakamaye.Pending.PseudoeffectiveOriginalCurrentEquivalence.classPseudoeffectiveViaPotentials_iff_exists_original_current,TranscendentalNakamaye.Pending.NefOriginalCurrentCalabiYau.originalCurrentACVolume_nef_eq_intersection,TranscendentalNakamaye.Pending.FiniteGeodesicsFullMassGeneral.bigWithRespectToViaPotentials_of_isBigClass}
\leanok
Let \((X,\omega)\) be a compact K\"ahler manifold and let \(\alpha\) be a
pseudoeffective real \((1,1)\)-class. Then \(\alpha\) is big, that is, contains a
K\"ahler current, if and only if \(v(\alpha)>0\). The case of a nef class
\(\alpha\) with \(\int_X\alpha^{n}>0\) is Demailly--P\u{a}un's theorem.
\end{theorem}
\begin{proof}
The standing convention supplies a nonempty connected compact K\"ahler manifold of dimension $n>0$ and a fixed K\"ahler form $\omega$. Write $n=m+1$. The original volume is the supremum over ALL original closed positive $(1,1)$-currents $T$ in the given smooth Bott--Chern class of the total mass of their canonical coefficientwise absolutely continuous top measures. Thus it is exactly $v_{\rm cur}$ in the original-current volume construction. The potential-pseudoeffective predicate agrees with the original current cone by the genuine original-current equivalence; no analytic presentation or supplied potential is needed for this dictionary.

We first check the fixed-background meaning of bigness. A genuine K\"ahler-current witness internally supplies a representative $\theta'$, a potential $u$ and a K\"ahler form $\Omega$ with $\theta'+dd^cu\geq\Omega$. Compact comparison of the two K\"ahler forms produces $A>0$ with $A\Omega-\omega\geq0$. Consequently the SAME current dominates $A^{-1}\omega$, with $A^{-1}>0$. For any prescribed smooth representative $\theta$ of the class, the actual smooth Bott--Chern relation transports $u$ by its smooth correction, preserving the full current and this inequality. Conversely a closed positive current in $\alpha$ dominating $\delta\omega$ has a genuine potential representation; after subtracting $\delta\omega$ the residual current is closed and positive. Its actual potential representation and the smooth class correction give a K\"ahler-current witness with lower bound $\delta\omega$. Hence the witness used below and the original convention $T\geq\delta\omega$ are equivalent on precisely this standing domain.

Suppose first that $\alpha$ is big. Unpack its genuine K\"ahler-current witness as $\theta+dd^cu\geq\Omega$, where $[\theta]=\alpha$ and $\Omega$ is K\"ahler. Form the original residual current $S=\theta-\Omega+dd^cu$ and the zero-potential original smooth current $D=\Omega+dd^c0$. Both are closed, of type $(1,1)$ and fully positive. Linearity of the actual full-test class map gives $[D+S]=[\Omega]+[\theta-\Omega]=\alpha$. The zero potential has analytic type internally, using the constant holomorphic generator $1$ and zero remainder in every chart. The genuine analytic zero-potential AC identification therefore gives the ENTIRE measure equality $\mu_D^{\rm AC}=\Omega^n$. For the SAME original positive currents, coefficientwise AC addition on a common conull set and positivity of their Hermitian density matrices imply determinant monotonicity and the ENTIRE measure inequality $\mu_D^{\rm AC}\leq\mu_{D+S}^{\rm AC}$. Nonemptiness and strict positivity of the smooth K\"ahler density give $\int_X\Omega^n>0$. Compactness makes that benchmark finite. The defining supremum now gives
\[
 0<\int_X\Omega^n=\mu_D^{\rm AC}(X)
       \leq\mu_{D+S}^{\rm AC}(X)\leq v(\alpha).
\]
Neither an incoming AC calibration nor a chosen admissible family is used.

Conversely assume $v(\alpha)>0$. The defining supremum internally selects an ORIGINAL closed fully positive current $T$ of actual class $\alpha$ with $\mu_T^{\rm AC}(X)>0$. The genuine all-current representation produces a smooth closed $\theta$ and a genuine $\theta$-PSH potential $u$ such that the SAME full-test current is $T=\theta+dd^cu$. Actual class calibration gives $[\theta]=\alpha$. The genuine global AC-domination theorem for arbitrary, possibly unbounded, $\theta$-PSH potentials gives
\[
 0<\mu_T^{\rm AC}(X)
          \leq\int_X\langle(\theta+dd^cu)^n\rangle=:M.
\]
This non-pluripolar mass is finite by the genuine global bounded-mass theorem. Thus $\eta=M/2$ is strictly positive and less than $M$.

Here is the proved positive-mass bigness route, rather than an invocation of the theorem being established. Genuine Demailly approximants $u_j\geq u$ have analytic singularities with remainders smooth off their poles and curvature errors $\delta_j\to0$. For EVERY $\epsilon>0$, select $j$ with $\delta_j<\epsilon$. Semipositive enlargement makes $u_j$ a genuine $(\theta+\epsilon\omega)$-PSH potential, and the actual potential mass comparison bounds its perturbed ambient mass below by $M>\eta$. Nondegeneracy internally supplies a nonpolar point, so this approximant is an admissible presentation on the entire manifold. Hence the perturbed ambient envelope exceeds the SAME $\eta>0$ at every positive error. The existing resolution step and uniform nef core are genuine producers: they internally resolve each such analytic presentation and produce the uniform positive K\"ahler lower bound, using the proved uniform mass-concentration argument, not a supplied resolution, nef-core or bigness certificate. The actual resolution assembly and uniform-perturbation argument then produce a genuine K\"ahler current in the UNPERTURBED class $[\theta]=\alpha$. The compact comparison above converts its lower bound to $\delta\omega$ for some $\delta>0$, exactly as in the original statement.

Finally let $\alpha$ be nef and $I=\int_X\alpha^n>0$. Choose a smooth closed representative $\theta$ internally from the actual quotient. Representative invariance of nefness gives the natural $\omega$-nef predicate for $\theta$, and the actual smooth intersection is the slot integral $I_\omega(\theta)=I$. The genuine ALL-original-current nef volume theorem proves
\[
 v(\alpha)=\operatorname{ofReal}(I_\omega(\theta))
             =\operatorname{ofReal}(I)>0.
\]
It proves this equality by the sharp AC bound for EVERY original admissible current and by internally constructing one nef current with the entire calibrated AC measure attaining that bound, including the zero-intersection case. No maximizing-current or calibration certificate is an input. Applying the already proved positive-volume implication yields a genuine K\"ahler current in $\alpha$, and compact comparison again gives $T\geq\delta\omega$. This proves the Demailly--P\u{a}un case mentioned in the original final sentence as well as both directions of the general pseudoeffective criterion.
\end{proof}

\begin{lemma}[Bigness is a uniform K\"ahler lower bound on the perturbed classes]
\label{ana:uniform-perturbed-lower-bound}
\uses{ana:nonkahler-locus,gnpp:kahler-domination}
\source{bou02-volume:4.8}
\lean{TranscendentalNakamaye.Pending.UniformPerturbedKahlerLowerBound.isBigClass_iff_hasUniformPerturbedKahlerLowerBound}
\leanok
Let \((X,\omega)\) be a compact K\"ahler manifold and let \(\alpha\) be a real
\((1,1)\)-class. Then \(\alpha\) is big if and only if there is a constant \(c>0\) such
that for arbitrarily small \(\epsilon>0\) the class \(\alpha+\epsilon\{\omega\}\)
contains a closed real \((1,1)\)-current \(T\) with \(T\ge c\,\omega\).

Because this is an \emph{equivalence}, the uniform bound must not be taken as a separate
hypothesis on the way to bigness: it is the conclusion restated. In particular the last step
of Boucksom's proof of Theorem~\ref{ana:volume-big-criterion} carries none of that theorem's
content, which lies entirely upstream of it --- the concentration of mass, the
Aubin--Calabi--Yau theorem, and the nef classes on the products \(X_k\times X_k\).
\end{lemma}
\begin{proof}\leanok
\uses{gnpp:kahler-domination}
Suppose such a \(c\) exists. Choose \(\epsilon\) with \(0<\epsilon<c\) and a current
\(T=\theta+\epsilon\omega+dd^cu\ge c\,\omega\) in \(\alpha+\epsilon\{\omega\}\), where
\(\{\theta\}=\alpha\). Then \(T-\epsilon\omega=\theta+dd^cu\) lies in \(\alpha\) and
\[
 T-\epsilon\omega\ \ge\ c\,\omega-\epsilon\omega=(c-\epsilon)\,\omega>0 ,
\]
so \(\alpha\) contains a K\"ahler current and is big. The potential \(u\) is untouched:
only the smooth representative and the K\"ahler lower bound move.

Conversely, let \(T=\theta+dd^cu\ge\omega'\) be a K\"ahler current in \(\alpha\), with
\(\omega'\) a K\"ahler form. Two K\"ahler forms on a compact manifold are comparable: by
K\"ahler domination (Lemma~\ref{gnpp:kahler-domination}) applied to the closed real
\((1,1)\)-form \(-\omega\) with background \(\omega'\) there is \(A>0\) with
\(-\omega+A\omega'\ge0\), that is, \(\omega'\ge A^{-1}\omega\). Put \(c=A^{-1}\). For
every \(\epsilon>0\) the current \(T+\epsilon\omega\) lies in
\(\alpha+\epsilon\{\omega\}\) and satisfies
\[
 T+\epsilon\omega\ \ge\ \omega'+\epsilon\omega\ \ge\ c\,\omega ,
\]
because \(\omega'-c\,\omega\) and \(\epsilon\omega\) are both semipositive. So the bound
holds at every \(\epsilon>0\), with one and the same constant \(c\).
\end{proof}

\begin{proposition}[Analytic singularities: the Siu and Lebesgue decompositions agree]
\label{ana:analytic-singularities-lebesgue}
\uses{ana:geometric-setting,as:smooth-off-pole-type,ana:analytic-type-pole-set,ana:coefficient-lebesgue-decomposition,gnpp:global-measure,ana:actual-original-native-siu-exterior-mass-mother}
\source{bou02-volume:2.1,bou02-volume:2.5}
\lean{TranscendentalNakamaye.Pending.ActualOriginalNativeSiuExteriorMassMother.exists_original_native_siu_exterior_mass_mother}
\leanok
In the compact positive-dimensional setting of Convention~\ref{ana:geometric-setting},
let \(T\) be a closed real \((1,1)\)-current with genuine local quasi-psh
potentials having analytic singularities \emph{with a bounded remainder smooth
off the poles}, in the sense of Definition~\ref{as:smooth-off-pole-type}.
Thus every point has a local native potential
\[
 u=c\log\!\left(\sum_{\ell=1}^{N}|f_\ell|^2\right)+b,\qquad c>0,
\]
where the \(f_\ell\) are holomorphic, their common zero germ is proper,
and \(b\) is locally bounded and smooth wherever the original \(f_\ell\)
do not all vanish. No smoothness of \(b\) across those zeros is assumed.
The native convention is \(dd^c_N=\frac{i}{2\pi}\partial\bar\partial\), so
\(dd^c_N\log|g|^2=[\operatorname{div}(g)]\).

Use the integration-current normalization of the Siu coefficients: if
\(\lambda_T(x)\) is the ordinary-log growth of a native local potential,
put \(\nu_T(x)=\lambda_T(x)/2\). Equivalently, \(\nu_T(x)\) is the
ordinary-log growth of the source local potential \(u/2\) for the
independently calibrated operator \(dd^c_B=\frac{i}{\pi}\partial\bar\partial\)
reading the \emph{same} \(T\). Write the Siu decomposition in this convention as
\[
 T=R_{\rm Siu}+D,\qquad D=\sum_{j=1}^{K}q_j[V_j],\qquad
 q_j=\inf_{x\in V_j}\nu_T(x)=\frac12\inf_{x\in V_j}\lambda_T(x).
\]
Then its divisorial sum is finite, consists of the codimension-one global
components of the common polar set \(P\), and is the divisorial part of the
local defining ideals with their analytic coefficients: at a generic point
of \(V_j\), \(q_j=c\,\operatorname{ord}_{V_j}(f_1,\ldots,f_N)\).
The coefficients of \(R_{\rm Siu}\) are locally \(L^1\), and the canonical
coefficientwise Lebesgue decomposition is
\[
 T_{\rm ac}=R_{\rm Siu},\qquad T_{\rm sg}=D.
\]
The current \(T\) is smooth on \(X\setminus P\), and
\(T_{\rm ac}=\mathbf1_{X\setminus P}T\), as a locally integrable
coefficient current. These assertions concern the original signed \(T\);
positivity is not an assumption here.

If this same \(T\) is semipositive, its actual non-pluripolar product and
canonical absolutely continuous top-density measure agree, and
\[
 \operatorname{NP}(T)(X)
   =\int_{X\setminus P}T^n
   =\int_X\bigl(T_{\rm ac}\bigr)^n.
\]
The integrals are nonnegative extended integrals. In particular this includes
the case with no divisorial part from the original mass assertion.
\end{proposition}
\begin{proof}
We keep the original current and its complete local coefficient tests
throughout. Let \(u=c\log\sum|f_\ell|^2+b\) be one of its genuine native
potentials. Its literal half \(w=u/2\) is a genuine local quasi-psh potential
with the same holomorphic family, coefficient \(c/2\), and bounded remainder
\(b/2\) smooth only off the same original common zeros. The source covector
and its exterior differential are independently defined on coordinate
tests. Their full distributional calibration gives \(dd^c_B=2dd^c_N\);
linearity of the locally integrable regular distribution then gives
\(dd^c_B\operatorname{Reg}(w)=dd^c_N\operatorname{Reg}(u)\) on every
two-form coefficient and every complementary current test. Thus both
potentials read this same \(T\), rather than \(T\) and \(2T\).
The source ordinary-log reader of \(w\) is the squared-log reader of \(u\),
and native ordinary-log growth is twice it. This identity holds pointwise
and for whole-set infima; it justifies exactly the normalization in the
statement.

The original native-input mother theorem constructs the common analytic
polar set \(P\), its genuine global components, and the whole-prime
weights from these original local readers. Here compactness makes the
locally finite component family finite. To explain the local factor and
its coefficient, choose the common holomorphic factor \(g\) and a
primitive residual family \(h_\ell\), so \(f_\ell=g h_\ell\) on a smaller
neighbourhood. The scalar regular distributions and all coefficient tests
satisfy
\[
 T=c[\operatorname{div}(g)]
       +dd^c_N\operatorname{Reg}
          \left(c\log\sum_\ell|h_\ell|^2+b\right).
\]
This is the native Poincar\'e--Lelong formula for \(\log|g|^2\), with the
original \(c\) and \(b\) retained. For a prime factor \(q\) of \(g\), the
actual original finite-family factor theorem produces a nonempty open
piece of its assigned global prime on which the remaining factor and one
residual generator are units. The point growth there is
\(2c\,\operatorname{ord}_{q}(g)\). Analytic Lelong superlevel sets and
irreducibility identify the infimum on this open piece with the infimum
on the whole global prime. Consequently the actual integration-current
weight is \(c\,\operatorname{ord}_{q}(g)\), not twice that number.
These assignments use the original family and identify genuine global
components, so they agree on overlaps. They produce a fixed global finite
\(D\) and the literal global current \(R=T-D\), with equality on all
variable complementary tests.

The absolute continuity of this signed residual does not follow from
Haar-nullity of \(P\) alone. The local residual producer subtracts the
holomorphic logarithm off its divisor, where it is pluriharmonic, and
internally makes a smooth quadratic correction so that the residual
potential is psh. Boundedness of the original \(b\), together with the
primitive generator bounds, gives a locally bounded-above psh extension
across that divisor. On the complement of the primitive common zeros it
is locally bounded. The bounded-locus current noncharging theorem excludes
mass there on the pluripolar divisor; the primitive zero set is handled
by the genuine codimension-two coefficient/trace-null producer. On the
complement of the original \(P\), the original family is nonzero and
\(b\) is smooth, so the same residual is smooth. Its positive corrected
current therefore has absolutely continuous coefficients, and removing
the smooth correction gives locally \(L^1\) two-form coefficients for
the original signed \(R\). This argument uses off-original-pole smoothness
and boundedness; it does not require smoothness across a pole or positivity
of the original total current.

More precisely, the mother produces in each chosen coordinate neighbourhood
a locally integrable real \((1,1)\)-form \(\alpha\) whose regular
distribution is every coefficient of this same \(R\), and whose complete
complementary reconstruction is the restricted current \(R\). Its polar
mask \(\mathbf1_{P^c}\alpha\) agrees with \(\alpha\) almost everywhere,
is locally integrable, and reads the identical coefficients and current.
On \(X\setminus P\) these coefficients agree with the smooth native
\(dd^c_N(u^{\rm real})\); the same fixed \(D\) is zero there. The
divisor coefficients are singular with respect to coordinate Haar measure,
since their supports are analytic hypersurfaces. The actually constructed
finite signed coefficient measures have canonical singular-part pairings
equal to \(D\) and absolutely continuous-part pairings equal to \(R\),
for every compact scalar test and every constant complementary form in
the full form span. Coefficient uniqueness and complete variable-test
overlap compatibility identify these as global coefficient-current
identities, giving \(T_{\rm sg}=D\), \(T_{\rm ac}=R\), and the literal
polar restriction assertion. No trace equality is substituted for a
complete current identity.

To compare with the given Siu decomposition, its divisor coefficients
are the same whole generic integration-current coefficients just computed.
A prime not contained in \(P\) meets a region where \(T\) is smooth and
has generic coefficient zero; a codimension-one prime contained in \(P\)
is one of the global components catalogued above. Hence its divisorial
current is this same \(D\), and subtraction from the same \(T\) gives
\(R_{\rm Siu}=T-D=R\). This comparison uses the given Siu decomposition
and the independently produced whole generic weights; it does not assert
that the outer mother type supplies a new residual superlevel-dimension
field or infer such a field from \(L^1\) alone.

For the last assertion suppose this same \(T\) is semipositive on its full
complementary test cone. The mother internally produces genuine psh
local potentials of \(T\), the actual diagonal non-pluripolar truncation
measures, and their inverse-chart transports. These measures give zero
mass to the same \(P\). On the selected measurable pieces of its finite
disjoint atlas outside \(P\), the same original potential is smooth and
the local non-pluripolar measure is its genuine native top exterior density
against coordinate Haar measure. The actual chart transports and disjoint
gluing give the global NP measure. Independently, the canonical RN top
measure is constructed from the absolutely continuous coefficient forms
of the same \(T\). The proved equality of these entire measures, their
vanishing on \(P\), and the actual exterior-density piece integrals give
\[
 \operatorname{NP}(T)(X)
 =\int_{X\setminus P}T^n
 =\int_X(T_{\rm ac})^n.
\]
All quantities may be infinite. Neither the NP measure nor the RN top
measure is defined by the other, and neither a desired mass identity nor
a selected divisor/atlas is an input. This proves the no-divisor special
case as well.
\end{proof}

The finite global sum here uses the explicitly standing compact
positive-dimensional \(X\). On an arbitrary noncompact manifold the
corresponding divisorial sum is in general only locally finite, and that
broader assertion is not claimed here. Boucksom's local smooth-remainder
argument is used in the stated bounded-remainder smooth-off-poles
convention. In the native convention one has
\(dd^c_N\log|g|=\tfrac12[\operatorname{div}(g)]\); the proof above therefore
uses \(\log|g|^2\) with its original coefficient instead of the previous
half-log formula.

\begin{theorem}[A uniformly positive perturbed ambient volume produces a K\"ahler current]
\label{ana:ambient-volume-uniform-big}
\uses{ana:restricted-volume,ana:nonkahler-locus,gnpp:global-measure}
\source{bou02-volume:4.8,bou02-volume:page-0014}
\lean{TranscendentalNakamaye.Pending.UniformAmbientVolumeBigness.isBigClass_of_hasUniformAmbientMassLowerBound}
\leanok
Let \((X,\omega)\) be a compact K\"ahler manifold of dimension \(n\ge1\), let
\(\alpha=\{\theta\}\) be a pseudoeffective real \((1,1)\)-class, and write, as in
Lemma~\ref{ana:ambient-volume-uniform-bound},
\[
 E(\epsilon)=\sup_{T\in\mathcal A_\epsilon}\operatorname{NP}(T+\epsilon\omega)(X),
\]
the supremum of the ambient diagonal non-pluripolar masses over the
\(\epsilon\)-admissible presentations of \(\alpha\) along \(X\). If there is
\(\eta>0\) with \(\eta<E(\epsilon)\) for \emph{every} \(\epsilon>0\), then \(\alpha\)
is big.
\end{theorem}
\begin{proof}
\uses{bou:resolution-assembly,bou:resolution-step,bou:uniform-nef-core,bou:uniform-nef-core-constant}
\leanok
This is Boucksom's proof of his Grauert--Riemenschneider criterion, entered at the point where it has currents with analytic singularities
and masses bounded away from \(0\) in the classes \(\alpha+\epsilon\{\omega\}\); the hypothesis supplies them. The presentations are taken along \(X\) itself, an
irreducible analytic subset of \(X\), so \(X\) is connected (Lemma~\ref{bou:uniform-nef-core-constant}). On a compact connected K\"ahler manifold of positive
dimension the resolution step (R) holds by Lemma~\ref{bou:resolution-step} and the uniform nef core (N) by Lemma~\ref{bou:uniform-nef-core}; part (2) of
Lemma~\ref{bou:resolution-assembly} then gives the bigness of \(\alpha\) under exactly the hypothesis of the theorem. Pseudoeffectivity of \(\alpha\) is not used.

Positivity of the dimension is used in the uniform nef core (the graph of a blowup has codimension \(n\ge1\) in \(Y\times X\)); in dimension \(0\) every class is big
outright, and that case is treated separately.
\end{proof}
This route differs from Boucksom's text in three places, none of which changes the mathematics of his argument. Boucksom first produces, from
\(v(\alpha)>0\), currents with analytic singularities whose absolutely continuous parts have volume bounded away from \(0\); here they come from the
hypothesis, so neither the volume function \(v\) nor its semicontinuity (Lemma~\ref{ana:volume-semicontinuity}) nor
Proposition~\ref{ana:analytic-singularities-lebesgue} is needed. Boucksom's resolved currents have a smooth part; here the remainders are only bounded,
and the nefness of the resolved class comes from Lemma~\ref{ana:residual-nef}, its top intersection number from the non-pluripolar mass. Finally
Boucksom concludes with the bigness of \(\alpha\) directly, while here the uniform bound is handed to Lemma~\ref{ana:uniform-perturbed-lower-bound}.
Boucksom's Lemma~4.10 is run on \(Y\times X\) along the graph of the blowup \(\mu\colon Y\to X\) (Lemma~\ref{bou:uniform-nef-core}) rather than on
\(Y\times Y\) along the diagonal: the preimage of the diagonal of \(X\times X\) under \(\mu\times\mathrm{id}\) is exactly the graph, a connected submanifold, so no
exceptional components occur and the lower bound is carried as \(c\,\mu^*\omega\), which descends exactly.
Boucksom's statement itself, Theorem~\ref{ana:volume-big-criterion}, together with Lemma~\ref{ana:volume-semicontinuity} and
Proposition~\ref{ana:analytic-singularities-lebesgue}, gives the theorem as well: for an \(\epsilon\)-admissible presentation the non-pluripolar mass is the
integral of the absolutely continuous part, so \(E(\epsilon)\le v(\alpha+\epsilon\{\omega\})\), and \(v(\alpha)\ge\eta\) in the limit.

The lemmas used above follow.

\begin{lemma}[Potentials of global analytic type carrying the mass]
\label{bou:global-type-potentials}
\uses{ana:demailly-algebraic-approximation,bc:class-equality,posclass:potential-shift,gpt:global-measure,mmass:monotonicity,mmass:diagonal-background-monotone}
\source{bou02-volume:2.6,bou02-volume:4.8}
\lean{TranscendentalNakamaye.Pending.BoucksomGlobalTypePotentials.exists_potential_of_presentation,TranscendentalNakamaye.Pending.BoucksomGlobalTypePotentials.exists_globalTypePotential}
\leanok
In the situation of Theorem~\ref{ana:ambient-volume-uniform-big}, let \(\theta\) be a smooth representative of \(\alpha\) and \(\epsilon>0\). There are
\(0<\epsilon'\le2\epsilon\) and a \((\theta+\epsilon'\omega)\)-psh potential \(\psi\) with analytic singularities of type \((c,\mathfrak a)\), \(c\in\mathbf Q_{>0}\) and
\(\mathfrak a\) coherent, whose remainders are bounded, such that \(\int_X\langle(\theta+\epsilon'\omega)_\psi^n\rangle>\eta\). One may take \(\epsilon'=2\epsilon\) and
\(\psi\) smooth off its poles.
\end{lemma}
\begin{proof}
\leanok
As \(\eta<E(\epsilon)\), there is an \(\epsilon\)-admissible presentation of \(\alpha\) along \(X\), with representative \(\theta_p\) and potential \(u\), such that
\(\int_X\langle(\theta_p+\epsilon\omega)_u^n\rangle>\eta\). As \(\{\theta_p\}=\{\theta\}\), \(\theta=\theta_p+dd^cf\) with \(f\) smooth (Lemma~\ref{bc:class-equality}), so
\(u-f\) is a \((\theta+\epsilon\omega)\)-psh potential (Lemma~\ref{posclass:potential-shift}) with the same non-pluripolar measure (Theorem~\ref{gpt:global-measure}).
Demailly's approximation of \(u-f\) (Theorem~\ref{ana:demailly-algebraic-approximation}, for the form \(\theta+\epsilon\omega\)) gives potentials \(\psi_j\ge u-f\) of
type \((c_j,\mathfrak a_j)\), \(c_j\in\mathbf Q_{>0}\), with bounded remainders, which are \((\theta+(\epsilon+\delta_j)\omega)\)-psh with \(\delta_j\to0\). Take \(j\)
with \(\delta_j<\epsilon\) and \(\psi=\psi_j\), which is smooth off its poles and \((\theta+(\epsilon+\delta_j)\omega)\)-psh, hence \((\theta+\epsilon'\omega)\)-psh for
\(\epsilon'=2\epsilon\). Since \(u-f\le\psi\), monotonicity of the mass
(Theorem~\ref{mmass:monotonicity}) and of the background (Theorem~\ref{mmass:diagonal-background-monotone}) give
\(\int_X\langle(\theta+\epsilon'\omega)_\psi^n\rangle\ge\int_X\langle(\theta+\epsilon\omega)_{u-f}^n\rangle>\eta\).
\end{proof}
The analytic singularities of the presentation itself are not used: they are local, with exponents and remainders varying from point to point,
whereas the resolution below needs one coherent ideal and one rational exponent, which Demailly's approximation provides.

\begin{lemma}[Boucksom 4.8 assembled from the resolution step and the uniform nef core, both as hypotheses]
\label{bou:resolution-assembly}
\uses{bou:global-type-potentials,pc:blowup-descent,ana:uniform-perturbed-lower-bound,gnpp:kahler-domination,ana:blowups}
\source{bou02-volume:4.8,bou02-volume:4.10}
\lean{TranscendentalNakamaye.Pending.BoucksomResolutionAssembly.pullback_smul,TranscendentalNakamaye.Pending.BoucksomResolutionAssembly.exists_kahlerCurrent_of_pullback_thetaPsh,TranscendentalNakamaye.Pending.BoucksomResolutionAssembly.ResolutionStep,TranscendentalNakamaye.Pending.BoucksomResolutionAssembly.UniformNefCore,TranscendentalNakamaye.Pending.BoucksomResolutionAssembly.exists_kahlerCurrent_of_globalTypePotential,TranscendentalNakamaye.Pending.BoucksomResolutionAssembly.hasUniformPerturbedKahlerLowerBound_of_resolutionStep,TranscendentalNakamaye.Pending.BoucksomResolutionAssembly.isBigClass_of_resolutionStep_of_uniformNefCore}
\leanok
Let \((X,\omega)\) be a compact K\"ahler manifold of dimension \(n\).
(1) If \(\mu\colon Y\to X\) is a finite sequence of blowups with smooth centres, \(c>0\), and \(\mu^*(\beta-c\,\omega)\) has a psh potential on \(Y\), then \(\{\beta\}\)
contains a K\"ahler current with lower bound exactly \(c\,\omega\).
Consider the two statements:
\begin{itemize}
\item[(R)] (\emph{resolution step}) for every smooth closed real \((1,1)\)-form \(\beta\), every \(\beta\)-psh potential \(\psi\) with analytic singularities, smooth off its
poles, and every \(C\) with \(C\omega-\beta\ge0\), there are a blowup sequence \(\mu\colon Y\to X\), a K\"ahler form \(\omega'\) on \(Y\) and a smooth closed real
\((1,1)\)-form \(b\) on \(Y\), nef, with \(\int_Yb^n\ge\int_X\langle\beta_\psi^n\rangle\), \(\int_Yb^i\wedge(\mu^*\omega)^{n-i}\le C^i\int_X\omega^n\) for \(i\le n\), and such that
for every smooth closed \(\gamma\) on \(X\) a \((b-\mu^*\gamma)\)-psh potential on \(Y\) yields a \(\mu^*(\beta-\gamma)\)-psh potential on \(Y\);
\item[(N)] (\emph{uniform nef core}) for \(\eta>0\) and \(C\) there is \(c>0\) such that for every blowup sequence \(\mu\colon Y\to X\), K\"ahler form \(\omega'\) on \(Y\)
and nef smooth closed \(b\) on \(Y\) with \(\int_Yb^n\ge\eta\) and \(\int_Yb^i\wedge(\mu^*\omega)^{n-i}\le C\) for \(i\le n\), the form \(b-c\,\mu^*\omega\) has a psh potential.
\end{itemize}
(2) \emph{Assume (R) and (N) (hypotheses).} Under the hypothesis of Theorem~\ref{ana:ambient-volume-uniform-big} (a uniform lower bound \(\eta>0\) on the perturbed
ambient volumes) there is \(c>0\) such that for arbitrarily small \(\epsilon>0\) the class \(\alpha+\epsilon\{\omega\}\) contains a K\"ahler current with lower bound
exactly \(c\,\omega\); hence \(\alpha\) is big.
\end{lemma}
\begin{proof}
\leanok
(1) A \(\mu^*(\beta-c\omega)\)-psh potential on \(Y\) descends to a \((\beta-c\omega)\)-psh potential \(s\) on \(X\) (Proposition~\ref{pc:blowup-descent}); then
\(\beta+dd^cs\ge c\,\omega\). (2) Fix \(C\ge1\) with \((C-1)\omega-\theta\ge0\) (Lemma~\ref{gnpp:kahler-domination}) and a real level \(0<\eta_0\le\eta\); let \(c\) be the
constant of (N) for \(\eta_0\) and \(C^n\int_X\omega^n\). For small \(\epsilon_0\), Lemma~\ref{bou:global-type-potentials} gives a \((\theta+\epsilon'\omega)\)-psh
potential \(\psi\) of global analytic type, smooth off its poles, with mass \(>\eta_0\), \(\epsilon'\le2\epsilon_0\). (R) for \(\beta=\theta+\epsilon'\omega\) gives \(\mu\), \(b\) with
\(\int b^n\ge\eta_0\) and mixed bounds \(\le C^i\int\omega^n\le C^n\int\omega^n\); (N) gives a \((b-c\mu^*\omega)\)-psh potential, the transfer of (R) a
\(\mu^*(\beta-c\omega)\)-psh potential, and (1) the K\"ahler current. Bigness follows by Lemma~\ref{ana:uniform-perturbed-lower-bound}.
\end{proof}
(R) is, in Boucksom's text, Hironaka's principalization of the ideal of \(\psi\) followed by Lemma~\ref{bou:resolved-nef-part}; it holds on every compact
connected K\"ahler manifold of positive dimension by Lemma~\ref{bou:resolution-step}, which uses the log resolution of
Proposition~\ref{pc:domination-log-resolution} (through the principalization Lemma~\ref{ana:relative-principalization}) in place of that lemma. (N) is Boucksom's
Lemma~4.10, uniform in the blowup sequence; it holds on every compact connected K\"ahler manifold of positive dimension (Lemma~\ref{bou:uniform-nef-core},
through Yau's theorem, Theorem~\ref{bou:yau}). The lemma itself takes both as hypotheses.

\begin{lemma}[The constant of the resolution step is nonnegative]
\label{bou:resolution-step-constant}
\uses{bou:nef-mixed-mass,mmass:monotonicity,bou:mixed-mass-class-invariance}
\lean{TranscendentalNakamaye.Pending.BoucksomResolutionStep.pullback_sub,TranscendentalNakamaye.Pending.BoucksomResolutionStep.slotIntegral_smul_prefix,TranscendentalNakamaye.Pending.BoucksomResolutionStep.nonneg_of_isThetaPshPotential}
\leanok
Let \((X,\omega)\) be a nonempty compact K\"ahler manifold of dimension \(n\ge1\), \(\beta\) a smooth closed real \((1,1)\)-form carrying a \(\beta\)-psh potential, and
\(C\in\mathbf R\) with \(C\omega-\beta\ge0\). Then \(C\ge0\). Moreover, for every smooth closed real \((1,1)\)-form \(a\), \(C\in\mathbf R\) and \(i\in\mathbf N\),
\(\int_X(Ca)^{\min(n,i)}\wedge a^{n-\min(n,i)}=C^{\min(n,i)}\int_Xa^n\), and pull-back of forms by a holomorphic map commutes with subtraction.
\end{lemma}
\begin{proof}
\leanok
If \(C<0\), a \(\beta\)-psh \(\psi\) is also psh for the semipositive form \(\beta+(C\omega-\beta)+|C|\omega=0\), so the mixed mass of \((0,\psi)\) is at most that of
\((0,0)\), which is \(0\) (Lemma~\ref{mmass:monotonicity}); adding \(|C|\omega\) one slot at a time bounds the same mass from below by \(|C|^n\int_X\omega^n>0\), a contradiction.
The scaling identity is multilinearity of the mixed top-degree density (Lemma~\ref{bou:nef-mixed-mass}); pull-back is linear.
\end{proof}

\begin{lemma}[Boucksom's resolution step]
\label{bou:resolution-step}
\uses{bou:resolution-assembly,pc:domination-log-resolution,ana:relative-principalization,bou:nef-bounded-mixed-mass,bou:divisor-mixed-mass,bou:resolution-step-constant,mmass:slot-background-monotone,mmass:monotonicity,pc:domination-modification-mass,bou:mixed-mass-class-invariance}
\source{bou02-volume:4.8,bou02-volume:2.4}
\lean{TranscendentalNakamaye.Pending.BoucksomResolutionStep.resolutionStep}
\leanok
Let \((X,\omega)\) be a compact connected K\"ahler manifold of dimension \(n\ge1\). Then statement (R) of Lemma~\ref{bou:resolution-assembly} holds.
\end{lemma}
\begin{proof}
\leanok
Let \(\beta\), \(\psi\) (with analytic singularities, smooth off its poles) and \(C\) with \(C\omega-\beta\ge0\) be given. Proposition~\ref{pc:domination-log-resolution}
(which uses the principalization statement Lemma~\ref{ana:relative-principalization} as stated) gives a blowup sequence \(\rho\colon Y\to X\) with \(Y\) compact, connected and
K\"ahler (\(\omega'\)), an effective \(\mathbf Q\)-divisor \(F\) with a smooth metric, \(\psi\circ\rho=\psi_F+b'\) off \(\operatorname{Supp}F\) with \(b'\) bounded and
\((\rho^*\beta-\Theta_F)\)-psh, the form \(b=\rho^*\beta-\Theta_F\) nef, and \(\int_Y\langle b_{b'}^n\rangle=\int_X\langle\beta_\psi^n\rangle\). Then
\(\int_X\langle\beta_\psi^n\rangle=\int_Yb^n\) by Lemma~\ref{bou:nef-bounded-mixed-mass}. Let \(i\le n\) and \(m=\rho^*\omega\), semipositive, hence nef; \(C\ge0\) by
Lemma~\ref{bou:resolution-step-constant}. With \(b\) in the first \(i\) slots and \(m\) in the others,
\[
 \int_Yb^i\wedge m^{n-i}=\int_Y\langle b_{b'}^i\wedge m_0^{n-i}\rangle=\int_Y\langle(\rho^*\beta)_{\psi\circ\rho}^i\wedge m_0^{n-i}\rangle
 \le\int_Y\langle(Cm)_{\psi\circ\rho}^i\wedge m_0^{n-i}\rangle\le\int_Y(Cm)^i\wedge m^{n-i}=C^i\int_X\omega^n ,
\]
by Lemma~\ref{bou:nef-bounded-mixed-mass} (bounded potentials \(b'\) and \(0\) in nef classes), Lemma~\ref{bou:divisor-mixed-mass} (the first \(i\) slots corrected),
Lemma~\ref{mmass:slot-background-monotone} (the backgrounds grow by \(\rho^*(C\omega-\beta)\ge0\)), Lemma~\ref{mmass:monotonicity} (\(\psi\circ\rho\) is bounded above, so more
singular than \(0\)) with Lemma~\ref{bou:mixed-mass-class-invariance}, and the scaling identity of Lemma~\ref{bou:resolution-step-constant} with
\(\int_Ym^n=\int_X\omega^n\) (Lemma~\ref{pc:domination-modification-mass}). Finally, if \(v\) is \((b-\rho^*\gamma)\)-psh, then \(v+\psi_F\) is \(\rho^*(\beta-\gamma)\)-psh,
since \(\psi_F\) is \(\Theta_F\)-psh.
\end{proof}

\begin{lemma}[The curvature form of a metered divisor represents the divisor]
\label{bou:divisor-class}
\uses{ana:ideal-sheaves,ana:exceptional-metric,ana:blowups,ana:geometric-setting}
\source{demailly-cadg:III.2.15,bou02-volume:2.9}
\lean{TranscendentalNakamaye.Pending.LelongPoincareDivisor.lelongPoincare,TranscendentalNakamaye.Pending.LelongPoincareDivisor.exists_exceptional_correction_class,TranscendentalNakamaye.Pending.LelongPoincareDivisor.ddCWeight_eq,TranscendentalNakamaye.Pending.LelongPoincareDivisor.integrationCurrent,TranscendentalNakamaye.Pending.LelongPoincareDivisor.chartLocallyIntegrable_weight,TranscendentalNakamaye.Pending.LelongPoincareDivisor.restrictForm_coefficients_eq_sum,TranscendentalNakamaye.Pending.LelongPoincareDivisor.globalCurrent_ext,TranscendentalNakamaye.Pending.LelongPoincareDivisor.smoothFormDistribution_sub}
\leanok
Let \(Y\) be a complex manifold and \(E\) an effective \(\mathbf Q\)-divisor on \(Y\) with a smooth Hermitian metric on \(\mathcal O_Y(E)\)
(Definition~\ref{ana:ideal-sheaves}): a global function \(\psi_E\) which near every point is \(\sum_ia_i\log|g_i|^2+h\), with \(a_i\in\mathbf Q_{\ge0}\), local
equations \(g_i\) vanishing identically on no open set and \(h\) smooth, and the curvature form \(\Theta_E=-dd^ch\). Then \(\psi_E\) is locally integrable and, as
currents on \(Y\),
\[
 dd^c\psi_E=[E]-\Theta_E ,
\]
where \([E]\) is the closed positive \((1,1)\)-current equal to \(\sum_ia_i[\operatorname{div}g_i]\) on the domain of each local presentation. So \([E]\) and the
smooth closed form \(\Theta_E\) differ by \(dd^c\) of a global function: they define the same class, \(\{E\}=\{\Theta_E\}\). In particular, if
\(\mu\colon\widetilde X\to X\) is a modification of a compact K\"ahler manifold \((X,\omega)\) and \(F\) is the effective \(\mathbf Q\)-divisor of
Lemma~\ref{ana:exceptional-metric}, supported in \(\operatorname{Exc}(\mu)\), then \(\widetilde\omega=\mu^*\omega-\Theta_F\) is a K\"ahler form with
\(\{\widetilde\omega\}=\mu^*\{\omega\}-\{F\}\).
\end{lemma}
\begin{proof}
\leanok
\uses{ana:metered-divisor-lelong-poincare,ana:exceptional-metric}
Let \(g\) be holomorphic near a point, vanishing identically on no open set. Then \(\log|g|\) is locally integrable and the Lelong--Poincar\'e equation
(Demailly, III.2.15) reads \(\frac{\mathrm i}{\pi}d'd''\log|g|=[\operatorname{div}g]\). With the normalisation \(dd^c=(2\pi)^{-1}\mathrm i\partial\bar\partial\) of
Convention~\ref{ana:geometric-setting}, \(dd^c\log|g|^2=\frac{\mathrm i}{\pi}\partial\bar\partial\log|g|\), so \(dd^c\log|g|^2=[\operatorname{div}g]\). On the domain
of a local presentation, \(\psi_E\) is therefore locally integrable (a finite sum of such logarithms and a smooth function) and
\(dd^c\psi_E=\sum_ia_i[\operatorname{div}g_i]+dd^ch=\sum_ia_i[\operatorname{div}g_i]-\Theta_E\). The left-hand side \(dd^c\psi_E+\Theta_E\) is a globally defined current,
because \(\psi_E\) and \(\Theta_E\) are global; it is locally \(\sum_ia_i[\operatorname{div}g_i]\), a closed positive current, and this local description is \([E]\).
Hence \([E]-\Theta_E=dd^c\psi_E\) with \(\psi_E\in L^1_{\mathrm{loc}}(Y)\).

For the last assertion, Lemma~\ref{ana:exceptional-metric} gives \(F\), effective, with a smooth metric, supported in \(\operatorname{Exc}(\mu)\), such that
\(\widetilde\omega=\mu^*\omega-\Theta_F\) is K\"ahler; its class is \(\mu^*\{\omega\}-\{\Theta_F\}=\mu^*\{\omega\}-\{F\}\) by the first part. This is Boucksom's Lemma~2.9
applied blowup by blowup: a smooth closed form cohomologous to the exceptional divisor whose small negative multiple corrects \(\mu^*\omega\) to a K\"ahler form.
Formalized as Theorem~\ref{ana:metered-divisor-lelong-poincare} (all assertions but the last; for \(\dim Y=0\) all two-forms vanish) and,
for the last one, with \([F]\) the current of Theorem~\ref{ana:metered-divisor-lelong-poincare}: as currents,
\(\widetilde\omega=\mu^*\omega-[F]+dd^c\psi_F\).
\end{proof}

\begin{lemma}[Resolving a potential of analytic type: the nef part and its mass]
\label{bou:resolved-nef-part}
\uses{ana:blowups,hir:principalization,ana:divisor-metric,psh:holomorphic-pullback,pc:analytic-type-restriction,ana:residual-nef,pc:domination-residual,pc:domination-modification-mass,bou:nef-bounded-mixed-mass,bou:divisor-mixed-mass,ana:exceptional-metric,ana:blowup-kahler-correction,gnpp:noncharging,gnpp:global-measure,gpt:global-measure}
\source{bou02-volume:2.4,bou02-volume:2.9,bou02-volume:4.8}
\lean{TranscendentalNakamaye.Pending.ResolvedNefPart.isNonzeroAlongAt_of_hasAnalyticSingularitiesOfType,TranscendentalNakamaye.Pending.ResolvedNefPart.eq_bot_or_exists_coe,TranscendentalNakamaye.Pending.ResolvedNefPart.bot_or_coe_of_le_of_le,TranscendentalNakamaye.Pending.ResolvedNefPart.logSingularitiesOfSameSingularityType,TranscendentalNakamaye.Pending.ResolvedNefPart.weight_ne_top,TranscendentalNakamaye.Pending.ResolvedNefPart.exists_regularizedResidual,TranscendentalNakamaye.Pending.ResolvedNefPart.coe_add_add_coe_sub,TranscendentalNakamaye.Pending.ResolvedNefPart.exists_patch_pair,TranscendentalNakamaye.Pending.ResolvedNefPart.diagonalPotentialMass_sub_curvature_eq,TranscendentalNakamaye.Pending.ResolvedNefPart.compactSpace_and_kahler_of_isBlowupSequence,TranscendentalNakamaye.Pending.ResolvedNefPart.exists_logResolution_of_pullback_eq,TranscendentalNakamaye.Pending.ResolvedNefPart.exists_resolvedNefPart}
\leanok
Let \((X,\omega)\) be a connected compact K\"ahler manifold of dimension \(n\ge1\), \(\beta\) a smooth closed real \((1,1)\)-form and \(\psi\) a \(\beta\)-psh potential with
analytic singularities of type \((c,\mathfrak a)\), \(c\in\mathbf Q_{>0}\), \(\mathfrak a\) coherent, whose remainders are bounded. There are a modification
\(\mu\colon\widetilde X\to X\) (Definition~\ref{ana:blowups}), a K\"ahler form \(\widetilde\omega\) on \(\widetilde X\), an effective \(\mathbf Q\)-divisor \(D\) and an
effective \(\mathbf R\)-divisor \(F\) on \(\widetilde X\), every component of \(F\) contained in \(\operatorname{Exc}(\mu)\), such that
\begin{enumerate}
\item \(\mu^*\psi\) has logarithmic singularities along \(D\);
\item the class \(b=\mu^*\{\beta\}-\{D\}\) is nef and \(b^n=\int_X\langle\beta_\psi^n\rangle\);
\item \(\{\widetilde\omega\}=\mu^*\{\omega\}-\{F\}\).
\end{enumerate}
\end{lemma}
\begin{proof}
\leanok
\uses{bou:divisor-class}
Here, as in Proposition~\ref{pc:domination-log-resolution}, the class \(\{E\}\) of an effective \(\mathbf Q\)-divisor with a smooth metric is the class of its
curvature form \(\Theta_E\) (Definition~\ref{ana:ideal-sheaves}); that it is also the class of the integration current \([E]\) is Lemma~\ref{bou:divisor-class}.
Since \(\psi\) is not identically \(-\infty\) near any point, at every point one of the local generators of \(\mathfrak a\) in the definition of the type does
not vanish identically nearby, so \(\mathfrak a\) is nonzero at every point. The principalization theorem (Theorem~\ref{hir:principalization}; Boucksom's
Corollary~2.4) gives a modification \(\mu\colon\widetilde X\to X\), a finite sequence of blowups with smooth centres, with
\(\mathfrak a\,\mathcal O_{\widetilde X}=\mathcal O_{\widetilde X}(-G)\) for an effective divisor \(G\); \(\widetilde X\) is compact, connected and K\"ahler
(Lemma~\ref{ana:blowup-kahler-correction}, by induction). Give \(G\) a smooth metric (Lemma~\ref{ana:divisor-metric}) and put \(D=cG\). Then \(\psi_D\) and
\(\psi\circ\mu\) are two weights of the same analytic type \((c,\mathfrak a\mathcal O_{\widetilde X})\) (Lemma~\ref{pc:analytic-type-restriction}) on a compact
manifold, so their difference is bounded: \(\psi\circ\mu\) has logarithmic singularities along \(D\), which is (1). As \(\psi\circ\mu\) is a
\(\mu^*\beta\)-psh potential (Theorem~\ref{psh:holomorphic-pullback} in charts), Lemma~\ref{ana:residual-nef} shows that \(\mu^*\beta-\Theta_D\) is nef, and the
upper regularization \(b^\star\) of the residual across \(\operatorname{Supp}D\) is a bounded \((\mu^*\beta-\Theta_D)\)-psh potential equal to the residual off
\(\operatorname{Supp}D\) (Lemma~\ref{pc:domination-residual}). Hence \(b^n=\int_{\widetilde X}\langle(\mu^*\beta-\Theta_D)_{b^\star}^n\rangle\)
(Lemma~\ref{bou:nef-bounded-mixed-mass}), which equals \(\int_{\widetilde X}\langle(\mu^*\beta)_{\psi\circ\mu}^n\rangle\) by the change of representative off
\(\operatorname{Supp}D\) (Lemma~\ref{bou:divisor-mixed-mass} with every slot corrected: the two measures agree off \(\operatorname{Supp}D\), which is pluripolar and
charged by neither, Theorem~\ref{gnpp:noncharging}), and this is \(\int_X\langle\beta_\psi^n\rangle\) (Lemma~\ref{pc:domination-modification-mass}); in particular
the mass is finite. This is (2). For (3), Lemma~\ref{ana:exceptional-metric} gives an effective \(\mathbf Q\)-divisor \(F\) with a smooth metric, supported in
\(\operatorname{Exc}(\mu)\), such that \(\widetilde\omega=\mu^*\omega-\Theta_F\) is a K\"ahler form; so \(\{\widetilde\omega\}=\mu^*\{\omega\}-\{F\}\), and a
\(\mathbf Q\)-divisor is in particular an \(\mathbf R\)-divisor.
\end{proof}

\begin{lemma}[Bounded potentials in nef classes: the mixed mass is the intersection number]
\label{bou:nef-bounded-mixed-mass}
\uses{bou:nef-mixed-mass,posclass:nef,mmass:mass,mmass:same-type,mmass:slot-linearity,mmass:slot-background-monotone,bou:mixed-mass-class-invariance}
\source{transcendental-nakamaye-copy4-20260911:lem:domination,bou02-volume:4.8}
\lean{TranscendentalNakamaye.Pending.NefBoundedMixedMass.isSemipositive_smul,TranscendentalNakamaye.Pending.NefBoundedMixedMass.thetaPlurisubharmonic_add_of_isSemipositive,TranscendentalNakamaye.Pending.NefBoundedMixedMass.globalMixedMeasure_update_add_smul_eq,TranscendentalNakamaye.Pending.NefBoundedMixedMass.globalMixedMeasure_le_of_add_smul_on,TranscendentalNakamaye.Pending.NefBoundedMixedMass.bumpOn,TranscendentalNakamaye.Pending.NefBoundedMixedMass.thetaPlurisubharmonic_bumpOn,TranscendentalNakamaye.Pending.NefBoundedMixedMass.exists_globalMixedMass_bumpOn_le,TranscendentalNakamaye.Pending.NefBoundedMixedMass.globalMixedMass_add_smul_toReal,TranscendentalNakamaye.Pending.NefBoundedMixedMass.globalMixedMass_toReal_eq_slotIntegral}
\leanok
Let \((M,\omega)\) be a compact K\"ahler manifold of dimension \(N\ge1\), \(\theta_1,\dots,\theta_N\) smooth closed real \((1,1)\)-forms, each nef
(Definition~\ref{posclass:nef}), and \(u_i\) a bounded \(\theta_i\)-psh potential for each \(i\). Then
\[
 \int_M\bigl\langle(\theta_1+dd^cu_1)\wedge\dots\wedge(\theta_N+dd^cu_N)\bigr\rangle=\int_M\theta_1\wedge\dots\wedge\theta_N .
\]
\end{lemma}
\begin{proof}
\leanok
Let \(\varepsilon>0\). Nefness gives smooth \(f_i\) with \(\theta_i+\varepsilon\omega+dd^cf_i\ge0\). In the class of \(\theta_i+\varepsilon\omega\) the potentials \(u_i\)
and \(f_i\) are both bounded, hence of the same singularity type, so the mixed masses of \((\theta_i+\varepsilon\omega,u_i)_i\) and \((\theta_i+\varepsilon\omega,f_i)_i\)
agree (Lemma~\ref{mmass:same-type}); the latter is \(\int_M\prod_i(\theta_i+\varepsilon\omega+dd^cf_i)=\int_M\prod_i(\theta_i+\varepsilon\omega)\)
(Lemmas~\ref{bou:mixed-mass-class-invariance} and~\ref{bou:nef-mixed-mass}). Write \(m(\varepsilon)\) for the mixed mass of \((\theta_i+\varepsilon\omega,u_i)_i\).
Enlarging backgrounds by semipositive forms does not decrease the mixed measure (Lemma~\ref{mmass:slot-background-monotone}), so \(m(0)\le m(\varepsilon)\).
Conversely, by slot linearity (Lemma~\ref{mmass:slot-linearity}) adding \(\varepsilon\omega\) to one slot \(j\) adds exactly \(\varepsilon\) times the mixed measure with
\(\omega\) and the potential \(0\) in slot \(j\); adding it slot by slot gives \(K<\infty\) with \(m(\varepsilon)\le m(0)+\varepsilon K\) for \(0<\varepsilon\le1\).
Since \(\varepsilon\mapsto\int_M\prod_i(\theta_i+\varepsilon\omega)\) is a polynomial (Lemma~\ref{bou:nef-mixed-mass}), letting \(\varepsilon\to0\) in
\(m(0)\le\int_M\prod_i(\theta_i+\varepsilon\omega)\le m(0)+\varepsilon K\) gives the claim.
\end{proof}

\begin{lemma}[Change of representative off a divisor, mixed form]
\label{bou:divisor-mixed-mass}
\uses{ana:ideal-sheaves,mpp:qpsh-poles,gnpp:noncharging,gnpp:global-characterization,mnpp:eqon-bounded,mmass:mass}
\source{transcendental-nakamaye-copy4-20260911:lem:domination,bou02-volume:4.8}
\lean{TranscendentalNakamaye.Pending.DivisorMixedMass.subCurvatureOn,TranscendentalNakamaye.Pending.DivisorMixedMass.exists_patch_pair_on,TranscendentalNakamaye.Pending.DivisorMixedMass.globalMixedMeasure_univ_subCurvatureOn_eq}
\leanok
Let \(M\) be a compact complex manifold of dimension \(n\), \(D\) an effective \(\mathbf Q\)-divisor with a smooth metric (weight \(\psi_D\), curvature \(\Theta_D\),
Definition~\ref{ana:ideal-sheaves}) whose weight is not identically \(-\infty\) on any open set, \(\theta_1,\dots,\theta_n\) smooth closed real \((1,1)\)-forms and
\(S\subset\{1,\dots,n\}\). Put \(\theta'_i=\theta_i-\Theta_D\) for \(i\in S\) and \(\theta'_i=\theta_i\) for \(i\notin S\). Let \(u_i\) be \(\theta'_i\)-psh and \(w_i\) be
\(\theta_i\)-psh with \(w_i=\psi_D+u_i\) (\(i\in S\)) and \(w_i=u_i\) (\(i\notin S\)) off \(\operatorname{Supp}D\). Then the global mixed non-pluripolar measures of
\((\theta'_i,u_i)_i\) and of \((\theta_i,w_i)_i\) have the same total mass.
\end{lemma}
\begin{proof}
\leanok
Off \(\operatorname{Supp}D\) the weight \(\psi_D\) is smooth with \(dd^c\psi_D=-\Theta_D\), so near each such point the corrected local potentials of the two families agree
on a common chart patch, and so do the local mixed measures (Lemma~\ref{mnpp:eqon-bounded}); by the local characterization of the global measure
(Lemma~\ref{gnpp:global-characterization}) the two measures agree on \(M\setminus\operatorname{Supp}D\). \(\operatorname{Supp}D\) is the pole set of the
quasi-psh, locally nondegenerate \(\psi_D\), hence locally pluripolar (Lemma~\ref{mpp:qpsh-poles}), and charged by neither measure (Lemma~\ref{gnpp:noncharging}).
\end{proof}

\begin{lemma}[Demailly--P\u{a}un's Levi estimate for \(\log(\sum_jc_j^2|G_j|^2+s)\)]
\label{bou:concentration-levi}
\uses{psh:local}
\source{dp04-kahler-cone:2.1,bou02-volume:4.11}
\lean{TranscendentalNakamaye.Pending.ConcentrationPotentialLevi.realHessian_mul,TranscendentalNakamaye.Pending.ConcentrationPotentialLevi.realHessian_log_of_pos,TranscendentalNakamaye.Pending.ConcentrationPotentialLevi.leviForm_mul,TranscendentalNakamaye.Pending.ConcentrationPotentialLevi.leviForm_log_of_pos,TranscendentalNakamaye.Pending.ConcentrationPotentialLevi.leviForm_sum,TranscendentalNakamaye.Pending.ConcentrationPotentialLevi.leviForm_re_eq_zero,TranscendentalNakamaye.Pending.ConcentrationPotentialLevi.leviForm_im_eq_zero,TranscendentalNakamaye.Pending.ConcentrationPotentialLevi.leviForm_sq_mul_normSq,TranscendentalNakamaye.Pending.ConcentrationPotentialLevi.fderiv_sq_mul_normSq,TranscendentalNakamaye.Pending.ConcentrationPotentialLevi.normSq_sum_mul_le,TranscendentalNakamaye.Pending.ConcentrationPotentialLevi.le_leviForm_log_sum,TranscendentalNakamaye.Pending.ConcentrationPotentialLevi.neg_le_leviForm_log_sum}
\leanok
Let \(c_1,\dots,c_r\) be real \(\mathcal C^2\) functions and \(G_1,\dots,G_r\) holomorphic functions near a point \(z\) of a finite-dimensional
complex vector space, \(s>0\), and \(S=\sum_jc_j^2|G_j|^2+s\). Write \(L_f(v)=\sum_{a,b}\partial_a\bar\partial_bf\,v_a\bar v_b\) for the Levi form,
\(\partial c(v)\) for the \((1,0)\)-derivative and \(\gamma_j=c_j\,\partial G_j+2G_j\,\partial c_j\). Then for every vector \(v\)
\[
 L_{c_j^2|G_j|^2}(v)=|\gamma_j(v)|^2+2|G_j|^2\bigl(c_jL_{c_j}(v)-|\partial c_j(v)|^2\bigr)
\]
and
\[
 L_{\log S}(v)\ \ge\ \frac{s\sum_j|\gamma_j(v)|^2}{S^2}+\frac{2\sum_j|G_j|^2\bigl(c_jL_{c_j}(v)-|\partial c_j(v)|^2\bigr)}{S}
 \ \ge\ -\frac{2\sum_j|G_j|^2\bigl(|\partial c_j(v)|^2-c_jL_{c_j}(v)\bigr)}{S}.
\]
\end{lemma}
\begin{proof}
\leanok
The real and imaginary parts of a holomorphic function are pluriharmonic, so the product rule for the Levi form gives the first identity. The
logarithm rule \(L_{\log S}=L_S/S-|\partial S|^2/S^2\), the sum rule, and \(\partial(c^2|G|^2)=c\,\bar G\,\gamma\) reduce the second to the
Cauchy--Schwarz inequality \(|\sum_jc_j\bar G_j\gamma_j(v)|^2\le\bigl(\sum_jc_j^2|G_j|^2\bigr)\bigl(\sum_j|\gamma_j(v)|^2\bigr)\), exactly as in
Demailly--P\u{a}un (2.2)--(2.3).
\end{proof}

\begin{lemma}[Demailly--P\u{a}un's Levi estimate off the zero set]
\label{bou:concentration-levi-off-zero}
\uses{bou:concentration-levi}
\source{dp04-kahler-cone:2.1}
\lean{TranscendentalNakamaye.Pending.ConcentrationPotentialLevi.le_leviForm_log_sum_of_pos,TranscendentalNakamaye.Pending.ConcentrationPotentialLevi.neg_le_leviForm_log_sum_of_nonneg}
\leanok
In the notation of Lemma~\ref{bou:concentration-levi}, let \(s\) be real and \(z\) a point with \(S(z)=\sum_jc_j(z)^2|G_j(z)|^2+s>0\). Then the first
inequality of Lemma~\ref{bou:concentration-levi} holds at \(z\), and if \(s\ge0\) so does the second. In particular, for \(s=0\) both hold at every point
where some \(c_jG_j\) does not vanish.
\end{lemma}
\begin{proof}
\leanok
The proof of Lemma~\ref{bou:concentration-levi} uses \(s>0\) only through \(S(z)>0\), which the logarithm rule needs, and \(s\ge0\) only to drop the
nonnegative first term \(s\sum_j|\gamma_j(v)|^2/S^2\).
\end{proof}

\begin{lemma}[Demailly--P\u{a}un's concentration potentials]
\label{bou:concentration-potentials}
\uses{ana:geometric-setting,ana:ideal-sheaves,bou:concentration-levi,bou:concentration-levi-off-zero}
\source{dp04-kahler-cone:2.1,bou02-volume:4.11}
\lean{TranscendentalNakamaye.Pending.SubmanifoldConcentrationPotential.ConcentrationPartition,TranscendentalNakamaye.Pending.SubmanifoldConcentrationPotential.ConcentrationPartition.genNormSq,TranscendentalNakamaye.Pending.SubmanifoldConcentrationPotential.ConcentrationPartition.weight,TranscendentalNakamaye.Pending.SubmanifoldConcentrationPotential.ConcentrationPartition.weight_nonneg,TranscendentalNakamaye.Pending.SubmanifoldConcentrationPotential.ConcentrationPartition.chartSmooth_weight,TranscendentalNakamaye.Pending.SubmanifoldConcentrationPotential.ConcentrationPartition.chartSmooth_log_weight,TranscendentalNakamaye.Pending.SubmanifoldConcentrationPotential.ConcentrationPartition.continuous_weight,TranscendentalNakamaye.Pending.SubmanifoldConcentrationPotential.ConcentrationPartition.weight_eq_zero_of_mem,TranscendentalNakamaye.Pending.SubmanifoldConcentrationPotential.ConcentrationPartition.weight_pos_of_notMem,TranscendentalNakamaye.Pending.SubmanifoldConcentrationPotential.ConcentrationPartition.exists_nhds_genNormSq_le_weight,TranscendentalNakamaye.Pending.SubmanifoldConcentrationPotential.ConcentrationPartition.exists_local_levi_bound,TranscendentalNakamaye.Pending.SubmanifoldConcentrationPotential.ConcentrationPartition.ddc_log_weight_lineValue,TranscendentalNakamaye.Pending.SubmanifoldConcentrationPotential.ConcentrationPartition.exists_pos_isSemipositive_ddc_log_weight,TranscendentalNakamaye.Pending.SubmanifoldConcentrationPotential.ConcentrationPartition.exists_pos_levi_log_weight,TranscendentalNakamaye.Pending.SubmanifoldConcentrationPotential.ConcentrationPartition.potential,TranscendentalNakamaye.Pending.SubmanifoldConcentrationPotential.ConcentrationPartition.chartSmooth_potential,TranscendentalNakamaye.Pending.SubmanifoldConcentrationPotential.ConcentrationPartition.isSemipositive_half_add_ddc_potential,TranscendentalNakamaye.Pending.SubmanifoldConcentrationPotential.ConcentrationPartition.tube,TranscendentalNakamaye.Pending.SubmanifoldConcentrationPotential.ConcentrationPartition.isOpen_tube,TranscendentalNakamaye.Pending.SubmanifoldConcentrationPotential.ConcentrationPartition.tube_mono,TranscendentalNakamaye.Pending.SubmanifoldConcentrationPotential.ConcentrationPartition.iInter_closure_tube_subset,TranscendentalNakamaye.Pending.SubmanifoldConcentrationPotential.flatGen,TranscendentalNakamaye.Pending.SubmanifoldConcentrationPotential.exists_adapted,TranscendentalNakamaye.Pending.SubmanifoldConcentrationPotential.exists_of_isCoherent}
\leanok
Let \(M\) be a compact complex manifold of dimension \(N\) with a K\"ahler form \(\omega\), and \(Y\subseteq M\) a closed set whose ideal sheaf
\(\mathcal I_Y\) (germs of holomorphic functions vanishing on \(Y\)) is coherent.
\begin{enumerate}
\item There are finitely many open sets \(D_a\), holomorphic functions \(g_{a,1},\dots,g_{a,r_a}\) on \(D_a\) generating \(\mathcal I_{Y,y}\) at every
      \(y\in D_a\), and smooth functions \(\theta_a\ge0\) with \(\operatorname{supp}\theta_a\subseteq D_a\) and \(\sum_a\theta_a=1\).
\item For such data the weight \(U=\sum_a\theta_a^2\sum_i|g_{a,i}|^2\) (each inner sum extended by \(0\) off \(D_a\)) is smooth, \(U=0\) on \(Y\) and
      \(U>0\) off \(Y\); near every point \(y_0\) there is \(C>0\) with \(\sum_i|g_{a,i}|^2\le C\,U\) for every \(a\) with \(y_0\in\operatorname{supp}\theta_a\).
\item There is \(A>0\) with \(dd^c\log(U+s)+A\omega\ge0\) for every \(s>0\), and \(dd^c\log U+A\omega\ge0\) at every point where \(U>0\).
\item For \(\varepsilon\ne0\) put \(f_\varepsilon=\log(U+\varepsilon^2)/(2A)\). Then \(\frac12\omega+dd^cf_\varepsilon\ge0\), i.e.\ the form
      \(\omega_\varepsilon=\omega+dd^cf_\varepsilon\) satisfies \(\omega_\varepsilon\ge\frac12\omega\).
\item The tubes \(V_\varepsilon=\{U<\varepsilon^2\}\) are open, increase with \(\varepsilon>0\), and \(\bigcap_{0<\varepsilon\le1}\overline{V_\varepsilon}\subseteq Y\).
\item If \(Y\) is a closed complex submanifold and \(Y\) has dimension \(k\) at \(x\in Y\), there are holomorphic coordinates \(\zeta\) centred at \(x\)
      in which \(Y=\{\zeta_{k+1}=\dots=\zeta_N=0\}\), and data as in (1) with \(U=\sum_{i>k}|\zeta_i|^2\) near \(x\).
\end{enumerate}
\end{lemma}
\begin{proof}
\leanok
(1) Coherence gives, near every point, an open set on which finitely many holomorphic functions generate \(\mathcal I_Y\) at every point; by compactness
finitely many such sets cover \(M\), and a smooth partition of unity subordinate to them gives the \(\theta_a\).
(2) \(U\) is smooth because \(\theta_a\) vanishes near the boundary of \(D_a\). A germ of \(\mathcal I_Y\) at \(y\in Y\) vanishes at \(y\), so \(U(y)=0\).
Off \(Y\) the unit germ lies in \(\mathcal I_{Y,y}\) (\(Y\) is closed), hence in the span of the generators of a piece with \(\theta_a(y)>0\), so one of
them does not vanish at \(y\) and \(U(y)>0\). Near \(y_0\), fix \(a_0\) with \(\theta_{a_0}(y_0)>0\); every generator of a piece with
\(y_0\in\operatorname{supp}\theta_a\subseteq D_a\) is near \(y_0\) a holomorphic combination of the \(g_{a_0,i}\) (both systems generate
\(\mathcal I_{Y,y_0}\); this is Demailly--P\u{a}un's ``by Cartan's theorem A''), so \(\sum_i|g_{a,i}|^2\le C\sum_i|g_{a_0,i}|^2\le C\theta_{a_0}^{-2}U\)
near \(y_0\).
(3) Near a point \(y_0\), in a chart, \(U+s=\sum_{a,i}\theta_a^2|g_{a,i}|^2+s\) with only the pieces whose support contains \(y_0\) contributing, and
Lemma~\ref{bou:concentration-levi} (for \(c=\theta_a\), \(G=g_{a,i}\)) bounds the Levi form of \(\log(U+s)\) below by
\(-2\sum_{a,i}|g_{a,i}|^2(|\partial\theta_a(v)|^2-\theta_aL_{\theta_a}(v))/(U+s)\). By (2) the factors \(|g_{a,i}|^2/(U+s)\) are bounded independently of
\(s>0\), and the first derivatives of \(\theta_a\) and \(L_{\theta_a}\) are bounded by a multiple of \(\omega(v,Jv)\) near \(y_0\); this gives a local
constant, independent of \(s\), and compactness one constant \(A\) on \(M\). At a point where \(U>0\) the Levi form of \(\log(U+s)\) is an explicit
continuous function of \(s\ge0\) (logarithm rule), so the bound passes to \(s=0\) as \(s\to0^+\). In a chart the coefficient of \(dd^c\) on a line is the Levi form divided by
\(\pi\).
(4) \(dd^cf_\varepsilon=(2A)^{-1}dd^c\log(U+\varepsilon^2)\), so \(\frac12\omega+dd^cf_\varepsilon=(2A)^{-1}\bigl(dd^c\log(U+\varepsilon^2)+A\omega\bigr)\ge0\)
by (3) with \(s=\varepsilon^2\).
(5) \(U\) is continuous, so \(V_\varepsilon\) is open and \(\overline{V_\varepsilon}\subseteq\{U\le\varepsilon^2\}\); a point \(y\notin Y\) has \(U(y)>0\), hence lies
outside \(\overline{V_\varepsilon}\) as soon as \(\varepsilon^2<U(y)\).
(6) The submanifold \(Y\) is flattened near \(x\) by holomorphic coordinates \(\zeta\), in which the normal coordinates
\(\zeta_{k+1},\dots,\zeta_N\) generate \(\mathcal I_Y\). Take them as the generators of one piece \(a_0\), with \(\theta_{a_0}=1\) near \(x\) and every
other \(\theta_a\) vanishing near \(x\) (a partition of unity subordinate to a cover in which \(x\) lies only in \(D_{a_0}\)); then
\(U=\sum_{i>k}|\zeta_i|^2\) near \(x\).
\end{proof}
The ideal sheaf of a closed complex submanifold is coherent, and so is the ideal sheaf of a closed analytic set (Cartan); dp04 normalise
\(\sum_a\theta_a^2=1\), which is not needed here because the weight carries the squares. Statement (6) is the one adaptation of the source: the mass
estimate near a point of \(Y\) is then computed with the exact function \(\log(\sum_{i>k}|\zeta_i|^2+\varepsilon^2)\) instead of dp04's eigenvalue
estimates.

\begin{lemma}[Demailly--P\u{a}un: concentration of mass near a submanifold]
\label{bou:mass-concentration}
\uses{ana:geometric-setting,ana:ideal-sheaves,bou:concentration-levi,bou:concentration-potentials,bou:smooth-mixed-density,bou:product-manifold,fpmma:geometric-density,fpmma:coefficient-positive}
\source{dp04-kahler-cone:2.1,bou02-volume:4.11}
\lean{TranscendentalNakamaye.Pending.SubmanifoldMassConcentration.re_mixedCoefficient_mono,TranscendentalNakamaye.Pending.SubmanifoldMassConcentration.modelSlot,TranscendentalNakamaye.Pending.SubmanifoldMassConcentration.pow_mul_pow_le_re_mixedCoefficient_modelSlot,TranscendentalNakamaye.Pending.SubmanifoldMassConcentration.posSemidef_normalizedHermitianMatrix_sub_diagonal,TranscendentalNakamaye.Pending.SubmanifoldMassConcentration.le_actualMixedTopDensity_of_model,TranscendentalNakamaye.Pending.SubmanifoldMassConcentration.scaleNormal,TranscendentalNakamaye.Pending.SubmanifoldMassConcentration.normalSq,TranscendentalNakamaye.Pending.SubmanifoldMassConcentration.normalSq_scaleNormal,TranscendentalNakamaye.Pending.SubmanifoldMassConcentration.le_leviForm_log_normalSq,TranscendentalNakamaye.Pending.SubmanifoldMassConcentration.det_restrictScalars_nonneg,TranscendentalNakamaye.Pending.SubmanifoldMassConcentration.lineValue_comp,TranscendentalNakamaye.Pending.SubmanifoldMassConcentration.exists_flat_ball,TranscendentalNakamaye.Pending.SubmanifoldMassConcentration.concentrationSlot,TranscendentalNakamaye.Pending.SubmanifoldMassConcentration.exists_mass_lower_bound,TranscendentalNakamaye.Pending.SubmanifoldMassConcentration.thetaPlurisubharmonic_concentrationSlot,TranscendentalNakamaye.Pending.SubmanifoldMassConcentration.exists_massConcentration,TranscendentalNakamaye.Pending.SubmanifoldMassConcentration.isClosedSubmanifold_diagonal,TranscendentalNakamaye.Pending.SubmanifoldMassConcentration.exists_diagonal_massConcentration}
\leanok
\begin{enumerate}
\item Let \((M,\omega)\) be a compact K\"ahler manifold of dimension \(N\), \(Y\subset M\) a closed complex submanifold and \(x_0\in Y\) a point at which \(Y\)
      has dimension \(N-p\). There are smooth functions \(f_\varepsilon\) and open sets \(V_\varepsilon\), increasing with \(\varepsilon>0\), with
      \(\bigcap_{0<\varepsilon\le1}\overline{V_\varepsilon}\subseteq Y\), such that for \(0<\varepsilon\le1\) the forms \(\omega_\varepsilon=\omega+dd^cf_\varepsilon\),
      cohomologous to \(\omega\), satisfy \(\omega_\varepsilon\ge\frac12\omega\), and for every neighbourhood \(U\) of \(x_0\) there is \(\delta(U)>0\) with
      \[
       \int_{U\cap V_\varepsilon}\omega_\varepsilon^p\wedge\omega^{N-p}\ge\delta(U)\qquad(0<\varepsilon\le1),
      \]
      the left-hand side being the global mixed measure of \((\omega_\varepsilon,\dots,\omega_\varepsilon,\omega,\dots,\omega)\) (\(p\) times \(\omega_\varepsilon\)) with the
      potentials \(0\), of \(U\cap V_\varepsilon\).
\item In particular, if \(X\) is a compact K\"ahler manifold of dimension \(n\) with K\"ahler form \(\omega\), then (1) holds on \(X\times X\) with the K\"ahler
      form \(p_1^*\omega+p_2^*\omega\), for the diagonal \(\Delta\), a closed submanifold of dimension \(n\), at every point \((x,x)\), with \(p=n\).
\end{enumerate}
\end{lemma}
\begin{proof}
\leanok
(1) Take the data of Lemma~\ref{bou:concentration-potentials} adapted to \(x_0\) (its part (6), with \(k=N-p\)), its constant \(A\) and
\(f_\varepsilon=\log(W+\varepsilon^2)/(2A)\), where \(W\) is the weight; then \(\omega_\varepsilon\ge\frac12\omega\) by part (4), and the tubes
\(V_\varepsilon=\{W<\varepsilon^2\}\) have the stated properties by part (5). In Demailly--P\u{a}un's notation, \(\psi=\frac12\log W\) has logarithmic poles along
\(Y\), \(\psi_\varepsilon=\frac12\log(W+\varepsilon^2)\) decreases to \(\psi\), and \(\omega_\varepsilon=\omega+(2A')^{-1}dd^c\psi_\varepsilon\) with \(A'=A/2\).

For the mass, work in the flat coordinates \(\zeta=(\zeta',\zeta'')\) of part (6) on a ball \(B\) around \(x_0\) inside \(U\), on which \(W=|\zeta''|^2\) and
\(\omega\ge c\,|v|^2\) on tangent vectors \(v\) (written in \(\zeta\)). Let \(T_\varepsilon(\xi',\xi'')=(\xi',\varepsilon\xi'')\) and \(F_\varepsilon=\zeta^{-1}\circ T_\varepsilon\). On the
fixed set \(B_0=\{|\xi'|<r,\ |\xi''|<1\}\) (with \(r\) small, so that \(F_\varepsilon(B_0)\subseteq B\cap V_\varepsilon\) for \(0<\varepsilon\le1\)) one has
\(f_\varepsilon\circ F_\varepsilon=\bigl(\log\varepsilon^2+\log(|\xi''|^2+1)\bigr)/(2A)\), so \(F_\varepsilon^*dd^cf_\varepsilon=(2A)^{-1}dd^c\log(|\xi''|^2+1)\), whose Levi form is
\(\ge|v''|^2/(8A)\) on \(B_0\) independently of \(\varepsilon\); and \(F_\varepsilon^*\omega\ge c\,|v'|^2\) because \(T_\varepsilon\) does not move the tangential directions.
By the change of variables formula and Lemma~\ref{bou:smooth-mixed-density}, the mass of \(U\cap V_\varepsilon\) is at least the integral over \(B_0\) of the density
of \((F_\varepsilon^*\omega_\varepsilon)^p\wedge(F_\varepsilon^*\omega)^{N-p}\) (the Jacobian of the complex-linear \(T_\varepsilon\) is \(\ge0\) and is absorbed by the
pull-back). The mixed coefficient is monotone in each positive semidefinite slot (Lemmas~\ref{fpmma:geometric-density} and~\ref{fpmma:coefficient-positive}), and
the slots dominate the model diagonal matrices \(a\,\mathbf 1''\) (\(p\) times, \(a=1/(8A)\)) and \(b\,\mathbf 1'\) (\(N-p\) times, \(b=c\)), whose mixed coefficient is
\(\ge a^pb^{N-p}\), up to the normalising constant. So \(\delta(U)\) can be taken a fixed positive multiple of \(a^pb^{N-p}\) times the volume of \(B_0\).

(2) The diagonal is a closed connected submanifold of dimension \(n\) (Lemma~\ref{bou:product-manifold}), so (1) applies at \((x,x)\) with \(N=2n\) and \(p=n\).
\end{proof}
This replaces Demailly--P\u{a}un's eigenvalue estimates by the exact rescaling of the model \(\log(|\zeta''|^2+\varepsilon^2)\); the price is that the data
depend on the point \(x_0\), which is harmless because the uniform nef core (Lemma~\ref{bou:uniform-nef-core}) fixes the point first.

\begin{lemma}[Bounded H\"older spaces]
\label{bou:holder-spaces}
\lean{TranscendentalNakamaye.Pending.BoundedHolderMap,TranscendentalNakamaye.Pending.BoundedHolderMap.instCompleteSpace,TranscendentalNakamaye.Pending.BoundedHolderMap.holderWith_mul,TranscendentalNakamaye.Pending.BoundedHolderMap.norm_mul_le',TranscendentalNakamaye.Pending.BoundedHolderMap.instNormedRing,TranscendentalNakamaye.Pending.BoundedHolderMap.instNormedAlgebra,TranscendentalNakamaye.Pending.BoundedHolderMap.holderWith_interpolate,TranscendentalNakamaye.Pending.BoundedHolderMap.ofLE,TranscendentalNakamaye.Pending.BoundedHolderMap.ofLECLM,TranscendentalNakamaye.Pending.BoundedHolderMap.isCompactOperator_ofLE,TranscendentalNakamaye.Pending.BoundedHolderMap.instNormedSpace,TranscendentalNakamaye.Pending.BoundedHolderMap.instNormedAddCommGroup}
\leanok
For \(r\ge0\), a metric space \(X\) and a normed space \(F\), the bounded \(r\)-H\"older continuous maps \(X\to F\) with the norm \(\|f\|_\infty+[f]_r\)
(\([f]_r\) the least H\"older constant) form a normed space \(C^{0,r}_b(X,F)\), complete when \(F\) is; for a normed algebra \(R\), \(C^{0,r}_b(X,R)\) is a normed algebra,
\([fg]_r\le\|f\|_\infty[g]_r+[f]_r\|g\|_\infty\); the interpolation inequality \([f]_{rt}\le[f]_r^{\,t}(2\|f\|_\infty)^{1-t}\) holds for \(0\le t\le1\); and for \(X\)
compact, \(F\) proper (e.g.\ finite-dimensional) and \(s<r\), the inclusion \(C^{0,r}(X,F)\to C^{0,s}(X,F)\) is a compact operator.
\end{lemma}
\begin{proof}
\leanok
Completeness: a Cauchy sequence converges uniformly, and the H\"older bound passes to the pointwise limit. The product estimate is the triangle inequality on
\(f(x)g(x)-f(y)g(y)\). Interpolation: \(|f(x)-f(y)|\le\min([f]_rd^r,2\|f\|_\infty)\le([f]_rd^r)^t(2\|f\|_\infty)^{1-t}\). Compactness: a bounded set of
\(C^{0,r}\) is equicontinuous and uniformly bounded, so relatively compact in \(C^0\) (Arzel\`a--Ascoli), and by interpolation \(C^0\)-convergence of a sequence
bounded in \(C^{0,r}\) implies \(C^{0,s}\)-convergence.
\end{proof}
This is the coefficient space of the Schauder theory planned for Theorem~\ref{bou:yau}.

\begin{lemma}[Bounded H\"older spaces of differentiable maps]
\label{bou:holder-spaces-higher}
\uses{bou:holder-spaces}
\lean{TranscendentalNakamaye.Pending.hasFDerivAt_of_tendstoUniformly_iteratedFDeriv,TranscendentalNakamaye.Pending.iteratedFDeriv_eq_of_tendstoUniformly,TranscendentalNakamaye.Pending.contDiff_of_tendstoUniformly_iteratedFDeriv,TranscendentalNakamaye.Pending.BoundedContDiffHolderMap,TranscendentalNakamaye.Pending.BoundedContDiffHolderMap.instNormedAddCommGroup,TranscendentalNakamaye.Pending.BoundedContDiffHolderMap.instNormedSpace,TranscendentalNakamaye.Pending.BoundedContDiffHolderMap.instCompleteSpace,TranscendentalNakamaye.Pending.BoundedContDiffHolderMap.iteratedFDerivBCF,TranscendentalNakamaye.Pending.BoundedContDiffHolderMap.iteratedFDerivHolder,TranscendentalNakamaye.Pending.BoundedContDiffHolderMap.holderWith_iteratedFDeriv,TranscendentalNakamaye.Pending.BoundedContDiffHolderMap.bilinear,TranscendentalNakamaye.Pending.BoundedContDiffHolderMap.norm_bilinear_le,TranscendentalNakamaye.Pending.BoundedHolderMap.holderWith_of_le,TranscendentalNakamaye.Pending.BoundedContDiffHolderMap.ofLE,TranscendentalNakamaye.Pending.BoundedContDiffHolderMap.ofLECLM,TranscendentalNakamaye.Pending.BoundedContDiffHolderMap.norm_ofLE_le,TranscendentalNakamaye.Pending.BoundedContDiffHolderMap.holderWith_iteratedFDeriv_interpolate,TranscendentalNakamaye.Pending.finiteDimensional_continuousMultilinearMap,TranscendentalNakamaye.Pending.properSpace_continuousMultilinearMap,TranscendentalNakamaye.Pending.iteratedFDeriv_eq_zero_of_notMem,TranscendentalNakamaye.Pending.BoundedContDiffHolderMap.exists_subseq_tendsto_ofLE,TranscendentalNakamaye.Pending.BoundedContDiffHolderMap.supportedIn,TranscendentalNakamaye.Pending.BoundedContDiffHolderMap.mem_supportedIn,TranscendentalNakamaye.Pending.BoundedContDiffHolderMap.isCompactOperator_ofLE_supportedIn,TranscendentalNakamaye.Pending.BoundedContDiffHolderMap.ofContDiff,TranscendentalNakamaye.Pending.BoundedContDiffHolderMap.norm_ofContDiff_le,TranscendentalNakamaye.Pending.BoundedContDiffHolderMap.ofHasCompactSupport,TranscendentalNakamaye.Pending.BoundedContDiffHolderMap.fderivMap,TranscendentalNakamaye.Pending.BoundedContDiffHolderMap.norm_fderivMap_le,TranscendentalNakamaye.Pending.BoundedContDiffHolderMap.ofFDeriv,TranscendentalNakamaye.Pending.BoundedContDiffHolderMap.norm_ofFDeriv_le,TranscendentalNakamaye.Pending.BoundedContDiffHolderMap.exists_comp,TranscendentalNakamaye.Pending.BoundedContDiffHolderMap.comp,TranscendentalNakamaye.Pending.BoundedContDiffHolderMap.compCLM,TranscendentalNakamaye.Pending.BoundedContDiffHolderMap.exists_smul_comp_of_contDiffOn}
\leanok
(1) If \(C^k\) maps \(f_j\colon E\to F\) converge pointwise to \(f\) and each derivative \(D^if_j\), \(i\le k\), converges uniformly, then \(f\) is \(C^k\) and \(D^if\) is the
limit of the \(D^if_j\). (2) For \(k\in\mathbf N\), \(r\ge0\) and real normed spaces \(E,F\), the \(C^k\) maps with bounded derivatives \(D^if\) (\(i\le k\)) and
\(r\)-H\"older \(D^kf\) form a normed space \(C^{k,r}_b(E,F)\) with \(\|f\|=\sum_{i\le k}\|D^if\|_\infty+[D^kf]_r\), complete when \(F\) is. (3) For \(r\le1\) and a continuous
bilinear \(B\), \(\|B(f,g)\|\le2^{k+3}\|B\|\,\|f\|\,\|g\|\). (4) For \(s\le r\) the inclusion \(C^{k,r}_b\to C^{k,s}_b\) is continuous of norm \(\le3\), and the top derivative
satisfies the interpolation inequality. (5) For finite-dimensional \(E\), \(F\), \(s<r\le1\) and a compact \(K\subset E\), the inclusion
\(C^{k,r}_b\to C^{k,s}_b\) restricted to the maps supported in \(K\) is a compact operator. (6) For \(r\le1\) and \(\tau\) of class \(C^{k+1}\) with bounded
derivatives of orders \(1,\dots,k+1\), \(g\mapsto g\circ\tau\) is a bounded linear operator \(C^{k,r}_b\to C^{k,r}_b\); if \(\tau\) is \(C^{k+1}\) only on an open set \(U\)
of a finite-dimensional space and \(\sigma\) is a \(C^{k+1}\) cut-off with compact support in \(U\), then \(\sigma\cdot(g\circ\tau)\in C^{k,r}_b\) with norm \(\le C\|g\|\).
\end{lemma}
\begin{proof}
\leanok
(1) Induction on \(i\) from the first-derivative case of uniform limits. (2) A Cauchy sequence has uniformly convergent derivatives and a H\"older bound that passes
to the limit; (1) identifies the limit. (3) Leibniz bound for \(B(f\circ\tau_v-f,g\circ\tau_v)+B(f,g\circ\tau_v-g)\), \(\tau_v(z)=z+v\), with
\(\|D^jf(x+v)-D^jf(x)\|\le2\|f\|\,|v|^r\). (4) \(d^r\le d^s\) for \(d\le1\) and the bound \(2\|D^kf\|_\infty\) for \(d\ge1\); interpolation as in Lemma~\ref{bou:holder-spaces}. (5) The derivatives of a bounded set of maps supported in \(K\) are equicontinuous and
vanish off \(K\); Arzel\`a--Ascoli gives a subsequence converging uniformly with all derivatives up to order \(k\), and interpolation upgrades the convergence of
\(D^kf\) to \(C^{0,s}\). (6) Induction on \(k\) through \(D(g\circ\tau)=B((Dg)\circ\tau,D\tau)\) with \(B\) the composition of linear maps and (3); the case \(k=0\) is
\([g\circ\tau]_r\le[g]_r\operatorname{Lip}(\tau)^r\).
\end{proof}

\begin{lemma}[Interpolation inequalities on the whole space]
\label{bou:holder-interpolation}
\uses{bou:holder-spaces-higher}
\lean{TranscendentalNakamaye.Pending.HolderInterpolation.norm_sub_sub_fderiv_le_of_holderWith,TranscendentalNakamaye.Pending.HolderInterpolation.norm_fderiv_le_of_holderWith,TranscendentalNakamaye.Pending.HolderInterpolation.holderWith_one_fderiv_of_bound,TranscendentalNakamaye.Pending.HolderInterpolation.norm_fderiv_fderiv_le_interpolation,TranscendentalNakamaye.Pending.HolderInterpolation.norm_fderiv_le_interpolation,TranscendentalNakamaye.Pending.HolderInterpolation.iteratedFDeriv_two_apply_eq,TranscendentalNakamaye.Pending.HolderInterpolation.exists_iteratedFDeriv_two_eq,TranscendentalNakamaye.Pending.HolderInterpolation.norm_iteratedFDeriv_two,TranscendentalNakamaye.Pending.HolderInterpolation.norm_iteratedFDeriv_two_sub,TranscendentalNakamaye.Pending.HolderInterpolation.holderWith_fderiv_fderiv_of_iteratedFDeriv,TranscendentalNakamaye.Pending.HolderInterpolation.norm_iteratedFDeriv_two_le_interpolation,TranscendentalNakamaye.Pending.HolderInterpolation.norm_iteratedFDeriv_one_le_interpolation}
\leanok
Let \(E,F\) be real normed spaces and \(f\colon E\to F\) differentiable.
(1) If \(Df\) is \(r\)-H\"older with constant \(H\), then \(\|f(x+h)-f(x)-Df(x)h\|\le H\|h\|^{1+r}\), and if moreover \(\|f\|\le A\) on \(E\), then
\(\|Df(x)\|\le2A/\rho+H\rho^r\) for every \(\rho>0\).
(2) If \(f\) is \(C^2\) with \(\|f\|\le A\), \(D^2f\) bounded and \(\alpha\)-H\"older with constant \(H\), then for every \(\rho>0\) and \(x\in E\)
\[
 \|D^2f(x)\|\le\frac{32A}{\rho^2}+2H\rho^\alpha,\qquad \|Df(x)\|\le\frac{34A}{\rho}+2H\rho^{1+\alpha}.
\]
(3) In particular, for \(u\in C^{2,\alpha}_b(E,F)\) (Lemma~\ref{bou:holder-spaces-higher}) the same bounds hold with \(A=\|u\|_\infty\) and \(H=[D^2u]_\alpha\).
\end{lemma}
\begin{proof}
\leanok
(1) The map \(g(y)=f(y)-Df(x)(y-x)\) has derivative \(Df(y)-Df(x)\), of norm \(\le H\|y-x\|^r\), so the mean value inequality on the segment from \(x\) to \(x+h\)
gives the Taylor bound. For \(\|h\|=\rho\), \(\|Df(x)h\|\le\|f(x+h)\|+\|f(x)\|+H\rho^{1+r}\le2A+H\rho^{1+r}\); divide by \(\rho\).
(2) Let \(S=\sup\|D^2f\|<\infty\). As \(Df\) is \(S\)-Lipschitz, (1) at the scale \(\rho/4\) gives \(\|Df\|\le8A/\rho+S\rho/4\); (1) applied to \(Df\)
(with the \(\alpha\)-H\"older \(D^2f\)) at the scale \(\rho\) gives \(\|D^2f\|\le2(8A/\rho+S\rho/4)/\rho+H\rho^\alpha=16A/\rho^2+S/2+H\rho^\alpha\); taking the supremum,
\(S\le32A/\rho^2+2H\rho^\alpha\). Then (1) for \(f\), whose derivative is Lipschitz with this constant, gives \(\|Df\|\le2A/\rho+S\rho\le34A/\rho+2H\rho^{1+\alpha}\).
(3) The second derivative as a bilinear map and as a symmetric multilinear map differ by a linear isometry, so norms and H\"older constants agree.
\end{proof}
These are the whole-space interpolation inequalities \(|D^jf|_0\le\varepsilon[D^2f]_\alpha+C(\varepsilon)|f|_0\), \(j=1,2\). In the Schauder estimate they pass
from a bound with the \(C^2\)-norm on the right to one with the \(C^0\)-norm; Sz\'ekelyhidi's sketch of the Schauder estimate (his Theorem~2.8) does this
step with a second blow-up argument instead.

\begin{lemma}[Mollification in bounded H\"older spaces]
\label{bou:holder-mollification}
\uses{bou:holder-spaces,bou:holder-spaces-higher}
\lean{TranscendentalNakamaye.Pending.HolderMollification.integrable_smul_comp_sub,TranscendentalNakamaye.Pending.HolderMollification.norm_convolution_le,TranscendentalNakamaye.Pending.HolderMollification.dist_convolution_le,TranscendentalNakamaye.Pending.HolderMollification.holderWith_convolution,TranscendentalNakamaye.Pending.HolderMollification.contDiff_convolution,TranscendentalNakamaye.Pending.HolderMollification.hasFDerivAt_convolution_of_bounded,TranscendentalNakamaye.Pending.HolderMollification.convolution_comp_linearIsometryEquiv,TranscendentalNakamaye.Pending.HolderMollification.iteratedFDeriv_convolution,TranscendentalNakamaye.Pending.HolderMollification.norm_convolution_sub_le,TranscendentalNakamaye.Pending.HolderMollification.BoundedContDiffHolderMap.mollify,TranscendentalNakamaye.Pending.HolderMollification.BoundedContDiffHolderMap.coe_mollify,TranscendentalNakamaye.Pending.HolderMollification.BoundedContDiffHolderMap.iteratedFDeriv_mollify,TranscendentalNakamaye.Pending.HolderMollification.BoundedContDiffHolderMap.norm_mollify_le,TranscendentalNakamaye.Pending.HolderMollification.BoundedContDiffHolderMap.norm_iteratedFDeriv_mollify_sub_le_of_lt,TranscendentalNakamaye.Pending.HolderMollification.BoundedContDiffHolderMap.norm_iteratedFDeriv_mollify_sub_le}
\leanok
Let \(E\) be a finite-dimensional real normed space with a Haar measure, \(F\) a real Banach space, \(\varphi\ge0\) a continuous integrable function on \(E\)
with \(\int_E\varphi=1\), \(u\colon E\to F\) continuous, and \(\varphi\star u(x)=\int_E\varphi(t)\,u(x-t)\,dt\).
\begin{enumerate}
\item If \(\|u\|\le A\), then \(\|\varphi\star u\|\le A\). If \(u\) is bounded and \(r\)-H\"older with constant \(H\), then so is \(\varphi\star u\), and if \(\varphi\) vanishes
      outside the ball of radius \(R\), then \(\|\varphi\star u-u\|\le HR^r\).
\item If \(\varphi\) is of class \(C^k\) with compact support, so is \(\varphi\star u\) (for bounded \(u\)); if \(u\) is of class \(C^k\) with bounded derivatives of
      order \(\le k\), then \(D^j(\varphi\star u)=\varphi\star D^ju\) for \(j\le k\).
\item If \(\varphi\) is smooth with compact support and \(u\in C^{k,r}_b(E,F)\) (Lemma~\ref{bou:holder-spaces-higher}), then \(\varphi\star u\in C^{k,r}_b(E,F)\) with
      \(\|\varphi\star u\|\le\|u\|\); if moreover \(\varphi\) vanishes outside the ball of radius \(R\ge0\), then \(\|D^i(\varphi\star u)-D^iu\|_\infty\le\|u\|R\) for
      \(i<k\) and \(\|D^k(\varphi\star u)-D^ku\|_\infty\le\|u\|R^r\).
\end{enumerate}
\end{lemma}
\begin{proof}
\leanok
(1) \(\varphi\star u(x)\) is an average of values of \(u\), and \(\varphi\star u(x)-\varphi\star u(y)=\int\varphi(t)(u(x-t)-u(y-t))\,dt\); similarly
\(\varphi\star u(x)-u(x)=\int\varphi(t)(u(x-t)-u(x))\,dt\) with \(\|t\|\le R\) on the support. (2) Differentiation under the integral sign, the derivatives of \(u\)
being bounded; for the first claim differentiate the compactly supported factor. (3) By (2) the derivatives of \(\varphi\star u\) are the mollified derivatives
of \(u\), so (1) bounds their sup norms and the H\"older seminorm of the top one; for \(i<k\), \(D^iu\) is \(\|u\|\)-Lipschitz.
\end{proof}
These are the regularity inputs of the Schauder estimates (mollifying a solution of a constant-coefficient equation).

\begin{lemma}[The mean value identity for constant-coefficient elliptic operators]
\label{bou:elliptic-mean-value}
\uses{bou:holder-spaces-higher}
\source{szekelyhidi-extremal:page-0041}
\lean{TranscendentalNakamaye.Pending.EllipticMeanValue.hessTrace,TranscendentalNakamaye.Pending.EllipticMeanValue.contDiff_deriv_smoothTransition,TranscendentalNakamaye.Pending.EllipticMeanValue.deriv_smoothTransition_nonneg,TranscendentalNakamaye.Pending.EllipticMeanValue.deriv_smoothTransition_of_neg,TranscendentalNakamaye.Pending.EllipticMeanValue.exists_deriv_smoothTransition_pos,TranscendentalNakamaye.Pending.EllipticMeanValue.kernelPotential,TranscendentalNakamaye.Pending.EllipticMeanValue.meanValueKernel,TranscendentalNakamaye.Pending.EllipticMeanValue.meanValueKernel_nonneg,TranscendentalNakamaye.Pending.EllipticMeanValue.contDiff_quadForm,TranscendentalNakamaye.Pending.EllipticMeanValue.contDiff_meanValueKernel,TranscendentalNakamaye.Pending.EllipticMeanValue.contDiff_kernelPotential,TranscendentalNakamaye.Pending.EllipticMeanValue.hasFDerivAt_quadForm,TranscendentalNakamaye.Pending.EllipticMeanValue.hasFDerivAt_kernelPotential,TranscendentalNakamaye.Pending.EllipticMeanValue.meanValueKernel_eq_zero,TranscendentalNakamaye.Pending.EllipticMeanValue.kernelPotential_eq_zero,TranscendentalNakamaye.Pending.EllipticMeanValue.norm_le_of_quadForm_le_one,TranscendentalNakamaye.Pending.EllipticMeanValue.meanValueKernel_eq_zero_of_norm,TranscendentalNakamaye.Pending.EllipticMeanValue.kernelPotential_eq_zero_of_norm,TranscendentalNakamaye.Pending.EllipticMeanValue.hasCompactSupport_meanValueKernel,TranscendentalNakamaye.Pending.EllipticMeanValue.hasCompactSupport_kernelPotential,TranscendentalNakamaye.Pending.EllipticMeanValue.integral_meanValueKernel_pos,TranscendentalNakamaye.Pending.EllipticMeanValue.integral_comp_sub_smul_mul_meanValueKernel}
\leanok
Let \(E\) be a finite-dimensional real inner product space with an orthonormal basis \((b_i)\) and a Haar measure, and \(A\), \(B\) operators on \(E\) with \(AB=1\),
\(B\) symmetric and \(\langle By,y\rangle\ge c|y|^2\) for some \(c>0\). Write \(A:D^2G(x)=\sum_iD^2G(x)(b_i,Ab_i)\). There is an explicit smooth, nonnegative, compactly
supported kernel \(\eta=\eta_B\), \(\eta(y)=2\psi(1-\langle By,y\rangle)\) with \(\psi\) the derivative of a smooth transition function, such that: if \(G\) is \(C^2\) with
\(A:D^2G\equiv0\), then \(\int_EG(x-ry)\,\eta(y)\,dy=\bigl(\int_E\eta\bigr)G(x)\) for every \(x\in E\) and \(r\in\mathbf R\); and \(\int_E\eta>0\) when \(E\ne0\).
\end{lemma}
\begin{proof}
\leanok
With \(H(y)=\psi_0(1-\langle By,y\rangle)\), \(\psi_0\) the smooth transition, one has \(y\,\eta(y)=-A\nabla H(y)\). Hence
\(\frac{d}{dr}\int G(x-ry)\eta(y)\,dy=-\int DG(x-ry)(y)\,\eta(y)\,dy=\sum_i\int DG(x-ry)(Ab_i)\,\partial_iH(y)\,dy=r\int(A:D^2G)(x-ry)H(y)\,dy=0\), by one integration by
parts per basis vector; at \(r=0\) both sides agree.
\end{proof}
This is Sz\'ekelyhidi's Corollary~2.2 for the operator \(A:D^2\), obtained with an \(A\)-adapted kernel instead of spheres.

\begin{lemma}[Liouville's theorem with H\"older growth; rigidity of the blow-up limit]
\label{bou:elliptic-liouville}
\uses{bou:elliptic-mean-value}
\source{szekelyhidi-extremal:page-0042,szekelyhidi-extremal:page-0046}
\lean{TranscendentalNakamaye.Pending.EllipticLiouville.add_eq_of_meanValue_of_holder,TranscendentalNakamaye.Pending.EllipticLiouville.hasFDerivAt_comp_const_add,TranscendentalNakamaye.Pending.EllipticLiouville.abs_secondDiff_sub_le,TranscendentalNakamaye.Pending.EllipticLiouville.fderiv_fderiv_eq_of_hessTrace_eq_zero}
\leanok
Let \(E\) be a finite-dimensional real inner product space with a Haar measure.
(1) Let \(\eta\) be continuous, Lipschitz and compactly supported with \(\int\eta\ne0\), and \(w\colon E\to\mathbf R\) continuous with
\(\int w(x-ry)\eta(y)\,dy=(\int\eta)\,w(x)\) for all \(x\) and \(r>0\), and \(|w(x)-w(x')|\le K|x-x'|^\alpha\) with \(0\le\alpha<1\). Then \(w\) is constant.
(2) If \(G\) is \(C^2\), \(A:D^2G\equiv0\) for a symmetric \(A\) with \(\langle Av,v\rangle\ge\lambda|v|^2\), \(\lambda>0\), and \(D^2G\) is \((K,\alpha)\)-H\"older with
\(0\le\alpha<1\), then \(D^2G\) is constant.
\end{lemma}
\begin{proof}
\leanok
(1) Translating the kernel by \(v/r\), \((\int\eta)(w(x+v)-w(x))=\int(w(x-ry)-w(x))(\eta(y+v/r)-\eta(y))\,dy=O(r^{\alpha-1})\to0\) as \(r\to\infty\). (2) The second
differences \(G(x+h+k)-G(x+h)-G(x+k)+G(x)\) have the mean value property (Lemma~\ref{bou:elliptic-mean-value} with \(B=A^{-1}\)) and H\"older growth
\(K|x-x'|^\alpha|h||k|\), so they are constant in \(x\) by (1); differentiating twice in \(h\), \(k\) gives the claim.
\end{proof}
This is Sz\'ekelyhidi's Corollary~2.4 in the form used at the end of his proof of Theorem~2.8.

\begin{lemma}[Compactness for blow-up sequences]
\label{bou:holder-blowup-compactness}
\uses{bou:holder-spaces-higher}
\source{szekelyhidi-extremal:page-0044,szekelyhidi-extremal:page-0046}
\lean{TranscendentalNakamaye.Pending.HolderBlowupCompactness.tendstoLocallyUniformly_of_equicontinuous,TranscendentalNakamaye.Pending.HolderBlowupCompactness.exists_tendsto_of_norm_le,TranscendentalNakamaye.Pending.HolderBlowupCompactness.exists_contDiff_two_tendsto_fderiv_fderiv}
\leanok
Let \(E\) be a finite-dimensional real normed space and \(w_i\colon E\to\mathbf R\) functions of class \(C^2\) with \(D^2w_i(0)=0\) and \(D^2w_i\) uniformly
\((K,\alpha)\)-H\"older, \(\alpha>0\). Along every ultrafilter on the index set there is a \(C^2\) function \(V\) with \(D^2w_i(x)\to D^2V(x)\) for every \(x\). (Also: pointwise
convergence of an equicontinuous family on a locally compact space is locally uniform, and bounded families in a finite-dimensional space converge
along ultrafilters.)
\end{lemma}
\begin{proof}
\leanok
Subtract the first-order Taylor polynomials at \(0\): then \(|D^2v_i(x)|\le K|x|^\alpha\), \(|Dv_i(x)|\le K|x|^{1+\alpha}\), \(|v_i(x)|\le K|x|^{2+\alpha}\), so the pointwise
limits along the ultrafilter exist (closed balls are compact), the families \(Dv_i\), \(D^2v_i\) are equicontinuous, so the convergence is locally uniform
(Arzel\`a--Ascoli), and limits of derivatives are derivatives of the limit.
\end{proof}
Ultrafilter limits replace the diagonal subsequences of Sz\'ekelyhidi's proof of Theorem~2.8 (and his Theorem~2.7).

\begin{lemma}[H\"older spaces on a compact complex manifold]
\label{bou:holder-spaces-manifold}
\uses{bou:holder-spaces-higher,ana:geometric-setting}
\lean{TranscendentalNakamaye.Pending.HolderAtlas,TranscendentalNakamaye.Pending.HolderAtlas.nonempty,TranscendentalNakamaye.Pending.HolderAtlas.localRep,TranscendentalNakamaye.Pending.HolderAtlas.localRep_chart,TranscendentalNakamaye.Pending.HolderAtlas.eq_sum_localRep,TranscendentalNakamaye.Pending.HolderAtlas.isCompact_image_tsupport,TranscendentalNakamaye.Pending.HolderAtlas.localRep_eq_zero_of_notMem,TranscendentalNakamaye.Pending.HolderAtlas.HolderMap,TranscendentalNakamaye.Pending.HolderAtlas.HolderMap.rep,TranscendentalNakamaye.Pending.HolderAtlas.HolderMap.norm_def,TranscendentalNakamaye.Pending.HolderAtlas.HolderMap.norm_rep_le,TranscendentalNakamaye.Pending.HolderAtlas.HolderMap.instNormedAddCommGroup,TranscendentalNakamaye.Pending.HolderAtlas.HolderMap.instNormedSpace,TranscendentalNakamaye.Pending.HolderAtlas.HolderMap.exists_tendsto_of_tendsto_rep,TranscendentalNakamaye.Pending.HolderAtlas.HolderMap.instCompleteSpace,TranscendentalNakamaye.Pending.HolderAtlas.HolderMap.localRep_chartTransition,TranscendentalNakamaye.Pending.exists_contDiff_eq_one_of_isCompact,TranscendentalNakamaye.Pending.HolderAtlas.HolderMap.contDiff_indicator_comp_symm,TranscendentalNakamaye.Pending.HolderAtlas.HolderMap.exists_smul_comp_symm,TranscendentalNakamaye.Pending.HolderAtlas.HolderMap.exists_toAtlas,TranscendentalNakamaye.Pending.HolderAtlas.HolderMap.toAtlas,TranscendentalNakamaye.Pending.HolderAtlas.HolderMap.coe_toAtlas,TranscendentalNakamaye.Pending.HolderAtlas.HolderMap.ofChartSmooth,TranscendentalNakamaye.Pending.HolderAtlas.HolderMap.coe_ofChartSmooth,TranscendentalNakamaye.Pending.HolderAtlas.HolderMap.ofLE,TranscendentalNakamaye.Pending.HolderAtlas.HolderMap.norm_ofLE_le,TranscendentalNakamaye.Pending.HolderAtlas.HolderMap.ofLECLM,TranscendentalNakamaye.Pending.HolderAtlas.HolderMap.exists_subseq_tendsto_ofLE,TranscendentalNakamaye.Pending.HolderAtlas.HolderMap.isCompactOperator_ofLE}
\leanok
Let \(M\) be a compact complex manifold of dimension \(n\). (1) There are finitely many charts \(e_a\) of the atlas and chart-smooth functions \(\chi_a\ge0\) with
\(\operatorname{tsupp}\chi_a\) compact inside the source of \(e_a\) and \(\sum_a\chi_a=1\) (a \emph{H\"older atlas}). (2) For a real Banach space \(F\), \(k\in\mathbf N\), \(r\ge0\) and a H\"older
atlas, the maps \(u\colon M\to F\) with \((\chi_au)\circ e_a^{-1}\in C^{k,r}_b(\mathbf C^n,F)\) for every \(a\) (extended by \(0\) off the chart target) form a Banach space
\(C^{k,r}(M,F)\) with the norm \(\sum_a\|(\chi_au)\circ e_a^{-1}\|\). (3) For \(r\le1\), every chart \(e\) of the atlas and every \(C^{k+1}\) function \(\sigma\) with compact support in
the target of \(e\), \(\sigma\cdot(u\circ e^{-1})\in C^{k,r}_b\) with norm \(\le C\|u\|\); consequently two H\"older atlases give the same space with equivalent norms, and
chart-smooth real functions lie in \(C^{k,r}(M,\mathbf R)\). (4) For \(s\le r\) the inclusion \(C^{k,r}(M,F)\to C^{k,s}(M,F)\) is continuous, and for \(s<r\le1\) and \(F\)
finite-dimensional it is a compact operator.
\end{lemma}
\begin{proof}
\leanok
(1) A chart-smooth partition of unity subordinate to finitely many chart sources covering the compact \(M\). (2) Each \(u\) is recovered from its local pieces,
\(u=\sum_a((\chi_au)\circ e_a^{-1})\circ e_a\) on the chart sources; a Cauchy sequence has convergent local pieces (Lemma~\ref{bou:holder-spaces-higher}(2)), whose
pointwise limits assemble to a limit in the space. (3) \(\sigma\cdot(u\circ e^{-1})=\sum_a\sigma\rho_a\cdot\bigl((\chi_au)\circ e_a^{-1}\bigr)\circ(e_a\circ e^{-1})\) with cut-offs
\(\rho_a\) equal to \(1\) on \(\operatorname{tsupp}\sigma\cap e(\operatorname{tsupp}\chi_a)\) and supported where the chart transition is defined; apply the localized composition
estimate (Lemma~\ref{bou:holder-spaces-higher}(6)). Applied to the pieces of a second atlas this bounds its norm by a multiple of the first. (4) The local pieces live
on the fixed compact sets \(e_a(\operatorname{tsupp}\chi_a)\); Lemma~\ref{bou:holder-spaces-higher}(4),(5).
\end{proof}

\begin{lemma}[Simon's Schauder estimate on the whole space]
\label{bou:schauder-blowup-estimate}
\uses{bou:holder-spaces-higher,bou:holder-blowup-compactness,bou:elliptic-liouville}
\source{szekelyhidi-extremal:page-0044,szekelyhidi-extremal:page-0045,szekelyhidi-extremal:page-0046,szekelyhidi-extremal:page-0047}
\lean{TranscendentalNakamaye.Pending.SchauderBlowup.hessTrace_sub_left,TranscendentalNakamaye.Pending.SchauderBlowup.hessTrace_sub_right,TranscendentalNakamaye.Pending.SchauderBlowup.hessTrace_smul_right,TranscendentalNakamaye.Pending.SchauderBlowup.abs_hessTrace_le,TranscendentalNakamaye.Pending.SchauderBlowup.continuous_hessTrace,TranscendentalNakamaye.Pending.SchauderBlowup.continuous_hessTrace_left,TranscendentalNakamaye.Pending.SchauderBlowup.hasFDerivAt_rescale,TranscendentalNakamaye.Pending.SchauderBlowup.hasFDerivAt_fderiv_rescale,TranscendentalNakamaye.Pending.SchauderBlowup.hasFDerivAt_half_quad,TranscendentalNakamaye.Pending.SchauderBlowup.norm_smul_real,TranscendentalNakamaye.Pending.SchauderBlowup.tendsto_norm_aux,TranscendentalNakamaye.Pending.SchauderBlowup.tendsto_norm_sub_aux,TranscendentalNakamaye.Pending.SchauderBlowup.eq_of_tendsto_aux,TranscendentalNakamaye.Pending.SchauderBlowup.norm_le_of_sub_eq_zero,TranscendentalNakamaye.Pending.SchauderBlowup.norm_zero_aux,TranscendentalNakamaye.Pending.SchauderBlowup.norm_sub_self_aux,TranscendentalNakamaye.Pending.SchauderBlowup.IsEllipticCoeff,TranscendentalNakamaye.Pending.SchauderBlowup.blowup_contradiction,TranscendentalNakamaye.Pending.SchauderBlowup.fderiv_fderiv_const_mul,TranscendentalNakamaye.Pending.SchauderBlowup.exists_holder_hessian_le}
\leanok
Let \(E\) be a finite-dimensional real inner product space with an orthonormal basis \((b_i)\), and write \(a:Q=\sum_iQ(b_i,ab_i)\) for an operator \(a\) and a
bilinear form \(Q\). Let \(0<\lambda\), \(0<\alpha<1\) and \(\Lambda\in\mathbf R\). There is \(C=C(E,\lambda,\Lambda,\alpha)\ge0\) with the following property. Let
\(a\colon E\to\operatorname{End}(E)\) be symmetric, \(\lambda\)-elliptic (\(\langle a(x)v,v\rangle\ge\lambda|v|^2\)), bounded by \(\Lambda\) and \((\Lambda,\alpha)\)-H\"older, and
let \(u\in C^2(E)\) have \(\alpha\)-H\"older second derivative, \(|D^2u|\le B\) and \(|a:D^2u(x)-a:D^2u(y)|\le M|x-y|^\alpha\). Then
\(|D^2u(x)-D^2u(y)|\le C(M+B)|x-y|^\alpha\) for all \(x,y\).
\end{lemma}
\begin{proof}
\leanok
If not, there are coefficient fields \(a_n\) and functions \(u_n\) in the class with \(C=n\); normalise so that \([D^2u_n]_\alpha=1\), with a pair of points
\(p_n\), \(p_n+h_n\) where \(|D^2u_n(p_n+h_n)-D^2u_n(p_n)|>\frac12|h_n|^\alpha\), and \(\varepsilon_n=\sup|D^2u_n|+[a_n:D^2u_n]_\alpha\to0\). Rescale:
\(w_n(y)=|h_n|^{-2-\alpha}\bigl(u_n(p_n+|h_n|y)-\text{(first-order Taylor polynomial at }p_n)\bigr)-\frac12Q_n(y,y)\) with \(Q_n=|h_n|^{-\alpha}D^2u_n(p_n)\). Then
\(D^2w_n(0)=0\), \([D^2w_n]_\alpha\le1\), so along an ultrafilter \(D^2w_n\to D^2V\) pointwise for a \(C^2\) function \(V\)
(Lemma~\ref{bou:holder-blowup-compactness}); the frozen coefficients \(a_n(p_n)\) converge to a symmetric elliptic \(A\), and the bounds on \(\varepsilon_n\),
\(|a_n(p_n+|h_n|y)-a_n(p_n)|\le\Lambda|h_n|^\alpha|y|^\alpha\) and \(|h_n|^\alpha|Q_n|=|D^2u_n(p_n)|\le\varepsilon_n\) give \(A:D^2V\equiv\text{const}\); subtracting a
quadratic polynomial, \(A:D^2V=0\) with \(D^2V\) \(\alpha\)-H\"older, so \(D^2V\) is constant (Lemma~\ref{bou:elliptic-liouville}). But at
\(e=\lim h_n/|h_n|\), \(|D^2V(e)-D^2V(0)|\ge\frac12\), a contradiction.
\end{proof}
This is the first part of Sz\'ekelyhidi's sketch of Theorem~2.8 (the blow-up of Simon), in global form on \(E\); its \(C^0\) form is
Lemma~\ref{bou:schauder-interior}.

\begin{lemma}[Schauder estimates on the whole space]
\label{bou:schauder-interior}
\uses{bou:schauder-blowup-estimate,bou:holder-interpolation}
\source{szekelyhidi-extremal:page-0044,szekelyhidi-extremal:page-0045}
\lean{TranscendentalNakamaye.Pending.SchauderWholeSpace.norm_sub_le_rpow_of_lipschitz_of_bound,TranscendentalNakamaye.Pending.SchauderWholeSpace.IsLowerCoeff,TranscendentalNakamaye.Pending.SchauderWholeSpace.ellipticOp,TranscendentalNakamaye.Pending.SchauderWholeSpace.isLowerCoeff_zero,TranscendentalNakamaye.Pending.SchauderWholeSpace.ellipticOp_zero,TranscendentalNakamaye.Pending.SchauderWholeSpace.nonneg_of_isEllipticCoeff,TranscendentalNakamaye.Pending.SchauderWholeSpace.norm_fderiv_sub_le,TranscendentalNakamaye.Pending.SchauderWholeSpace.norm_fderiv_le_two_mul_add,TranscendentalNakamaye.Pending.SchauderWholeSpace.abs_sub_le_of_fderiv,TranscendentalNakamaye.Pending.SchauderWholeSpace.norm_fderiv_sub_le_rpow,TranscendentalNakamaye.Pending.SchauderWholeSpace.abs_sub_le_rpow,TranscendentalNakamaye.Pending.SchauderWholeSpace.abs_lowerOrder_sub_le,TranscendentalNakamaye.Pending.SchauderWholeSpace.exists_schauder_estimate,TranscendentalNakamaye.Pending.SchauderWholeSpace.exists_holder_hessian_le_of_sup}
\leanok
Let \(E\), \((b_i)\), \(\lambda,\Lambda,\alpha\) be as in Lemma~\ref{bou:schauder-blowup-estimate}. There is \(C=C(E,\lambda,\Lambda,\alpha)\ge0\) such that for
\(Lu=a:D^2u+Du(\beta)+cu\), with \(a\) as there and \(\beta\colon E\to E\), \(c\colon E\to\mathbf R\) bounded by \(\Lambda\) and \((\Lambda,\alpha)\)-H\"older, every
\(u\in C^2(E)\) with bounded \(\alpha\)-H\"older \(D^2u\), \(|u|\le A\) and \([Lu]_\alpha\le M\) (\(M\ge0\)) satisfies
\[
 [D^2u]_\alpha\le C(M+A),\qquad \sup|D^2u|\le C(M+A),\qquad \sup|Du|\le C(M+A).
\]
\end{lemma}
\begin{proof}
\leanok
Write \(a:D^2u=Lu-Du(\beta)-cu\). With \(\sup|Du|\le2\sup|u|+\sup|D^2u|\) and \([Du]_\alpha\le3\sup|D^2u|+4\sup|u|\) (mean value inequality, splitting
\(|x-y|\le1\) and \(|x-y|>1\)), the product rule for H\"older seminorms gives \([Du(\beta)+cu]_\alpha\le\Lambda(11\sup|u|+5\sup|D^2u|)\). Lemma~\ref{bou:schauder-blowup-estimate}
bounds \([D^2u]_\alpha\) by \(C_0(M+\Lambda(11A+5\sup|D^2u|)+\sup|D^2u|)\), and the interpolation inequality \(\sup|D^2u|\le32A/\rho^2+2[D^2u]_\alpha\rho^\alpha\)
(Lemma~\ref{bou:holder-interpolation}) at the scale \(\rho^\alpha=1/(4(C_0+1)(5|\Lambda|+1))\) absorbs \(\sup|D^2u|\); the bound on \(Du\) follows.
\end{proof}
This is Sz\'ekelyhidi's estimate (2.3) for \(k=0\) in global form on \(E\), with \([Lu]_\alpha\) on the right; local and manifold versions are obtained with
cut-offs (Lemma~\ref{bou:schauder-global}).

\begin{lemma}[The global Schauder estimate, \(k=0\)]
\label{bou:schauder-global-zero}
\uses{bou:schauder-interior,bou:holder-spaces-manifold,bou:holder-interpolation}
\source{szekelyhidi-extremal:page-0048,szekelyhidi-extremal:page-0049}
\lean{TranscendentalNakamaye.Pending.ManifoldSchauder.abs_apply_sub_apply_le,TranscendentalNakamaye.Pending.ManifoldSchauder.abs_crossTerm_le,TranscendentalNakamaye.Pending.ManifoldSchauder.abs_crossTerm_sub_le,TranscendentalNakamaye.Pending.ManifoldSchauder.abs_cross_sub_le,TranscendentalNakamaye.Pending.ManifoldSchauder.abs_hessTrace_sub_le,TranscendentalNakamaye.Pending.ManifoldSchauder.abs_mul_sub_mul_le,TranscendentalNakamaye.Pending.ManifoldSchauder.abs_remainder_sub_le,TranscendentalNakamaye.Pending.ManifoldSchauder.abs_sub_le_of_zero,TranscendentalNakamaye.Pending.ManifoldSchauder.crossTerm,TranscendentalNakamaye.Pending.ManifoldSchauder.ellipticOp_congr,TranscendentalNakamaye.Pending.ManifoldSchauder.ellipticOp_mul,TranscendentalNakamaye.Pending.ManifoldSchauder.ellipticOp_mul_eq,TranscendentalNakamaye.Pending.ManifoldSchauder.exists_bounds_of_hasCompactSupport,TranscendentalNakamaye.Pending.ManifoldSchauder.exists_chart_estimate,TranscendentalNakamaye.Pending.ManifoldSchauder.exists_lipschitz_of_hasCompactSupport,TranscendentalNakamaye.Pending.ManifoldSchauder.exists_norm_le,TranscendentalNakamaye.Pending.ManifoldSchauder.fderiv_fderiv_mul_apply,TranscendentalNakamaye.Pending.ManifoldSchauder.remainder,TranscendentalNakamaye.Pending.ManifoldSchauder.remainderConst}
\leanok
Let \(M\) be a complex manifold of dimension \(n\) with a H\"older atlas \((e_i,\chi_i)\) (Lemma~\ref{bou:holder-spaces-manifold}), \((b_j)\) an orthonormal basis of
the real space \(\mathbf C^n\), \(\lambda>0\), \(\Lambda\in\mathbf R\) and \(0<r<1\). There is \(C\ge0\) with the following property. Let, for each \(i\), \(L_i=a_i:D^2+D(\cdot)(\beta_i)+c_i\)
be an operator on \(\mathbf C^n\) of the class of Lemma~\ref{bou:schauder-interior} with constants \((\lambda,\Lambda,r)\); let \(S\ge0\), \(u\in C^{2,r}(M)\) with \(|u|\le S\), and
\(g\in C^{0,r}(M)\) with \(g(e_i^{-1}(z))=L_i(u\circ e_i^{-1})(z)\) for \(z\in e_i(\operatorname{tsupp}\chi_i)\) and all \(i\). Then
\(\|u\|_{C^{2,r}(M)}\le C(\|g\|_{C^{0,r}(M)}+S)\).
\end{lemma}
\begin{proof}
\leanok
In one chart, \((\chi_iu)\circ e_i^{-1}=\tilde\chi v\) with \(v=\sigma\cdot(u\circ e_i^{-1})\) (\(\sigma=1\) near the support, Lemma~\ref{bou:holder-spaces-manifold}~(3)), and by the product
rule \(L_i(\tilde\chi v)=\tilde\chi\,(g\circ e_i^{-1})+R\), where \(R\) contains only \(v\), \(Dv\) and derivatives of \(\tilde\chi\), so \([R]_r\le K(\sup|v|+\sup|D^2v|)\).
Lemma~\ref{bou:schauder-interior} and the interpolation inequality (Lemma~\ref{bou:holder-interpolation}) give, for every \(\rho>0\),
\(\|(\chi_iu)\circ e_i^{-1}\|_{C^{2,r}_b}\le K_g\|g\|+K_s(1+\rho^{-2})S+K_u\rho^r\|u\|\), with constants independent of the operator in the class. Summing over the
finitely many charts and choosing \(\rho^r=1/(2(\sum_iK_u+1))\) absorbs \(\|u\|\).
\end{proof}

\begin{lemma}[Schauder estimates of every order on the whole space]
\label{bou:schauder-whole-space-higher}
\uses{bou:schauder-interior,bou:holder-spaces-higher,bou:linear-elliptic-regularity,bou:holder-germ}
\source{szekelyhidi-extremal:page-0044,szekelyhidi-extremal:page-0045}
\lean{TranscendentalNakamaye.Pending.BoundedContDiffHolderMap.lowerMap,TranscendentalNakamaye.Pending.BoundedContDiffHolderMap.coe_lowerMap,TranscendentalNakamaye.Pending.BoundedContDiffHolderMap.norm_lowerMap_le,TranscendentalNakamaye.Pending.BoundedContDiffHolderMap.constMap,TranscendentalNakamaye.Pending.BoundedContDiffHolderMap.coe_constMap,TranscendentalNakamaye.Pending.BoundedContDiffHolderMap.norm_constMap_le,TranscendentalNakamaye.Pending.BoundedContDiffHolderMap.clmMap,TranscendentalNakamaye.Pending.BoundedContDiffHolderMap.coe_clmMap,TranscendentalNakamaye.Pending.BoundedContDiffHolderMap.norm_clmMap_le,TranscendentalNakamaye.Pending.BoundedContDiffHolderMap.dirMap,TranscendentalNakamaye.Pending.BoundedContDiffHolderMap.coe_dirMap,TranscendentalNakamaye.Pending.BoundedContDiffHolderMap.norm_dirMap_le,TranscendentalNakamaye.Pending.SchauderHigherOrder.abs_hessTrace_le,TranscendentalNakamaye.Pending.SchauderHigherOrder.pairBil,TranscendentalNakamaye.Pending.SchauderHigherOrder.norm_apply_le',TranscendentalNakamaye.Pending.SchauderHigherOrder.hessTraceMap,TranscendentalNakamaye.Pending.SchauderHigherOrder.coe_hessTraceMap,TranscendentalNakamaye.Pending.SchauderHigherOrder.norm_hessTraceMap_le,TranscendentalNakamaye.Pending.SchauderHigherOrder.eq_sum_apply_smul,TranscendentalNakamaye.Pending.SchauderHigherOrder.principalRhs,TranscendentalNakamaye.Pending.SchauderHigherOrder.coe_principalRhs,TranscendentalNakamaye.Pending.SchauderHigherOrder.norm_principalRhs_le,TranscendentalNakamaye.Pending.SchauderHigherOrder.diffRhs,TranscendentalNakamaye.Pending.SchauderHigherOrder.coe_diffRhs,TranscendentalNakamaye.Pending.SchauderHigherOrder.norm_diffRhs_le,TranscendentalNakamaye.Pending.SchauderHigherOrder.coe_sum,TranscendentalNakamaye.Pending.SchauderHigherOrder.norm_le_add_norm_fderivMap,TranscendentalNakamaye.Pending.SchauderHigherOrder.norm_fderivMap_le_sum,TranscendentalNakamaye.Pending.SchauderHigherOrder.exists_norm_le}
\leanok
Let \(E\) be a finite-dimensional real inner product space with an orthonormal basis \((b_i)\), \(0<r<1\), \(k\in\mathbf N\), \(\lambda>0\), and
\(Lu=a:D^2u+Du(\beta)+cu\) with coefficients \(a\in C^{k,r}_b(E,\operatorname{End}E)\), \(\beta\in C^{k,r}_b(E,E)\), \(c\in C^{k,r}_b(E)\) (Lemma~\ref{bou:holder-spaces-higher}),
\(a(x)\) symmetric and \(\lambda\)-elliptic for every \(x\). There is \(C\) such that every \(u\in C^{k+2,r}_b(E)\) with \(|u|\le S\) satisfies
\[
 \|u\|_{C^{k+2,r}_b}\ \le\ C\bigl(\|Lu\|_{C^{k,r}_b}+S\bigr).
\]
\end{lemma}
\begin{proof}
\leanok
Induction on \(k\). For \(k=0\) the coefficients are bounded and \(r\)-H\"older with constants \(2\|\cdot\|_{C^{0,r}_b}\), and Lemma~\ref{bou:schauder-interior} bounds
\([D^2u]_r\), \(\sup|D^2u|\) and \(\sup|Du|\) by \(C([Lu]_r+S)\). For \(k+1\): the case \(k\), applied to \(u\) with the same coefficients (in \(C^{k,r}_b\)), bounds
\(\|u\|_{C^{k+2,r}_b}\le C(\|Lu\|_{C^{k,r}_b}+S)\). The function \(u\) solves \(a:D^2u=\tilde g\) with \(\tilde g=Lu-Du(\beta)-cu\in C^{k+1,r}_b\), so by the symmetry of third
derivatives (as in Lemma~\ref{bou:linear-elliptic-regularity}) \(w_i=\partial_{b_i}u\) solves
\[
 Lw_i=\partial_{b_i}\tilde g-(\partial_{b_i}a):D^2u+Dw_i(\beta)+c\,w_i ,
\]
and products, derivatives and inclusions in the spaces \(C^{j,r}_b\) (Lemma~\ref{bou:holder-germ}) bound the right-hand side in \(C^{k,r}_b\) by a constant times
\(\|Lu\|_{C^{k+1,r}_b}+\|u\|_{C^{k+2,r}_b}\). The case \(k\) applied to \(w_i\) (with \(|w_i|\le\|u\|_{C^{k+2,r}_b}\)) bounds \(\|w_i\|_{C^{k+2,r}_b}\), and
\(Du=\sum_iw_i\langle b_i,\cdot\rangle\) together with \(\|u\|_{C^{k+3,r}_b}\le S+\|Du\|_{C^{k+2,r}_b}\) gives the estimate for \(k+1\).
\end{proof}
This is the case of the whole space of Sz\'ekelyhidi's Theorem~2.8, estimate (2.3), for every \(k\); on the whole space the bound for \(k+1\) follows
from the case \(k\) by differentiating the equation, without a further interpolation.

\begin{lemma}[Global Schauder estimates on a compact manifold]
\label{bou:schauder-global}
\uses{bou:schauder-global-zero,bou:schauder-whole-space-higher,bou:holder-spaces-manifold,bou:holder-germ,bou:linear-elliptic-regularity}
\source{szekelyhidi-extremal:page-0048,szekelyhidi-extremal:page-0049}
\lean{TranscendentalNakamaye.Pending.HolderAtlas.HolderMap.lowerMap,TranscendentalNakamaye.Pending.HolderAtlas.HolderMap.coe_lowerMap,TranscendentalNakamaye.Pending.HolderAtlas.HolderMap.norm_lowerMap_le,TranscendentalNakamaye.Pending.ManifoldSchauderHigher.crossMap,TranscendentalNakamaye.Pending.ManifoldSchauderHigher.coe_crossMap,TranscendentalNakamaye.Pending.ManifoldSchauderHigher.norm_clmMap_apply_le,TranscendentalNakamaye.Pending.ManifoldSchauderHigher.norm_crossTerm_le,TranscendentalNakamaye.Pending.ManifoldSchauderHigher.exists_norm_crossMap_le,TranscendentalNakamaye.Pending.ManifoldSchauderHigher.remainderMap,TranscendentalNakamaye.Pending.ManifoldSchauderHigher.coe_remainderMap,TranscendentalNakamaye.Pending.ManifoldSchauderHigher.exists_norm_remainderMap_le,TranscendentalNakamaye.Pending.ManifoldSchauderHigher.exists_chart_estimate_succ,TranscendentalNakamaye.Pending.ManifoldSchauderHigher.exists_norm_le,TranscendentalNakamaye.Pending.ManifoldSchauderRegularity.exists_of_holderAt,TranscendentalNakamaye.Pending.ManifoldSchauderRegularity.exists_holderMap_of_holderAt,TranscendentalNakamaye.Pending.ManifoldSchauderRegularity.holderAt_of_ellipticOp_eq,TranscendentalNakamaye.Pending.ManifoldSchauderRegularity.exists_holderMap_of_ellipticOp_eq}
\leanok
Let \(M\) be a compact complex manifold with a H\"older atlas (Lemma~\ref{bou:holder-spaces-manifold}), \(k\ge0\), \(0<\alpha<1\), and \(L\) a second-order
operator which in every chart of the atlas reads \(a^{ij}\partial_i\partial_j+b^i\partial_i+c\) (real coordinates) with \((a^{ij})\) symmetric and
\(\lambda\)-elliptic and all coefficients in \(C^{k,\alpha}\). There is \(C\) such that every \(f\in C^{k+2,\alpha}(M)\) satisfies
\(\|f\|_{C^{k+2,\alpha}(M)}\le C(\|Lf\|_{C^{k,\alpha}(M)}+\sup_M|f|)\); and if \(f\in C^{2,\alpha}(M)\) and \(Lf\in C^{k,\alpha}(M)\), then \(f\in C^{k+2,\alpha}(M)\).
\end{lemma}
\begin{proof}
\leanok
In the chart \(e_i\) the operator is \(L_iv=a_i:D^2v+Dv(\beta_i)+c_iv\) with coefficients extended to the whole model space inside the class \(C^{k,\alpha}_b\),
\(a_i\) symmetric and \(\lambda\)-elliptic; the estimate only uses \(L_i(f\circ e_i^{-1})=(Lf)\circ e_i^{-1}\) on \(e_i(\operatorname{tsupp}\chi_i)\).

Estimate, by induction on \(k\). For \(k=0\) this is Lemma~\ref{bou:schauder-global-zero}. For \(k+1\), let \(\sigma_i\) be a cut-off equal to \(1\) near
\(e_i(\operatorname{tsupp}\chi_i)\) with compact support in the chart target, \(\tilde\chi_i=\chi_i\circ e_i^{-1}\) and \(v_i=\sigma_i\cdot(f\circ e_i^{-1})\), so that the local
piece \(f_i=(\chi_if)\circ e_i^{-1}=\tilde\chi_iv_i\) and \(\|v_i\|_{C^{k+2,\alpha}_b}\le C\|f\|_{C^{k+2,\alpha}(M)}\) (Lemma~\ref{bou:holder-spaces-manifold}). By the
product rule \(L_if_i=\tilde\chi_i\,(Lf)\circ e_i^{-1}+R_i\), where \(R_i=v_i(a_i:D^2\tilde\chi_i+D\tilde\chi_i(\beta_i))+a_i:(D\tilde\chi_i\otimes Dv_i+Dv_i\otimes D\tilde\chi_i)\)
only involves \(v_i\) and \(Dv_i\); products and derivatives in the spaces \(C^{j,\alpha}_b\) (Lemma~\ref{bou:holder-germ}) give
\(\|R_i\|_{C^{k+1,\alpha}_b}\le K\|v_i\|_{C^{k+2,\alpha}_b}\). The whole-space estimate of order \(k+1\) (Lemma~\ref{bou:schauder-whole-space-higher}) bounds
\(\|f_i\|_{C^{k+3,\alpha}_b}\) by \(C(\|Lf\|_{C^{k+1,\alpha}(M)}+\|f\|_{C^{k+2,\alpha}(M)}+\sup|f|)\); summing over the charts and bounding \(\|f\|_{C^{k+2,\alpha}(M)}\) by the case \(k\)
(applied to the same operator, whose coefficients are in \(C^{k,\alpha}\)) gives the estimate. No interpolation is needed beyond \(k=0\).

Regularity. In the chart \(e_i\), \(v=f\circ e_i^{-1}\) is \(C^{2,\alpha}\) near every point of the chart target (Lemma~\ref{bou:holder-germ}) and solves
\(a_i:D^2v=\tilde g\) with \(\tilde g=(Lf)\circ e_i^{-1}-Dv(\beta_i)-c_iv\). If \(v\) is \(C^{j+2,\alpha}\) there (\(j<k\)), then \(\tilde g\) is \(C^{j+1,\alpha}\), and
Lemma~\ref{bou:linear-elliptic-regularity} makes \(v\) of class \(C^{j+3,\alpha}\); so \(v\) is \(C^{k+2,\alpha}\) near every point of every chart target. A function that is
\(C^{m,\alpha}\) near every point and has compact support lies in \(C^{m,\alpha}_b\) (cover the support by finitely many neighbourhoods on which it agrees with elements of
\(C^{m,\alpha}_b\); for a Lebesgue number \(\delta\) of this cover, pairs of points at distance \(<\delta\) lie in one neighbourhood, and pairs at distance \(\ge\delta\) are controlled
by the sup norm). Applied to the local pieces \((\chi_if)\circ e_i^{-1}\), this gives \(f\in C^{k+2,\alpha}(M)\).
\end{proof}
Sz\'ekelyhidi's Theorem~2.10 has the \(L^1\)-norm of \(f\) on the right; the sup norm suffices here, since the a priori estimates of Yau's theorem bound
\(\sup|\varphi|\) first (Lemma~\ref{bou:yau-c0-estimate}).

\begin{lemma}[Second-order condition at an interior extremum]
\label{bou:max-principle}
\lean{IsLocalMin.fderiv_fderiv_nonneg,IsLocalMax.fderiv_fderiv_nonpos}
\leanok
Let \(E\) be a real normed space and \(f\colon E\to\mathbf R\) of class \(C^2\) near \(x\). If \(f\) has a local minimum at \(x\), then \(D^2f(x)(h,h)\ge0\) for
every \(h\in E\); at a local maximum \(D^2f(x)(h,h)\le0\).
\end{lemma}
\begin{proof}
\leanok
Restrict to the line \(t\mapsto x+th\): \(g(t)=f(x+th)\) has a local minimum at \(0\), \(g'(0)=0\) and \(g''(0)=D^2f(x)(h,h)\); if \(g''(0)<0\), Taylor's formula
gives \(g(t)<g(0)\) for small \(t\ne0\). The maximum case is the minimum case for \(-f\).
\end{proof}
At an extremum of a function on a complex manifold the complex Hessian \(\partial_j\partial_{\bar k}f\) is then semidefinite of the same sign (it is the
average of \(D^2f(h,h)\) and \(D^2f(ih,ih)\)); this is the maximum principle used in Lemmas~\ref{bou:yau-c2-estimate} and~\ref{bou:yau-c3-estimate}.

\begin{lemma}[Hadamard's inequality and Gram determinants]
\label{bou:gram-determinant}
\source{blocki05-uniform-estimate:page-0003}
\lean{TranscendentalNakamaye.Pending.GramDeterminant.apply_eq_zero_of_diag_eq_zero,TranscendentalNakamaye.Pending.GramDeterminant.det_le_prod_diag,TranscendentalNakamaye.Pending.GramDeterminant.riesz,TranscendentalNakamaye.Pending.GramDeterminant.riesz_apply,TranscendentalNakamaye.Pending.GramDeterminant.inner_riesz,TranscendentalNakamaye.Pending.GramDeterminant.riesz_innerSL,TranscendentalNakamaye.Pending.GramDeterminant.det_riesz_comp,TranscendentalNakamaye.Pending.GramDeterminant.posSemidef_gram,TranscendentalNakamaye.Pending.GramDeterminant.det_gram_nonneg,TranscendentalNakamaye.Pending.GramDeterminant.det_gram_le_prod}
\leanok
(1) If \(M\) is a real positive semidefinite matrix, then \(\det M\le\prod_kM_{kk}\). (2) Let \(Q\) be a symmetric positive semidefinite bilinear form on a
finite-dimensional real inner product space \(E\). The Gram determinant \(\det(Q(b_i,b_j))\) does not depend on the orthonormal basis \((b_i)\), it is \(\ge0\),
and \(\det(Q(b_i,b_j))\le\prod_kQ(c_k,c_k)\) for every orthonormal basis \((c_k)\).
\end{lemma}
\begin{proof}
\leanok
(1) If some \(M_{kk}=0\), the \(k\)-th row vanishes (Cauchy--Schwarz for the semidefinite form) and \(\det M=0\). Otherwise \(D^{-1/2}MD^{-1/2}\), \(D=\operatorname{diag}(M_{kk})\),
is positive semidefinite with unit diagonal, so its eigenvalues are \(\ge0\) with sum \(\dim\), and by the arithmetic--geometric mean inequality its
determinant is \(\le1\). (2) \(Q(b_i,b_j)\) is the matrix of the operator \(R_Q\) with \(\langle R_Qx,y\rangle=Q(x,y)\) in the basis \((b_i)\), whose determinant does not
depend on the basis; apply (1) to the matrix of \(R_Q\) in \((c_k)\).
\end{proof}

\begin{lemma}[The Alexandrov--Bakelman--Pucci estimate, measure form]
\label{bou:yau-abp}
\uses{bou:max-principle,bou:gram-determinant}
\source{blocki05-uniform-estimate:page-0002,lnm2038-monge-ampere:page-0219,lnm2038-monge-ampere:page-0220}
\lean{TranscendentalNakamaye.Pending.AlexandrovBakelmanPucci.isClosed_setOf_nonneg,TranscendentalNakamaye.Pending.AlexandrovBakelmanPucci.measurableSet_contact,TranscendentalNakamaye.Pending.AlexandrovBakelmanPucci.exists_mem_contact_riesz_eq,TranscendentalNakamaye.Pending.AlexandrovBakelmanPucci.addHaar_ball_le_lintegral_det}
\leanok
Let \(E\) be a finite-dimensional real inner product space with a Haar measure \(\mu\) and an orthonormal basis \((b_i)\), \(U\subset E\) open, \(v\in C^2(U)\),
\(K\subset U\) compact, \(x_0\in K\) with \(v(x_0)\le-a\), \(v\ge0\) on \(K\setminus\operatorname{int}K\), and \(|x-x_0|\le d\) on \(K\). Then
\[
 \mu\bigl(B(0,a/d)\bigr)\le\int_\Gamma\det\bigl(D^2v(x)(b_i,b_j)\bigr)\,d\mu(x),\qquad \Gamma=\{x\in\operatorname{int}K:\ v(x)<0,\ D^2v(x)\ge0\}.
\]
\end{lemma}
\begin{proof}
\leanok
For \(|p|<a/d\), \(w=v-\langle p,\cdot\rangle\) attains its minimum on \(K\) at some \(x\); on \(K\setminus\operatorname{int}K\),
\(w\ge-\langle p,x\rangle>-a-\langle p,x_0\rangle\ge w(x_0)\), so \(x\in\operatorname{int}K\), \(v(x)<0\), \(\nabla v(x)=p\) and \(D^2v(x)\ge0\)
(Lemma~\ref{bou:max-principle}). Hence \(B(0,a/d)\) lies in the image of \(\Gamma\) under the gradient map \(\nabla v\), whose derivative has determinant
\(\det(D^2v(b_i,b_j))\ge0\) on \(\Gamma\) (Lemma~\ref{bou:gram-determinant}), and the area inequality bounds the measure of the image.
\end{proof}
This is Gilbarg--Trudinger's Lemma~9.2 in the form used by B\l{}ocki (proof of his Theorem~1; Theorem~5.10 of the Lecture Notes chapter).

\begin{lemma}[B\l{}ocki's determinant inequality]
\label{bou:yau-hessian-det}
\uses{bou:gram-determinant}
\source{blocki05-uniform-estimate:page-0002,blocki05-uniform-estimate:page-0003}
\lean{TranscendentalNakamaye.Pending.ComplexHessianDeterminant.isHermitian_complexHessianOfBilin,TranscendentalNakamaye.Pending.ComplexHessianDeterminant.I_smul_coordCombination,TranscendentalNakamaye.Pending.ComplexHessianDeterminant.dotProduct_mulVec_complexHessianOfBilin,TranscendentalNakamaye.Pending.ComplexHessianDeterminant.coordCombination_eq_toLp,TranscendentalNakamaye.Pending.ComplexHessianDeterminant.real_inner_eq_re,TranscendentalNakamaye.Pending.ComplexHessianDeterminant.unitPair,TranscendentalNakamaye.Pending.ComplexHessianDeterminant.re_conj_unitPair_mul,TranscendentalNakamaye.Pending.ComplexHessianDeterminant.orthonormal_unitPair_smul,TranscendentalNakamaye.Pending.ComplexHessianDeterminant.det_le_four_pow_mul_sq,TranscendentalNakamaye.Pending.ComplexHessianDeterminant.re_det_nonneg,TranscendentalNakamaye.Pending.ComplexHessianDeterminant.complexHessian_eq_complexHessianOfBilin}
\leanok
Let \(Q\) be a symmetric positive semidefinite real bilinear form on \(\mathbf C^n\), and \(H(Q)_{jk}=\tfrac14\bigl(Q(e_j,e_k)+Q(ie_j,ie_k)\bigr)+\tfrac i4\bigl(Q(e_j,ie_k)-Q(ie_j,e_k)\bigr)\)
its complex Hessian. Then \(H(Q)\) is Hermitian, \(\operatorname{Re}\det H(Q)\ge0\), and for every real orthonormal basis \((b_i)\) of \(\mathbf C^n\)
\[
 \det\bigl(Q(b_i,b_j)\bigr)\le4^n\bigl(\operatorname{Re}\det H(Q)\bigr)^2 .
\]
For \(u\in C^2\), \(H(D^2u(z))\) is the complex Hessian \((u_{j\bar k}(z))\).
\end{lemma}
\begin{proof}
\leanok
Let \(w^1,\dots,w^n\) be unit eigenvectors of \(H(Q)\); then \(w^1,\dots,w^n,iw^1,\dots,iw^n\) is a real orthonormal basis, and
\(\det H(Q)=\prod_l\tfrac14\bigl(Q(w^l,w^l)+Q(iw^l,iw^l)\bigr)\ge\prod_l\tfrac12\sqrt{Q(w^l,w^l)Q(iw^l,iw^l)}\) (arithmetic--geometric mean), whose square is
\(\ge4^{-n}\det(Q(b_i,b_j))\) by Hadamard's inequality in that basis and basis independence (Lemma~\ref{bou:gram-determinant}).
\end{proof}

\begin{lemma}[Local stability of the complex Monge--Amp\`ere equation]
\label{bou:yau-local-stability}
\uses{bou:yau-abp,bou:yau-hessian-det}
\source{blocki05-uniform-estimate:page-0003,lnm2038-monge-ampere:page-0220}
\lean{TranscendentalNakamaye.Pending.MongeAmpereLocalStability.finrank_complexModel,TranscendentalNakamaye.Pending.MongeAmpereLocalStability.volume_ball_mul_le_lintegral,TranscendentalNakamaye.Pending.MongeAmpereLocalStability.ofReal_mul_le_lintegral}
\leanok
Let \(n\ge1\), \(\Omega\subset\mathbf C^n\) be open, \(u\in C^2(\Omega)\) real with \(\operatorname{Re}\det(u_{j\bar k})\le F\) on \(\Omega\), \(z_0\in\Omega\), \(a>0\), \(d\) with
\(|z-z_0|\le d\) on \(\Omega\), and assume that \(K=\overline{\{z\in\Omega:u(z)<u(z_0)+a\}}\) is a compact subset of \(\Omega\). Then
\[
 (a/d)^{2n}\,\lambda_{2n}\,\bigl(-u(z_0)-a\bigr)\le4^nF^2\int_\Omega u^-\,d\lambda ,
\]
with \(\lambda_{2n}\) the volume of the unit ball and \(u^-=\max(-u,0)\).
\end{lemma}
\begin{proof}
\leanok
If \(-u(z_0)-a\le0\) there is nothing to prove. Otherwise apply Lemma~\ref{bou:yau-abp} to \(v=u-u(z_0)-a\) on \(K\): \(v(z_0)=-a\), and a point of
\(K\setminus\operatorname{int}K\) is not in the open set \(\{u<u(z_0)+a\}\), so \(v\ge0\) there. On \(\Gamma\), \(D^2v\ge0\), so
\(\det D^2v\le4^n(\operatorname{Re}\det(u_{j\bar k}))^2\le4^nF^2\) (Lemma~\ref{bou:yau-hessian-det}); hence \(\lambda_{2n}(a/d)^{2n}\le4^nF^2\lambda(\Gamma)\). On
\(\Gamma\), \(u<u(z_0)+a<0\), so \(u^->-u(z_0)-a\) and \(\lambda(\Gamma)(-u(z_0)-a)\le\int_\Omega u^-\).
\end{proof}
This is B\l{}ocki's Proposition~3 (Proposition~5.12 of the Lecture Notes chapter) with the \(L^2\)-stability replaced by its \(L^\infty\) consequence, as there.

\begin{lemma}[A uniformly normalised local potential]
\label{bou:kahler-potential-normalization}
\source{blocki05-uniform-estimate:page-0004,lnm2038-monge-ampere:page-0221}
\lean{TranscendentalNakamaye.Pending.KahlerPotentialNormalization.mulI,TranscendentalNakamaye.Pending.KahlerPotentialNormalization.antiHalf,TranscendentalNakamaye.Pending.KahlerPotentialNormalization.antiHalf_apply,TranscendentalNakamaye.Pending.KahlerPotentialNormalization.pluriharmonicPart,TranscendentalNakamaye.Pending.KahlerPotentialNormalization.contDiff_pluriharmonicPart,TranscendentalNakamaye.Pending.KahlerPotentialNormalization.complexHessian_pluriharmonicPart,TranscendentalNakamaye.Pending.KahlerPotentialNormalization.complexHessian_sub_pluriharmonicPart,TranscendentalNakamaye.Pending.KahlerPotentialNormalization.leviForm_smul_right,TranscendentalNakamaye.Pending.KahlerPotentialNormalization.exists_pos_mul_norm_sq_le_leviForm,TranscendentalNakamaye.Pending.KahlerPotentialNormalization.exists_pos_forall_norm_sub_lt,TranscendentalNakamaye.Pending.KahlerPotentialNormalization.neg_mul_le_taylor_remainder,TranscendentalNakamaye.Pending.KahlerPotentialNormalization.exists_normalized_potential}
\leanok
Let \(n\ge1\), \(W\subset\mathbf C^n\) open, \(\rho\in C^2(W)\) strictly plurisubharmonic (positive definite Levi form at every point), and
\(\overline B(c,3R)\subset W\), \(R>0\). There are \(0<r\le R\) and \(a>0\) such that for every \(z_0\in\overline B(c,R)\), with \(P_{z_0}\) the pluriharmonic part of the
second-order Taylor polynomial of \(\rho\) at \(z_0\) (the constant and linear terms and the \(J\)-anti-invariant half of the Hessian), \(g=\rho-P_{z_0}\) satisfies
\(g(z_0)=0\), \(g\ge0\) on \(B(z_0,2r)\), and \(g\ge a\) on \(B(z_0,2r)\setminus B(z_0,r)\); moreover \(g\) has the same complex Hessian as \(\rho\).
\end{lemma}
\begin{proof}
\leanok
\(P_{z_0}\) is pluriharmonic, so \(g\) has the complex Hessian of \(\rho\). Up to the Taylor remainder, \(g(z_0+h)\) is the \(J\)-invariant part of
\(\frac12D^2\rho(z_0)(h,h)\), a positive multiple of the Levi form of \(\rho\) at \(z_0\), which is \(\ge l|h|^2\) with \(l>0\) uniform on the compact
\(\overline B(c,3R)\); by uniform continuity of \(D^2\rho\) there, the remainder is absorbed for \(|h|<2r\) with \(r\) small, giving \(g(z_0+h)\ge\frac l2|h|^2\) there,
so \(a=lr^2/2\) works.
\end{proof}

\begin{theorem}[The uniform estimate]
\label{bou:yau-c0-estimate}
\uses{bou:yau-local-stability,bou:kahler-potential-normalization,bou:theta-psh-l1-bound,ana:geometric-setting}
\source{blocki05-uniform-estimate:page-0003,blocki05-uniform-estimate:page-0004,lnm2038-monge-ampere:page-0217,lnm2038-monge-ampere:page-0221}
\lean{TranscendentalNakamaye.Pending.YauUniformEstimate.exists_ofReal_neg_le,TranscendentalNakamaye.Pending.YauUniformEstimate.continuous_of_chartSmooth,TranscendentalNakamaye.Pending.YauUniformEstimate.re_det_complexHessian_eq,TranscendentalNakamaye.Pending.YauUniformEstimate.exists_local_bound,TranscendentalNakamaye.Pending.YauUniformEstimate.eq_of_dim_zero,TranscendentalNakamaye.Pending.YauUniformEstimate.exists_sub_le_of_mongeAmpere_le}
\leanok
Let \((Y,A)\) be a compact connected K\"ahler manifold of dimension \(N\) and \(B\in\mathbf R\). There is \(C\) such that every smooth \(\varphi\) with
\(A+dd^c\varphi\ge0\) and \((A+dd^c\varphi)^N\le B\,A^N\) satisfies \(\varphi(x)-\varphi(y)\le C\) for all \(x,y\in Y\).
\end{theorem}
\begin{proof}
\leanok
If \(N=0\), the connected \(Y\) is a point. Otherwise subtract \(\max\varphi\), so \(\varphi\le0\) with \(\varphi(x_1)=0\). Then \(\varphi\) is \(A\)-psh, and Lemma~\ref{bou:theta-psh-l1-bound} bounds
\(\int|\varphi\circ e^{-1}|\) on fixed compact subsets of the charts of a finite cover. Let \(z_0\) be a minimum point. In a chart around \(z_0\) (from
a fixed finite family covering \(Y\)), let \(g\) be a local potential of \(A\) (\(A=dd^cg\)); Lemma~\ref{bou:kahler-potential-normalization} gives, with \(r\), \(a\)
depending only on the cover, \(\tilde g=g-P_{z_0}\) with \(\tilde g(z_0)=0\), \(\tilde g\ge0\) on \(B(z_0,2r)\) and \(\tilde g\ge a\) on \(B(z_0,2r)\setminus B(z_0,r)\). Apply Lemma~\ref{bou:yau-local-stability} on \(\Omega=B(z_0,2r)\) to \(u=\tilde g+\varphi\): the
complex Hessian of \(u\) is that of \(A+dd^c\varphi\), so \(\det(u_{j\bar k})\le BC''\); \(u(z_0)=\min\varphi\) and \(u\ge a+\min\varphi\) on the annulus, so \(\{u<u(z_0)+a\}\subset B(z_0,r)\) is relatively compact; and
\(\int u^-\le C(1+\int|\varphi|)\). This bounds \(-u(z_0)\), i.e.\ \(-\min\varphi\), by a constant depending only on \((Y,A)\) and \(B\).
\end{proof}
Sz\'ekelyhidi proves this estimate by Moser iteration (his Proposition~3.15); the route here is B\l{}ocki's
second proof (Theorem~4 of B\l{}ocki 2005, Theorem~5.6 of the Lecture Notes chapter), which needs no Sobolev inequality.

\begin{lemma}[The continuity path]
\label{bou:yau-continuity-path}
\uses{ana:geometric-setting}
\source{szekelyhidi-extremal:page-0054,szekelyhidi-extremal:page-0067}
\lean{TranscendentalNakamaye.Pending.YauContinuityMethod.add_ddcClosedForm_zero}
\leanok
Let \((Y,A)\), \(N\ge1\), and \(F\) be as in Theorem~\ref{bou:yau}, and \(0<\varepsilon\le1\). For \(t\in[0,1]\) let \((\ast)_{t,\varepsilon}\) be the problem: find a smooth
\(\varphi\) with \(A_\varphi=A+dd^c\varphi\) K\"ahler and \(A_\varphi^N=e^{tF+\varepsilon\varphi}A^N\). Then \(\varphi=0\) solves \((\ast)_{0,\varepsilon}\). Let
\(S_\varepsilon\subset[0,1]\) be the set of \(t\) for which \((\ast)_{t,\varepsilon}\) has a solution.
\end{lemma}
\begin{proof}
\leanok
For \(\varphi=0\), \(A_0=A\) and \(e^{0}=1\).
\end{proof}
The term \(\varepsilon\varphi\) makes the linearisation \(\Delta_{A_\varphi}-\varepsilon\) of the equation invertible (Lemma~\ref{bou:kahler-laplacian-realizable}), so no
normalising constant and no Fredholm alternative are needed; the passage \(\varepsilon\to0\) is Lemma~\ref{bou:yau-eps-limit}. Sz\'ekelyhidi's path for
\(c_1(M)<0\) is \(e^{tF+\varphi}\).

\begin{lemma}[The Helmholtz equation on the whole space]
\label{bou:helmholtz-whole-space}
\uses{bou:elliptic-mean-value}
\lean{TranscendentalNakamaye.Pending.SchwartzHelmholtz.postcomp,TranscendentalNakamaye.Pending.SchwartzHelmholtz.postcomp_apply,TranscendentalNakamaye.Pending.SchwartzHelmholtz.hasTemperateGrowth_inv_one_add,TranscendentalNakamaye.Pending.SchwartzHelmholtz.fourierMultiplierCLM_add_apply,TranscendentalNakamaye.Pending.SchwartzHelmholtz.laplacian_apply_eq_sum,TranscendentalNakamaye.Pending.SchwartzHelmholtz.fderiv_fderiv_comp_left,TranscendentalNakamaye.Pending.SchwartzHelmholtz.exists_laplacian_sub_eq,TranscendentalNakamaye.Pending.SchwartzHelmholtz.exists_contDiff_laplacian_sub_eq}
\leanok
Let \(E\) be a finite-dimensional real inner product space with an orthonormal basis \((b_i)\), \(\Delta u=\sum_iD^2u(b_i,b_i)\), and \(\mu>0\). For every real
Schwartz function \(f\) on \(E\) there is a real Schwartz function \(u\) with \(\Delta u-\mu u=f\); in particular, for \(f\) smooth with compact support there
is a smooth \(u\), with all derivatives bounded, solving \(\Delta u-\mu u=f\).
\end{lemma}
\begin{proof}
\leanok
On Schwartz functions \(\Delta\) is the Fourier multiplier with symbol \(-4\pi^2|\xi|^2\); the symbol \(-(\mu+4\pi^2|\xi|^2)^{-1}\) has temperate growth, so the
corresponding multiplier maps \(f\) to a complex Schwartz function \(v\) with \(\Delta v-\mu v=f\), and its real part is a real solution. A smooth compactly
supported function is a Schwartz function, and Schwartz functions have bounded derivatives of every order.
\end{proof}

\begin{lemma}[The maximum principle for \(L-\mu\) on the whole space]
\label{bou:whole-space-max-principle}
\uses{bou:max-principle,bou:schauder-interior}
\lean{TranscendentalNakamaye.Pending.WholeSpaceMaxPrinciple.innerBilin,TranscendentalNakamaye.Pending.WholeSpaceMaxPrinciple.innerBilin_apply,TranscendentalNakamaye.Pending.WholeSpaceMaxPrinciple.hessTrace_eq_of_orthonormalBasis,TranscendentalNakamaye.Pending.WholeSpaceMaxPrinciple.hessTrace_nonpos,TranscendentalNakamaye.Pending.WholeSpaceMaxPrinciple.hessTrace_innerBilin_le,TranscendentalNakamaye.Pending.WholeSpaceMaxPrinciple.hasFDerivAt_norm_sub_sq,TranscendentalNakamaye.Pending.WholeSpaceMaxPrinciple.hasFDerivAt_two_smul_innerBilin_sub,TranscendentalNakamaye.Pending.WholeSpaceMaxPrinciple.ellipticOp_neg,TranscendentalNakamaye.Pending.WholeSpaceMaxPrinciple.mul_le_of_ellipticOp_sub,TranscendentalNakamaye.Pending.WholeSpaceMaxPrinciple.mul_abs_le_of_ellipticOp_sub}
\leanok
Let \(E\) be a finite-dimensional real inner product space, \(Lu=a:D^2u+Du(\beta)+cu\) with \(a(x)\) symmetric positive semidefinite, \(|a(x)|,|\beta(x)|\le\Lambda\) and
\(c\le0\), and \(\mu>0\). Every bounded \(u\in C^2(E)\) with \(|Lu-\mu u|\le F\) on \(E\) satisfies \(\mu|u|\le F\) on \(E\).
\end{lemma}
\begin{proof}
\leanok
It suffices to show \(\mu u(x_0)\le F\) (apply it to \(-u\) as well). For \(\varepsilon>0\) the function \(u-\varepsilon|x-x_0|^2\) attains its maximum at some \(x\) (\(u\) is
bounded); there \(Du(x)=2\varepsilon\langle x-x_0,\cdot\rangle\) and \(D^2u(x)\le2\varepsilon\langle\cdot,\cdot\rangle\) (Lemma~\ref{bou:max-principle}), so
\(a:D^2u(x)\le2\varepsilon\operatorname{tr}a\), \(|Du(x)(\beta)|\le2\varepsilon\Lambda|x-x_0|\), and if \(u(x)\ge0\), \(cu\le0\); hence
\(\mu u(x)\le F+\varepsilon C(1+|x-x_0|)\); and \(u(x_0)\le u(x)-\varepsilon|x-x_0|^2\), so \(\varepsilon|x-x_0|^2\le2\sup|u|\) and
\(\varepsilon|x-x_0|\le\sqrt{2\varepsilon\sup|u|}\). Hence \(\mu u(x_0)\le\mu u(x)\le F+o(1)\) as \(\varepsilon\to0\) (if \(u(x)<0\), then \(u(x_0)<0\) directly).
\end{proof}

\begin{lemma}[The Schauder estimate with a large parameter]
\label{bou:schauder-weighted}
\uses{bou:schauder-interior,bou:whole-space-max-principle}
\lean{TranscendentalNakamaye.Pending.SchauderWeighted.norm_sub_le_rpow_of_min,TranscendentalNakamaye.Pending.SchauderWeighted.norm_smul_sub_smul,TranscendentalNakamaye.Pending.SchauderWeighted.fderiv_comp_smul,TranscendentalNakamaye.Pending.SchauderWeighted.fderiv_fderiv_comp_smul,TranscendentalNakamaye.Pending.SchauderWeighted.ellipticOp_rescale,TranscendentalNakamaye.Pending.SchauderWeighted.exists_weighted_estimate}
\leanok
Let \(E\), \((b_i)\), \(\lambda,\Lambda,\alpha\) and \(L\) be as in Lemma~\ref{bou:schauder-interior}, with moreover \(c\le0\). There is \(C=C(E,\lambda,\Lambda,\alpha)\ge0\) such that
for every \(\mu\ge1\) and every bounded \(u\in C^2(E)\) with bounded \(\alpha\)-H\"older \(D^2u\), \(f=Lu-\mu u\) with \(|f|\le F\), \([f]_\alpha\le M\) and
\(K=M+\mu^{\alpha/2}F\):
\[
 [D^2u]_\alpha\le CK,\qquad \mu^{\alpha/2}\sup|D^2u|\le CK,\qquad \mu^{(1+\alpha)/2}\sup|Du|\le CK,\qquad \mu\sup|u|\le F .
\]
\end{lemma}
\begin{proof}
\leanok
The last bound is Lemma~\ref{bou:whole-space-max-principle}. For the others rescale: \(v(y)=u(y/\sqrt\mu)\) solves an equation of the same class with
\(-u\) replaced by \(-v\) and data \(f(y/\sqrt\mu)/\mu\); apply Lemma~\ref{bou:schauder-interior} to \(v\) and scale back, using \(\sup|u|\le F/\mu\).
\end{proof}

\begin{lemma}[The Helmholtz equation for bounded H\"older data]
\label{bou:helmholtz-holder}
\uses{bou:helmholtz-whole-space,bou:whole-space-max-principle,bou:schauder-interior,bou:holder-mollification,bou:holder-blowup-compactness}
\lean{TranscendentalNakamaye.Pending.HelmholtzHolder.exists_contDiff_two_tendsto_of_bounded,TranscendentalNakamaye.Pending.HelmholtzHolder.exists_lipschitz_cutoff,TranscendentalNakamaye.Pending.HelmholtzHolder.exists_smooth_approx,TranscendentalNakamaye.Pending.HelmholtzHolder.isEllipticCoeff_id,TranscendentalNakamaye.Pending.HelmholtzHolder.ellipticOp_id_const,TranscendentalNakamaye.Pending.HelmholtzHolder.exists_holder_fderiv_fderiv,TranscendentalNakamaye.Pending.HelmholtzHolder.exists_laplacian_sub_eq_of_holder}
\leanok
Let \(E\), \((b_i)\) and \(\Delta\) be as in Lemma~\ref{bou:helmholtz-whole-space}, \(0<\alpha<1\) and \(\mu>0\). For every \(f\colon E\to\mathbf R\) with \(|f|\le F\) and
\(|f(x)-f(y)|\le M|x-y|^\alpha\) (\(M\ge0\)) there is \(u\in C^2(E)\) with \(D^2u\) bounded and \(\alpha\)-H\"older, \(\mu|u|\le F\) on \(E\), and \(\Delta u-\mu u=f\).
\end{lemma}
\begin{proof}
\leanok
A smooth cut-off \(\chi\) (\(\chi=1\) on the unit ball) is Lipschitz, so \(f_n=\chi(\cdot/n)\,(\phi_{1/n}\ast f)\) (Lemma~\ref{bou:holder-mollification}) are smooth with
compact support, \(|f_n|\le F\), \([f_n]_\alpha\le M+KF\) with \(K\) depending only on \(E\), and \(f_n\to f\) pointwise. Lemma~\ref{bou:helmholtz-whole-space}
gives smooth \(u_n\) with bounded derivatives and \(\Delta u_n-\mu u_n=f_n\); in particular \(D^2u_n\) is \(\alpha\)-H\"older (bounded second and third
derivatives). Lemma~\ref{bou:whole-space-max-principle} gives \(\mu|u_n|\le F\), and Lemma~\ref{bou:schauder-interior} (for \(a=\mathrm{id}\),
\(c=-\mu\)) bounds \(Du_n\), \(D^2u_n\) and \([D^2u_n]_\alpha\) uniformly in \(n\). By the compactness argument of Lemma~\ref{bou:holder-blowup-compactness}
(Arzel\`a--Ascoli on compact sets, now with uniformly bounded families), along an ultrafilter \(u_n\), \(Du_n\), \(D^2u_n\) converge pointwise to \(u\), \(Du\),
\(D^2u\) for some \(u\in C^2(E)\); the bounds and the equation pass to the limit.
\end{proof}

\begin{lemma}[The method of continuity]
\label{bou:method-of-continuity}
\lean{TranscendentalNakamaye.Pending.MethodOfContinuity.surjective_step,TranscendentalNakamaye.Pending.MethodOfContinuity.surjective_of_continuity}
\leanok
Let \(X\) be a real Banach space, \(Y\) a real normed space, and \(T_0,T_1\colon X\to Y\) linear maps with \(\|T_1u-T_0u\|\le K\|u\|\) (\(K\ge0\)). Put
\(T_t=T_0+t(T_1-T_0)\). If \(\|u\|\le C\|T_tu\|\) for all \(t\in[0,1]\) and \(u\in X\) (\(C\ge0\)), and \(T_0\) is onto, then \(T_1\) is onto.
\end{lemma}
\begin{proof}
\leanok
If \(T_s\) is onto and \(|t-s|\,CK\le\frac12\), then for \(f\in Y\) and a right inverse \(R\) of \(T_s\) the map \(u\mapsto R\bigl(f-(t-s)(T_1-T_0)u\bigr)\) is a
contraction of \(X\) with constant \(\frac12\) (the estimate for \(T_s\) applied to differences), and its fixed point \(u\) satisfies \(T_tu=f\). Steps of
length \(\delta=1/(2(CK+1))\) lead from \(0\) to \(1\) in finitely many steps.
\end{proof}

\begin{lemma}[Solvability of \(Lu-\mu u=f\) on the whole space]
\label{bou:whole-space-solvability}
\uses{bou:schauder-interior,bou:whole-space-max-principle,bou:helmholtz-holder,bou:method-of-continuity,bou:holder-spaces,bou:holder-spaces-higher}
\lean{TranscendentalNakamaye.Pending.WholeSpaceSolvability.contDiff_two,TranscendentalNakamaye.Pending.WholeSpaceSolvability.abs_apply_le,TranscendentalNakamaye.Pending.WholeSpaceSolvability.norm_fderiv_le,TranscendentalNakamaye.Pending.WholeSpaceSolvability.norm_fderiv_fderiv_le,TranscendentalNakamaye.Pending.WholeSpaceSolvability.norm_fderiv_fderiv_sub_le,TranscendentalNakamaye.Pending.WholeSpaceSolvability.norm_le_of_bounds,TranscendentalNakamaye.Pending.WholeSpaceSolvability.ofBounds,TranscendentalNakamaye.Pending.WholeSpaceSolvability.coe_ofBounds,TranscendentalNakamaye.Pending.WholeSpaceSolvability.abs_apply_le_holder,TranscendentalNakamaye.Pending.WholeSpaceSolvability.abs_sub_le_holder,TranscendentalNakamaye.Pending.WholeSpaceSolvability.ofHolder,TranscendentalNakamaye.Pending.WholeSpaceSolvability.coe_ofHolder,TranscendentalNakamaye.Pending.WholeSpaceSolvability.norm_ofHolder_le,TranscendentalNakamaye.Pending.WholeSpaceSolvability.hessTrace_add_right,TranscendentalNakamaye.Pending.WholeSpaceSolvability.ellipticOp_add,TranscendentalNakamaye.Pending.WholeSpaceSolvability.ellipticOp_smul,TranscendentalNakamaye.Pending.WholeSpaceSolvability.abs_ellipticOp_le,TranscendentalNakamaye.Pending.WholeSpaceSolvability.abs_ellipticOp_sub_le,TranscendentalNakamaye.Pending.WholeSpaceSolvability.opY,TranscendentalNakamaye.Pending.WholeSpaceSolvability.coe_opY,TranscendentalNakamaye.Pending.WholeSpaceSolvability.norm_opY_le,TranscendentalNakamaye.Pending.WholeSpaceSolvability.hessTrace_combo,TranscendentalNakamaye.Pending.WholeSpaceSolvability.exists_ellipticOp_sub_eq}
\leanok
Let \(E\), \((b_i)\), \(\lambda,\Lambda,\alpha\) and \(L\) be as in Lemma~\ref{bou:schauder-interior}, with moreover \(c\le0\), and let \(\mu>0\). For every
\(f\colon E\to\mathbf R\) with \(|f|\le F\) and \(|f(x)-f(y)|\le M|x-y|^\alpha\) (\(M\ge0\)) there is a bounded \(u\in C^2(E)\) with \(D^2u\) bounded and
\(\alpha\)-H\"older and \(Lu-\mu u=f\).
\end{lemma}
\begin{proof}
\leanok
Let \(X=C^{2,\alpha}_b(E)\) (a Banach space, Lemma~\ref{bou:holder-spaces-higher}) and \(Y=C^{0,\alpha}_b(E)\) (Lemma~\ref{bou:holder-spaces}). For \(t\in[0,1]\) let
\(L_t\) be the operator with coefficients \((1-t)\,\mathrm{id}+ta\), \(t\beta\), \(tc-\mu\); these are bounded linear maps \(X\to Y\), \(L_t=L_0+t(L_1-L_0)\),
and they all belong to the class of Lemma~\ref{bou:schauder-interior} with constants \((\min(1,\lambda),\Lambda+1+\mu,\alpha)\). That lemma and
Lemma~\ref{bou:whole-space-max-principle} (\(\mu\sup|u|\le\sup|L_tu|\)) give \(\|u\|_X\le C\|L_tu\|_Y\) with \(C\) independent of \(t\). \(L_0=\Delta-\mu\) is onto by
Lemma~\ref{bou:helmholtz-holder}, so \(L_1=L-\mu\) is onto by Lemma~\ref{bou:method-of-continuity}.
\end{proof}

\begin{lemma}[Push-forward from a chart; the inclusion into \(C^{0,r}\)]
\label{bou:holder-pushforward}
\uses{bou:holder-spaces-manifold,bou:holder-spaces-higher}
\lean{TranscendentalNakamaye.Pending.BoundedContDiffHolderMap.coe_toZero,TranscendentalNakamaye.Pending.BoundedContDiffHolderMap.norm_toZero_le,TranscendentalNakamaye.Pending.BoundedContDiffHolderMap.toZero,TranscendentalNakamaye.Pending.BoundedContDiffHolderMap.toZeroCLM,TranscendentalNakamaye.Pending.BoundedContDiffHolderMap.toZeroCLM_apply,TranscendentalNakamaye.Pending.HolderAtlas.HolderMap.coe_toZero,TranscendentalNakamaye.Pending.HolderAtlas.HolderMap.coe_toZeroCLM,TranscendentalNakamaye.Pending.HolderAtlas.HolderMap.continuous,TranscendentalNakamaye.Pending.HolderAtlas.HolderMap.continuous_χ_smul,TranscendentalNakamaye.Pending.HolderAtlas.HolderMap.exists_pushforward,TranscendentalNakamaye.Pending.HolderAtlas.HolderMap.norm_toZero_le,TranscendentalNakamaye.Pending.HolderAtlas.HolderMap.pushforward,TranscendentalNakamaye.Pending.HolderAtlas.HolderMap.pushforward_eq_zero_of_notMem,TranscendentalNakamaye.Pending.HolderAtlas.HolderMap.pushforward_of_mem,TranscendentalNakamaye.Pending.HolderAtlas.HolderMap.pushforward_of_notMem,TranscendentalNakamaye.Pending.HolderAtlas.HolderMap.pushforward_symm,TranscendentalNakamaye.Pending.HolderAtlas.HolderMap.rep_toZero,TranscendentalNakamaye.Pending.HolderAtlas.HolderMap.toZero,TranscendentalNakamaye.Pending.HolderAtlas.HolderMap.toZeroCLM}
\leanok
Let \(M\) be a complex manifold of dimension \(n\) with a H\"older atlas, \(F\) a real normed space, \(k\in\mathbf N\) and \(r\le1\).
\begin{enumerate}
\item Every element of \(C^{k,r}(M,F)\) is continuous.
\item For \(k\ge1\) the inclusions \(C^{k,r}_b(E,F)\to C^{0,r}_b(E,F)\) (\(E\) a real normed space) and \(C^{k,r}(M,F)\to C^{0,r}(M,F)\) are continuous linear maps of norm \(\le3\).
\item For a chart \(e\) of \(M\) and a compact \(K\) inside its target there is \(C\) such that every \(w\in C^{k,r}_b(\mathbf C^n,F)\) with \(\operatorname{tsupp}w\subset K\) defines
\(W\in C^{k,r}(M,F)\), equal to \(w\circ e\) on the source of \(e\) and to \(0\) elsewhere, with \(\|W\|\le C\|w\|\).
\end{enumerate}
\end{lemma}
\begin{proof}
\leanok
(1) \(u=\sum_a\chi_au\), and each \(\chi_au\) is \(((\chi_au)\circ e_a^{-1})\circ e_a\) on the source of \(e_a\) and vanishes near the complement of \(\operatorname{tsupp}\chi_a\).
(2) The \(C^{0,r}\) norm of a \(C^{k,r}_b\) function is bounded by its sup norm, the sup norm of its derivative (mean value theorem) and its \(r\)-H\"older seminorm.
(3) In each atlas chart \(e_a\) the local representative of \(W\) is \(\sigma_a\cdot(w\circ\tau_a)\), with \(\tau_a=e\circ e_a^{-1}\) the chart transition and \(\sigma_a\) a cut-off
equal to \(1\) on the compact set \(e_a(\operatorname{tsupp}\chi_a\cap e^{-1}(K))\) and supported in the domain of \(\tau_a\) (\(M\) is Hausdorff); Lemma~\ref{bou:holder-spaces-manifold}~(3).
\end{proof}

\begin{lemma}[The maximum principle for \(L-\mu\) on a compact manifold]
\label{bou:manifold-max-principle}
\uses{bou:max-principle,bou:holder-spaces-manifold,bou:holder-pushforward}
\lean{TranscendentalNakamaye.Pending.ManifoldMaxPrinciple.exists_local_rep,TranscendentalNakamaye.Pending.ManifoldMaxPrinciple.mul_abs_le_of_chartOp_sub,TranscendentalNakamaye.Pending.ManifoldMaxPrinciple.mul_le_of_chartOp_sub}
\leanok
Let \(M\) be a compact complex manifold of dimension \(n\) with a H\"older atlas \((e_i,\chi_i)\), \((b_j)\) an orthonormal basis of the real space \(\mathbf C^n\), \(r\le1\), and for
each \(i\) let \(L_i=a_i:D^2+D(\cdot)(\beta_i)+c_i\) be an operator on \(\mathbf C^n\) with \(a_i(z)\) symmetric positive semidefinite and \(c_i\le0\). Let \(u\in C^{2,r}(M)\),
\(g\colon M\to\mathbf R\) with \(g(e_i^{-1}(z))=L_i(u\circ e_i^{-1})(z)\) for \(z\in e_i(\operatorname{tsupp}\chi_i)\) and all \(i\), and \(\mu>0\). If \(|g-\mu u|\le F\) on \(M\), then
\(\mu|u|\le F\) on \(M\).
\end{lemma}
\begin{proof}
\leanok
\(u\) is continuous (Lemma~\ref{bou:holder-pushforward}), so it attains its maximum at some \(y_0\); choose \(i\) with \(\chi_i(y_0)>0\). Then \(u\circ e_i^{-1}\) has a local maximum at
\(z_0=e_i(y_0)\), so its derivative vanishes and its Hessian is \(\le0\) there (Lemma~\ref{bou:max-principle}), whence \(a_i:D^2(u\circ e_i^{-1})(z_0)\le0\); if \(u(y_0)\ge0\) also
\(c_iu(y_0)\le0\), so \(g(y_0)\le0\) and \(\mu u\le\mu u(y_0)\le g(y_0)+F\le F\) on \(M\). Applying this to \(-u\) gives the claim.
\end{proof}

\begin{lemma}[The parametrix for a large parameter]
\label{bou:manifold-parametrix}
\uses{bou:whole-space-solvability,bou:whole-space-max-principle,bou:schauder-weighted,bou:holder-pushforward,bou:holder-spaces-manifold}
\lean{TranscendentalNakamaye.Pending.ManifoldParametrix.abs_remainder_le,TranscendentalNakamaye.Pending.ManifoldParametrix.abs_remainder_sub_le',TranscendentalNakamaye.Pending.ManifoldParametrix.coe_ofHolderZero,TranscendentalNakamaye.Pending.ManifoldParametrix.coeffOne,TranscendentalNakamaye.Pending.ManifoldParametrix.coeffZero,TranscendentalNakamaye.Pending.ManifoldParametrix.crossTerm_sub_right,TranscendentalNakamaye.Pending.ManifoldParametrix.crossTerm_zero_left,TranscendentalNakamaye.Pending.ManifoldParametrix.ellipticOp_add_of_contDiffAt,TranscendentalNakamaye.Pending.ManifoldParametrix.ellipticOp_const_zero,TranscendentalNakamaye.Pending.ManifoldParametrix.ellipticOp_finset_sum,TranscendentalNakamaye.Pending.ManifoldParametrix.ellipticOp_sub,TranscendentalNakamaye.Pending.ManifoldParametrix.eq_of_ellipticOp_sub_eq,TranscendentalNakamaye.Pending.ManifoldParametrix.eq_push_identity,TranscendentalNakamaye.Pending.ManifoldParametrix.exists_ellipticOp_sub_eq_of_large,TranscendentalNakamaye.Pending.ManifoldParametrix.exists_mul_cutoff,TranscendentalNakamaye.Pending.ManifoldParametrix.exists_solve,TranscendentalNakamaye.Pending.ManifoldParametrix.exists_solve_small,TranscendentalNakamaye.Pending.ManifoldParametrix.norm_ofHolderZero_le,TranscendentalNakamaye.Pending.ManifoldParametrix.norm_sub_le_scale,TranscendentalNakamaye.Pending.ManifoldParametrix.ofHolderZero,TranscendentalNakamaye.Pending.ManifoldParametrix.pushforward_sub,TranscendentalNakamaye.Pending.ManifoldParametrix.remainder_eq_zero_of_notMem,TranscendentalNakamaye.Pending.ManifoldParametrix.remainder_sub,TranscendentalNakamaye.Pending.ManifoldParametrix.tsupport_remainder_subset}
\leanok
Let \(M\) be a complex manifold of dimension \(n\) with a H\"older atlas \((e_i,\chi_i)\), \((b_j)\) an orthonormal basis of the real space \(\mathbf C^n\), \(\lambda>0\), \(\Lambda\in\mathbf R\)
and \(0<r<1\). For each \(i\) let \(L_i\) be an operator on \(\mathbf C^n\) of the class of Lemma~\ref{bou:schauder-interior} with constants \((\lambda,\Lambda,r)\) and \(c_i\le0\),
and \(U_i\) open with \(e_i(\operatorname{tsupp}\chi_i)\subset U_i\subset e_i(\operatorname{dom}e_i)\); assume the \(L_i\) come from one operator on \(M\): for every
\(u\in C^{2,r}(M)\) there is \(g\colon M\to\mathbf R\) with \(g(e_i^{-1}(z))=L_i(u\circ e_i^{-1})(z)\) for \(z\in U_i\) and all \(i\). Then there is \(\mu_0\ge1\) such that for every
\(\mu\ge\mu_0\) and every \(f\in C^{0,r}(M)\) some \(u\in C^{2,r}(M)\) satisfies \(L_i(u\circ e_i^{-1})-\mu\,u\circ e_i^{-1}=f\circ e_i^{-1}\) on every \(U_i\).
\end{lemma}
\begin{proof}
\leanok
For \(h\in C^{0,r}(M)\) let \(w_i=(L_i-\mu)^{-1}((\chi_ih)\circ e_i^{-1})\) on \(\mathbf C^n\) (Lemma~\ref{bou:whole-space-solvability}; unique, hence additive in \(h\), by
Lemma~\ref{bou:whole-space-max-principle}), \(\psi_i\) cut-offs equal to \(1\) on \(e_i(\operatorname{tsupp}\chi_i)\) with support in \(U_i\), and
\(Ph=\sum_i\) (push-forward of \(\psi_iw_i\), Lemma~\ref{bou:holder-pushforward}). By the product rule, \((L-\mu)Ph=h+Rh\) with \(Rh\) the push-forward of first-order terms
in \(w_i\) times derivatives of \(\psi_i\); Lemma~\ref{bou:schauder-weighted} bounds \(w_i\), \(Dw_i\) and their H\"older seminorms by \(C\mu^{(r-1)/2}\|h\|\), so
\(\|Rh\|\le\Theta\mu^{(r-1)/2}\|h\|\le\frac12\|h\|\) for \(\mu\ge\mu_0\). Then \(h\mapsto f-Rh\) is a contraction of the Banach space \(C^{0,r}(M)\), and \(u=Ph\) for its fixed
point \(h\).
\end{proof}

\begin{lemma}[\(L-s\) is onto for every \(s>0\)]
\label{bou:manifold-continuity}
\uses{bou:manifold-parametrix,bou:method-of-continuity,bou:schauder-global-zero,bou:manifold-max-principle,bou:holder-pushforward}
\lean{TranscendentalNakamaye.Pending.ManifoldContinuity.abs_apply_le_norm,TranscendentalNakamaye.Pending.ManifoldContinuity.contDiffAt_comp_symm,TranscendentalNakamaye.Pending.ManifoldContinuity.ellipticOp_smul_of_contDiffAt,TranscendentalNakamaye.Pending.ManifoldContinuity.eq_of_chartRel,TranscendentalNakamaye.Pending.ManifoldContinuity.exists_chartOp,TranscendentalNakamaye.Pending.ManifoldContinuity.exists_ellipticOp_sub_eq,TranscendentalNakamaye.Pending.ManifoldContinuity.exists_mem_of_chartRel,TranscendentalNakamaye.Pending.ManifoldContinuity.funext_chart}
\leanok
Let \(M\) be a compact complex manifold and \((e_i,\chi_i)\), \((b_j)\), \(\lambda,\Lambda,r\), \(L_i\) (with \(c_i\le0\)) and \(U_i\) as in Lemma~\ref{bou:manifold-parametrix}, the \(L_i\)
coming from one operator \(L\) on \(M\). Then \(L\) is a bounded linear map \(C^{2,r}(M)\to C^{0,r}(M)\), and for every \(s>0\) and \(f\in C^{0,r}(M)\) there is \(u\in C^{2,r}(M)\)
with \(L_i(u\circ e_i^{-1})-s\,u\circ e_i^{-1}=f\circ e_i^{-1}\) on every \(U_i\).
\end{lemma}
\begin{proof}
\leanok
\(L-\mu_1\) is onto for \(\mu_1=\max(\mu_0,s)\) (Lemma~\ref{bou:manifold-parametrix}). For \(s'\in[s,\mu_1]\) the operators \(L-s'\) lie in one Schauder class, so
Lemma~\ref{bou:schauder-global-zero} and Lemma~\ref{bou:manifold-max-principle} (\(s\sup|u|\le\sup|(L-s')u|\)) give \(\|u\|_{C^{2,r}}\le C(1+1/s)\|(L-s')u\|_{C^{0,r}}\)
uniformly; Lemma~\ref{bou:method-of-continuity} carries surjectivity from \(L-\mu_1\) to \(L-s\).
\end{proof}

\begin{lemma}[Stokes for the K\"ahler Laplacian]
\label{bou:kahler-laplacian-stokes}
\uses{fpmma:hessian-ibp,fpmma:cofactor-divergence,fpmma:geometric-density,btgc:coefficient-congruence,bou:closed-top-currents,bou:smooth-current-closed,bou:smooth-mixed-density,hgeom:chart-partition,bou:integration-current,bou:kahler-laplacian-realizable}
\lean{TranscendentalNakamaye.Pending.KahlerLaplacianStokes.contDiff_wirt_of_le,TranscendentalNakamaye.Pending.KahlerLaplacianStokes.mixedCoefficient_cons_eq_sum_wirt,TranscendentalNakamaye.Pending.KahlerLaplacianStokes.integral_mixedCoefficient_cons_eq_zero,TranscendentalNakamaye.Pending.KahlerLaplacianStokes.mixedCoefficient_cons_const,TranscendentalNakamaye.Pending.KahlerLaplacianStokes.re_mixedCoefficient_const_mul_re_trace,TranscendentalNakamaye.Pending.KahlerLaplacianStokes.fderiv_fderiv_sum,TranscendentalNakamaye.Pending.KahlerLaplacianStokes.laplacian_sum,TranscendentalNakamaye.Pending.KahlerLaplacianStokes.laplacian_eq_zero_of_notMem_tsupport,TranscendentalNakamaye.Pending.KahlerLaplacianStokes.continuousOn_chartLaplacian,TranscendentalNakamaye.Pending.KahlerLaplacianStokes.continuous_laplacian,TranscendentalNakamaye.Pending.KahlerLaplacianStokes.contDiff_indicator_comp_symm_two,TranscendentalNakamaye.Pending.KahlerLaplacianStokes.chartDensity_mul_chartLaplacian,TranscendentalNakamaye.Pending.KahlerLaplacianStokes.integral_laplacian_eq_zero_of_slotChart,TranscendentalNakamaye.Pending.KahlerLaplacianStokes.integral_laplacian_eq_zero,TranscendentalNakamaye.Pending.KahlerLaplacianStokes.kahlerVolume_pos_of_isOpen,TranscendentalNakamaye.Pending.KahlerLaplacianStokes.isOpenPosMeasure_kahlerVolume}
\leanok
Let \((M,\omega)\) be a compact K\"ahler manifold of dimension \(N\), \(\omega^N\) the measure of Lemma~\ref{bou:integration-current}, and
\(\Delta_\omega u=\operatorname{Re}\operatorname{tr}\bigl(H_e^{-1}(\partial_j\partial_{\bar k}(u\circ e^{-1}))\bigr)\) the Laplacian of \(\omega\) (Lemma~\ref{bou:kahler-laplacian-realizable}),
\(H_e\) the coefficient matrix of \(\omega\) in the chart \(e\). Let \(u\) be of class \(C^2\) in every chart.
(1) In every chart, \(\Delta_\omega u\,\omega^N=N\,dd^cu\wedge\omega^{N-1}\) at the level of densities.
(2) \(\Delta_\omega u\) is continuous and \(\int_M\Delta_\omega u\,\omega^N=0\).
(3) \(\omega^N\) charges every nonempty open subset of \(M\).
\end{lemma}
\begin{proof}
\leanok
(1) For an invertible matrix \(H\) and any \(K\), the mixed coefficient satisfies \(C(K,H,\dots,H)=(N-1)!\,\det H\,\operatorname{tr}(H^{-1}K)\): write
\(K=H\cdot H^{-1}K\) and \(H=H\cdot1\) and use the congruence \(C(HK_1,\dots,HK_N)=\det H\,C(K_1,\dots,K_N)\) (Lemma~\ref{btgc:coefficient-congruence}) and
\(C(A,1,\dots,1)=(N-1)!\operatorname{tr}A\) (Lemma~\ref{bou:closed-top-currents}); also \(C(H,\dots,H)=N!\det H\). For \(H\) positive definite \(\det H>0\), and the
densities of Lemma~\ref{fpmma:geometric-density} give (1).
(2) In a chart \(\Delta_\omega u=a:D^2(u\circ e^{-1})\) with continuous coefficients, so \(\Delta_\omega u\) is continuous. Take a partition of unity subordinate to
charts on which \(\omega=dd^c\rho\) (Proposition~\ref{hgeom:chart-partition}); \(\Delta_\omega\) is additive on \(C^2\) functions, so it suffices to treat \(u\) with compact
support in such a chart. By (1) and the change of variables of Lemma~\ref{bou:smooth-mixed-density} (as in Lemma~\ref{bou:smooth-current-closed}), the integral is
\(N\pi^{-N}\operatorname{Re}\int C(\partial\bar\partial v,\partial\bar\partial\rho,\dots,\partial\bar\partial\rho)\), with \(v\) the chart representative of \(u\) (of class
\(C^2\) with compact support) and \(\rho\) cut off away from the support of \(v\). The cofactor field of the smooth slots is \(\partial\)-divergence free
(Lemma~\ref{fpmma:cofactor-divergence}), so \(C(\partial\bar\partial v,\dots)=\sum_{j,k}\partial_j(\bar\partial_kv\cdot\mathrm{cof}_{jk})\) is the divergence of a compactly
supported \(C^1\) field and has integral \(0\). This is the computation of Theorem~\ref{fpmma:hessian-ibp}, which needs only \(v\in C^2\) in the first slot.
(3) In a chart the density of \(\omega^N\) is continuous and positive (Lemma~\ref{bou:integration-current}), and closed balls have positive Lebesgue measure.
\end{proof}

\begin{lemma}[Solvability of the Laplace equation on a compact K\"ahler manifold]
\label{bou:kahler-laplacian-solvability}
\uses{bou:kahler-laplacian-stokes,bou:kahler-laplacian-realizable,bou:schauder-global,bou:linear-elliptic-regularity,bou:holder-spaces-manifold,bou:holder-germ,bou:manifold-continuity,bou:integration-current}
\source{szekelyhidi-extremal:page-0050,szekelyhidi-extremal:page-0051}
\lean{TranscendentalNakamaye.Pending.KahlerLaplacianStokes.fderiv_fderiv_mul_self,TranscendentalNakamaye.Pending.KahlerLaplacianStokes.hessTrace_smulRight_self,TranscendentalNakamaye.Pending.KahlerLaplacianStokes.laplacian_mul_self_sub,TranscendentalNakamaye.Pending.KahlerLaplacianStokes.exists_eq_const_of_laplacian_eq_zero,TranscendentalNakamaye.Pending.KahlerLaplacianSolvability.toZeroOne,TranscendentalNakamaye.Pending.KahlerLaplacianSolvability.coe_toZeroOne,TranscendentalNakamaye.Pending.KahlerLaplacianSolvability.norm_toZeroOne_le,TranscendentalNakamaye.Pending.KahlerLaplacianSolvability.holderToZeroOne,TranscendentalNakamaye.Pending.KahlerLaplacianSolvability.coe_holderToZeroOne,TranscendentalNakamaye.Pending.KahlerLaplacianSolvability.norm_holderToZeroOne_le,TranscendentalNakamaye.Pending.KahlerLaplacianSolvability.holderToZeroOneCLM,TranscendentalNakamaye.Pending.KahlerLaplacianSolvability.isCompactOperator_toZeroCLM,TranscendentalNakamaye.Pending.KahlerLaplacianSolvability.isCompactOperator_of_real,TranscendentalNakamaye.Pending.KahlerLaplacianSolvability.lowerCLM,TranscendentalNakamaye.Pending.KahlerLaplacianSolvability.lowerToTwo,TranscendentalNakamaye.Pending.KahlerLaplacianSolvability.coe_lowerToTwo,TranscendentalNakamaye.Pending.KahlerLaplacianSolvability.evalCLM,TranscendentalNakamaye.Pending.KahlerLaplacianSolvability.evalCLM_apply,TranscendentalNakamaye.Pending.KahlerLaplacianSolvability.one,TranscendentalNakamaye.Pending.KahlerLaplacianSolvability.one_apply,TranscendentalNakamaye.Pending.KahlerLaplacianSolvability.contDiffOn_comp_symm,TranscendentalNakamaye.Pending.KahlerLaplacianSolvability.integrable_holderMap,TranscendentalNakamaye.Pending.KahlerLaplacianSolvability.integralCLM,TranscendentalNakamaye.Pending.KahlerLaplacianSolvability.integralCLM_apply,TranscendentalNakamaye.Pending.KahlerLaplacianSolvability.laplacian_const,TranscendentalNakamaye.Pending.KahlerLaplacianSolvability.exists_holderMap_of_laplacian_eq,TranscendentalNakamaye.Pending.KahlerLaplacianSolvability.exists_holderMap_laplacian,TranscendentalNakamaye.Pending.KahlerLaplacianSolvability.laplacian_eq_laplacianCLM,TranscendentalNakamaye.Pending.KahlerLaplacianSolvability.laplacianLMk,TranscendentalNakamaye.Pending.KahlerLaplacianSolvability.laplacianLMk_apply,TranscendentalNakamaye.Pending.KahlerLaplacianSolvability.continuous_laplacianLMk,TranscendentalNakamaye.Pending.KahlerLaplacianSolvability.laplacianCLMk,TranscendentalNakamaye.Pending.KahlerLaplacianSolvability.laplacianCLMk_apply,TranscendentalNakamaye.Pending.KahlerLaplacianSolvability.kahlerVolume_real_univ_pos,TranscendentalNakamaye.Pending.KahlerLaplacianSolvability.integral_laplacian_holderMap,TranscendentalNakamaye.Pending.KahlerLaplacianSolvability.eq_zero_of_laplacian_add_integral,TranscendentalNakamaye.Pending.KahlerLaplacianSolvability.bijective_laplacianCLM_add,TranscendentalNakamaye.Pending.KahlerLaplacianSolvability.exists_laplacian_eq_of_integral_eq_zero_zero,TranscendentalNakamaye.Pending.KahlerLaplacianSolvability.exists_laplacian_eq_of_integral_eq_zero,TranscendentalNakamaye.Pending.KahlerLaplacianSolvability.exists_eq_add_const_of_laplacian_eq,TranscendentalNakamaye.Pending.KahlerLaplacianSolvability.laplacianConstCLM,TranscendentalNakamaye.Pending.KahlerLaplacianSolvability.laplacianConstCLM_apply,TranscendentalNakamaye.Pending.KahlerLaplacianSolvability.bijective_laplacianConstCLM,TranscendentalNakamaye.Pending.KahlerLaplacianSolvability.laplacianConstEquiv,TranscendentalNakamaye.Pending.KahlerLaplacianSolvability.laplacianConstEquiv_apply}
\leanok
Let \((Y,\omega)\) be a compact connected K\"ahler manifold of dimension \(N\), \(\Delta_\omega u=N\,dd^cu\wedge\omega^{N-1}/\omega^N\), \(k\ge0\), \(0<\alpha<1\).
(1) If \(u\in C^2\) and \(\Delta_\omega u=0\), then \(u\) is constant. (2) For every \(h\in C^{k,\alpha}(Y)\) with \(\int_Yh\,\omega^N=0\) there is \(u\in C^{k+2,\alpha}(Y)\),
unique up to an additive constant, with \(\Delta_\omega u=h\). Consequently \((v,\sigma)\mapsto\Delta_\omega v-\sigma\) is an isomorphism from
\(\{v\in C^{k+2,\alpha}:\int v\,\omega^N=0\}\times\mathbf R\) onto \(C^{k,\alpha}(Y)\).
\end{lemma}
\begin{proof}
\leanok
Here \(\Delta_\omega\) is the Laplacian of Lemma~\ref{bou:kahler-laplacian-realizable}; the identity \(\Delta_\omega u\,\omega^N=N\,dd^cu\wedge\omega^{N-1}\) is
Lemma~\ref{bou:kahler-laplacian-stokes}(1). \(C^{k,\alpha}(Y)\) is defined through a H\"older atlas (Lemma~\ref{bou:holder-spaces-manifold}) and \(\omega^N\) is the measure of
Lemma~\ref{bou:integration-current}.

(1) In a chart, with \(v=u\circ e^{-1}\) and \(a\) the real coefficient of \(\Delta_\omega\) (symmetric and positive definite,
Lemma~\ref{bou:kahler-laplacian-realizable}), \(D^2(v^2)=2(v\,D^2v+Dv\otimes Dv)\) gives \(\Delta_\omega(u^2)-2u\,\Delta_\omega u=2\langle a\nabla v,\nabla v\rangle\ge0\), with
equality only where \(Dv=0\). If \(\Delta_\omega u=0\), then \(\Delta_\omega(u^2)\) is continuous and nonnegative with \(\int\Delta_\omega(u^2)\,\omega^N=0\)
(Lemma~\ref{bou:kahler-laplacian-stokes}(2)); since \(\omega^N\) charges open sets (Lemma~\ref{bou:kahler-laplacian-stokes}(3)), \(\Delta_\omega(u^2)=0\). So \(du=0\) in every
chart, \(u\) is locally constant, and \(u\) is constant on the connected \(Y\).

(2) First \(k=0\). Let \(Tu=\Delta_\omega u+\int u\,\omega^N\) (a constant function), \(C^{2,\alpha}(Y)\to C^{0,\alpha}(Y)\). \(T\) is one-to-one: if \(Tu=0\), integrating
with Lemma~\ref{bou:kahler-laplacian-stokes}(2) gives \(\int u\,\omega^N=0\), so \(\Delta_\omega u=0\), \(u\) is constant by (1), and \(u=0\). Write \(T=(\Delta_\omega-1)(I+K)\) with
\(K=(\Delta_\omega-1)^{-1}(\iota+P)\). Here \(\Delta_\omega-1\) is an isomorphism (Lemma~\ref{bou:kahler-laplacian-realizable}); the inclusion \(\iota\colon C^{2,\alpha}\subset C^{0,\alpha}\)
is compact, since it factors through \(C^{0,1}\subset C^{0,\alpha}\), which is compact for \(\alpha<1\) (Lemma~\ref{bou:holder-spaces-manifold}); and \(Pu=\int u\,\omega^N\) has
rank one. So \(K\) is compact. By the Fredholm alternative for compact operators (mathlib), either \(-1\) is an eigenvalue of \(K\), which would contradict the
injectivity of \(T\), or \(I+K\) is invertible; hence \(T\) is bijective. If \(\int h\,\omega^N=0\), the solution of \(Tu=h\) has \(\int u\,\omega^N=0\) by Stokes, so
\(\Delta_\omega u=h\).

Regularity, for \(k\ge1\). In every chart \(\Delta_\omega=a:D^2\) with \(a\) smooth, symmetric and positive definite on the chart target. If \(u\in C^{2,\alpha}\) and
\(\Delta_\omega u=h\in C^{k,\alpha}\), Lemma~\ref{bou:linear-elliptic-regularity} bootstraps \(u\circ e^{-1}\) to \(C^{k+2,\alpha}\) near every point. The local-to-global step of
Lemma~\ref{bou:schauder-global} (with Lemma~\ref{bou:holder-germ}) then gives \(u\in C^{k+2,\alpha}(Y)\). Uniqueness up to an additive constant is (1) applied to the
difference of two solutions.

The isomorphism. \(\Delta_\omega\) maps \(C^{k+2,\alpha}(Y)\) into \(C^{k,\alpha}(Y)\) (the chart expression \(a:D^2v\) is a \(C^{k,\alpha}\) germ, Lemma~\ref{bou:holder-germ}). It is
bounded by the closed graph theorem: \(\Delta_\omega\colon C^{2,\alpha}\to C^{0,\alpha}\) is bounded, and point evaluations are bounded
(Lemma~\ref{bou:manifold-continuity}). The map \((v,\sigma)\mapsto\Delta_\omega v-\sigma\) on \(\{v:\int v\,\omega^N=0\}\times\mathbf R\) is one-to-one (integrate, then use (1)) and
onto (solve \(\Delta_\omega u=h+\sigma\) with \(\sigma=-\int h\,\omega^N/\int_Y\omega^N\), then subtract the mean of \(u\)). By the open mapping theorem it is an
isomorphism of Banach spaces.
\end{proof}
Sz\'ekelyhidi's Theorem~2.12 (variational proof) and Theorem~2.13 (mapping properties) state these facts; the route here (a parametrix, the method of
continuity and the Fredholm alternative) avoids Sobolev spaces.

\begin{lemma}[The Laplacian of a Hermitian form]
\label{bou:kahler-laplacian-realizable}
\uses{bou:manifold-parametrix,bou:manifold-continuity,bou:manifold-max-principle,bou:holder-spaces-manifold,bou:schauder-interior}
\source{szekelyhidi-extremal:page-0050}
\lean{TranscendentalNakamaye.Pending.KahlerLaplacianChart.bijective_laplacianCLM_sub,TranscendentalNakamaye.Pending.KahlerLaplacianChart.bilin_smul_add,TranscendentalNakamaye.Pending.KahlerLaplacianChart.chartLaplacian,TranscendentalNakamaye.Pending.KahlerLaplacianChart.chartLaplacian_eq_hessTrace,TranscendentalNakamaye.Pending.KahlerLaplacianChart.chartLaplacian_eq_of_mem,TranscendentalNakamaye.Pending.KahlerLaplacianChart.coeff_eq_compLinearMap,TranscendentalNakamaye.Pending.KahlerLaplacianChart.complexHessian_congr',TranscendentalNakamaye.Pending.KahlerLaplacianChart.contDiffOn_apply,TranscendentalNakamaye.Pending.KahlerLaplacianChart.contDiffOn_lapCoeff,TranscendentalNakamaye.Pending.KahlerLaplacianChart.contDiffOn_matCLM_map_conj,TranscendentalNakamaye.Pending.KahlerLaplacianChart.contDiff_smul_of_tsupport,TranscendentalNakamaye.Pending.KahlerLaplacianChart.eq_zero_of_laplacian_sub_eq_zero,TranscendentalNakamaye.Pending.KahlerLaplacianChart.exists_atlas_coeff,TranscendentalNakamaye.Pending.KahlerLaplacianChart.exists_chart_coeff,TranscendentalNakamaye.Pending.KahlerLaplacianChart.exists_cutoff,TranscendentalNakamaye.Pending.KahlerLaplacianChart.exists_lapCoeff_ge,TranscendentalNakamaye.Pending.KahlerLaplacianChart.exists_laplacianCLM,TranscendentalNakamaye.Pending.KahlerLaplacianChart.exists_laplacian_sub_eq,TranscendentalNakamaye.Pending.KahlerLaplacianChart.im_map_conj_normalizedHermitianMatrix,TranscendentalNakamaye.Pending.KahlerLaplacianChart.image_tsupport_subset_target,TranscendentalNakamaye.Pending.KahlerLaplacianChart.inner_matCLM_of_isHermitian,TranscendentalNakamaye.Pending.KahlerLaplacianChart.inner_matCLM_self_pos,TranscendentalNakamaye.Pending.KahlerLaplacianChart.isEllipticCoeff_mono,TranscendentalNakamaye.Pending.KahlerLaplacianChart.isHermitian_map_conj,TranscendentalNakamaye.Pending.KahlerLaplacianChart.isUnit_det_jacobianMatrix,TranscendentalNakamaye.Pending.KahlerLaplacianChart.isUnit_matCLM_map_conj,TranscendentalNakamaye.Pending.KahlerLaplacianChart.lapCoeff,TranscendentalNakamaye.Pending.KahlerLaplacianChart.lapCoeff_eq_inverse,TranscendentalNakamaye.Pending.KahlerLaplacianChart.lapCoeff_pos,TranscendentalNakamaye.Pending.KahlerLaplacianChart.lapCoeff_symm,TranscendentalNakamaye.Pending.KahlerLaplacianChart.laplacian,TranscendentalNakamaye.Pending.KahlerLaplacianChart.laplacianCLM,TranscendentalNakamaye.Pending.KahlerLaplacianChart.laplacianCLM_apply,TranscendentalNakamaye.Pending.KahlerLaplacianChart.laplacianSubEquiv,TranscendentalNakamaye.Pending.KahlerLaplacianChart.laplacianSubEquiv_apply,TranscendentalNakamaye.Pending.KahlerLaplacianChart.laplacian_eq,TranscendentalNakamaye.Pending.KahlerLaplacianChart.laplacian_eq_ellipticOp,TranscendentalNakamaye.Pending.KahlerLaplacianChart.map_conj_eq_transpose,TranscendentalNakamaye.Pending.KahlerLaplacianChart.matCLM,TranscendentalNakamaye.Pending.KahlerLaplacianChart.matCLM_apply,TranscendentalNakamaye.Pending.KahlerLaplacianChart.matCLM_coordinateVector,TranscendentalNakamaye.Pending.KahlerLaplacianChart.matCLM_eq_sum,TranscendentalNakamaye.Pending.KahlerLaplacianChart.matCLM_map_conj_inv,TranscendentalNakamaye.Pending.KahlerLaplacianChart.matCLM_map_conj_mul,TranscendentalNakamaye.Pending.KahlerLaplacianChart.matCLM_map_conj_one,TranscendentalNakamaye.Pending.KahlerLaplacianChart.matCLM_mul,TranscendentalNakamaye.Pending.KahlerLaplacianChart.matCLM_one,TranscendentalNakamaye.Pending.KahlerLaplacianChart.matCLM_smul_complex,TranscendentalNakamaye.Pending.KahlerLaplacianChart.matCLMₗ,TranscendentalNakamaye.Pending.KahlerLaplacianChart.orthonormal_coordinateVector,TranscendentalNakamaye.Pending.KahlerLaplacianChart.posDef_coeff,TranscendentalNakamaye.Pending.KahlerLaplacianChart.posDef_map_conj,TranscendentalNakamaye.Pending.KahlerLaplacianChart.re_map_conj_normalizedHermitianMatrix,TranscendentalNakamaye.Pending.KahlerLaplacianChart.re_trace_inv_mul_complexHessian,TranscendentalNakamaye.Pending.KahlerLaplacianChart.re_trace_mul_complexHessian,TranscendentalNakamaye.Pending.KahlerLaplacianChart.realBasis,TranscendentalNakamaye.Pending.KahlerLaplacianChart.realBasis_apply,TranscendentalNakamaye.Pending.KahlerLaplacianChart.standardMatrix_transitionLinear,TranscendentalNakamaye.Pending.KahlerLaplacianChart.trace_inv_mul_congr,TranscendentalNakamaye.Pending.KahlerLaplacianChart.transitionLinear,TranscendentalNakamaye.Pending.HolomorphicHessianTwo.complexHessian_comp_holomorphic_two,TranscendentalNakamaye.Pending.HolomorphicHessianTwo.differentiableAt_smoothDc,TranscendentalNakamaye.Pending.HolomorphicHessianTwo.hessForm_comp_holomorphic_two,TranscendentalNakamaye.Pending.HolomorphicHessianTwo.smoothDDC_comp_two}
\leanok
Let \(M\) be a compact complex manifold of dimension \(n\) with a H\"older atlas, \(0<r<1\), and \(g\) a smooth real \((1,1)\)-form which is Hermitian:
its coefficient matrix \(H_e\) is positive definite in every chart \(e\). For \(u\) of class \(C^2\) in every chart,
\(\Delta_gu=\operatorname{Re}\operatorname{tr}\bigl(H_e^{-1}(\partial_j\partial_{\bar k}(u\circ e^{-1}))\bigr)\) does not depend on the chart; \(\Delta_g\) is a bounded
operator \(C^{2,r}(M)\to C^{0,r}(M)\), and for every \(s>0\), \(\Delta_g-s\colon C^{2,r}(M)\to C^{0,r}(M)\) is an isomorphism.
\end{lemma}
\begin{proof}
\leanok
Under a holomorphic change of coordinates with Jacobian \(J\), both the complex Hessian of a \(C^2\) function and \(H_e\) transform by \(J^{T}(\cdot)\bar J\), and
\(\operatorname{tr}(H^{-1}K)\) is invariant. In a chart, \(\operatorname{Re}\operatorname{tr}(P\,\partial\bar\partial v)=(\tfrac14\bar P):D^2v\); with \(P=H_e^{-1}\) this is
a real elliptic operator with smooth, symmetric, positive definite coefficients. Cut off and extended, these give whole-space coefficients of the class of
Lemma~\ref{bou:schauder-interior} that agree with \(\Delta_g\) near the supports of a partition of unity of the atlas: \(\Delta_g\) is realizable in the sense of
Lemma~\ref{bou:manifold-parametrix}. Lemma~\ref{bou:manifold-continuity} gives that \(\Delta_g-s\) is onto, and the maximum principle
(Lemma~\ref{bou:manifold-max-principle}) that it is one-to-one.
\end{proof}

\begin{lemma}[Nemytskii operators on H\"older spaces]
\label{bou:holder-nemytskii}
\uses{bou:holder-spaces}
\lean{TranscendentalNakamaye.Pending.HolderNemytskii.abs_mul_sub_le,TranscendentalNakamaye.Pending.HolderNemytskii.abs_remainder_le,TranscendentalNakamaye.Pending.HolderNemytskii.abs_remainder_sub_le,TranscendentalNakamaye.Pending.HolderNemytskii.abs_remainder_sub_le_of_le_one,TranscendentalNakamaye.Pending.HolderNemytskii.norm_interp_sub_le,TranscendentalNakamaye.Pending.HolderNemytskii.norm_interp_sub_sub_le,TranscendentalNakamaye.Pending.HolderNemytskii.norm_mixed_sub_le,TranscendentalNakamaye.Pending.HolderNemytskii.norm_sub_le_of_hasFDerivAt}
\leanok
Let \(X,E\) be normed spaces, \(\Gamma\colon X\to\mathbf R\) differentiable with derivative \(g\), and \(g\) differentiable with derivative bounded by \(B_2\) and
\(B_3\)-Lipschitz. For maps \(P_0,P_1,P_2\colon E\to X\) let \(R(z)=\Gamma(P_1z)-\Gamma(P_2z)-g(P_0z)(P_1z-P_2z)\). If \(|P_j-P_0|\le\delta\) and the
\(r\)-H\"older seminorm of \(P_j-P_0\) is \(\le\delta\) (\(j=1,2\)), \(|P_1-P_2|\le d\) with \(r\)-H\"older seminorm \(\le d\), and
\(|P_0z-P_0z'|\le M_0|z-z'|^r\) for \(|z-z'|\le1\), then \(|R|\le B_2\delta d\) and \(|R(z)-R(z')|\le(2B_2+B_3(M_0+\delta))\,\delta d\,|z-z'|^r\).
\end{lemma}
\begin{proof}
\leanok
The mean value theorem along \(s\mapsto P_2+s(P_1-P_2)\), and for the mixed second difference of \(g\).
\end{proof}

\begin{lemma}[Local H\"older germs]
\label{bou:holder-germ}
\uses{bou:holder-spaces,bou:holder-spaces-manifold}
\lean{TranscendentalNakamaye.Pending.HolderGerm.exists_cutoff_nhds,TranscendentalNakamaye.Pending.HolderGerm.contDiff_smul_of_contDiffAt,TranscendentalNakamaye.Pending.HolderGerm.exists_norm_le_of_hasCompactSupport,TranscendentalNakamaye.Pending.HolderGerm.exists_lipschitz_of_hasCompactSupport,TranscendentalNakamaye.Pending.HolderGerm.HolderAt,TranscendentalNakamaye.Pending.HolderGerm.HolderAt.congr,TranscendentalNakamaye.Pending.HolderGerm.holderAt_congr,TranscendentalNakamaye.Pending.HolderGerm.BoundedContDiffHolderMap.holderAt,TranscendentalNakamaye.Pending.HolderGerm.HolderAt.contDiffAt,TranscendentalNakamaye.Pending.HolderGerm.HolderAt.eventually,TranscendentalNakamaye.Pending.HolderGerm.HolderAt.eventually_differentiableAt,TranscendentalNakamaye.Pending.HolderGerm.HolderAt.add,TranscendentalNakamaye.Pending.HolderGerm.HolderAt.sub,TranscendentalNakamaye.Pending.HolderGerm.HolderAt.const_smul,TranscendentalNakamaye.Pending.HolderGerm.holderAt_zero,TranscendentalNakamaye.Pending.HolderGerm.HolderAt.sum,TranscendentalNakamaye.Pending.HolderGerm.HolderAt.bilinear,TranscendentalNakamaye.Pending.HolderGerm.HolderAt.mul,TranscendentalNakamaye.Pending.HolderGerm.HolderAt.smul,TranscendentalNakamaye.Pending.HolderGerm.HolderAt.derivative,TranscendentalNakamaye.Pending.HolderGerm.holderAt_of_contDiffAt,TranscendentalNakamaye.Pending.HolderGerm.holderAt_const,TranscendentalNakamaye.Pending.HolderGerm.HolderAt.of_succ,TranscendentalNakamaye.Pending.HolderGerm.HolderAt.of_le,TranscendentalNakamaye.Pending.HolderGerm.HolderAt.clm,TranscendentalNakamaye.Pending.HolderGerm.HolderAt.of_fderiv,TranscendentalNakamaye.Pending.HolderGerm.HolderAt.fderiv_of_apply,TranscendentalNakamaye.Pending.HolderGerm.contDiffAt_infty_of_holderAt,TranscendentalNakamaye.Pending.HolderGerm.ofHolderWith,TranscendentalNakamaye.Pending.HolderGerm.HolderAt.comp_contDiff,TranscendentalNakamaye.Pending.HolderGerm.holderAt_comp_symm}
\leanok
Let \(E\) be a finite-dimensional real normed space, \(k\in\mathbf N\) and \(0\le r\le1\). Say that \(u\colon E\to F\) is \(C^{k,r}\) at \(z\) if it agrees near \(z\)
with a function of \(C^{k,r}_b(E,F)\). This property is local, holds for \(C^{k+1}\) germs, is preserved by sums, products, bilinear and continuous linear maps and
by composition \(x\mapsto G(x,w(x))\) with a smooth \(G\) (Nemytskii); if \(u\) is \(C^{k+1,r}\) at \(z\) then \(Du\) is \(C^{k,r}\) at \(z\), and conversely if
\(u\) is differentiable near \(z\) and \(u\), \(Du\) are \(C^{k,r}\) at \(z\) then \(u\) is \(C^{k+1,r}\) at \(z\) (it suffices that \(Du\,b_i\) is \(C^{k,r}\) for a
basis \((b_i)\)). A germ that is \(C^{k,r}\) for every \(k\) is smooth, and for \(u\in C^{k,r}(M)\) on a manifold with a H\"older atlas, \(u\circ e^{-1}\) is
\(C^{k,r}\) at every point of the target of every chart \(e\).
\end{lemma}
\begin{proof}
\leanok
Multiply by a cut-off \(\theta\equiv1\) near \(z\): \(\theta u\in C^{k,r}_b\) when \(u\) is \(C^{k+1}\) near the support, and the algebraic operations are those of
\(C^{k,r}_b\). For the derivative criterion, \(D(\theta u)=\theta Du+D\theta\otimes u\in C^{k,r}_b\). For the Nemytskii operator cut \(G\) off so that it becomes
globally Lipschitz (case \(k=0\)) and induct on \(k\) by the chain rule and the derivative criterion.
\end{proof}

\begin{lemma}[One derivative more from difference quotients]
\label{bou:difference-quotient-regularity}
\uses{bou:holder-spaces,bou:schauder-interior}
\lean{TranscendentalNakamaye.Pending.DifferenceQuotient.diffQuot,TranscendentalNakamaye.Pending.DifferenceQuotient.norm_diffQuot_le,TranscendentalNakamaye.Pending.DifferenceQuotient.norm_diffQuot_sub_le,TranscendentalNakamaye.Pending.DifferenceQuotient.tendsto_diffQuot,TranscendentalNakamaye.Pending.DifferenceQuotient.fderiv_diffQuot,TranscendentalNakamaye.Pending.DifferenceQuotient.differentiable_diffQuot,TranscendentalNakamaye.Pending.DifferenceQuotient.translate,TranscendentalNakamaye.Pending.DifferenceQuotient.coe_translate,TranscendentalNakamaye.Pending.DifferenceQuotient.diffQuotMap,TranscendentalNakamaye.Pending.DifferenceQuotient.coe_diffQuotMap,TranscendentalNakamaye.Pending.DifferenceQuotient.abs_remainder_sub_le',TranscendentalNakamaye.Pending.DifferenceQuotient.fderiv_fderiv_sub_map,TranscendentalNakamaye.Pending.DifferenceQuotient.exists_mul_fderiv_mem}
\leanok
Let \(E\) be a finite-dimensional real inner product space, \(0<r<1\), \(w\in C^{2,r}_b(E)\), \(\theta\) smooth with compact support and \(\zeta\in E\). Suppose that
for \(0<h\le h_0\) the difference quotient \(D_hw=(w(\cdot+h\zeta)-w)/h\) solves \(a_h:D^2(D_hw)=g_h\) on the support of \(\theta\), where the coefficients
\(a_h\) are elliptic of class \(C^{0,r}\) with bounds independent of \(h\), and \(|g_h|\le M\), \([g_h]_r\le M\). Then \(\theta\,\partial_\zeta w\in C^{2,r}_b(E)\).
\end{lemma}
\begin{proof}
\leanok
\(V_h=\theta D_hw\) solves \(a_h:D^2V_h=\theta g_h+R_h\), where \(R_h\) involves only \(D_hw\) and \(D(D_hw)=D_h(Dw)\) and is bounded in \(C^{0,r}\) uniformly in
\(h\) by the mean value theorem. The whole-space Schauder estimate (Lemma~\ref{bou:schauder-interior}) bounds \(\|V_h\|_{C^{2,r}_b}\) uniformly; along
\(h_j\to0\) a subsequence converges in \(C^{2,r/2}_b\), the limit is \(\theta\,\partial_\zeta w\), and the uniform \(r\)-H\"older bound of \(D^2V_h\) passes to the
limit.
\end{proof}

\begin{lemma}[Interior regularity for linear elliptic equations]
\label{bou:linear-elliptic-regularity}
\uses{bou:difference-quotient-regularity,bou:holder-germ,bou:schauder-interior}
\lean{TranscendentalNakamaye.Pending.LinearRegularity.hessTrace_add_left,TranscendentalNakamaye.Pending.LinearRegularity.hessTrace_smul_left,TranscendentalNakamaye.Pending.LinearRegularity.hessTrace_eq_sum,TranscendentalNakamaye.Pending.LinearRegularity.norm_sub_le_of_zero,TranscendentalNakamaye.Pending.LinearRegularity.norm_fderiv_le_of_one,TranscendentalNakamaye.Pending.LinearRegularity.norm_fderiv_sub_le_of_one,TranscendentalNakamaye.Pending.LinearRegularity.differentiable_of_one,TranscendentalNakamaye.Pending.LinearRegularity.holderAt_second_partial,TranscendentalNakamaye.Pending.LinearRegularity.holderAt_coeff,TranscendentalNakamaye.Pending.LinearRegularity.holderAt_hessTrace,TranscendentalNakamaye.Pending.LinearRegularity.exists_inner_ge,TranscendentalNakamaye.Pending.LinearRegularity.exists_ellipticExtension,TranscendentalNakamaye.Pending.LinearRegularity.holderAt_fderiv_apply_of_hessTrace_eq,TranscendentalNakamaye.Pending.LinearRegularity.holderAt_three_of_hessTrace_eq,TranscendentalNakamaye.Pending.LinearRegularity.hessTrace_fderiv_apply_eq,TranscendentalNakamaye.Pending.LinearRegularity.holderAt_of_hessTrace_eq}
\leanok
Let \(E\) be a finite-dimensional real inner product space, \(0<r<1\), \(U\subset E\) open, and \(a\colon E\to\operatorname{End}(E)\) symmetric and positive definite
on \(U\). Let \(k\ge1\). If \(a\) and \(g\) are \(C^{k,r}\) and \(v\) is \(C^{k+1,r}\) near every point of \(U\), and \(a:D^2v=g\) on \(U\), then \(v\) is \(C^{k+2,r}\)
near every point of \(U\).
\end{lemma}
\begin{proof}
\leanok
For \(k=1\): the difference quotients satisfy \(a:D^2(D_hv)=D_hg-(D_ha):D^2v(\cdot+h\zeta)\), whose right-hand side is bounded and \(r\)-H\"older uniformly in
\(h\); extend \(a\) to whole-space elliptic coefficients of Schauder class and apply Lemma~\ref{bou:difference-quotient-regularity}. For \(k\to k+1\):
\(w=\partial_\zeta v\) solves \(a:D^2w=\partial_\zeta g-(\partial_\zeta a):D^2v\), whose right-hand side is \(C^{k,r}\) (Lemma~\ref{bou:holder-germ}).
\end{proof}

\begin{lemma}[The Monge--Amp\`ere operator on \(C^{2,r}\)]
\label{bou:monge-ampere-holder}
\uses{bou:holder-nemytskii,bou:kahler-laplacian-realizable,bou:holder-spaces-manifold}
\lean{TranscendentalNakamaye.Pending.MongeAmpereHolder.abs_sub_le_of_local,TranscendentalNakamaye.Pending.MongeAmpereHolder.coeffMatrix_transition,TranscendentalNakamaye.Pending.MongeAmpereHolder.complexHessianOfBilin_sub,TranscendentalNakamaye.Pending.MongeAmpereHolder.complexHessian_transition,TranscendentalNakamaye.Pending.MongeAmpereHolder.contDiffOn_det,TranscendentalNakamaye.Pending.MongeAmpereHolder.contDiffOn_entries,TranscendentalNakamaye.Pending.MongeAmpereHolder.contDiffOn_maLocal,TranscendentalNakamaye.Pending.MongeAmpereHolder.continuous_kfun,TranscendentalNakamaye.Pending.MongeAmpereHolder.exists_chart_estimates,TranscendentalNakamaye.Pending.MongeAmpereHolder.exists_cutoffModel,TranscendentalNakamaye.Pending.MongeAmpereHolder.exists_lipschitz_of_hasCompactSupport',TranscendentalNakamaye.Pending.MongeAmpereHolder.fderiv_fderiv_sub,TranscendentalNakamaye.Pending.MongeAmpereHolder.fderiv_maLocal_apply,TranscendentalNakamaye.Pending.MongeAmpereHolder.hasDerivAt_det_add_smul,TranscendentalNakamaye.Pending.MongeAmpereHolder.hermPart,TranscendentalNakamaye.Pending.MongeAmpereHolder.hermPart_add,TranscendentalNakamaye.Pending.MongeAmpereHolder.hermPart_kmat,TranscendentalNakamaye.Pending.MongeAmpereHolder.hermPart_smul,TranscendentalNakamaye.Pending.MongeAmpereHolder.isHermitian_complexHessian_of_contDiffAt,TranscendentalNakamaye.Pending.MongeAmpereHolder.isHermitian_hermPart,TranscendentalNakamaye.Pending.MongeAmpereHolder.isOpen_maDomain,TranscendentalNakamaye.Pending.MongeAmpereHolder.kfun,TranscendentalNakamaye.Pending.MongeAmpereHolder.kfun_sub,TranscendentalNakamaye.Pending.MongeAmpereHolder.kfun_sub',TranscendentalNakamaye.Pending.MongeAmpereHolder.kfun_sub_kfun,TranscendentalNakamaye.Pending.MongeAmpereHolder.kmat,TranscendentalNakamaye.Pending.MongeAmpereHolder.kmat_sub,TranscendentalNakamaye.Pending.MongeAmpereHolder.le_rpow_self_of_le_one,TranscendentalNakamaye.Pending.MongeAmpereHolder.lipschitzWith_kmat_complexHessianOfBilin,TranscendentalNakamaye.Pending.MongeAmpereHolder.ma,TranscendentalNakamaye.Pending.MongeAmpereHolder.maChart,TranscendentalNakamaye.Pending.MongeAmpereHolder.maChart_eq_of_mem,TranscendentalNakamaye.Pending.MongeAmpereHolder.maDomain,TranscendentalNakamaye.Pending.MongeAmpereHolder.maLocal,TranscendentalNakamaye.Pending.MongeAmpereHolder.ma_eq,TranscendentalNakamaye.Pending.MongeAmpereHolder.norm_I_smul_coordinateVector,TranscendentalNakamaye.Pending.MongeAmpereHolder.norm_coordinateVector,TranscendentalNakamaye.Pending.MongeAmpereHolder.norm_kfun_holder,TranscendentalNakamaye.Pending.MongeAmpereHolder.norm_kfun_le,TranscendentalNakamaye.Pending.MongeAmpereHolder.norm_kfun_sub_le',TranscendentalNakamaye.Pending.MongeAmpereHolder.norm_kmat_complexHessianOfBilin_le,TranscendentalNakamaye.Pending.MongeAmpereHolder.norm_kmat_le,TranscendentalNakamaye.Pending.MongeAmpereHolder.re_det_congr}
\leanok
Let \(a,g\) be Hermitian smooth real \((1,1)\)-forms on a complex manifold \(M\) with a H\"older atlas, \(0<r<1\), \(\Phi_0\in C^{2,r}(M)\) with
\(H_g=H_a+\partial\bar\partial\Phi_0\) in every chart, and \(\mathrm{ma}_a\psi=\log\det(H_a+\partial\bar\partial\psi)-\log\det H_a\) (in charts; it does not depend on
the chart). For each chart of the atlas there are \(\rho,C>0\) such that for \(\|\psi-\Phi_0\|_{C^{2,r}}<\rho\) the local representative of \(\mathrm{ma}_a\psi\) is in
\(C^{0,r}\), and for \(\psi_1,\psi_2\) in that ball the local representative of \(\mathrm{ma}_a\psi_1-\mathrm{ma}_a\psi_2-\Delta_g(\psi_1-\psi_2)\) has \(C^{0,r}\) norm
\(\le C(\|\psi_1-\Phi_0\|+\|\psi_2-\Phi_0\|)\|\psi_1-\psi_2\|\).
\end{lemma}
\begin{proof}
\leanok
Locally \(\mathrm{ma}_a\psi=\Gamma(z,\partial\bar\partial\psi(z))\) with \(\Gamma(z,K)=\log\det(H(z)+K)-\log\det H(z)\), smooth where the determinants are positive, and by
Jacobi's formula \(D_K\Gamma(z,K)L=\operatorname{Re}\operatorname{tr}((H+K)^{-1}L)\), the Laplacian of \(a+dd^c\psi\). Cut \(\Gamma\) off near the compact graph of
\(\partial\bar\partial\Phi_0\) over the chart piece and apply Lemma~\ref{bou:holder-nemytskii} with \(P_j(z)=(z,\partial\bar\partial\psi_j(z))\).
\end{proof}

\begin{lemma}[Positivity near a K\"ahler potential]
\label{bou:monge-ampere-positivity}
\uses{bou:holder-spaces-manifold,bou:yau-openness}
\lean{TranscendentalNakamaye.Pending.MongeAmperePositivity.conjTranspose_map_conj,TranscendentalNakamaye.Pending.MongeAmperePositivity.continuousOn_matCLM,TranscendentalNakamaye.Pending.MongeAmperePositivity.exists_chart_posDef,TranscendentalNakamaye.Pending.MongeAmperePositivity.exists_inner_ge_of_isCompact,TranscendentalNakamaye.Pending.MongeAmperePositivity.exists_ma_eq_near_posDef,TranscendentalNakamaye.Pending.MongeAmperePositivity.exists_posDef_near,TranscendentalNakamaye.Pending.MongeAmperePositivity.im_normalizedHermitianMatrix,TranscendentalNakamaye.Pending.MongeAmperePositivity.matCLM_add,TranscendentalNakamaye.Pending.MongeAmperePositivity.norm_matCLM_le,TranscendentalNakamaye.Pending.MongeAmperePositivity.re_normalizedHermitianMatrix,TranscendentalNakamaye.Pending.MongeAmperePositivity.re_star_dotProduct_mulVec}
\leanok
Let \(M\) be compact with a H\"older atlas and \(a+dd^c\Phi_0\) Hermitian. Then \(H_a+\partial\bar\partial\psi\) is positive definite in every chart for every
\(\psi\in C^{2,r}(M)\) close enough to \(\Phi_0\). Consequently the solutions \(\psi\) of Lemma~\ref{bou:yau-openness} can be taken with \(A+dd^c\psi>0\).
\end{lemma}
\begin{proof}
\leanok
On the compact pieces of the atlas there is a uniform margin of positivity, and a perturbation of \(\partial\bar\partial\psi\) small in \(C^{0}\) keeps it; in an
arbitrary chart use the congruence \(H_e+K_e=J^T(H_i+K_i)\bar J\) along the chart transition.
\end{proof}

\begin{lemma}[Openness]
\label{bou:yau-openness}
\uses{bou:yau-continuity-path,bou:kahler-laplacian-realizable,bou:monge-ampere-holder,bou:holder-spaces-manifold}
\source{szekelyhidi-extremal:page-0054,szekelyhidi-extremal:page-0055,szekelyhidi-extremal:page-0056,szekelyhidi-extremal:page-0067}
\lean{TranscendentalNakamaye.Pending.YauOpenness.exists_atlas_estimates,TranscendentalNakamaye.Pending.YauOpenness.exists_eq_smul_near,TranscendentalNakamaye.Pending.YauOpenness.exists_ma_eq_near,TranscendentalNakamaye.Pending.YauOpenness.exists_ma_eq_near_of_smooth,TranscendentalNakamaye.Pending.YauOpenness.ma_eq_of_density,TranscendentalNakamaye.Pending.YauOpenness.normalizedHermitianMatrix_add_ddc}
\leanok
Let \((Y,A)\) be a compact K\"ahler manifold of dimension \(N\) with a H\"older atlas, \(0<r<1\), \(F\) smooth and \(\varepsilon>0\). If \(\varphi_0\) is smooth,
\(A_{\varphi_0}\) is K\"ahler and \(A_{\varphi_0}^N=e^{t_0F+\varepsilon\varphi_0}A^N\), then for every \(\rho>0\) there is \(\delta>0\) such that for \(|t-t_0|<\delta\) there is
\(\psi\in C^{2,r}(Y)\) with \(\|\psi-\varphi_0\|_{C^{2,r}}<\rho\) and \(\mathrm{ma}_A\psi=tF+\varepsilon\psi\).
\end{lemma}
\begin{proof}
\leanok
Summing Lemma~\ref{bou:monge-ampere-holder} over the atlas, \(\psi\mapsto\mathrm{ma}_A\psi\) is strictly differentiable at \(\varphi_0\) from \(C^{2,r}\) to \(C^{0,r}\) with
derivative \(\Delta_{A_{\varphi_0}}\); so \(\psi\mapsto\mathrm{ma}_A\psi-\varepsilon\psi\) has the derivative \(\Delta_{A_{\varphi_0}}-\varepsilon\), an isomorphism
(Lemma~\ref{bou:kahler-laplacian-realizable}). The inverse function theorem for \((\psi,t)\mapsto(\mathrm{ma}_A\psi-\varepsilon\psi-tF,t)\) at \((\varphi_0,t_0)\)
solves \(\mathrm{ma}_A\psi-\varepsilon\psi=tF\) for \(t\) near \(t_0\) with \(\psi\) near \(\varphi_0\).
\end{proof}

\begin{lemma}[Regularity]
\label{bou:yau-regularity}
\uses{bou:holder-spaces-manifold,bou:monge-ampere-holder,bou:yau-openness,bou:monge-ampere-positivity,bou:holder-germ,bou:difference-quotient-regularity,bou:linear-elliptic-regularity}
\source{szekelyhidi-extremal:page-0056}
\lean{TranscendentalNakamaye.Pending.MongeAmpereRegularity.contDiffOn_lapCoeffField,TranscendentalNakamaye.Pending.MongeAmpereRegularity.continuousOn_matCLM_field,TranscendentalNakamaye.Pending.MongeAmpereRegularity.exists_ellipticModel,TranscendentalNakamaye.Pending.MongeAmpereRegularity.exists_inner_ge',TranscendentalNakamaye.Pending.MongeAmpereRegularity.hessTraceLeft,TranscendentalNakamaye.Pending.MongeAmpereRegularity.hessTraceLeft_apply,TranscendentalNakamaye.Pending.MongeAmpereRegularity.holderAt_fderiv_apply_of_maLocal_eq,TranscendentalNakamaye.Pending.MongeAmpereRegularity.holderAt_three_of_maLocal_eq,TranscendentalNakamaye.Pending.MongeAmpereRegularity.isOpen_pdDomain,TranscendentalNakamaye.Pending.MongeAmpereRegularity.kfun_diffQuot,TranscendentalNakamaye.Pending.MongeAmpereRegularity.norm_sub_le_rpow_of_le,TranscendentalNakamaye.Pending.MongeAmpereRegularity.pdDomain,TranscendentalNakamaye.Pending.MongeAmpereRegularity.pdDomain_subset_maDomain,TranscendentalNakamaye.Pending.MongeAmpereRegularity.posDef_of_inner_pos,TranscendentalNakamaye.Pending.MongeAmpereRegularity.chartSmooth_of_ma_eq,TranscendentalNakamaye.Pending.MongeAmpereRegularity.contDiffOn_of_maLocal_eq,TranscendentalNakamaye.Pending.MongeAmpereRegularity.fderiv_fderiv_fderiv_apply,TranscendentalNakamaye.Pending.MongeAmpereRegularity.hasDerivAt_kfun_line,TranscendentalNakamaye.Pending.MongeAmpereRegularity.hessTrace_lapCoeff_eq,TranscendentalNakamaye.Pending.MongeAmpereRegularity.holderAt_comp_of_contDiffOn,TranscendentalNakamaye.Pending.MongeAmpereRegularity.holderAt_complex,TranscendentalNakamaye.Pending.MongeAmpereRegularity.holderAt_kfun,TranscendentalNakamaye.Pending.MongeAmpereRegularity.holderAt_of_maLocal_eq,TranscendentalNakamaye.Pending.MongeAmpereRegularity.holderAt_pi,TranscendentalNakamaye.Pending.MongeAmpereRegularity.kfun_apply,TranscendentalNakamaye.Pending.YauSmoothOpenness.density_eq_of_ma_eq,TranscendentalNakamaye.Pending.YauSmoothOpenness.exists_smooth_ma_eq_near,TranscendentalNakamaye.Pending.YauSmoothOpenness.isKahler_of_posDef,TranscendentalNakamaye.Pending.YauSmoothOpenness.isStrictlyPositive_of_posDef}
\leanok
Let \((Y,A)\), \(F\) and \(\varepsilon\) be as in Lemma~\ref{bou:yau-openness}. A solution \(\psi\in C^{2,r}(Y)\) of \(\mathrm{ma}_A\psi=tF+\varepsilon\psi\) with
\(A+dd^c\psi>0\) is smooth; in particular, the set \(S_\varepsilon\) of Lemma~\ref{bou:yau-continuity-path} is open (Lemmas~\ref{bou:yau-openness}
and~\ref{bou:monge-ampere-positivity}).
\end{lemma}
\begin{proof}
\leanok
In a chart \(e\), \(u=\psi\circ e^{-1}\) is \(C^{2,r}\) near every point and solves \(\log\det(H_A+\partial\bar\partial u)-\log\det H_A=tF\circ e^{-1}+\varepsilon u\)
with \(H_A+\partial\bar\partial u\) positive definite. A difference quotient \(\delta_hu\) solves a linear elliptic equation
\(\operatorname{Re}\operatorname{tr}(Q_h\,\partial\bar\partial\delta_hu)=\delta_h(tF\circ e^{-1})+\varepsilon\delta_hu\) whose coefficients
\(Q_h=\int_0^1(H_A+\partial\bar\partial(u+sh\delta_hu))^{-1}ds\) are elliptic and bounded in \(C^{0,r}\) uniformly in \(h\) on compact pieces of the chart, so
\(u\) is \(C^{3,r}\) near every point (Lemma~\ref{bou:difference-quotient-regularity}, after cutting off). If \(u\) is \(C^{k+2,r}\), \(k\ge1\), every
directional derivative \(v=\partial_\zeta u\) solves the linear equation obtained by differentiating the equation along \(\zeta\), whose coefficients
\((H_A+\partial\bar\partial u)^{-1}\) and right-hand side are \(C^{k,r}\) (Lemma~\ref{bou:holder-germ}); so \(v\) is \(C^{k+2,r}\)
(Lemma~\ref{bou:linear-elliptic-regularity}) and \(u\) is \(C^{k+3,r}\). A germ that is \(C^{k,r}\) for every \(k\) is smooth.

For the openness of \(S_\varepsilon\): if \(t_0\in S_\varepsilon\) with the smooth solution \(\varphi_0\), Lemmas~\ref{bou:yau-openness}
and~\ref{bou:monge-ampere-positivity} give, for \(t\) near \(t_0\), \(\psi\in C^{2,r}(Y)\) with \(\mathrm{ma}_A\psi=tF+\varepsilon\psi\) and \(H_A+\partial\bar\partial\psi\)
positive definite in every chart. By the first part \(\psi\) is smooth, \(A+dd^c\psi\) is K\"ahler (its coefficient matrix is positive definite in every chart),
and \(\mathrm{ma}_A\psi=tF+\varepsilon\psi\) is the equation \(A_\psi^N=e^{tF+\varepsilon\psi}A^N\) read in the charts; so \(t\in S_\varepsilon\).
\end{proof}

\begin{lemma}[The maximum principle in trace form]
\label{bou:complex-hessian-trace-max}
\uses{bou:max-principle}
\source{szekelyhidi-extremal:page-0060}
\lean{TranscendentalNakamaye.Pending.AubinYauMatrixBounds.re_dotProduct_mulVec_complexHessian_nonpos,TranscendentalNakamaye.Pending.AubinYauMatrixBounds.re_trace_mul_complexHessian_nonpos}
\leanok
Let \(u\colon\mathbf C^n\to\mathbf R\) be of class \(C^2\) near \(z\) with a local maximum at \(z\), and \(H=(\partial_j\partial_{\bar k}u(z))\) its complex Hessian.
Then \(y^*Hy\le0\) for every \(y\in\mathbf C^n\), and \(\operatorname{Re}\operatorname{tr}(PH)\le0\) for every positive semidefinite Hermitian \(n\times n\) matrix \(P\).
\end{lemma}
\begin{proof}
\leanok
\(y^*Hy\) is \(\frac14\bigl(D^2u(z)(v,v)+D^2u(z)(iv,iv)\bigr)\) for the real vector \(v=\sum_jy_je_j\), and both terms are \(\le0\) by Lemma~\ref{bou:max-principle}.
Write \(P=\sum_iv_iv_i^*\); then \(\operatorname{tr}(PH)=\sum_iv_i^*Hv_i\).
\end{proof}
With \(P=(g'_{j\bar k})^{-1}\) this is \(\Delta'H(z)\le0\) at a maximum point, the form in which Lemma~\ref{bou:yau-c2-estimate} uses the maximum principle.

\begin{lemma}[Eigenvalue bounds for two positive definite matrices]
\label{bou:yau-c2-matrix-endgame}
\source{szekelyhidi-extremal:page-0060}
\lean{TranscendentalNakamaye.Pending.AubinYauMatrixBounds.le_of_sum_inv_le_of_prod_le,TranscendentalNakamaye.Pending.AubinYauMatrixBounds.le_of_sum_le_of_le_prod,TranscendentalNakamaye.Pending.AubinYauMatrixBounds.star_mulVec_dotProduct_mulVec,TranscendentalNakamaye.Pending.AubinYauMatrixBounds.re_dotProduct_diagonal_mulVec,TranscendentalNakamaye.Pending.AubinYauMatrixBounds.exists_frame,TranscendentalNakamaye.Pending.AubinYauMatrixBounds.quad_le_of_trace_inv_le,TranscendentalNakamaye.Pending.AubinYauMatrixBounds.quad_le_of_trace_le}
\leanok
Let \(G,G'\) be positive definite Hermitian \(n\times n\) matrices and write \(q(x)=x^*Gx\), \(q'(x)=x^*G'x\) (real numbers).
\begin{enumerate}
\item There are \(\lambda_1,\dots,\lambda_n>0\) with \(\operatorname{tr}(G'^{-1}G)=\sum_k\lambda_k^{-1}\), \(\operatorname{tr}(G^{-1}G')=\sum_k\lambda_k\),
\(\det G'=\prod_k\lambda_k\det G\); if \(\lambda_k\le c\) for all \(k\), then \(q'\le cq\), and if \(1\le c\lambda_k\) for all \(k\), then \(q\le cq'\).
\item If \(\operatorname{Re}\operatorname{tr}(G'^{-1}G)\le K\) and \(\det G'\le D\det G\), then \(q'\le DK^{n-1}q\).
\item If \(\operatorname{Re}\operatorname{tr}(G^{-1}G')\le T\) and \(d\det G\le\det G'\) with \(d>0\), then \(q'\le Tq\) and \(q\le(T^{n-1}/d)\,q'\).
\end{enumerate}
\end{lemma}
\begin{proof}
\leanok
(1) Simultaneous diagonalisation: there is an invertible \(R\) with \(R^*GR=I\) and \(R^*G'R=\operatorname{diag}(\lambda_k)\), \(\lambda_k>0\); then
\(q'(Ry)=\sum_k\lambda_k|y_k|^2\) and \(q(Ry)=\sum_k|y_k|^2\), and the traces and determinants are computed in this frame.
(2) Each \(\lambda_j^{-1}\le\sum_k\lambda_k^{-1}\le K\), so \(\lambda_j\ge K^{-1}\), and \(\lambda_k\cdot K^{-(n-1)}\le\prod_j\lambda_j\le D\); apply (1).
(3) Each \(\lambda_k\le\sum_j\lambda_j\le T\), and \(d\le\prod_j\lambda_j\le\lambda_kT^{n-1}\); apply (1).
\end{proof}

\begin{lemma}[The largest relative eigenvalue]
\label{bou:relative-eigenvalue-max}
\uses{ana:geometric-setting}
\source{lnm2038-monge-ampere:page-0222}
\lean{TranscendentalNakamaye.Pending.RelativeEigenvalueMax.chartTransitionDerivative_bijective,TranscendentalNakamaye.Pending.RelativeEigenvalueMax.chartTransition_apply_chart,TranscendentalNakamaye.Pending.RelativeEigenvalueMax.coeff_line_smul,TranscendentalNakamaye.Pending.RelativeEigenvalueMax.coeff_line_transfer,TranscendentalNakamaye.Pending.RelativeEigenvalueMax.exists_max_log_ratio_add,TranscendentalNakamaye.Pending.RelativeEigenvalueMax.line,TranscendentalNakamaye.Pending.RelativeEigenvalueMax.mem_chartTransition_source,TranscendentalNakamaye.Pending.RelativeEigenvalueMax.ratio_chart_transfer,TranscendentalNakamaye.Pending.RelativeEigenvalueMax.ratio_normalize,TranscendentalNakamaye.Pending.RelativeEigenvalueMax.ratio_smul}
\leanok
Let \(Y\) be a nonempty compact complex manifold of dimension \(N\ge1\) and \(a,b\) smooth real \(2\)-forms on \(Y\) with \(a(v,iv)>0\) and \(b(v,iv)>0\) for \(v\ne0\) in every chart.
(1) For two charts \(e,f\) at \(x\) and \(v\ne0\), the ratio \(b(v,iv)/a(v,iv)\) at \(x\) read in \(e\) equals the ratio at \(x\) read in \(f\) for some unit vector. (2) For every
continuous \(f\colon Y\to\mathbf R\) there are \(y\in Y\) and \(M\) such that \(\log\bigl(b(v,iv)/a(v,iv)\bigr)+f\le M\) in every chart, at every point, for every \(v\ne0\), and in
every chart at \(y\) some unit vector \(\zeta\) gives equality. In particular \(\eta=\max_{v\ne0}b(v,iv)/a(v,iv)\) is independent of the chart and \(\log\eta+f\) attains its maximum.
\end{lemma}
\begin{proof}
\leanok
(1) The derivative of the chart transition is \(\mathbf C\)-linear and bijective and transports \(a\) and \(b\); the ratio is invariant under real rescaling of \(v\).
(2) Cover \(Y\) by finitely many closed chart balls; the function \((z,v)\mapsto\log(b/a)(z)(v,iv)+f\) is continuous on the compact set of pairs (point of a ball, unit
vector) and attains its maximum there; by (1) that maximum bounds the values in every chart.
\end{proof}

\begin{lemma}[Derivatives of \(\det\) and of \(\log\det\) of a complex Hessian]
\label{bou:log-det-derivatives}
\source{lnm2038-monge-ampere:page-0223}
\lean{TranscendentalNakamaye.Pending.WirtingerLogDet.adjugate_eq_det_smul_inv',TranscendentalNakamaye.Pending.WirtingerLogDet.contDiff_complexHessian_apply',TranscendentalNakamaye.Pending.WirtingerLogDet.contDiff_det_complexHessian,TranscendentalNakamaye.Pending.WirtingerLogDet.det_eq_sum,TranscendentalNakamaye.Pending.WirtingerLogDet.det_updateCol_eq_sum,TranscendentalNakamaye.Pending.WirtingerLogDet.differentiableAt_adjugate_apply,TranscendentalNakamaye.Pending.WirtingerLogDet.differentiableAt_det,TranscendentalNakamaye.Pending.WirtingerLogDet.differentiableAt_inv_apply,TranscendentalNakamaye.Pending.WirtingerLogDet.differentiable_complexHessian_apply,TranscendentalNakamaye.Pending.WirtingerLogDet.inv_apply_eq,TranscendentalNakamaye.Pending.WirtingerLogDet.isHermitian_complexHessian,TranscendentalNakamaye.Pending.WirtingerLogDet.leviForm_log_det_complexHessian,TranscendentalNakamaye.Pending.WirtingerLogDet.ofReal_re_det_of_isHermitian,TranscendentalNakamaye.Pending.WirtingerLogDet.wirtM,TranscendentalNakamaye.Pending.WirtingerLogDet.wirtM_I_complexHessian,TranscendentalNakamaye.Pending.WirtingerLogDet.wirtM_apply,TranscendentalNakamaye.Pending.WirtingerLogDet.wirt_I_conj,TranscendentalNakamaye.Pending.WirtingerLogDet.wirt_congr,TranscendentalNakamaye.Pending.WirtingerLogDet.wirt_det,TranscendentalNakamaye.Pending.WirtingerLogDet.wirt_det_eq_det_mul_trace,TranscendentalNakamaye.Pending.WirtingerLogDet.wirt_eq_zero_of_eventuallyEq_const,TranscendentalNakamaye.Pending.WirtingerLogDet.wirt_inv_apply,TranscendentalNakamaye.Pending.WirtingerLogDet.wirt_ofReal_log}
\leanok
Let \(\partial_v\) be a directional Wirtinger derivative on a complex normed space, applied entrywise to matrix-valued functions. (1) If \(M\) is differentiable at \(x\),
then \(\partial_v\det M=\operatorname{tr}(\operatorname{adj}M\,\partial_vM)\) at \(x\), \(=\det M\operatorname{tr}(M^{-1}\partial_vM)\) when \(M(x)\) is invertible, and then
\(\partial_v(M^{-1})=-M^{-1}(\partial_vM)M^{-1}\). (2) Let \(u\colon\mathbf C^n\to\mathbf R\) be smooth, \(H=(u_{j\bar k})\) its complex Hessian, positive definite at \(x\), and
\(\zeta\in\mathbf C^n\). Then at \(x\)
\[
 (\log\det H)_{\zeta\bar\zeta}=\operatorname{tr}\bigl(H^{-1}\,\partial_{\bar\zeta}\partial_\zeta H\bigr)-\operatorname{tr}\bigl(H^{-1}\,\partial_{\bar\zeta}H\,H^{-1}\,\partial_\zeta H\bigr).
\]
\end{lemma}
\begin{proof}
\leanok
(1) Expand \(\det M=\sum_\sigma\operatorname{sgn}\sigma\prod_iM_{\sigma(i)i}\) and differentiate the products; the cofactor expansion along a column identifies the result with the
trace against the adjugate. Differentiating \(MM^{-1}=I\) gives the derivative of the inverse. (2) \(H\) is Hermitian, so \(\det H\) is real and positive near \(x\); apply (1)
twice to \(\log\det H\) (\(\partial_\zeta\log\det H=\operatorname{tr}(H^{-1}\partial_\zeta H)\)).
\end{proof}

\begin{lemma}[B\l{}ocki's local inequality]
\label{bou:yau-blocki-lemma}
\uses{bou:log-det-derivatives}
\source{lnm2038-monge-ampere:page-0223}
\lean{TranscendentalNakamaye.Pending.BlockiLocalEstimate.blocki_local_estimate,TranscendentalNakamaye.Pending.BlockiLocalEstimate.complexHessian_apply_eq_wirt,TranscendentalNakamaye.Pending.BlockiLocalEstimate.complexHessian_log,TranscendentalNakamaye.Pending.BlockiLocalEstimate.contDiff_leviForm,TranscendentalNakamaye.Pending.BlockiLocalEstimate.contDiff_ofReal_comp,TranscendentalNakamaye.Pending.BlockiLocalEstimate.coordinateVector_sum,TranscendentalNakamaye.Pending.BlockiLocalEstimate.leviForm_eq_re_wirt,TranscendentalNakamaye.Pending.BlockiLocalEstimate.norm_sum_mul_sq_le,TranscendentalNakamaye.Pending.BlockiLocalEstimate.ofReal_leviForm_eq_dotProduct,TranscendentalNakamaye.Pending.BlockiLocalEstimate.ofReal_leviForm_eq_wirt,TranscendentalNakamaye.Pending.BlockiLocalEstimate.re_star_dotProduct_inv_mulVec_le,TranscendentalNakamaye.Pending.BlockiLocalEstimate.re_star_dotProduct_self,TranscendentalNakamaye.Pending.BlockiLocalEstimate.re_star_mulVec_dotProduct_le,TranscendentalNakamaye.Pending.BlockiLocalEstimate.re_trace_conjTranspose_mul_self,TranscendentalNakamaye.Pending.BlockiLocalEstimate.wirt_I_ofReal_eq_conj,TranscendentalNakamaye.Pending.BlockiLocalEstimate.wirt_comm_fun,TranscendentalNakamaye.Pending.BlockiLocalEstimate.wirt_inv',TranscendentalNakamaye.Pending.BlockiLocalEstimate.wirt_leviForm,TranscendentalNakamaye.Pending.BlockiLocalEstimate.wirt_wirt_leviForm}
\leanok
Let \(u\colon\mathbf C^n\to\mathbf R\) be smooth, \(x\) a point at which \(H=(u_{j\bar k})\) is positive definite, \(F=\det(u_{j\bar k})\), and \((u^{j\bar k})=H(x)^{-1}\). For every
\(\zeta\in\mathbf C^n\setminus\{0\}\), with \(u_{\zeta\bar\zeta}=\sum_{j,k}\zeta_j\bar\zeta_ku_{j\bar k}\), at \(x\)
\[
 \frac{(\log F)_{\zeta\bar\zeta}}{u_{\zeta\bar\zeta}}\le u^{j\bar k}(\log u_{\zeta\bar\zeta})_{j\bar k} .
\]
\end{lemma}
\begin{proof}
\leanok
By Lemma~\ref{bou:log-det-derivatives}, \(u^{j\bar k}u_{j\bar k\zeta}=(\log F)_\zeta\) and
\(u^{j\bar k}u_{j\bar k\zeta\bar\zeta}=(\log F)_{\zeta\bar\zeta}+u^{j\bar l}u^{m\bar k}u_{j\bar k\zeta}u_{m\bar l\bar\zeta}\). Since fourth derivatives commute,
\(u_{\zeta\bar\zeta j\bar k}=u_{j\bar k\zeta\bar\zeta}\), and the Hessian of a logarithm gives
\[
 u_{\zeta\bar\zeta}\,u^{j\bar k}(\log u_{\zeta\bar\zeta})_{j\bar k}=(\log F)_{\zeta\bar\zeta}+u^{j\bar l}u^{m\bar k}u_{j\bar k\zeta}u_{m\bar l\bar\zeta}
 -\frac1{u_{\zeta\bar\zeta}}u^{j\bar k}u_{\zeta\bar\zeta j}u_{\zeta\bar\zeta\bar k}.
\]
With \(X=(u_{j\bar k\zeta})\) and \(w=\bar\zeta\), \(u_{\zeta\bar\zeta j}=(Xw)_j\), and the Cauchy--Schwarz inequality in the form
\((Xw)^*H^{-1}(Xw)\le(w^*Hw)\operatorname{tr}(H^{-1}X^*H^{-1}X)\) (in a frame where \(H=I\)) shows that the last term is at most the middle one.
\end{proof}

\begin{lemma}[Uniform constants for a local potential]
\label{bou:yau-c2-uniform-constants}
\source{lnm2038-monge-ampere:page-0224}
\lean{TranscendentalNakamaye.Pending.AubinYauUniformConstants.complexHessian_congr,TranscendentalNakamaye.Pending.AubinYauUniformConstants.contDiff_fourthCoeff,TranscendentalNakamaye.Pending.AubinYauUniformConstants.exists_abs_leviForm_le,TranscendentalNakamaye.Pending.AubinYauUniformConstants.exists_contDiff_eqOn_ball,TranscendentalNakamaye.Pending.AubinYauUniformConstants.exists_uniform_constants,TranscendentalNakamaye.Pending.AubinYauUniformConstants.fourthCoeff,TranscendentalNakamaye.Pending.AubinYauUniformConstants.leviForm_congr,TranscendentalNakamaye.Pending.AubinYauUniformConstants.leviForm_leviForm_le,TranscendentalNakamaye.Pending.AubinYauUniformConstants.ofReal_leviForm_leviForm_eq_sum,TranscendentalNakamaye.Pending.AubinYauUniformConstants.wirt_I_sum_dir,TranscendentalNakamaye.Pending.AubinYauUniformConstants.wirt_I_wirt_negI_expand,TranscendentalNakamaye.Pending.AubinYauUniformConstants.wirt_negI_eq_dDir,TranscendentalNakamaye.Pending.AubinYauUniformConstants.wirt_negI_sum_dir,TranscendentalNakamaye.Pending.AubinYauUniformConstants.wirt_wirt_sum_const_mul}
\leanok
(1) Let \(n\ge1\), \(\rho\colon\mathbf C^n\to\mathbf R\) smooth with \(\rho_{w\bar w}(z)>0\) for \(w\ne0\) and \(z\) in an open set \(V\), and \(K\subset V\) compact. There are \(c>0\) and
\(C,M,K_1\ge0\) such that for \(z\in K\), all \(w\) and all \(|\zeta|\le1\): \(c|w|^2\le\rho_{w\bar w}(z)\le C|w|^2\), \((\rho_{\zeta\bar\zeta})_{w\bar w}(z)\le M|w|^2\) and
\((\log\det(\rho_{j\bar k}))_{\zeta\bar\zeta}(z)\ge-K_1\). (2) A function smooth on an open set \(W\) agrees, on a ball \(B(c,r)\) with \(\overline B(c,R)\subset W\), \(r<R\), with a smooth
function on \(\mathbf C^n\).
\end{lemma}
\begin{proof}
\leanok
(1) The lower bound is a compactness argument on \(K\times\{|w|=1\}\); the upper bounds hold because \((\rho_{\zeta\bar\zeta})_{w\bar w}\) is a combination of fourth derivatives of
\(\rho\) with coefficients \(\bar w_aw_b\bar\zeta_c\zeta_d\), and \((\log\det(\rho_{j\bar k}))_{\zeta\bar\zeta}\) is continuous in \((z,\zeta)\) on a compact set. (2) Multiply by a smooth
bump function equal to \(1\) on \(B(c,r)\) with support in \(\overline B(c,R)\).
\end{proof}

\begin{lemma}[The local step of the second-order estimate]
\label{bou:yau-c2-local-step}
\uses{bou:complex-hessian-trace-max,bou:yau-blocki-lemma}
\source{lnm2038-monge-ampere:page-0224}
\lean{TranscendentalNakamaye.Pending.AubinYauLocalStep.complexHessian_const_mul,TranscendentalNakamaye.Pending.AubinYauLocalStep.complexHessian_fun_add,TranscendentalNakamaye.Pending.AubinYauLocalStep.contDiffAt_two,TranscendentalNakamaye.Pending.AubinYauLocalStep.coordCombination_ne_zero,TranscendentalNakamaye.Pending.AubinYauLocalStep.le_re_trace_mul_complexHessian,TranscendentalNakamaye.Pending.AubinYauLocalStep.norm_sq_coordCombination,TranscendentalNakamaye.Pending.AubinYauLocalStep.one_le_leviForm_mul_re_trace_inv,TranscendentalNakamaye.Pending.AubinYauLocalStep.posDef_complexHessian,TranscendentalNakamaye.Pending.AubinYauLocalStep.re_trace_inv_le_of_isLocalMax,TranscendentalNakamaye.Pending.AubinYauLocalStep.re_trace_mul_complexHessian_eq_sum,TranscendentalNakamaye.Pending.AubinYauLocalStep.re_trace_mul_complexHessian_le,TranscendentalNakamaye.Pending.AubinYauLocalStep.re_trace_mul_complexHessian_log_le,TranscendentalNakamaye.Pending.AubinYauLocalStep.re_trace_nonneg,TranscendentalNakamaye.Pending.AubinYauLocalStep.re_trace_sum_vecMulVec,TranscendentalNakamaye.Pending.AubinYauLocalStep.trace_mul_vecMulVec,TranscendentalNakamaye.Pending.AubinYauLocalStep.trace_vecMulVec_mul}
\leanok
Let \(u,\rho\colon\mathbf C^n\to\mathbf R\) be smooth, \(|\zeta|=1\), \(z_0\) a point where \(H=(u_{j\bar k}(z_0))\) is positive definite, and \(c_0>0\), \(M_0\ge0\), \(K_0\ge0\), \(\kappa\)
constants with \(c_0|w|^2\le\rho_{w\bar w}(z_0)\) and \((\rho_{\zeta\bar\zeta})_{w\bar w}(z_0)\le M_0|w|^2\) for all \(w\), \((\log\det(u_{j\bar k}))_{\zeta\bar\zeta}(z_0)\ge-K_0\), and
\(\kappa c_0\ge K_0+M_0/c_0+1\). If \(z_0\) is a local maximum of \(\log u_{\zeta\bar\zeta}-\log\rho_{\zeta\bar\zeta}-\kappa(u-\rho)\), then \(\operatorname{tr}H^{-1}\le\kappa n\).
\end{lemma}
\begin{proof}
\leanok
By Lemma~\ref{bou:complex-hessian-trace-max} with \(P=H^{-1}\), the trace of \(H^{-1}\) against the complex Hessian of the function is \(\le0\) at \(z_0\). The \(u\)-term is
\(\ge(\log\det(u_{j\bar k}))_{\zeta\bar\zeta}/u_{\zeta\bar\zeta}\ge-K_0/u_{\zeta\bar\zeta}\) (Lemma~\ref{bou:yau-blocki-lemma}); the \(\rho\)-term is at most
\((M_0/c_0)\operatorname{tr}H^{-1}\) (the gradient term of a logarithm has a sign); \(-\kappa\) times the trace of \(H^{-1}(u-\rho)_{j\bar k}\) is
\(-\kappa n+\kappa\operatorname{tr}(H^{-1}\rho_{j\bar k})\ge-\kappa n+\kappa c_0\operatorname{tr}H^{-1}\); and \(u_{\zeta\bar\zeta}\ge\lambda_{\min}(H)\ge1/\operatorname{tr}H^{-1}\). Hence
\(0\ge(\kappa c_0-K_0-M_0/c_0)\operatorname{tr}H^{-1}-\kappa n\ge\operatorname{tr}H^{-1}-\kappa n\).
\end{proof}

\begin{lemma}[The second-order estimate]
\label{bou:yau-c2-estimate}
\uses{bou:relative-eigenvalue-max,bou:complex-hessian-trace-max,bou:yau-c2-matrix-endgame,bou:yau-blocki-lemma,bou:yau-c2-local-step,bou:yau-c2-uniform-constants}
\source{lnm2038-monge-ampere:page-0222,lnm2038-monge-ampere:page-0223,lnm2038-monge-ampere:page-0224}
\lean{TranscendentalNakamaye.Pending.AubinYauEstimate.coeff_add_line_eq_leviForm,TranscendentalNakamaye.Pending.AubinYauEstimate.coeff_line_eq_leviForm,TranscendentalNakamaye.Pending.AubinYauEstimate.contDiffAt_two_of_contDiffOn,TranscendentalNakamaye.Pending.AubinYauEstimate.ddc_coeff_line_eq_leviForm,TranscendentalNakamaye.Pending.AubinYauEstimate.exists_le_smul_of_mongeAmpere_eq_exp,TranscendentalNakamaye.Pending.AubinYauEstimate.exists_piece,TranscendentalNakamaye.Pending.AubinYauEstimate.leviForm_pos_of_isKahler,TranscendentalNakamaye.Pending.AubinYauEstimate.re_det_complexHessian_eq_exp,TranscendentalNakamaye.Pending.AubinYauEstimate.re_det_complexHessian_potential,TranscendentalNakamaye.Pending.AubinYauEstimate.re_quad_complexHessian,TranscendentalNakamaye.Pending.AubinYauEstimate.re_quad_complexHessian_star,TranscendentalNakamaye.Pending.AubinYauEstimate.re_trace_inv_mul_le_of_quad_le}
\leanok
Let \((Y,A)\) be a compact connected K\"ahler manifold of dimension \(N\) and \(B_\varphi,B_G,L\in\mathbf R\). There is \(C>0\) such that every pair of smooth functions
\(\varphi,G\) with \(A_\varphi=A+dd^c\varphi\) K\"ahler, \(A_\varphi^N=e^GA^N\), \(\varphi(x)-\varphi(y)\le B_\varphi\) for all \(x,y\), \(|G|\le B_G\) and \(dd^cG\ge-LA\) satisfies
\(C^{-1}A\le A_\varphi\le CA\).
\end{lemma}
\begin{proof}
\leanok
Cover \(Y\) by finitely many charts in which \(A=dd^c\rho\) with \(\rho\) smooth up to the boundary, with constants uniform over the cover
(Lemma~\ref{bou:yau-c2-uniform-constants}); write \(g_{j\bar k}=\rho_{j\bar k}\) and
\(u=\rho+\varphi\), so \(u_{j\bar k}=g_{j\bar k}+\varphi_{j\bar k}\) is positive definite and \(\det(u_{j\bar k})=e^G\det(g_{j\bar k})\). The function
\(\eta=\max_{\zeta\ne0}u_{\zeta\bar\zeta}/g_{\zeta\bar\zeta}\), the largest eigenvalue of \(A_\varphi\) relative to \(A\), does not depend on the chart, and \(\alpha=\log\eta-K\varphi\), with \(K\) fixed below,
attains its maximum at some \(y\), where \(\eta(y)=u_{\zeta\bar\zeta}/g_{\zeta\bar\zeta}\) for a unit vector \(\zeta\) (Lemma~\ref{bou:relative-eigenvalue-max}). Then \(\tilde\alpha=\log u_{\zeta\bar\zeta}-\log g_{\zeta\bar\zeta}-K\varphi\le\alpha\) near \(y\), with equality at \(y\), so \(\tilde\alpha\) has a local maximum at
\(y\), and \(u^{j\bar k}\tilde\alpha_{j\bar k}(y)\le0\) (Lemma~\ref{bou:complex-hessian-trace-max} with \(P=(u^{j\bar k})\)). By Lemma~\ref{bou:yau-blocki-lemma} (applied to a smooth function on \(\mathbf C^n\) equal to \(u\) near \(y\)) with
\(\log F=G+\log\det(g_{j\bar k})\), and \(u^{j\bar k}\varphi_{j\bar k}=N-u^{j\bar k}g_{j\bar k}\),
\[
 0\ge u^{j\bar k}\tilde\alpha_{j\bar k}\ge\frac{G_{\zeta\bar\zeta}+(\log\det g)_{\zeta\bar\zeta}}{u_{\zeta\bar\zeta}}-u^{j\bar k}(\log g_{\zeta\bar\zeta})_{j\bar k}-KN+K\,u^{j\bar k}g_{j\bar k}
\]
at \(y\). The hypothesis gives \(G_{\zeta\bar\zeta}\ge-Lg_{\zeta\bar\zeta}\); the charts give \((\log\det g)_{\zeta\bar\zeta}\ge-C_1g_{\zeta\bar\zeta}\) and
\(|u^{j\bar k}(\log g_{\zeta\bar\zeta})_{j\bar k}|\le C_0\,u^{j\bar k}g_{j\bar k}\); and \(g_{\zeta\bar\zeta}/u_{\zeta\bar\zeta}=1/\eta(y)\le u^{j\bar k}g_{j\bar k}\) (the
smallest eigenvalue of \(A\) relative to \(A_\varphi\) is at most their sum). With \(K=L+C_0+C_1+1\) this gives \(\operatorname{tr}_{A_\varphi}A=u^{j\bar k}g_{j\bar k}\le KN\) at \(y\) (Lemma~\ref{bou:yau-c2-local-step} is the chart form of this step, with \(\rho\) the local
potential: it bounds \(\operatorname{tr}(u_{j\bar k})^{-1}\) at \(y\), and \(\operatorname{tr}_{A_\varphi}A\le C\operatorname{tr}(u_{j\bar k})^{-1}\) with \(C\) bounding \(g\) on the chart);
with \(\det(u_{j\bar k})\le e^{\sup|G|}\det(g_{j\bar k})\), Lemma~\ref{bou:yau-c2-matrix-endgame}~(2) bounds \(\eta(y)\le e^{\sup|G|}(KN)^{N-1}\). Hence
\(\log\eta\le\log\eta(y)+K\operatorname{osc}\varphi\) on \(Y\), i.e.\ \(A_\varphi\le CA\); and \(\operatorname{tr}_AA_\varphi\le N\eta\) with \(\det(u_{j\bar k})\ge e^{-\sup|G|}\det(g_{j\bar k})\)
gives \(A\le CA_\varphi\) (Lemma~\ref{bou:yau-c2-matrix-endgame}~(3)).
\end{proof}
This is B\l{}ocki's proof of the second-order estimate (Theorem~5.13 of his chapter, in the nondegenerate case \(f=e^G\)); it uses ordinary derivatives in
charts and a fixed direction \(\zeta\), so neither normal coordinates nor the curvature of \(A\) occur. In the nondegenerate case his term in
\(f^{1/(n-1)}\) and the arithmetic--geometric mean step are not needed: \(u_{\zeta\bar\zeta}\ge1/\operatorname{tr}(u_{j\bar k})^{-1}\) for \(|\zeta|=1\) suffices. Sz\'ekelyhidi's Lemmas~3.7--3.8 prove the same bound
with a lower bound for \(\Delta_AG\) in place of \(dd^cG\ge-LA\), through holomorphic normal coordinates.

\begin{lemma}[The trace identity of the third-order estimate]
\label{bou:yau-c3-trace-identity}
\uses{bou:log-det-derivatives}
\source{pss07-multiplier-krf:page-0008}
\lean{TranscendentalNakamaye.Pending.ThirdOrderTraceIdentity.complexHessian_trace,TranscendentalNakamaye.Pending.ThirdOrderTraceIdentity.re_trace_inv_mul_conjTranspose_mul_nonneg,TranscendentalNakamaye.Pending.ThirdOrderTraceIdentity.re_trace_mul_complexHessian_trace}
\leanok
Let \(u\) be a smooth real function on \(\mathbf C^n\) whose complex Hessian \(H=(u_{j\bar k})\) is positive definite at \(x\), \(P=H(x)^{-1}\) and \(X_i=\partial_iH(x)\).
Then \(\operatorname{Re}\operatorname{tr}\bigl(P\,\partial\bar\partial(\operatorname{tr}H)\bigr)(x)=\sum_i(\log\det H)_{i\bar\imath}(x)+\sum_i\operatorname{Re}\operatorname{tr}(PX_i^*PX_i)\), and the last
sum is \(\ge0\).
\end{lemma}
\begin{proof}
\leanok
Differentiate \(\log\det H\) twice (Jacobi's formula): \((\log\det H)_{i\bar\imath}=\operatorname{tr}(P\,\partial_i\partial_{\bar\imath}H)-\operatorname{tr}(PX_i^*PX_i)\) at \(x\), and
\(\partial_i\partial_{\bar\imath}H\) is the complex Hessian of \(u_{i\bar\imath}\); \(\operatorname{tr}(PX^*PX)=|P^{1/2}XP^{1/2}|^2\ge0\).
\end{proof}

\begin{lemma}[Calabi's identity, flat coordinates]
\label{bou:calabi-flat-identity}
\uses{bou:log-det-derivatives,bou:yau-c3-trace-identity}
\source{pss07-multiplier-krf:page-0008,pss07-multiplier-krf:page-0009}
\lean{TranscendentalNakamaye.Pending.CalabiFlatIdentity.Mat,TranscendentalNakamaye.Pending.CalabiFlatIdentity.calA,TranscendentalNakamaye.Pending.CalabiFlatIdentity.calB,TranscendentalNakamaye.Pending.CalabiFlatIdentity.calabiFlatSum,TranscendentalNakamaye.Pending.CalabiFlatIdentity.calabiFlatSum_eq,TranscendentalNakamaye.Pending.CalabiFlatIdentity.calabiForm,TranscendentalNakamaye.Pending.CalabiFlatIdentity.calabiForm_K_left,TranscendentalNakamaye.Pending.CalabiFlatIdentity.calabiForm_K_mid,TranscendentalNakamaye.Pending.CalabiFlatIdentity.calabiForm_K_right,TranscendentalNakamaye.Pending.CalabiFlatIdentity.calabiForm_SW,TranscendentalNakamaye.Pending.CalabiFlatIdentity.calabiForm_one,TranscendentalNakamaye.Pending.CalabiFlatIdentity.calabiForm_sub_left,TranscendentalNakamaye.Pending.CalabiFlatIdentity.calabiForm_sub_mid,TranscendentalNakamaye.Pending.CalabiFlatIdentity.calabiForm_sub_right,TranscendentalNakamaye.Pending.CalabiFlatIdentity.kappa4,TranscendentalNakamaye.Pending.CalabiFlatIdentity.kappa4_eq_cross,TranscendentalNakamaye.Pending.CalabiFlatIdentity.lterm,TranscendentalNakamaye.Pending.CalabiFlatIdentity.mDiag,TranscendentalNakamaye.Pending.CalabiFlatIdentity.ricTerm,TranscendentalNakamaye.Pending.CalabiFlatIdentity.star_calabiForm,TranscendentalNakamaye.Pending.CalabiFlatIdentity.sum_A4,TranscendentalNakamaye.Pending.CalabiFlatIdentity.sum_A5,TranscendentalNakamaye.Pending.CalabiFlatIdentity.sum_R1,TranscendentalNakamaye.Pending.CalabiFlatIdentity.sum_R2,TranscendentalNakamaye.Pending.CalabiFlatIdentity.sum_R3,TranscendentalNakamaye.Pending.CalabiFlatIdentity.sum_T12,TranscendentalNakamaye.Pending.CalabiFlatIdentity.sum_T14,TranscendentalNakamaye.Pending.CalabiFlatIdentity.sum_T24,TranscendentalNakamaye.Pending.CalabiFlatIdentity.sum_T31,TranscendentalNakamaye.Pending.CalabiFlatIdentity.sum_T34,TranscendentalNakamaye.Pending.CalabiFlatIdentity.sum_T35,TranscendentalNakamaye.Pending.CalabiFlatIdentity.sum_T43,TranscendentalNakamaye.Pending.CalabiFlatIdentity.sum_T53,TranscendentalNakamaye.Pending.CalabiFlatIdentity.sum_calabiForm_left,TranscendentalNakamaye.Pending.CalabiFlatIdentity.sum_calabiForm_mid,TranscendentalNakamaye.Pending.CalabiFlatIdentity.sum_calabiForm_right,TranscendentalNakamaye.Pending.CalabiFlatIdentity.sum_diag_eq_trace,TranscendentalNakamaye.Pending.CalabiFlatIdentity.sum_mDiag,TranscendentalNakamaye.Pending.CalabiFlatIdentity.sum_mul_conjTranspose_self,TranscendentalNakamaye.Pending.CalabiFlatIdentity.sw2_0,TranscendentalNakamaye.Pending.CalabiFlatIdentity.sw3_0,TranscendentalNakamaye.Pending.CalabiFlatIdentity.sw3_1,TranscendentalNakamaye.Pending.CalabiFlatIdentity.sw4_0,TranscendentalNakamaye.Pending.CalabiFlatIdentity.sw4_1,TranscendentalNakamaye.Pending.CalabiFlatIdentity.sw4_2,TranscendentalNakamaye.Pending.CalabiFlatIdentity.sw5_0,TranscendentalNakamaye.Pending.CalabiFlatIdentity.sw5_1,TranscendentalNakamaye.Pending.CalabiFlatIdentity.sw5_2,TranscendentalNakamaye.Pending.CalabiFlatIdentity.sw5_3,TranscendentalNakamaye.Pending.CalabiFlatIdentity.sw6_0,TranscendentalNakamaye.Pending.CalabiFlatIdentity.sw6_1,TranscendentalNakamaye.Pending.CalabiFlatIdentity.sw6_2,TranscendentalNakamaye.Pending.CalabiFlatIdentity.sw6_3,TranscendentalNakamaye.Pending.CalabiFlatIdentity.sw6_4,TranscendentalNakamaye.Pending.CalabiFlatIdentity.sw7_0,TranscendentalNakamaye.Pending.CalabiFlatIdentity.sw7_1,TranscendentalNakamaye.Pending.CalabiFlatIdentity.sw7_2,TranscendentalNakamaye.Pending.CalabiFlatIdentity.sw7_3,TranscendentalNakamaye.Pending.CalabiFlatIdentity.sw7_4,TranscendentalNakamaye.Pending.CalabiFlatIdentity.sw7_5,TranscendentalNakamaye.Pending.CalabiFlatIdentity.sw8_0,TranscendentalNakamaye.Pending.CalabiFlatIdentity.sw8_1,TranscendentalNakamaye.Pending.CalabiFlatIdentity.sw8_2,TranscendentalNakamaye.Pending.CalabiFlatIdentity.sw8_3,TranscendentalNakamaye.Pending.CalabiFlatIdentity.sw8_4,TranscendentalNakamaye.Pending.CalabiFlatIdentity.sw8_5,TranscendentalNakamaye.Pending.CalabiFlatIdentity.sw8_6,TranscendentalNakamaye.Pending.CalabiFlatIdentity.trace_mul_conjTranspose,TranscendentalNakamaye.Pending.CalabiFlatIdentity.trace_sq_sub,TranscendentalNakamaye.Pending.CalabiFlatIdentity.trace_sq_sub_sub,TranscendentalNakamaye.Pending.CalabiFlatLaplacian.MDiffAt,TranscendentalNakamaye.Pending.CalabiFlatLaplacian.MDiffAt.add,TranscendentalNakamaye.Pending.CalabiFlatLaplacian.MDiffAt.conjTranspose,TranscendentalNakamaye.Pending.CalabiFlatLaplacian.MDiffAt.differentiableAt_trace,TranscendentalNakamaye.Pending.CalabiFlatLaplacian.MDiffAt.inv,TranscendentalNakamaye.Pending.CalabiFlatLaplacian.MDiffAt.mul,TranscendentalNakamaye.Pending.CalabiFlatLaplacian.MDiffAt.neg,TranscendentalNakamaye.Pending.CalabiFlatLaplacian.MSmooth,TranscendentalNakamaye.Pending.CalabiFlatLaplacian.MSmooth.conjTranspose,TranscendentalNakamaye.Pending.CalabiFlatLaplacian.MSmooth.mdiffAt,TranscendentalNakamaye.Pending.CalabiFlatLaplacian.MSmooth.wirtM,TranscendentalNakamaye.Pending.CalabiFlatLaplacian.aOne,TranscendentalNakamaye.Pending.CalabiFlatLaplacian.aOne_eq,TranscendentalNakamaye.Pending.CalabiFlatLaplacian.calabiFlatSum_eq_sum,TranscendentalNakamaye.Pending.CalabiFlatLaplacian.calabiS,TranscendentalNakamaye.Pending.CalabiFlatLaplacian.complexHessian_entry,TranscendentalNakamaye.Pending.CalabiFlatLaplacian.contDiff_wirt_I_ofReal,TranscendentalNakamaye.Pending.CalabiFlatLaplacian.dH,TranscendentalNakamaye.Pending.CalabiFlatLaplacian.dH_entry,TranscendentalNakamaye.Pending.CalabiFlatLaplacian.dH_symm,TranscendentalNakamaye.Pending.CalabiFlatLaplacian.dV,TranscendentalNakamaye.Pending.CalabiFlatLaplacian.dV_symm,TranscendentalNakamaye.Pending.CalabiFlatLaplacian.dV_symm',TranscendentalNakamaye.Pending.CalabiFlatLaplacian.dW,TranscendentalNakamaye.Pending.CalabiFlatLaplacian.dZ,TranscendentalNakamaye.Pending.CalabiFlatLaplacian.dZ_conjTranspose,TranscendentalNakamaye.Pending.CalabiFlatLaplacian.dZ_entry,TranscendentalNakamaye.Pending.CalabiFlatLaplacian.dZ_symm,TranscendentalNakamaye.Pending.CalabiFlatLaplacian.dZ_symm',TranscendentalNakamaye.Pending.CalabiFlatLaplacian.det_ne_zero_of_eq_one,TranscendentalNakamaye.Pending.CalabiFlatLaplacian.differentiableAt_calabiForm,TranscendentalNakamaye.Pending.CalabiFlatLaplacian.flatSummand,TranscendentalNakamaye.Pending.CalabiFlatLaplacian.invH,TranscendentalNakamaye.Pending.CalabiFlatLaplacian.invH_eq_one,TranscendentalNakamaye.Pending.CalabiFlatLaplacian.mSmooth_complexHessian,TranscendentalNakamaye.Pending.CalabiFlatLaplacian.mSmooth_dH,TranscendentalNakamaye.Pending.CalabiFlatLaplacian.mSmooth_dV,TranscendentalNakamaye.Pending.CalabiFlatLaplacian.mSmooth_dZ,TranscendentalNakamaye.Pending.CalabiFlatLaplacian.mdiffAt_aOne,TranscendentalNakamaye.Pending.CalabiFlatLaplacian.mdiffAt_invH,TranscendentalNakamaye.Pending.CalabiFlatLaplacian.stageOne,TranscendentalNakamaye.Pending.CalabiFlatLaplacian.sum_calabiForm_C,TranscendentalNakamaye.Pending.CalabiFlatLaplacian.sum_calabiForm_Y,TranscendentalNakamaye.Pending.CalabiFlatLaplacian.sum_wirt_wirt_calabiS,TranscendentalNakamaye.Pending.CalabiFlatLaplacian.wirtM_I_conjTranspose,TranscendentalNakamaye.Pending.CalabiFlatLaplacian.wirtM_I_dV,TranscendentalNakamaye.Pending.CalabiFlatLaplacian.wirtM_I_invH,TranscendentalNakamaye.Pending.CalabiFlatLaplacian.wirtM_congr,TranscendentalNakamaye.Pending.CalabiFlatLaplacian.wirtM_inv,TranscendentalNakamaye.Pending.CalabiFlatLaplacian.wirtM_invH,TranscendentalNakamaye.Pending.CalabiFlatLaplacian.wirtM_mul,TranscendentalNakamaye.Pending.CalabiFlatLaplacian.wirtM_neg,TranscendentalNakamaye.Pending.CalabiFlatLaplacian.wirtM_negI_aOne,TranscendentalNakamaye.Pending.CalabiFlatLaplacian.wirtM_negI_conjTranspose,TranscendentalNakamaye.Pending.CalabiFlatLaplacian.wirtM_negI_dH_conjTranspose,TranscendentalNakamaye.Pending.CalabiFlatLaplacian.wirtM_negI_invH,TranscendentalNakamaye.Pending.CalabiFlatLaplacian.wirt_add5,TranscendentalNakamaye.Pending.CalabiFlatLaplacian.wirt_calabiForm,TranscendentalNakamaye.Pending.CalabiFlatLaplacian.wirt_calabiS,TranscendentalNakamaye.Pending.CalabiFlatLaplacian.wirt_negI_conj,TranscendentalNakamaye.Pending.CalabiFlatLaplacian.wirt_stageOne,TranscendentalNakamaye.Pending.CalabiFlatLaplacian.wirt_trace,TranscendentalNakamaye.Pending.CalabiFlatEquation.calabi_identity,TranscendentalNakamaye.Pending.CalabiFlatEquation.complexHessian_at,TranscendentalNakamaye.Pending.CalabiFlatEquation.eventually_complexHessian_eq,TranscendentalNakamaye.Pending.CalabiFlatEquation.logDetHess,TranscendentalNakamaye.Pending.CalabiFlatEquation.trace_dW,TranscendentalNakamaye.Pending.CalabiFlatEquation.wirtM_complexHessian_at,TranscendentalNakamaye.Pending.CalabiFlatEquation.wirt_I_log_det,TranscendentalNakamaye.Pending.CalabiFlatEquation.wirt_logDetHess,TranscendentalNakamaye.Pending.CalabiFlatEquation.wirt_sub,TranscendentalNakamaye.Pending.CalabiFlatEquation.wirt_wirt_log_det}
\leanok
Let \(u\) be smooth on \(\mathbf C^n\), \(H=(u_{j\bar k})\), \(P=H^{-1}\), \(X_m=\partial_mH\), and let \(S=\sum_{m,t}P_{tm}\operatorname{tr}(PX_mPX_t^*)\) be Calabi's
quantity. Suppose \(\log\det H=F\) near a point \(x\) with \(H(x)=1\), \(F\) smooth. Then at \(x\),
\(\sum_p\partial_p\partial_{\bar p}S=|\nabla\Gamma|^2+|\bar\nabla\Gamma|^2+2\operatorname{Re}\sum_m\operatorname{tr}(F_{\cdot\bar\cdot m}X_m^*)-\operatorname{Ric}(F_{\cdot\bar\cdot})\), where
\(\nabla_p\Gamma_m=\partial_pX_m-X_pX_m-\sum_l(X_p)_{ml}X_l\), \(\bar\nabla_q\Gamma_m=\partial_{\bar q}X_m-X_mX_q^*\), \(F_{\cdot\bar\cdot}=(F_{a\bar b})\),
\(F_{\cdot\bar\cdot m}=\partial_m(F_{a\bar b})\), and \(\operatorname{Ric}(F_{\cdot\bar\cdot})\) is linear in \(F_{\cdot\bar\cdot}\) and quadratic in the \(X_m\).
\end{lemma}
\begin{proof}
\leanok
At \(x\), by the Leibniz rule (with \(\partial P=-PXP\) and \(P(x)=1\)), \(\sum_p\partial_p\partial_{\bar p}S\) is a sum of 25 terms in the derivatives of \(H\) of orders one
to three at \(x\). The twice and thrice differentiated equation \(\log\det H=F\) and the symmetries of the derivatives of \(u\) reduce this sum, by reindexing finite
sums, to the stated expression (the flat-reference form of Phong--\v{S}e\v{s}um--Sturm (2.43)/(2.44)).
\end{proof}

\begin{lemma}[Calabi's inequality at a point]
\label{bou:calabi-flat-inequality}
\uses{bou:calabi-flat-identity}
\source{pss07-multiplier-krf:page-0009}
\lean{TranscendentalNakamaye.Pending.CalabiFlatInequality.calabiS_at_one,TranscendentalNakamaye.Pending.CalabiFlatInequality.calabi_inequality_at_one,TranscendentalNakamaye.Pending.CalabiFlatInequality.complexHessian_add_smul,TranscendentalNakamaye.Pending.CalabiFlatInequality.frob,TranscendentalNakamaye.Pending.CalabiFlatInequality.frob_conjTranspose,TranscendentalNakamaye.Pending.CalabiFlatInequality.frob_mul_le,TranscendentalNakamaye.Pending.CalabiFlatInequality.frob_nonneg,TranscendentalNakamaye.Pending.CalabiFlatInequality.norm_entry_le,TranscendentalNakamaye.Pending.CalabiFlatInequality.norm_ricTerm_le,TranscendentalNakamaye.Pending.CalabiFlatInequality.norm_sum_mul_sq_le',TranscendentalNakamaye.Pending.CalabiFlatInequality.norm_sum_trace_mul_le,TranscendentalNakamaye.Pending.CalabiFlatInequality.norm_trace_mul_conjTranspose_le,TranscendentalNakamaye.Pending.CalabiFlatInequality.norm_trace_mul_conjTranspose_sq_le,TranscendentalNakamaye.Pending.CalabiFlatInequality.norm_trace_mul_le,TranscendentalNakamaye.Pending.CalabiFlatInequality.sqrt_frob_conjTranspose_mul_le,TranscendentalNakamaye.Pending.CalabiFlatInequality.sqrt_frob_mul_conjTranspose_le,TranscendentalNakamaye.Pending.CalabiFlatInequality.trace_mul_conjTranspose_self,TranscendentalNakamaye.Pending.CalabiFlatInequality.wirtM_complexHessian_add_smul,TranscendentalNakamaye.Pending.CalabiFlatInequality.wirt_negI_calabiS}
\leanok
Let \(u\) be smooth on \(\mathbf C^n\) with \(\log\det(u_{j\bar k})=f+\varepsilon u\) near a point \(x\) with \((u_{j\bar k})(x)=1\), \(f\) smooth, \(\varepsilon\ge0\). There is
\(G\ge0\) with \(\sum_p|\partial_pS|^2\le2SG\) and
\(\sum_p\partial_p\partial_{\bar p}S\ge G-2|\partial\partial\bar\partial f|\sqrt S-(n+2)|\partial\bar\partial f+\varepsilon|\,S\) at \(x\) (Frobenius norms of the derivatives of
\(f\) at \(x\)).
\end{lemma}
\begin{proof}
\leanok
Take \(G=|\nabla\Gamma|^2+|\bar\nabla\Gamma|^2\). The gradient identity \(\partial_pS=\sum_m\langle\nabla_p\Gamma_m,X_m\rangle+\sum_m\langle X_m,\bar\nabla_p\Gamma_m\rangle\) and
Cauchy--Schwarz give the first bound; Lemma~\ref{bou:calabi-flat-identity} for \(F=f+\varepsilon u\) gives the second (the \(\varepsilon\)-part of the third
derivatives contributes \(2\varepsilon S\ge0\)).
\end{proof}

\begin{lemma}[A complex-linear change of coordinates]
\label{bou:calabi-linear-change}
\uses{bou:calabi-flat-identity}
\lean{TranscendentalNakamaye.Pending.CalabiLinearChange.frob_transpose,TranscendentalNakamaye.Pending.CalabiLinearChange.frob_map_conj,TranscendentalNakamaye.Pending.CalabiLinearChange.frob_sandwich_le,TranscendentalNakamaye.Pending.CalabiLinearChange.frob_sum_smul_le,TranscendentalNakamaye.Pending.CalabiLinearChange.transform,TranscendentalNakamaye.Pending.CalabiLinearChange.sum_frob_transform_le,TranscendentalNakamaye.Pending.CalabiLinearChange.transform_inv_transform,TranscendentalNakamaye.Pending.CalabiLinearChange.sum4_swap,TranscendentalNakamaye.Pending.CalabiLinearChange.conjTranspose_transpose_eq,TranscendentalNakamaye.Pending.CalabiLinearChange.det_map_conj,TranscendentalNakamaye.Pending.CalabiLinearChange.calabiForm_transform,TranscendentalNakamaye.Pending.CalabiLinearChange.det_sandwich,TranscendentalNakamaye.Pending.CalabiLinearChange.trace_eq_trace_mul_sandwich,TranscendentalNakamaye.Pending.CalabiLinearChange.frob_inv_map_conj,TranscendentalNakamaye.Pending.CalabiLinearChange.frob_le_mul_re_trace,TranscendentalNakamaye.Pending.CalabiLinearChange.norm_trace_inv_mul_le,TranscendentalNakamaye.Pending.CalabiLinearChange.exists_normalizing,TranscendentalNakamaye.Pending.CalabiLinearChange.lmat,TranscendentalNakamaye.Pending.CalabiLinearChange.fderiv_clm,TranscendentalNakamaye.Pending.CalabiLinearChange.contDiff_comp_clm,TranscendentalNakamaye.Pending.CalabiLinearChange.wirt_comp_clm,TranscendentalNakamaye.Pending.CalabiLinearChange.wirtM_comp_clm,TranscendentalNakamaye.Pending.CalabiLinearChange.wirtM_negI_clm_apply,TranscendentalNakamaye.Pending.CalabiLinearChange.wirtM_const,TranscendentalNakamaye.Pending.CalabiLinearChange.wirtM_const_sandwich,TranscendentalNakamaye.Pending.CalabiLinearChange.complexHessian_comp_clm,TranscendentalNakamaye.Pending.CalabiLinearChange.mdiffAt_comp_clm,TranscendentalNakamaye.Pending.CalabiLinearChange.dH_comp_clm,TranscendentalNakamaye.Pending.CalabiLinearChange.calabiS_comp_clm,TranscendentalNakamaye.Pending.CalabiLinearChange.trace_complexHessian_comp_clm,TranscendentalNakamaye.Pending.CalabiLinearChange.det_complexHessian_comp_clm,TranscendentalNakamaye.Pending.CalabiLinearChange.lmat_toEuclideanCLM,TranscendentalNakamaye.Pending.CalabiLinearChange.toEuclideanCLM_apply_inv,TranscendentalNakamaye.Pending.CalabiLinearChange.exists_normalizing_clm}
\leanok
Let \(L\colon\mathbf C^n\to\mathbf C^n\) be complex linear with invertible matrix, also denoted \(L\). For \(g\) smooth, the complex Hessian transforms by
\(H(g\circ L)(y)=L^TH(g)(Ly)\bar L\), its first derivatives by \(\partial_mH(g\circ L)(y)=L^T\bigl(\sum_rL_{rm}\partial_rH(g)(Ly)\bigr)\bar L\), Calabi's quantity
is invariant, \(S(u\circ L)=S(u)\circ L\), and \(\det H(u\circ L)=|\det L|^2\det H(u)\circ L\). If \(H(u)(x)>0\) there are \(L\) and \(y\) with \(Ly=x\) and
\(H(u\circ L)(y)=1\), \(|L|^2=\operatorname{tr}H(u)(x)^{-1}\), \(|L^{-1}|^2=\operatorname{tr}H(u)(x)\) (Frobenius norms), together with the corresponding bounds for
the transformed third derivatives.
\end{lemma}
\begin{proof}
\leanok
The Wirtinger derivatives satisfy \(\partial_w(f\circ L)=(\partial_{Lw}f)\circ L\), which gives the transformation rules; \(S\) is built from \(H^{-1}\) and the
\(\partial_mH\) by contractions invariant under \(X\mapsto L^TX\bar L\). For the normalisation, diagonalise \(H(u)(x)\) simultaneously with itself: there
is \(R\) with \(R^*H(u)(x)R=1\), so \(H(u)(x)^{-1}=RR^*\) and \(H(u)(x)=R^{*-1}R^{-1}\); take for \(L\) the map with matrix \(\bar R\), so that
\(L^TH(u)(x)\bar L=R^*H(u)(x)R=1\).
\end{proof}

\begin{lemma}[Calabi's estimate on a ball]
\label{bou:calabi-cutoff-estimate}
\uses{bou:calabi-linear-change,bou:calabi-flat-inequality,bou:max-principle}
\source{pss07-multiplier-krf:page-0008,pss07-multiplier-krf:page-0009}
\lean{TranscendentalNakamaye.Pending.CalabiCutoffEstimate.MC2At,TranscendentalNakamaye.Pending.CalabiCutoffEstimate.MC2At.mul,TranscendentalNakamaye.Pending.CalabiCutoffEstimate.MC2At.conjTranspose,TranscendentalNakamaye.Pending.CalabiCutoffEstimate.contDiffAt_det,TranscendentalNakamaye.Pending.CalabiCutoffEstimate.MC2At.inv,TranscendentalNakamaye.Pending.CalabiCutoffEstimate.MC2At.trace,TranscendentalNakamaye.Pending.CalabiCutoffEstimate.contDiffAt_calabiForm,TranscendentalNakamaye.Pending.CalabiCutoffEstimate.mC2At_of_mSmooth,TranscendentalNakamaye.Pending.CalabiCutoffEstimate.contDiffAt_calabiS,TranscendentalNakamaye.Pending.CalabiCutoffEstimate.star_calabiS,TranscendentalNakamaye.Pending.CalabiCutoffEstimate.ofReal_re_calabiS,TranscendentalNakamaye.Pending.CalabiCutoffEstimate.ofReal_re_trace_complexHessian,TranscendentalNakamaye.Pending.CalabiCutoffEstimate.differentiableAt_wirt,TranscendentalNakamaye.Pending.CalabiCutoffEstimate.eventually_differentiableAt,TranscendentalNakamaye.Pending.CalabiCutoffEstimate.wirt_wirt_add,TranscendentalNakamaye.Pending.CalabiCutoffEstimate.wirt_wirt_const_mul,TranscendentalNakamaye.Pending.CalabiCutoffEstimate.wirt_wirt_mul,TranscendentalNakamaye.Pending.CalabiCutoffEstimate.wirt_I_ofReal,TranscendentalNakamaye.Pending.CalabiCutoffEstimate.frob_add_le,TranscendentalNakamaye.Pending.CalabiCutoffEstimate.frob_smul_one,TranscendentalNakamaye.Pending.CalabiCutoffEstimate.complexHessian_add_const,TranscendentalNakamaye.Pending.CalabiCutoffEstimate.norm_trace_sandwich_le,TranscendentalNakamaye.Pending.CalabiCutoffEstimate.sum_norm_wirt_comp_clm_sq_le,TranscendentalNakamaye.Pending.CalabiCutoffEstimate.re_trace_sandwich_inv_nonneg,TranscendentalNakamaye.Pending.CalabiCutoffEstimate.neg_le_re_sum_conj_mul,TranscendentalNakamaye.Pending.CalabiCutoffEstimate.sum_wirt_wirt_trace_mul,TranscendentalNakamaye.Pending.CalabiCutoffEstimate.sum_dZ_diag,TranscendentalNakamaye.Pending.CalabiCutoffEstimate.real_endgame,TranscendentalNakamaye.Pending.CalabiCutoffEstimate.calabiS_le_of_isLocalMax,TranscendentalNakamaye.Pending.CalabiCutoffEstimate.sum_frob_dH_le,TranscendentalNakamaye.Pending.CalabiCutoffEstimate.exists_norm_le_of_hasCompactSupport,TranscendentalNakamaye.Pending.CalabiCutoffEstimate.exists_sum_frob_dH_le}
\leanok
Let \(c\in\mathbf C^n\), \(r>0\) and \(\Lambda,K_0\in\mathbf R\). There is \(C\) such that for smooth \(u,f\) on \(\mathbf C^n\) and \(0\le\varepsilon\le1\), if on the ball
\(B(c,2r)\) the complex Hessian \(H=H(u)\) is positive definite with \(\operatorname{tr}H\le\Lambda\), \(\operatorname{tr}H^{-1}\le\Lambda\), \(\log\det H=f+\varepsilon u\),
\(|\partial\bar\partial f|^2\le K_0\) and \(\sum_m|\partial_m\partial\bar\partial f|^2\le K_0\) (Frobenius norms), then \(\sum_m|\partial_mH|^2\le C\) on
\(\bar B(c,r)\).
\end{lemma}
\begin{proof}
\leanok
Let \(\eta\) be a bump equal to \(1\) on \(\bar B(c,r)\) and \(0\) off \(B(c,3r/2)\), \(S\) Calabi's quantity of \(u\), and \(Q=\eta^2S+K\operatorname{tr}H\) for a
large constant \(K\); \(Q\) attains its maximum on \(\bar B(c,3r/2)\) at some \(x_0\). If \(\eta(x_0)\ne0\), \(x_0\) is an interior local maximum; normalise
\(H(x_0)=1\) by a complex-linear change of coordinates \(L\) (Lemma~\ref{bou:calabi-linear-change}: \(S\) is invariant, \(\operatorname{tr}H\) becomes
\(\operatorname{tr}(MH)\) with the constant matrix \(M=(L^T\bar L)^{-1}\)). There \(\sum_p\partial_p\partial_{\bar p}(Q\circ L)\le0\) (Lemma~\ref{bou:max-principle});
Calabi's inequality (Lemma~\ref{bou:calabi-flat-inequality}) bounds \(\Delta S\) from below and \(|\partial S|^2\le2SG\), the cross term is absorbed by
\(\eta^2G\), and the twice differentiated equation gives \(\operatorname{tr}(M(\partial\bar\partial F+\sum_pX_pX_p^*))\ge S/\Lambda-\Lambda|\partial\bar\partial F|\).
For \(K\) large this bounds \(S(x_0)\). Hence \(S\le Q(x_0)\le C'\) on \(\bar B(c,r)\), and \(\sum_m|\partial_mH|^2\le(\operatorname{tr}H)^3S\).
\end{proof}

\begin{lemma}[The third-order estimate]
\label{bou:yau-c3-estimate}
\uses{bou:calabi-cutoff-estimate,bou:calabi-linear-change,fpglobal:theta-psh}
\source{szekelyhidi-extremal:page-0057,szekelyhidi-extremal:page-0058,szekelyhidi-extremal:page-0061,szekelyhidi-extremal:page-0062,szekelyhidi-extremal:page-0063}
\lean{TranscendentalNakamaye.Pending.CalabiThirdOrderEstimate.norm_fderiv_complexHessian_le,TranscendentalNakamaye.Pending.CalabiThirdOrderEstimate.frob_smul_le,TranscendentalNakamaye.Pending.CalabiThirdOrderEstimate.frob_comb_le,TranscendentalNakamaye.Pending.CalabiThirdOrderEstimate.complexHessian_comb,TranscendentalNakamaye.Pending.CalabiThirdOrderEstimate.dH_comb,TranscendentalNakamaye.Pending.CalabiThirdOrderEstimate.exists_frob_bound,TranscendentalNakamaye.Pending.CalabiThirdOrderEstimate.re_star_dotProduct_self,TranscendentalNakamaye.Pending.CalabiThirdOrderEstimate.exists_chart_piece,TranscendentalNakamaye.Pending.CalabiThirdOrderEstimate.ball_hyps,TranscendentalNakamaye.Pending.CalabiThirdOrderEstimate.exists_norm_fderiv_complexHessian_le}
\leanok
Let \((Y,A)\) be a K\"ahler manifold of dimension \(N\), \(F\) smooth, \(C_0\in\mathbf R\), \(e\) a chart and \(K\subset e.\mathrm{target}\) compact. There is \(C\) such
that for every smooth \(\varphi\) with \(A_\varphi=A+dd^c\varphi\) K\"ahler, every \(t,\varepsilon\in[0,1]\) and \(\kappa\in\mathbf R\), if
\(A_\varphi^N=e^{tF+\varepsilon\varphi+\kappa}A^N\) and \(C_0^{-1}A\le A_\varphi\le C_0A\) in the chart \(e\), then the real derivative of every entry of the complex
Hessian of \(\varphi\circ e^{-1}\) is bounded by \(C\) on \(K\).
\end{lemma}
\begin{proof}
\leanok
Cover \(K\) by finitely many balls \(B(p,2R)\) on which \(A=dd^c\rho\) with \(\rho\) smooth (local potentials, Definition~\ref{fpglobal:theta-psh}). On such a ball
\(u=\rho+\varphi\circ e^{-1}\) satisfies \(\log\det H(u)=f+\varepsilon u\) with \(f=tF\circ e^{-1}-\varepsilon\rho+\kappa+\log\det H(\rho)\); the two-sided bound
gives \(\operatorname{tr}H(u)\le\Lambda\) and \(\operatorname{tr}H(u)^{-1}\le\Lambda\), and the derivatives of \(\partial\bar\partial f\) up to order one are bounded
independently of \(\varphi,t,\varepsilon,\kappa\) (\(\kappa\) disappears). Lemma~\ref{bou:calabi-cutoff-estimate} bounds \(\partial_mH(u)\) on \(\bar B(p,R)\), and
\(\partial_mH(\varphi\circ e^{-1})=\partial_mH(u)-\partial_mH(\rho)\).
\end{proof}
This is Sz\'ekelyhidi's third-order estimate (his Lemmas~3.7--3.10, after Phong--\v{S}e\v{s}um--Sturm) in a local form: an interior estimate by a cut-off in
flat coordinates of a chart instead of the global maximum principle with the reference metric, for the whole family
\(A_\varphi^N=e^{tF+\varepsilon\varphi+\kappa}A^N\) (\(\kappa\) arbitrary, as needed for the normalisation in Lemma~\ref{bou:yau-eps-limit}).

\begin{lemma}[The \(C^0\) estimate at fixed \(\varepsilon\)]
\label{bou:yau-sup-bound}
\uses{bou:complex-hessian-trace-max,bou:max-principle}
\lean{TranscendentalNakamaye.Pending.MongeAmpereSupBound.abs_le_of_mongeAmpere_eq_exp,TranscendentalNakamaye.Pending.MongeAmpereSupBound.exists_potential_hessian,TranscendentalNakamaye.Pending.MongeAmpereSupBound.isUnit_det_of_posDef,TranscendentalNakamaye.Pending.MongeAmpereSupBound.le_zero_of_isMax,TranscendentalNakamaye.Pending.MongeAmpereSupBound.one_le_prod_of_sum_inv_le,TranscendentalNakamaye.Pending.MongeAmpereSupBound.prod_le_one_of_sum_le,TranscendentalNakamaye.Pending.MongeAmpereSupBound.re_det_le_of_trace_inv_le,TranscendentalNakamaye.Pending.MongeAmpereSupBound.re_det_le_of_trace_le,TranscendentalNakamaye.Pending.MongeAmpereSupBound.zero_le_of_isMin}
\leanok
Let \(Y\) be a compact complex manifold of dimension \(N\) with a K\"ahler form \(A\), \(\varepsilon>0\), \(B\in\mathbf R\) and \(H\colon Y\to\mathbf R\) with
\(|H|\le B\). If \(\varphi\) is smooth, \(A_\varphi=A+dd^c\varphi\) is K\"ahler and \(A_\varphi^N=e^{H+\varepsilon\varphi}A^N\), then \(\varepsilon|\varphi|\le B\) on \(Y\).
\end{lemma}
\begin{proof}
\leanok
Let \(y_1\) be a maximum point of \(\varphi\). In a chart at \(y_1\) the complex Hessian of \(\varphi\circ e^{-1}\) is \(\le0\) (Lemma~\ref{bou:complex-hessian-trace-max}),
so with \(P=H_A\) positive definite and \(G=H_A+\partial\bar\partial(\varphi\circ e^{-1})\) positive definite, \(\operatorname{Re}\operatorname{tr}(P^{-1}G)\le N\), and the
arithmetic--geometric mean inequality for the eigenvalues of \(P^{-1}G\) gives \(\det G\le\det P\); hence \(e^{H(y_1)+\varepsilon\varphi(y_1)}\le1\). At a minimum point
\(y_2\) the reverse inequalities give \(H(y_2)+\varepsilon\varphi(y_2)\ge0\). So \(-B\le-H(y_2)\le\varepsilon\varphi(y_2)\le\varepsilon\varphi\le\varepsilon\varphi(y_1)\le-H(y_1)\le B\).
\end{proof}

\begin{lemma}[The uniform \(C^{2,r}\) bound along the path]
\label{bou:yau-holder-bound}
\uses{bou:max-principle,bou:yau-c2-estimate,bou:yau-c3-estimate,bou:schauder-global-zero,bou:holder-spaces-manifold,bou:kahler-laplacian-realizable,gnpp:kahler-domination}
\lean{TranscendentalNakamaye.Pending.YauHolderBound.abs_sub_le_of_norm_fderiv_le,TranscendentalNakamaye.Pending.YauHolderBound.bounds_of_eventuallyEq,TranscendentalNakamaye.Pending.YauHolderBound.chart_bound,TranscendentalNakamaye.Pending.YauHolderBound.coeff_line_le,TranscendentalNakamaye.Pending.YauHolderBound.coeff_line_nonneg,TranscendentalNakamaye.Pending.YauHolderBound.contDiffOn_inv_apply,TranscendentalNakamaye.Pending.YauHolderBound.differentiableAt_complexHessian_apply,TranscendentalNakamaye.Pending.YauHolderBound.exists_bound_of_contDiffOn,TranscendentalNakamaye.Pending.YauHolderBound.exists_norm_ofChartSmooth_le,TranscendentalNakamaye.Pending.YauHolderBound.hasFDerivAt_re_trace_mul,TranscendentalNakamaye.Pending.YauHolderBound.invEntryCLM,TranscendentalNakamaye.Pending.YauHolderBound.invEntryCLM_lapCoeff,TranscendentalNakamaye.Pending.YauHolderBound.inv_apply_eq_lapCoeff,TranscendentalNakamaye.Pending.YauHolderBound.laplacian_symm_eq,TranscendentalNakamaye.Pending.YauHolderBound.leviForm_lincomb,TranscendentalNakamaye.Pending.YauHolderBound.norm_complexHessian_entry_le_of_leviForm,TranscendentalNakamaye.Pending.YauHolderBound.norm_complexHessian_entry_le_of_two_sided,TranscendentalNakamaye.Pending.YauHolderBound.norm_le_of_lipschitz,TranscendentalNakamaye.Pending.YauHolderBound.norm_re_trace_mul_le,TranscendentalNakamaye.Pending.YauHolderBound.norm_sum_reCLM_comp_le,TranscendentalNakamaye.Pending.YauHolderBound.re_trace_mul_eq}
\leanok
Let \(Y\) be a compact connected complex manifold of dimension \(N\) with a K\"ahler form \(A\) and a H\"older atlas, \(0<r<1\), \(F\) smooth and
\(S,K\in\mathbf R\). There is \(C\) such that every smooth \(\varphi\) with \(A_\varphi\) K\"ahler, \(A_\varphi^N=e^{tF+\varepsilon\varphi+\kappa}A^N\) for some
\(t,\varepsilon\in[0,1]\) and \(|\kappa|\le K\), and \(|\varphi|\le S\) satisfies \(\|\varphi\|_{C^{2,r}(Y)}\le C\).
\end{lemma}
\begin{proof}
\leanok
\(G=tF+\varepsilon\varphi+\kappa\) satisfies \(|G|\le\sup|F|+S+K\) and \(dd^cG=t\,dd^cF+\varepsilon(A_\varphi-A)\ge-(L_F+1)A\), where \(dd^cF\ge-L_FA\) and
\(\varepsilon\le1\) (such an \(L_F\) exists by Lemma~\ref{gnpp:kahler-domination}); Lemma~\ref{bou:yau-c2-estimate} gives \(C_0^{-1}A\le A_\varphi\le C_0A\),
which bounds the entries of the complex Hessian of \(\varphi\circ e^{-1}\) on the charts (polarisation of the Levi form), and Lemma~\ref{bou:yau-c3-estimate}
bounds their derivatives on compact pieces of the charts of the H\"older atlas. In every chart \(\Delta_A\varphi=\operatorname{Re}\operatorname{tr}(H_A^{-1}\,
\partial\bar\partial(\varphi\circ e^{-1}))\), so the local representatives of \(\chi_i\Delta_A\varphi\) are bounded and uniformly Lipschitz, hence
\(\Delta_A\varphi\) is bounded in \(C^{0,r}(Y)\); and the Schauder estimate
(Lemma~\ref{bou:schauder-global-zero}, for the operator \(\Delta_A\) of Lemma~\ref{bou:kahler-laplacian-realizable}) together with \(|\varphi|\le S\) bounds \(\varphi\) in
\(C^{2,r}(Y)\).
\end{proof}

\begin{lemma}[Limits of bounded families of solutions]
\label{bou:yau-limit-solutions}
\uses{bou:holder-spaces-manifold,bou:monge-ampere-holder,bou:yau-regularity}
\lean{TranscendentalNakamaye.Pending.YauClosedness.det_eq_of_density,TranscendentalNakamaye.Pending.YauClosedness.exists_smooth_of_tendsto,TranscendentalNakamaye.Pending.YauClosedness.limit_of_tendsto,TranscendentalNakamaye.Pending.YauClosedness.ma_eq_of_det_eq,TranscendentalNakamaye.Pending.YauClosedness.posDef_of_isKahler,TranscendentalNakamaye.Pending.YauClosedness.posDef_of_tendsto,TranscendentalNakamaye.Pending.YauClosedness.tendsto_apply,TranscendentalNakamaye.Pending.YauClosedness.tendsto_complexHessian}
\leanok
Let \(Y\) be a complex manifold of dimension \(N\) with a K\"ahler form \(A\) and a H\"older atlas, \(0<r<1\) and \(F\) smooth. Let \(\varphi_j\) be smooth with
\(A_{\varphi_j}\) K\"ahler and \(A_{\varphi_j}^N=e^{t_jF+\varepsilon_j\varphi_j+\kappa_j}A^N\), where \(t_j\to t_0\), \(\varepsilon_j\to\varepsilon_0\), \(\kappa_j\to\kappa_0\), and
\(\|\varphi_j\|_{C^{2,r}(Y)}\le C\). Then there is a smooth \(\varphi_0\) with \(A_{\varphi_0}\) K\"ahler and \(A_{\varphi_0}^N=e^{t_0F+\varepsilon_0\varphi_0+\kappa_0}A^N\).
\end{lemma}
\begin{proof}
\leanok
By the compactness of \(C^{2,r}(Y)\subset C^{2,s}(Y)\), \(s=r/2\) (Lemma~\ref{bou:holder-spaces-manifold}), a subsequence converges in \(C^{2,s}(Y)\) to some \(U\); so the
values converge, and so do the complex Hessians in every chart (cut off near the point and compare with the charts of the atlas). In determinant form the equation
reads \(\det(H_A+\partial\bar\partial\varphi_j)=e^{G_j}\det H_A\) and passes to the limit; the limit of the positive definite matrices
\(H_A+\partial\bar\partial\varphi_j\) is positive semidefinite with positive determinant \(e^{G_0}\det H_A\), hence positive definite. So \(U\in C^{2,s}(Y)\)
solves \(\mathrm{ma}_AU=(t_0F+\kappa_0)+\varepsilon_0U\) (Lemma~\ref{bou:monge-ampere-holder}) with \(H_A+\partial\bar\partial U\) positive definite in every chart, and
it is smooth by Lemma~\ref{bou:yau-regularity} (with \(t_0F+\kappa_0\) in place of \(tF\)); \(A_U\) is K\"ahler and the equation is the density equation.
\end{proof}

\begin{lemma}[The normalising constant]
\label{bou:yau-normalizing-constant}
\uses{bou:mixed-mass-class-invariance,bou:yau-volume-ratio}
\lean{TranscendentalNakamaye.Pending.YauNormalizingConstant.eq_zero_of_mongeAmpere_eq_exp_add,TranscendentalNakamaye.Pending.YauNormalizingConstant.integral_topRatio_formPow_add_ddc,TranscendentalNakamaye.Pending.YauNormalizingConstant.topRatio_formPow_eq_exp}
\leanok
Let \(Y\) be a compact connected complex manifold of dimension \(N\ge1\) with a K\"ahler form \(A\), and \(F\) a function with \(\int_Ye^FA^N=\int_YA^N\). If
\(\varphi\) is smooth, \(A_\varphi\) is K\"ahler and \(A_\varphi^N=e^{F+\kappa}A^N\) for a constant \(\kappa\), then \(\kappa=0\).
\end{lemma}
\begin{proof}
\leanok
\(\int_YA_\varphi^N=\int_YA^N\) (Lemma~\ref{bou:mixed-mass-class-invariance}), and the density of \(A_\varphi^N\) against \(A^N\) is \(e^{F+\kappa}\) (Lemma~\ref{bou:yau-volume-ratio});
so \(e^\kappa\int_Ye^FA^N=\int_YA^N\), i.e.\ \((e^\kappa-1)\int_YA^N=0\), and \(\int_YA^N>0\) (Lemma~\ref{bou:yau-volume-ratio}). Hence \(e^\kappa=1\).
\end{proof}

\begin{lemma}[Closedness]
\label{bou:yau-closedness}
\uses{bou:yau-continuity-path,bou:yau-sup-bound,bou:yau-holder-bound,bou:yau-limit-solutions}
\source{szekelyhidi-extremal:page-0056,szekelyhidi-extremal:page-0063,szekelyhidi-extremal:page-0067}
\lean{TranscendentalNakamaye.Pending.YauContinuityMethod.exists_smooth_of_mem_Icc,TranscendentalNakamaye.Pending.YauContinuityMethod.exists_abs_le}
\leanok
For \(0<\varepsilon\le1\), the set \(S_\varepsilon\) of Lemma~\ref{bou:yau-continuity-path} is closed in \([0,1]\).
\end{lemma}
\begin{proof}
\leanok
Let \(t_i\in S_\varepsilon\), \(t_i\to t\), with solutions \(\varphi_i\). Lemma~\ref{bou:yau-sup-bound} (with \(H=t_iF\), \(B=\sup|F|\)) gives
\(\varepsilon|\varphi_i|\le\sup|F|\); Lemma~\ref{bou:yau-holder-bound} (with \(S=\sup|F|/\varepsilon\), \(K=0\), \(\kappa=0\)) bounds \(\varphi_i\) in
\(C^{2,r}(Y)\) uniformly in \(i\); and Lemma~\ref{bou:yau-limit-solutions} (with \(\varepsilon_i=\varepsilon\), \(\kappa_i=0\)) gives a smooth solution of
\((\ast)_{t,\varepsilon}\). So \(t\in S_\varepsilon\).
\end{proof}

\begin{lemma}[The limit \(\varepsilon\to0\)]
\label{bou:yau-eps-limit}
\uses{bou:yau-continuity-path,bou:yau-sup-bound,bou:yau-c0-estimate,bou:yau-holder-bound,bou:yau-limit-solutions,bou:yau-normalizing-constant}
\source{szekelyhidi-extremal:page-0064}
\lean{TranscendentalNakamaye.Pending.YauContinuityMethod.exists_smooth_eq_exp_add,TranscendentalNakamaye.Pending.YauContinuityMethod.coeff_add_ddcClosedForm_sub_const,TranscendentalNakamaye.Pending.YauContinuityMethod.actualTopPowerDensity_pos,TranscendentalNakamaye.Pending.YauContinuityMethod.exists_kahler_mongeAmpere_eq_exp_of_estimates}
\leanok
Let \((Y,A)\), \(N\) and \(F\) be as in Theorem~\ref{bou:yau}. If \((\ast)_{1,\varepsilon}\) has a solution for every \(0<\varepsilon\le1\), then there is a smooth
\(\psi\) with \(A_\psi\) K\"ahler and \(A_\psi^N=e^FA^N\).
\end{lemma}
\begin{proof}
\leanok
Let \(\varphi_\varepsilon\) solve \((\ast)_{1,\varepsilon}\), \(\psi_\varepsilon=\varphi_\varepsilon-\sup\varphi_\varepsilon\) and \(\kappa_\varepsilon=\varepsilon\sup\varphi_\varepsilon\), so
that \(A_{\psi_\varepsilon}^N=e^{F+\varepsilon\psi_\varepsilon+\kappa_\varepsilon}A^N\). Lemma~\ref{bou:yau-sup-bound} (with \(H=F\)) gives \(\varepsilon|\varphi_\varepsilon|\le\sup|F|\),
so \(|\kappa_\varepsilon|\le\sup|F|\) and \(A_{\psi_\varepsilon}^N\le e^{2\sup|F|}A^N\); Lemma~\ref{bou:yau-c0-estimate} bounds the oscillation of \(\psi_\varepsilon\), and
\(\sup\psi_\varepsilon=0\), so \(|\psi_\varepsilon|\le S\) with \(S\) independent of \(\varepsilon\). Lemma~\ref{bou:yau-holder-bound} (with \(t=1\), \(K=\sup|F|\)) bounds
\(\psi_\varepsilon\) in \(C^{2,r}(Y)\) uniformly in \(\varepsilon\). Take \(\varepsilon_i\to0\) with \(\kappa_{\varepsilon_i}\to\kappa\) (a bounded sequence); Lemma~\ref{bou:yau-limit-solutions}
gives a smooth \(\psi\) with \(A_\psi\) K\"ahler and \(A_\psi^N=e^{F+\kappa}A^N\), and \(\kappa=0\) by Lemma~\ref{bou:yau-normalizing-constant}.
\end{proof}

\begin{theorem}[Yau: the complex Monge--Amp\`ere equation in a K\"ahler class]
\label{bou:yau}
\uses{ana:geometric-setting,bou:yau-continuity-path,bou:yau-openness,bou:yau-regularity,bou:yau-closedness,bou:yau-eps-limit}
\source{szekelyhidi-extremal:page-0064,dp04-kahler-cone:2.5,bou02-volume:3.4}
\lean{TranscendentalNakamaye.Pending.YauTheorem.exists_kahler_mongeAmpere_eq_exp,TranscendentalNakamaye.Pending.YauContinuityMethod.exists_kahler_mongeAmpere_eq_exp_of_holderBound}
\leanok
Let \((M,\omega)\) be a compact connected K\"ahler manifold of dimension \(N\ge1\) and \(F\) a smooth function with \(\int_Me^F\omega^N=\int_M\omega^N\). There is a
smooth function \(\varphi\) with \(\omega+dd^c\varphi>0\) and \((\omega+dd^c\varphi)^N=e^F\omega^N\).
\end{theorem}
\begin{proof}
\leanok
Fix \(0<\varepsilon\le1\) and let \(S_\varepsilon\subset[0,1]\) be the set of Lemma~\ref{bou:yau-continuity-path}: \(0\in S_\varepsilon\), \(S_\varepsilon\) is
open (Lemmas~\ref{bou:yau-openness} and~\ref{bou:yau-regularity}) and closed (Lemma~\ref{bou:yau-closedness}), so \(S_\varepsilon=[0,1]\) as \([0,1]\) is
connected. Thus \((\ast)_{1,\varepsilon}\) is solvable for every \(\varepsilon\), and Lemma~\ref{bou:yau-eps-limit} gives the solution.
\end{proof}
This is Yau's solution of Calabi's conjecture (S.-T. Yau, 1978), Theorem~3.14 of Sz\'ekelyhidi, quoted by Demailly--P\u{a}un and by Boucksom; the proof above is Sz\'ekelyhidi's continuity method (his Section~3.4), along the path \(e^{tF+\varepsilon\varphi}\) followed by
\(\varepsilon\to0\), with B\l{}ocki's uniform and second-order estimates. The source writes \(\sqrt{-1}\,\partial\bar\partial\varphi\), which differs from \(dd^c\varphi\) by a positive constant absorbed in \(\varphi\). A smooth positive
volume form \(f\) with \(\int_Mf=\int_M\omega^N\) is \(e^F\omega^N\) with \(F=\log(f/\omega^N)\). Connectedness is needed: on \(M_1\sqcup M_2\) the normalisation must hold
on each component.

\begin{lemma}[Volume ratios of K\"ahler forms]
\label{bou:yau-volume-ratio}
\uses{bou:integration-current,bou:mixed-mass-class-invariance,bou:nef-core,fpglobal:smooth-ddc,bc:local-linearity}
\lean{TranscendentalNakamaye.Pending.YauTheorem.actualTopPowerDensity_pos,TranscendentalNakamaye.Pending.YauTheorem.topRatio_formPow_symm_eq,TranscendentalNakamaye.Pending.YauTheorem.integral_topRatio_formPow_eq,TranscendentalNakamaye.Pending.YauTheorem.kahlerVolume_div_pos,TranscendentalNakamaye.Pending.YauTheorem.ddcClosedForm_coeff_add}
\leanok
Let \(A,\omega'\) be K\"ahler forms on a compact complex manifold \(M\) of dimension \(N\). In every chart the top-degree density of \(A^N\) is positive, and the
ratio \(\omega'^N/A^N\) of top-degree forms, read in a chart, is the quotient of the two chart densities; \(\int_M(\omega'^N/A^N)\,A^N=\int_M\omega'^N\); and if \(M\) is
connected and \(N\ge1\), then \(\int_MA^N/\int_M\omega'^N>0\). Moreover \(dd^c(f+g)=dd^cf+dd^cg\) for smooth functions \(f,g\) on \(M\).
\end{lemma}
\begin{proof}
\leanok
A K\"ahler form has positive definite coefficient matrix in every chart, so its top power has positive density; the ratio of two top-degree forms is the
quotient of their densities in any chart. The integral identity is the class invariance of the mixed mass (Lemma~\ref{bou:mixed-mass-class-invariance}) for the
smooth potential \(0\), written against the volume form of \(A\). Both volumes are positive on a nonempty connected compact manifold (Lemma~\ref{bou:nef-core}).
Additivity of \(dd^c\) is chartwise linearity of the smooth \(dd^c\) (Lemma~\ref{bc:local-linearity}).
\end{proof}

\begin{corollary}[The forms of Yau's theorem used here]
\label{bou:yau-consequences}
\uses{bou:yau,bou:yau-volume-ratio,bou:mixed-mass-class-invariance,bou:product-manifold,posclass:nef,bou:nef-core}
\source{bou02-volume:3.4,dp04-kahler-cone:2.12}
\lean{TranscendentalNakamaye.Pending.YauTheorem.exists_kahler_mongeAmpere_eq_smul,TranscendentalNakamaye.Pending.YauTheorem.yauHypothesis_of_nef,TranscendentalNakamaye.Pending.YauTheorem.exists_kahlerCurrent_of_nef,TranscendentalNakamaye.Pending.YauTheorem.isBigClass_of_nef}
\leanok
(1) For K\"ahler forms \(A,\omega'\) on a compact connected complex manifold \(M\) of dimension \(N\ge1\) there is a smooth \(\varphi\) with \(A+dd^c\varphi\) K\"ahler
and \((A+dd^c\varphi)^N=c\,\omega'^N\), \(c=\int_MA^N/\int_M\omega'^N\).
(2) If \(b\) is a nef smooth closed real \((1,1)\)-form on a compact connected K\"ahler manifold \((X,\omega)\), the hypothesis of Lemma~\ref{bou:nef-core} holds;
hence if moreover \(\int_Xb^n>0\), then \(b\) represents a K\"ahler current and \(\{b\}\) is big.
\end{corollary}
\begin{proof}
\leanok
(1) The ratio \(\omega'^N/A^N\) of top-degree forms does not depend on the chart and is smooth and positive (Lemma~\ref{bou:yau-volume-ratio}); Theorem~\ref{bou:yau} with \(F=\log(c\,\omega'^N/A^N)\)
applies, since \(\int_Me^FA^N=c\int_M\omega'^N=\int_MA^N\).
(2) For \(0<\varepsilon\le1\) nefness gives a smooth \(h\) with \(b+dd^ch+\tfrac\varepsilon2\omega\ge0\), so \(B=b+\varepsilon\omega+dd^ch\ge\tfrac\varepsilon2\omega\) is
K\"ahler. Apply (1) on \(X\times X\) to \(p_1^*B+p_2^*B\) and a K\"ahler form \(\omega'\), obtaining \(\varphi\); since \(p_i^*dd^ch=dd^c(h\circ p_i)\), the function
\(g(x,y)=h(x)+h(y)+\varphi(x,y)\) gives \(p_1^*(b+\varepsilon\omega)+p_2^*(b+\varepsilon\omega)+dd^cg=p_1^*B+p_2^*B+dd^c\varphi\), as Lemma~\ref{bou:nef-core}
requires. The last statement is Lemma~\ref{bou:nef-core}.
\end{proof}
Both parts are deduced from Theorem~\ref{bou:yau}.

\begin{lemma}[The pointwise estimate behind the concentration of mass]
\label{bou:concentration-pointwise}
\uses{fpmma:coefficient,fpmma:coefficient-laws,fpmma:coefficient-positive,btgc:coefficient-congruence}
\source{dp04-kahler-cone:2.6,bou02-volume:4.12}
\lean{TranscendentalNakamaye.Pending.ConcentrationEigenvalueBound.concentration_bound,TranscendentalNakamaye.Pending.ConcentrationEigenvalueBound.diagonal_concentration_bound,TranscendentalNakamaye.Pending.ConcentrationEigenvalueBound.exists_simultaneous_diagonalization,TranscendentalNakamaye.Pending.ConcentrationEigenvalueBound.mixedCoefficient_conj_ite,TranscendentalNakamaye.Pending.ConcentrationEigenvalueBound.mixedCoefficient_unitDiag,TranscendentalNakamaye.Pending.ConcentrationEigenvalueBound.mixedCoefficient_eq_zero_of_unitDiag,TranscendentalNakamaye.Pending.ConcentrationEigenvalueBound.mixedCoefficient_ite_diagonal,TranscendentalNakamaye.Pending.ConcentrationEigenvalueBound.le_re_mixedCoefficient_ite_diagonal,TranscendentalNakamaye.Pending.ConcentrationEigenvalueBound.prod_le_re_mixedCoefficient_ite_diagonal,TranscendentalNakamaye.Pending.ConcentrationEigenvalueBound.prod_le_re_mixedCoefficient_ite_diagonal'}
\leanok
Let \(A\) and \(B\) be positive definite and \(C\) positive semidefinite Hermitian \(n\times n\) matrices, \(0\le p\le n\) and \(t>0\). For the mixed coefficient
\(C(H_1,\dots,H_n)\) of Definition~\ref{fpmma:coefficient}, write \(C(A[k],B[n-k])\) for its value with \(A\) in the first \(k\) slots and \(B\) in the others. If
\[
 \operatorname{Re}C(A[n-p],B[p])<t\,\operatorname{Re}\det B ,
\]
then
\[
 \frac{\operatorname{Re}\det A}{t\,\operatorname{Re}\det B}\,\operatorname{Re}C(B[p],C[n-p])\ \le\ \operatorname{Re}C(A[p],C[n-p]).
\]
\end{lemma}
\begin{proof}
\leanok
By the spectral theorem, applied to \(B\) and then to the positive definite matrix \(B^{-1/2}AB^{-1/2}\), there are an invertible \(R\) and \(\lambda_1,\dots,\lambda_n>0\) with
\(R^*BR=1\) and \(R^*AR=\operatorname{diag}(\lambda)\). A congruence \(H\mapsto R^*HR\) of every slot multiplies the mixed coefficient by \(|\det R|^2>0\)
(Lemma~\ref{btgc:coefficient-congruence}) and keeps \(C\) positive semidefinite, and \(\det A=\prod_k\lambda_k/|\det R|^2\), \(\det B=1/|\det R|^2\); so it suffices to
treat \(B=1\), \(A=\operatorname{diag}(\lambda)\), and to show: if \(\operatorname{Re}C(\operatorname{diag}(\lambda)[n-p],1[p])<t\), then
\(t^{-1}\prod_k\lambda_k\operatorname{Re}C(1[p],C[n-p])\le\operatorname{Re}C(\operatorname{diag}(\lambda)[p],C[n-p])\).

Write \(\operatorname{diag}(\lambda)=\sum_k\lambda_kE_k\) with the diagonal unit matrices \(E_k\) and expand the diagonal slots by multilinearity. With every slot
diagonal, \(C(E_{j_1},\dots,E_{j_n})\) is \(1\) if \(j\) is a bijection and \(0\) otherwise, so
\(C(\operatorname{diag}(\lambda)[n-p],1[p])=\sum_\sigma\prod_{i<n-p}\lambda_{\sigma(i)}\), the sum over the permutations \(\sigma\); every product
\(\prod_{k\in T}\lambda_k\) with \(|T|=n-p\) occurs in it, all terms being positive, so each such product is \(<t\). Expanding the first \(p\) slots of
\(C(\operatorname{diag}(\lambda)[p],C[n-p])\) gives \(\sum_r\prod_{i<p}\lambda_{r(i)}\,C(E_{r(0)},\dots,E_{r(p-1)},C[n-p])\) over the maps \(r\) from the first \(p\)
slots to the indices; a term with \(r\) not injective vanishes, because two slots carrying the same \(E_k\) kill the mixed coefficient, and every remaining
term has nonnegative real part (Lemma~\ref{fpmma:coefficient-positive}). For injective \(r\) with image \(S\), \(\prod_{k\in S}\lambda_k=\prod_k\lambda_k/\prod_{k\notin S}\lambda_k
>t^{-1}\prod_k\lambda_k\), since \(|S^c|=n-p\). The same expansion of \(C(1[p],C[n-p])\) has the same terms with all coefficients \(1\). Comparing term by term
gives the claim.
\end{proof}
This replaces the eigenvalue inequalities \(\alpha_\varepsilon^p\ge\lambda_1\cdots\lambda_p\,\omega_\varepsilon^p\) and
\(\alpha_\varepsilon^{N-p}\wedge\omega_\varepsilon^p\ge\frac1{N!}\lambda_{p+1}\cdots\lambda_N\,\omega_\varepsilon^N\) of Demailly--P\u{a}un and Boucksom by one inequality without
eigenvalues: the set on which the product of the \(N-p\) largest eigenvalues is large is replaced by the set on which the mixed coefficient \(C(A[N-p],B[p])\)
is large compared with \(\det B\), which is defined by a continuous function of the point. The constant is the same up to the factor \(p!\,(n-p)!\), which
the argument does not see.

\begin{definition}[Currents of bidimension \((q,q)\) on test forms of product type; positivity, mass, push-forward, wedge]
\label{bou:product-currents}
\uses{fpgeom:real-form,hgeom:form-pullback,fpglobal:quasi-psh,posclass:background}
\source{demailly-cadg:III.1.13,demailly-cadg:III.1.4,demailly-cadg:III.1.21}
\lean{TranscendentalNakamaye.Pending.StronglyPositiveForms.spCone,TranscendentalNakamaye.Pending.ProductTestForms.testSpace,TranscendentalNakamaye.Pending.ProductTestForms.positiveCone,TranscendentalNakamaye.Pending.PositiveProductCurrents.ProductCurrent,TranscendentalNakamaye.Pending.PositiveProductCurrents.IsPositive,TranscendentalNakamaye.Pending.PositiveProductCurrents.kahlerTest,TranscendentalNakamaye.Pending.PositiveProductCurrents.mass,TranscendentalNakamaye.Pending.PositiveProductCurrents.pushforward,TranscendentalNakamaye.Pending.PositiveProductCurrents.wedge}
\leanok
Let \(M\) be a complex manifold of dimension \(n\) and \(q\ge0\). At a point, the strongly positive forms of degree \(2q\) are the convex cone generated by the
products \(\beta_1\wedge\dots\wedge\beta_q\) of semipositive real \((1,1)\)-forms (for \(q=0\): the nonnegative constants). The test forms of degree \(2q\)
are the real linear span of the forms \(h\,\beta_1\wedge\dots\wedge\beta_q\), with \(h\) a smooth function and \(\beta_j\) smooth real \((1,1)\)-forms on \(M\); a test
form is positive if it is strongly positive at every point. A current of bidimension \((q,q)\) is a real linear functional \(T\) on the test forms of degree
\(2q\); it is positive if \(T(\varphi)\ge0\) for every positive test form \(\varphi\). For a K\"ahler form \(\omega\), the mass of \(T\) is \(T(\omega^q)\). For a
holomorphic map \(f\colon V\to M\) the push-forward of a current \(T\) on \(V\) is \(f_*T=T\circ f^*\), and for a smooth real \((1,1)\)-form \(\beta\) on \(M\) the
wedge of a current \(T\) of bidimension \((q+1,q+1)\) with \(\beta\) is the current \((T\wedge\beta)(\varphi)=T(\varphi\wedge\beta)\) of bidimension \((q,q)\).
\end{definition}
On a compact manifold no support condition is needed. The products \(h\,\beta_1\wedge\dots\wedge\beta_q\) span all smooth real \((q,q)\)-forms, by
Demailly's basis of strongly positive forms (III.1.4) and a partition of unity, so these are the currents of bidimension \((q,q)\) of order \(0\) in the usual
sense; nothing below uses that identification. Demailly's trace measure (III.1.21) is \(\frac1{2^qq!}T\wedge\omega^q\); the mass here omits the factor.

\begin{lemma}[Mass, weak compactness, push-forward and wedge of positive currents]
\label{bou:positive-currents}
\uses{bou:product-currents}
\source{demailly-cadg:III.1.14,demailly-cadg:III.1.23,dp04-kahler-cone:2.12}
\lean{TranscendentalNakamaye.Pending.StronglyPositiveForms.orderUnit_step,TranscendentalNakamaye.Pending.ProductTestForms.exists_orderUnit_bound,TranscendentalNakamaye.Pending.ProductTestForms.pullbackTestLM_mem_positiveCone,TranscendentalNakamaye.Pending.OrderUnitWeakCompactness.isCompact_positive_le,TranscendentalNakamaye.Pending.OrderUnitWeakCompactness.exists_ultrafilter_tendsto,TranscendentalNakamaye.Pending.PositiveProductCurrents.mass_nonneg,TranscendentalNakamaye.Pending.PositiveProductCurrents.exists_abs_apply_le_mul_mass,TranscendentalNakamaye.Pending.PositiveProductCurrents.isCompact_positive_mass_le,TranscendentalNakamaye.Pending.PositiveProductCurrents.exists_ultrafilter_tendsto,TranscendentalNakamaye.Pending.PositiveProductCurrents.tendsto_mass,TranscendentalNakamaye.Pending.PositiveProductCurrents.IsPositive.pushforward,TranscendentalNakamaye.Pending.PositiveProductCurrents.exists_mass_pushforward_le,TranscendentalNakamaye.Pending.PositiveProductCurrents.tendsto_pushforward_apply,TranscendentalNakamaye.Pending.PositiveProductCurrents.IsPositive.wedge,TranscendentalNakamaye.Pending.PositiveProductCurrents.mass_wedge_kahler,TranscendentalNakamaye.Pending.PositiveProductCurrents.tendsto_wedge_apply}
\leanok
Let \(M\) be a compact complex manifold of dimension \(n\) with a K\"ahler form \(\omega\), and \(q\ge0\). For currents of bidimension \((q,q)\)
(Definition~\ref{bou:product-currents}):
\begin{enumerate}
\item a positive current has nonnegative mass, and for every test form \(\varphi\) there is \(c\ge0\) with \(|T(\varphi)|\le c\,\operatorname{mass}(T)\) for every positive
      \(T\);
\item for every \(m\), the positive currents of mass \(\le m\) form a compact set for pointwise convergence on test forms; consequently every family of
      positive currents whose masses are eventually \(\le m\) along a nontrivial filter converges pointwise, along some ultrafilter finer than that filter, to a positive
      current of mass \(\le m\); and the mass is continuous for pointwise convergence;
\item for a holomorphic map \(f\colon V\to M\), \(f_*\) maps positive currents to positive currents and is continuous for pointwise convergence; if \(V\) is
      compact with a K\"ahler form \(\omega_V\), there is \(c\ge0\) with \(\operatorname{mass}(f_*T)\le c\,\operatorname{mass}(T)\) for every positive \(T\) on \(V\);
\item for a smooth real \((1,1)\)-form \(\beta\), \(T\mapsto T\wedge\beta\) is continuous for pointwise convergence, \(T\wedge\beta\) is positive when \(T\) is positive
      and \(\beta\) is semipositive, and \(\operatorname{mass}(T\wedge\omega)=\operatorname{mass}(T)\).
\end{enumerate}
\end{lemma}
\begin{proof}
\leanok
(1) \(\omega^q\) is an order unit for the positive test forms: for every test form \(\varphi\) there is \(c\ge0\) with \(c\,\omega^q\pm\varphi\) positive. By
linearity it suffices to see this for \(\varphi=h\,\beta_1\wedge\dots\wedge\beta_q\), and by induction on \(q\) it follows from the pointwise step: if
\(c\,\omega^q\pm P\) and \(K\omega\pm\beta\) are strongly positive, so are \(cK\,\omega^{q+1}\pm P\wedge\beta\) (expand
\((c\omega^q+P)\wedge(K\omega+\beta)+(c\omega^q-P)\wedge(K\omega-\beta)\)); the constants are uniform on the compact \(M\) because the coefficients are continuous and
\(\omega\) is positive definite. Then \(T(c\,\omega^q\pm\varphi)\ge0\) gives \(|T(\varphi)|\le c\,T(\omega^q)\), and \(T(\omega^q)\ge0\) since \(\omega^q\) is positive.
(2) By (1) the positive currents of mass \(\le m\) lie in the product \(\prod_\varphi[-c_\varphi m,c_\varphi m]\), compact by Tychonoff, and form a closed subset of it:
linearity, positivity and the mass bound are closed conditions for pointwise convergence. In a compact space every ultrafilter converges, which gives the
weak limits; the mass is the value at \(\omega^q\). (3) The pull-back of a test form by a holomorphic map is a test form, and of a positive test form a positive
one, because pull-backs of semipositive \((1,1)\)-forms by complex-linear maps are semipositive; so \(f_*T\) is positive, and \(f_*\) is continuous since it is
evaluation at \(f^*\varphi\). The mass bound is (1) on \(V\) for the test form \(f^*\omega^q\). (4) \(\varphi\wedge\beta\) is a test form, positive when \(\varphi\) is
positive and \(\beta\) semipositive, and \(\omega^q\wedge\omega=\omega^{q+1}\).
\end{proof}

\begin{lemma}[The trace measure of a positive current; masses of closed sets under weak limits]
\label{bou:trace-measure}
\uses{bou:product-currents,bou:positive-currents}
\source{demailly-cadg:III.1.21,demailly-cadg:III.1.14,bou02-volume:4.12}
\lean{TranscendentalNakamaye.Pending.CurrentTraceMeasure.exists_chartSmooth_bump,TranscendentalNakamaye.Pending.CurrentTraceMeasure.chartSmoothSubalgebra_topologicalClosure,TranscendentalNakamaye.Pending.CurrentTraceMeasure.traceExt_nonneg,TranscendentalNakamaye.Pending.CurrentTraceMeasure.traceMeasure,TranscendentalNakamaye.Pending.CurrentTraceMeasure.integral_traceMeasure_chartSmooth,TranscendentalNakamaye.Pending.CurrentTraceMeasure.traceMeasure_univ,TranscendentalNakamaye.Pending.CurrentTraceMeasure.tendsto_integral_traceMeasure,TranscendentalNakamaye.Pending.CurrentTraceMeasure.le_traceMeasure_of_tendsto}
\leanok
Let \(M\) be a compact complex manifold with a K\"ahler form \(\omega\) and \(T\) a positive current of bidimension \((q,q)\). There is a finite regular Borel
measure \(\sigma_T\), the trace measure \(T\wedge\omega^q\), with \(\int_Mh\,d\sigma_T=T(h\,\omega^q)\) for every smooth function \(h\) and
\(\sigma_T(M)=\operatorname{mass}(T)\). If positive currents \(T_i\) converge pointwise along a filter to a positive current \(S\) and their masses are eventually
bounded, then \(\int_Mf\,d\sigma_{T_i}\to\int_Mf\,d\sigma_S\) for every continuous \(f\); and if moreover \(F\subseteq M\) is closed and
\(\sigma_{T_i}(F)\ge\delta\) eventually (along a filter which is not trivial), then \(\sigma_S(F)\ge\delta\).
\end{lemma}
\begin{proof}
\leanok
Smooth functions form a subalgebra of \(C(M,\mathbf R)\) containing the constants and separating points (bump functions in charts), so they are uniformly
dense by Stone--Weierstrass. The functional \(h\mapsto T(h\,\omega^q)\) on smooth functions is positive, since \(h\,\omega^q\) is a positive test form when
\(h\ge0\), hence bounded by \(\sup|h|\,T(\omega^q)\); it extends to a positive continuous functional on \(C(M,\mathbf R)\), and the Riesz--Markov--Kakutani theorem gives
\(\sigma_T\), with total mass \(T(\omega^q)\). For the convergence, approximate \(f\) uniformly within \(\epsilon\) by a smooth \(g\): \(\int g\,d\sigma_{T_i}=T_i(g\,\omega^q)\to
S(g\,\omega^q)=\int g\,d\sigma_S\), and the errors are at most \(\epsilon\) times the masses, which are bounded. For the closed set: if \(\sigma_S(F)<\delta\), outer
regularity gives an open \(U\supseteq F\) with \(\sigma_S(U)<\delta\), and Urysohn's lemma a continuous \(f\) with \(\mathbf 1_F\le f\le\mathbf 1_U\); then
\(\delta\le\sigma_{T_i}(F)\le\int f\,d\sigma_{T_i}\to\int f\,d\sigma_S\le\sigma_S(U)<\delta\) eventually, a contradiction.
\end{proof}

\begin{lemma}[Top-degree densities and the integration current of \(M\)]
\label{bou:integration-current}
\uses{bou:product-currents,fpmma:geometric-density,fpmma:coefficient-positive,gnpp:global-measure}
\source{demailly-cadg:III.1.13,demailly-cadg:III.1.4}
\lean{TranscendentalNakamaye.Pending.TopDegreeDensity.topDensity,TranscendentalNakamaye.Pending.TopDegreeDensity.topDensity_compContinuousLinearMap,TranscendentalNakamaye.Pending.TopDegreeDensity.wedgeProd_apply,TranscendentalNakamaye.Pending.TopDegreeDensity.topDensity_wedgeProd,TranscendentalNakamaye.Pending.TopDegreeDensity.normalizedHermitianMatrix_posSemidef,TranscendentalNakamaye.Pending.TopDegreeDensity.actualMixedTopDensity_nonneg,TranscendentalNakamaye.Pending.TopDegreeDensity.spCone_le_prodCone,TranscendentalNakamaye.Pending.TopDegreeDensity.topDensity_nonneg_of_mem_spCone,TranscendentalNakamaye.Pending.IntegrationCurrent.chartRatio_transition,TranscendentalNakamaye.Pending.IntegrationCurrent.topRatio,TranscendentalNakamaye.Pending.IntegrationCurrent.continuous_topRatio,TranscendentalNakamaye.Pending.IntegrationCurrent.topRatio_formPow,TranscendentalNakamaye.Pending.IntegrationCurrent.topRatio_nonneg,TranscendentalNakamaye.Pending.IntegrationCurrent.kahlerVolume,TranscendentalNakamaye.Pending.IntegrationCurrent.integrationCurrent,TranscendentalNakamaye.Pending.IntegrationCurrent.integrationCurrent_apply,TranscendentalNakamaye.Pending.IntegrationCurrent.isPositive_integrationCurrent,TranscendentalNakamaye.Pending.IntegrationCurrent.mass_integrationCurrent}
\leanok
Let \(M\) be a compact complex manifold of dimension \(n\) with a K\"ahler form \(\omega\). At a point, a strongly positive form of degree \(2n\) is a
nonnegative combination of products of \(n\) semipositive real \((1,1)\)-forms and has nonnegative density; the density of a top-degree form is
multiplied by \(\det L\) under a real-linear change of coordinates \(L\). Consequently, for a smooth form \(\varphi\) of degree \(2n\) the ratio
\(\varphi/\omega^n\) of densities is a well-defined continuous function on \(M\), nonnegative when \(\varphi\) is a positive test form and equal to \(1\) for
\(\varphi=\omega^n\). With \(dV_\omega\) the volume measure \(\omega^n\) (the non-pluripolar measure of the zero potential in \(\{\omega\}\),
Definition~\ref{gnpp:global-measure}), the integration current
\[
 [M](\varphi)=\int_M\frac{\varphi}{\omega^n}\,dV_\omega
\]
is a positive current of bidimension \((n,n)\) (Definition~\ref{bou:product-currents}) of mass \(\int_MdV_\omega\).
\end{lemma}
\begin{proof}
\leanok
By induction on the degree, a product of strongly positive forms with a semipositive \((1,1)\)-form lies in the cone generated by products of semipositive
\((1,1)\)-forms. The density of a product \(\beta_1\wedge\dots\wedge\beta_n\) of real \((1,1)\)-forms is their mixed density, which is
\(\pi^{-n}\operatorname{Re}C(H_{\beta_1},\dots,H_{\beta_n})\) (Lemma~\ref{fpmma:geometric-density}); the Hermitian matrix of a semipositive form is positive
semidefinite, so the density is \(\ge0\) (Lemma~\ref{fpmma:coefficient-positive}). A top-degree form evaluated on the image of the standard frame under
\(L\) is \(\det L\) times its value on the frame, so both \(\varphi\) and \(\omega^n\) change by the same Jacobian between charts and their ratio is
chart-independent; it is continuous because the coefficients are, and \(\omega^n\) has positive density. Linearity of \([M]\) is linearity of the ratio and
of the integral, positivity is the nonnegativity of the ratio, and \([M](\omega^n)=\int_M1\,dV_\omega\).
\end{proof}

\begin{lemma}[The global mixed measure of smooth semipositive forms has the smooth density]
\label{bou:smooth-mixed-density}
\uses{gnpp:global-measure,fpmma:geometric-density,bou:integration-current}
\source{begz-big-classes:equ:nonpp}
\lean{TranscendentalNakamaye.Pending.SmoothMixedMeasureDensity.isSemipositiveOneOne_coeff,TranscendentalNakamaye.Pending.SmoothMixedMeasureDensity.chartDensity,TranscendentalNakamaye.Pending.SmoothMixedMeasureDensity.continuousOn_chartDensity,TranscendentalNakamaye.Pending.SmoothMixedMeasureDensity.exists_ball_globalMixedMeasure_image_eq,TranscendentalNakamaye.Pending.SmoothMixedMeasureDensity.globalMixedMeasure_image_eq_lintegral,TranscendentalNakamaye.Pending.SmoothMixedMeasureDensity.restrict_image_eq_map,TranscendentalNakamaye.Pending.SmoothMixedMeasureDensity.setLIntegral_image_eq}
\leanok
Let \(M\) be a compact complex manifold of dimension \(n\) and \(\theta_1,\dots,\theta_n\) smooth closed semipositive real \((1,1)\)-forms. The global
mixed non-pluripolar measure of the \(\theta_i\) with the potentials \(0\) (Definition~\ref{gnpp:global-measure}) is, in every chart \(e\) and on
every measurable subset \(S\) of its target, \(\int_SD\,d\lambda\), where \(D\) is the density of \(\theta_1\wedge\dots\wedge\theta_n\) in the chart
(the mixed density of Lemma~\ref{fpmma:geometric-density}).
\end{lemma}
\begin{proof}
\leanok
Near a point each \(\theta_i\) has a smooth plurisubharmonic local potential, and on small balls the global measure is the smooth mixed
Monge--Amp\`ere measure of these potentials, whose density is \(D\). A measurable subset of the chart target is a disjoint countable union of
pieces inside such balls.
\end{proof}

\begin{lemma}[The current of smooth closed semipositive forms and its trace measure]
\label{bou:smooth-form-current}
\uses{bou:smooth-mixed-density,bou:product-currents,bou:integration-current,bou:positive-currents,bou:trace-measure,fpmma:geometric-density,gnpp:global-measure}
\source{demailly-cadg:III.1.13,demailly-cadg:III.1.21}
\lean{TranscendentalNakamaye.Pending.SmoothCurrentTraceMeasure.chainTestLM,TranscendentalNakamaye.Pending.SmoothCurrentTraceMeasure.chainTestLM_coeff,TranscendentalNakamaye.Pending.SmoothCurrentTraceMeasure.chainTestLM_mem_positiveCone,TranscendentalNakamaye.Pending.SmoothCurrentTraceMeasure.measure_eq_of_integral_chartSmooth_eq,TranscendentalNakamaye.Pending.SmoothCurrentTraceMeasure.traceMeasure_eq_withDensity,TranscendentalNakamaye.Pending.SmoothCurrentTraceMeasure.withDensity_topRatio_eq_globalMixedMeasure,TranscendentalNakamaye.Pending.SmoothCurrentTraceMeasure.smoothCurrent,TranscendentalNakamaye.Pending.SmoothCurrentTraceMeasure.isPositive_smoothCurrent,TranscendentalNakamaye.Pending.SmoothCurrentTraceMeasure.traceMeasure_smoothCurrent,TranscendentalNakamaye.Pending.SmoothCurrentTraceMeasure.mass_smoothCurrent,TranscendentalNakamaye.Pending.SmoothCurrentTraceMeasure.traceMeasure_smoothCurrent_eq_of_slots,TranscendentalNakamaye.Pending.SmoothCurrentTraceMeasure.smoothCurrentOf,TranscendentalNakamaye.Pending.SmoothCurrentTraceMeasure.isPositive_smoothCurrentOf,TranscendentalNakamaye.Pending.SmoothCurrentTraceMeasure.traceMeasure_smoothCurrentOf,TranscendentalNakamaye.Pending.SmoothCurrentTraceMeasure.mass_smoothCurrentOf}
\leanok
Let \(M\) be a compact complex manifold of dimension \(N\) with a K\"ahler form \(\omega\), \(0\le p\le N\), and \(\gamma_1,\dots,\gamma_p\) smooth closed
semipositive real \((1,1)\)-forms. Then \(T=[M]\wedge\gamma_1\wedge\dots\wedge\gamma_p\), \(T(\varphi)=\int_M\gamma_1\wedge\dots\wedge\gamma_p\wedge\varphi\), is a positive
current of bidimension \((N-p,N-p)\), and its trace measure \(T\wedge\omega^{N-p}\) is the global mixed non-pluripolar measure of the forms
\((\gamma_1,\dots,\gamma_p,\omega,\dots,\omega)\) with the potentials \(0\) (Definition~\ref{gnpp:global-measure}); in particular its mass is the
total mass of that measure, \(\int_M\gamma_1\wedge\dots\wedge\gamma_p\wedge\omega^{N-p}\), and in every chart its density is the mixed density of
Lemma~\ref{fpmma:geometric-density}.
\end{lemma}
\begin{proof}
\leanok
\(T=[M]\circ(\varphi\mapsto\varphi\wedge\gamma_1\wedge\dots\wedge\gamma_p)\) with the integration current of Lemma~\ref{bou:integration-current}; the map
sends positive test forms to positive test forms (products with semipositive forms, pointwise), so \(T\) is positive. For a smooth \(h\),
\(T(h\,\omega^{N-p})=\int_Mh\,R\,dV_\omega\) with \(R\) the ratio of the densities of \(\gamma_1\wedge\dots\wedge\gamma_p\wedge\omega^{N-p}\) and \(\omega^N\), a continuous
nonnegative function, so the trace measure is \(R\,dV_\omega\) (Lemma~\ref{bou:trace-measure}: finite Borel measures on \(M\) are determined by the
integrals of smooth functions, which are dense in \(C(M,\mathbf R)\)). In every chart \(R\,dV_\omega\) has the mixed density of the slots, which is
also the density of the global mixed measure (Lemma~\ref{bou:smooth-mixed-density}); covering \(M\) by charts identifies the two measures. The order of
the slots does not matter (the global mixed measure is symmetric).
\end{proof}

\begin{definition}[Closed currents of bidimension \((q,q)\)]
\label{bou:closed-currents}
\uses{bou:product-currents,fpgeom:real-form}
\source{demailly-cadg:III.1.13,demailly-cadg:III.2.3,demailly-cadg:III.2.13}
\lean{TranscendentalNakamaye.Pending.ClosedProductCurrents.complexJ,TranscendentalNakamaye.Pending.ClosedProductCurrents.pi11,TranscendentalNakamaye.Pending.ClosedProductCurrents.pi11CLM,TranscendentalNakamaye.Pending.ClosedProductCurrents.pi11Form,TranscendentalNakamaye.Pending.ClosedProductCurrents.oneWedge,TranscendentalNakamaye.Pending.ClosedProductCurrents.exactPair,TranscendentalNakamaye.Pending.ClosedProductCurrents.closedTestForm,TranscendentalNakamaye.Pending.ClosedProductCurrents.IsClosedCurrent}
\leanok
Let \(M\) be a complex manifold. For a real two-form \(A\) at a point, its \((1,1)\)-part is \(\Pi_{1,1}A(u,v)=\frac12\bigl(A(u,v)+A(Ju,Jv)\bigr)\); for a smooth
real two-form on \(M\) this is computed in charts. A current \(T\) of bidimension \((k+1,k+1)\) (Definition~\ref{bou:product-currents}) is \emph{closed} if
\[
 T\bigl(\Pi_{1,1}d(f\,dg)\wedge\beta_1\wedge\dots\wedge\beta_k\bigr)=0
\]
for all smooth functions \(f,g\) on \(M\) and all smooth closed real \((1,1)\)-forms \(\beta_1,\dots,\beta_k\). Every current of bidimension \((0,0)\) is
closed.
\end{definition}
If \(T\) extends to a current \(\widehat T\) on all smooth forms of degree \(2(k+1)\), vanishing on the forms whose \((k+1,k+1)\)-part is \(0\), and
\(d\widehat T=0\) in the usual sense, then \(\widehat T(d(f\,dg\wedge\beta))=\widehat T(d(f\,dg)\wedge\beta)=T(\Pi_{1,1}d(f\,dg)\wedge\beta)\) for a product
\(\beta\) of closed \((1,1)\)-forms, because the \((2,0)+(0,2)\)-part of \(d(f\,dg)\) wedged with a \((k,k)\)-form has no \((k+1,k+1)\)-part; so closed currents
in the sense of CADG~III.1 are closed here. The definition tests exactly what the Skoda--El Mir theorem and the support theorems use (CADG~III.2.3,
III.2.10, III.2.13): \(g\) a cut-off coordinate function and \(\beta\) the forms \(dd^c\) of cut-off quadratic functions, which are globally closed. On a
compact manifold every one-form is a finite sum of forms \(f\,dg\). Demailly's basis lemma (III.1.4) is not used.

\begin{lemma}[Elementary properties of closed currents]
\label{bou:closed-currents-limits}
\uses{bou:closed-currents,bou:product-currents}
\source{demailly-cadg:III.1.13}
\lean{TranscendentalNakamaye.Pending.ClosedProductCurrents.pi11_apply,TranscendentalNakamaye.Pending.ClosedProductCurrents.isOneOne_pi11,TranscendentalNakamaye.Pending.ClosedProductCurrents.pi11_eq_self,TranscendentalNakamaye.Pending.ClosedProductCurrents.isTypeOneOne_pi11Form,TranscendentalNakamaye.Pending.ClosedProductCurrents.oneWedge_apply_two,TranscendentalNakamaye.Pending.ClosedProductCurrents.extDeriv_smul_eq,TranscendentalNakamaye.Pending.ClosedProductCurrents.extDeriv_zeroForm_eq,TranscendentalNakamaye.Pending.ClosedProductCurrents.exactPair_coeff_apply,TranscendentalNakamaye.Pending.ClosedProductCurrents.closedTestForm_snoc,TranscendentalNakamaye.Pending.ClosedProductCurrents.isClosedCurrent_succ_iff,TranscendentalNakamaye.Pending.ClosedProductCurrents.isClosedCurrent_zero,TranscendentalNakamaye.Pending.ClosedProductCurrents.IsClosedCurrent.add,TranscendentalNakamaye.Pending.ClosedProductCurrents.IsClosedCurrent.smul,TranscendentalNakamaye.Pending.ClosedProductCurrents.IsClosedCurrent.sub,TranscendentalNakamaye.Pending.ClosedProductCurrents.isClosedCurrent_of_tendsto,TranscendentalNakamaye.Pending.ClosedProductCurrents.IsClosedCurrent.wedge}
\leanok
Let \(M\) be a complex manifold.
\begin{enumerate}
\item \(\Pi_{1,1}A\) is of type \((1,1)\), and \(\Pi_{1,1}A=A\) when \(A\) is of type \((1,1)\); \(\Pi_{1,1}\) commutes with holomorphic changes of
      coordinates, so it maps smooth real two-forms on \(M\) to smooth real \((1,1)\)-forms.
\item In a chart, \(d(F\alpha)=dF\wedge\alpha+F\,d\alpha\) for a smooth function \(F\) and a smooth form \(\alpha\), and
      \(d(f\,dg)(u,v)=df(u)\,dg(v)-df(v)\,dg(u)\).
\item The closed currents of bidimension \((q,q)\) form a linear subspace; a pointwise limit, along a nontrivial filter, of closed currents is closed;
      and if \(T\) is closed of bidimension \((k+2,k+2)\) and \(\gamma\) is a smooth closed real \((1,1)\)-form, then \(T\wedge\gamma\) is closed.
\end{enumerate}
\end{lemma}
\begin{proof}
\leanok
(1) \(J\) is complex linear with \(J^2=-1\), so \(\Pi_{1,1}A(Ju,Jv)=\Pi_{1,1}A(u,v)\), which is the \((1,1)\) condition; chart changes have complex
linear differentials, so they commute with \(J\). (2) The Leibniz rule for the exterior derivative of a product with a function, and
\(d(dg)=0\). (3) Closedness is a family of linear equations \(T(\varphi)=0\) for fixed test forms \(\varphi\), preserved by linear combinations and by
pointwise limits; and \((T\wedge\gamma)(\Pi_{1,1}d(f\,dg)\wedge\beta)=T(\Pi_{1,1}d(f\,dg)\wedge\beta\wedge\gamma)\), a test of the closedness of \(T\)
since \(\gamma\) is closed.
\end{proof}

\begin{lemma}[Closedness of a positive current may be tested locally, with semipositive slots]
\label{bou:closed-currents-local}
\uses{bou:closed-currents,bou:closed-currents-limits,bou:positive-currents}
\source{demailly-cadg:III.2.3,demailly-cadg:III.2.13}
\lean{TranscendentalNakamaye.Pending.ClosedCurrentLocalTests.wedgeProd_update_add,TranscendentalNakamaye.Pending.ClosedCurrentLocalTests.wedgeProd_update_smul,TranscendentalNakamaye.Pending.ClosedCurrentLocalTests.wedgeProd_update_zero,TranscendentalNakamaye.Pending.ClosedCurrentLocalTests.wedgeProd_update_sum,TranscendentalNakamaye.Pending.ClosedCurrentLocalTests.slot_update_coeff,TranscendentalNakamaye.Pending.ClosedCurrentLocalTests.coeff_sum,TranscendentalNakamaye.Pending.ClosedCurrentLocalTests.closedTestForm_update_add,TranscendentalNakamaye.Pending.ClosedCurrentLocalTests.closedTestForm_update_smul,TranscendentalNakamaye.Pending.ClosedCurrentLocalTests.exactPair_coeff_sum,TranscendentalNakamaye.Pending.ClosedCurrentLocalTests.closedTestForm_coeff_sum,TranscendentalNakamaye.Pending.ClosedCurrentLocalTests.apply_closedTestForm_eq_zero_of_semipositive,TranscendentalNakamaye.Pending.ClosedCurrentLocalTests.isClosedCurrent_of_local}
\leanok
Let \(M\) be a compact complex manifold with a K\"ahler form \(\omega\).
\begin{enumerate}
\item The test form \(\Pi_{1,1}d(f\,dg)\wedge\beta_1\wedge\dots\wedge\beta_k\) is linear in each closed slot \(\beta_j\), and, at every chart point, it is the sum
      of the test forms of \(\chi_af\) for a smooth partition of unity \((\chi_a)\).
\item If a current \(S\) of bidimension \((k+1,k+1)\) vanishes on these test forms whenever all \(\beta_j\) are closed and semipositive, it vanishes on
      all of them.
\item A positive current \(S\) of bidimension \((k+1,k+1)\) is closed as soon as, for some open neighbourhood \(U_x\) of every point \(x\),
      \(S(\Pi_{1,1}d(f\,dg)\wedge\beta_1\wedge\dots\wedge\beta_k)=0\) for all smooth \(f\) with compact support in some \(U_x\), all smooth \(g\) and all closed
      semipositive \(\beta_j\).
\end{enumerate}
\end{lemma}
\begin{proof}
\leanok
(1) Products of two-forms are multilinear, and \(d(\sum_a\chi_af\,dg)=d(f\,dg)\) since \(\sum_a\chi_a=1\). (2) On the compact \(M\), \(\beta+A\omega\ge0\) for some
\(A>0\), and \(\beta=(\beta+A\omega)-A\omega\) is a difference of closed semipositive forms; by (1) and induction on the number of slots that are not
semipositive, the value vanishes. (3) Take a partition of unity subordinate to the \(U_x\); a positive current only sees the coefficients of a test
form at chart points (Lemma~\ref{bou:positive-currents}), so by (1) its value on a test form of \(f\) is the sum of its values on those of the \(\chi_af\),
each zero by hypothesis and (2).
\end{proof}

\begin{lemma}[Smooth currents of closed forms are closed]
\label{bou:smooth-current-closed}
\uses{bou:closed-currents,bou:closed-currents-limits,bou:smooth-form-current,bou:integration-current,fpmma:geometric-density,lpot:existence-global,fpmma:hessian-ibp}
\source{demailly-cadg:III.1.13}
\lean{TranscendentalNakamaye.Pending.SmoothCurrentClosedness.pairMatrix,TranscendentalNakamaye.Pending.SmoothCurrentClosedness.pairMatrix_cofactor_eq_sum_wirt,TranscendentalNakamaye.Pending.SmoothCurrentClosedness.integral_pairMatrix_cofactor_eq_zero,TranscendentalNakamaye.Pending.SmoothCurrentClosedness.normalizedHermitianMatrix_pi11_eq,TranscendentalNakamaye.Pending.SmoothCurrentClosedness.actualMixedTopDensity_update_eq,TranscendentalNakamaye.Pending.SmoothCurrentClosedness.actualMixedTopDensity_closedSlots_eq,TranscendentalNakamaye.Pending.SmoothCurrentClosedness.contDiff_mul_of_tsupport_subset,TranscendentalNakamaye.Pending.SmoothCurrentClosedness.exactPair_coeff_eq_zero,TranscendentalNakamaye.Pending.SmoothCurrentClosedness.setIntegral_closedSlots_eq_zero,TranscendentalNakamaye.Pending.SmoothCurrentClosedness.setIntegral_image_eq,TranscendentalNakamaye.Pending.SmoothCurrentClosedness.integral_topRatio_eq_zero_of_slotChart,TranscendentalNakamaye.Pending.SmoothCurrentClosedness.integral_topRatio_eq_zero,TranscendentalNakamaye.Pending.SmoothCurrentClosedness.closedTestForm_coeff,TranscendentalNakamaye.Pending.SmoothCurrentClosedness.isClosedCurrent_smoothCurrent,TranscendentalNakamaye.Pending.SmoothCurrentClosedness.isClosedCurrent_smoothCurrentOf}
\leanok
Let \(M\) be a compact complex manifold of dimension \(N\) with a K\"ahler form \(\omega\), \(0\le p\le N\), and \(\gamma_1,\dots,\gamma_p\) smooth closed
real \((1,1)\)-forms. The current \([M]\wedge\gamma_1\wedge\dots\wedge\gamma_p\) of Lemma~\ref{bou:smooth-form-current} is closed.
\end{lemma}
\begin{proof}
\leanok
For \(p=N\) there is nothing to prove. Otherwise we must show \(\int_M\Pi_{1,1}d(f\,dg)\wedge P=0\) for smooth \(f,g\) and a product \(P\) of \(N-1\) smooth
closed real \((1,1)\)-forms. By a partition of unity we may assume that \(f\) is supported in a chart ball \(B\) on which every factor of \(P\) is
\(dd^c\rho_i\) for a smooth \(\rho_i\) (Corollary~\ref{lpot:existence-global}). In the chart the density of a product of \(N\) real \((1,1)\)-forms is
\(\pi^{-N}\operatorname{Re}C(H_1,\dots,H_N)\) of their Hermitian matrices (Lemma~\ref{fpmma:geometric-density}); the matrix of \(\Pi_{1,1}d(f\,dg)\) is, by
Lemma~\ref{bou:closed-currents-limits}(2), a constant multiple of \(\bigl(\partial_jf\,\bar\partial_kg-\partial_jg\,\bar\partial_kf\bigr)_{j,k}\). Expanding the
mixed coefficient in its first slot with the cofactors of \((\operatorname{Hess}\rho_1,\dots,\operatorname{Hess}\rho_{N-1})\), which form a divergence-free
field in the \(\partial\) and in the \(\bar\partial\) directions, and integrating by parts in \(f\) (Theorem~\ref{fpmma:hessian-ibp}) turns both halves of the
matrix into \(\mp\int f\,C(\operatorname{Hess}g,\operatorname{Hess}\rho_1,\dots)\), which cancel. (Equivalently: in top degree \(\Pi_{1,1}\) is invisible,
\(d(f\,dg)\wedge P=d(f\,dg\wedge P)\) because \(P\) is closed, and Stokes's theorem on the compact \(M\).)
\end{proof}

\begin{lemma}[Total mixed masses of smooth closed forms depend only on their classes]
\label{bou:mixed-mass-class-invariance}
\uses{bou:smooth-mixed-density,bou:integration-current,lpot:existence-global,fpmma:hessian-ibp,fpmma:geometric-density,gnpp:global-measure}
\source{demailly-cadg:III.1.13}
\lean{TranscendentalNakamaye.Pending.MixedMassClassInvariance.smoothDDC_sum,TranscendentalNakamaye.Pending.MixedMassClassInvariance.normalizedHermitianMatrix_zero,TranscendentalNakamaye.Pending.MixedMassClassInvariance.complexHessian_const,TranscendentalNakamaye.Pending.MixedMassClassInvariance.setIntegral_ddcSlots_eq_zero,TranscendentalNakamaye.Pending.MixedMassClassInvariance.integral_topRatio_eq_zero_of_local,TranscendentalNakamaye.Pending.MixedMassClassInvariance.integral_topRatio_eq_zero_of_slot,TranscendentalNakamaye.Pending.MixedMassClassInvariance.integral_topRatio_ddc_eq_zero,TranscendentalNakamaye.Pending.MixedMassClassInvariance.formProd,TranscendentalNakamaye.Pending.MixedMassClassInvariance.formProd_coeff,TranscendentalNakamaye.Pending.MixedMassClassInvariance.slotForm,TranscendentalNakamaye.Pending.MixedMassClassInvariance.topDensity_slotForm_coeff,TranscendentalNakamaye.Pending.MixedMassClassInvariance.update_coeff,TranscendentalNakamaye.Pending.MixedMassClassInvariance.integral_topRatio_slotForm_update_add_ddc,TranscendentalNakamaye.Pending.MixedMassClassInvariance.integral_topRatio_slotForm_add_ddc,TranscendentalNakamaye.Pending.MixedMassClassInvariance.globalMixedMeasure_real_univ_eq_integral,TranscendentalNakamaye.Pending.MixedMassClassInvariance.globalMixedMeasure_real_univ_add_ddc}
\leanok
Let \(M\) be a compact K\"ahler manifold of dimension \(N\) and \(\theta_1,\dots,\theta_N\) smooth closed real \((1,1)\)-forms. For every smooth \(f\),
\(\int_Mdd^cf\wedge\theta_2\wedge\dots\wedge\theta_N=0\); hence for smooth \(f_1,\dots,f_N\)
\[
 \int_M(\theta_1+dd^cf_1)\wedge\dots\wedge(\theta_N+dd^cf_N)=\int_M\theta_1\wedge\dots\wedge\theta_N .
\]
If the forms \(\theta_i\) and \(\theta_i+dd^cf_i\) are semipositive, the global mixed measures of the two families with the potentials \(0\)
(Definition~\ref{gnpp:global-measure}) have the same total mass.
\end{lemma}
\begin{proof}
\leanok
In a chart where every \(\theta_i\) is \(dd^c\rho_i\) (Corollary~\ref{lpot:existence-global}) and for \(f\) supported there, the density of
\(dd^cf\wedge\theta_2\wedge\dots\) is \(\pi^{-N}\operatorname{Re}C(\operatorname{Hess}f,\operatorname{Hess}\rho_2,\dots)\) (Lemma~\ref{fpmma:geometric-density}), and the
exchange of Hessians (Theorem~\ref{fpmma:hessian-ibp}) gives \(\int C(\operatorname{Hess}f,\dots)=\int f\,C(\operatorname{Hess}1,\dots)=0\); a partition of unity
subordinate to such charts and the change of variables in charts globalise this. Replacing the slots one at a time gives the identity. With
semipositive slots and potentials \(0\) the total mass of the global mixed measure is the integral of the smooth density
(Lemma~\ref{bou:smooth-mixed-density}).
\end{proof}

\begin{lemma}[The integration step of the concentrated current]
\label{bou:concentrated-current-integration}
\uses{bou:concentration-pointwise,bou:smooth-mixed-density,fpmma:geometric-density,fpmma:coefficient-laws,fpmma:coefficient-positive,gnpp:global-measure}
\source{bou02-volume:4.12,dp04-kahler-cone:2.6}
\lean{TranscendentalNakamaye.Pending.ConcentrationMassComparison.re_det_pos_of_posDef,TranscendentalNakamaye.Pending.ConcentrationMassComparison.mixedCoefficient_smul_family,TranscendentalNakamaye.Pending.ConcentrationMassComparison.re_mixedCoefficient_le_of_mongeAmpere,TranscendentalNakamaye.Pending.ConcentrationMassComparison.isOneOne_add,TranscendentalNakamaye.Pending.ConcentrationMassComparison.isOneOne_smul,TranscendentalNakamaye.Pending.ConcentrationMassComparison.normalizedHermitianMatrix_posDef,TranscendentalNakamaye.Pending.ConcentrationMassComparison.actualTopPowerDensity_eq_det,TranscendentalNakamaye.Pending.ConcentrationMassComparison.actualMixedTopDensity_eq_re,TranscendentalNakamaye.Pending.ConcentrationMassComparison.chartDensity_le_of_mongeAmpere,TranscendentalNakamaye.Pending.ConcentrationMassComparison.measure_le_of_forall_chart,TranscendentalNakamaye.Pending.ConcentrationMassComparison.mixedMeasure_le_of_mongeAmpere,TranscendentalNakamaye.Pending.ConcentrationMassComparison.le_mixedMeasure_of_mongeAmpere}
\leanok
Let \((M,\omega)\) be a compact K\"ahler manifold of dimension \(N\), \(0\le p\le N\), and \(\omega_\varepsilon\), \(\alpha\) smooth closed real \((1,1)\)-forms with
\(\omega_\varepsilon\ge\frac12\omega\), \(\alpha\) K\"ahler and \(\alpha^N=c\,\omega_\varepsilon^N\) for a constant \(c>0\). Write \(\int_S\theta_1\wedge\dots\wedge\theta_N\) for
the global mixed measure of the semipositive slots \(\theta_i\) with the potentials \(0\) (Definition~\ref{gnpp:global-measure}) of a set \(S\). Then for every
\(t>0\)
\[
 \int_S\omega_\varepsilon^p\wedge\omega^{N-p}\ \le\ \frac tc\int_S\alpha^p\wedge\omega^{N-p}+\frac{2^{N-p}N!}t\int_S\alpha^{N-p}\wedge\omega_\varepsilon^p ,
\]
and if \(\int_S\omega_\varepsilon^p\wedge\omega^{N-p}\ge\delta>0\) and \(\int_S\alpha^{N-p}\wedge\omega_\varepsilon^p\le M_0\) with \(M_0>0\), then
\[
 \int_S\alpha^p\wedge\omega^{N-p}\ \ge\ \frac{c\,\delta^2}{4\cdot2^{N-p}N!\,M_0}.
\]
\end{lemma}
\begin{proof}
\leanok
Pointwise, with \(A\), \(B\), \(C\) the Hermitian matrices of \(\alpha\), \(\omega_\varepsilon\), \(\omega\) in a chart (\(A,B\) positive definite, \(0\le C\le2B\),
\(\det A=c\det B\)): \(\operatorname{Re}C(B[p],C[N-p])\le\frac tc\operatorname{Re}C(A[p],C[N-p])+\frac{2^{N-p}N!}t\operatorname{Re}C(A[N-p],B[p])\) --- where
\(\operatorname{Re}C(A[N-p],B[p])\ge t\,N!\det B\) the last term alone dominates (\(\operatorname{Re}C(B[p],C[N-p])\le2^{N-p}\operatorname{Re}C(B[N])=2^{N-p}N!\det B\)), and
elsewhere Lemma~\ref{bou:concentration-pointwise} gives the first. The densities of the mixed measures are these mixed coefficients divided by \(\pi^N\)
(Lemmas~\ref{fpmma:geometric-density} and~\ref{bou:smooth-mixed-density}), so the inequality holds for the measures of the images of measurable subsets of
chart targets, hence for all sets (finitely many charts cover \(M\)). For the second claim take \(t=2\cdot2^{N-p}N!\,M_0/\delta\): then
\(\delta\le\frac tc X+\frac\delta2\), i.e.\ \(X\ge\frac{c\,\delta}{2t}\).
\end{proof}

\begin{lemma}[The concentrated current with an explicit constant, given the Monge--Amp\`ere data]
\label{bou:concentrated-current-explicit}
\uses{bou:concentrated-current-integration,bou:mixed-mass-class-invariance,bou:smooth-form-current,bou:trace-measure,bou:positive-currents,bou:smooth-current-closed,bou:closed-currents-limits,bou:closed-set-mass-limit}
\source{bou02-volume:4.12,dp04-kahler-cone:2.6}
\lean{TranscendentalNakamaye.Pending.ConcentratedCurrent.mul_vol_eq_slotIntegral,TranscendentalNakamaye.Pending.ConcentratedCurrent.isKahler_of_half,TranscendentalNakamaye.Pending.ConcentratedCurrent.le_mixedMeasure_concentrated,TranscendentalNakamaye.Pending.ConcentratedCurrent.concentratedSlots,TranscendentalNakamaye.Pending.ConcentratedCurrent.thetaPlurisubharmonic_concentratedSlots,TranscendentalNakamaye.Pending.ConcentratedCurrent.traceMeasure_smoothCurrent_eq,TranscendentalNakamaye.Pending.ConcentratedCurrent.mass_smoothCurrent_eq,TranscendentalNakamaye.Pending.ConcentratedCurrent.le_traceMeasure_smoothCurrent,TranscendentalNakamaye.Pending.ConcentratedCurrent.exists_concentratedCurrent}
\leanok
Let \((M,\omega)\) be a compact K\"ahler manifold of dimension \(N=d+p\). (1) If \(A\) is a smooth closed real \((1,1)\)-form, \(g\), \(f\) chart-smooth and
\((A+dd^cg)^N=c\,(\omega+dd^cf)^N\) pointwise, then \(c\int_M\omega^N=\int_MA^N\). (2) If moreover \(\alpha=A+dd^cg\) is K\"ahler, \(c>0\), \(\omega+dd^cf\ge\frac12\omega\),
\(\int_S(\omega+dd^cf)^p\wedge\omega^d\ge\delta>0\), \(\int_MA^N\ge P_0\) and \(\int_MA^d\wedge\omega^p\le M_0\) with \(M_0>0\), then
\[
 \int_S\alpha^p\wedge\omega^d\ \ge\ \frac{P_0}{\int_M\omega^N}\,\frac{\delta^2}{4\cdot2^d\,N!\,M_0},
\]
and the same bound holds for the trace measure of the positive current \([M]\wedge\alpha^p\), whose mass is \(\int_MA^p\wedge\omega^d\). (3) If such data
\(A_k,g_k,f_k,c_k,S_k\) are given for every \(k\) with the same \(\delta,P_0,M_0\), \(\int_MA_k^p\wedge\omega^d\le M_1\), and \(F_j\) are closed sets, decreasing, with
\(S_k\subset F_j\) for \(j\le k\), then along an ultrafilter finer than the cofinite filter the currents \([M]\wedge\alpha_k^p\) converge weakly to a closed positive current
\(\Theta\) of mass \(\le M_1\) whose trace measure of \(\bigcap_jF_j\) is at least the constant of (2).
\end{lemma}
\begin{proof}
\leanok
(1) Integrate the pointwise identity against \(\omega^N\); top-degree integrals do not change when \(dd^c\) of a smooth function is added to a slot
(Lemma~\ref{bou:mixed-mass-class-invariance}). (2) Lemma~\ref{bou:concentrated-current-integration} gives \(\int_S\alpha^p\wedge\omega^d\ge c\,\delta^2/(4\cdot2^dN!\,M_0)\), and
\(c=\int_MA^N/\int_M\omega^N\ge P_0/\int_M\omega^N\) by (1); the trace measure of \([M]\wedge\alpha^p\) is the global mixed measure of the slots
\((\alpha,\dots,\alpha,\omega,\dots,\omega)\) (Lemmas~\ref{bou:smooth-form-current} and~\ref{bou:trace-measure}), and its mass is a class invariant. (3) The masses are bounded by \(M_1\);
weak compactness of positive currents (Lemma~\ref{bou:positive-currents}), closedness of the limit (Lemmas~\ref{bou:smooth-current-closed} and~\ref{bou:closed-currents-limits}), and
the passage of the lower bound to \(\bigcap_jF_j\) (Lemma~\ref{bou:closed-set-mass-limit}).
\end{proof}
Yau's theorem (Theorem~\ref{bou:yau}) supplies the data of (2) for every K\"ahler class \(\{A\}\) and every \(f\); Lemma~\ref{bou:concentrated-current} below is the case \(P_0=a^N\).

\begin{lemma}[A current in \(a^p\) with mass on \(Y\)]
\label{bou:concentrated-current}
\uses{bou:yau,bou:concentrated-current-explicit,pc:domination-intersection,bou:concentration-pointwise,bou:product-currents,bou:positive-currents,bou:trace-measure,bou:smooth-form-current,fpmma:geometric-density,fpmma:coefficient-laws,bou:closed-currents-limits,bou:smooth-current-closed,bou:mixed-mass-class-invariance,bou:concentrated-current-integration,bou:closed-set-mass-limit}
\source{bou02-volume:4.12,dp04-kahler-cone:2.6}
\notready
Let \((M,\omega)\) be a compact K\"ahler manifold of dimension \(N\), \(Y\subset M\) a closed subset, \(1\le p\le N\), and for \(0<\varepsilon\le1\) let
\(\omega_\varepsilon\) be smooth closed real \((1,1)\)-forms cohomologous to \(\omega\) with \(\omega_\varepsilon\ge\frac12\omega\) and \(V_\varepsilon\) open sets, decreasing as
\(\varepsilon\downarrow0\), with \(\bigcap_\varepsilon\overline{V_\varepsilon}\subseteq Y\). Let \(K\subset M\) be compact with
\(\int_{K\cap V_\varepsilon}\omega_\varepsilon^p\wedge\omega^{N-p}\ge\delta>0\) for every \(\varepsilon\). Then for every nef class \(a\) with \(a^N>0\) there is a closed
positive \((p,p)\)-current \(T\) in \(a^p\) with
\[
 \int_{K\cap Y}T\wedge\omega^{N-p}\ \ge\ \frac{1}{2^{N+1}N!}\,\frac{a^N}{\int_M\omega^N}\,\frac{\delta^2}{(a+\{\omega\})^{N-p}\cdot\{\omega\}^p}.
\]
\end{lemma}
\begin{proof}
For \(0<\varepsilon\le1\) the class \(a+\varepsilon\{\omega\}\) is K\"ahler, and \(\omega_\varepsilon^N\) is a smooth volume form of total mass \(\int_M\omega^N\). By
Theorem~\ref{bou:yau} there is a K\"ahler form \(\alpha_\varepsilon\) in \(a+\varepsilon\{\omega\}\) with \(\alpha_\varepsilon^N=C_\varepsilon\,\omega_\varepsilon^N\), where
\(C_\varepsilon=(a+\varepsilon\{\omega\})^N/\int_M\omega^N\ge a^N/\int_M\omega^N\) (expand; products of nef classes are \(\ge0\), Lemma~\ref{pc:domination-intersection}(2)). Put
\(m=(a+\{\omega\})^{N-p}\cdot\{\omega\}^p\) and \(s=2^{N-p+1}m/\delta\).

At a point, let \(A\), \(B\), \(C\) be the Hermitian matrices of \(\alpha_\varepsilon\), \(\omega_\varepsilon\), \(\omega\) in a chart. The density of a top-degree product of
real \((1,1)\)-forms is \(\pi^{-N}\operatorname{Re}C(H_1,\dots,H_N)\) of their matrices (Lemma~\ref{fpmma:geometric-density}), and \(C(H[N])=N!\det H\)
(Lemma~\ref{fpmma:coefficient-laws}); so \(\alpha_\varepsilon^N=C_\varepsilon\omega_\varepsilon^N\) reads \(\det A=C_\varepsilon\det B\). Let
\(E=\{\alpha_\varepsilon^{N-p}\wedge\omega_\varepsilon^p\ge s\,\omega_\varepsilon^N\}\), a closed set. At a point off \(E\),
\(\operatorname{Re}C(A[N-p],B[p])<s\,N!\det B\), and Lemma~\ref{bou:concentration-pointwise} with \(t=sN!\) gives
\(\alpha_\varepsilon^p\wedge\omega^{N-p}\ge(C_\varepsilon/sN!)\,\omega_\varepsilon^p\wedge\omega^{N-p}\) there. On \(E\),
\(s\int_E\omega_\varepsilon^N\le\int_M\alpha_\varepsilon^{N-p}\wedge\omega_\varepsilon^p=(a+\varepsilon\{\omega\})^{N-p}\cdot\{\omega\}^p\le m\)
(Lemma~\ref{pc:domination-intersection}(2),(3)); as \(\omega\le2\omega_\varepsilon\), \(\omega_\varepsilon^p\wedge\omega^{N-p}\le2^{N-p}\omega_\varepsilon^N\), so
\(\int_E\omega_\varepsilon^p\wedge\omega^{N-p}\le2^{N-p}m/s=\delta/2\). As \(\alpha_\varepsilon^p\wedge\omega^{N-p}\ge0\),
\[
 \int_{K\cap V_\varepsilon}\alpha_\varepsilon^p\wedge\omega^{N-p}\ \ge\ \frac{C_\varepsilon}{sN!}\int_{(K\cap V_\varepsilon)\setminus E}\omega_\varepsilon^p\wedge\omega^{N-p}
 \ \ge\ \frac{C_\varepsilon}{sN!}\,\frac\delta2\ =\ \frac{C_\varepsilon\,\delta^2}{2^{N-p+2}N!\,m}
 \ \ge\ \frac{1}{2^{N+1}N!}\,\frac{a^N}{\int_M\omega^N}\,\frac{\delta^2}{m},
\]
using \(p\ge1\). (This integration step is Lemma~\ref{bou:concentrated-current-integration}, with \(c=C_\varepsilon\) and \(S=K\cap V_\varepsilon\).)

By Lemma~\ref{bou:smooth-form-current} the forms \(\alpha_\varepsilon^p\) define positive currents \(T_\varepsilon\) of bidimension \((N-p,N-p)\) whose trace measures are
\(\alpha_\varepsilon^p\wedge\omega^{N-p}\), of mass \((a+\varepsilon\{\omega\})^p\cdot\{\omega\}^{N-p}\le(a+\{\omega\})^p\cdot\{\omega\}^{N-p}\) (such total masses depend only on
the classes, Lemma~\ref{bou:mixed-mass-class-invariance}). By
Lemma~\ref{bou:positive-currents}(2), along \(\varepsilon\to0\) they have a weak limit \(T\), positive; it is closed, because the \(T_\varepsilon\) are closed
(Lemma~\ref{bou:smooth-current-closed}) and weak limits of closed currents are closed (Lemma~\ref{bou:closed-currents-limits}(3)), and it lies in
\(a^p\) because the classes \((a+\varepsilon\{\omega\})^p\) of the \(T_\varepsilon\) tend to \(a^p\). For fixed \(\varepsilon_0\) and \(\varepsilon\le\varepsilon_0\),
\(K\cap V_\varepsilon\subseteq F_{\varepsilon_0}=K\cap\overline{V_{\varepsilon_0}}\), a closed set, so \(\sigma_{T_\varepsilon}(F_{\varepsilon_0})\) is at least the bound above, and
by Lemma~\ref{bou:trace-measure} so is \(\sigma_T(F_{\varepsilon_0})=\int_{F_{\varepsilon_0}}T\wedge\omega^{N-p}\). The sets \(F_{\varepsilon_0}\) decrease as \(\varepsilon_0\downarrow0\)
to \(K\cap\bigcap_\varepsilon\overline{V_\varepsilon}\subseteq K\cap Y\), and \(\sigma_T\) is finite, so the bound holds for \(\int_{K\cap Y}T\wedge\omega^{N-p}\). (These last two steps are Lemma~\ref{bou:closed-set-mass-limit}, with \(\varepsilon_0=1/k\).)
\end{proof}
This is the argument of Demailly--P\u{a}un and Boucksom with the eigenvalue inequalities replaced by Lemma~\ref{bou:concentration-pointwise}: their set
\(\{\lambda_{p+1}\cdots\lambda_N\ \text{large}\}\) becomes \(E\), on which the mixed coefficient \(C(A[N-p],B[p])\) is large, and the constants come out the same.

\emph{Superseded.} As stated, without connectedness of \(M\), the lemma is false. Let \(M=M_1\sqcup M_2\) be the disjoint union of two compact K\"ahler manifolds
of dimension \(N\), \(p=N\), \(a\) the class equal to \(0\) on \(M_1\) and to \(\{\omega\}\) on \(M_2\) (nef, \(a^N=\int_{M_2}\omega^N>0\)), \(Y=\{y_0\}\) a point of \(M_1\), \(K=M_1\),
and \(\omega_\varepsilon\), \(V_\varepsilon\) on \(M_1\) the forms and sets of Demailly--P\u{a}un's concentration of mass at \(y_0\) (and \(\omega_\varepsilon=\omega\) on \(M_2\)). A
closed positive current \(T\) in \(a^N\) has mass \(\int_{M_1}a^N=0\) on \(M_1\), so \(\int_{K\cap Y}T=0\), below the positive bound. The proof uses Yau's theorem
(Theorem~\ref{bou:yau}), which needs \(M\) connected. Moreover the library has no cohomology class of a \((p,p)\)-current, so ``\(T\) in \(a^p\)'' has no Lean
meaning yet. Lemma~\ref{bou:concentrated-current-connected} is the corrected statement, with \(M\) connected and the class of \(T\) stated through its intersection
numbers, and it is proved; the statement above is kept unchanged.

\begin{lemma}[A current with the intersection numbers of \(a^p\) and mass on \(Y\), corrected form]
\label{bou:concentrated-current-connected}
\uses{bou:yau,bou:yau-consequences,bou:concentrated-current-explicit,bou:nef-mixed-mass,bou:nef-core,bou:mixed-mass-class-invariance,bou:smooth-form-current,bou:diagonal-tangential-reduction,bou:product-currents,bou:positive-currents,bou:trace-measure,bou:closed-currents-limits}
\source{bou02-volume:4.12,dp04-kahler-cone:2.6}
\lean{TranscendentalNakamaye.Pending.ConcentratedCurrentNef.slotIntegral_add_smul_mono,TranscendentalNakamaye.Pending.ConcentratedCurrentNef.slotIntegral_const_form_eq,TranscendentalNakamaye.Pending.ConcentratedCurrentNef.nefWithRespectTo_zero,TranscendentalNakamaye.Pending.ConcentratedCurrentNef.nefWithRespectTo_form,TranscendentalNakamaye.Pending.ConcentratedCurrentNef.kahlerVolume_real_univ_pos,TranscendentalNakamaye.Pending.ConcentratedCurrentNef.prodTest,TranscendentalNakamaye.Pending.ConcentratedCurrentNef.smoothCurrent_prodTest,TranscendentalNakamaye.Pending.ConcentratedCurrentNef.slotIntegral_append_add_ddc,TranscendentalNakamaye.Pending.ConcentratedCurrentNef.tendsto_slotIntegral_append,TranscendentalNakamaye.Pending.ConcentratedCurrentNef.thetaPsh_concSlots,TranscendentalNakamaye.Pending.ConcentratedCurrentNef.HasClassIntersections,TranscendentalNakamaye.Pending.ConcentratedCurrentNef.exists_concentratedCurrent_of_nef}
\leanok
Let \((M,\omega)\) be a compact connected K\"ahler manifold of dimension \(N\), \(Y\subset M\) a subset, \(1\le p\le N\), and for \(0<\varepsilon\le1\) let
\(\omega_\varepsilon=\omega+dd^cf_\varepsilon\), \(f_\varepsilon\) smooth, with \(\omega_\varepsilon\ge\frac12\omega\), and \(V_\varepsilon\subset M\) sets, decreasing as \(\varepsilon\downarrow0\), with
\(\bigcap_\varepsilon\overline{V_\varepsilon}\subseteq Y\). Let \(K\subset M\) be compact with \(\int_{K\cap V_\varepsilon}\omega_\varepsilon^p\wedge\omega^{N-p}\ge\delta>0\) for every \(\varepsilon\).
Then for every nef class \(a=\{A\}\) there is a closed positive current \(T\) of bidimension \((N-p,N-p)\) (Definition~\ref{bou:product-currents}) with the
intersection numbers of \(a^p\),
\[
 T(\beta_1\wedge\dots\wedge\beta_{N-p})=\int_M\beta_1\wedge\dots\wedge\beta_{N-p}\wedge A^p\quad\text{for all closed real \((1,1)\)-forms }\beta_i,
\]
and
\[
 \int_{K\cap Y}T\wedge\omega^{N-p}\ \ge\ \frac{1}{2^{N+1}N!}\,\frac{a^N}{\int_M\omega^N}\,\frac{\delta^2}{(a+\{\omega\})^{N-p}\cdot\{\omega\}^p}.
\]
\end{lemma}
\begin{proof}
\leanok
Put \(d=N-p\) and \(\varepsilon_k=1/(k+1)\). Mixed integrals of nef classes are \(\ge0\) (Lemma~\ref{bou:nef-mixed-mass}); expanding a mixed integral of the classes
\(\theta_i+t\beta_i\), with \(\theta_i\), \(\beta_i\) nef, as a polynomial in \(t\) with nonnegative coefficients shows that it increases with \(t\ge0\). Hence
\(a^N\le(a+\varepsilon\{\omega\})^N\), \((a+\varepsilon\{\omega\})^d\cdot\{\omega\}^p\le m:=(a+\{\omega\})^d\cdot\{\omega\}^p\) and \((a+\varepsilon\{\omega\})^p\cdot\{\omega\}^d\le
(a+\{\omega\})^p\cdot\{\omega\}^d\) for \(0<\varepsilon\le1\); and \(m>0\), since one term of the expansion of \(m\) at \(t=1\) is \(\int_M\omega^N>0\).

As \(a\) is nef there is a smooth \(h_k\) with \(A+dd^ch_k+\frac{\varepsilon_k}2\omega\ge0\), so \(B_k=A+\varepsilon_k\omega+dd^ch_k\) is K\"ahler; \(\omega_{\varepsilon_k}\) is
K\"ahler. By Yau's theorem in the form of Corollary~\ref{bou:yau-consequences} (here \(M\) connected is used) there is \(\varphi_k\) with \(\alpha_k=B_k+dd^c\varphi_k\) K\"ahler and
\(\alpha_k^N=c_k\,\omega_{\varepsilon_k}^N\), \(c_k>0\). Lemma~\ref{bou:concentrated-current-explicit}(3), applied with \(S_k=K\cap V_{\varepsilon_k}\), the closed sets
\(F_j=K\cap\overline{V_{\varepsilon_j}}\), \(P_0=a^N\) (mixed integrals are class invariants, Lemma~\ref{bou:mixed-mass-class-invariance}, so \(\int B_k^N=(a+\varepsilon_k\{\omega\})^N\ge a^N\)), \(M_0=m\) and
\(M_1=(a+\{\omega\})^p\cdot\{\omega\}^d\), gives along an ultrafilter finer than the cofinite filter a weak limit \(T\) of the currents \([M]\wedge\alpha_k^p\), closed and positive,
with \(\int_{\bigcap_jF_j}T\wedge\omega^d\ge\frac{a^N}{\int_M\omega^N}\frac{\delta^2}{4\cdot2^dN!\,m}\). Since \(\bigcap_jF_j\subseteq K\cap\bigcap_\varepsilon\overline{V_\varepsilon}\subseteq K\cap Y\)
(every \(\varepsilon\in(0,1]\) is \(\ge\) some \(\varepsilon_j\)) and \(4\cdot2^d=2^{d+2}\le2^{N+1}\) (\(p\ge1\)), the bound follows.

For closed real \((1,1)\)-forms \(\beta_1,\dots,\beta_d\), the test form \(\beta_1\wedge\dots\wedge\beta_d\) (Lemma~\ref{bou:diagonal-tangential-reduction}) is sent by \([M]\wedge\alpha_k^p\) to
\(\int_M\beta\wedge\alpha_k^p\) (Lemma~\ref{bou:smooth-form-current}), which equals \(\int_M\beta\wedge(A+\varepsilon_k\omega)^p\) by class invariance (Lemma~\ref{bou:nef-mixed-mass}, applied in the
last \(p\) slots, as in Lemma~\ref{bou:nef-core}); this tends to \(\int_M\beta\wedge A^p\) as \(k\to\infty\), the mixed integral being a polynomial in \(\varepsilon_k\). So
\(T(\beta_1\wedge\dots\wedge\beta_d)=\int_M\beta\wedge A^p\).
\end{proof}
This is the case of Lemma~\ref{bou:concentrated-current} that Boucksom's argument uses: there \(M=X\times X\) is connected, and the push-forward
\(p_{1*}(T\wedge p_2^*\omega)\) to \(X\) pairs \(T\) only with products of closed \((1,1)\)-forms. The openness of \(V_\varepsilon\) and the closedness of \(Y\) are not needed.

\begin{lemma}[Products of complex manifolds, the product K\"ahler form and the diagonal]
\label{bou:product-manifold}
\uses{ana:geometric-setting,ana:blowups,posclass:background}
\source{dp04-kahler-cone:2.12,bou02-volume:4.8}
\lean{TranscendentalNakamaye.Pending.ProductComplexManifold.chartedSpace,TranscendentalNakamaye.Pending.ProductComplexManifold.contDiffOn_transition,TranscendentalNakamaye.Pending.ProductComplexManifold.contMDiff_fst,TranscendentalNakamaye.Pending.ProductComplexManifold.contMDiff_snd,TranscendentalNakamaye.Pending.ProductComplexManifold.contMDiff_prodMk,TranscendentalNakamaye.Pending.ProductComplexManifold.contMDiff_prodMap,TranscendentalNakamaye.Pending.ProductComplexManifold.contMDiff_diag,TranscendentalNakamaye.Pending.ProductComplexManifold.prodForm,TranscendentalNakamaye.Pending.ProductComplexManifold.prodForm_isKahler,TranscendentalNakamaye.Pending.ProductComplexManifold.prodKahlerBackground,TranscendentalNakamaye.Pending.ProductComplexManifold.regularAtDim_diagonal,TranscendentalNakamaye.Pending.ProductComplexManifold.isClosedConnectedSubmanifold_diagonal}
\leanok
Let \(X\) and \(Y\) be complex manifolds of dimensions \(n\) and \(m\). With the products of charts, followed by the linear identification
\(\mathbf C^n\times\mathbf C^m\cong\mathbf C^{n+m}\), \(X\times Y\) is a complex manifold of dimension \(n+m\); the projections, the pair of two
holomorphic maps into \(X\) and \(Y\), the product of two holomorphic maps and the diagonal map \(X\to X\times X\) are holomorphic. If \(a\) and \(b\) are
K\"ahler forms on \(X\) and \(Y\), then \(p_1^*a+p_2^*b\) is a K\"ahler form on \(X\times Y\). If \(X\) is nonempty and connected, the diagonal
\(\Delta\subset X\times X\) is a closed connected complex submanifold of dimension \(n\).
\end{lemma}
\begin{proof}
\leanok
The transition maps of product charts are products of transition maps, hence holomorphic, and the projections are linear in product charts.
\(p_1^*a+p_2^*b\) is closed, real and of type \((1,1)\), and positive: on a tangent vector \(v\) it is \(a(dp_1v,J\,dp_1v)+b(dp_2v,J\,dp_2v)\ge0\), and a vector
killed by both projections is \(0\). In the product chart at \((p,p)\), the diagonal is \(\{(u,u)\}\), which the linear change
\((u,u')\mapsto(u-u_0,u'-u)\) flattens to the coordinate plane \(\mathbf C^n\times0\); it is closed as \(X\) is Hausdorff, and connected as the image of \(X\).
\end{proof}

\begin{lemma}[Mixed measures on a product: Fubini, and the top power of \(p_1^*a+p_2^*b\)]
\label{bou:product-mixed-measure}
\uses{bou:product-manifold,bou:smooth-mixed-density,gnpp:global-measure,fpmma:geometric-density,fpmma:coefficient-laws}
\source{bou02-volume:4.8,dp04-kahler-cone:2.12}
\lean{TranscendentalNakamaye.Pending.ProductMixedMeasure.mixedCoeff,TranscendentalNakamaye.Pending.ProductMixedMeasure.mixedCoefficient_eq_mixedCoeff,TranscendentalNakamaye.Pending.ProductMixedMeasure.blockFamily,TranscendentalNakamaye.Pending.ProductMixedMeasure.blockFamily_inl_inr,TranscendentalNakamaye.Pending.ProductMixedMeasure.blockFamily_inr_inl,TranscendentalNakamaye.Pending.ProductMixedMeasure.mixedCoeff_blockFamily,TranscendentalNakamaye.Pending.ProductMixedMeasure.mixedCoefficientImMultilinear,TranscendentalNakamaye.Pending.ProductMixedMeasure.mixedCoefficient_im_eq_zero,TranscendentalNakamaye.Pending.ProductMixedMeasure.projFst,TranscendentalNakamaye.Pending.ProductMixedMeasure.projSnd,TranscendentalNakamaye.Pending.ProductMixedMeasure.projFst_apply,TranscendentalNakamaye.Pending.ProductMixedMeasure.projSnd_apply,TranscendentalNakamaye.Pending.ProductMixedMeasure.projFst_smul,TranscendentalNakamaye.Pending.ProductMixedMeasure.projSnd_smul,TranscendentalNakamaye.Pending.ProductMixedMeasure.prodModelEquiv_symm_coordinateVector_castAdd,TranscendentalNakamaye.Pending.ProductMixedMeasure.prodModelEquiv_symm_coordinateVector_natAdd,TranscendentalNakamaye.Pending.ProductMixedMeasure.comp_vec,TranscendentalNakamaye.Pending.ProductMixedMeasure.apply_vec_zero_right,TranscendentalNakamaye.Pending.ProductMixedMeasure.apply_vec_zero_left,TranscendentalNakamaye.Pending.ProductMixedMeasure.projFst_castAdd,TranscendentalNakamaye.Pending.ProductMixedMeasure.projFst_natAdd,TranscendentalNakamaye.Pending.ProductMixedMeasure.projSnd_castAdd,TranscendentalNakamaye.Pending.ProductMixedMeasure.projSnd_natAdd,TranscendentalNakamaye.Pending.ProductMixedMeasure.normalizedHermitianMatrix_comp_projFst,TranscendentalNakamaye.Pending.ProductMixedMeasure.normalizedHermitianMatrix_comp_projSnd,TranscendentalNakamaye.Pending.ProductMixedMeasure.mixedCoeff_eq_mixedCoefficient,TranscendentalNakamaye.Pending.ProductMixedMeasure.isOneOne_comp_projFst,TranscendentalNakamaye.Pending.ProductMixedMeasure.isOneOne_comp_projSnd,TranscendentalNakamaye.Pending.ProductMixedMeasure.actualMixedTopDensity_append_comp,TranscendentalNakamaye.Pending.ProductMixedMeasure.actualMixedTopDensity_const_add_comp,TranscendentalNakamaye.Pending.ProductMixedMeasure.prodComplexChart,TranscendentalNakamaye.Pending.ProductMixedMeasure.prodComplexChart_val,TranscendentalNakamaye.Pending.ProductMixedMeasure.mem_prodComplexChart_target,TranscendentalNakamaye.Pending.ProductMixedMeasure.prodComplexChart_symm_apply,TranscendentalNakamaye.Pending.ProductMixedMeasure.coordinateMap_fst_eventuallyEq,TranscendentalNakamaye.Pending.ProductMixedMeasure.coordinateMap_snd_eventuallyEq,TranscendentalNakamaye.Pending.ProductMixedMeasure.pullback_fst_coeff,TranscendentalNakamaye.Pending.ProductMixedMeasure.pullback_snd_coeff,TranscendentalNakamaye.Pending.ProductMixedMeasure.measurePreserving_prodModelEquiv,TranscendentalNakamaye.Pending.ProductMixedMeasure.isFiniteMeasure_globalMixedMeasure,TranscendentalNakamaye.Pending.ProductMixedMeasure.exists_chart_partition,TranscendentalNakamaye.Pending.ProductMixedMeasure.measurableSet_chart_image,TranscendentalNakamaye.Pending.ProductMixedMeasure.globalMixedMeasure_rect,TranscendentalNakamaye.Pending.ProductMixedMeasure.measure_eq_tsum_inter,TranscendentalNakamaye.Pending.ProductMixedMeasure.measure_prod_eq_of_chart,TranscendentalNakamaye.Pending.ProductMixedMeasure.prodSlots,TranscendentalNakamaye.Pending.ProductMixedMeasure.thetaPlurisubharmonic_prodSlots,TranscendentalNakamaye.Pending.ProductMixedMeasure.chartDensity_prodSlots,TranscendentalNakamaye.Pending.ProductMixedMeasure.globalMixedMeasure_prodSlots,TranscendentalNakamaye.Pending.ProductMixedMeasure.globalMixedMeasure_prodSlots_univ,TranscendentalNakamaye.Pending.ProductMixedMeasure.chartDensity_prodForm,TranscendentalNakamaye.Pending.ProductMixedMeasure.globalMixedMeasure_prodForm,TranscendentalNakamaye.Pending.ProductMixedMeasure.globalMixedMeasure_prodForm_univ}
\leanok
Let \(X\) and \(Y\) be compact complex manifolds of dimensions \(n\) and \(m\), \(p_1,p_2\) the projections of \(X\times Y\) (Lemma~\ref{bou:product-manifold}),
and write \(\theta_1\wedge\dots\wedge\theta_N\) for the global mixed measure of smooth closed semipositive real \((1,1)\)-forms \(\theta_i\) with the
potentials \(0\) (Definition~\ref{gnpp:global-measure}).
\begin{enumerate}
\item For smooth closed semipositive \(A_1,\dots,A_n\) on \(X\) and \(B_1,\dots,B_m\) on \(Y\),
      \(p_1^*A_1\wedge\dots\wedge p_1^*A_n\wedge p_2^*B_1\wedge\dots\wedge p_2^*B_m=(A_1\wedge\dots\wedge A_n)\otimes(B_1\wedge\dots\wedge B_m)\) as measures on
      \(X\times Y\); in particular the total masses multiply.
\item For smooth closed semipositive \(a\) on \(X\) and \(b\) on \(Y\), \((p_1^*a+p_2^*b)^{n+m}=\binom{n+m}n\,a^n\otimes b^m\) as measures on \(X\times Y\); in
      particular \(\int_{X\times Y}(p_1^*a+p_2^*b)^{n+m}=\binom{n+m}n\int_Xa^n\int_Yb^m\), and for \(Y=X\), \(b=a\) this is \(\binom{2n}n\bigl(\int_Xa^n\bigr)^2\).
\end{enumerate}
\end{lemma}
\begin{proof}
\leanok
In the product of two charts followed by \(\mathbf C^n\times\mathbf C^m\cong\mathbf C^{n+m}\), the Hermitian coefficient matrix of \(p_1^*\alpha\) is
\(H(\alpha)\oplus0\) and that of \(p_2^*\beta\) is \(0\oplus H(\beta)\). The mixed coefficient \(C(H_1,\dots,H_N)=\sum_\sigma\det\bigl(H_{\sigma(j)}\bigr)_{j\text{-th row}}\)
of block-diagonal slots factors, \(C(H_1\oplus0,\dots,H_n\oplus0,0\oplus K_1,\dots,0\oplus K_m)=C(H_1,\dots,H_n)\,C(K_1,\dots,K_m)\): each matrix of the sum is
block diagonal, and its first block has a zero row unless \(\sigma\) preserves the two blocks (the mixed analogue of
\(\det\bigl(\begin{smallmatrix}A&0\\0&B\end{smallmatrix}\bigr)=\det A\det B\)); the mixed coefficient of real \((1,1)\)-forms is real (polarization). For (2),
\(C(H[N])=N!\det H\) and \(\det(H(a)\oplus H(b))=\det H(a)\det H(b)\), so the density of \((p_1^*a+p_2^*b)^{n+m}\) is \(\frac{(n+m)!}{n!\,m!}\) times the product
of the densities of \(a^n\) and \(b^m\). By Lemma~\ref{bou:smooth-mixed-density} the measures have these densities in charts; the identification
\(\mathbf C^n\times\mathbf C^m\cong\mathbf C^{n+m}\) preserves Lebesgue measure, so Tonelli gives both identities on measurable rectangles \(S\times T\) with
\(S\), \(T\) inside chart domains. Finitely many charts cover the compact factors; a measurable partition of each factor subordinate to them reduces an
arbitrary measurable rectangle to these, and finite measures that agree on rectangles are equal (uniqueness of the product measure).
\end{proof}

\begin{lemma}[Restriction of a positive current to a Borel set]
\label{bou:current-restriction}
\uses{bou:product-currents,bou:positive-currents,bou:trace-measure}
\source{demailly-cadg:III.1.14,demailly-cadg:III.2.4}
\lean{TranscendentalNakamaye.Pending.CurrentRestriction.zeroFun,TranscendentalNakamaye.Pending.CurrentRestriction.chartSmooth_zeroFun,TranscendentalNakamaye.Pending.CurrentRestriction.functionTestLM,TranscendentalNakamaye.Pending.CurrentRestriction.testCurrent,TranscendentalNakamaye.Pending.CurrentRestriction.isPositive_testCurrent,TranscendentalNakamaye.Pending.CurrentRestriction.testMeasure,TranscendentalNakamaye.Pending.CurrentRestriction.integral_testMeasure,TranscendentalNakamaye.Pending.CurrentRestriction.testMeasure_real_univ,TranscendentalNakamaye.Pending.CurrentRestriction.testMeasure_kahlerTest,TranscendentalNakamaye.Pending.CurrentRestriction.testMeasure_add,TranscendentalNakamaye.Pending.CurrentRestriction.testMeasure_smul,TranscendentalNakamaye.Pending.CurrentRestriction.testMeasure_le_smul_traceMeasure,TranscendentalNakamaye.Pending.CurrentRestriction.testMeasure_smulTest,TranscendentalNakamaye.Pending.CurrentRestriction.exists_const,TranscendentalNakamaye.Pending.CurrentRestriction.restrictVal,TranscendentalNakamaye.Pending.CurrentRestriction.restrictVal_eq,TranscendentalNakamaye.Pending.CurrentRestriction.restrictCurrent,TranscendentalNakamaye.Pending.CurrentRestriction.restrictCurrent_apply,TranscendentalNakamaye.Pending.CurrentRestriction.restrictCurrent_apply_of_mem,TranscendentalNakamaye.Pending.CurrentRestriction.isPositive_restrictCurrent,TranscendentalNakamaye.Pending.CurrentRestriction.mass_restrictCurrent,TranscendentalNakamaye.Pending.CurrentRestriction.restrictCurrent_add_compl,TranscendentalNakamaye.Pending.CurrentRestriction.restrictCurrent_univ,TranscendentalNakamaye.Pending.CurrentRestriction.restrictCurrent_eq_zero,TranscendentalNakamaye.Pending.CurrentRestriction.restrictCurrent_smulTest,TranscendentalNakamaye.Pending.CurrentRestriction.traceMeasure_restrictCurrent}
\leanok
Let \(M\) be a compact complex manifold with a K\"ahler form \(\omega\) and \(T\) a positive current of bidimension \((q,q)\) with trace measure \(\sigma_T\).
\begin{enumerate}
\item For every positive test form \(\psi\) there is a finite Borel measure \(T\wedge\psi\) with \(\int_Mh\,d(T\wedge\psi)=T(h\psi)\) for every smooth \(h\); it is
      additive and positively homogeneous in \(\psi\), \(T\wedge\omega^q=\sigma_T\), \(T\wedge\psi\le c\,\sigma_T\) when \(c\,\omega^q-\psi\) is positive, and
      \(T\wedge(g\psi)=g\,(T\wedge\psi)\) for smooth \(g\ge0\).
\item For a Borel set \(E\) and a test form \(\varphi\), the number \(\mathbf 1_ET(\varphi)=(T\wedge(c\,\omega^q+\varphi))(E)-c\,\sigma_T(E)\) does not depend on the
      constant \(c\ge0\) with \(c\,\omega^q\pm\varphi\) positive. \(\mathbf 1_ET\) is a positive current of mass \(\sigma_T(E)\), \(\mathbf 1_ET+\mathbf 1_{M\setminus E}T=T\),
      \(\mathbf 1_MT=T\), \(\mathbf 1_ET=0\) if \(\sigma_T(E)=0\), \(\mathbf 1_ET(h\psi)=\int_Eh\,d(T\wedge\psi)\) for positive \(\psi\) and smooth \(h\), and the trace
      measure of \(\mathbf 1_ET\) is \(\sigma_T\) restricted to \(E\).
\end{enumerate}
\end{lemma}
\begin{proof}
\leanok
(1) For positive \(\psi\), \(h\mapsto T(h\psi)\) is a positive current of bidimension \((0,0)\) (\(h\psi\) is positive for \(h\ge0\)); its trace measure
(Lemma~\ref{bou:trace-measure}) is \(T\wedge\psi\). Additivity and homogeneity follow from uniqueness of the measure with given integrals of smooth functions; for
\(c\,\omega^q-\psi\) positive, \(c\,\sigma_T-T\wedge\psi=T\wedge(c\,\omega^q-\psi)\ge0\); and both sides of \(T\wedge(g\psi)=g\,(T\wedge\psi)\) integrate \(h\) to
\(T(hg\psi)\).
(2) By Lemma~\ref{bou:positive-currents}(1) such \(c\) exist, and for two of them, \(c\le c'\), additivity gives
\(T\wedge(c'\omega^q+\varphi)=T\wedge(c\,\omega^q+\varphi)+(c'-c)\sigma_T\), so the value is independent of \(c\); the same additivity makes \(\mathbf 1_ET\) linear. For
positive \(\varphi\) take \(c=0\): \(\mathbf 1_ET(\varphi)=(T\wedge\varphi)(E)\ge0\); with \(\varphi=\omega^q\) this is \(\sigma_T(E)\). Additivity of measures over \(E\) and
\(M\setminus E\) gives the decomposition, and \((T\wedge\varphi')(M)=T(\varphi')\) gives \(\mathbf 1_MT=T\). If \(\sigma_T(E)=0\) then \((T\wedge\psi)(E)=0\) for every
positive \(\psi\le c\,\omega^q\), so \(\mathbf 1_ET=0\). The formula for \(\mathbf 1_ET(h\psi)\) follows from (1) by writing \(h\) as a difference of two
nonnegative smooth functions (\(h=(h+C)-C\)), and with \(\psi=\omega^q\) it identifies the trace measure of \(\mathbf 1_ET\) with \(\sigma_T|_E\).
\end{proof}

\begin{lemma}[Slot test forms: the Leibniz rule and integration by parts against a closed positive current]
\label{bou:slot-ibp}
\uses{bou:closed-currents,bou:closed-currents-local,bou:positive-currents,bou:product-currents,fpmma:hessian-ibp}
\source{demailly-cadg:III.2.3}
\lean{TranscendentalNakamaye.Pending.SlotTestForms.I_smul_I_smul,TranscendentalNakamaye.Pending.SlotTestForms.pi11_apply_line,TranscendentalNakamaye.Pending.SlotTestForms.isSemipositiveOneOne_pi11,TranscendentalNakamaye.Pending.SlotTestForms.pi11_add,TranscendentalNakamaye.Pending.SlotTestForms.pi11_smul,TranscendentalNakamaye.Pending.SlotTestForms.pi11_sub,TranscendentalNakamaye.Pending.SlotTestForms.pi11_sum,TranscendentalNakamaye.Pending.SlotTestForms.pi11_zero,TranscendentalNakamaye.Pending.SlotTestForms.wedgeProd_cons_zero,TranscendentalNakamaye.Pending.SlotTestForms.pi11Form_add,TranscendentalNakamaye.Pending.SlotTestForms.pi11Form_smul,TranscendentalNakamaye.Pending.SlotTestForms.pi11Form_smulFun,TranscendentalNakamaye.Pending.SlotTestForms.slotForm,TranscendentalNakamaye.Pending.SlotTestForms.closedTestForm_eq_slotForm,TranscendentalNakamaye.Pending.SlotTestForms.slotForm_add,TranscendentalNakamaye.Pending.SlotTestForms.slotForm_smul,TranscendentalNakamaye.Pending.SlotTestForms.slotLM,TranscendentalNakamaye.Pending.SlotTestForms.slotForm_sub,TranscendentalNakamaye.Pending.SlotTestForms.slotForm_sum,TranscendentalNakamaye.Pending.SlotTestForms.slotForm_smulFun,TranscendentalNakamaye.Pending.SlotTestForms.slotForm_coeff,TranscendentalNakamaye.Pending.SlotTestForms.slotForm_mem_positiveCone,TranscendentalNakamaye.Pending.SlotTestForms.apply_slotForm_congr,TranscendentalNakamaye.Pending.SlotTestForms.apply_slotForm_congr',TranscendentalNakamaye.Pending.SlotTestForms.apply_slotForm_nonneg,TranscendentalNakamaye.Pending.SlotTestForms.exactPair_mul_coeff,TranscendentalNakamaye.Pending.SlotTestForms.exactPair_coeff_eq_deriv_smul,TranscendentalNakamaye.Pending.SlotTestForms.apply_smulTest_closedTestForm,TranscendentalNakamaye.Pending.SlotTestForms.coordFun,TranscendentalNakamaye.Pending.SlotTestForms.sum_coordFun_smul,TranscendentalNakamaye.Pending.SlotTestForms.ibpCoeff,TranscendentalNakamaye.Pending.SlotTestForms.ibpForm,TranscendentalNakamaye.Pending.SlotTestForms.sum_ibpCoeff_mul,TranscendentalNakamaye.Pending.SlotTestForms.hasFDerivAt_fderiv_apply,TranscendentalNakamaye.Pending.SlotTestForms.hasFDerivAt_ibpForm,TranscendentalNakamaye.Pending.SlotTestForms.hasFDerivAt_ibpCoeff,TranscendentalNakamaye.Pending.SlotTestForms.sum_fderiv_ibpCoeff,TranscendentalNakamaye.Pending.SlotTestForms.ibp_identity,TranscendentalNakamaye.Pending.SlotTestForms.coeff_eq_of_chart,TranscendentalNakamaye.Pending.SlotTestForms.chartExtend_symm,TranscendentalNakamaye.Pending.SlotTestForms.chartExtend_eventuallyEq,TranscendentalNakamaye.Pending.SlotTestForms.tsupport_chartExtend_subset,TranscendentalNakamaye.Pending.SlotTestForms.image_symm_subset_source,TranscendentalNakamaye.Pending.SlotTestForms.ddcForm_coeff_eq_zero,TranscendentalNakamaye.Pending.SlotTestForms.slotForm_ibp}
\leanok
Let \(M\) be a complex manifold, \(\beta_1,\dots,\beta_k\) smooth closed real \((1,1)\)-forms, and for a smooth real two-form \(\gamma\) put
\(\sigma(\gamma)=\Pi_{1,1}\gamma\wedge\beta_1\wedge\dots\wedge\beta_k\), a test form of bidegree \((k+1,k+1)\).
\begin{enumerate}
\item \(\sigma\) is linear, \(\sigma(h\gamma)=h\,\sigma(\gamma)\) for smooth \(h\), \(\sigma(d(f\,dg))\) is the closed test form of Definition~\ref{bou:closed-currents},
      and if the \(\beta_j\) are semipositive and \(\gamma(v,Jv)\ge0\) for all \(v\), then \(\sigma(\gamma)\) is positive. For a positive current \(T\) of
      bidimension \((k+1,k+1)\), \(T(\sigma(\gamma))\) depends only on the values of \(\Pi_{1,1}\gamma\) at points.
\item (Leibniz) If \(T\) is closed and positive, then for smooth \(U\), \(f\), \(g\):
      \(T(U\,\sigma(d(f\,dg)))=-T(f\,\sigma(dU\wedge dg))\); and \(d(\theta(v)\,dg)=\theta'(v)\,d(v\,dg)\) for smooth \(\theta\colon\mathbf R\to\mathbf R\).
\item (Integration by parts) If \(T\) is closed and positive, \(h\) is smooth with compact support in one chart and \(u\) is smooth, then
      \(T(h\,dd^cu\wedge\beta_1\wedge\dots\wedge\beta_k)=T(u\,dd^ch\wedge\beta_1\wedge\dots\wedge\beta_k)\).
\end{enumerate}
\end{lemma}
\begin{proof}
\leanok
(1) Pointwise linear algebra; \(\Pi_{1,1}\gamma(v,Jv)=\gamma(v,Jv)\), so positivity on complex lines is that of \(\gamma\), and two test forms with the same
values at chart points have the same pairing with a positive current. (2) \(d(Uf)=U\,df+f\,dU\) gives \(d((Uf)\,dg)=U\,d(f\,dg)+f\,dU\wedge dg\), and
\(T(\sigma(d((Uf)\,dg)))=0\) by closedness; the chain rule gives the second identity. (3) In the chart, with real coordinates \(\ell_j\) and a cut-off
\(\chi=1\) near the support of \(h\), \(h\,d^cu-u\,d^ch=\sum_jc_j\,d(\chi\ell_j)\) with smooth compactly supported \(c_j\), so
\(d(h\,d^cu-u\,d^ch)=h\,dd^cu-u\,dd^ch+(dh\wedge d^cu-du\wedge d^ch)\) has \(T\)-pairing \(0\) by closedness, and \(\Pi_{1,1}(dh\wedge d^cu-du\wedge d^ch)=0\)
(the model identity of Theorem~\ref{fpmma:hessian-ibp}).
\end{proof}

\begin{lemma}[Integration by parts against a closed positive current on a compact manifold]
\label{bou:slot-ibp-compact}
\uses{bou:slot-ibp,bou:closed-currents,bou:positive-currents}
\lean{TranscendentalNakamaye.Pending.SlotFormCompactIBP.smoothDDC_finset_sum,TranscendentalNakamaye.Pending.SlotFormCompactIBP.contDiff_of_contDiffOn_of_tsupport_subset,TranscendentalNakamaye.Pending.SlotFormCompactIBP.exists_eq_chartExtend,TranscendentalNakamaye.Pending.SlotFormCompactIBP.slotForm_ibp_of_tsupport_subset,TranscendentalNakamaye.Pending.SlotFormCompactIBP.slotForm_ibp_of_compact}
\leanok
Let \(T\) be a closed positive current of bidimension \((k+1,k+1)\) on a compact complex manifold \(M\), \(\beta_1,\dots,\beta_k\) closed real \((1,1)\)-forms and
\(h,u\) chart-smooth functions. Then \(T(h\,dd^cu\wedge\beta)=T(u\,dd^ch\wedge\beta)\).
\end{lemma}
\begin{proof}
\leanok
For \(h\) with compact support in one chart this is Lemma~\ref{bou:slot-ibp}. In general take a chart-smooth partition of unity \((\chi_a)\) subordinate to finitely many
chart domains; \(h=\sum_a\chi_ah\), each \(\chi_ah\) is the extension by zero of a smooth function with compact support in one chart target, and both sides are
additive in \(h\) (\(dd^c\) is additive and \(T\) only sees the coefficients at points of the charts).
\end{proof}

\begin{lemma}[The Skoda--El Mir lemma in one chart]
\label{bou:skoda-elmir-chart}
\uses{bou:slot-ibp,bou:current-restriction,bou:closed-currents,bou:closed-currents-local,bou:positive-currents,bou:trace-measure}
\source{demailly-cadg:III.2.2,demailly-cadg:III.2.3}
\lean{TranscendentalNakamaye.Pending.SkodaElMirLimit.apply_smulTest_eq,TranscendentalNakamaye.Pending.SkodaElMirLimit.measureReal_withDensity_ofReal,TranscendentalNakamaye.Pending.SkodaElMirLimit.integral_withDensity_ofReal,TranscendentalNakamaye.Pending.SkodaElMirLimit.restrictCurrent_eq_setIntegral,TranscendentalNakamaye.Pending.SkodaElMirLimit.tendsto_apply_smulTest,TranscendentalNakamaye.Pending.SkodaElMirEstimate.elMirCutoff,TranscendentalNakamaye.Pending.SkodaElMirEstimate.contDiff_elMirCutoff,TranscendentalNakamaye.Pending.SkodaElMirEstimate.elMirCutoff_nonneg,TranscendentalNakamaye.Pending.SkodaElMirEstimate.elMirCutoff_le_one,TranscendentalNakamaye.Pending.SkodaElMirEstimate.elMirCutoff_zero,TranscendentalNakamaye.Pending.SkodaElMirEstimate.elMirCutoff_one,TranscendentalNakamaye.Pending.SkodaElMirEstimate.deriv_elMirCutoff_zero,TranscendentalNakamaye.Pending.SkodaElMirEstimate.deriv_elMirCutoff_one,TranscendentalNakamaye.Pending.SkodaElMirEstimate.contDiff_deriv_elMirCutoff,TranscendentalNakamaye.Pending.SkodaElMirEstimate.continuous_deriv_elMirCutoff,TranscendentalNakamaye.Pending.SkodaElMirEstimate.exists_abs_deriv_elMirCutoff_le,TranscendentalNakamaye.Pending.SkodaElMirEstimate.gradForm,TranscendentalNakamaye.Pending.SkodaElMirEstimate.contDiffAt_two_of_chartSmooth,TranscendentalNakamaye.Pending.SkodaElMirEstimate.gradForm_coeff_line,TranscendentalNakamaye.Pending.SkodaElMirEstimate.gradForm_coeff_line_nonneg,TranscendentalNakamaye.Pending.SkodaElMirEstimate.exactPair_coeff_line,TranscendentalNakamaye.Pending.SkodaElMirEstimate.sq_apply_slotForm_le,TranscendentalNakamaye.Pending.SkodaElMirEstimate.apply_slotForm_gradForm_le,TranscendentalNakamaye.Pending.SkodaElMirEstimate.exists_ddcForm_le_smul,TranscendentalNakamaye.Pending.SkodaElMirEstimate.tendsto_apply_slotForm_smulFun,TranscendentalNakamaye.Pending.SkodaElMirEstimate.chartSmooth_comp_real,TranscendentalNakamaye.Pending.SkodaElMirEstimate.tendsto_zero_of_sq_le,TranscendentalNakamaye.Pending.SkodaElMirEstimate.restrictCurrent_closedTestForm_eq_zero}
\leanok
Let \(M\) be a compact complex manifold with a K\"ahler form \(\omega\), \(T\) a closed positive current of bidimension \((k+1,k+1)\), \(\beta_1,\dots,\beta_k\) smooth
closed semipositive \((1,1)\)-forms, \(\Phi(\gamma)=T(\Pi_{1,1}\gamma\wedge\beta_1\wedge\dots\wedge\beta_k)\), \(s(v)=\pi(dd^c(v^2)-2v\,dd^cv)\) (on complex lines
\((dv(w)^2+dv(Jw)^2)/2\)) and \(E\subseteq M\) a Borel set.
\begin{enumerate}
\item (Limit) If \(\varphi\) is a test form, \(h\le1\) is smooth with \(h\varphi=\varphi\), and \(U_i\) are smooth, uniformly bounded, with \(U_i\to\mathbf 1_Eh\)
      pointwise, then \(T(U_i\varphi)\to\mathbf 1_ET(\varphi)\) (Lemma~\ref{bou:current-restriction}).
\item (Cauchy--Schwarz) If \(h_1\ge0\) and \(h_1=1\) where \(a\ne0\), then \(\Phi(a\,\Pi_{1,1}d(v\,dg))^2\le4\,\Phi(h_1s(v))\,\Phi(a^2s(g))\).
\item (Uniform bound) If \(0\le v\le1\), \(dd^cv\ge0\) where \(h_1\ne0\), \(h_1\ge0\) is smooth with compact support in a chart and \(B\omega-dd^ch_1\ge0\), then
      \(\Phi(h_1s(v))\le\pi B\,\Phi(\omega)\); such a \(B\) exists for every such \(h_1\).
\item (El Mir in a chart) Let \(\theta(t)=\rho(3t-1)\), \(\rho\) the smooth transition function (\(0\) on \((-\infty,0]\), \(1\) on \([1,\infty)\)), so that \(\theta\) is
      smooth, \(0\le\theta\le1\), \(\theta(0)=0\), \(\theta(1)=1\), \(\theta'(0)=\theta'(1)=0\) and \(\theta'\) is bounded; let \(f\), \(g\) be
      smooth, \(h_1=\chi\circ e\) with \(0\le\chi\le1\) smooth with compact support in a chart \(e\) and \(h_1=1\) near \(\operatorname{supp}f\), and \(v_i\) smooth with
      \(0\le v_i\le1\) and \(dd^cv_i\ge0\) where \(h_1\ne0\). If \(\theta(v_i)\to\mathbf 1_E\) where \(h_1\ne0\) and \(\theta'(v_i)\to0\) where \(f\ne0\), then
      \(\mathbf 1_ET(\Pi_{1,1}d(f\,dg)\wedge\beta_1\wedge\dots\wedge\beta_k)=0\).
\end{enumerate}
\end{lemma}
\begin{proof}
\leanok
(1) \(T(U\varphi)=\int U\,d(T\wedge(c\,\omega^{k+1}+\varphi))-c\int U\,d\sigma_T\) (Lemma~\ref{bou:current-restriction}); as \(h\varphi=\varphi\), \((1-h)\) kills the
signed measure \(T\wedge\varphi\), so \(\mathbf 1_ET(\varphi)=\int_Eh\,d(T\wedge(c\,\omega^{k+1}+\varphi))-c\int_Eh\,d\sigma_T\), and dominated convergence applies to
both finite measures. (2) For every real \(t\) the two-form \(h_1s(v)+t\,a\,\Pi_{1,1}(dv\wedge dg)+t^2a^2s(g)\) is nonnegative on complex lines (it is
\(((p+tas)^2+(q-tar)^2)/2\) there, \(p,q,r,s\) the derivatives of \(v\), \(g\) along \(w\), \(Jw\), where \(h_1=1\), and \(h_1s(v)\ge0\) elsewhere), so by
Lemma~\ref{bou:slot-ibp}(1) \(\Phi\) of it is \(\ge0\): a nonnegative quadratic polynomial in \(t\) has nonpositive discriminant. (3)
\(s(v)\le\pi\,dd^c(v^2)\) where \(dd^cv\ge0\) and \(v\ge0\), so \(\Phi(h_1s(v))\le\pi\Phi(h_1dd^c(v^2))=\pi\Phi(v^2dd^ch_1)\le\pi B\Phi(\omega)\) by
Lemma~\ref{bou:slot-ibp}(3) and \(v^2\le1\). (4) By (1) with \(U_i=h_1\theta(v_i)\), \(T(h_1\theta(v_i)\,\Pi_{1,1}d(f\,dg)\wedge\beta)\to\mathbf 1_ET(\dots)\); by
Lemma~\ref{bou:slot-ibp}(2) and \(h_1f=f\) it equals \(-\Phi(f\theta'(v_i)\,\Pi_{1,1}d(v_i\,dg))\), whose square is at most
\(4\pi B\,\Phi(\omega)\,\Phi((f\theta'(v_i))^2s(g))\) by (2), (3); the last factor tends to \(0\) by dominated convergence.
\end{proof}

\begin{lemma}[Skoda--El Mir for a linearly flat closed set]
\label{bou:skoda-elmir-submanifold}
\uses{bou:closed-currents,bou:closed-currents-limits,bou:current-restriction,bou:positive-currents,bou:product-manifold,bou:closed-currents-local,bou:slot-ibp,bou:skoda-elmir-chart,bou:concentration-potentials}
\source{demailly-cadg:III.2.2,demailly-cadg:III.2.3,demailly-cadg:III.2.4}
\lean{TranscendentalNakamaye.Pending.LogCutoffLevi.cutoffProfile,TranscendentalNakamaye.Pending.LogCutoffLevi.contDiff_cutoffProfile,TranscendentalNakamaye.Pending.LogCutoffLevi.continuous_cutoffProfile,TranscendentalNakamaye.Pending.LogCutoffLevi.cutoffProfile_of_nonpos,TranscendentalNakamaye.Pending.LogCutoffLevi.cutoffProfile_of_pos,TranscendentalNakamaye.Pending.LogCutoffLevi.cutoffProfile_nonneg,TranscendentalNakamaye.Pending.LogCutoffLevi.cutoffProfile_one,TranscendentalNakamaye.Pending.LogCutoffLevi.cutoffProfile_le_one,TranscendentalNakamaye.Pending.LogCutoffLevi.hasDerivAt_profileExp,TranscendentalNakamaye.Pending.LogCutoffLevi.deriv_cutoffProfile_of_pos,TranscendentalNakamaye.Pending.LogCutoffLevi.deriv_cutoffProfile_of_neg,TranscendentalNakamaye.Pending.LogCutoffLevi.contDiff_deriv_cutoffProfile,TranscendentalNakamaye.Pending.LogCutoffLevi.contDiff_deriv2_cutoffProfile,TranscendentalNakamaye.Pending.LogCutoffLevi.deriv2_cutoffProfile_of_pos,TranscendentalNakamaye.Pending.LogCutoffLevi.deriv_cutoffProfile_of_neg',TranscendentalNakamaye.Pending.LogCutoffLevi.eq_zero_at_zero,TranscendentalNakamaye.Pending.LogCutoffLevi.deriv_cutoffProfile_nonneg,TranscendentalNakamaye.Pending.LogCutoffLevi.deriv2_cutoffProfile_nonneg,TranscendentalNakamaye.Pending.LogCutoffLevi.realHessian_comp,TranscendentalNakamaye.Pending.LogCutoffLevi.leviForm_comp,TranscendentalNakamaye.Pending.LogCutoffLevi.leviForm_affine,TranscendentalNakamaye.Pending.LogCutoffLevi.contDiff_normSq_comp,TranscendentalNakamaye.Pending.LogCutoffLevi.contDiffAt_log_normSq,TranscendentalNakamaye.Pending.LogCutoffLevi.leviForm_log_normSq_nonneg,TranscendentalNakamaye.Pending.LogCutoffLevi.logCutoff,TranscendentalNakamaye.Pending.LogCutoffLevi.logCutoff_of_eq_zero,TranscendentalNakamaye.Pending.LogCutoffLevi.logCutoff_eventuallyEq_zero,TranscendentalNakamaye.Pending.LogCutoffLevi.logCutoff_eventuallyEq,TranscendentalNakamaye.Pending.LogCutoffLevi.contDiff_logCutoff,TranscendentalNakamaye.Pending.LogCutoffLevi.logCutoff_nonneg,TranscendentalNakamaye.Pending.LogCutoffLevi.logCutoff_le_one,TranscendentalNakamaye.Pending.LogCutoffLevi.tendsto_logCutoff,TranscendentalNakamaye.Pending.LogCutoffLevi.smoothDDC_logCutoff_nonneg,TranscendentalNakamaye.Pending.SkodaElMirSubmanifold.IsLinearlyFlat,TranscendentalNakamaye.Pending.SkodaElMirSubmanifold.exists_cutoff,TranscendentalNakamaye.Pending.SkodaElMirSubmanifold.chartExtend_ne_zero,TranscendentalNakamaye.Pending.SkodaElMirSubmanifold.exists_chart_bump,TranscendentalNakamaye.Pending.SkodaElMirSubmanifold.local_vanishing_of_notMem,TranscendentalNakamaye.Pending.SkodaElMirSubmanifold.local_vanishing_of_mem,TranscendentalNakamaye.Pending.SkodaElMirSubmanifold.isClosedCurrent_restrictCurrent_compl,TranscendentalNakamaye.Pending.SkodaElMirSubmanifold.isClosedCurrent_restrictCurrent,TranscendentalNakamaye.Pending.SkodaElMirSubmanifold.diagDifference,TranscendentalNakamaye.Pending.SkodaElMirSubmanifold.isLinearlyFlat_diagonal,TranscendentalNakamaye.Pending.SkodaElMirSubmanifold.isClosedCurrent_restrictCurrent_diagonal}
\leanok
Let \(M\) be a compact complex manifold with a K\"ahler form and \(Y\subseteq M\) a closed set which is \emph{linearly flat}: every point of \(Y\) lies in the
source of a chart \(e\) of the atlas such that, on that source, \(Y=\{L\circ e=0\}\) for a \(\mathbf C\)-linear map \(L\). For every closed positive current
\(T\) of bidimension \((q,q)\), \(\mathbf 1_YT\) and \(\mathbf 1_{M\setminus Y}T\) (Lemma~\ref{bou:current-restriction}) are closed. The diagonal of \(X\times X\) is
closed and linearly flat in the product charts (Lemma~\ref{bou:product-manifold}), with \(L(u,u')=u'-u\).
\end{lemma}
\begin{proof}
\leanok
As \(\mathbf 1_YT=T-\mathbf 1_{M\setminus Y}T\), it suffices to treat \(S=\mathbf 1_{M\setminus Y}T\), and closedness of the positive current \(S\) may be tested locally,
with closed semipositive slots (Lemma~\ref{bou:closed-currents-local}): it suffices that \(S(\Pi_{1,1}d(f\,dg)\wedge\beta)=0\) for \(f\) supported in a small
neighbourhood of a point. Near a point of \(Y\) take the flattening chart, \(f\) supported in \(\{|L\circ e|^2<1\}\), and a chart bump \(h_1=1\) near
\(\operatorname{supp}f\). With \(\rho(t)=e^2\exp(-2/t)\) for \(t>0\) and \(\rho=0\) on \((-\infty,0]\) --- smooth, \(\rho(1)=1\), \(0\le\rho\le1\), \(\rho'\ge0\) and
\(\rho''\ge0\) on \((-\infty,1]\) --- the functions \(\rho(\log|L|^2/(i+1)+1)\) vanish near \(\{L=0\}\), lie in \([0,1]\) on \(\{|L|^2<1\}\), tend to \(1\) off
\(\{L=0\}\) and are plurisubharmonic there: \(\log|L|^2\) is plurisubharmonic off \(\{L=0\}\) (the case \(s=0\) of Lemma~\ref{bou:concentration-potentials}) and
\(\rho\) is convex nondecreasing on \((-\infty,1]\). Multiplied by a cut-off equal to \(1\) near \(\operatorname{supp}h_1\) they are cut-offs \(v_i\) as in
Lemma~\ref{bou:skoda-elmir-chart}(4), with \(E=M\setminus Y\): \(\theta(v_i)\to\mathbf 1_{M\setminus Y}\) where \(h_1\ne0\), and \(\theta'(v_i)\to\theta'(0)\) or \(\theta'(1)\), both
\(0\). Near a point off \(Y\) take \(h_1\) supported off \(Y\) and \(v_i=1\). In both cases Lemma~\ref{bou:skoda-elmir-chart}(4) gives
\(S(\Pi_{1,1}d(f\,dg)\wedge\beta)=0\). For the diagonal, in the product chart at \((p,p)\) the diagonal is \(\{u'-u=0\}\).
\end{proof}

\begin{lemma}[Skoda--El Mir for zero sets of holomorphic maps; graphs]
\label{bou:skoda-elmir-holomorphic}
\uses{bou:skoda-elmir-submanifold,bou:concentration-levi,bou:current-restriction,bou:closed-currents-local,bou:product-manifold}
\source{demailly-cadg:III.2.3}
\lean{TranscendentalNakamaye.Pending.SkodaElMirHolomorphic.leviForm_log_normSq_nonneg_of_holo,TranscendentalNakamaye.Pending.SkodaElMirHolomorphic.holoLogCutoff,TranscendentalNakamaye.Pending.SkodaElMirHolomorphic.holoLogCutoff_of_eq_zero,TranscendentalNakamaye.Pending.SkodaElMirHolomorphic.holoLogCutoff_eventuallyEq_zero,TranscendentalNakamaye.Pending.SkodaElMirHolomorphic.holoLogCutoff_eventuallyEq,TranscendentalNakamaye.Pending.SkodaElMirHolomorphic.contDiffAt_log_normSq_of_contDiffAt,TranscendentalNakamaye.Pending.SkodaElMirHolomorphic.contDiffAt_holoLogCutoff,TranscendentalNakamaye.Pending.SkodaElMirHolomorphic.holoLogCutoff_nonneg,TranscendentalNakamaye.Pending.SkodaElMirHolomorphic.holoLogCutoff_le_one,TranscendentalNakamaye.Pending.SkodaElMirHolomorphic.tendsto_holoLogCutoff,TranscendentalNakamaye.Pending.SkodaElMirHolomorphic.smoothDDC_holoLogCutoff_nonneg,TranscendentalNakamaye.Pending.SkodaElMirHolomorphic.IsHoloFlat,TranscendentalNakamaye.Pending.SkodaElMirHolomorphic.IsLinearlyFlat.isHoloFlat,TranscendentalNakamaye.Pending.SkodaElMirHolomorphic.isOpen_inter_normSq_lt,TranscendentalNakamaye.Pending.SkodaElMirHolomorphic.local_vanishing_of_mem,TranscendentalNakamaye.Pending.SkodaElMirHolomorphic.isClosedCurrent_restrictCurrent_compl,TranscendentalNakamaye.Pending.SkodaElMirHolomorphic.isClosedCurrent_restrictCurrent,TranscendentalNakamaye.Pending.SkodaElMirHolomorphic.graphDefining,TranscendentalNakamaye.Pending.SkodaElMirHolomorphic.isHoloFlat_graph,TranscendentalNakamaye.Pending.SkodaElMirHolomorphic.isClosedCurrent_restrictCurrent_graph}
\leanok
Let \(M\) be a compact complex manifold with a K\"ahler form and \(\Delta\subseteq M\) a closed set which is \emph{holomorphically flat}: every point of \(\Delta\) lies in
the source of a chart \(e\) of the atlas and there are an open \(W\) in the chart target containing its image and a map \(Y\), holomorphic on \(W\), with
\(\Delta=\{Y\circ e=0\}\) over \(W\). For every closed positive current \(T\) of bidimension \((q,q)\), \(\mathbf 1_\Delta T\) and \(\mathbf 1_{M\setminus\Delta}T\) are closed.
Linearly flat sets (Lemma~\ref{bou:skoda-elmir-submanifold}) are holomorphically flat. If \(f\colon V\to X\) is a holomorphic map between compact complex manifolds of
the same dimension, its graph \(\Gamma_f=\{(v,x):f(v)=x\}\) is closed and holomorphically flat in the product charts of \(V\times X\) (with \(Y(z,w)=w-F(z)\), \(F\)
the coordinate representative of \(f\)); so \(\mathbf 1_{\Gamma_f}T\) and \(\mathbf 1_{V\times X\setminus\Gamma_f}T\) are closed for every closed positive \(T\) on \(V\times X\).
\end{lemma}
\begin{proof}
\leanok
As in Lemma~\ref{bou:skoda-elmir-submanifold}, with the cut-offs \(V_i=\rho(\log|Y|^2/(i+1)+1)\) on \(W\): they are smooth, vanish near \(\{Y=0\}\), tend to \(1\) off it, and
are plurisubharmonic on \(W\cap\{|Y|^2<1\}\), because \(\log|Y|^2\) is plurisubharmonic off \(\{Y=0\}\) for holomorphic \(Y\) (the Levi estimate of
Lemma~\ref{bou:concentration-levi} with \(G_j=Y_j\)). The local vanishing argument and its assembly are unchanged.
\end{proof}

\begin{lemma}[What a positive current carried by a set does not see]
\label{bou:carried-current-tests}
\uses{bou:current-restriction,bou:positive-currents,bou:trace-measure,bou:closed-currents,bou:slot-ibp}
\source{demailly-cadg:III.2.10}
\lean{TranscendentalNakamaye.Pending.CarriedCurrentTests.integral_eq_zero_of_carried,TranscendentalNakamaye.Pending.CarriedCurrentTests.apply_smulTest_eq_zero,TranscendentalNakamaye.Pending.CarriedCurrentTests.apply_smulTest_eq_of_eq_one,TranscendentalNakamaye.Pending.CarriedCurrentTests.apply_slotForm_smulFun_exactPair_eq_zero}
\leanok
Let \(M\) be a compact complex manifold with a K\"ahler form, \(E\subseteq M\) a Borel set and \(T\) a positive current whose trace measure is carried by \(E\),
\(\sigma_T(M\setminus E)=0\). Then \(T(g\varphi)=0\) for every test form \(\varphi\) and every smooth \(g\) vanishing on \(E\), and \(T(g\varphi)=T(\varphi)\) if \(g=1\) on
\(E\). If moreover \(T\) is closed, then \(T(h\,\Pi_{1,1}(dg\wedge dF)\wedge\beta_1\wedge\dots\wedge\beta_k)=0\) for smooth \(h\), \(F\), smooth \(g\) vanishing on
\(E\) and smooth closed real \((1,1)\)-forms \(\beta_j\).
\end{lemma}
\begin{proof}
\leanok
\(T(g\varphi)=\int g\,d(T\wedge(c\,\omega^q+\varphi))-c\int g\,d\sigma_T\) (Lemma~\ref{bou:current-restriction}), and both measures are dominated by multiples of
\(\sigma_T\), which does not charge \(M\setminus E\); \(g=0\) on \(E\) gives \(0\). For the second claim, \(d((hg)\,dF)=h\,dg\wedge dF+g\,dh\wedge dF\): the closed
test form \(\Pi_{1,1}d((hg)\,dF)\wedge\beta\) has \(T\)-pairing \(0\) (closedness), and the second term is a multiple of \(g\) (CADG~III.2.10:
\(dg\wedge T=d(gT)-g\,dT=0\)).
\end{proof}

\begin{lemma}[A measure with vanishing distributional derivatives is a multiple of Lebesgue measure]
\label{bou:zero-derivative-measure}
\source{demailly-cadg:III.2.13}
\lean{TranscendentalNakamaye.Pending.ZeroDerivativeMeasure.HasZeroDerivOn,TranscendentalNakamaye.Pending.ZeroDerivativeMeasure.tsupport_comp_sub,TranscendentalNakamaye.Pending.ZeroDerivativeMeasure.hasCompactSupport_comp_sub,TranscendentalNakamaye.Pending.ZeroDerivativeMeasure.integral_comp_sub_eq,TranscendentalNakamaye.Pending.ZeroDerivativeMeasure.integral_mul_eq_double,TranscendentalNakamaye.Pending.ZeroDerivativeMeasure.integral_mul_integral_comm,TranscendentalNakamaye.Pending.ZeroDerivativeMeasure.exists_integral_eq_const_mul}
\leanok
Let \(\nu\) be a finite measure on a finite-dimensional real normed space \(E\) with an additive Haar measure \(\lambda\), \(x_0\in E\) and \(r>0\), such that
\(\int d\psi(v)\,d\nu=0\) for every smooth \(\psi\) with compact support in \(B(x_0,3r)\) and every \(v\in E\). Then there is a constant \(c\) with
\(\int\psi\,d\nu=c\int\psi\,d\lambda\) for every smooth \(\psi\) with compact support in \(B(x_0,r)\).
\end{lemma}
\begin{proof}
\leanok
(1) For such \(\psi\) and \(v\), \(s\mapsto\int\psi(x-sv)\,d\nu(x)\) has derivative \(-\int d\psi(x-sv)(v)\,d\nu=0\) (differentiation under the integral sign), so it
is constant. (2) For \(\psi\), \(\rho\) supported in \(B(x_0,r)\), Fubini and (1) give \((\int\psi\,d\nu)(\int\rho\,d\lambda)=\iint\rho(y)\psi(x-y+x_0)\,d\lambda(y)\,d\nu(x)\),
and the substitution \(y\mapsto x+x_0-y\) exchanges \(\psi\) and \(\rho\): \((\int\psi\,d\nu)(\int\rho\,d\lambda)=(\int\rho\,d\nu)(\int\psi\,d\lambda)\). (3) Take for \(\rho\) a
bump with \(\int\rho\,d\lambda=1\) and \(c=\int\rho\,d\nu\).
\end{proof}

\begin{lemma}[Positive currents of top bidimension are measures]
\label{bou:top-bidimension-currents}
\uses{bou:positive-currents,bou:trace-measure,bou:integration-current}
\source{demailly-cadg:III.1.14}
\lean{TranscendentalNakamaye.Pending.TopBidimensionCurrents.topForm_eq_smul,TranscendentalNakamaye.Pending.TopBidimensionCurrents.chartSmooth_topRatio,TranscendentalNakamaye.Pending.TopBidimensionCurrents.coeff_eq_topRatio_smul,TranscendentalNakamaye.Pending.TopBidimensionCurrents.apply_eq_integral_topRatio,TranscendentalNakamaye.Pending.TopBidimensionCurrents.eq_smul_integrationCurrent}
\leanok
Let \(M\) be a compact complex manifold of dimension \(n\) with a K\"ahler form \(\omega\). A top-degree real form on \(\mathbf C^n\) is a multiple of \(\omega^n\), and the
ratio \(\varphi/\omega^n\) of a smooth top-degree form is smooth. A positive current \(S\) of bidimension \((n,n)\) satisfies \(S(\varphi)=\int\varphi/\omega^n\,d\sigma_S\);
in particular, if \(\sigma_S=c\cdot\mathrm{vol}_\omega\) with \(c\ge0\), then \(S=c\,[M]\).
\end{lemma}
\begin{proof}
\leanok
Top-degree real forms on \(\mathbf C^n\) form a line and \(\omega^n\) has positive density, so \(\varphi=(\varphi/\omega^n)\,\omega^n\) pointwise, with a smooth ratio (quotient of
smooth coefficients by a positive one). A positive current depends only on the pointwise values of its test forms, so \(S(\varphi)=S((\varphi/\omega^n)\omega^n)\),
which is the integral of the ratio against the trace measure (Lemma~\ref{bou:trace-measure}); for \(\sigma_S=c\,\mathrm{vol}_\omega\) this is \(c\,[M](\varphi)\).
\end{proof}

\begin{lemma}[Closed positive currents carried by a set do not see normal directions]
\label{bou:normal-direction-vanishing}
\uses{bou:carried-current-tests,bou:slot-ibp,bou:closed-currents,bou:product-currents}
\source{demailly-cadg:III.2.10}
\lean{TranscendentalNakamaye.Pending.NormalDirectionVanishing.wedgeProd_comp_perm,TranscendentalNakamaye.Pending.NormalDirectionVanishing.wedgeProd_snoc_eq_cons,TranscendentalNakamaye.Pending.NormalDirectionVanishing.IsHoloPairOn,TranscendentalNakamaye.Pending.NormalDirectionVanishing.exactPair_coeff_line_holo,TranscendentalNakamaye.Pending.NormalDirectionVanishing.kahlerUnit,TranscendentalNakamaye.Pending.NormalDirectionVanishing.smulTest_sub_right,TranscendentalNakamaye.Pending.NormalDirectionVanishing.smulTest_smul_right,TranscendentalNakamaye.Pending.NormalDirectionVanishing.smulTest_wedge_mem_positiveCone,TranscendentalNakamaye.Pending.NormalDirectionVanishing.apply_kahlerUnit_eq_zero,TranscendentalNakamaye.Pending.NormalDirectionVanishing.apply_wedge_holoPair_eq_zero,TranscendentalNakamaye.Pending.NormalDirectionVanishing.apply_wedge_normalPair_eq_zero_of_pos,TranscendentalNakamaye.Pending.NormalDirectionVanishing.apply_wedge_normalPair_eq_zero}
\leanok
Let \(T\) be a closed positive current on a compact complex manifold with a K\"ahler form whose trace measure is carried by a Borel set \(E\). Let \(g\), \(F\),
\(G\), \(\rho\ge0\) be smooth with \(g=0\) on \(E\) and \(g+iF\) holomorphic where \(\rho\ne0\), and \(H\) smooth with \(\rho=1\) where \(H\ne0\). Then
\(T(H\,R\wedge\Pi_{1,1}(dg\wedge dG))=0\) for every test form \(R\) of the complementary bidegree.
\end{lemma}
\begin{proof}
\leanok
(1) For \(R=\omega^k\) this is the case of closed slots (Lemma~\ref{bou:carried-current-tests}; the slots commute). (2) For a positive \(R\) and \(P\ge0\)
supported where \(\rho\ne0\), the form \(N=\Pi_{1,1}(dg\wedge dF)\) is \(dg^2+dF^2\) on complex lines (holomorphy of \(g+iF\)), hence semipositive, so
\(0\le T(P\,R\wedge N)\le c\,T(P\,\omega^k\wedge N)=0\), \(\omega^k\) being an order unit. (3) For every real \(t\), \(N+t\,\Pi_{1,1}(dg\wedge dG)+t^2\,\Pi_{1,1}s(G)\ge0\) on
complex lines, so by (2) and Cauchy--Schwarz the linear term \(T(P\,R\wedge\Pi_{1,1}(dg\wedge dG))\) vanishes. (4) Linearity in \(R\) and \(H\).
\end{proof}

\begin{lemma}[{Closed positive currents of top bidimension are multiples of \([X]\)}]
\label{bou:closed-top-currents}
\uses{bou:top-bidimension-currents,bou:zero-derivative-measure,bou:closed-currents,bou:integration-current,hgeom:chart-partition}
\source{demailly-cadg:III.2.13,demailly-cadg:III.2.14}
\lean{TranscendentalNakamaye.Pending.ClosedTopCurrents.mixedCoefficient_cons_one,TranscendentalNakamaye.Pending.ClosedTopCurrents.re_trace_pairMatrix,TranscendentalNakamaye.Pending.ClosedTopCurrents.imInner,TranscendentalNakamaye.Pending.ClosedTopCurrents.imInner_apply,TranscendentalNakamaye.Pending.ClosedTopCurrents.sum_imInner,TranscendentalNakamaye.Pending.ClosedTopCurrents.topDensity_slot_eq,TranscendentalNakamaye.Pending.ClosedTopCurrents.topDensity_closedTestForm_eq,TranscendentalNakamaye.Pending.ClosedTopCurrents.volDensity,TranscendentalNakamaye.Pending.ClosedTopCurrents.volDensity_pos,TranscendentalNakamaye.Pending.ClosedTopCurrents.contDiffOn_volDensity,TranscendentalNakamaye.Pending.ClosedTopCurrents.continuousOn_volDensity,TranscendentalNakamaye.Pending.ClosedTopCurrents.setIntegral_kahlerVolume_image,TranscendentalNakamaye.Pending.ClosedTopCurrents.chartMeasure,TranscendentalNakamaye.Pending.ClosedTopCurrents.measurable_indicator_inv_volDensity,TranscendentalNakamaye.Pending.ClosedTopCurrents.isFiniteMeasure_chartMeasure,TranscendentalNakamaye.Pending.ClosedTopCurrents.integral_chartMeasure,TranscendentalNakamaye.Pending.ClosedTopCurrents.hasZeroDerivOn_chartMeasure,TranscendentalNakamaye.Pending.ClosedTopCurrents.exists_local_const,TranscendentalNakamaye.Pending.ClosedTopCurrents.integral_kahlerVolume_pos,TranscendentalNakamaye.Pending.ClosedTopCurrents.eq_smul_kahlerVolume_of_isClosedCurrent,TranscendentalNakamaye.Pending.ClosedTopCurrents.eq_smul_integrationCurrent_of_isClosedCurrent}
\leanok
Let \(X\) be a compact \emph{connected} complex manifold of dimension \(m+1\) with a K\"ahler form \(\omega\). Every closed positive current \(S\) of bidimension
\((m+1,m+1)\) on \(X\) is \(c\,[X]\) for a constant \(c\ge0\).
\end{lemma}
\begin{proof}
\leanok
By Lemma~\ref{bou:top-bidimension-currents} \(S(\varphi)=\int\varphi/\omega^{m+1}\,d\sigma_S\). In a chart, with the flat K\"ahler form, the top density of
\(\Pi_{1,1}(dF\wedge dG)\wedge(dd^c|z|^2)^m\) for \(G=\operatorname{Im}\langle v,\cdot\rangle\) is \(\frac{m!}{\pi^m}\,dF(v)\), so the closedness of \(S\) says that the measure
\(\sigma_S/\omega^{m+1}\) has vanishing distributional derivatives in the chart: by Lemma~\ref{bou:zero-derivative-measure} it is a constant multiple of Lebesgue measure
near every point. The local constants agree on overlaps and \(X\) is connected, so \(\sigma_S=c\,\mathrm{vol}_\omega\) (a partition of unity,
Lemma~\ref{hgeom:chart-partition}), and \(S=c\,[X]\) by Lemma~\ref{bou:top-bidimension-currents}.
\end{proof}

\begin{lemma}[The tangential reduction on the diagonal]
\label{bou:diagonal-tangential-reduction}
\uses{bou:normal-direction-vanishing,bou:carried-current-tests,bou:product-manifold,bou:product-currents,bou:closed-currents,bou:positive-currents,hgeom:chart-partition}
\source{demailly-cadg:III.2.10}
\lean{TranscendentalNakamaye.Pending.DiagonalTangentialReduction.elemForm,TranscendentalNakamaye.Pending.DiagonalTangentialReduction.elemForm_succ,TranscendentalNakamaye.Pending.DiagonalTangentialReduction.elemForm_snoc,TranscendentalNakamaye.Pending.DiagonalTangentialReduction.elemForm_coeff,TranscendentalNakamaye.Pending.DiagonalTangentialReduction.elemForm_mem,TranscendentalNakamaye.Pending.DiagonalTangentialReduction.elemTest,TranscendentalNakamaye.Pending.DiagonalTangentialReduction.coe_elemTest,TranscendentalNakamaye.Pending.DiagonalTangentialReduction.coe_wedgeTestLM,TranscendentalNakamaye.Pending.DiagonalTangentialReduction.elemSet,TranscendentalNakamaye.Pending.DiagonalTangentialReduction.wedgeTwo_zero_left,TranscendentalNakamaye.Pending.DiagonalTangentialReduction.testSpace_le_span_elemSet,TranscendentalNakamaye.Pending.DiagonalTangentialReduction.span_elemSet_le_testSpace,TranscendentalNakamaye.Pending.DiagonalTangentialReduction.linearMap_ext_elem,TranscendentalNakamaye.Pending.DiagonalTangentialReduction.pullback_elemForm,TranscendentalNakamaye.Pending.DiagonalTangentialReduction.twoFormBilin,TranscendentalNakamaye.Pending.DiagonalTangentialReduction.twoFormBilin_apply,TranscendentalNakamaye.Pending.DiagonalTangentialReduction.twoForm_eq_sum,TranscendentalNakamaye.Pending.DiagonalTangentialReduction.twoForm_sub_comp,TranscendentalNakamaye.Pending.DiagonalTangentialReduction.diagProj,TranscendentalNakamaye.Pending.DiagonalTangentialReduction.diagProj_apply,TranscendentalNakamaye.Pending.DiagonalTangentialReduction.diagDifference_apply,TranscendentalNakamaye.Pending.DiagonalTangentialReduction.normalCoord,TranscendentalNakamaye.Pending.DiagonalTangentialReduction.normalRe,TranscendentalNakamaye.Pending.DiagonalTangentialReduction.normalVec,TranscendentalNakamaye.Pending.DiagonalTangentialReduction.normalRe_apply,TranscendentalNakamaye.Pending.DiagonalTangentialReduction.sub_diagProj_eq,TranscendentalNakamaye.Pending.DiagonalTangentialReduction.pairForm,TranscendentalNakamaye.Pending.DiagonalTangentialReduction.pairForm_apply,TranscendentalNakamaye.Pending.DiagonalTangentialReduction.halfVec,TranscendentalNakamaye.Pending.DiagonalTangentialReduction.halfVec_sum,TranscendentalNakamaye.Pending.DiagonalTangentialReduction.sub_compDiagProj_eq,TranscendentalNakamaye.Pending.DiagonalTangentialReduction.diagProjMap,TranscendentalNakamaye.Pending.DiagonalTangentialReduction.contMDiff_diagProjMap,TranscendentalNakamaye.Pending.DiagonalTangentialReduction.diagChart,TranscendentalNakamaye.Pending.DiagonalTangentialReduction.diagChart_val,TranscendentalNakamaye.Pending.DiagonalTangentialReduction.diagChart_mapsTo,TranscendentalNakamaye.Pending.DiagonalTangentialReduction.diagProj_mem_target,TranscendentalNakamaye.Pending.DiagonalTangentialReduction.coordinateMap_diagProjMap,TranscendentalNakamaye.Pending.DiagonalTangentialReduction.pullback_diagProjMap_coeff,TranscendentalNakamaye.Pending.DiagonalTangentialReduction.diagonal_chart,TranscendentalNakamaye.Pending.DiagonalTangentialReduction.normalIm,TranscendentalNakamaye.Pending.DiagonalTangentialReduction.isHoloPairOn_normal,TranscendentalNakamaye.Pending.DiagonalTangentialReduction.chartSmooth_chartExtend_mul,TranscendentalNakamaye.Pending.DiagonalTangentialReduction.exactPair_cutoff_coeff,TranscendentalNakamaye.Pending.DiagonalTangentialReduction.isTypeOneOne_sub,TranscendentalNakamaye.Pending.DiagonalTangentialReduction.isTypeOneOne_sub_pullback,TranscendentalNakamaye.Pending.DiagonalTangentialReduction.coordinateWedge_smul_zero_left,TranscendentalNakamaye.Pending.DiagonalTangentialReduction.exists_nhds_apply_wedge_sub_pullback_eq_zero,TranscendentalNakamaye.Pending.DiagonalTangentialReduction.pullSlots,TranscendentalNakamaye.Pending.DiagonalTangentialReduction.apply_elemTest_comp_perm,TranscendentalNakamaye.Pending.DiagonalTangentialReduction.apply_elemTest_update_pullback,TranscendentalNakamaye.Pending.DiagonalTangentialReduction.elemTest_congr,TranscendentalNakamaye.Pending.DiagonalTangentialReduction.apply_elemTest_pullSlots,TranscendentalNakamaye.Pending.DiagonalTangentialReduction.apply_elemTest_comp_diagProjMap,TranscendentalNakamaye.Pending.DiagonalTangentialReduction.apply_pullback_diagProjMap_eq}
\leanok
Let \(X\) be a compact complex manifold of dimension \(n\), \(\widetilde\omega\) a K\"ahler form on \(X\times X\), and \(\Theta\) a closed positive current on
\(X\times X\) acting on test forms of bidegree \((p,p)\), \(p\ge1\), whose trace measure is carried by the diagonal \(\Delta\). With \(\Phi=\operatorname{diag}\circ p_1\),
\(\Phi(x,y)=(x,x)\),
\[
 \Theta(\Phi^*\varphi)=\Theta(\varphi)\qquad\text{for every test form }\varphi ,
\]
i.e.\ \(\Theta\) sees only the restriction of \(\varphi\) to the tangent spaces of \(\Delta\) along \(\Delta\).
\end{lemma}
\begin{proof}
\leanok
Test forms are spanned by the elementary forms \(h\,b_1\wedge\dots\wedge b_p\) with \(b_i\) real \((1,1)\)-forms, and
\(\Phi^*(h\,b_1\wedge\dots\wedge b_p)=(h\circ\Phi)\,\Phi^*b_1\wedge\dots\wedge\Phi^*b_p\). Localise \(h\) by a partition of unity subordinate to product charts at diagonal
points and to \(X\times X\setminus\Delta\), where \(\Theta\) sees nothing (Lemma~\ref{bou:carried-current-tests}). In the product chart at \((x_0,x_0)\) the map \(\Phi\)
is the linear projection \(\pi(u,u')=(u,u)\), and for a real two-form \(A\) and \(\nu=1-\pi\),
\(A(\xi,\eta)-A(\pi\xi,\pi\eta)=A(\nu\xi,(\pi+\nu/2)\eta)-A(\nu\eta,(\pi+\nu/2)\xi)\), where \(\nu\xi=\sum_a\operatorname{Re}L_a(\xi)\,N_a\) with \(L_a\) the
\(\mathbf C\)-linear coordinates of \(u'-u\). Hence \(b-\Phi^*b\) is a sum of terms \(\gamma\,\Pi_{1,1}(d\operatorname{Re}L_a\wedge ds_c)\) and
\(\eta\,\Pi_{1,1}(ds_c\wedge ds_d)\) with coefficients \(\eta\) vanishing on \(\Delta\). The first are killed against every other factor by
Lemma~\ref{bou:normal-direction-vanishing} (\(\operatorname{Re}L_a+i\operatorname{Im}L_a\) is holomorphic), the second by Lemma~\ref{bou:carried-current-tests};
replacing the factors \(b_i\) by \(\Phi^*b_i\) one at a time (two-forms commute) and \(h\) by \(h\circ\Phi\) (equal on \(\Delta\)) proves the claim.
\end{proof}

\begin{lemma}[Closed positive currents carried by the diagonal]
\label{bou:diagonal-support}
\uses{bou:closed-currents,bou:closed-currents-limits,bou:current-restriction,bou:product-manifold,bou:integration-current,bou:positive-currents,bou:trace-measure,bou:closed-currents-local,bou:fibre-pushforward,bou:carried-current-tests,bou:zero-derivative-measure,bou:top-bidimension-currents,bou:normal-direction-vanishing,bou:diagonal-tangential-reduction,bou:closed-top-currents}
\source{demailly-cadg:III.2.10,demailly-cadg:III.2.13,demailly-cadg:III.2.14}
\lean{TranscendentalNakamaye.Pending.DiagonalSupport.wedgeProd_compCLM,TranscendentalNakamaye.Pending.DiagonalSupport.comp_two',TranscendentalNakamaye.Pending.DiagonalSupport.pi11_compCLM,TranscendentalNakamaye.Pending.DiagonalSupport.pullback_exactPair_coeff,TranscendentalNakamaye.Pending.DiagonalSupport.pullback_closedTestForm_coeff,TranscendentalNakamaye.Pending.DiagonalSupport.IsClosedCurrent.pushforward,TranscendentalNakamaye.Pending.DiagonalSupport.eq_smul_diagonalCurrent}
\leanok
Let \(X\) be a connected compact K\"ahler manifold of dimension \(n\ge1\), \(\operatorname{diag}\colon X\to X\times X\) the diagonal map and
\(\Delta=\operatorname{diag}(X)\). If \(\Theta\) is a closed positive current of bidimension \((n,n)\) on \(X\times X\) whose trace measure (for a K\"ahler form \(\widetilde\omega\) on
\(X\times X\)) is carried by \(\Delta\), then \(\Theta=\lambda\,\operatorname{diag}_*[X]\) for some \(\lambda\ge0\), \(\operatorname{diag}_*[X]\) being the diagonal current of Lemma~\ref{bou:fibre-pushforward}.
\end{lemma}
\begin{proof}
\leanok
(1) \emph{Push-forward.} \(\mu=p_{1*}\Theta\), \(\mu(\psi)=\Theta(p_1^*\psi)\), is a positive current of top bidimension \((n,n)\) on \(X\), and it is closed: pull-back
commutes with \(\Pi_{1,1}\,d(f\,dg)\), so every closedness test of \(\mu\) is a closedness test of \(\Theta\). By Lemma~\ref{bou:closed-top-currents} (\(X\) is
connected), \(\mu=c\,[X]\) with \(c\ge0\). (2) \emph{Tangential reduction.} By Lemma~\ref{bou:diagonal-tangential-reduction},
\(\Theta(\varphi)=\Theta(\Phi^*\varphi)\) with \(\Phi=\operatorname{diag}\circ p_1\), and \(\Phi^*\varphi\) and \(p_1^*\operatorname{diag}^*\varphi\) agree pointwise (a positive
current depends only on the pointwise values of its test forms). Hence
\(\Theta(\varphi)=\mu(\operatorname{diag}^*\varphi)=c\,[X](\operatorname{diag}^*\varphi)=c\,\operatorname{diag}_*[X](\varphi)\).
(Step (i) of the earlier sketch --- \(\Theta\) does not see the normal directions --- is Lemmas~\ref{bou:carried-current-tests}
and~\ref{bou:normal-direction-vanishing}, assembled in Lemma~\ref{bou:diagonal-tangential-reduction}; steps (ii)--(iii) are
Lemma~\ref{bou:closed-top-currents} applied to \(p_{1*}\Theta\).)
\end{proof}

\begin{corollary}[The diagonal part of a closed positive current]
\label{bou:diagonal-part}
\uses{bou:diagonal-support,bou:skoda-elmir-submanifold,bou:current-restriction}
\lean{TranscendentalNakamaye.Pending.DiagonalSupport.exists_restrictCurrent_diagonal_eq_smul}
\leanok
Let \(X\) be a connected compact K\"ahler manifold of dimension \(n\ge1\) and \(\widetilde\omega\) a K\"ahler form on \(X\times X\). For every closed positive
current \(\Theta\) of bidimension \((n,n)\) on \(X\times X\) there is \(\lambda\ge0\) with \(\mathbf 1_\Delta\Theta=\lambda\operatorname{diag}_*[X]\), and then
\(\Theta-\lambda\operatorname{diag}_*[X]\) is positive.
\end{corollary}
\begin{proof}
\leanok
\(\mathbf 1_\Delta\Theta\) is positive, closed (Lemma~\ref{bou:skoda-elmir-submanifold}) and carried by \(\Delta\), so Lemma~\ref{bou:diagonal-support} gives
\(\lambda\); and \(\Theta-\mathbf 1_\Delta\Theta=\mathbf 1_{X\times X\setminus\Delta}\Theta\) is positive (Lemma~\ref{bou:current-restriction}).
\end{proof}

\begin{lemma}[Mass on a closed set passes to weak limits along a shrinking family]
\label{bou:closed-set-mass-limit}
\uses{bou:trace-measure}
\lean{TranscendentalNakamaye.Pending.DiagonalSupport.le_traceMeasure_iInter_of_tendsto}
\leanok
Let \(M\) be a compact complex manifold with a K\"ahler form, \(T_i\) positive currents with bounded masses converging weakly (along a nontrivial filter) to a
positive current \(S\), and \(F_1\supseteq F_2\supseteq\cdots\) closed sets with \(\sigma_{T_i}(F_k)\ge\delta\) eventually, for every \(k\). Then
\(\sigma_S\bigl(\bigcap_kF_k\bigr)\ge\delta\).
\end{lemma}
\begin{proof}
\leanok
\(\sigma_S(F_k)\ge\delta\) for every \(k\) by the upper semicontinuity of the mass of closed sets under weak limits (Lemma~\ref{bou:trace-measure}), and
\(\sigma_S(F_k)\to\sigma_S(\bigcap_kF_k)\) by continuity from above, \(\sigma_S\) being finite.
\end{proof}

\begin{corollary}[The diagonal part is bounded below by the mass on the diagonal]
\label{bou:diagonal-lower-bound}
\uses{bou:diagonal-part,bou:diagonal-current-mass,bou:closed-set-mass-limit,bou:current-restriction,bou:trace-measure,bou:closed-currents-limits}
\source{bou02-volume:4.10,dp04-kahler-cone:2.12}
\lean{TranscendentalNakamaye.Pending.DiagonalLowerBound.exists_isPositive_sub_smul_diagonalCurrent_of_le,TranscendentalNakamaye.Pending.DiagonalLowerBound.exists_isPositive_sub_smul_diagonalCurrent_of_tendsto}
\leanok
Let \(X\) be a connected compact complex manifold of dimension \(n\ge1\) with a K\"ahler form \(\omega\), \(\widetilde\omega=p_1^*\omega+p_2^*\omega\), and
\(\sigma_\Theta\) the trace measure of a positive current \(\Theta\) on \(X\times X\) for \(\widetilde\omega\).
\begin{enumerate}
\item If \(\Theta\) is a closed positive current of bidimension \((n,n)\) on \(X\times X\) and \(\delta\le\sigma_\Theta(\Delta)\), there is
\(c\ge\max\bigl(0,\delta/(2^n\int_X\omega^n)\bigr)\) with \(\Theta-c\operatorname{diag}_*[X]\) positive.
\item The same conclusion holds for a positive current \(\Theta\) that is the weak limit (along a nontrivial filter) of closed positive currents \(T_i\) of
bidimension \((n,n)\) with eventually bounded masses, provided there are closed sets \(F_1\supseteq F_2\supseteq\cdots\) with \(\bigcap_kF_k=\Delta\) and, for every
\(k\), eventually \(\sigma_{T_i}(F_k)\ge\delta\).
\end{enumerate}
\end{corollary}
\begin{proof}
\leanok
(1) Corollary~\ref{bou:diagonal-part} gives \(c\ge0\) with \(\mathbf 1_\Delta\Theta=c\operatorname{diag}_*[X]\) and \(\Theta-c\operatorname{diag}_*[X]\) positive. Taking masses,
\(\sigma_\Theta(\Delta)=\operatorname{mass}(\mathbf 1_\Delta\Theta)=c\operatorname{mass}(\operatorname{diag}_*[X])=c\,2^n\int_X\omega^n\)
(Lemmas~\ref{bou:current-restriction} and~\ref{bou:diagonal-current-mass}), so \(\delta\le c\,2^n\int_X\omega^n\); if \(\int_X\omega^n=0\) the bound reads
\(0\le c\). (2) \(\Theta\) is closed as a weak limit of closed currents (Lemma~\ref{bou:closed-currents-limits}(3)), and
\(\sigma_\Theta(\Delta)\ge\delta\) by Lemma~\ref{bou:closed-set-mass-limit}; apply (1).
\end{proof}


\begin{lemma}[Integration along the fibres of the first projection; the diagonal current]
\label{bou:fibre-pushforward}
\uses{bou:product-manifold,bou:product-currents,bou:positive-currents,bou:integration-current,bou:smooth-form-current}
\source{dp04-kahler-cone:2.12,bou02-volume:4.8}
\lean{TranscendentalNakamaye.Pending.DiagonalFibreIntegration.pullback_comp_geometricallyEq,TranscendentalNakamaye.Pending.DiagonalFibreIntegration.pullback_id_geometricallyEq,TranscendentalNakamaye.Pending.DiagonalFibreIntegration.wedgeTwo_geometricallyEq,TranscendentalNakamaye.Pending.DiagonalFibreIntegration.IsPositive.apply_eq_of_geometricallyEq,TranscendentalNakamaye.Pending.DiagonalFibreIntegration.fibrePushforward,TranscendentalNakamaye.Pending.DiagonalFibreIntegration.fibrePushforward_apply,TranscendentalNakamaye.Pending.DiagonalFibreIntegration.coe_fibreTest,TranscendentalNakamaye.Pending.DiagonalFibreIntegration.fibrePushforward_add,TranscendentalNakamaye.Pending.DiagonalFibreIntegration.fibrePushforward_smul,TranscendentalNakamaye.Pending.DiagonalFibreIntegration.fibrePushforward_sub,TranscendentalNakamaye.Pending.DiagonalFibreIntegration.IsPositive.fibrePushforward,TranscendentalNakamaye.Pending.DiagonalFibreIntegration.tendsto_fibrePushforward_apply,TranscendentalNakamaye.Pending.DiagonalFibreIntegration.diagonalCurrent,TranscendentalNakamaye.Pending.DiagonalFibreIntegration.diagonalCurrent_apply,TranscendentalNakamaye.Pending.DiagonalFibreIntegration.isPositive_diagonalCurrent,TranscendentalNakamaye.Pending.DiagonalFibreIntegration.pullback_diag_fibreTest_geometricallyEq,TranscendentalNakamaye.Pending.DiagonalFibreIntegration.fibrePushforward_diagonalCurrent,TranscendentalNakamaye.Pending.DiagonalFibreIntegration.fibrePushforward_diagonalCurrent_eq_smoothCurrent,TranscendentalNakamaye.Pending.DiagonalFibreIntegration.isPositive_fibrePushforward_sub}
\leanok
Let \(X\) and \(Y\) be compact complex manifolds, \(p_1,p_2\) the projections of \(X\times Y\), \(\beta\) a smooth real \((1,1)\)-form on \(Y\) and \(\Theta\) a
current of bidimension \((q+1,q+1)\) on \(X\times Y\). The current \(p_{1*}(\Theta\wedge p_2^*\beta)\) of bidimension \((q,q)\) on \(X\)
(Definition~\ref{bou:product-currents}) is linear and weakly continuous in \(\Theta\), and positive when \(\Theta\) is positive and \(\beta\) semipositive. For a
K\"ahler form \(\omega\) on \(X\) let \([\Delta]=\operatorname{diag}_*[X]\), a positive current on \(X\times X\) (Lemma~\ref{bou:integration-current}). Then
\(p_{1*}([\Delta]\wedge p_2^*\beta)=[X]\wedge\beta\) (for closed \(\beta\), the current of Lemma~\ref{bou:smooth-form-current}), and if \(\Theta-\delta[\Delta]\) is
positive and \(\beta\) is semipositive, then \(p_{1*}(\Theta\wedge p_2^*\beta)-\delta\,[X]\wedge\beta\) is positive.
\end{lemma}
\begin{proof}
\leanok
Linearity, weak continuity and positivity are those of push-forward and wedge (Lemma~\ref{bou:positive-currents}(3),(4); \(p_2^*\beta\) is semipositive when
\(\beta\) is). Since \(p_1\circ\operatorname{diag}=p_2\circ\operatorname{diag}=\mathrm{id}\), \(\operatorname{diag}^*(p_1^*\varphi\wedge p_2^*\beta)=\varphi\wedge\beta\) pointwise
(pull-backs compose), and a positive current depends only on the pointwise values of its test forms; this gives the identity, and the last claim is
positivity of \(p_{1*}((\Theta-\delta[\Delta])\wedge p_2^*\beta)\) rewritten with it.
\end{proof}

\begin{lemma}[Integration along the fibres over a graph]
\label{bou:graph-fibre-integration}
\uses{bou:fibre-pushforward,bou:integration-current,bou:smooth-form-current,bou:product-manifold,bou:positive-currents}
\source{dp04-kahler-cone:2.12,bou02-volume:4.8}
\lean{TranscendentalNakamaye.Pending.GraphFibreIntegration.contMDiff_graph,TranscendentalNakamaye.Pending.GraphFibreIntegration.pullback_graph_fibreTest_geometricallyEq,TranscendentalNakamaye.Pending.GraphFibreIntegration.fibrePushforward_graphCurrent,TranscendentalNakamaye.Pending.GraphFibreIntegration.fibrePushforward_graphCurrent_eq_smoothCurrent,TranscendentalNakamaye.Pending.GraphFibreIntegration.isPositive_fibrePushforward_sub_graph,TranscendentalNakamaye.Pending.GraphFibreIntegration.isKahler_add_of_isSemipositive,TranscendentalNakamaye.Pending.GraphFibreIntegration.graphKahler,TranscendentalNakamaye.Pending.GraphFibreIntegration.pullback_graph_prodForm_geometricallyEq,TranscendentalNakamaye.Pending.GraphFibreIntegration.mass_graphCurrent}
\leanok
Let \(f\colon V\to X\) be a holomorphic map between complex manifolds of the same dimension \(n=q+1\), \(V\) compact with a K\"ahler form \(w\), and
\([\Gamma_f]=(\mathrm{id},f)_*[V]\) the current of integration over the graph. For a smooth real \((1,1)\)-form \(\beta\) on \(X\),
\(p_{1*}([\Gamma_f]\wedge p_2^*\beta)=[V]\wedge f^*\beta\) (for closed \(\beta\), the current of the closed form \(f^*\beta\)); if \(\Theta-\delta[\Gamma_f]\) is positive and \(\beta\) is
semipositive, then \(p_{1*}(\Theta\wedge p_2^*\beta)-\delta\,[V]\wedge f^*\beta\) is positive. For a K\"ahler form \(\omega\) on \(X\), \(w+f^*\omega\) is K\"ahler and the mass
of \([\Gamma_f]\) for the product K\"ahler form \(p_1^*w+p_2^*\omega\) is \(\int_V(w+f^*\omega)^n\).
\end{lemma}
\begin{proof}
\leanok
As for the diagonal (Lemma~\ref{bou:fibre-pushforward}, the case \(f=\mathrm{id}\)): the pull-back of \(p_1^*\varphi\wedge p_2^*\beta\) by \(v\mapsto(v,f(v))\) is
\(\varphi\wedge f^*\beta\), and the pull-back of \((p_1^*w+p_2^*\omega)^n\) is \((w+f^*\omega)^n\), as forms acting on the tangent spaces of \(V\); positivity of
\(p_{1*}(\,\cdot\wedge p_2^*\beta)\) for semipositive \(\beta\) and linearity give the second claim.
\end{proof}

\begin{lemma}[The normal decomposition for an arbitrary linear projection]
\label{bou:graph-tangential-algebra}
\uses{bou:diagonal-tangential-reduction}
\source{demailly-cadg:III.2.10}
\lean{TranscendentalNakamaye.Pending.GraphTangentialAlgebra.halfVecOf,TranscendentalNakamaye.Pending.GraphTangentialAlgebra.halfVecOf_sum,TranscendentalNakamaye.Pending.GraphTangentialAlgebra.sub_compDiagProj_eq',TranscendentalNakamaye.Pending.GraphTangentialAlgebra.sub_comp_eq}
\leanok
Let \(\pi\) be a \(\mathbf C\)-linear endomorphism of \(\mathbf C^N\) with a normal decomposition \(\xi-\pi\xi=\sum_a\ell_a(\xi)N_a\) (real linear forms \(\ell_a\), vectors \(N_a\)).
For real \((1,1)\)-forms \(A\), \(B\) on \(\mathbf C^N\),
\[
 A-B\circ\pi=\sum_a\sum_cA(N_a,h_c)\,\Pi_{1,1}(\ell_a\wedge s_c)+\sum_{c,d}\tfrac12(A-B)(\pi E_c,\pi E_d)\,\Pi_{1,1}(s_c\wedge s_d),\qquad h_c=\pi E_c+\tfrac12(E_c-\pi E_c),
\]
with \((E_c)\) the real basis and \((s_c)\) its dual basis. For the diagonal projection this is the decomposition used in Lemma~\ref{bou:diagonal-tangential-reduction}.
\end{lemma}
\begin{proof}
\leanok
The proof of the diagonal case never uses the particular \(\pi\): expand \(A(\xi,\eta)-B(\pi\xi,\pi\eta)\) with \(\xi=\pi\xi+\sum_a\ell_a(\xi)N_a\) in both slots and collect.
\end{proof}

\begin{lemma}[The graph retraction in product charts]
\label{bou:graph-tangential-chart}
\uses{bou:graph-tangential-algebra,bou:normal-direction-vanishing,bou:product-manifold,bou:skoda-elmir-holomorphic}
\source{demailly-cadg:III.2.10}
\lean{TranscendentalNakamaye.Pending.GraphTangentialChart.baseCoord,TranscendentalNakamaye.Pending.GraphTangentialChart.baseCoordCLM,TranscendentalNakamaye.Pending.GraphTangentialChart.baseCoordCLM_apply,TranscendentalNakamaye.Pending.GraphTangentialChart.baseCoord_apply,TranscendentalNakamaye.Pending.GraphTangentialChart.chartSmooth_chartExtend_mul_of_isOpen,TranscendentalNakamaye.Pending.GraphTangentialChart.contMDiff_graphProjMap,TranscendentalNakamaye.Pending.GraphTangentialChart.continuous_baseCoord,TranscendentalNakamaye.Pending.GraphTangentialChart.coordinateMap_graphProjMap,TranscendentalNakamaye.Pending.GraphTangentialChart.differentiableAt_graphNormalFun,TranscendentalNakamaye.Pending.GraphTangentialChart.differentiableAt_graphRep,TranscendentalNakamaye.Pending.GraphTangentialChart.exactPair_cutoff_coeff_of_hasFDerivAt,TranscendentalNakamaye.Pending.GraphTangentialChart.graphChart,TranscendentalNakamaye.Pending.GraphTangentialChart.graphChart_mapsTo,TranscendentalNakamaye.Pending.GraphTangentialChart.graphChart_val,TranscendentalNakamaye.Pending.GraphTangentialChart.graphDiff,TranscendentalNakamaye.Pending.GraphTangentialChart.graphDiff_apply,TranscendentalNakamaye.Pending.GraphTangentialChart.graphDomain,TranscendentalNakamaye.Pending.GraphTangentialChart.graphDomain_subset_target,TranscendentalNakamaye.Pending.GraphTangentialChart.graphModel,TranscendentalNakamaye.Pending.GraphTangentialChart.graphModel_eq_of_mem_graph,TranscendentalNakamaye.Pending.GraphTangentialChart.graphModel_mem_target,TranscendentalNakamaye.Pending.GraphTangentialChart.graphNormalCoord,TranscendentalNakamaye.Pending.GraphTangentialChart.graphNormalFun,TranscendentalNakamaye.Pending.GraphTangentialChart.graphNormalFun_eq_zero_of_graphModel_eq,TranscendentalNakamaye.Pending.GraphTangentialChart.graphNormalRe,TranscendentalNakamaye.Pending.GraphTangentialChart.graphNormalRe_apply,TranscendentalNakamaye.Pending.GraphTangentialChart.graphProj,TranscendentalNakamaye.Pending.GraphTangentialChart.graphProjMap,TranscendentalNakamaye.Pending.GraphTangentialChart.graphProjMap_eq_self,TranscendentalNakamaye.Pending.GraphTangentialChart.graphProj_apply,TranscendentalNakamaye.Pending.GraphTangentialChart.graphRep,TranscendentalNakamaye.Pending.GraphTangentialChart.hasFDerivAt_graphModel,TranscendentalNakamaye.Pending.GraphTangentialChart.hasFDerivAt_graphNormalFun,TranscendentalNakamaye.Pending.GraphTangentialChart.hasFDerivAt_re_graphNormalFun,TranscendentalNakamaye.Pending.GraphTangentialChart.isHoloPairOn_re_im,TranscendentalNakamaye.Pending.GraphTangentialChart.isOpen_graphDomain,TranscendentalNakamaye.Pending.GraphTangentialChart.pullback_graphProjMap_coeff,TranscendentalNakamaye.Pending.GraphTangentialChart.sub_graphProj_eq}
\leanok
Let \(f\colon V\to X\) be holomorphic between complex manifolds of the same dimension \(n\) and \(\Phi(v,x)=(v,f(v))\) the holomorphic retraction of \(V\times X\) onto
the graph \(\Gamma_f\). In the product chart at a graph point, with \(F\) the coordinate representative of \(f\), \(\Phi\) reads \((z,w)\mapsto(z,F(z))\), its derivative is the
\(\mathbf C\)-linear projection \(\pi_\zeta(\xi,\eta)=(\xi,DF(z)\xi)\), so \((\Phi^*b)(\zeta)=b(\Phi\zeta)\circ\pi_\zeta\); \(\xi-\pi_\zeta\xi=\sum_a\ell_a(\xi)N_a\) with the normal vectors
of the diagonal case and \(\ell_a=\operatorname{Re}dL_a(\zeta)\), where \(L_a\) are holomorphic normal coordinates of the graph (the components of \(w-F(z)\) times \(1\) and
\(-i\)); on the graph \(\Phi\zeta=\zeta\) and \(L_a(\zeta)=0\). Moreover, in any chart, cut-offs of the real and imaginary parts of a holomorphic function form a holomorphic
pair (Lemma~\ref{bou:normal-direction-vanishing}), and \(d(\chi g\wedge d\chi k)=dg\wedge dk\) near a point where the cut-off \(\chi\) equals \(1\).
\end{lemma}
\begin{proof}
\leanok
Direct computation in the product chart; the derivative of \((z,w)\mapsto(z,F(z))\) is \((\xi,\eta)\mapsto(\xi,DF(z)\xi)\), and \(\xi-\pi_\zeta\xi=(0,\eta-DF(z)\xi)=\sum_a\ell_a(\xi)N_a\)
with \(\ell_a\) the real parts of the differentials of the components of \(w-F(z)\).
\end{proof}

\begin{lemma}[\(g\times\mathrm{id}\) is an isomorphism over \(N\times W\)]
\label{bou:product-map-isomorphism-over}
\uses{bou:product-manifold,ana:blowups}
\lean{TranscendentalNakamaye.Pending.ProductMapIsomorphismOver.contMDiffAt_prodMk,TranscendentalNakamaye.Pending.ProductMapIsomorphismOver.isIsomorphismOver_prodMap,TranscendentalNakamaye.Pending.ProductMapIsomorphismOver.not_mem_exceptionalLocus_prodMap}
\leanok
Let \(g\colon Z\to Y\) be holomorphic between complex manifolds of the same dimension. If \(dg\) is bijective at \(z\), then \(d(g\times\mathrm{id}_X)\) is bijective at
\((z,w)\); hence if \(g\) is an isomorphism over an open \(N\subset Y\), then \(g\times\mathrm{id}_X\) is an isomorphism over \(N\times W\) for every open \(W\subset X\).
\end{lemma}
\begin{proof}
\leanok
A holomorphic local inverse \(s\) of \(g\) near \(g(z)\) gives the local inverse \(s\times\mathrm{id}\) of \(g\times\mathrm{id}\), so \(d(g\times\mathrm{id})\) is onto at \((z,w)\), hence bijective
(equal dimensions); the remaining clauses of ``isomorphism over'' hold factorwise.
\end{proof}

\begin{lemma}[The tangential reduction on a graph]
\label{bou:graph-tangential-reduction}
\uses{bou:graph-tangential-chart,bou:graph-tangential-algebra,bou:diagonal-tangential-reduction,bou:normal-direction-vanishing,bou:carried-current-tests,bou:product-manifold}
\source{demailly-cadg:III.2.10}
\lean{TranscendentalNakamaye.Pending.GraphTangentialReduction.apply_elemTest_comp_graphProjMap,TranscendentalNakamaye.Pending.GraphTangentialReduction.apply_elemTest_comp_perm,TranscendentalNakamaye.Pending.GraphTangentialReduction.apply_elemTest_pullSlots,TranscendentalNakamaye.Pending.GraphTangentialReduction.apply_elemTest_update_pullback,TranscendentalNakamaye.Pending.GraphTangentialReduction.apply_pullback_graphProjMap_eq,TranscendentalNakamaye.Pending.GraphTangentialReduction.baseCoord_mem_coordinateDomain,TranscendentalNakamaye.Pending.GraphTangentialReduction.contDiffOn_fderiv_graphRep_apply,TranscendentalNakamaye.Pending.GraphTangentialReduction.contDiffOn_graphModel,TranscendentalNakamaye.Pending.GraphTangentialReduction.contDiffOn_graphProj_apply,TranscendentalNakamaye.Pending.GraphTangentialReduction.contDiffOn_im_graphNormalFun,TranscendentalNakamaye.Pending.GraphTangentialReduction.contDiffOn_re_graphNormalFun,TranscendentalNakamaye.Pending.GraphTangentialReduction.contDiffOn_twoForm_apply,TranscendentalNakamaye.Pending.GraphTangentialReduction.elemTest_congr,TranscendentalNakamaye.Pending.GraphTangentialReduction.exists_nhds_apply_wedge_sub_pullback_eq_zero,TranscendentalNakamaye.Pending.GraphTangentialReduction.graphProj_apply',TranscendentalNakamaye.Pending.GraphTangentialReduction.isClosed_graph,TranscendentalNakamaye.Pending.GraphTangentialReduction.isTypeOneOne_sub_pullback_graph,TranscendentalNakamaye.Pending.GraphTangentialReduction.mem_graphDomain_self,TranscendentalNakamaye.Pending.GraphTangentialReduction.pullSlotsGraph,TranscendentalNakamaye.Pending.GraphTangentialReduction.twoForm_apply_eq_sum}
\leanok
Let \(f\colon V\to X\) be holomorphic between compact complex manifolds of the same dimension, \(\Gamma_f\) its graph and \(\Phi(v,x)=(v,f(v))\). For a closed positive current
\(\Theta\) on \(V\times X\) whose trace measure is carried by \(\Gamma_f\), \(\Theta(\Phi^*\varphi)=\Theta(\varphi)\) for every test form \(\varphi\); that is,
\(\Theta=(\mathrm{id},f)_*p_{1*}\Theta\).
\end{lemma}
\begin{proof}
\leanok
As in Lemma~\ref{bou:diagonal-tangential-reduction}, with the point-dependent projection \(\pi_\zeta\) of Lemma~\ref{bou:graph-tangential-chart}: by
Lemma~\ref{bou:graph-tangential-algebra}, near a graph point \(b-\Phi^*b\) is a sum of terms \(\gamma\,\Pi_{1,1}d(\operatorname{Re}L_a\wedge ds_c)\), with holomorphic normal
coordinates \(L_a\), killed by Lemma~\ref{bou:normal-direction-vanishing}, and terms whose coefficients vanish on \(\Gamma_f\), killed by Lemma~\ref{bou:carried-current-tests};
the slots are pulled back one at a time and the function \(h\) is replaced by \(h\circ\Phi\), which agrees with \(h\) on \(\Gamma_f\).
\end{proof}

\begin{lemma}[The support theorem on the graph of a holomorphic map]
\label{bou:graph-support}
\uses{bou:graph-tangential-reduction,bou:skoda-elmir-holomorphic,bou:closed-top-currents,bou:fibre-pushforward,bou:integration-current,bou:current-restriction}
\source{demailly-cadg:III.2.10,demailly-cadg:III.2.13,demailly-cadg:III.2.14}
\lean{TranscendentalNakamaye.Pending.GraphSupport.eq_smul_graphCurrent,TranscendentalNakamaye.Pending.GraphSupport.exists_restrictCurrent_graph_eq_smul,TranscendentalNakamaye.Pending.GraphSupport.graphCurrent,TranscendentalNakamaye.Pending.GraphSupport.graphCurrent_apply,TranscendentalNakamaye.Pending.GraphSupport.isPositive_graphCurrent}
\leanok
Let \(V\), \(X\) be compact complex manifolds of dimension \(n\), \(V\) connected with a K\"ahler form, \(f\colon V\to X\) holomorphic and \([\Gamma_f]=(\mathrm{id},f)_*[V]\) (a
positive current). For a closed positive current \(\Theta\) of bidimension \((n,n)\) on \(V\times X\): if \(\Theta\) is carried by \(\Gamma_f\), then \(\Theta=c\,[\Gamma_f]\) with
\(c\ge0\); in general \(\mathbf 1_{\Gamma_f}\Theta=c\,[\Gamma_f]\) with \(c\ge0\) and \(\Theta-c\,[\Gamma_f]=\mathbf 1_{V\times X\setminus\Gamma_f}\Theta\) is positive.
\end{lemma}
\begin{proof}
\leanok
\(p_{1*}\Theta\) is closed positive of top bidimension on the connected \(V\), hence \(c\,[V]\) (Lemma~\ref{bou:closed-top-currents}); by
Lemma~\ref{bou:graph-tangential-reduction}, \(\Theta(\varphi)=\Theta(p_1^*(\mathrm{id},f)^*\varphi)=c\,[V]((\mathrm{id},f)^*\varphi)=c\,[\Gamma_f](\varphi)\). For a general \(\Theta\),
\(\mathbf 1_{\Gamma_f}\Theta\) is closed by Lemma~\ref{bou:skoda-elmir-holomorphic} and carried by \(\Gamma_f\).
\end{proof}

\begin{lemma}[The mass of the diagonal current]
\label{bou:diagonal-current-mass}
\uses{bou:fibre-pushforward,bou:product-manifold,bou:integration-current}
\source{bou02-volume:4.10,dp04-kahler-cone:2.12}
\lean{TranscendentalNakamaye.Pending.DiagonalFibreIntegration.formPower_smul,TranscendentalNakamaye.Pending.DiagonalFibreIntegration.formPow_geometricallyEq,TranscendentalNakamaye.Pending.DiagonalFibreIntegration.pullback_formPow,TranscendentalNakamaye.Pending.DiagonalFibreIntegration.pullback_diag_prodForm_geometricallyEq,TranscendentalNakamaye.Pending.DiagonalFibreIntegration.mass_diagonalCurrent}
\leanok
Let \(X\) be a complex manifold of dimension \(n\) and \(\operatorname{diag}\colon X\to X\times X\) the diagonal map. For smooth closed real \((1,1)\)-forms \(a\), \(b\)
on \(X\), \(\operatorname{diag}^*(p_1^*a+p_2^*b)=a+b\). If \(X\) is compact and \(\omega\) is a K\"ahler form on \(X\), the mass of the diagonal current
\([\Delta]=\operatorname{diag}_*[X]\) for the K\"ahler form \(\widetilde\omega=p_1^*\omega+p_2^*\omega\) is
\([\Delta](\widetilde\omega^{\,n})=2^n\int_X\omega^n\).
\end{lemma}
\begin{proof}
\leanok
Pull-backs compose and \(p_1\circ\operatorname{diag}=p_2\circ\operatorname{diag}=\mathrm{id}\), which gives the first identity pointwise; pull-back commutes with
wedge powers, so \(\operatorname{diag}^*\widetilde\omega^{\,n}=(2\omega)^n=2^n\omega^n\) pointwise. A positive current depends only on the pointwise values of
its test forms (Lemma~\ref{bou:fibre-pushforward}), so \([\Delta](\widetilde\omega^{\,n})=[X](2^n\omega^n)=2^n\int_X\omega^n\).
\end{proof}

\begin{lemma}[Integration along the fibres of a smooth product current]
\label{bou:product-fibre-integration}
\uses{bou:fibre-pushforward,bou:product-mixed-measure,bou:smooth-form-current,bou:integration-current,bou:product-manifold,bou:smooth-mixed-density}
\source{dp04-kahler-cone:2.12,bou02-volume:4.8}
\lean{TranscendentalNakamaye.Pending.FibrePushforwardProductForm.prodModelLinearEquiv,TranscendentalNakamaye.Pending.FibrePushforwardProductForm.scaleFstLinear,TranscendentalNakamaye.Pending.FibrePushforwardProductForm.scaleFst,TranscendentalNakamaye.Pending.FibrePushforwardProductForm.finrank_complexModel,TranscendentalNakamaye.Pending.FibrePushforwardProductForm.det_scaleFst,TranscendentalNakamaye.Pending.FibrePushforwardProductForm.scaleFst_apply,TranscendentalNakamaye.Pending.FibrePushforwardProductForm.projFst_comp_scaleFst,TranscendentalNakamaye.Pending.FibrePushforwardProductForm.projSnd_comp_scaleFst,TranscendentalNakamaye.Pending.FibrePushforwardProductForm.compContinuousLinearMap_smul_two,TranscendentalNakamaye.Pending.FibrePushforwardProductForm.compContinuousLinearMap_comp,TranscendentalNakamaye.Pending.FibrePushforwardProductForm.comp_projFst_comp_scaleFst,TranscendentalNakamaye.Pending.FibrePushforwardProductForm.comp_projSnd_comp_scaleFst,TranscendentalNakamaye.Pending.FibrePushforwardProductForm.actualMixedTopDensity_eq_zero_of_scaleFst,TranscendentalNakamaye.Pending.FibrePushforwardProductForm.append_update_right,TranscendentalNakamaye.Pending.FibrePushforwardProductForm.appendRight,TranscendentalNakamaye.Pending.FibrePushforwardProductForm.appendRight_apply,TranscendentalNakamaye.Pending.FibrePushforwardProductForm.actualMixedTopDensity_fibreFamily_single,TranscendentalNakamaye.Pending.FibrePushforwardProductForm.actualMixedTopDensity_fibreFamily,TranscendentalNakamaye.Pending.FibrePushforwardProductForm.coordChain,TranscendentalNakamaye.Pending.FibrePushforwardProductForm.coordChain_succ_apply,TranscendentalNakamaye.Pending.FibrePushforwardProductForm.coordChain_wedgeProd,TranscendentalNakamaye.Pending.FibrePushforwardProductForm.chainTestLM_coeff_eq,TranscendentalNakamaye.Pending.FibrePushforwardProductForm.wedgeProd_compContinuousLinearMap,TranscendentalNakamaye.Pending.FibrePushforwardProductForm.coordinateWedge_zero_left,TranscendentalNakamaye.Pending.FibrePushforwardProductForm.wedgeProdSpan,TranscendentalNakamaye.Pending.FibrePushforwardProductForm.wedgeProd_mem_wedgeProdSpan,TranscendentalNakamaye.Pending.FibrePushforwardProductForm.coordinateWedge_mem_wedgeProdSpan,TranscendentalNakamaye.Pending.FibrePushforwardProductForm.coeff_mem_wedgeProdSpan,TranscendentalNakamaye.Pending.FibrePushforwardProductForm.topDensity_fibreChain,TranscendentalNakamaye.Pending.FibrePushforwardProductForm.prodForm_coeff,TranscendentalNakamaye.Pending.FibrePushforwardProductForm.topDensity_formPow_prodForm,TranscendentalNakamaye.Pending.FibrePushforwardProductForm.topRatio_fibreChain,TranscendentalNakamaye.Pending.FibrePushforwardProductForm.fibreIntegral,TranscendentalNakamaye.Pending.FibrePushforwardProductForm.kahlerVolume_prodKahlerBackground,TranscendentalNakamaye.Pending.FibrePushforwardProductForm.fibreIntegral_eq_globalMixedMeasure,TranscendentalNakamaye.Pending.FibrePushforwardProductForm.fibrePushforward_smoothCurrent_prodForm}
\leanok
Let \(X\), \(Y\) be compact complex manifolds of dimensions \(q+1\) and \(p+1\) with K\"ahler forms, \(a\), \(b\) smooth closed real \((1,1)\)-forms on \(X\), \(Y\) (no
positivity) and \(\beta\) a smooth real \((1,1)\)-form on \(Y\). For the smooth current \(\Theta=[X\times Y]\wedge(p_1^*a+p_2^*b)^{p+1}\) of bidimension \((q+1,q+1)\),
\[
 p_{1*}(\Theta\wedge p_2^*\beta)\ =\ (p+1)\Bigl(\int_Yb^p\wedge\beta\Bigr)\,[X]\wedge a
\]
as currents on \(X\) (Lemma~\ref{bou:fibre-pushforward}); for semipositive closed \(b\), \(\beta\) the constant \(\int_Yb^p\wedge\beta\) is the total mixed mass of
\((b,\dots,b,\beta)\). With \(Y=X\), \(b=a\) and \(\beta=\widetilde\omega\): \(R=p_{1*}(\Theta\wedge p_2^*\widetilde\omega)=n\,(a^{n-1}\cdot\{\widetilde\omega\})\,[X]\wedge a\).
\end{lemma}
\begin{proof}
\leanok
In product coordinates \(\mathbf C^{q+1}\times\mathbf C^{p+1}\), the density of \(p_1^*\alpha_1\wedge\dots\wedge p_1^*\alpha_q\wedge p_2^*\beta\wedge(p_1^*A+p_2^*B)^{p+1}\) is
\((p+1)\,D(\alpha,A)\,D(B^p,\beta)\): expand the \(p+1\) sum slots multilinearly; a family with \(k\ne q+1\) slots pulled back from the first factor has density
\(0\) (pull back along \((x,y)\mapsto(tx,y)\): the density is multiplied by \(t^{2(q+1)}\) and by \(t^{2k}\)); the \(p+1\) terms with exactly one first-factor
sum slot are equal by a transposition of slots and factor as block families (Lemma~\ref{bou:product-mixed-measure}). Every test form has, at every point,
coefficients in the span of products of real \((1,1)\)-forms, so the identity of densities holds on all test forms; the K\"ahler volume of \(X\times Y\) is
\(\binom{n+m}n\,\mathrm{vol}_X\otimes\mathrm{vol}_Y\) (Lemma~\ref{bou:product-mixed-measure}), and Fubini gives the identity of currents.
\end{proof}

\begin{lemma}[A uniform \(L^1\) bound for normalized \(\theta\)-psh functions]
\label{bou:theta-psh-l1-bound}
\uses{posclass:potential-shift,dist:psh-representative}
\source{demailly-cadg:I.4.17}
\lean{TranscendentalNakamaye.Pending.ThetaPshUniformL1Bound.HarnackBall,TranscendentalNakamaye.Pending.ThetaPshUniformL1Bound.HarnackBall.set,TranscendentalNakamaye.Pending.ThetaPshUniformL1Bound.HarnackBall.isOpen_set,TranscendentalNakamaye.Pending.ThetaPshUniformL1Bound.HarnackBall.ball_subset_closedBall,TranscendentalNakamaye.Pending.ThetaPshUniformL1Bound.HarnackBall.ball_subset_W,TranscendentalNakamaye.Pending.ThetaPshUniformL1Bound.HarnackBall.ball_subset_target,TranscendentalNakamaye.Pending.ThetaPshUniformL1Bound.HarnackBall.symm_mem_set,TranscendentalNakamaye.Pending.ThetaPshUniformL1Bound.HarnackBall.set_nonempty,TranscendentalNakamaye.Pending.ThetaPshUniformL1Bound.HarnackBall.exists_bounds,TranscendentalNakamaye.Pending.ThetaPshUniformL1Bound.HarnackBall.corrected,TranscendentalNakamaye.Pending.ThetaPshUniformL1Bound.HarnackBall.discPshOn_corrected,TranscendentalNakamaye.Pending.ThetaPshUniformL1Bound.HarnackBall.deficit_le,TranscendentalNakamaye.Pending.ThetaPshUniformL1Bound.HarnackBall.neg_le_deficit,TranscendentalNakamaye.Pending.ThetaPshUniformL1Bound.HarnackBall.exists_point,TranscendentalNakamaye.Pending.ThetaPshUniformL1Bound.HarnackBall.transfer,TranscendentalNakamaye.Pending.ThetaPshUniformL1Bound.exists_harnackBall,TranscendentalNakamaye.Pending.ThetaPshUniformL1Bound.exists_forall_mem_set_le,TranscendentalNakamaye.Pending.ThetaPshUniformL1Bound.HarnackBall.exists_lintegral_le,TranscendentalNakamaye.Pending.ThetaPshUniformL1Bound.lintegral_biUnion_finset_le,TranscendentalNakamaye.Pending.ThetaPshUniformL1Bound.exists_lintegral_compact_le}
\leanok
Let \(M\) be a compact \emph{connected} complex manifold, \(\theta\) a smooth real two-form with smooth local potentials, \(e\) a chart and \(K\) a compact subset
of its target, and \(c\in\mathbf R\). There is \(B<\infty\) such that every \(\theta\)-psh function \(u\le0\) with \(u(x_0)\ge-c\) at some point \(x_0\in M\) satisfies
\(\int_K|u\circ e^{-1}|\,d\lambda\le B\).
\end{lemma}
\begin{proof}
\leanok
In a chart ball \(B(z,r)\) with a smooth local potential \(\rho\) of \(\theta\) on \(\overline B(z,3r)\), \(v=u\circ e^{-1}+\rho\) is psh and \(\le A=\max\rho\), and the
Harnack comparison between the mean and the supremum bounds \(\int_{B(z,r)}(A-v)^+\) by \(\lambda(B(0,2r))\,(A+c-\min\rho)\) as soon as \(u\ge-c\) at one point of the
ball. A bound on that integral gives, in every nonempty open subset of the ball, a point where \(u\ge-c'\) (first moment), so control in one ball gives control in
every ball meeting it. The relation ``control near \(x\) gives control near \(y\)'' is therefore transitive and holds for \(y\) near \(x\); \(M\) is connected, so it
holds for all \(x\), \(y\), uniformly by a finite subcover. (Without connectedness the statement fails: \(u=0\) on one component and \(u=-k\) on another.)
\end{proof}

\begin{lemma}[Weak limits of locally \(L^1\)-bounded plurisubharmonic functions]
\label{bou:psh-weak-compactness}
\uses{dist:regular,dist:psh-representative,dist:psh-representative-unique}
\source{demailly-cadg:I.4.20,demailly-cadg:I.5.8}
\lean{TranscendentalNakamaye.Pending.PshUltrafilterLimit.supNorm,TranscendentalNakamaye.Pending.PshUltrafilterLimit.abs_apply_le_supNorm,TranscendentalNakamaye.Pending.PshUltrafilterLimit.abs_integral_le,TranscendentalNakamaye.Pending.PshUltrafilterLimit.exists_tendsto_regularDistribution,TranscendentalNakamaye.Pending.PshUltrafilterLimit.tendsto_distributionalLeviCLM,TranscendentalNakamaye.Pending.PshUltrafilterLimit.HasNonnegativeDistributionalLevi.of_tendsto,TranscendentalNakamaye.Pending.PshUltrafilterLimit.isPshRepresentative_pshRegularDistribution,TranscendentalNakamaye.Pending.PshUltrafilterLimit.exists_isPshRepresentative_of_tendsto}
\leanok
Let \(\Omega\) be an open subset of a finite-dimensional space. A family of locally integrable functions \(f_i\) with \(\sup_i\int_K|f_i|<\infty\) for every compact
\(K\subseteq\Omega\) has, along every ultrafilter on the index set, a limit distribution \(T\): \(\int\varphi f_i\to T(\varphi)\) for every test function \(\varphi\). If the
\(f_i\) are plurisubharmonic on a preconnected \(\Omega\) and not identically \(-\infty\), the limit has nonnegative Levi form and hence a plurisubharmonic
representative (CADG~I.5.8).
\end{lemma}
\begin{proof}
\leanok
Bounded real families converge along ultrafilters, so \(T(\varphi)=\lim\int\varphi f_i\) exists, is linear, and \(|T(\varphi)|\le B_K\|\varphi\|_\infty\) on test functions
supported in \(K\), which is the continuity of a distribution. Nonnegativity of \(\int f_i\,L\varphi\) for nonnegative \(\varphi\) (the Levi form of a psh function) passes to
the limit, and a distribution with nonnegative Levi form is represented by a plurisubharmonic function.
\end{proof}

\begin{lemma}[Chart weak limits glue into one \(\theta\)-psh potential]
\label{bou:chart-limit-gluing}
\uses{bou:psh-weak-compactness,dist:psh-representative,dist:psh-representative-unique,posclass:potential-shift}
\source{demailly-cadg:I.5.8}
\lean{TranscendentalNakamaye.Pending.ThetaPshChartLimitGluing.integrable_testFunction_mul,TranscendentalNakamaye.Pending.ThetaPshChartLimitGluing.regularDistribution_add,TranscendentalNakamaye.Pending.ThetaPshChartLimitGluing.regularDistribution_sub,TranscendentalNakamaye.Pending.ThetaPshChartLimitGluing.restrictDistribution_add,TranscendentalNakamaye.Pending.ThetaPshChartLimitGluing.distribution_eq_of_restrict_eq,TranscendentalNakamaye.Pending.ThetaPshChartLimitGluing.ereal_add_coe_neg_add_coe,TranscendentalNakamaye.Pending.ThetaPshChartLimitGluing.ereal_eq_add_coe_neg_of_add_coe_eq,TranscendentalNakamaye.Pending.ThetaPshChartLimitGluing.ereal_add_coe_finite,TranscendentalNakamaye.Pending.ThetaPshChartLimitGluing.ereal_toReal_add_coe,TranscendentalNakamaye.Pending.ThetaPshChartLimitGluing.contDiff_det,TranscendentalNakamaye.Pending.ThetaPshChartLimitGluing.transitionJac,TranscendentalNakamaye.Pending.ThetaPshChartLimitGluing.contDiffOn_transitionJac,TranscendentalNakamaye.Pending.ThetaPshChartLimitGluing.pullFun,TranscendentalNakamaye.Pending.ThetaPshChartLimitGluing.pullFun_of_mem,TranscendentalNakamaye.Pending.ThetaPshChartLimitGluing.pullFun_of_notMem,TranscendentalNakamaye.Pending.ThetaPshChartLimitGluing.image_tsupport_subset_source,TranscendentalNakamaye.Pending.ThetaPshChartLimitGluing.isCompact_image_tsupport,TranscendentalNakamaye.Pending.ThetaPshChartLimitGluing.pullFun_eq_zero_of_notMem_image,TranscendentalNakamaye.Pending.ThetaPshChartLimitGluing.tsupport_pullFun_subset,TranscendentalNakamaye.Pending.ThetaPshChartLimitGluing.contDiff_pullFun,TranscendentalNakamaye.Pending.ThetaPshChartLimitGluing.pullTest,TranscendentalNakamaye.Pending.ThetaPshChartLimitGluing.pullTest_apply,TranscendentalNakamaye.Pending.ThetaPshChartLimitGluing.image_tsupport_subset_chartTarget,TranscendentalNakamaye.Pending.ThetaPshChartLimitGluing.integral_mul_eq_integral_pullFun,TranscendentalNakamaye.Pending.ThetaPshChartLimitGluing.IsChartLimit,TranscendentalNakamaye.Pending.ThetaPshChartLimitGluing.chartTransition_apply_eq,TranscendentalNakamaye.Pending.ThetaPshChartLimitGluing.symm_chartTransition_apply,TranscendentalNakamaye.Pending.ThetaPshChartLimitGluing.IsChartLimit.apply_pullTest,TranscendentalNakamaye.Pending.ThetaPshChartLimitGluing.chartTransition_target_eq_source,TranscendentalNakamaye.Pending.ThetaPshChartLimitGluing.chartTransition_symm_eq,TranscendentalNakamaye.Pending.ThetaPshChartLimitGluing.IsChartLimit.apply_eq_integral_comp,TranscendentalNakamaye.Pending.ThetaPshChartLimitGluing.isSmoothLocalPotential_comp_transition,TranscendentalNakamaye.Pending.ThetaPshChartLimitGluing.plurisubharmonicOn_transport,TranscendentalNakamaye.Pending.ThetaPshChartLimitGluing.LocalRep,TranscendentalNakamaye.Pending.ThetaPshChartLimitGluing.LocalRep.restrict_eq,TranscendentalNakamaye.Pending.ThetaPshChartLimitGluing.LocalRep.eq_transport,TranscendentalNakamaye.Pending.ThetaPshChartLimitGluing.IsLocallyThetaPshDistribution,TranscendentalNakamaye.Pending.ThetaPshChartLimitGluing.IsLocallyThetaPshDistribution.exists_localRep,TranscendentalNakamaye.Pending.ThetaPshChartLimitGluing.exists_isThetaPshPotential_of_isChartLimit,TranscendentalNakamaye.Pending.ThetaPshChartLimitGluing.exists_isChartLimit,TranscendentalNakamaye.Pending.ThetaPshChartLimitGluing.exists_kahlerCurrent_of_isChartLimit}
\leanok
Let \(M\) be a complex manifold, \(\theta\) a smooth closed real \((1,1)\)-form, \(v_i\) functions on \(M\) and \(l\) a nontrivial filter on the indices. Suppose that in
every chart \(e\) the functions \(v_i\circ e^{-1}\) converge weakly along \(l\) to a distribution \(T_e\) which is locally \(\theta\)-psh (near every point,
\(T_e+\rho\) has nonnegative Levi form for a smooth local potential \(\rho\) of \(\theta\)). Then there is a \(\theta\)-psh potential \(V\) on \(M\) with \(T_e=V\circ e^{-1}\)
(as distributions) in every chart. If \(\theta=\theta_0-\omega_\delta\) for a K\"ahler form \(\omega_\delta\), this is a K\"ahler current in \(\{\theta_0\}\) with lower bound
\(\omega_\delta\). Uniform bounds of \(\int_K|v_i\circ e^{-1}|\) on the compact subsets of every chart give the chart limits along any ultrafilter.
\end{lemma}
\begin{proof}
\leanok
Change of variables along a chart transition \(F\) (\(\int\psi\,h=\int(J\psi\circ F)(h\circ F)\), \(J=|\det DF|\) smooth) makes the chart limits compatible,
\(T_f(\psi)=T_e(J\psi\circ F)\). Near every point \(T_e+\rho\) has a plurisubharmonic representative \(u\) (CADG~I.5.8); for two charts the functions
\(u_d\circ F-\rho_d\circ F+\rho_e\) and \(u_e\) are psh and represent the same distribution (\(\rho_e-\rho_d\circ F\) is pluriharmonic), so they agree everywhere on a ball
(uniqueness of psh representatives). Hence \(V=u_x-\rho_x\) in the preferred chart at \(x\) is well defined and \(\theta\)-psh.
\end{proof}

\begin{lemma}[The integration current read in one chart]
\label{bou:integration-current-chart}
\uses{bou:integration-current,bou:closed-top-currents}
\lean{TranscendentalNakamaye.Pending.IntegrationCurrentChart.integrationCurrent_apply_eq_setIntegral,TranscendentalNakamaye.Pending.IntegrationCurrentChart.integrationCurrent_apply_eq_integral}
\leanok
Let \(X\) be a compact complex manifold of dimension \(n\ge1\) with a K\"ahler form \(\omega\), \(e\) a chart and \(K\) a measurable subset of the chart target.
If \(\varphi\) is a test form of degree \(2n\) with \(\varphi/\omega^n=0\) off \(e^{-1}(K)\), then \([X](\varphi)=\int_K D(\varphi_e)\,d\lambda\), \(D(\varphi_e)\) the top density of the
coefficient of \(\varphi\) in \(e\); if moreover \(D(\varphi_e)=0\) off \(K\), the integral may be taken over all of \(\mathbf C^n\).
\end{lemma}
\begin{proof}
\leanok
\([X](\varphi)=\int\varphi/\omega^n\,dV_\omega\) is an integral over \(e^{-1}(K)\); the change of variables for the K\"ahler volume in a chart
(the volume density \(D(\omega^n_e)>0\), proof of Lemma~\ref{bou:closed-top-currents}) turns it into \(\int_K(\varphi/\omega^n)(e^{-1}z)\,D(\omega^n_e)(z)\,d\lambda(z)\),
and \((\varphi/\omega^n)\,D(\omega^n_e)=D(\varphi_e)\) pointwise.
\end{proof}

\begin{lemma}[The directional Levi form as a top density]
\label{bou:directional-levi-density}
\uses{btgc:coefficient-congruence,fpma:quadratic-potential,fpma:geometric-density,levi:real-hessian,bou:closed-top-currents}
\lean{TranscendentalNakamaye.Pending.DirectionalLeviDensity.perpDiag,TranscendentalNakamaye.Pending.DirectionalLeviDensity.mixedCoefficient_cons_perpDiag,TranscendentalNakamaye.Pending.DirectionalLeviDensity.exists_unitary_col_zero,TranscendentalNakamaye.Pending.DirectionalLeviDensity.mixedCoefficient_cons_projection,TranscendentalNakamaye.Pending.DirectionalLeviDensity.perpQuadratic,TranscendentalNakamaye.Pending.DirectionalLeviDensity.contDiff_normSq_inner,TranscendentalNakamaye.Pending.DirectionalLeviDensity.contDiff_perpQuadratic,TranscendentalNakamaye.Pending.DirectionalLeviDensity.realHessian_normSq_inner,TranscendentalNakamaye.Pending.DirectionalLeviDensity.complexHessian_normSq_inner,TranscendentalNakamaye.Pending.DirectionalLeviDensity.complexHessian_perpQuadratic,TranscendentalNakamaye.Pending.DirectionalLeviDensity.topDensity_smoothDDC_perpQuadratic}
\leanok
Let \(a\in\mathbf C^{m+1}\) be a unit vector and \(q_a(z)=|z|^2-|\langle a,z\rangle|^2\). For \(f\) of class \(C^2\) near \(z\), the top density of
\(dd^cf\wedge(dd^cq_a)^m\) at \(z\) is \(\frac{m!}{\pi^{m+1}}\,L_f(z;a)\), \(L_f(z;a)\) the Levi form of \(f\) at \(z\) in the direction \(a\).
\end{lemma}
\begin{proof}
\leanok
The complex Hessian of \(q_a\) is the orthogonal projection \(P=1-y\,y^*\), \(y=\bar a\), onto \(y^\perp\). For \(D=\operatorname{diag}(0,1,\dots,1)\) the mixed
coefficient is \(C(K,D,\dots,D)=m!\,K_{00}\); a unitary matrix \(W\) with first column \(y\) conjugates \(D\) to \(P\), and the mixed coefficient is invariant under
unitary congruence, so \(C(H,P,\dots,P)=m!\,y^*Hy\), which for \(H\) the complex Hessian of \(f\) is \(m!\,L_f(z;a)\); the normalisation of \(dd^c\) gives the
factor \(\pi^{-(m+1)}\).
\end{proof}
These two lemmas are the flat core of the positivity-to-Levi step (Lemma~\ref{bou:positivity-levi-link}): tested on
\(\chi\,(dd^cq_a)^{n-1}\) in a chart, a current \(R=\theta+dd^cV\) gives \(\int\chi\,\theta\wedge(dd^cq_a)^{n-1}\) plus a multiple of \(\int V\,L_a\chi\,d\lambda\).

\begin{lemma}[From positivity of a \((1,1)\)-current to the Levi form of its chart potentials]
\label{bou:positivity-levi-link}
\uses{bou:directional-levi-density,bou:integration-current-chart,bou:chart-limit-gluing,bou:smooth-form-current,bou:positive-currents,bou:slot-ibp,dist:levi}
\source{demailly-cadg:I.5.8}
\lean{TranscendentalNakamaye.Pending.PositivityLeviLink.distributionalLeviCLM_restrictDistribution,TranscendentalNakamaye.Pending.PositivityLeviLink.distributionalLeviCLM_smul_apply,TranscendentalNakamaye.Pending.PositivityLeviLink.contDiff_cutQuadratic,TranscendentalNakamaye.Pending.PositivityLeviLink.tsupport_cutQuadratic_subset,TranscendentalNakamaye.Pending.PositivityLeviLink.chartSmooth_cutQuadratic,TranscendentalNakamaye.Pending.PositivityLeviLink.leviSlot,TranscendentalNakamaye.Pending.PositivityLeviLink.leviSlot_coeff_eq,TranscendentalNakamaye.Pending.PositivityLeviLink.leviForm_perpQuadratic,TranscendentalNakamaye.Pending.PositivityLeviLink.coefficientSemipositive_smoothDDC_perpQuadratic,TranscendentalNakamaye.Pending.PositivityLeviLink.isSemipositiveOneOne_leviSlot_coeff,TranscendentalNakamaye.Pending.PositivityLeviLink.chartSmooth_testFunction,TranscendentalNakamaye.Pending.PositivityLeviLink.leviSlotPow,TranscendentalNakamaye.Pending.PositivityLeviLink.leviTest,TranscendentalNakamaye.Pending.PositivityLeviLink.leviTest_coeff,TranscendentalNakamaye.Pending.PositivityLeviLink.leviTest_mem_positiveCone,TranscendentalNakamaye.Pending.PositivityLeviLink.actualMixedTopDensity_append_one,TranscendentalNakamaye.Pending.PositivityLeviLink.leviTest_eq_smulTest,TranscendentalNakamaye.Pending.PositivityLeviLink.smoothCurrent_leviTest,TranscendentalNakamaye.Pending.PositivityLeviLink.LeviTestIdentity,TranscendentalNakamaye.Pending.PositivityLeviLink.isLocallyThetaPshDistribution_of_leviTest,TranscendentalNakamaye.Pending.PositivityLeviLink.smoothCurrent_sub_apply,TranscendentalNakamaye.Pending.PositivityLeviLink.isLocallyThetaPshDistribution_sub_of_leviTest,TranscendentalNakamaye.Pending.PositivityLeviLink.exists_kahlerCurrent_of_leviTest,TranscendentalNakamaye.Pending.PositivityLeviLink.LeviTestIdentity.of_tendsto,TranscendentalNakamaye.Pending.PositivityLeviLink.LeviTestIdentity.of_isChartLimit}
\leanok
Let \(X\) be a compact complex manifold of dimension \(n=q+1\) with a K\"ahler form, \(\theta\) a smooth closed real \((1,1)\)-form with smooth local potentials,
\(P\) a positive \((1,1)\)-current on \(X\) and \(T_e\) a distribution on the target of every chart \(e\). For a chart \(e\), test functions \(\psi,\varphi\) on its target
with \(\varphi\ge0\) and \(\psi\equiv1\) near \(\operatorname{supp}\varphi\), and a unit vector \(a\), put \(\beta_a=dd^c((\psi q_a)\circ e)\) and
\(\eta=(\varphi\circ e)\,\beta_a^{\,q}\), a positive test form. Suppose (the \emph{Levi identity})
\[
 P(\eta)=\int_X\theta\wedge\eta+\frac{q!}{\pi^{q+1}}\,(L_aT_e)(\varphi)
\]
for all such \((e,\psi,\varphi,a)\), \(L_a\) the distributional Levi form in the direction \(a\). Then every \(T_e\) is locally \(\theta\)-psh: near every point
\(T_e+\rho\) has nonnegative Levi form for a smooth local potential \(\rho\) of \(\theta\). The same holds for \(P=R-[X]\wedge\sigma\) and \(\theta=\theta_0-\sigma\) when
\(R\) satisfies the identity with \(\theta_0\); if moreover \(\sigma=\omega_\delta\) is K\"ahler and the \(T_e\) are chart weak limits of \(v_i\circ e^{-1}\), there is a
K\"ahler current in \(\{\theta_0\}\) with lower bound \(\omega_\delta\) whose potential represents every \(T_e\). Finally, the Levi identity passes to limits: if
\(R_i\to R\) weakly, \(T_{e,i}\to T_e\) (or \(T_{e,i}=[v_i\circ e^{-1}]\) with chart limits \(T_e\)) and each \((R_i,T_{e,i})\) satisfies it with \(\theta\), so does \((R,T_e)\).
\end{lemma}
\begin{proof}
\leanok
Where \(\varphi\circ e\ne0\) the cut-off is \(\equiv1\), so \(\beta_a=dd^cq_a\) there, semipositive in every chart (its complex Hessian is a projection and chart
transitions are \(\mathbf C\)-linear to first order): \(\eta\) is positive. If \(\theta=dd^c\rho\) on \(\operatorname{supp}\varphi\), Lemmas~\ref{bou:integration-current-chart}
and~\ref{bou:directional-levi-density} give \(\int\theta\wedge\eta=\frac{q!}{\pi^{q+1}}\int\varphi\,L_\rho(a)\,d\lambda=\frac{q!}{\pi^{q+1}}(L_a\rho)(\varphi)\). Near a point take a local
potential \(\rho\) of \(\theta\), cut off; then \(0\le P(\eta)=\frac{q!}{\pi^{q+1}}\,L_a(T_e+\rho)(\varphi)\) for every unit \(a\) and every \(\varphi\ge0\) supported nearby,
which is the claim. For \(R-[X]\wedge\sigma\) subtract \(\int\sigma\wedge\eta\); the K\"ahler current is Lemma~\ref{bou:chart-limit-gluing}. Both sides of the
Levi identity are limits of the corresponding sides for \((R_i,T_{e,i})\) at every fixed test form, so the identity passes to the limit.
\end{proof}

\begin{lemma}[Fubini along the fibres of the first projection]
\label{bou:vertical-density-factorization}
\uses{bou:product-manifold,bou:product-mixed-measure,bou:integration-current,bou:smooth-form-current}
\source{dp04-kahler-cone:2.12,bou02-volume:4.8}
\lean{TranscendentalNakamaye.Pending.VerticalDensityFactorization.inclFst,TranscendentalNakamaye.Pending.VerticalDensityFactorization.inclSnd,TranscendentalNakamaye.Pending.VerticalDensityFactorization.inclFst_apply,TranscendentalNakamaye.Pending.VerticalDensityFactorization.inclSnd_apply,TranscendentalNakamaye.Pending.VerticalDensityFactorization.inclSnd_smul,TranscendentalNakamaye.Pending.VerticalDensityFactorization.scaleFst_eq,TranscendentalNakamaye.Pending.VerticalDensityFactorization.scaleFst_zero,TranscendentalNakamaye.Pending.VerticalDensityFactorization.compContinuousLinearMap_smul_add,TranscendentalNakamaye.Pending.VerticalDensityFactorization.isOneOne_comp_inclSnd,TranscendentalNakamaye.Pending.VerticalDensityFactorization.scaledFamily,TranscendentalNakamaye.Pending.VerticalDensityFactorization.continuous_scaledFamily,TranscendentalNakamaye.Pending.VerticalDensityFactorization.actualMixedTopDensity_scaledFamily,TranscendentalNakamaye.Pending.VerticalDensityFactorization.actualMixedTopDensity_append_comp_inclSnd,TranscendentalNakamaye.Pending.VerticalDensityFactorization.topDensity_coordChain_comp_projFst,TranscendentalNakamaye.Pending.VerticalDensityFactorization.vertForm,TranscendentalNakamaye.Pending.VerticalDensityFactorization.vertRatio,TranscendentalNakamaye.Pending.VerticalDensityFactorization.continuous_vertRatio,TranscendentalNakamaye.Pending.VerticalDensityFactorization.coeff_chain_pullback,TranscendentalNakamaye.Pending.VerticalDensityFactorization.topRatio_chain_pullback,TranscendentalNakamaye.Pending.VerticalDensityFactorization.fibrePotential,TranscendentalNakamaye.Pending.VerticalDensityFactorization.continuous_fibrePotential,TranscendentalNakamaye.Pending.VerticalDensityFactorization.integrationCurrent_chain_pullback}
\leanok
(1) In \(\mathbf C^{n+m}=\mathbf C^n\times\mathbf C^m\), for real two-forms \(P_1,\dots,P_n\) on \(\mathbf C^n\) and real \((1,1)\)-forms \(G_1,\dots,G_m\) on \(\mathbf C^{n+m}\),
the top-degree density satisfies \(D(p_1^*P_1,\dots,p_1^*P_n,G_1,\dots,G_m)=D(P)\cdot D(\iota_2^*G_1,\dots,\iota_2^*G_m)\), \(\iota_2(w)=(0,w)\).
(2) Let \(X\), \(Y\) be compact complex manifolds of dimensions \(n\), \(m\) with K\"ahler forms \(\omega_X\), \(\omega_Y\), \(G_1,\dots,G_m\) closed real \((1,1)\)-forms on
\(X\times Y\), \(\Lambda\) a top-degree test form on \(X\) and \(g\) a chart-smooth function on \(X\times Y\). With the vertical ratio
\(\operatorname{vert}=(G_1\wedge\dots\wedge G_m)|_{\text{fibre}}/\omega_Y^m\) and the fibre potential \(v(x)=\binom{n+m}n\int_Yg(x,y)\operatorname{vert}(x,y)\,dV_Y(y)\),
\[
 [X\times Y]\bigl(g\,p_1^*\Lambda\wedge G_1\wedge\dots\wedge G_m\bigr)=\int_Xv\,\Lambda,
\]
and \(v\) is continuous when \(g\) is.
\end{lemma}
\begin{proof}
\leanok
(1) The scaling \(S_t(x,y)=(tx,y)\) multiplies the left side by \(t^{2n}\) and each \(p_1^*P_i\) by \(t^2\), so \(t\mapsto D(p_1^*P,S_t^*G)\) is constant for \(t\ne0\);
it is continuous, and \(S_0=\iota_2\circ p_2\) gives a block family whose density factors (Lemma~\ref{bou:product-mixed-measure}). (2) By (1) the density of the
integrand in product charts is \(g\) times the density of \(\Lambda\) times the vertical ratio; Fubini on \(X\times Y\) with the product volume
(Lemma~\ref{bou:product-mixed-measure}) gives the identity, and continuity of \(v\) is continuity of a parametric integral over the compact \(Y\).
\end{proof}

\begin{lemma}[The pairing with \(dd^c\) of a chart test function is a Levi distribution]
\label{bou:fibre-potential-levi}
\uses{bou:directional-levi-density,bou:integration-current-chart,dist:levi}
\source{demailly-cadg:I.5.8}
\lean{TranscendentalNakamaye.Pending.FibrePotentialLevi.ddcTest,TranscendentalNakamaye.Pending.FibrePotentialLevi.leviDDCTest,TranscendentalNakamaye.Pending.FibrePotentialLevi.ddcTest_coeff_eq,TranscendentalNakamaye.Pending.FibrePotentialLevi.ddcTest_coeff_eq_zero,TranscendentalNakamaye.Pending.FibrePotentialLevi.topDensity_leviDDCTest_coeff,TranscendentalNakamaye.Pending.FibrePotentialLevi.leviDDCTest_coeff_eq_zero,TranscendentalNakamaye.Pending.FibrePotentialLevi.integral_mul_topRatio_leviDDCTest}
\leanok
Let \(X\) be a compact complex manifold of dimension \(n=q+1\) with a K\"ahler form, \(e\) a chart, \(\psi,\varphi\) test functions on its target with \(\psi\equiv1\)
near \(\operatorname{supp}\varphi\), \(a\) a unit vector, \(\beta_a=dd^c((\psi q_a)\circ e)\), and \(v\) a continuous function on \(X\). Then
\[
 \int_Xv\;dd^c(\varphi\circ e)\wedge\beta_a^{\,q}=\frac{q!}{\pi^{q+1}}\,L_a\bigl([v\circ e^{-1}]\bigr)(\varphi).
\]
\end{lemma}
\begin{proof}
\leanok
In the chart the density of \(dd^c(\varphi\circ e)\wedge\beta_a^{\,q}\) is \(\frac{q!}{\pi^{q+1}}L_\varphi(a)\) where \(\psi\equiv1\) (Lemma~\ref{bou:directional-levi-density} with
\(f=\varphi\)), and it vanishes off \(\operatorname{supp}\varphi\); Lemma~\ref{bou:integration-current-chart} reads the integral in the chart, and
\(\int(v\circ e^{-1})L_\varphi(a)\,d\lambda=L_a([v\circ e^{-1}])(\varphi)\) by definition of the distributional Levi form.
\end{proof}

\begin{lemma}[Stokes on \(X\times X\): the Levi identity for the fibre push-forward of a smooth product current]
\label{bou:product-stokes-levi}
\uses{bou:slot-ibp-compact,bou:vertical-density-factorization,bou:fibre-potential-levi,bou:positivity-levi-link,bou:product-fibre-integration,bou:fibre-pushforward,bou:smooth-form-current,bou:product-mixed-measure,bou:chart-limit-gluing}
\source{dp04-kahler-cone:2.12,bou02-volume:4.8}
\lean{TranscendentalNakamaye.Pending.ProductStokesLevi.teleSlots,TranscendentalNakamaye.Pending.ProductStokesLevi.actualMixedTopDensity_cons_removeNth,TranscendentalNakamaye.Pending.ProductStokesLevi.removeNth_append_update,TranscendentalNakamaye.Pending.ProductStokesLevi.actualMixedTopDensity_telescope,TranscendentalNakamaye.Pending.ProductStokesLevi.append_cons_cons,TranscendentalNakamaye.Pending.ProductStokesLevi.isClosedCurrent_integrationCurrent,TranscendentalNakamaye.Pending.ProductStokesLevi.liftTest,TranscendentalNakamaye.Pending.ProductStokesLevi.chartSmooth_liftTest,TranscendentalNakamaye.Pending.ProductStokesLevi.fibreFam,TranscendentalNakamaye.Pending.ProductStokesLevi.teleForms,TranscendentalNakamaye.Pending.ProductStokesLevi.fibreLeviTest,TranscendentalNakamaye.Pending.ProductStokesLevi.fibreFam_coeff,TranscendentalNakamaye.Pending.ProductStokesLevi.fibreLeviTest_coeff,TranscendentalNakamaye.Pending.ProductStokesLevi.topDensity_chain_fibreLeviTest,TranscendentalNakamaye.Pending.ProductStokesLevi.topDensity_slot_fibre,TranscendentalNakamaye.Pending.ProductStokesLevi.smoothCurrent_sub_eq_sum_slot,TranscendentalNakamaye.Pending.ProductStokesLevi.smoothDDC_liftTest,TranscendentalNakamaye.Pending.ProductStokesLevi.vertSlots,TranscendentalNakamaye.Pending.ProductStokesLevi.map_cons_comp,TranscendentalNakamaye.Pending.ProductStokesLevi.integrationCurrent_slot_ibp_eq_chain,TranscendentalNakamaye.Pending.ProductStokesLevi.fibreV,TranscendentalNakamaye.Pending.ProductStokesLevi.continuous_fibreV,TranscendentalNakamaye.Pending.ProductStokesLevi.locallyIntegrableOn_fibreV,TranscendentalNakamaye.Pending.ProductStokesLevi.smoothCurrent_smul_apply,TranscendentalNakamaye.Pending.ProductStokesLevi.leviTestIdentity_fibrePushforward,TranscendentalNakamaye.Pending.ProductStokesLevi.isLocallyThetaPshDistribution_fibreV}
\leanok
Let \(X\) be a compact complex manifold of dimension \(n=q+1\) with a K\"ahler form \(\omega\), \(b\) a closed real \((1,1)\)-form on \(X\), \(A=p_1^*b+p_2^*b\), \(g\) a
chart-smooth function on \(X\times X\) and \(\alpha=A+dd^cg\). Put \(R=p_{1*}([X\times X]\wedge\alpha^n\wedge p_2^*\omega)\), \(c=n\int_Xb^{n-1}\wedge\omega\), and let
\(v=\sum_{k=0}^{q}v_k\) be the continuous fibre potential, \(v_k\) the fibre potential (Lemma~\ref{bou:vertical-density-factorization}) of \(g\) against the vertical forms
\(p_2^*\omega,\alpha^k,A^{q-k}\). Then \(R\) satisfies the Levi identity of Lemma~\ref{bou:positivity-levi-link} with \(\theta=c\,b\) and \(T_e=[v\circ e^{-1}]\), with no positivity
assumption. If moreover \(\alpha\ge0\) and \(c\,b\) has smooth local potentials, every \([v\circ e^{-1}]\) is locally \(c\,b\)-psh.
\end{lemma}
\begin{proof}
\leanok
Pointwise, \(\alpha^n-A^n=\sum_{k=0}^{q}dd^cg\wedge\alpha^k\wedge A^{q-k}\) (telescoping in the mixed density, one slot at a time). Against the lifted Levi test form,
the \(A^n\) term gives \(c\,[X]\wedge b\) (Lemma~\ref{bou:product-fibre-integration}); each difference term is moved by integration by parts against the closed
positive current \([X\times X]\) (Lemma~\ref{bou:slot-ibp-compact}) to \(g\,dd^c(\varphi\circ e\circ p_1)\wedge\dots\), which Fubini
(Lemma~\ref{bou:vertical-density-factorization}) turns into \(\int_Xv_k\,dd^c(\varphi\circ e)\wedge\beta_a^{\,q}\), a Levi distribution by
Lemma~\ref{bou:fibre-potential-levi}. Summing gives the identity. Under \(\alpha\ge0\) the product current is positive and Lemma~\ref{bou:positivity-levi-link} applies.
\end{proof}

\begin{lemma}[A continuous function that is locally \(\theta\)-psh as a distribution is \(\theta\)-psh]
\label{bou:continuous-theta-psh}
\uses{bou:chart-limit-gluing,bou:psh-weak-compactness}
\source{demailly-cadg:I.5.8}
\lean{TranscendentalNakamaye.Pending.ContinuousThetaPsh.le_of_ae_le_of_semicontinuousOn,TranscendentalNakamaye.Pending.ContinuousThetaPsh.thetaPlurisubharmonic_of_continuous_of_isLocallyThetaPshDistribution}
\leanok
Let \(v\) be a continuous function on a complex manifold such that in every chart \(e\) the distribution \([v\circ e^{-1}]\) is locally \(\theta\)-psh (near every
point \([v\circ e^{-1}]+\rho\) has nonnegative distributional Levi form for a smooth local potential \(\rho\) of \(\theta\)). Then \(v\) is \(\theta\)-psh.
\end{lemma}
\begin{proof}
\leanok
Near a point, \([v\circ e^{-1}]+\rho\) is represented by a psh function \(u\) with \(u=v\circ e^{-1}+\rho\) almost everywhere (Lemma~\ref{bou:chart-limit-gluing}).
A psh function below an upper semicontinuous function almost everywhere is below it everywhere, and a lower semicontinuous function below an upper
semicontinuous one almost everywhere is below it everywhere (the failure set would contain a nonempty open set); as \(v\circ e^{-1}+\rho\) is continuous,
\(u=v\circ e^{-1}+\rho\) everywhere, so \(v\circ e^{-1}+\rho\) is psh.
\end{proof}

\begin{lemma}[Sup-normalized potentials: uniform bounds, chart limits, the limit K\"ahler current]
\label{bou:sup-normalized-limit}
\uses{bou:theta-psh-l1-bound,bou:psh-weak-compactness,bou:chart-limit-gluing,bou:positivity-levi-link,bou:continuous-theta-psh}
\source{dp04-kahler-cone:2.12,bou02-volume:4.8}
\lean{TranscendentalNakamaye.Pending.SupNormalizedPotentialLimit.supNormalize,TranscendentalNakamaye.Pending.SupNormalizedPotentialLimit.continuous_supNormalize,TranscendentalNakamaye.Pending.SupNormalizedPotentialLimit.supNormalize_nonpos,TranscendentalNakamaye.Pending.SupNormalizedPotentialLimit.exists_supNormalize_eq_zero,TranscendentalNakamaye.Pending.SupNormalizedPotentialLimit.thetaPlurisubharmonic_of_sub_isSemipositive,TranscendentalNakamaye.Pending.SupNormalizedPotentialLimit.thetaPlurisubharmonic_supNormalize,TranscendentalNakamaye.Pending.SupNormalizedPotentialLimit.locallyIntegrableOn_comp_symm,TranscendentalNakamaye.Pending.SupNormalizedPotentialLimit.exists_integral_abs_supNormalize_le,TranscendentalNakamaye.Pending.SupNormalizedPotentialLimit.exists_isSemipositive_smul_sub,TranscendentalNakamaye.Pending.SupNormalizedPotentialLimit.exists_isChartLimit_supNormalize,TranscendentalNakamaye.Pending.SupNormalizedPotentialLimit.distributionalLeviCLM_regularDistribution_const,TranscendentalNakamaye.Pending.SupNormalizedPotentialLimit.distributionalLeviCLM_regularDistribution_sub_const,TranscendentalNakamaye.Pending.SupNormalizedPotentialLimit.LeviTestIdentity.supNormalize,TranscendentalNakamaye.Pending.SupNormalizedPotentialLimit.exists_kahlerCurrent_of_tendsto_leviTest}
\leanok
Let \(X\) be a compact connected complex manifold of dimension \(n=q+1\) with a K\"ahler form \(\omega\). For continuous \(v\) put \(\hat v=v-\sup_Xv\).
(1) For a closed real \((1,1)\)-form \(\Theta\), a chart \(e\) and a compact \(K\subset e(\text{chart})\), one constant bounds \(\int_K|\hat v\circ e^{-1}|\) for every continuous
\(\theta\)-psh \(v\) and every \(\theta\) with \(\Theta-\theta\ge0\). (2) Along any ultrafilter, the sup-normalized potentials of such a family have weak limits in every
chart. (3) Let currents \(P_i\) satisfy the Levi identity (Lemma~\ref{bou:positivity-levi-link}) with backgrounds \(\theta_i\le\Theta\) and continuous \(\theta_i\)-psh
potentials \(v_i\); if \(P_i\to R\) and \([X]\wedge\theta_i\to[X]\wedge\theta\) weakly along an ultrafilter and \(R-[X]\wedge\omega_\delta\) is positive for a K\"ahler form
\(\omega_\delta\), then \(\{\theta\}\) contains a K\"ahler current with representative \(\theta\) and lower bound \(\omega_\delta\).
\end{lemma}
\begin{proof}
\leanok
(1) is the uniform \(L^1\) bound of Lemma~\ref{bou:theta-psh-l1-bound} at level \(0\) (\(\hat v\le0\) and \(\hat v\) vanishes somewhere), the \(\theta\)-psh functions being
\(\Theta\)-psh. (2) follows from (1) and Lemma~\ref{bou:psh-weak-compactness}. (3) The Levi identity is unchanged by subtracting a constant, passes to the chart limits
of (2) (Lemma~\ref{bou:positivity-levi-link}), and Lemma~\ref{bou:positivity-levi-link} with \(\sigma=\omega_\delta\) gives the K\"ahler current.
\end{proof}

\begin{lemma}[The limit of the fibre integrals is a K\"ahler current]
\label{bou:fibre-potential-kahler-limit}
\uses{bou:product-stokes-levi,bou:sup-normalized-limit,bou:continuous-theta-psh,bou:fibre-pushforward,bou:diagonal-current-mass,bou:nef-mixed-mass}
\source{dp04-kahler-cone:2.12,bou02-volume:4.8}
\lean{TranscendentalNakamaye.Pending.FibrePotentialKahlerLimit.thetaPlurisubharmonic_fibreV,TranscendentalNakamaye.Pending.FibrePotentialKahlerLimit.smoothCurrent_add_apply,TranscendentalNakamaye.Pending.FibrePotentialKahlerLimit.baseSlots,TranscendentalNakamaye.Pending.FibrePotentialKahlerLimit.pertSlots,TranscendentalNakamaye.Pending.FibrePotentialKahlerLimit.topDensity_fibreForm_coeff,TranscendentalNakamaye.Pending.FibrePotentialKahlerLimit.topRatio_fibreForm,TranscendentalNakamaye.Pending.FibrePotentialKahlerLimit.fibreIntegral_add_smul_eq_sum,TranscendentalNakamaye.Pending.FibrePotentialKahlerLimit.continuous_fibreIntegral_add_smul,TranscendentalNakamaye.Pending.FibrePotentialKahlerLimit.tendsto_smoothCurrent_fibreConst,TranscendentalNakamaye.Pending.FibrePotentialKahlerLimit.exists_kahlerCurrent_of_tendsto_fibrePushforward,TranscendentalNakamaye.Pending.FibrePotentialKahlerLimit.exists_kahlerCurrent_of_tendsto_prodForm}
\leanok
Let \(X\) be a compact connected complex manifold of dimension \(n=q+1\) with a K\"ahler form \(\omega\), \(b\) a closed real \((1,1)\)-form, \(\varepsilon_i\in[0,1]\)
tending to \(0\) along an ultrafilter, \(b_i=b+\varepsilon_i\omega\), and \(g_i\) chart-smooth on \(X\times X\) with \(\alpha_i=p_1^*b_i+p_2^*b_i+dd^cg_i\ge0\). If
\([X\times X]\wedge\alpha_i^{\,n}\to\Theta\) weakly and \(\Theta-\delta[\Delta]\) is positive with \(\delta>0\), then the class of \(c\,b\), \(c=n\int_Xb^{n-1}\wedge\omega\),
contains a K\"ahler current with representative \(c\,b\) and lower bound \(\delta\omega\). (Also in the form with general backgrounds \(b_i\) whose multiples
\(c_ib_i\) are dominated by one closed form and converge on test forms.)
\end{lemma}
\begin{proof}
\leanok
By Lemma~\ref{bou:product-stokes-levi} the fibre push-forwards \(R_i=p_{1*}([X\times X]\wedge\alpha_i^{\,n}\wedge p_2^*\omega)\) satisfy the Levi identity with
\(\theta_i=c_ib_i\) and continuous fibre potentials, which are \(\theta_i\)-psh by Lemmas~\ref{bou:product-stokes-levi} and~\ref{bou:continuous-theta-psh}.
\(c_i=n\int_Xb_i^{n-1}\wedge\omega\) is a polynomial in \(\varepsilon_i\) (Lemma~\ref{bou:nef-mixed-mass}), so \(c_ib_i\to c\,b\) and the \(c_ib_i\) are dominated by one
multiple of \(\omega\). \(R_i\to R=p_{1*}(\Theta\wedge p_2^*\omega)\) and \(R-\delta\omega\) is positive (Lemma~\ref{bou:fibre-pushforward}); Lemma~\ref{bou:sup-normalized-limit}(3)
with \(\omega_\delta=\delta\omega\) concludes.
\end{proof}

\begin{lemma}[The nef core, given Yau's theorem as a hypothesis]
\label{bou:nef-core}
\uses{bou:mass-concentration,bou:concentrated-current-integration,bou:diagonal-lower-bound,bou:closed-set-mass-limit,bou:fibre-potential-kahler-limit,bou:nef-mixed-mass,bou:product-manifold}
\source{dp04-kahler-cone:2.12,bou02-volume:4.8}
\lean{TranscendentalNakamaye.Pending.NefCoreKahlerCurrent.pullback_add_smul,TranscendentalNakamaye.Pending.NefCoreKahlerCurrent.prodForm_add_smul,TranscendentalNakamaye.Pending.NefCoreKahlerCurrent.isKahler_add_smul_of_isSemipositive,TranscendentalNakamaye.Pending.NefCoreKahlerCurrent.topRatio_formPow_eq_slotForm',TranscendentalNakamaye.Pending.NefCoreKahlerCurrent.topRatio_formPow_pos,TranscendentalNakamaye.Pending.NefCoreKahlerCurrent.topRatio_formPow_eq_mul,TranscendentalNakamaye.Pending.NefCoreKahlerCurrent.ddcClosedForm_zero,TranscendentalNakamaye.Pending.NefCoreKahlerCurrent.slotIntegral_ite_add_ddc,TranscendentalNakamaye.Pending.NefCoreKahlerCurrent.YauHypothesis,TranscendentalNakamaye.Pending.NefCoreKahlerCurrent.volPoly,TranscendentalNakamaye.Pending.NefCoreKahlerCurrent.massPoly,TranscendentalNakamaye.Pending.NefCoreKahlerCurrent.prodForm_add_smul_form,TranscendentalNakamaye.Pending.NefCoreKahlerCurrent.continuous_volPoly,TranscendentalNakamaye.Pending.NefCoreKahlerCurrent.continuous_massPoly,TranscendentalNakamaye.Pending.NefCoreKahlerCurrent.volPoly_zero,TranscendentalNakamaye.Pending.NefCoreKahlerCurrent.kahlerVolume_real_univ_pos,TranscendentalNakamaye.Pending.NefCoreKahlerCurrent.integral_pos_of_forall_pos,TranscendentalNakamaye.Pending.NefCoreKahlerCurrent.kahlerVolume_prod_univ_ne_zero,TranscendentalNakamaye.Pending.NefCoreKahlerCurrent.volPoly_pos,TranscendentalNakamaye.Pending.NefCoreKahlerCurrent.exists_pos_le_volPoly,TranscendentalNakamaye.Pending.NefCoreKahlerCurrent.mul_vol_eq_volPoly,TranscendentalNakamaye.Pending.NefCoreKahlerCurrent.exists_massPoly_le,TranscendentalNakamaye.Pending.NefCoreKahlerCurrent.exists_kahlerCurrent_of_nef,TranscendentalNakamaye.Pending.NefCoreKahlerCurrent.isBigClass_of_nef}
\leanok
Let \((X,\omega)\) be a compact connected K\"ahler manifold of dimension \(n=q+1\) and \(b\) a smooth closed real \((1,1)\)-form which is nef with \(\int_Xb^n>0\).
Assume (\emph{hypothesis}, the form of Theorem~\ref{bou:yau} used here): for \(0<\varepsilon\le1\) and every K\"ahler form \(\omega'\) on \(X\times X\) there are a chart-smooth
\(g\) and \(c>0\) with \(\alpha=p_1^*(b+\varepsilon\omega)+p_2^*(b+\varepsilon\omega)+dd^cg\) K\"ahler and \(\alpha^{2n}=c\,\omega'^{2n}\). Then \(b\) is the representative of a
K\"ahler current; in particular \(\{b\}\) is big.
\end{lemma}
\begin{proof}
\leanok
On \(X\times X\) with \(\widetilde\omega=p_1^*\omega+p_2^*\omega\): concentration of mass on the diagonal (Lemma~\ref{bou:mass-concentration}) gives \(f_\varepsilon\), tubes
\(V_\varepsilon\) and \(\omega_\varepsilon=\widetilde\omega+dd^cf_\varepsilon\ge\widetilde\omega/2\) with \(\int_{V_\varepsilon}\omega_\varepsilon^n\wedge\widetilde\omega^n\ge\delta\). The hypothesis at
\(\varepsilon_k=1/(k+1)\), \(\omega'=\omega_{\varepsilon_k}\) gives K\"ahler \(\alpha_k\) with \(\alpha_k^{2n}=c_k\omega_{\varepsilon_k}^{2n}\); integrating, \(c_k\int\widetilde\omega^{2n}\) is the
polynomial \(P(\varepsilon_k)=\int(p_1^*b_{\varepsilon_k}+p_2^*b_{\varepsilon_k})^{2n}\), bounded below on \([0,1]\) (\(P(0)=\binom{2n}n(\int b^n)^2>0\), Lemma~\ref{bou:nef-mixed-mass}).
The integration step (Lemma~\ref{bou:concentrated-current-integration}) gives \(\int_{V_{\varepsilon_k}}\alpha_k^n\wedge\widetilde\omega^n\ge\delta'>0\); the currents
\([X\times X]\wedge\alpha_k^n\) have bounded masses, a weak limit \(\Theta\) along an ultrafilter satisfies \(\Theta\ge\lambda[\Delta]\) with \(\lambda>0\)
(Lemmas~\ref{bou:closed-set-mass-limit} and~\ref{bou:diagonal-lower-bound}), and Lemma~\ref{bou:fibre-potential-kahler-limit} gives a K\"ahler current in
\(\{c\,b\}\), \(c=n\int b^{n-1}\wedge\omega>0\) (Lemma~\ref{bou:nef-mixed-mass}); divide by \(c\).
\end{proof}

\begin{lemma}[Mixed integrals of closed \((1,1)\)-forms: polynomiality, nef positivity, products]
\label{bou:nef-mixed-mass}
\uses{bou:product-mixed-measure,posclass:nef,gnpp:kahler-domination}
\source{dp04-kahler-cone:2.12,bou02-volume:4.8}
\lean{TranscendentalNakamaye.Pending.NefMixedMass.topDensity_slotForm_add_smul,TranscendentalNakamaye.Pending.NefMixedMass.topRatio_slotForm_add_smul,TranscendentalNakamaye.Pending.NefMixedMass.topRatio_slotForm_nonneg,TranscendentalNakamaye.Pending.NefMixedMass.exists_isSemipositive_of_nef,TranscendentalNakamaye.Pending.NefMixedMass.nefWithRespectTo_of_isSemipositive,TranscendentalNakamaye.Pending.NefMixedMass.slotIntegral,TranscendentalNakamaye.Pending.NefMixedMass.integrable_topRatio,TranscendentalNakamaye.Pending.NefMixedMass.slotIntegral_add_smul_eq_sum,TranscendentalNakamaye.Pending.NefMixedMass.continuous_slotIntegral_add_smul,TranscendentalNakamaye.Pending.NefMixedMass.slotIntegral_nonneg,TranscendentalNakamaye.Pending.NefMixedMass.slotIntegral_add_ddc,TranscendentalNakamaye.Pending.NefMixedMass.slotIntegral_nonneg_of_nef,TranscendentalNakamaye.Pending.NefMixedMass.topDensity_slotForm_snoc_sub_smul,TranscendentalNakamaye.Pending.NefMixedMass.slotIntegral_snoc_sub_smul,TranscendentalNakamaye.Pending.NefMixedMass.topRatio_formPow_eq_slotForm,TranscendentalNakamaye.Pending.NefMixedMass.integral_formPow_eq_slotIntegral,TranscendentalNakamaye.Pending.NefMixedMass.fibreIntegral_eq_slotIntegral,TranscendentalNakamaye.Pending.NefMixedMass.fibreIntegral_pos_of_nef,TranscendentalNakamaye.Pending.NefMixedMass.topRatio_formPow_prodForm,TranscendentalNakamaye.Pending.NefMixedMass.integral_formPow_prodForm}
\leanok
Let \((M,\omega)\) be a compact K\"ahler manifold of dimension \(N\) and \(\theta_1,\dots,\theta_N\) smooth closed real \((1,1)\)-forms. Then
\(t\mapsto\int_M(\theta_1+t\eta_1)\wedge\dots\wedge(\theta_N+t\eta_N)\) is a polynomial in \(t\); \(\int_M\theta_1\wedge\dots\wedge\theta_N\ge0\) if all \(\theta_i\) are semipositive,
and also if all are nef; it is unchanged by adding \(dd^c\) of smooth functions. For a closed real \((1,1)\)-form \(a\) on \(X\) of dimension \(n\),
\(\int_{X\times X}(p_1^*a+p_2^*a)^{2n}=\binom{2n}n\bigl(\int_Xa^n\bigr)^2\); and for nef \(b\) with \(\int_Xb^n>0\), \(\int_Xb^{n-1}\wedge\omega>0\).
\end{lemma}
\begin{proof}
\leanok
Polynomiality is multilinearity of the mixed density. For nef forms, \(\theta_i+t\omega+dd^ch_i\ge0\) for every \(t>0\); class invariance and nonnegativity for
semipositive slots give \(\int\prod(\theta_i+t\omega)\ge0\), and \(t\to0\). The product formula is Lemma~\ref{bou:product-mixed-measure}(2) for smooth representatives.
Finally, with \(C\omega-b\ge0\), \(C\int b^{n-1}\wedge\omega-\int b^n=\int b^{n-1}\wedge(C\omega-b)\ge0\) (nef slots), so \(\int b^{n-1}\wedge\omega\ge C^{-1}\int b^n>0\).
\end{proof}

\begin{lemma}[Mixed measures pulled back along \(\mu\times\mathrm{id}\) do not lose mass]
\label{bou:product-map-mixed-measure}
\uses{ana:blowups,bou:product-map-isomorphism-over,pc:domination-modification-mass,gnpp:noncharging,bou:product-manifold}
\source{bou02-volume:4.10}
\lean{TranscendentalNakamaye.Pending.ProductMapMixedMeasure.thetaPlurisubharmonic_zero_of_isSemipositive,TranscendentalNakamaye.Pending.ProductMapMixedMeasure.globalMixedMeasure_prod_univ_eq_zero,TranscendentalNakamaye.Pending.ProductMapMixedMeasure.le_globalMixedMeasure_preimage_prodMap_of_isBlowupAlong,TranscendentalNakamaye.Pending.ProductMapMixedMeasure.le_globalMixedMeasure_preimage_prodMap_of_isBlowupSequence}
\leanok
Let \(W\) be a complex manifold. (1) If \(C\) is an analytic subset with empty interior of a compact complex manifold \(Y\), then \(C\times W\) is charged by no global
mixed non-pluripolar measure on \(Y\times W\). (2) Let \(\mu\colon Y\to X\) be a finite sequence of blowups with smooth centres of compact complex manifolds, \(\nu=\mu\times
\mathrm{id}\colon Y\times W\to X\times W\), and \(\theta_1,\dots,\theta_N\) semipositive closed real \((1,1)\)-forms on \(X\times W\), \(N=\dim(X\times W)\). For the mixed
measures with the potential \(0\), \(\mu_\theta(S)\le\mu_{\nu^*\theta}(\nu^{-1}S)\) for every set \(S\subset X\times W\).
\end{lemma}
\begin{proof}
\leanok
(1) \(C\times W\) is analytic with empty interior, hence locally pluripolar, and such sets are not charged (Lemma~\ref{gnpp:noncharging}). (2) For one blowup along \(C\),
\(\nu\) is an isomorphism over \((Y\setminus C)\times W\) (Lemma~\ref{bou:product-map-isomorphism-over}), where the measures correspond
(Lemma~\ref{pc:domination-modification-mass}), and \(C\times W\) carries no mass by (1); a blowup sequence follows by induction.
\end{proof}

\begin{lemma}[Stokes on \(V\times X\): the Levi identity of the fibre integral]
\label{bou:two-factor-stokes-levi}
\uses{bou:product-stokes-levi,bou:fibre-pushforward,bou:product-fibre-integration,bou:positivity-levi-link,bou:product-manifold}
\source{dp04-kahler-cone:2.12,bou02-volume:4.8}
\lean{TranscendentalNakamaye.Pending.TwoFactorStokesLevi.liftTest,TranscendentalNakamaye.Pending.TwoFactorStokesLevi.chartSmooth_liftTest,TranscendentalNakamaye.Pending.TwoFactorStokesLevi.fibreFam,TranscendentalNakamaye.Pending.TwoFactorStokesLevi.fibreLeviTest,TranscendentalNakamaye.Pending.TwoFactorStokesLevi.fibreFam_coeff,TranscendentalNakamaye.Pending.TwoFactorStokesLevi.fibreLeviTest_coeff,TranscendentalNakamaye.Pending.TwoFactorStokesLevi.topDensity_chain_fibreLeviTest,TranscendentalNakamaye.Pending.TwoFactorStokesLevi.topDensity_slot_fibre,TranscendentalNakamaye.Pending.TwoFactorStokesLevi.smoothCurrent_sub_eq_sum_slot,TranscendentalNakamaye.Pending.TwoFactorStokesLevi.smoothDDC_liftTest,TranscendentalNakamaye.Pending.TwoFactorStokesLevi.vertSlots,TranscendentalNakamaye.Pending.TwoFactorStokesLevi.integrationCurrent_slot_ibp_eq_chain,TranscendentalNakamaye.Pending.TwoFactorStokesLevi.fibreV,TranscendentalNakamaye.Pending.TwoFactorStokesLevi.continuous_fibreV,TranscendentalNakamaye.Pending.TwoFactorStokesLevi.locallyIntegrableOn_fibreV,TranscendentalNakamaye.Pending.TwoFactorStokesLevi.leviTestIdentity_fibrePushforward,TranscendentalNakamaye.Pending.TwoFactorStokesLevi.isLocallyThetaPshDistribution_fibreV,TranscendentalNakamaye.Pending.TwoFactorStokesLevi.thetaPlurisubharmonic_fibreV}
\leanok
Let \(V\) and \(X\) be compact complex manifolds of dimension \(n\ge1\), \(\omega_V,\omega_X\) K\"ahler forms, \(a\) and \(c\) closed real \((1,1)\)-forms on \(V\) and \(X\), \(g\) a
smooth function on \(V\times X\), \(\alpha=p_1^*a+p_2^*c+dd^cg\) and \(R=p_{1*}([V\times X]\wedge\alpha^n\wedge p_2^*\omega_X)\). There is a continuous function \(v\) on \(V\)
(the fibre potential) with \(R=m\,a+dd^cv\) on the Levi test forms of \(V\), \(m=n\int_Xc^{n-1}\wedge\omega_X\); and if \(\alpha\ge0\), then \(v\) is \((m\,a)\)-psh.
\end{lemma}
\begin{proof}
\leanok
This is Lemma~\ref{bou:product-stokes-levi} (the case \(V=X\), \(a=c\)) with two different factors: the vertical density factorization, the fibre push-forward of
product forms and the integration by parts in the slots are stated for two factors; positivity of \(\alpha\) makes \(R\) positive, and the Levi identity then gives the
psh property of \(v\) (Lemma~\ref{bou:positivity-levi-link}).
\end{proof}

\begin{lemma}[The limit on \(V\times X\) with a semipositive lower bound]
\label{bou:graph-fibre-potential-limit}
\uses{bou:two-factor-stokes-levi,bou:fibre-potential-kahler-limit,bou:sup-normalized-limit,bou:positivity-levi-link,bou:graph-fibre-integration,bou:graph-support}
\source{dp04-kahler-cone:2.12,bou02-volume:4.8,bou02-volume:4.10}
\lean{TranscendentalNakamaye.Pending.GraphFibrePotentialLimit.exists_isThetaPshPotential_of_tendsto_leviTest,TranscendentalNakamaye.Pending.GraphFibrePotentialLimit.fibreConst,TranscendentalNakamaye.Pending.GraphFibrePotentialLimit.tendsto_smoothCurrent_smul_add_smul,TranscendentalNakamaye.Pending.GraphFibrePotentialLimit.smul_wedge_pullback_eq,TranscendentalNakamaye.Pending.GraphFibrePotentialLimit.exists_isThetaPshPotential_of_tendsto_graph}
\leanok
Let \(V\) be a connected compact complex manifold of dimension \(n\ge1\) and \(X\) a compact complex manifold of dimension \(n\), with K\"ahler forms
\(\omega_V,\omega_X\), \(b\) a closed real \((1,1)\)-form on \(V\) and \(\mu\colon V\to X\) holomorphic. Along an ultrafilter let \(\varepsilon_i\in[0,1]\) tend to \(0\) and
let \(g_i\) be smooth functions on \(V\times X\) with \(\alpha_i=p_1^*(b+\varepsilon_i\omega_V)+p_2^*\omega_X+dd^cg_i\ge0\), such that \([V\times X]\wedge\alpha_i^n\) converges
weakly to \(\Theta\) and \(\Theta-\lambda[\Gamma_\mu]\) is positive, \([\Gamma_\mu]\) the integration current of the graph. Then \(m\,b-\lambda\,\mu^*\omega_X\) has a psh potential
on \(V\), \(m=n\int_X\omega_X^n\). More generally, for Levi-identity data with continuous potentials and backgrounds converging to \(\theta'\), a limit \(R\) with
\(R-[V]\wedge\sigma\) positive gives a \((\theta'-\sigma)\)-psh potential, for every closed real \((1,1)\)-form \(\sigma\).
\end{lemma}
\begin{proof}
\leanok
The limit argument of Lemma~\ref{bou:fibre-potential-kahler-limit} uses the K\"ahler property of the lower bound only to package the result as a K\"ahler current; it
gives a psh potential for \(\theta'-\sigma\) for any closed \(\sigma\) with \(R-[V]\wedge\sigma\ge0\). In the graph case, \(R=p_{1*}(\Theta\wedge p_2^*\omega_X)\), the Levi
identity is Lemma~\ref{bou:two-factor-stokes-levi} with \(m=n\int_X\omega_X^n\), and \(p_{1*}([\Gamma_\mu]\wedge p_2^*\omega_X)=[V]\wedge\mu^*\omega_X\)
(Lemma~\ref{bou:graph-fibre-integration}), so \(\sigma=\lambda\mu^*\omega_X\).
\end{proof}

\begin{lemma}[Tools for the uniform nef core on \(Y\times X\)]
\label{bou:uniform-nef-core-support}
\uses{bou:nef-mixed-mass,bou:mixed-mass-class-invariance,bou:product-mixed-measure,bou:graph-support,bou:diagonal-lower-bound,bou:product-map-mixed-measure,bou:mass-concentration,mmass:slot-background-monotone,posclass:nef}
\source{bou02-volume:4.8,bou02-volume:4.10,dp04-kahler-cone:2.12}
\lean{TranscendentalNakamaye.Pending.UniformNefCoreSupport.withDensity_topRatio_formPow,TranscendentalNakamaye.Pending.UniformNefCoreSupport.slotIntegral_eq_of_background,TranscendentalNakamaye.Pending.UniformNefCoreSupport.slotIntegral_comp_perm,TranscendentalNakamaye.Pending.UniformNefCoreSupport.exists_perm_prefix,TranscendentalNakamaye.Pending.UniformNefCoreSupport.slotIntegral_piecewise_const,TranscendentalNakamaye.Pending.UniformNefCoreSupport.slotIntegral_add_le_two_pow,TranscendentalNakamaye.Pending.UniformNefCoreSupport.slotIntegral_const_smul,TranscendentalNakamaye.Pending.UniformNefCoreSupport.slotIntegral_form,TranscendentalNakamaye.Pending.UniformNefCoreSupport.topRatio_formPow_prodForm₂,TranscendentalNakamaye.Pending.UniformNefCoreSupport.slotIntegral_prodForm,TranscendentalNakamaye.Pending.UniformNefCoreSupport.mass_graphCurrent_eq,TranscendentalNakamaye.Pending.UniformNefCoreSupport.exists_isPositive_sub_smul_graphCurrent_of_le,TranscendentalNakamaye.Pending.UniformNefCoreSupport.pullback_add,TranscendentalNakamaye.Pending.UniformNefCoreSupport.exists_isKahler_add_smul_of_nef,TranscendentalNakamaye.Pending.UniformNefCoreSupport.contMDiff_prodMap_id,TranscendentalNakamaye.Pending.UniformNefCoreSupport.pullback_prodForm_geometricallyEq,TranscendentalNakamaye.Pending.UniformNefCoreSupport.prodForm_sub_coeff,TranscendentalNakamaye.Pending.UniformNefCoreSupport.isSemipositive_half_add_ddc_comp,TranscendentalNakamaye.Pending.UniformNefCoreSupport.isSemipositive_concentrationSlot,TranscendentalNakamaye.Pending.UniformNefCoreSupport.le_globalMixedMeasure_concentrationSlot_comp}
\leanok
Let \(M\) be a compact complex manifold of dimension \(N\) with K\"ahler forms \(\omega,\omega'\), and \(V,X\) compact complex manifolds of dimension \(n\).
\begin{enumerate}
\item \(\int_M\theta_1\wedge\dots\wedge\theta_N\) does not depend on the K\"ahler form used to integrate and is symmetric in the slots; if
      \(\int_Ma^i\wedge c^{N-i}\le C\) for all \(i\le N\), then \(\int_M(c+a)^N\le2^NC\); and \(\int_{V\times X}(p_1^*a+p_2^*c)^{2n}=\binom{2n}n\int_Va^n\int_Xc^n\).
\item For \(\mu\colon V\to X\) holomorphic, the graph current has mass \(\int_V(\omega_V+\mu^*\omega_X)^n\); if \(V\) is connected and \(\Theta\) is a closed positive current of
      bidimension \((n,n)\) on \(V\times X\) whose trace measure gives the graph mass \(\ge\delta\), there is \(c\ge\max(0,\delta/\operatorname{mass}[\Gamma_\mu])\) with
      \(\Theta-c[\Gamma_\mu]\) positive.
\item If \(b\) is nef with respect to \(\omega'\) and \(\omega_V-\rho\omega'\ge0\) with \(\rho>0\), then for every \(\varepsilon>0\) some \(b+\varepsilon\omega_V+dd^ch\) is K\"ahler.
\item Let \(\mu\colon V\to X\) be a finite sequence of blowups, \(\nu=\mu\times\mathrm{id}\colon V\times X\to X\times X\), \(\omega_V\ge\mu^*\omega_X\) a K\"ahler form on \(V\),
      \(W=p_1^*\omega_X+p_2^*\omega_X\), \(W_V=p_1^*\omega_V+p_2^*\omega_X\), and \(f\) smooth on \(X\times X\) with \(\tfrac12W+dd^cf\ge0\). Then \(\tfrac12W_V+dd^c(f\circ\nu)\ge0\),
      and for every \(p\) and every set \(S\subset X\times X\) the mixed measure of \(S\) with \(p\) slots \(W+dd^cf\) and the others \(W\) is at most the mixed measure of
      \(\nu^{-1}S\) with \(p\) slots \(W_V+dd^c(f\circ\nu)\) and the others \(W_V\) (potentials \(0\)).
\end{enumerate}
\end{lemma}
\begin{proof}
\leanok
(1) Class invariance of slot integrals (Lemma~\ref{bou:mixed-mass-class-invariance}) and the product formula (Lemma~\ref{bou:product-mixed-measure}); the binomial bound
expands \((c+a)^N\) and bounds each of the \(2^N\) terms, after a permutation of the slots, by a prefix integral. (2) The graph is a copy of \(V\) embedded by
\((\mathrm{id},\mu)\); the support theorem on a graph (Lemma~\ref{bou:graph-support}) and the mass comparison of Corollary~\ref{bou:diagonal-lower-bound}.
(3) Nefness gives \(b+\tfrac{\varepsilon\rho}2\omega'+dd^ch\ge0\), and \(\varepsilon\omega_V\ge\varepsilon\rho\,\omega'\). (4) \(\nu^*(W+dd^cf)\ge\tfrac12\nu^*W\) and
\(W_V\ge\nu^*W\); enlarge the slots (Lemma~\ref{mmass:slot-background-monotone}) and use Lemma~\ref{bou:product-map-mixed-measure}.
\end{proof}

\begin{corollary}[Yau's theorem on \(V\times X\)]
\label{bou:yau-product}
\uses{bou:yau-consequences,bou:yau-volume-ratio,bou:product-manifold}
\source{bou02-volume:3.4,dp04-kahler-cone:2.12}
\lean{TranscendentalNakamaye.Pending.UniformNefCoreSupport.exists_yau_prodForm}
\leanok
Let \(V\) and \(X\) be connected compact complex manifolds of dimension \(n\ge1\), \(\omega_X\) a K\"ahler form on \(X\), and \(a\) a closed real \((1,1)\)-form on \(V\) with
\(a+dd^ch\) K\"ahler for some smooth \(h\). For every K\"ahler form \(\omega'\) on \(V\times X\) there are a smooth \(g\) and \(c>0\) with \(p_1^*a+p_2^*\omega_X+dd^cg\) K\"ahler and
\((p_1^*a+p_2^*\omega_X+dd^cg)^{2n}=c\,\omega'^{2n}\).
\end{corollary}
\begin{proof}
\leanok
Corollary~\ref{bou:yau-consequences}(1) for the K\"ahler form \(p_1^*(a+dd^ch)+p_2^*\omega_X\) on the connected manifold \(V\times X\), and \(g=h\circ p_1+\varphi\).
\end{proof}
This is deduced from Theorem~\ref{bou:yau}.

\begin{lemma}[Constants of the uniform nef core]
\label{bou:uniform-nef-core-constant}
\uses{bou:nef-mixed-mass,bou:graph-fibre-potential-limit,fpgeom:complex-analytic-subvariety}
\lean{TranscendentalNakamaye.Pending.UniformNefCore.exists_small_pos,TranscendentalNakamaye.Pending.UniformNefCore.prodForm_add_prodForm,TranscendentalNakamaye.Pending.UniformNefCore.coreConst,TranscendentalNakamaye.Pending.UniformNefCore.fibreConst_pos,TranscendentalNakamaye.Pending.UniformNefCore.coreConst_pos,TranscendentalNakamaye.Pending.UniformNefCore.connectedSpace_of_isComplexAnalyticallyIrreducible_univ}
\leanok
Let \((X,\omega)\) be a nonempty compact K\"ahler manifold of dimension \(n\ge1\). For \(\delta,\eta>0\) and \(C\ge0\) the explicit constant
\[
 c(\delta,\eta,C)=\frac{P_0\,\delta^2}{V_0\cdot4\cdot2^n(2n)!\,M_0\,(2^nC+1)\,m},\qquad B=\tbinom{2n}n,\ v=\int_X\omega^n,\ m=n\,v,
\]
\(P_0=B\tfrac\eta2v\), \(V_0=B(C+1)v\), \(M_0=B(2^nC+2)2^nv\), is positive. A complex manifold whose underlying set is an irreducible analytic subset of itself is connected.
\end{lemma}
\begin{proof}
\leanok
Every factor is positive (\(v>0\), Lemma~\ref{bou:nef-mixed-mass}). A disconnected manifold is the union of two nonempty open and closed sets, each an analytic subset.
\end{proof}

\begin{lemma}[Boucksom's uniform nef core]
\label{bou:uniform-nef-core}
\uses{bou:resolution-assembly,bou:uniform-nef-core-constant,bou:uniform-nef-core-support,bou:yau-product,bou:mass-concentration,bou:concentrated-current-explicit,bou:graph-support,bou:graph-fibre-integration,bou:graph-fibre-potential-limit,bou:product-map-mixed-measure,bou:nef-mixed-mass,ana:blowups}
\source{bou02-volume:4.10,bou02-volume:4.8,dp04-kahler-cone:2.12,dp04-kahler-cone:2.6}
\lean{TranscendentalNakamaye.Pending.UniformNefCore.exists_thetaPsh_of_blowup,TranscendentalNakamaye.Pending.UniformNefCore.uniformNefCore,TranscendentalNakamaye.Pending.UniformNefCore.uniformNefCore_of_pos}
\leanok
Let \((X,\omega)\) be a compact connected K\"ahler manifold of dimension \(n\ge1\). Then statement (N) of Lemma~\ref{bou:resolution-assembly} holds, with the constant
\(c=c(\delta,\eta,C)\) of Lemma~\ref{bou:uniform-nef-core-constant}, \(\delta\) the concentration constant of the diagonal of \(X\times X\).
\end{lemma}
\begin{proof}
\leanok
Fix the concentration data \((f_\varepsilon,V_\varepsilon,\delta)\) of Lemma~\ref{bou:mass-concentration} for the diagonal \(\Delta\) of \(X\times X\) with
\(W=p_1^*\omega+p_2^*\omega\). Let \(\mu\colon Y\to X\) be a blowup sequence (so \(Y\) is connected), \(\omega'\) a K\"ahler form on \(Y\) and \(b\) nef with \(\int_Yb^n\ge\eta\) and
\(\int_Yb^i\wedge(\mu^*\omega)^{n-i}\le C\). Work on \(Y\times X\) with \(\nu=\mu\times\mathrm{id}\), so that \(\nu^{-1}(\Delta)=\Gamma_\mu\), the graph of \(\mu\), and with
\(\omega_Y=\mu^*\omega+\rho\,\omega'\), \(W_Y=p_1^*\omega_Y+p_2^*\omega\), \(\rho>0\) small depending on \(Y\). The concentration pulls back
(Lemma~\ref{bou:uniform-nef-core-support}(4)). For \(\varepsilon_k\to0\), \(b+\varepsilon_k\omega_Y\) has a K\"ahler representative
(Lemma~\ref{bou:uniform-nef-core-support}(3)), and Yau's theorem on \(Y\times X\) (Corollary~\ref{bou:yau-product}) with the pulled-back targets, followed by the
concentration lemma with explicit constant (Lemma~\ref{bou:concentrated-current-explicit}), gives closed positive currents in
\((p_1^*(b+\varepsilon_k\omega_Y)+p_2^*\omega)^n\) whose trace measures give \(\Gamma_\mu\) mass bounded below in terms of \(\delta\), \(\eta\), \(C\) and \(\int_X\omega^n\) only
(Lemma~\ref{bou:uniform-nef-core-support}(1): product formula, background independence, binomial bound). A weak limit \(\Theta\) satisfies \(\Theta\ge\lambda[\Gamma_\mu]\) with
\(\lambda\) bounded below likewise (Lemma~\ref{bou:uniform-nef-core-support}(2)), and Lemma~\ref{bou:graph-fibre-potential-limit} gives a psh potential for
\(m\,b-\lambda\,\mu^*\omega\), \(m=n\int_X\omega^n\). Dividing by \(m\) and shrinking \(\lambda/m\) to \(c\), \(b-c\,\mu^*\omega\) has a psh potential.
\end{proof}
This uses Yau's theorem (Theorem~\ref{bou:yau}), through Corollary~\ref{bou:yau-product}.

\begin{lemma}[The canonical potential of a non-big class has zero mass]
\label{ana:nonbig-canonical-mass}
\uses{ana:nonkahler-locus,fpenv:envelopes,trace:diagonal-mass,pc:bigness-from-mass}
\source{bou02-volume:1.3,bou02-volume:4.8}
\lean{TranscendentalNakamaye.Pending.BoucksomNonBigSupremumMass.diagonalPotentialMass_thetaPshSupremum_eq_zero_of_not_isBigClass,TranscendentalNakamaye.Pending.NonBigSupremumMass.diagonalPotentialMass_thetaPshSupremum_eq_zero_of_not_isBigClass,TranscendentalNakamaye.Pending.BoucksomMassBigness.diagonalPotentialMass_eq_zero_of_not_isBigClass}
\leanok
Let \((X,\omega)\) be a connected compact K\"ahler manifold of dimension \(n>0\) and \(\theta\) a smooth
closed real \((1,1)\)-form such that \(V_\theta=\sup\{u\in\operatorname{PSH}(X,\theta):u\le0\}\)
is theta-PSH. If \(\{\theta\}\) is not big, then
\(\int_X\langle\theta^n_{V_\theta}\rangle=0\).
\end{lemma}
\begin{proof}
\leanok
If \(\int_X\langle\theta^n_{V_\theta}\rangle>0\), Proposition~\ref{pc:bigness-from-mass} with \(Z=X\) and \(u=V_\theta\) makes \(\{\theta\}\) big.
\end{proof}
This is the contrapositive of Boucksom's criterion (Theorem~\ref{ana:ambient-volume-uniform-big}).

\begin{lemma}[Positive limiting volume produces a K\"ahler current]
\label{ana:ambient-volume-vanishing}
\uses{ana:restricted-volume,ana:nonkahler-locus,gnpp:global-measure}
\source{bou02-volume:4.8}
\lean{TranscendentalNakamaye.Pending.FullSpaceAmbientVanishing.tendsto_perturbedDiagonalPotentialEnvelope_nhds_zero}
\leanok
Let \((X,\omega)\) be compact K\"ahler and let \(\alpha\) be pseudoeffective and not big.
Then, along every sequence \(\epsilon_k\downarrow0\), the supremum of the ambient diagonal
non-pluripolar masses \(\operatorname{NP}(\theta+\epsilon_k\omega+dd^cu)(X)\) over the
\(\epsilon_k\)-admissible potentials \(u\) tends to \(0\).
\end{lemma}
\begin{proof}
\uses{ana:ambient-volume-uniform-bound,ana:ambient-volume-uniform-big,ana:nonkahler-locus}
\leanok
In dimension \(0\) every class is big, contrary to the hypothesis on \(\alpha\); so
\(n\ge1\). Write
\[
 E(\epsilon)=\sup_{T\in\mathcal A_\epsilon}\operatorname{NP}(T+\epsilon\omega)(X)
\]
for the quantity in question, as in Lemma~\ref{ana:ambient-volume-uniform-bound}.

Suppose \(E(\epsilon_k)\not\to0\). Since the \(E(\epsilon_k)\) are nonnegative, this
means that some neighbourhood \([0,c)\) of \(0\), with \(c>0\), is left by infinitely
many terms: there are arbitrarily large \(k\) with \(c\le E(\epsilon_k)\). Choosing
\(\eta\) with \(0<\eta<c\) gives \(\eta<E(\epsilon_k)\) for arbitrarily large \(k\),
that is, a strictly positive cofinal cluster of the ambient volumes along
\((\epsilon_k)\).

Lemma~\ref{ana:ambient-volume-uniform-bound} turns that cluster into the bound
\(\eta<E(\epsilon)\) at \emph{every} error \(\epsilon>0\), and
Theorem~\ref{ana:ambient-volume-uniform-big}, applied to the pseudoeffective class
\(\alpha\), then makes \(\alpha\) big. This contradicts the hypothesis, so
\(E(\epsilon_k)\to0\).
\end{proof}

\begin{lemma}[The volume of a pseudoeffective non-big class vanishes]
\label{ana:rv-nonbig-vanishing}
\uses{ana:restricted-volume,ana:nonkahler-locus,ana:fullspace-ambient-comparison}
\source{ct22-restricted-volumes:volr}
\lean{TranscendentalNakamaye.Pending.LocusProperties.numericalRestrictedVolume_eq_zero_of_not_isBigClass}
\leanok
Let \(X\) be a connected compact K\"ahler manifold of positive dimension and let
\(\alpha\) be pseudoeffective and not big. Then
\(\operatorname{RV}_{X\mid X}(\alpha)=0\).
\end{lemma}
\begin{proof}
\uses{ana:fullspace-ambient-comparison,ana:ambient-volume-vanishing,ana:restricted-volume-monotone}
\leanok
If \(E_{\mathrm{nn}}(\alpha)=X\) the restricted volume vanishes by definition. Otherwise it
is the infimum over \(\epsilon>0\) of the admissible mass envelopes
(Lemma~\ref{ana:restricted-volume-monotone}), and that infimum is already attained along
the cofinal sequence \(\epsilon_k=1/(k+1)\). Each envelope at \(\epsilon_k\) is bounded, by
Lemma~\ref{ana:fullspace-ambient-comparison}, by the supremum of the ambient
non-pluripolar masses at \(\epsilon_k\), and those tend to \(0\) by
Lemma~\ref{ana:ambient-volume-vanishing}; a sequence dominated by a null sequence has zero
infimum.
\end{proof}

\begin{proposition}[Loci, positivity, and modifications]
\label{ana:locus-properties}
\uses{ana:restricted-volume,ana:nonkahler-locus,ana:ideal-sheaves,ana:blowups,ana:locus-properties-elementary,ana:rv-nonbig-vanishing}
\source{ct22-restricted-volumes:sect2,ct22-restricted-volumes:finite,ct22-restricted-volumes:volr,ct22-restricted-volumes:nonnef,ct22-restricted-volumes:pull,ct22-restricted-volumes:invar}
\source{transcendental-nakamaye-copy4-20260911:sec:nK,transcendental-nakamaye-copy4-20260911:it:rv-positive,transcendental-nakamaye-copy4-20260911:it:rv-birational,transcendental-nakamaye-copy4-20260911:thm:main}
\lean{TranscendentalNakamaye.Pending.LocusProperties.numericalRestrictedVolume_ne_top,TranscendentalNakamaye.Pending.LocusProperties.numericalRestrictedVolume_eq_zero_iff_not_isBigClass,TranscendentalNakamaye.Pending.LocusProperties.isComplexAnalyticSubset_nonKahlerLocus,TranscendentalNakamaye.Pending.LocusProperties.exists_kahlerCurrent_poleSet_eq_nonKahlerLocus,TranscendentalNakamaye.Pending.LocusProperties.numericalRestrictedVolume_pos_of_not_subset_nonKahlerLocus,TranscendentalNakamaye.Pending.LocusProperties.nonNefLocus_pullback,TranscendentalNakamaye.Pending.LocusProperties.exists_strictTransform_numericalRestrictedVolume_eq}
\leanok
Let \(X\) be a connected compact K\"ahler manifold of dimension \(n\), let
\(\omega\) be a K\"ahler form, and let \(\alpha\) be a real \((1,1)\)-class.
If \(\alpha\) is pseudoeffective, its numerical restricted volumes along
positive-dimensional subvarieties are finite and nonnegative, and, when \(n>0\),
\(\operatorname{RV}_{X\mid X}(\alpha)\) (the volume of \(\alpha\)) is zero exactly
when \(\alpha\) is not big. If \(\alpha\) is big and \(n>0\),
\(E_{\mathrm{nK}}(\alpha)\) is an analytic subset, and there is a
K\"ahler current in \(\alpha\) with analytic singularities whose polar set is
exactly \(E_{\mathrm{nK}}(\alpha)\). If \(\alpha\) is
pseudoeffective and \(V\) is a positive-dimensional closed irreducible analytic
subset not contained in \(E_{\mathrm{nK}}(\alpha)\), then
\(\operatorname{RV}_{X\mid V}(\alpha)>0\).

For every modification \(\mu:\widetilde X\to X\), the manifold
\(\widetilde X\) is compact K\"ahler and \(\mu^*\alpha\) is big whenever
\(\alpha\) is big. If \(\alpha\) is big, then
\[
 E_{\mathrm{nn}}(\mu^*\alpha)=\mu^{-1}(E_{\mathrm{nn}}(\alpha)),
 \qquad
 E_{\mathrm{nK}}(\mu^*\alpha)=\mu^{-1}(E_{\mathrm{nK}}(\alpha))
 \cup\operatorname{Exc}(\mu),
\]
and if moreover \(V\) is a positive-dimensional closed irreducible analytic
subset with \(V\not\subset\mu(\operatorname{Exc}(\mu))\), its strict transform
\(\widetilde V\) is a closed irreducible analytic subset of dimension
\(\dim V\) with
\(\operatorname{RV}_{\widetilde X\mid\widetilde V}(\mu^*\alpha)
=\operatorname{RV}_{X\mid V}(\alpha)\).
\end{proposition}
\begin{proof}
\uses{ana:restricted-volume-monotone,ana:pure-dimension,ana:rv-mass-bound,ana:nonkahler-locus-analytic,ana:nonkahler-locus-analytic-bounded,ana:rv-kahler-independence,ana:strict-transform-composition,ana:blowup-surjective,ana:rv-birational-blowup,ana:modification-big-nonkahler}
\leanok
Finiteness: if \(V\subset E_{\mathrm{nn}}(\alpha)\) the restricted volume is \(0\);
otherwise, by Lemma~\ref{ana:restricted-volume-monotone}, it is the infimum over
\(\epsilon>0\) of the admissible mass envelopes, hence at most the envelope at
\(\epsilon=1\), which is finite by Lemma~\ref{ana:rv-mass-bound}. For \(V=X\) the restricted volume is the volume
of \(\alpha\), which vanishes exactly when \(\alpha\) is not big. For a big class, Theorem~\ref{ana:nonkahler-locus-analytic-bounded} shows that \(E_{\mathrm{nK}}(\alpha)\) is analytic, and Theorem~\ref{ana:nonkahler-locus-analytic} gives a K\"ahler current with analytic singularities in both conventions whose polar set is \(E_{\mathrm{nK}}(\alpha)\). If
\(V\not\subset E_{\mathrm{nK}}(\alpha)\), then \(\alpha\) is big and, by the lemma of Collins--Tosatti on restricted volumes off
the non-K\"ahler locus, \(\operatorname{RV}_{X\mid V}(\alpha)\) is the
supremum of \(\int_{V_{\mathrm{reg}}}(T|_{V_{\mathrm{reg}}})_{\mathrm{ac}}^k\) over
K\"ahler currents with analytic singularities in the sense of
Definition~\ref{as:smooth-off-pole-type} not
containing \(V\) in their polar sets; such a current \(T\ge\delta\omega\) exists.
Its polar set is analytic, so by Theorem~\ref{ana:pure-dimension} \(V\) has a regular
point of dimension \(k=\dim V\) off it, and near that point the integral is at least
\(\delta^k\int\omega^k>0\). The statements on
modifications are the Collins--Tosatti transformation rules for the non-nef and
non-K\"ahler loci, together with the facts that blowups of compact K\"ahler
manifolds along smooth centers are compact K\"ahler and that bigness is preserved
by pullback. Bimeromorphic invariance of restricted volumes follows by induction
on the number of blowups: for no blowup it is
Lemma~\ref{ana:rv-kahler-independence}; for \(\mu=f\circ g\) with \(g\) one
blowup, Lemma~\ref{ana:blowup-surjective} gives \(V\not\subset f(\operatorname{Exc} f)\)
and \(\widetilde V^{f}\not\subset g(\operatorname{Exc} g)\), the induction
hypothesis and Lemma~\ref{ana:rv-birational-blowup} give the two equalities of
restricted volumes through any K\"ahler form on the intermediate manifold, and
Lemma~\ref{ana:strict-transform-composition} identifies
\(\widetilde{(\widetilde V^{f})}^{g}\) with \(\widetilde V^{\mu}\).
\end{proof}

\begin{lemma}[Modifications: compactness, K\"ahler forms, bigness and the non-K\"ahler locus]
\label{ana:modification-big-nonkahler}
\uses{ana:blowups,ana:nonkahler-locus,ana:blowup-compact,ana:blowup-kahler-correction}
\source{ct22-restricted-volumes:pull,transcendental-nakamaye-copy4-20260911:sec:nK}
\lean{TranscendentalNakamaye.Pending.LocusProperties.compactSpace_and_kahler_of_isBlowupSequence,TranscendentalNakamaye.Pending.DominationLogResolution.compactSpace_and_kahler_of_isBlowupSequence,TranscendentalNakamaye.Pending.LocusProperties.isBigClass_pullback,TranscendentalNakamaye.Pending.LocusProperties.nonKahlerLocus_pullback}
\leanok
Let \(X\) be a compact K\"ahler manifold and \(\mu:\widetilde X\to X\) a finite composition of
blowups with smooth centres. Then \(\widetilde X\) is compact K\"ahler. If
\(\alpha=\{\theta\}\) is big, then \(\mu^*\alpha\) is big and
\[
 E_{\mathrm{nK}}(\mu^*\alpha)=\mu^{-1}(E_{\mathrm{nK}}(\alpha))\cup\operatorname{Exc}(\mu).
\]
\end{lemma}
\begin{proof}
\uses{ana:blowup-compact,ana:blowup-kahler-correction}
\leanok
Induction on the number of blowups; for the identity there is nothing to prove, and for
\(\mu=f\circ g\) with \(g\) one blowup the statements compose, the exceptional loci through
\(\operatorname{Exc}(f\circ g)=\operatorname{Exc}(g)\cup g^{-1}(\operatorname{Exc}f)\). For one
blowup \(g:Z\to Y\) along a centre \(C\): \(Z\) is compact by
Lemma~\ref{ana:blowup-compact}, and \(g^*\omega+\varepsilon\Theta\) is K\"ahler for small
\(\varepsilon>0\), where \(\Theta=dd^c\log(U\circ g)\) is the form of the Blanchard weight \(U\)
of the centre (Lemma~\ref{ana:blowup-kahler-correction}). If \(\theta+dd^cu\ge\omega'\) is a
K\"ahler current, then
\(g^*\theta+dd^c(u\circ g+\varepsilon\log(U\circ g))\ge g^*\omega'+\varepsilon\Theta\), which is
a K\"ahler current for small \(\varepsilon\); so \(g^*\alpha\) is big. For the non-K\"ahler
locus, a K\"ahler current of \(\alpha\) with zero Lelong number at \(g(z)\), \(z\notin
\operatorname{Exc}g\), pulls back in this way to one of \(g^*\alpha\) with zero Lelong number at
\(z\), which gives \(\subseteq\). Conversely, a K\"ahler current of \(g^*\alpha\) has a potential
of the form \(\psi\circ g+c\log(U\circ g)\) plus a smooth function, with \(c>0\) and \(\psi\) the
potential of a K\"ahler current of \(\alpha\) obtained by pushing it forward across the centre,
of codimension at least two. At a point of \(\operatorname{Exc}g\) the term \(\log(U\circ g)\)
has a logarithmic pole, and at \(z\notin\operatorname{Exc}g\) with
\(g(z)\in E_{\mathrm{nK}}(\alpha)\) the Lelong number at \(z\) is at least
\(\nu(\psi,g(z))>0\); this gives \(\supseteq\).
\end{proof}

\begin{proposition}[A connected manifold is irreducible]
\label{ana:manifold-irreducible}
\uses{fpgeom:complex-analytic-subvariety,ana:nonkahler-locus,ana:blowups}
\source{demailly-cadg:II.4.2}
\lean{TranscendentalNakamaye.Pending.AnalyticIrreducibleComponents.isComplexAnalyticallyIrreducible_univ,TranscendentalNakamaye.Pending.AnalyticIrreducibleComponents.regularAtDim_univ,TranscendentalNakamaye.Pending.AnalyticIrreducibleComponents.exists_subvariety_univ,TranscendentalNakamaye.Pending.AnalyticIrreducibleComponents.eq_of_regularAtDim_univ,TranscendentalNakamaye.Pending.AnalyticIrreducibleComponents.dim_eq_of_carrier_eq_univ,TranscendentalNakamaye.Pending.AnalyticIrreducibleComponents.isOpen_of_forall_regularAtDim_self,TranscendentalNakamaye.Pending.AnalyticIrreducibleComponents.eq_univ_of_isClosedConnectedSubmanifold_self}
\leanok
Let \(X\) be a connected complex manifold of dimension \(n\).
\begin{enumerate}
\item \(X\) is an irreducible analytic subset of itself, of dimension \(n\), and
      \(n\) is the only dimension of a regular point of \(X\).
\item A closed connected complex submanifold of \(X\) of dimension \(n\) is \(X\).
\end{enumerate}
\end{proposition}
\begin{proof}
\leanok
(1) Every point of \(X\) is a regular point of dimension \(n\). If
\(X=B\cup C\) with \(B,C\) analytic and \(B\ne X\), then \(B\) has empty
interior: an analytic subset containing a nonempty open set contains, by the
identity theorem applied to its local equations along chains of coordinate balls,
the whole connected manifold. Hence the closed set \(C\) contains the dense set
\(X\setminus B\), so \(C=X\). A flattening of \(X\) of dimension \(k\) at a point
maps an open neighbourhood of that point onto an open set contained in
\(\mathbb C^k\times\{0\}\); a nonempty open subset of \(\mathbb C^n\) lies in that
plane only if \(k=n\).

(2) At each of its points a submanifold of dimension \(n\) contains a
neighbourhood of that point, so it is open. It is also closed and nonempty, hence
equal to the connected manifold \(X\).
\end{proof}

\begin{proposition}[Dimension strata are analytic]
\label{ana:dimension-strata}
\uses{fpgeom:complex-analytic-subvariety,lag:local-statement,lag:regular-point-basics,lag:regular-locus-dense}
\source{demailly-cadg:II.5.3}
\lean{TranscendentalNakamaye.Pending.PureDimension.regularStratum,TranscendentalNakamaye.Pending.AnalyticLocalStatement.LocalDecomposition.closure_inter_nbhd_eq,TranscendentalNakamaye.Pending.PureDimension.isComplexAnalyticSubset_closure_regularStratum}
\leanok
Let \(A\) be an analytic subset of a complex manifold \(M\) and, for each \(k\), let
\(R_k\) be the set of regular points of \(A\) of dimension \(k\). For every \(k\) the
closure \(\overline{R_k}\) is an analytic subset of \(M\); near a point \(p\in A\),
with the data of Theorem~\ref{lag:local-statement} in a chart, it equals the union of
those \(A_i\) with \(d_i=k\). For every \(m\), the closure of
\(\bigcup_{j\ne m}R_j\) is an analytic subset of \(M\).
\end{proposition}
\begin{proof}
\leanok
\(\overline{R_k}\subset A\), so near points outside the closed set \(A\) it is empty.
Let \(p\in A\) and take \(\Omega\), \(A_i\), \(d_i\) as in
Theorem~\ref{lag:local-statement} in a chart at \(p\). We show
\(\overline{R_k}\cap\Omega=\bigcup_{d_i=k}A_i\).

Let \(x\in R_k\cap\Omega\). By Theorem~\ref{lag:regular-locus-dense} there is an
open \(V\subset\Omega\) around \(x\) in which every point of \(A\) is regular of
dimension \(k\). Let \(V'\subset V\) be any open neighbourhood of \(x\). The nonempty
space \(A\cap V'\) is the finite union of its relatively closed subsets
\(A_i\cap V'\); a finite union of closed sets with empty interior has empty
interior, so some \(A_i\cap V'\) contains a nonempty relatively open subset of
\(A\cap V'\): there are \(y\) and an open \(V_y\subset V'\) around \(y\) with
\(\emptyset\ne A\cap V_y\subset A_i\), so \(A_i\cap V_y=A\cap V_y\). The point \(y\) is a
regular point of \(A\) of dimension \(k\), hence of \(A_i\)
(Lemma~\ref{lag:regular-point-basics}(1)), so \(d_i=k\) by
Theorem~\ref{lag:local-statement}(2). Thus every neighbourhood of \(x\) meets
\(\bigcup_{d_i=k}A_i\), which is closed in \(\Omega\), so \(x\) belongs to it. As
\(\Omega\) is open, \(\overline{R_k}\cap\Omega\subset\overline{R_k\cap\Omega}\cap\Omega
\subset\bigcup_{d_i=k}A_i\).

Conversely let \(d_i=k\), \(y\in A_i\), and \(V\subset\Omega\) an open neighbourhood
of \(y\). By Theorem~\ref{lag:local-statement}(2) there is \(z\in A_i\cap V\) regular
for \(A_i\) of dimension \(k\), and by Theorem~\ref{lag:regular-locus-dense} an open
\(V_z\subset V\) around \(z\) in which every point of \(A_i\) is regular for \(A_i\)
of dimension \(k\). By Theorem~\ref{lag:local-statement}(3), the finite union
\(\bigcup_{l\ne i}A_i\cap A_l\) has empty interior in \(A_i\), so the nonempty
relatively open set \(A_i\cap V_z\) contains a point \(z'\notin\bigcup_{l\ne i}A_l\).
The \(A_l\) are closed in \(\Omega\), so an open \(V''\subset V_z\) around \(z'\)
satisfies \(A\cap V''=A_i\cap V''\), and \(z'\in R_k\cap V\) by
Lemma~\ref{lag:regular-point-basics}(1). Hence \(y\in\overline{R_k}\).

So near every point of \(M\), \(\overline{R_k}\) coincides with a finite union of
analytic sets, and is analytic. Finally \(R_j=\emptyset\) for \(j>\dim M\), so the
closure of \(\bigcup_{j\ne m}R_j\) is the finite union \(\bigcup_{j\ne m,\,j\le\dim M}\overline{R_j}\),
again analytic.
\end{proof}

\begin{theorem}[Irreducible analytic sets are pure-dimensional]
\label{ana:pure-dimension}
\uses{fpgeom:complex-analytic-subvariety,ana:dimension-strata,lag:regular-locus-dense,lag:regular-point-basics}
\source{demailly-cadg:II.5.3}
\lean{TranscendentalNakamaye.Pending.PureDimension.eq_dim_of_regularAtDim,TranscendentalNakamaye.Pending.AnalyticIrreducibleComponents.eq_dim_of_regularAtDim,TranscendentalNakamaye.Pending.AnalyticIrreducibleComponents.exists_mem_regularLocus_not_mem_of_not_subset}
\leanok
Let \(V\) be an irreducible analytic subset of a complex manifold \(M\) with a
regular point of dimension \(m\). Then every regular point of \(V\) has dimension
\(m\). Consequently, if \(A\) is an analytic subset of \(M\) with \(V\not\subset A\),
then \(V\setminus A\) contains a regular point of \(V\) of dimension \(m\).
\end{theorem}
\begin{proof}
\leanok
Let \(R_j\) be the regular points of \(V\) of dimension \(j\). By
Theorem~\ref{lag:regular-locus-dense}, \(V=\overline{\bigcup_jR_j}\), and the union is
finite, so \(V=\overline{R_m}\cup T\) with \(T=\overline{\bigcup_{j\ne m}R_j}\).
Both sets are analytic (Proposition~\ref{ana:dimension-strata}) and contained in
\(V\), so by irreducibility \(V=\overline{R_m}\) or \(V=T\).

Around every \(x\in R_j\) there is an open \(U_x\) in which every point of \(V\) is
regular of dimension \(j\) (Theorem~\ref{lag:regular-locus-dense}), hence of no
other dimension (Lemma~\ref{lag:regular-point-basics}(2)); so \(U_x\) meets no
\(R_{j'}\) with \(j'\ne j\). Taking \(x\in R_m\ne\emptyset\) shows \(x\notin T\), so
\(V\ne T\) and \(V=\overline{R_m}\). If \(x\in R_k\), then \(x\in\overline{R_m}\)
forces \(U_x\cap R_m\ne\emptyset\), hence \(k=m\).

For the consequence, \(V\setminus A\) is a nonempty relatively open subset of
\(V\); by Theorem~\ref{lag:regular-locus-dense} it contains a regular point, whose
dimension is \(m\).
\end{proof}

\begin{proposition}[Closures of differences are analytic]
\label{ana:closure-difference-analytic}
\uses{fpgeom:complex-analytic-subvariety,lag:germ-decomposition,lag:thin-subgerm}
\source{demailly-cadg:II.5.3}
\lean{TranscendentalNakamaye.Pending.ClosureDifferenceAnalytic.isComplexAnalyticSubset_closure_diff}
\leanok
Let \(X,Y\) be analytic subsets of a complex manifold \(M\). Then
\(\overline{X\setminus Y}\) is an analytic subset of \(M\).
\end{proposition}
\begin{proof}
\leanok
Near points outside the closed set \(X\) the set is empty. Let \(p\in X\) and work in a
chart centred at \(p\). By Theorem~\ref{lag:germ-decomposition}, near \(p\),
\(X=X_1\cup\dots\cup X_r\) with irreducible analytic germs \(X_i\), each represented by
the common zero set of finitely many holomorphic functions. Let \(I\) be the set of
\(i\) with \(X_i\not\subset Y\) near \(p\). For \(i\in I\),
Proposition~\ref{lag:thin-subgerm} (with \(B=Y\)) gives \(\Omega_i\) such that
\(X_i\cap Y\cap\Omega_i\) has empty interior relative to \(X_i\cap\Omega_i\). Let
\(\Omega\) be an open ball around \(p\) contained in every \(\Omega_i\), in the
domains of the defining functions, in a neighbourhood where the decomposition holds,
and in a neighbourhood where \(X_i\subset Y\) for \(i\notin I\). We claim
\(\overline{X\setminus Y}\cap\Omega=\bigcup_{i\in I}X_i\cap\Omega\), a finite union of
analytic subsets of \(\Omega\).

If \(z\in(X\setminus Y)\cap\Omega\), then \(z\in X_i\) for some \(i\), and \(i\in I\)
since \(z\notin Y\). The right side is closed in \(\Omega\), so it contains
\(\overline{X\setminus Y}\cap\Omega\). Conversely let \(y\in X_i\cap\Omega\), \(i\in I\),
and \(O\subset\Omega\) an open neighbourhood of \(y\). The set \(O\cap X_i\) is a
nonempty relatively open subset of \(X_i\cap\Omega_i\), so it contains a point
\(z\notin Y\); then \(z\in(X\setminus Y)\cap O\), and \(y\in\overline{X\setminus Y}\).
\end{proof}

\begin{proposition}[Closures of connected components of the regular locus]
\label{ana:regular-component}
\uses{fpgeom:complex-analytic-subvariety,lag:local-statement-connected,lag:regular-locus-dense,lag:regular-point-basics}
\source{demailly-cadg:II.5.3}
\lean{TranscendentalNakamaye.Pending.RegularLocusComponent.regularPoints,TranscendentalNakamaye.Pending.RegularLocusComponent.component,TranscendentalNakamaye.Pending.RegularLocusComponent.exists_open_inter_subset_component,TranscendentalNakamaye.Pending.RegularLocusComponent.regularAtDim_of_mem_component,TranscendentalNakamaye.Pending.RegularLocusComponent.isComplexAnalyticSubset_closure_component,TranscendentalNakamaye.Pending.RegularLocusComponent.isComplexAnalyticallyIrreducible_closure_component,TranscendentalNakamaye.Pending.RegularLocusComponent.regularAtDim_closure_component}
\leanok
Let \(A\) be an analytic subset of a complex manifold \(M\), \(R\) the set of regular
points of \(A\) (of any dimension), and \(C\) a connected component of \(R\). Then:
\begin{enumerate}
\item every point \(c\in C\) has an open neighbourhood \(W_c\) with
      \(A\cap W_c\subset C\), and all points of \(C\) are regular of the same
      dimension \(k\);
\item \(\overline C\) is an analytic subset of \(M\);
\item \(\overline C\) is irreducible, and every point of \(C\) is a regular point of
      \(\overline C\) of dimension \(k\).
\end{enumerate}
\end{proposition}
\begin{proof}
\leanok
(1) Let \(c\in C\) be regular of dimension \(k\), with a flattening \(\Phi\) on a
neighbourhood \(W\). Choosing \(W=\Phi^{-1}(\text{ball})\), \(A\cap W\) is homeomorphic
to a ball in \(\mathbb C^k\), hence connected, and consists of regular points of
dimension \(k\) (Lemma~\ref{lag:regular-point-basics}(1)); so \(A\cap W\subset R\) is a
connected set meeting \(C\), whence \(A\cap W\subset C\). The dimension is therefore
locally constant on \(C\) (Lemma~\ref{lag:regular-point-basics}(2)), and constant as
\(C\) is connected.

(2) Near points outside the closed set \(A\) there is nothing to prove. Let \(p\in A\)
and take \(\Omega\), \(A_i\), \(d_i\), \(\Delta_i\), \(G_i\) as in
Theorem~\ref{lag:local-statement-connected} in a chart at \(p\); let
\(\Omega_1=\Omega\cap\bigcap_i\Delta_i\), an open neighbourhood of \(p\). Each \(G_i\) is a connected subset of \(R\), so it lies
in a single connected component. Let \(J=\{i: G_i\subset C\}\). We show
\(\overline C\cap\Omega_1=\bigcup_{i\in J}A_i\cap\Omega_1\), a finite union of analytic
sets. If \(i\in J\), then \(A_i\cap\Omega_1\subset A_i\cap\Delta_i\subset\overline{G_i}\subset
\overline C\). Conversely let \(c\in C\cap\Omega_1\) and \(O\subset\Omega_1\cap W_c\) an
open neighbourhood of \(c\). The nonempty space \(A\cap O\subset C\) is the finite union
of the relatively closed sets \(A_i\cap O\), so some \(A_i\cap O\) has nonempty
interior \(U\) relative to \(A\cap O\). As \(U\) is relatively open in
\(A_i\cap\Delta_i\) and \(G_i\) is dense there, \(G_i\) meets \(U\subset C\), so
\(i\in J\) and \(O\) meets \(A_i\). Hence \(c\) lies in the closure of
\(\bigcup_{i\in J}A_i\cap\Omega_1\), which is closed in \(\Omega_1\). Since \(\Omega_1\) is
open, \(\overline C\cap\Omega_1\subset\overline{C\cap\Omega_1}\cap\Omega_1\) lies in the
union too.

(3) By (1), near \(c\in C\) the sets \(\overline C\) and \(A\) coincide, so \(c\) is a
regular point of \(\overline C\) of dimension \(k\). Let \(\overline C=B_1\cup B_2\) with
analytic \(B_j\subset\overline C\). For an analytic set \(B\), the interior \(U\) of
\(B\cap C\) relative to \(C\) is closed in \(C\): if \(c\in C\) lies in the closure of
\(U\), take a flattening \(\Phi\) at \(c\) mapping \(A\cap W\) onto a ball \(D\) in
\(\mathbb C^k\times\{0\}\), and local equations \(f_1,\dots,f_s\) of \(B\) near \(c\);
the functions \(f_j\circ\Phi^{-1}\) are holomorphic on the connected \(D\) and vanish on
the nonempty open set \(\Phi(U\cap W)\), hence on \(D\) by the identity theorem, so
\(A\cap W\subset B\) and \(c\in U\). Thus \(U=C\) or \(U=\emptyset\). If \(B_1\cap C\) and
\(B_2\cap C\) both had empty interior in \(C\), then \(C\setminus B_1\) would be open and
dense in \(C\), hence meet the nonempty open set \(C\setminus B_2\), contradicting
\(C\subset B_1\cup B_2\). So some \(B_j\supset C\), and \(B_j\supset\overline C\) since
\(B_j\) is closed. As \(C\ne\emptyset\), \(\overline C\) is nonempty; it is irreducible.
\end{proof}

\begin{lemma}[Maximality from an open piece]
\label{ana:component-maximal}
\uses{ana:closure-difference-analytic}
\source{demailly-cadg:II.5.3}
\lean{TranscendentalNakamaye.Pending.ClosureDifferenceAnalytic.eq_of_subset_of_inter_open_subset}
\leanok
Let \(V\subset Z\) be analytic subsets of a complex manifold with \(Z\) irreducible, and
suppose that some open set \(O\) satisfies \(\emptyset\ne Z\cap O\subset V\). Then
\(V=Z\).
\end{lemma}
\begin{proof}
\leanok
\(Z=V\cup\overline{Z\setminus V}\), and both sets are analytic
(Proposition~\ref{ana:closure-difference-analytic}) and contained in \(Z\). By
irreducibility \(Z=V\) or \(Z=\overline{Z\setminus V}\). In the second case a point
\(y\in Z\cap O\) lies in \(\overline{Z\setminus V}\), but \(O\) is a neighbourhood of
\(y\) disjoint from \(Z\setminus V\), a contradiction.
\end{proof}

\begin{proposition}[Irreducible components]
\label{ana:irreducible-components}
\uses{fpgeom:complex-analytic-subvariety,ana:regular-component,ana:component-maximal,lag:local-statement-connected,ana:pure-dimension}
\source{demailly-cadg:II.5.3}
\lean{TranscendentalNakamaye.Pending.AnalyticIrreducibleComponents.exists_isIrreducibleComponent_pos_dim_of_accPt,TranscendentalNakamaye.Pending.PositiveDimensionalComponent.exists_isIrreducibleComponent_pos_dim_of_accPt}
\leanok
Let \(X\) be a complex manifold, \(A\subset X\) an analytic subset and \(x\in A\) a
point which is not isolated in \(A\). Then \(x\) lies on a positive-dimensional
irreducible component of \(A\): an irreducible closed analytic subset
\(V\subset A\) with a regular point of positive dimension \(\dim V\), containing \(x\),
such that every irreducible analytic \(Z\) with \(V\subset Z\subset A\) equals \(V\).
\end{proposition}
\begin{proof}
\leanok
Take \(\Omega\), \(A_i\), \(d_i\), \(\Delta_i\), \(G_i\) as in
Theorem~\ref{lag:local-statement-connected} at \(x\); every \(A_i\) contains \(x\), and
\(A\cap\Omega=\bigcup_iA_i\). If \(d_i=0\), the points of \(G_i\) are isolated in \(A\) (a regular
point of dimension \(0\) has a neighbourhood meeting \(A\) only in itself), so the
connected set \(G_i\) is a single point, and \(A_i\cap\Delta_i=\overline{G_i}\cap\Delta_i\)
is that point, namely \(x\). If all \(d_i\) were \(0\), then \(A\) would meet the
neighbourhood \(\Omega\cap\bigcap_i\Delta_i\) of \(x\) only in \(x\), and \(x\) would be isolated. So some \(i\) has
\(d_i>0\).

\(G_i\) is a connected set of regular points of \(A\), contained in a connected
component \(C\) of the regular locus, whose points have dimension \(d_i\)
(Proposition~\ref{ana:regular-component}(1)). Put \(V=\overline C\): it is closed,
analytic and irreducible, contained in \(A\), with regular points of dimension
\(d_i>0\) (Proposition~\ref{ana:regular-component}); by
Theorem~\ref{ana:pure-dimension} its dimension is \(d_i\). Also
\(x\in A_i\cap\Delta_i\subset\overline{G_i}\subset V\). Finally let \(Z\) be irreducible
analytic with \(V\subset Z\subset A\). For \(c\in C\) and \(W_c\) as in
Proposition~\ref{ana:regular-component}(1), \(Z\cap W_c\subset A\cap W_c\subset C\subset V\)
and \(c\in Z\cap W_c\), so \(Z=V\) by Lemma~\ref{ana:component-maximal}.
\end{proof}

\begin{proposition}[An irreducible component fills the set near some of its points]
\label{ana:component-local-filling}
\uses{ana:nonkahler-locus,ana:regular-component,lag:local-statement-connected,ana:irreducible-open-density,lag:regular-locus-dense,ana:pure-dimension}
\source{demailly-cadg:II.5.3}
\lean{TranscendentalNakamaye.Pending.IrreducibleComponentFilling.componentSubvariety,TranscendentalNakamaye.Pending.IrreducibleComponentFilling.componentSubvariety_carrier,TranscendentalNakamaye.Pending.IrreducibleComponentFilling.componentSubvariety_subset,TranscendentalNakamaye.Pending.IrreducibleComponentFilling.exists_isOpen_subset_inter_subset_of_mem,TranscendentalNakamaye.Pending.IrreducibleComponentFilling.exists_isOpen_inter_subset,TranscendentalNakamaye.Pending.IrreducibleComponentFilling.exists_mem_regularLocus_isOpen_inter_subset}
\leanok
Let \(A\) be an analytic subset of a complex manifold \(M\) and \(V\) an irreducible
component of \(A\): a closed irreducible analytic subset \(V\subset A\) such that every
closed irreducible analytic \(Z\) with \(V\subset Z\subset A\) equals \(V\). Then for every
\(p\in V\) and every open neighbourhood \(O\) of \(p\) there is an open set \(U\subset O\)
with
\[
 V\cap U\neq\emptyset\qquad\text{and}\qquad A\cap U\subset V .
\]
In particular some open set \(U\) meets \(V\) with \(A\cap U\subset V\), and every open set
\(O\) meeting \(V\) contains such a \(U\) around a point at which \(V\) is regular of
dimension \(\dim V\).
\end{proposition}
\begin{proof}
\leanok
As \(p\in V\subset A\), take \(\Omega\), \(A_i\), \(d_i\), \(\Delta_i\), \(G_i\)
(\(1\le i\le N\)) as in Theorem~\ref{lag:local-statement-connected} at \(p\): every \(A_i\)
contains \(p\), \(A\cap\Omega=\bigcup_iA_i\), and \(G_i\) is a connected set of regular
points of \(A\) of dimension \(d_i\), dense in \(A_i\cap\Delta_i\). In particular \(N\ge1\),
and \(p\in A_i\cap\Delta_i\subset\overline{G_i}\), so we may choose \(c_i\in G_i\cap O\).
Let \(C_i\) be the connected component of the regular locus of \(A\) through \(c_i\). It
contains the connected set \(G_i\) of regular points, and by
Proposition~\ref{ana:regular-component} its closure \(\overline{C_i}\) is a closed
irreducible analytic subset of \(A\), regular of dimension \(d_i\) at \(c_i\). Hence
\(A_i\cap\Delta_i\subset\overline{G_i}\subset\overline{C_i}\), and on the open
neighbourhood \(\Omega_1=\Omega\cap\bigcap_i\Delta_i\) of \(p\),
\[
 V\cap\Omega_1\subset A\cap\Omega_1\subset\bigcup_{i=1}^N\overline{C_i}.
\]
By Lemma~\ref{ana:irreducible-open-density} (finitely many analytic sets, listed as a
sequence with repetitions), \(V\subset\overline{C_i}\) for some \(i\). As
\(\overline{C_i}\subset A\) is closed irreducible analytic, maximality of \(V\) gives
\(\overline{C_i}=V\). By Proposition~\ref{ana:regular-component}(1), \(c_i\) has an open
neighbourhood \(W\) with \(A\cap W\subset C_i\subset V\). Then \(U=W\cap O\) is open,
contains \(c_i\in C_i\subset V\), and \(A\cap U\subset V\).

For the last assertion, take \(p\in V\cap O\) and \(U\subset O\) as above. The open set
\(U\) meets \(V\), so it contains a regular point \(x\) of \(V\)
(Theorem~\ref{lag:regular-locus-dense}), whose dimension is \(\dim V\) by
Theorem~\ref{ana:pure-dimension}.
\end{proof}
Demailly's Theorem~II.5.3 describes all global irreducible components at once (closures of the
connected components of the regular locus, a locally finite family covering \(A\)); the
proposition is the consequence of that description which is needed below, proved from the local
statement and maximality only.

\begin{lemma}[Maximality transported along a map, without images]
\label{ana:component-preimage-maximal}
\uses{ana:component-local-filling,ana:irreducible-open-density,ana:nonkahler-locus}
\lean{TranscendentalNakamaye.Pending.IrreducibleComponentFilling.isIrreducibleComponent_of_preimage}
\leanok
Let \(\mu:Y\to M\) be a continuous map of complex manifolds, \(V\) an irreducible component
of an analytic subset \(A\subset M\), \(E\subset Y\) closed and
\(S\subset\mu^{-1}(A)\cup E\). Let \(\widetilde V\subset S\) be a closed irreducible analytic
subset of \(Y\) with \(\mu^{-1}(V)\setminus E\subset\widetilde V\), such that every point at
which \(V\) is regular of dimension \(\dim V\) is the image of a point of
\(\widetilde V\setminus E\). Then \(\widetilde V\) is an irreducible component of \(S\).
\end{lemma}
\begin{proof}
\leanok
Let \(Z\) be closed irreducible analytic with \(\widetilde V\subset Z\subset S\).
Proposition~\ref{ana:component-local-filling} gives an open \(U\subset M\) with
\(A\cap U\subset V\) containing a point \(x\) at which \(V\) is regular of dimension
\(\dim V\); write \(x=\mu(y)\) with \(y\in\widetilde V\setminus E\). The set
\(N=\mu^{-1}(U)\setminus E\) is open and contains \(y\in Z\). If \(z\in Z\cap N\), then
\(z\in S\setminus E\subset\mu^{-1}(A)\), so \(\mu(z)\in A\cap U\subset V\) and
\(z\in\mu^{-1}(V)\setminus E\subset\widetilde V\). Thus \(\emptyset\ne Z\cap N\subset
\widetilde V\), and Lemma~\ref{ana:irreducible-open-density} gives \(Z\subset\widetilde V\),
so \(Z=\widetilde V\).
\end{proof}
For a modification \(\mu\) with exceptional set \(E\), the strict transform \(\widetilde V\) of a
component \(V\) of \(E_{\mathrm{nK}}(\alpha)\) and \(S=E_{\mathrm{nK}}(\mu^*\alpha)=
\mu^{-1}(E_{\mathrm{nK}}(\alpha))\cup E\), the lemma is the maximality half of
Lemma~\ref{ana:resolution}, proved there this way, without the image \(\mu(Z)\) and so without
Lemma~\ref{ana:proper-image-analytic}.

\begin{lemma}[The non-K\"ahler locus has no isolated points]
\label{ana:no-isolated-points}
\uses{ana:locus-properties}
\source{transcendental-nakamaye-copy4-20260911:lem:no-isolated-points}
\lean{TranscendentalNakamaye.Pending.NoIsolatedPoints.accPt_nonKahlerLocus}
\leanok
The non-K\"ahler locus of a big class on a connected compact K\"ahler
manifold of positive dimension has no isolated points.
\end{lemma}
\begin{proof}
\uses{ana:gluing,levi:norm-square}
\leanok
Suppose that $x$ is isolated in $E_{\mathrm{nK}}(\alpha)$. Choose a
K\"ahler current $R=\theta+dd^cF\ge\delta\omega$, with $\delta>0$,
whose analytic polar set is exactly $E_{\mathrm{nK}}(\alpha)$.
Choose a coordinate ball $B$ whose closure meets that polar set only at $x$.
Since $F$ has analytic singularities, it is locally bounded off the polar set, and $F(x)=-\infty$.
On the compact closure of a smaller coordinate ball, a sufficiently large
$A>0$ satisfies $\theta+dd^c(A|z|^2)\ge\delta'\omega$ for some
$\delta'>0$. Relabel this smaller ball as $B$.
Choose $x\in U_1\Subset U_2\Subset B$. On the compact collar
$\overline U_2\setminus U_1$, $F$ has a finite lower bound, whereas
$A|z|^2$ is bounded above. Thus the global branch $F$ and the local
quadratic branch satisfy the collar hypothesis of Proposition~
\ref{ana:gluing}. The other hypotheses follow from the two strict
curvature lower bounds, the smoothness of the quadratic branch, and
$F(x)=-\infty$. That proposition produces a K\"ahler current in
$\alpha$ which is smooth near $x$, contradicting
$x\in E_{\mathrm{nK}}(\alpha)$.
\end{proof}

\begin{theorem}[Trace mass from currents with minimal singularities]
\label{ana:trace-minimal-mass}
\lean{TranscendentalNakamaye.Pending.MinimalPotentialOrdinaryRestrictionRightLimit.exists_tendsto_ordinary_minimal_restriction_mass_of_not_mem_nonNefLocus}
\uses{ana:geometric-setting,trace:actual-minimal-ordinary-restriction-right-limit,ana:nonnef-lelong-locus,ana:submanifold-complex-structure,pc:kahler-restriction,trace:actual-positive-normalized-trace-domain-value,trace:actual-positive-normalized-trace-models,trace:actual-numerical-trace-mass,trace:diagonal-mass,trace:minimal-restriction-mass,trace:perturbative-mass-limit}
\source{transcendental-nakamaye-copy4-20260911:thm:trace-formula,xia-sgpt:page-0277,xia-sgpt:page-0286,xia-sgpt:page-0287,xia-sgpt:page-0284}
\leanok
Let $\alpha$ be big and choose a smooth representative $\theta\in\alpha$.
Let $V\subset X$ be a closed connected complex submanifold of positive
dimension $k$, with $V\not\subset E_{\mathrm{nn}}(\alpha)$.
For $T_{\min,\epsilon}=\theta+\epsilon\omega+dd^cV_{\theta+\epsilon\omega}$,
\[
 \lim_{\epsilon\downarrow0}\int_V
 \left\langle(T_{\min,\epsilon}|_V)^k\right\rangle
 =\int_V\left\langle
 (\theta|_V+dd^c\operatorname{Tr}^{\theta}_V(V_\theta))^k\right\rangle.
\]
The right-hand side is zero if the positive-mass trace is undefined.
\end{theorem}
\begin{proof}
\uses{trace:minimal-restriction-mass,trace:perturbative-mass-limit}
For each $\epsilon>0$, identify the mass of the ordinary restriction of the
minimal potential with the mass of its normalized trace.  The perturbative
trace-mass theorem then takes the limit and identifies it with the numerical
$\theta$-trace mass of $V_\theta$, including the undefined-trace convention.
\end{proof}

The integral of the ordinary restricted minimal current in this statement is read through its actual potential $V_{\theta+\epsilon\omega}\circ\iota$ and the actual pulled-back form.  The original big class and $V\not\subset E_{\rm nn}(\alpha)$ provide ambient generic-zero via the compiled non-nef Lelong-locus identity; they are not supplemented by an incoming trace, mass, or target-current certificate.  The actual smooth submanifold atlas and K\"ahler restriction supply the holomorphic inclusion and target background.  The attached theorem produces the entire genuine ordinary restricted family before all filters and positive real paths to zero, identifies its masses with the actual normalized numerical values, and proves the displayed limit including the original undefined-positive-trace zero convention.  The two proof dependencies below retain their exact mathematical statements and have separate approved whole-contract packets.  This statement is the second-to-last mass identity in Xia (8.12) used in Copy4's trace-formula proof; the restricted-volume equality and the separate CT22 minimal-mass representation remain outside this closure.

\begin{theorem}[Trace formula]
\label{ana:trace-formula}
\uses{ana:restricted-volume,ana:trace-minimal-mass,ana:actual-restricted-volume-trace-formula}
\source{transcendental-nakamaye-copy4-20260911:thm:trace-formula,xia-sgpt:page-0286,xia-sgpt:page-0287}
\lean{TranscendentalNakamaye.Pending.ActualRestrictedVolumeTraceFormula.numericalRestrictedVolume_eq_numericalTraceMass}
\leanok
Let $\alpha$ be big, $\theta\in\alpha$ smooth, and $V\subset X$ a
positive-dimensional closed connected complex submanifold with
$V\not\subset E_{\mathrm{nn}}(\alpha)$. Then
\[
 \operatorname{RV}_{X\mid V}(\alpha)=
 \int_V\left\langle
 (\theta|_V+dd^c\operatorname{Tr}^{\theta}_V(V_\theta))^k\right\rangle,
 \qquad k=\dim V,
\]
where the right-hand side is zero if the positive-mass trace is undefined.
\end{theorem}
\begin{proof}
\uses{ana:ct-minimal-mass,ana:trace-minimal-mass}
Theorem~\ref{ana:ct-minimal-mass} identifies the numerical restricted
volume with the limit of the restricted masses of $T_{\min,\epsilon}$.
Theorem~\ref{ana:trace-minimal-mass} identifies that same limit with the
stated trace mass, including the zero-mass convention. Combining the
two equalities gives the result.
\end{proof}

\begin{theorem}[Multiplier ideals of the trace]
\label{ana:trace-multiplier}
\lean{TranscendentalNakamaye.Pending.ActualTraceMultiplierInclusion.multiplierIdeal_le_restrictIdeal_of_isTraceRepresentativeAlong,TranscendentalNakamaye.Pending.NormalizedTraceMultiplierInclusion.multiplierIdeal_singularityEnvelope_le_restrictIdeal}
\uses{ana:geometric-setting,trace:representative,trace:actual-trace-representative-along-map,trace:actual-reader-multiplier-ideal-restriction,trace:actual-normalized-trace-multiplier-ideal-restriction,trace:closed-form-pullback,fpenv:envelopes,fpglobal:theta-psh,trace:diagonal-mass,ana:ideal-sheaves,ana:restricted-ideals}
\source{transcendental-nakamaye-copy4-20260911:thm:trace-multiplier,xia-sgpt:page-0283,xia-sgpt:page-0277}
\leanok
Let \(V\subset X\) be a closed connected complex submanifold and let
\(\varphi\) be quasi-plurisubharmonic with \(\nu(\varphi,V)=0\).
For every \(\lambda>0\) and every representative \(\tau\) of its trace,
\[
 \mathcal I(\lambda\tau)
 \subset\operatorname{Res}_V\mathcal I(\lambda\varphi).
\]
In particular this inclusion holds for
\(\tau=\operatorname{Tr}^{\theta}_V(\varphi)\), whenever that
positive-mass trace is defined.
\end{theorem}
\begin{proof}
\leanok
\uses{trace:qe-definition,trace:actual-antitone-native-global-p-order,trace:global-source-p-implies-multiplier-order,trace:i-order,trace:singular-global-p-order,trace:multiplier-restriction,trace:strong-openness,pc:trace-multiplier-inclusion,trace:singular-global-p-equivalence-original-envelope,trace:actual-reader-multiplier-ideal-restriction,trace:actual-normalized-trace-multiplier-ideal-restriction}
Choose a quasi-equisingular approximation $(\varphi_j)$ whose restrictions
converge to the trace.  The trace is more singular in $P$-order than every
$\varphi_j|_V$.  Given $\lambda'>\lambda>0$, quasi-equisingularity supplies
an index $j$ for which
$\mathcal I(\lambda'\varphi_j)\subseteq\mathcal I(\lambda\varphi)$.
$P$-order comparison and the restriction theorem give
\[
 \mathcal I(\lambda'\operatorname{Tr}_V(\varphi))
 \subseteq\mathcal I(\lambda'\varphi_j|_V)
 \subseteq\operatorname{Res}_V\mathcal I(\lambda'\varphi_j)
 \subseteq\operatorname{Res}_V\mathcal I(\lambda\varphi).
\]
Strong openness, followed by $P$-equivalence of representatives, yields the
claimed inclusion at exponent $\lambda$ and for the normalized trace.
\end{proof}

\begin{theorem}[Removable singularities of plurisubharmonic functions across analytic sets]
\label{ana:psh-removable-singularity}
\uses{ana:geometric-setting,as:logarithmic-psh,fpenv:germwise-disc-criterion,psh:decreasing-limit,fpenv:local-psh-supremum,lpot:existence-global}
\source{demailly-cadg:I.5.24,demailly-cadg:I.5.23,demailly-cadg:I.5.22}
\lean{TranscendentalNakamaye.Pending.PshRemovableSingularity.demaillyCadgI524_analyticRemovableThetaPsh,TranscendentalNakamaye.Pending.PshAnalyticRemovable.plurisubharmonicOn_extension_of_zeroLocus,TranscendentalNakamaye.Pending.PshAnalyticRemovable.thetaPshLocalClosure_of_isComplexAnalyticSubset,TranscendentalNakamaye.Pending.PshAnalyticRemovable.analyticRemovableThetaPsh,TranscendentalNakamaye.Pending.PshAnalyticRemovable.analyticRemovableThetaPsh_of_closed,TranscendentalNakamaye.Pending.PshAnalyticSetExtension.thetaPlurisubharmonic_extensionAcross_of_analyticRemovable,TranscendentalNakamaye.Pending.PshAnalyticSetExtension.thetaPlurisubharmonic_of_analyticRemovable,TranscendentalNakamaye.Pending.DemaillyRegularizationStep3Removable.thetaPlurisubharmonic_of_analyticPoleSet}
\leanok
Let \(X\) be a complex manifold, \(A\subset X\) a closed complex analytic subset, and
\(\theta\) a smooth real \(2\)-form on \(X\) which admits smooth local potentials; every
smooth closed real \((1,1)\)-form does (Corollary~\ref{lpot:existence-global}). Let \(v\) be
\(\theta\)-plurisubharmonic on \(X\setminus A\), and suppose every point of \(X\) has a
neighbourhood on which \(v\) is bounded above off \(A\). Let \(\widetilde v\) be the upper
semicontinuous regularization of the function equal to \(v\) on \(X\setminus A\) and to
\(-\infty\) on \(A\). Then \(\widetilde v\) is \(\theta\)-plurisubharmonic on a neighbourhood of
every point of \(A\): in every chart, around every point over \(A\), \(\widetilde v+\rho\) is
plurisubharmonic for a smooth local potential \(\rho\) of \(\theta\). If \(X\) is connected and
\(A\ne X\), then \(\widetilde v\in\operatorname{PSH}(X,\theta)\). In particular, on a connected
\(X\) with \(A\ne X\), a function \(u\) on \(X\) which satisfies the hypotheses on \(v\), equals
\(-\infty\) on \(A\) and tends to \(-\infty\) at every point of \(A\) from outside \(A\) lies in
\(\operatorname{PSH}(X,\theta)\).
\end{theorem}
\begin{proof}
\uses{as:logarithmic-psh,fpenv:germwise-disc-criterion,psh:decreasing-limit,fpenv:local-psh-supremum}
\leanok
\emph{Reduction to a coordinate domain.} Fix a chart and a point \(z\) over \(A\), and a
smooth potential \(\rho\) of \(\theta\) near \(z\). Shrink to an open set \(\Omega\ni z\) in
the chart on which \(\rho\) is defined and \(A\) is the common zero set of finitely many
holomorphic functions \(g_1,\dots,g_N\) on \(\Omega\). On \(\Omega\setminus A\) the function
\(\psi=v+\rho\) is plurisubharmonic, and near every point of \(\Omega\) it is bounded above
off \(A\), since \(\rho\) is continuous. Adding the continuous function \(\rho\) commutes with
upper semicontinuous regularization, so it suffices to show: the regularization
\(\widetilde\psi\) of (\(\psi\) on \(\Omega\setminus A\), \(-\infty\) on \(A\)) is
plurisubharmonic on \(\Omega\). Plurisubharmonicity is local, so we fix
\(x_0\in\Omega\) and work on a ball \(B\ni x_0\) in \(\Omega\) chosen below.

\emph{The functions \(v_k\).} Put \(w=\log\sum_i|g_i|^2\). It is plurisubharmonic on
\(\Omega\) (Lemma~\ref{as:logarithmic-psh}) and equals \(-\infty\) exactly on \(A\). Choose
\(B\) so small that \(w<c\) on \(B\) for a real \(c\) (upper semicontinuity) and
\(\psi\le C\) on \(B\setminus A\) for a real \(C\). For \(k\ge0\) let
\(\varepsilon_k=1/(k+1)\) and
\[
 v_k=\psi+\varepsilon_k\,(w-c)\ \text{ on }B\setminus A,\qquad v_k=-\infty\ \text{ on }A\cap B.
\]

\emph{Each \(v_k\) is plurisubharmonic on \(B\).} On \(B\setminus A\) it is a sum of
plurisubharmonic functions. At \(a\in A\cap B\), \(v_k(y)\to-\infty\) as \(y\to a\): off \(A\)
we have \(v_k\le C+\varepsilon_k(w-c)\), and \(w(y)\to-\infty\) because \(w\) is upper
semicontinuous with \(w(a)=-\infty\). Hence, for every \(m\in\mathbf N\), the truncation
\(\max(v_k,-m)\) is the constant \(-m\) on a neighbourhood of each point of \(A\cap B\), and is
plurisubharmonic near each point of \(B\setminus A\). Having a plurisubharmonic germ at every
point of \(B\), it satisfies the submean inequality on every closed complex affine disc in
\(B\) (Theorem~\ref{fpenv:germwise-disc-criterion}), so it is plurisubharmonic on \(B\). As
\(m\to\infty\) these truncations decrease to \(v_k\), which is therefore plurisubharmonic on
\(B\) (Theorem~\ref{psh:decreasing-limit}).

\emph{The envelope.} Since \(w-c<0\) on \(B\), we have \(v_k\le\psi\le C\) on
\(B\setminus A\) and \(v_k=-\infty\) on \(A\), so the family \((v_k)\) is uniformly bounded
above on \(B\), and the upper semicontinuous regularization of \(\sup_kv_k\) is
plurisubharmonic on \(B\) (Theorem~\ref{fpenv:local-psh-supremum}). At a point of
\(B\setminus A\), \(w\) is finite and \(\varepsilon_k(w-c)\) increases to \(0\), so
\(\sup_kv_k=\psi\) there; on \(A\cap B\) the supremum is \(-\infty\). So \(\sup_kv_k\) is
exactly the function (\(\psi\) off \(A\), \(-\infty\) on \(A\)) on \(B\), and its regularization
is \(\widetilde\psi\) on \(B\). This proves the local statement.

\emph{The global statement.} Off the closed set \(A\) the function \(v\) is
\(\theta\)-plurisubharmonic, hence upper semicontinuous, so \(\widetilde v=v\) on the open set
\(X\setminus A\) and \(\widetilde v\) is \(\theta\)-plurisubharmonic there; near the points of
\(A\) it is \(\theta\)-plurisubharmonic by the local statement. Being
\(\theta\)-plurisubharmonic is a local property. For the last claim, \(\widetilde u=u\): off \(A\)
as just said, and at a point of \(A\) the regularization is \(-\infty=u\) there, because \(u\)
tends to \(-\infty\) at that point from outside \(A\) and equals \(-\infty\) on \(A\).
\end{proof}
This is the analytic case of Demailly's Theorem~(5.24), proved by his argument for the
subharmonic Theorem~(5.23) (the supremum of \(v+\varepsilon u\), then regularization), with
\(u=w-c\) the function of Definition~(5.22) for which an analytic set with empty interior is
pluripolar. Demailly's theorem allows any closed pluripolar \(A\) and asserts that the
extension is \emph{unique}; uniqueness needs \(A\) to be Lebesgue-null and is not asserted
here. The local statement needs neither connectedness of \(X\) nor \(A\ne X\): at an interior
point of \(A\) the regularization is \(-\infty\) near the point, and the constant \(-\infty\)
is plurisubharmonic in the convention of this development.

\begin{theorem}[Smooth regularization of a bounded quasi-plurisubharmonic
potential]
\label{ana:bounded-potential-smoothing}
\uses{ana:geometric-setting,ana:demailly-algebraic-approximation}
\source{transcendental-nakamaye-copy4-20260911:lem:regularization,bou02-volume:2.6}
\lean{TranscendentalNakamaye.Pending.BoundedPotentialSmoothing.demailly1992_exists_chartSmooth_of_bounded_thetaPsh,TranscendentalNakamaye.Pending.BoundedPotentialSmoothingOfAlgebraicApproximation.exists_chartSmooth_of_bounded_thetaPsh}
\leanok
Let \((X,\omega)\) be a compact K\"ahler manifold, let \(\theta\) be a smooth
closed real \((1,1)\)-form and let \(u\in\operatorname{PSH}(X,\theta)\) satisfy
\(a\le u\le b\) for real constants \(a,b\). For every \(\varepsilon>0\) there is
a smooth \(h\colon X\to\mathbf R\) with
\[
 \theta+d d^ch+\varepsilon\omega\ge0,\qquad
 u\le h+\varepsilon,\qquad h\le b+\varepsilon .
\]
\end{theorem}
\begin{proof}
\uses{ana:demailly-algebraic-approximation}
\leanok
Apply Theorem~\ref{ana:demailly-algebraic-approximation} to \(\theta\) and the potential
\(u\) of \(\theta+dd^cu\ge0\), with \(\omega\) as the K\"ahler form. It gives \(u_k\) decreasing to \(u\),
with analytic singularities, and \(\delta_k\to0\) with \(\theta+\delta_k\omega+dd^cu_k\ge0\).
Since \(u_k\ge u\ge a\), no \(u_k\) has a pole, so each \(u_k\), being smooth off its
poles, is smooth; and \(u\le u_k\) everywhere. Choose \(k\) with \(\delta_k\le\varepsilon\);
then \(\theta+dd^cu_k+\varepsilon\omega\ge0\). Finally the continuous functions \(u_k\)
decrease to the upper semicontinuous \(u\le b<b+\varepsilon\) on the compact \(X\), so by
the one-sided Dini argument (for each \(x\) some \(u_k(x)<b+\varepsilon\); this persists on a
neighbourhood by continuity and for all later \(k\) by monotonicity, and finitely many
neighbourhoods cover \(X\)) some \(u_k<b+\varepsilon\) everywhere. Take \(h=u_k\) for \(k\)
large enough for both requirements.
\end{proof}
Only the first display is used downstream (in \ref{ana:residual-nef}); the two comparisons
come for free from clause \emph{(ii)} of the approximation theorem.

The first display alone is \emph{not} weaker than the \(D=0\) case of
\ref{ana:residual-nef}: taking \(D=0\), so that \(\psi_D=0\) and
\(\Theta_D=0\), a bounded \(u\in\operatorname{PSH}(X,\theta)\) has logarithmic
singularities along \(D\) with residual \(u\) itself, and nefness of
\(\{\theta\}\) is exactly that display. The two comparison inequalities are the
whole difference between the two statements, and a nefness assertion implies
neither of them.

\begin{lemma}[Nef residual class of logarithmic singularities]
\label{ana:residual-nef}
\uses{ana:geometric-setting}
\source{transcendental-nakamaye-copy4-20260911:lem:residual-nef,xia-sgpt:page-0048}
\lean{TranscendentalNakamaye.Pending.ResidualNef.nefWithRespectTo_sub_curvature,TranscendentalNakamaye.Pending.ResidualNef.exists_chartSmooth_isKahler_sub_curvature}
\leanok
Let \(X\) be a compact K\"ahler manifold, let \(\theta\) be a smooth
closed real \((1,1)\)-form, and suppose that
\(\varphi\in\operatorname{PSH}(X,\theta)\) has logarithmic
singularities along an effective rational divisor \(D\). Then
\(\{\theta\}-\{D\}\) is nef. If
\(\theta+d d^c\varphi\) is a K\"ahler current, then this residual
class is K\"ahler.
\end{lemma}
\begin{proof}
\uses{ana:psh-removable-singularity,ana:bounded-potential-smoothing,ana:divisor-weight-smooth-potential,ana:ideal-sheaves}
\leanok
Fix a smooth metric on \(\mathcal O(D)\), with weight \(\psi_D\) and curvature
\(\Theta_D\) (Definition~\ref{ana:ideal-sheaves}). That \(\varphi\) has
logarithmic singularities along \(D\) means exactly that
\(b:=\varphi-\psi_D\) is a bounded real function; say \(|b|\le K\).

\emph{Step 1: \(b\) is \((\theta-\Theta_D)\)-plurisubharmonic off
\(\operatorname{Supp}D\).} Away from the support, \(\psi_D\) is a genuine
smooth function and an exact local potential of \(-\Theta_D\)
(Lemma~\ref{ana:divisor-weight-smooth-potential}). So if \(\rho\) is a smooth local potential of \(\theta\) on a small
coordinate piece, then \(\psi_D+\rho\) is a smooth local potential of
\(\theta-\Theta_D\) there, and
\(b+(\psi_D+\rho)=\varphi+\rho\), which is plurisubharmonic because
\(\varphi\in\operatorname{PSH}(X,\theta)\).

\emph{Step 2: \(b\) has a global bounded representative.} The set
\(\operatorname{Supp}D\) is a closed complex analytic subset, and it is proper:
\(\varphi\) is not identically \(-\infty\), and where \(\varphi\) is finite so is
\(\psi_D\). By \ref{ana:psh-removable-singularity} the upper semicontinuous
regularization \(b^\star\) of \(b\) across \(\operatorname{Supp}D\) lies in
\(\operatorname{PSH}(X,\theta-\Theta_D)\). It still satisfies \(|b^\star|\le K\):
the upper bound because upper regularization does not exceed a pointwise bound,
and the lower bound because a proper analytic subset of a connected manifold has
empty interior, so every neighbourhood of a point of \(\operatorname{Supp}D\)
meets the complement, where \(b\ge-K\).

\emph{Step 3: nefness.} Apply \ref{ana:bounded-potential-smoothing} to
\(\theta-\Theta_D\) and \(b^\star\): for every \(\varepsilon>0\) there is a
smooth \(h\) with \(\theta-\Theta_D+d d^ch+\varepsilon\omega\ge0\). As
\(\theta-\Theta_D\) represents \(\{\theta\}-\{D\}\), that class is nef.

\emph{Step 4: the K\"ahler case.} Suppose \(\theta+d d^c\varphi\ge\omega'\) for
some K\"ahler form \(\omega'\). Then \(\varphi\in\operatorname{PSH}(X,
\theta-\omega')\), with the same logarithmic singularities along the same \(D\).
Steps 1--3 applied to \(\theta-\omega'\), and with \(\omega'\) itself used as the
reference K\"ahler form, give for \(\varepsilon=\tfrac12\) a smooth \(h\) with
\[
 (\theta-\omega')-\Theta_D+d d^ch+\tfrac12\omega'\ge0,
 \qquad\text{that is}\qquad
 \theta-\Theta_D+d d^ch\ \ge\ \tfrac12\omega'\ >\ 0 .
\]
So \(\{\theta\}-\{D\}\) contains a K\"ahler form. No comparison of two K\"ahler
forms is needed, because the reference form of the nefness statement was chosen
to be \(\omega'\).
\end{proof}

Here \(H^{1,1}(Z,\mathbb R)\) is modeled by the real Bott--Chern
quotient of smooth closed real \((1,1)\)-forms modulo global smooth
\(dd^c\) corrections, as in Definition~\ref{bc:quotient}. Its topology
is coinduced from all compact chart scalar coefficient derivatives,
as in Lemma~\ref{xiao:actual-pure-smooth-bottchern-quotient}; the
K\"ahler cone is the image of genuine K\"ahler forms. The class
intersection numbers are the descended wedge integrals of
Lemma~\ref{xiao:actual-class-top-and-mixed-intersection}, agreeing with
Lemma~\ref{pc:domination-intersection}.

\begin{theorem}[Bigness of a difference of nef classes (Xiao)]
\label{ana:nef-difference}
\uses{ana:geometric-setting,bc:quotient,bc:class-module,pc:domination-intersection,xiao:actual-class-top-and-mixed-intersection,xiao:actual-pure-smooth-closure-iff-nef,xiao:actual-pure-smooth-nef-difference-big,ana:nonkahler-locus}
\source{transcendental-nakamaye-copy4-20260911:thm:nef-difference,xiao15-weak-morse,dp04-kahler-cone:1.5,demailly-cadg}
\lean{TranscendentalNakamaye.Pending.XiaoPureSmoothBottChernClosure.isBigClass_sub_of_pureClosure_gap}
\leanok
Let \(Z\) be a connected compact K\"ahler manifold of dimension \(k>0\).
Let \(A,B\in H^{1,1}(Z,\mathbb R)\) be nef classes, meaning classes in
the closure of the K\"ahler cone. If
\(A^k-4kA^{k-1}\cdot B>0\), then \(A-B\) is big.
\end{theorem}
\begin{proof}
The pure smooth closure equivalence gives smooth \(\varepsilon\)-nef
approximations for each of the two classes. The simultaneous
Bott--Chern invariance of the class slot integrals identifies the
stated strict gap with the corresponding top and one-slot integrals
of any smooth representatives. The numerical difference criterion
of Theorem~\ref{xiao:actual-pure-smooth-nef-difference-big} then
produces a K\"ahler current in the exact difference class \(A-B\).
Consequently this class is big.
\end{proof}
This is the compact K\"ahler case of Xiao's Theorem~1.4 quoted by
Copy-4. It is used in Proposition~\ref{pc:domination-general},
Copy-4's proof of the domination lemma without bigness. The domination
lemma for big classes (Lemma~\ref{ana:domination}) and the main theorem
use their independently proved big-class argument.

\begin{lemma}[Uniform divisorial domination]
\label{ana:domination}
\uses{ana:geometric-setting,mmass:diagonal-background-monotone,mmass:monotonicity,xab:barrier,pc:positive-rooftop,trace:i-order}
\source{transcendental-nakamaye-copy4-20260911:lem:domination,xia-sgpt:page-0083}
\lean{TranscendentalNakamaye.Pending.PositiveCurrentAssembly.DominationProperty,TranscendentalNakamaye.Pending.PositiveCurrentAssembly.DominationLemma,TranscendentalNakamaye.Pending.PositiveCurrentAssembly.HasDominationLemma,TranscendentalNakamaye.Pending.DivisorialDomination.diagonalPotentialMass_congr,TranscendentalNakamaye.Pending.DivisorialDomination.exists_isThetaPshPotential_sub_smul_of_isBigClass,TranscendentalNakamaye.Pending.DivisorialDomination.exists_diagonalPotentialMass_le,TranscendentalNakamaye.Pending.DivisorialDomination.exists_one_lt_ofReal_pow_sub_one_mul_lt,TranscendentalNakamaye.Pending.DivisorialDomination.isAdmissibleWeight_of_le,TranscendentalNakamaye.Pending.DivisorialDomination.dominationLemma_of_isBigClass,TranscendentalNakamaye.Pending.DominationBigRestriction.hasDominationLemma_of_isBigClass}
\leanok
Let \((Z,\omega_Z)\) be a connected compact K\"ahler manifold of
dimension \(k>0\), let \(\gamma\) be a smooth closed real
\((1,1)\)-form with big class, and let \(h_0>0\). There is a
\(\delta_*>0\), depending only on \((Z,\omega_Z),\gamma,h_0\), such
that whenever \(0\le\eta\le1\) and
\[
 v\in\operatorname{PSH}(Z,\gamma+\eta\omega_Z),\qquad
 \int_Z\langle(\gamma+\eta\omega_Z+d d^c v)^k\rangle\ge h_0,
\]
there is a closed positive current
\(S_v\in\{\gamma+\eta\omega_Z\}\) satisfying
\[
 S_v\ge\delta_*\omega_Z,\qquad S_v\preceq_{\mathcal I}v.
\]
Here the second assertion means that a local potential \(s_v\) of \(S_v\)
satisfies
\(\mathcal I(\lambda s_v)\subset\mathcal I(\lambda v)\) for every
\(\lambda>0\).
\end{lemma}
\begin{proof}
\leanok
Since \(\{\gamma\}\) is big, it contains a K\"ahler current: there are \(\delta_0>0\) and a \((\gamma-\delta_0\omega_Z)\)-psh potential \(u_0\). The
masses in the classes \(\gamma+\eta\omega_Z\), \(\eta\le1\), are bounded by one constant \(M<\infty\): with \(\gamma+A\omega_Z\ge0\) and \(B=A+1\), every
\(\psi\in\operatorname{PSH}(Z,\gamma+\eta\omega_Z)\) satisfies \(\int\langle(\gamma+\eta\omega_Z)_\psi^k\rangle\le\int\langle(\gamma+B\omega_Z)_\psi^k\rangle\le
\int(\gamma+B\omega_Z)^k\) (Theorems~\ref{mmass:diagonal-background-monotone} and~\ref{mmass:monotonicity}, comparing with the potential \(0\)). Let
\(h_0>0\) and choose \(a>1\) with \((a^k-1)M<h_0\); put \(\delta_*=(1-a^{-1})\delta_0\), which depends only on \((Z,\omega_Z)\), \(\gamma\) and \(h_0\).

Let \(0\le\eta\le1\), \(\gamma_\eta=\gamma+\eta\omega_Z\) and \(v\in\operatorname{PSH}(Z,\gamma_\eta)\) with \(\int\langle(\gamma_\eta)_v^k\rangle\ge h_0\). The
potential \(u_0\) is also \(\gamma_\eta\)-psh, so \(\psi=\max(v,u_0)\in\operatorname{PSH}(Z,\gamma_\eta)\), \(v\le\psi\), and
\(h_0\le m_v\le m_\psi\le M\) for the masses \(m_v,m_\psi\) in \(\gamma_\eta\). Hence
\(a^k(m_\psi-m_v)=(m_\psi-m_v)+(a^k-1)(m_\psi-m_v)<(m_\psi-m_v)+h_0\le m_\psi\): \(a\) is an admissible weight for \(v\le\psi\)
(Definition~\ref{xab:barrier}). The class \(\{\gamma_\eta\}\) is big, and Xia's Lemma~2.4.4 for big classes (the barrier in the proof of
Lemma~\ref{pc:positive-rooftop}) gives \(w\in\operatorname{PSH}(Z,\gamma_\eta)\) with \(a^{-1}w+(1-a^{-1})\psi\le v\). Put
\(s=a^{-1}w+(1-a^{-1})u_0\). Then
\[
 \gamma_\eta+dd^cs=a^{-1}(\gamma_\eta+dd^cw)+(1-a^{-1})(\gamma-\delta_0\omega_Z+dd^cu_0)+(1-a^{-1})(\delta_0+\eta)\omega_Z\ \ge\ \delta_*\omega_Z,
\]
so \(s\) is a \((\gamma_\eta-\delta_*\omega_Z)\)-psh potential, and \(s\le a^{-1}w+(1-a^{-1})\psi\le v\) everywhere because \(u_0\le\psi\). A pointwise bound
gives \(\mathcal I(\lambda s)\subseteq\mathcal I(\lambda v)\) for every \(\lambda>0\), so \(S_v=\gamma_\eta+dd^cs\) has both properties.
\end{proof}
The attached declarations state the lemma (\(S_v=\gamma+\eta\omega_Z+dd^cs\) with \(s\) a potential of \(\gamma+\eta\omega_Z-\delta_*\omega_Z\) and
\(s\preceq_{\mathcal I}v\)) and prove it under the bigness of \(\{\gamma\}\), Copy-4's hypothesis. Copy-4's own proof (approximation, log
resolution, Xiao's criterion, descent and a limit) does not use the bigness and is recorded in Proposition~\ref{pc:domination-general}; for a
big class the argument above, which is Xia's Lemma~2.4.4 applied once, needs none of it, and the comparison \(S_v\preceq_{\mathcal I}v\) holds
pointwise.

\begin{lemma}[Quantitative regularization with a one-sided ideal bound]
\label{ana:regularization}
\uses{ana:ideal-sheaves,ana:geometric-setting}
\source{transcendental-nakamaye-copy4-20260911:lem:regularization}
\lean{TranscendentalNakamaye.Pending.QuantitativeRegularization.exists_regularization_multiplierIdeal_lowerBound,TranscendentalNakamaye.Pending.QuantitativeRegularization.demailly1992_exists_regularization_multiplierIdeal_lowerBound,TranscendentalNakamaye.Pending.QuantitativeRegularizationFixedCover.exists_regularization_multiplierIdeal_lowerBound_of_pos,TranscendentalNakamaye.Pending.QuantitativeRegularizationFixedCover.exists_regularization_multiplierIdeal_lowerBound_dim_zero,TranscendentalNakamaye.Pending.QuantitativeRegularizationFixedCover.exists_regularization_multiplierIdeal_lowerBound}
\leanok
For every connected compact K\"ahler manifold \((Z,\omega_Z)\) there
is a constant \(c_Z>0\), independent of the smooth closed real
\((1,1)\)-form \(\zeta\), with the following property. If
\(\varphi\in\operatorname{PSH}(Z,\zeta)\), then for every integer \(m\ge1\)
there exists a quasi-plurisubharmonic function \(\varphi_m\) such that
\[
 \zeta+d d^c\varphi_m\ge-\frac{c_Z}{m}\omega_Z,
 \qquad
 \varphi_m\ge\frac1m\psi_{\mathcal I(m\varphi)}+O(1)
 \quad\text{locally}.
\]
The bounded term in the second inequality may depend on the coordinate
neighborhood and on \(m\); the constant \(c_Z\) is independent of
\(\zeta,\varphi,m\). The second inequality is meant for every finite family of germs: near
every point \(x\), \(\varphi_m\ge m^{-1}\log\sum_i|g_i|^2-C\) for all finitely many germs
\(g_i\) of \(\mathcal I(m\varphi)\) at \(x\) (for local generators this is the displayed form,
and every finite family of germs is dominated by local generators).
\end{lemma}
\begin{proof}
\uses{reg:fixed-cover-branches,ana:regularization-gluing,gnpp:kahler-domination}
\leanok
If \(\dim Z=0\), the points of \(Z\) are open and every \((1,1)\)-form vanishes. Take
\(\varphi_m=\varphi\) and any \(c_Z>0\): \(\varphi\) is a \(\zeta\)-plurisubharmonic potential,
hence finite at every point (a point is an open set on which \(\varphi\) is not identically
\(-\infty\)), and near \(x\) (that is, at \(x\)) the finitely many values \(g_i(x)\) give
\(m^{-1}\log\sum_i|g_i(x)|^2-C\le\varphi(x)\) for \(C\) large.

If \(\dim Z\ge1\), Lemma~\ref{reg:fixed-cover-branches} gives a finite cut-off cover
\((U_j,V_j,\chi_j)\) and a constant \(C_B\ge0\), fixed before \(\zeta\), together with branch
data (i)--(iv) of Lemma~\ref{ana:regularization-gluing} for every \(\zeta\), \(\varphi\) and
\(m\). Lemma~\ref{gnpp:kahler-domination} gives \(C_\chi>0\) with
\(dd^c\chi_j\ge-C_\chi\omega_Z\) for all \(j\), and Lemma~\ref{ana:regularization-gluing}
gives the conclusion with \(c_Z=(C_B+1)C_\chi>0\).

This is Copy-4's proof: Step~1 is the fixed cover and the chart Bergman branches
(Lemmas~\ref{reg:chart-branch-properties} and~\ref{reg:fixed-cover-branches}), Step~2 the
overlap estimate (Lemma~\ref{reg:chart-overlap}), Steps~3--4 the gluing. It differs from
Copy-4 in three places: the overlap estimate needs no Lebesgue-number condition on the cover;
the gluing uses the plain maximum (Lemma~\ref{ana:regularization-gluing}; the regularized
maximum form is Corollary~\ref{ana:regularization-rmax}); and the lower bound is obtained for
every finite family of germs directly (Lemma~\ref{berg:germ-domination}), so neither the ideal
identity \(\mathfrak J_{j,m}=\mathcal I(m\varphi)\) of Copy-4's Step~4 nor coherence of \(\mathcal I(m\varphi)\) is used.
Connectedness of \(Z\) is not used.
\end{proof}

\begin{corollary}[The regularization dominates every multiplier-ideal weight globally]
\label{ana:regularization-global-lower-bound}
\uses{ana:regularization,ana:ideal-sheaves}
\source{transcendental-nakamaye-copy4-20260911:lem:regularization}
\lean{TranscendentalNakamaye.Pending.AmbientUMWeightLowerBound.moreSingularAt_smul_of_generatorLog_lowerBound,TranscendentalNakamaye.Pending.AmbientUMWeightLowerBound.moreSingularOn_smul_of_generatorLog_lowerBound,TranscendentalNakamaye.Pending.AmbientUMWeightLowerBound.exists_regularization_moreSingularOn_multiplierIdealWeight}
\leanok
In Lemma~\ref{ana:regularization}, if \(\psi\) has analytic singularities of type
\((1,\mathcal I(m\varphi))\), then \(\varphi_m\ge m^{-1}\psi-C\) on all of \(Z\) for one constant
\(C\).
\end{corollary}
\begin{proof}
\leanok
Near every point \(x\), \(\psi=\log\sum_i|g_i|^2+R\) with \(g_i\) generators of
\(\mathcal I(m\varphi)\) near \(x\) and \(R\) bounded; the lower bound of
Lemma~\ref{ana:regularization} for the finite family \((g_i)\) at \(x\) gives
\(m^{-1}\psi\le\varphi_m+C_x\) near \(x\). Finitely many such neighbourhoods cover the compact
\(Z\); take \(C\) the largest \(C_x\).
\end{proof}

\begin{lemma}[Gluing local branches on a fixed cover]
\label{ana:regularization-gluing}
\uses{ana:ideal-sheaves,psh:maximum,gnpp:kahler-domination}
\source{transcendental-nakamaye-copy4-20260911:lem:regularization,transcendental-nakamaye-copy4-20260911:eq:reg-glue}
\lean{TranscendentalNakamaye.Pending.QuantitativeRegularizationMaxGluing.CutoffCover,TranscendentalNakamaye.Pending.QuantitativeRegularizationMaxGluing.CutoffCover.glue,TranscendentalNakamaye.Pending.QuantitativeRegularizationMaxGluing.CutoffCover.branch_le_glue,TranscendentalNakamaye.Pending.QuantitativeRegularizationMaxGluing.CutoffCover.le_glue,TranscendentalNakamaye.Pending.QuantitativeRegularizationMaxGluing.CutoffCover.exists_nhds_glue_eq_filter,TranscendentalNakamaye.Pending.QuantitativeRegularizationMaxGluing.CutoffCover.thetaPlurisubharmonic_glue,TranscendentalNakamaye.Pending.QuantitativeRegularizationMaxGluing.CutoffCover.isThetaPshPotential_glue,TranscendentalNakamaye.Pending.QuantitativeRegularizationMaxGluing.exists_regularization_of_branches}
\leanok
Let \((Z,\omega_Z)\) be a compact K\"ahler manifold. Let \((V_j)_{j\in J}\) be a finite family
of open sets, \((U_j)_{j\in J}\) open sets covering \(Z\), and \(\chi_j\) smooth functions with
\(\chi_j\ge0\), \(\chi_j=1\) on \(U_j\) and \(\operatorname{supp}\chi_j\subset V_j\); let
\(C_\chi>0\) satisfy \(dd^c\chi_j\ge-C_\chi\omega_Z\) for all \(j\), and let \(C_B\ge0\). Suppose
that for every smooth closed real \((1,1)\)-form \(\zeta\), every
\(\varphi\in\operatorname{PSH}(Z,\zeta)\) and every integer \(m\ge1\) there are functions
\(b_j\), \(j\in J\), such that
\begin{enumerate}
\item \(b_j\) is \(\zeta\)-plurisubharmonic on \(V_j\);
\item \(b_j\le b_l+C_B/m\) on \(V_j\cap V_l\);
\item \(b_j\) is identically \(-\infty\) on no neighbourhood of a point of \(U_j\);
\item for every \(x\in U_j\) and all finitely many germs \(g_1,\dots,g_N\) of
\(\mathcal I(m\varphi)\) at \(x\), \(b_j\ge m^{-1}\log\sum_i|g_i|^2-C\) near \(x\) for a
constant \(C\).
\end{enumerate}
Then the conclusion of Lemma~\ref{ana:regularization} holds with \(c_Z=(C_B+1)C_\chi\): for
every \(\zeta\), \(\varphi\in\operatorname{PSH}(Z,\zeta)\) and \(m\ge1\) there is
\(\varphi_m\in\operatorname{PSH}(Z,\zeta+\tfrac{c_Z}m\omega_Z)\) such that, near every point
\(x\), \(\varphi_m\ge m^{-1}\log\sum_i|g_i|^2-C\) for all finitely many germs \(g_i\) of
\(\mathcal I(m\varphi)\) at \(x\).
\end{lemma}
\begin{proof}
\leanok
Put \(A=(C_B+1)/m\) and
\(\varphi_m(z)=\max\{b_j(z)+A\chi_j(z):j\in J,\ z\in V_j\}\), a maximum over a nonempty set
because the \(U_j\subset V_j\) cover \(Z\).

\emph{Near every point, \(\varphi_m\) is the maximum of finitely many branches defined on a
whole neighbourhood.} Fix \(q\in Z\), let \(J_q=\{j:q\in V_j\}\), and choose \(l\) with
\(q\in U_l\), so \(l\in J_q\). If \(q\notin\overline{V_j}\), a neighbourhood of \(q\) misses
\(V_j\). If \(q\in\overline{V_j}\setminus V_j\), then \(q\notin\operatorname{supp}\chi_j\), so
\(\chi_j=0\) near \(q\), while \(\chi_l=1\) near \(q\); at points \(z\in V_j\cap V_l\) near
\(q\), (ii) and \(A\ge C_B/m\) give
\(b_j(z)+A\chi_j(z)=b_j(z)\le b_l(z)+C_B/m\le b_l(z)+A\chi_l(z)\), so the \(j\)-th branch never
realizes the maximum there. Hence on a neighbourhood \(N\subset\bigcap_{j\in J_q}V_j\) of \(q\),
\(\varphi_m=\max_{j\in J_q}(b_j+A\chi_j)\).

\emph{Positivity.} On \(V_j\), \(b_j+A\chi_j\) is
\((\zeta+AC_\chi\omega_Z)\)-plurisubharmonic, since \(dd^c(A\chi_j)\ge-AC_\chi\omega_Z\): with a
smooth local potential \(\rho\) of \(\zeta\) and one, \(\sigma\), of \(C_\chi\omega_Z\), the sum
\((b_j+\rho)+A(\chi_j+\sigma)\) is plurisubharmonic. A maximum of finitely many plurisubharmonic
functions is plurisubharmonic (Proposition~\ref{psh:maximum}), so \(\varphi_m\) is
\((\zeta+\tfrac{c_Z}m\omega_Z)\)-plurisubharmonic near every point, hence on \(Z\).

\emph{Nondegeneracy and the lower bound.} On \(V_j\), \(\varphi_m\ge b_j+A\chi_j\ge b_j\). If
\(\varphi_m\) were \(-\infty\) on a nonempty open set \(O\), take \(x\in O\) and \(l\) with
\(x\in U_l\); then \(b_l=-\infty\) on the neighbourhood \(O\cap V_l\) of \(x\), against (iii).
For the lower bound at \(x\), take \(j\) with \(x\in U_j\) and use \(\varphi_m\ge b_j\) near
\(x\) and (iv).

A constant \(C_\chi\) as required exists for every finite family of cut-offs
(Lemma~\ref{gnpp:kahler-domination}), and depends only on the cover and \(\omega_Z\).
\end{proof}
These are Steps~3 and~4 of Copy-4's proof of Lemma~\ref{ana:regularization}, with one
difference: Copy-4 glues with the regularized maximum \(M_\delta\), whose output is smooth off
the zero locus of \(\mathcal I(m\varphi)\), and argues off that locus, so it needs the extension
theorem for plurisubharmonic functions across it and the ideal identity
\(\mathfrak J_{j,m}=\mathcal I(m\varphi)\) to control where the branches are \(-\infty\). The
statement of Lemma~\ref{ana:regularization} asks for no smoothness; with branches that are
\(\zeta\)-plurisubharmonic across their poles, neither ingredient is used, whether the gluing is
the plain maximum here or the regularized maximum of Lemma~\ref{ana:regularization-gluing-rmax}
(whose boundary argument is carried out in \([-\infty,\infty)\)). The lower bound is obtained for
every finite family of germs of \(\mathcal I(m\varphi)\), so coherence of \(\mathcal I(m\varphi)\)
is not used either. The branch data (i)--(iv) on a cover fixed in advance, with \(C_B\)
independent of \(\zeta\), \(\varphi\) and \(m\), are Lemma~\ref{reg:fixed-cover-branches}.

\begin{lemma}[Relative principalization]
\label{ana:relative-principalization}
\uses{ana:blowups,ana:geometric-setting,alg:restriction,alg:controlled-factor}
\source{transcendental-nakamaye-copy4-20260911:lem:relative-principalization}
\lean{TranscendentalNakamaye.Pending.RelativePrincipalization.exists_relativePrincipalization,TranscendentalNakamaye.Pending.RelativePrincipalization.exists_factorizationData,TranscendentalNakamaye.Pending.PrincipalizationStepOneInterface.RelativePrincipalizationHolds,TranscendentalNakamaye.Pending.PrincipalizationStepOneInterface.relativePrincipalizationHolds}
\leanok
Let \(Y\) be a connected compact complex manifold and let \(W\subsetneq Y\)
be a nonempty closed connected complex submanifold. Let
\(\mathfrak a\subset\mathcal O_Y\) be a coherent ideal with
\(\operatorname{Res}_W\mathfrak a\ne0\). There is a finite
composition of blowups \(f:Y'\to Y\) with smooth centers, each center
either contained in a proper analytic subset of the current strict
transform of \(W\), or disjoint from that strict transform. The manifold
\(Y'\) is connected and compact, the strict
transform \(W'\) is smooth and connected; \(f\) is an isomorphism over
a nonempty open subset of \(W\); and there are effective divisors
\(H,H'\) such that \(H\) is exceptional, \(H'\) is disjoint from
\(W'\), no component of \(H\) contains \(W'\), and
\[
 \begin{aligned}
 \mathcal I_W\mathcal O_{Y'}&=\mathcal O_{Y'}(-H)\mathcal I_{W'},\\
 \mathfrak a\mathcal O_{Y'}&=\mathcal O_{Y'}(-H-H'),\\
 (\operatorname{Res}_W\mathfrak a)\mathcal O_{W'}
 &=\mathcal O_{W'}(-H|_{W'}).
 \end{aligned}
\]
\end{lemma}
\begin{proof}
\leanok
\uses{ana:relative-principalization-from-core,mc:hironaka-from-seq-invariant,hir:canonical-centres}
The geometric core is Hironaka's principalization: a sequence of blowups \(f\) adapted to
\(W\), with connected strict transform \(W'\) regular of dimension \(k\) at each of its
points, and effective divisors \(H\) (exceptional) and \(H'\) (disjoint from \(W'\)) with
\(\mathcal I_W\mathcal O_{Y'}=\mathcal O_{Y'}(-H)\mathcal I_{W'}\) and
\(\mathfrak a\mathcal O_{Y'}=\mathcal O_{Y'}(-H-H')\). By Theorem~\ref{hir:canonical-centres} there is a resolution invariant in
dimension \(\le n\) on complex manifolds of dimension \(n=\dim Y\), and Theorem~\ref{mc:hironaka-from-seq-invariant}(2), with \(j=n\),
gives exactly these data. Lemma~\ref{ana:relative-principalization-from-core} derives everything else from them.
\end{proof}
The geometric core of this lemma is one of the two places where resolution of singularities enters the main theorem's
dependency cone, the other being Lemma~\ref{ana:embedded-resolution}; both are proved from the canonical resolution invariant
(Theorem~\ref{hir:canonical-centres}) through Theorem~\ref{mc:hironaka-from-seq-invariant}.

\begin{lemma}[Relative principalization from its geometric core]
\label{ana:relative-principalization-from-core}
\uses{ana:blowups}
\source{transcendental-nakamaye-copy4-20260911:lem:relative-principalization}
\lean{TranscendentalNakamaye.Pending.RelativePrincipalization.exists_relativePrincipalization_of_factorizationData}
\leanok
Let \(Y\), \(W\), \(k\) and the coherent ideal \(\mathfrak a\) be as in
Lemma~\ref{ana:relative-principalization}. Suppose given a sequence of blowups
\(f:Y'\to Y\) adapted to \(W\) whose strict transform \(W'\) is connected and regular of
dimension \(k\) at each of its points, and effective divisors \(H\) with support in
\(\operatorname{Exc}f\) and \(H'\) disjoint from \(W'\), such that
\(\mathcal I_W\mathcal O_{Y'}=\mathcal O_{Y'}(-H)\mathcal I_{W'}\) and
\(\mathfrak a\mathcal O_{Y'}=\mathcal O_{Y'}(-H-H')\). Then \(f\), \(H\), \(H'\) have all the
properties asserted in Lemma~\ref{ana:relative-principalization}.
\end{lemma}
\begin{proof}
\leanok
\(Y'\) is compact and connected, being
obtained from \(Y\) by blowups along smooth centres. \(W'\) is closed, as the closure of
\(f^{-1}(W)\setminus\operatorname{Exc}f\), so it is a closed connected submanifold of
dimension \(k\). Since \(W'\) is nonempty and the centres are adapted to \(W\), some open set
meeting \(W\) lies over the locus where \(f\) is an isomorphism; a point of \(W'\) above it is
off \(\operatorname{Exc}f\supseteq\operatorname{Supp}H\), so no component of \(H\) contains
\(W'\). Finally, at a point \(y\in W'\) the divisor \(H'\) is a unit near \(y\), so the second
factorization makes the restrictions to \(W'\) of \(\mathfrak a\mathcal O_{Y'}\) and of
\(\mathcal O_{Y'}(-H)\) span the same germs at \(y\), which is the third identity.
\end{proof}

\begin{theorem}[Positive mass for affine rooftop obstacles]
\label{ana:rooftop-positive-mass}
\uses{trace:p-order,trace:diagonal-mass,fprooftop:definition}
\source{transcendental-nakamaye-copy4-20260911:thm:rooftop}
\lean{TranscendentalNakamaye.Pending.POrderAffineRooftop.isThetaPshPotential_rooftopEnvelope_of_pOrderWith,TranscendentalNakamaye.Pending.POrderAffineRooftop.isAdmissibleWeight_of_lt_mul}
\leanok
Let \(\varphi,\psi\in\operatorname{PSH}(X,\theta)\) have positive
non-pluripolar mass and suppose \(\varphi\preceq_{P,\theta}\psi\).
If \(t>1\) and
\[
 t^n\int_X\langle\theta_\varphi^n\rangle
 >(t^n-1)\int_X\langle\theta_\psi^n\rangle,
\]
then \(P_\theta(t\varphi-(t-1)\psi)\) is a genuine theta-PSH
potential of positive mass.
\end{theorem}
\begin{proof}
\uses{trace:p-order-affine-rooftop,xab:p-mass-preserving,fpenv:p-psh,ana:mass-monotonicity,contact:diagonal-measure-regular}
\leanok
This is Xia's Theorem~6.1.1 (Theorem~\ref{trace:p-order-affine-rooftop}), which Copy-4 cites, with \(a=t\) and
\(P_\theta(t\varphi-(t-1)\psi)=P_\theta(t\varphi+(1-t)\psi)\). Write \(m_\varphi,m_\psi\) for the two masses. By
Proposition~\ref{xab:p-mass-preserving} they are the masses of \(P_\theta[\varphi]\le P_\theta[\psi]\), both theta-PSH
(Theorem~\ref{fpenv:p-psh}), so \(m_\varphi\le m_\psi\) by Theorem~\ref{ana:mass-monotonicity}, and \(m_\psi<\infty\)
(Lemma~\ref{contact:diagonal-measure-regular}). With \(m_\varphi\le m_\psi<\infty\) and \(m_\varphi>0\), the hypothesis
\(t^nm_\varphi>(t^n-1)m_\psi\) is \(t^n(m_\psi-m_\varphi)<m_\psi\), so \(t\) is an admissible weight, and
Theorem~\ref{trace:p-order-affine-rooftop} gives a theta-PSH envelope of positive mass. The positivity of \(m_\psi\) is not used.
\end{proof}

\begin{lemma}[Quasi-everywhere rooftop obstacle bound]
\label{ana:rooftop-obstacle-bound}
\uses{fprooftop:definition,mpp:qe,psh:ae-comparison,cap:qe-reduction,cap:countable-stability,cap:negligible-pluripolar}
\source{transcendental-nakamaye-copy4-20260911:sec:npp,xia-sgpt:page-0078}
\lean{TranscendentalNakamaye.Pending.PolydiskPatchBalayageContactQE.rooftopEnvelope_le_quasiEverywhere_of_kahler,TranscendentalNakamaye.Pending.RooftopObstacleQuasiEverywhere.rooftopEnvelope_le_quasiEverywhere,TranscendentalNakamaye.Pending.PluripolarOuterCapacity.countablyStablePluripolar,TranscendentalNakamaye.Pending.NegligibleSetsPluripolarProof.negligibleSetsPluripolar,TranscendentalNakamaye.FirstPass.le_rooftopEnvelope}
\leanok
If \(P_\theta(f)\) is a genuine theta-PSH potential, then
\(P_\theta(f)\le f\) quasi-everywhere. Every theta-PSH competitor
\(v\le f\) quasi-everywhere satisfies \(v\le P_\theta(f)\) everywhere.
\end{lemma}
\begin{proof}
The second statement holds by construction: such a \(v\) is a member of \(\mathcal R_\theta(f)\), so
\(v\le R_\theta(f)\le P_\theta(f)\). For the first, \(X\) is compact, connected and K\"ahler, so countable
stability holds (Theorem~\ref{cap:countable-stability}) and negligible sets are pluripolar for \(\theta\)
(Theorem~\ref{cap:negligible-pluripolar}); Theorem~\ref{cap:qe-reduction} concludes.
\end{proof}
The first attached declaration is the first statement; the next three are its ingredients (the conditional
reduction and the two pluripolarity statements it consumes), and the last is the second statement.

\begin{proposition}[Analytic approximations of the trace]
\label{ana:trace-approximants}
\uses{ana:trace-formula,ana:trace-multiplier,ana:mass-calculus,alg:restriction,trace:qe-existence,trace:ds-full-mixed-mass-continuity,ana:actual-restricted-trace-approximants,trace:actual-connected-submanifold-irreducible-carrier,ana:actual-restricted-trace-ideal-relations,ana:actual-restricted-trace-analytic-approximants}
\source{transcendental-nakamaye-copy4-20260911:lem:trace-approximants}
\lean{TranscendentalNakamaye.Pending.ActualRestrictedTraceApproximants.exists_restricted_trace_approximants}
\leanok
Let \(\alpha\) be big, let \(\theta\in\alpha\), and let
\(V\subset X\) be a closed connected complex submanifold of dimension
\(k>0\), not contained in \(E_{\mathrm{nn}}(\alpha)\), with
\(h=\operatorname{RV}_{X\mid V}(\alpha)>0\). Put
\(U=V_\theta\), \(\gamma=\theta|_V\), \(\omega_V=\omega|_V\), and
\[
 u=\operatorname{Tr}^{\theta}_V(U)\in\operatorname{PSH}(V,\gamma),
 \qquad \int_V\langle\gamma_u^k\rangle=h.
\]
Then \(\{\gamma\}\) is big. Set
\(\mathfrak B_q=\mathcal I(qU)\) and
\(\overline{\mathfrak B}_q=\operatorname{Res}_V\mathfrak B_q\).
These ideals are units on nonempty open subsets of \(V\), and
\(\mathcal I(qu)\subset\overline{\mathfrak B}_q\). For positive integers
\(m\mid M\), they satisfy
\[
 \mathfrak B_M\subset\mathfrak B_m^{M/m},\qquad
 \overline{\mathfrak B}_M\subset\overline{\mathfrak B}_m^{M/m}.
\]
There are numbers \(\eta_q\downarrow0\), functions
\(\phi_q\in\operatorname{PSH}(V,\gamma+\eta_q\omega_V)\) of type
\((q^{-1},\overline{\mathfrak B}_q)\), obtained by restricting a
decreasing dyadic quasi-equisingular approximation of \(U\), such that
\[
 \begin{gathered}
 \int_V\langle(\gamma+\eta_q\omega_V+d d^c\phi_q)^k\rangle
 \longrightarrow h,\\
 u\preceq_{P}\phi_q\quad(q\text{ sufficiently large}).
 \end{gathered}
\]
For every local section \(g\) of \(\overline{\mathfrak B}_q^{\,i}\),
\(i\ge1\), the function \(|g|^2e^{-qi\phi_q}\) is locally bounded
outside a pluripolar set.
\end{proposition}
\begin{proof}
Theorem~\ref{trace:actual-connected-submanifold-irreducible-carrier}
supplies the irreducible analytic carrier of the same submanifold $V$
with the same dimension $k$. By the trace formula, $h>0$ excludes the
undefined, zero-mass branch of the normalized trace. Hence the stated
$u$ is genuinely $\gamma$-PSH with mass $h$; positive non-pluripolar mass
on $V$ implies that $\{\gamma\}$ is big.

The ideal assertions concern every integer $q\ge1$.
Theorem~\ref{ana:actual-restricted-trace-ideal-relations} applies to this
same $U$ and normalized trace $u$. Indeed, a point of
$V\setminus E_{\mathrm{nn}}(\alpha)$ has $\nu(U,x)=0$; Skoda
integrability and coherence give a nonempty unit neighborhood for
each $\mathfrak B_q$, and restriction preserves the unit ideal.
The trace multiplier inclusion gives
$\mathcal I(qu)\subset\overline{\mathfrak B}_q$.
Repeated multiplier-ideal subadditivity gives
$\mathfrak B_M\subset\mathfrak B_m^{M/m}$ when $m\mid M$;
restriction of products gives the corresponding restricted inclusion.

For the approximation family write $q_j=2^j$.
Theorem~\ref{ana:actual-restricted-trace-analytic-approximants} constructs
a decreasing quasi-equisingular family $U_j$ of $U$, with positive
antitone errors $\eta_j\to0$ and analytic type
$(q_j^{-1},\mathfrak B_{q_j})$.
Its genuine restrictions
$\phi_{q_j}=U_j|_V$ lie in
$\operatorname{PSH}(V,\gamma+\eta_j\omega_V)$ and have type
$(q_j^{-1},\overline{\mathfrak B}_{q_j})$.
The shrinking-background mass theorem identifies their mass limit
with the mass of
$\inf_jP_{\gamma+\eta_j\omega_V}[\phi_{q_j}]$.
Normalized-reader compatibility identifies this function with the
same $u$, so the limit is $h$.
In a common admissible K\"ahler background, full trace-reader
convergence to this $u$, antitonicity and global order closure give
$u\preceq_P\phi_{q_j}$, in particular for sufficiently large $j$.

Finally, choose local generators $f_a$ of
$\overline{\mathfrak B}_{q_j}$. Its analytic type gives
\[
 \phi_{q_j}=q_j^{-1}\log\sum_a|f_a|^2+b_j
\]
with $b_j$ locally bounded. If
$g\in\overline{\mathfrak B}_{q_j}^{\,i}$, a finite expression in
products of $i$ generators, with holomorphic coefficients, gives
$|g|^2\le C(\sum_a|f_a|^2)^i$ after shrinking the neighborhood.
Outside the pluripolar pole set this yields
\[
 |g|^2e^{-q_ji\phi_{q_j}}
 \le C\exp(q_ji\|b_j\|_\infty),
\]
which is the required bound with the literal exponent $q_ji$.
\end{proof}

\begin{lemma}[Positive rooftop built from a trace approximant]
\label{ana:positive-rooftop}
\lean{TranscendentalNakamaye.Pending.ActualNormalizedTracePositiveRooftop.eventually_positive_rooftop_of_normalizedTraceAlong}
\uses{ana:trace-approximants,ana:rooftop-positive-mass,ana:rooftop-obstacle-bound,ana:mass-calculus,num:rooftop-gap,ana:actual-normalized-trace-positive-rooftop,trace:actual-positive-normalized-trace-domain-value,trace:qe-existence-global-type,pc:kahler-restriction}
\source{transcendental-nakamaye-copy4-20260911:lem:positive-rooftop}
\leanok
For every fixed integer \(C>1\), all sufficiently large dyadic \(m\)
satisfy
\[
 \begin{gathered}
 w:=P_{\gamma+\eta_m\omega_V}
       (2Cu-(2C-1)\phi_m)
 \in\operatorname{PSH}(V,\gamma+\eta_m\omega_V)_{>0},\\
 h_C:=\tfrac12(w+\phi_m),\qquad
 h_C\le Cu-(C-1)\phi_m\quad\text{quasi-everywhere},\\
 \int_V\langle(\gamma+\eta_m\omega_V+d d^c h_C)^k\rangle
 \ge2^{-k-1}h.
 \end{gathered}
\]
\end{lemma}

\begin{lemma}[Two-level multiplier-ideal comparison]
\label{ana:ideal-comparison}
\lean{TranscendentalNakamaye.FirstPass.Copy4IntegerLevels.integral_identities,TranscendentalNakamaye.FirstPass.Copy4IntegerLevels.m_dvd_M,TranscendentalNakamaye.FirstPass.Copy4IntegerLevels.p_eq_mul_ell_div_m,TranscendentalNakamaye.Pending.ActualTraceIdealComparison.twoLevel_multiplierIdeal_inclusions,TranscendentalNakamaye.Pending.ActualTracePrincipalizedIdealRemainders.exists_principalized_multiplierIdeal_remainders,TranscendentalNakamaye.FirstPass.Copy4IntegerLevels.divisor_coefficient_decomposition}
\uses{ana:trace-approximants,ana:positive-rooftop,alg:restriction,num:copy4-level-identities,num:divisor-decomposition,pc:ideal-comparison-first,ana:actual-trace-two-level-ideal-inclusions,ana:actual-principalized-multiplier-ideal-remainders,ana:ideal-sheaves,dem:del00-subadditivity,ana:actual-trace-free-normalized-reader,ana:trace-multiplier}
\source{transcendental-nakamaye-copy4-20260911:lem:ideal-comparison}
\leanok
In the notation of Proposition~\ref{ana:trace-approximants}, fix
\(m,C\) and let \(K=\gamma+d d^c\varphi_K\) satisfy
\(\mathcal I(\lambda\varphi_K)\subset\mathcal I(\lambda h_C)\)
for every \(\lambda>0\). Choose a dyadic integer \(\ell\ge m\), and put
\[
 M=C\ell,\qquad p=\frac{(C-1)\ell}{m},\qquad
 \mathfrak B=\mathcal I(\ell\varphi_K).
\]
Then \(M/m\) and \(p\) are integers, \(mp=\ell(C-1)\), and
\(M=C\ell\).
Moreover,
\[
 \mathfrak B\,\overline{\mathfrak B}_m^{\,p}
 \subset\mathcal I(Mu)\subset\overline{\mathfrak B}_M.
\]
On any modification of \(V\) principalizing these ideals, let their
effective divisors be \(G_B,G_m,G_M\). Then
\[
 G_M-\frac{M}{m}G_m\ge0,\qquad
 F_B:=G_B+pG_m-G_M\ge0,
\]
and
\[
 \frac{1}{\ell}G_B=
 \frac1M G_M+\frac{C-1}{M}\left(G_M-\frac{M}{m}G_m\right)
 +\frac{C}{M}F_B.
\]
\end{lemma}
\begin{proof}
\leanok
\uses{ana:trace-approximants,ana:positive-rooftop,pc:ideal-comparison-first,ana:actual-trace-two-level-ideal-inclusions,ana:trace-multiplier,ana:actual-trace-free-normalized-reader,dem:del00-subadditivity,ana:actual-holomorphic-ideal-product-pullback,ana:actual-principalized-multiplier-ideal-remainders,ana:actual-effective-divisor-difference,ana:ideal-sheaves,num:copy4-level-identities,num:divisor-decomposition}
Let \(f\) be a local section of \(\mathfrak B\) and let \(g\) be a
local section of \(\overline{\mathfrak B}_m^{\,p}\). The first
assumption on \(K\) implies that
\(|f|^2e^{-\ell h_C}\) is locally integrable. Proposition~
\ref{ana:trace-approximants} and \(mp=\ell(C-1)\) imply that
\(|g|^2e^{-\ell(C-1)\phi_m}\) is locally bounded outside a pluripolar
set, hence almost everywhere. Lemma~
\ref{ana:positive-rooftop} gives
\(h_C\le Cu-(C-1)\phi_m\) outside a pluripolar set, hence almost
everywhere. Since \(M=C\ell\), it gives
\[
 |fg|^2e^{-Mu}
 \le |f|^2e^{-\ell h_C}\,|g|^2e^{-\ell(C-1)\phi_m}
 \quad\text{almost everywhere}.
\]
The right-hand side is locally integrable. Thus \(fg\in\mathcal I(Mu)\),
and finite sums of these products give the first ideal inclusion.
The second inclusion is the trace multiplier-ideal inclusion at level
\(M\).

Pullback preserves inclusion and products of ideals. On the chosen
principalization, \(\overline{\mathfrak B}_M\subset
\overline{\mathfrak B}_m^{M/m}\) therefore says
\(G_M\ge(M/m)G_m\). The product inclusion just proved similarly
says \(G_B+pG_m\ge G_M\). Substituting
\(F_B=G_B+pG_m-G_M\) in the displayed decomposition, its right-hand
side becomes
\[
 \frac{C}{M}G_B+\left(\frac{Cp}{M}-\frac{C-1}{m}\right)G_m.
\]
The coefficient in parentheses is zero, and \(C/M=1/\ell\), which
proves the identity.
\end{proof}

\begin{proposition}[Principalization data for gluing]
\label{ana:principalization-data}
\uses{ana:blowups,ana:gluing-conditions,ana:nonkahler-locus,ana:locus-properties,ana:domination,ana:regularization,trace:qe-existence,ana:relative-principalization,ana:residual-nef,ana:exceptional-metric,ana:geometric-setting,num:divisor-decomposition,num:positivity}
\source{transcendental-nakamaye-copy4-20260911:sec:fixed-data,transcendental-nakamaye-copy4-20260911:lem:ideal-function,transcendental-nakamaye-copy4-20260911:prop:principalization}
\lean{TranscendentalNakamaye.Pending.PrincipalizationData.exists_principalizationGluingData,TranscendentalNakamaye.Pending.PrincipalizationStepOneInterface.exists_adaptedBlowup_polarBranch_stepFiveData,TranscendentalNakamaye.Pending.PrincipalizationStepOneInterface.exists_adaptedBlowup_principalizationGluingConclusions}
\leanok
Let \(\alpha\) be big with smooth representative \(\theta\), and let
\(V\subset X\) be a closed connected complex submanifold of dimension
\(0<k<n\) which is an irreducible component of
\(E_{\mathrm{nK}}(\alpha)\), is not contained in
\(E_{\mathrm{nn}}(\alpha)\), and satisfies
\(h=\operatorname{RV}_{X\mid V}(\alpha)>0\). There exist a sequence of
blowups \(\pi:Y\to X\) adapted to \(V\), with \(Y\) a connected compact
K\"ahler manifold, its nonempty closed connected smooth strict transform \(W\)
of dimension \(k\), a point \(y\in W\), an open neighborhood
\(U\) of \(W\), a quasi-plurisubharmonic function \(\varphi_G\) on
\(Y\), and a quasi-plurisubharmonic function \(\varphi_U\) on \(U\).
Writing \(\theta_Y=\pi^*\theta\), these data satisfy:
\begin{enumerate}
\item \(\theta_Y+d d^c\varphi_G\) is a K\"ahler current, and
      \(\varphi_G=-\infty\) at every point of \(W\);
\item \(\theta_Y+d d^c\varphi_U\) dominates a K\"ahler form on \(U\),
      and \(\varphi_U\) is smooth near \(y\);
\item for every compact \(\Lambda\subset U\setminus W\), there is
      \(C_\Lambda\in\mathbb R\) such that
      \(\varphi_U\le\varphi_G+C_\Lambda\) on \(\Lambda\).
\end{enumerate}
\end{proposition}
\begin{proof}
\uses{ana:relative-principalization,ana:modification-big-nonkahler,ana:copy4-steps-one-to-five,ana:copy4-steps-three-to-six,ana:polar-branch-pole,ana:step-five-local-branch}
\leanok
Proposition~\ref{ana:copy4-steps-one-to-five}, with relative principalization supplied by
Lemma~\ref{ana:relative-principalization}, gives a sequence of blowups \(\pi:Y\to X\)
adapted to \(V\), whose strict transform \(W\) is a closed connected submanifold of dimension
\(k\), carrying Step-1 data. \(Y\) is compact and connected, as a sequence of blowups of the
compact connected \(X\), and K\"ahler by Lemma~\ref{ana:modification-big-nonkahler}.
Proposition~\ref{ana:copy4-steps-three-to-six}
gives a K\"ahler-current potential \(\varphi_G\) of \(\theta_Y\) with Step-5 data;
\(\varphi_G=-\infty\) on \(W\) by Lemma~\ref{ana:polar-branch-pole}, which is (i); and
Lemma~\ref{ana:step-five-local-branch} gives \(y\), \(U\) and \(\varphi_U\) with (ii) and (iii).
\end{proof}
As in Copy-4, \(\pi\) is the modification built in Step~1 from the ideals of the positive
current, so that they become invertible along \(W\).

\begin{lemma}[Blanchard weight of the exceptional centre]
\label{ana:blanchard-weight}
\uses{ana:blowups,ana:blowup-exceptional-image,ana:geometric-setting}
\lean{TranscendentalNakamaye.Pending.BlanchardWeight.BlanchardData,TranscendentalNakamaye.Pending.BlanchardWeight.BlanchardData.weight,TranscendentalNakamaye.Pending.BlanchardWeight.BlanchardData.weight_eq_zero_of_mem,TranscendentalNakamaye.Pending.BlanchardWeight.BlanchardData.weight_pos_of_notMem,TranscendentalNakamaye.Pending.BlowupExceptionalCentre.exceptionalCentre,TranscendentalNakamaye.Pending.BlowupExceptionalCentre.exists_blanchardData_exceptionalCentre}
\leanok
Let \(g:Z\to Y\) be the blowup of a compact complex manifold \(Y\) of dimension
\(n\) along a nowhere dense closed submanifold \(C\), and let
\(C_2=g(\operatorname{Exc}g)\) be its exceptional centre. There are finitely
many smooth functions \(\chi_a\ge0\) on \(Y\) with \(\sum_a\chi_a=1\) and, for
each index \(a\) of a subset \(A_0\), a point \(p_a\in C_2\), an integer
\(k_a\) with \(k_a+2\le n\) and a holomorphic chart \(\Phi_a\) at \(p_a\) with
domain \(D_a\supset\operatorname{supp}\chi_a\) and coordinates
\(\zeta^{(a)}_0,\dots,\zeta^{(a)}_{n-1}\), such that
\[
 C\cap D_a=C_2\cap D_a=
 \{\zeta^{(a)}_{k_a}=\cdots=\zeta^{(a)}_{n-1}=0\},
\]
while \(\operatorname{supp}\chi_a\cap C_2=\emptyset\) for \(a\notin A_0\). The
associated Blanchard weight
\[
 U=\sum_a\chi_aN_a,\qquad
 N_a=\sum_{i\ge k_a}|\zeta^{(a)}_i|^2\ (a\in A_0),\qquad
 N_a=1\ (a\notin A_0),
\]
is a nonnegative smooth function on \(Y\) whose zero set is exactly \(C_2\).
\end{lemma}
\begin{proof}
\leanok
By Lemma~\ref{ana:blowup-exceptional-image}, \(C_2\) is a closed subset of
\(C\), namely the union of those components of \(C\) whose codimension in \(Y\)
is at least two. Let \(p\in C_2\) and let \(k\) be the dimension of the
component of \(C\) through \(p\). Since \(C\) is a closed submanifold, there is
a holomorphic chart at \(p\) with domain \(D\) in which \(C\cap D\) becomes the
coordinate plane \(\{\zeta_k=\cdots=\zeta_{n-1}=0\}\); shrinking \(D\) we may
assume that \(C\cap D\) meets only the component of \(C\) through \(p\). That
component has codimension at least two, since \(p\in C_2\), so \(k+2\le n\),
and every point of \(C\cap D\) lies in a component of codimension at least two,
hence in \(C_2\) by Lemma~\ref{ana:blowup-exceptional-image}. Thus
\(C\cap D=C_2\cap D\).

The domains \(D\) so obtained, together with \(Y\setminus C_2\), form an open
cover of the compact manifold \(Y\). Take a finite smooth partition of unity
\((\chi_a)\) subordinate to it, let \(A_0\) collect the indices whose member of
the cover is one of the domains \(D\), and define \(N_a\) and \(U\) as in the
statement; each \(N_a\) is smooth and nonnegative on \(D_a\), so \(U\) is
smooth and \(U\ge0\).

Let \(y\in C_2\) and \(\chi_a(y)\ne0\). Then \(a\in A_0\), because the other
partition functions vanish on \(C_2\), and \(y\in D_a\cap C_2\), so
\(\zeta^{(a)}_i(y)=0\) for all \(i\ge k_a\) and \(N_a(y)=0\); hence \(U(y)=0\).
Let now \(y\notin C_2\) and choose \(a\) with \(\chi_a(y)>0\), which is
possible because \(\sum_a\chi_a=1\). If \(a\notin A_0\) then \(N_a(y)=1\); if
\(a\in A_0\) then \(y\in D_a\setminus C_2=D_a\setminus C\), so some
\(\zeta^{(a)}_i(y)\) with \(i\ge k_a\) is nonzero and \(N_a(y)>0\). In both
cases \(U(y)>0\).
\end{proof}

\begin{lemma}[Factorization of the Blanchard weight on the blowup]
\label{ana:blanchard-form}
\uses{ana:blanchard-weight,ana:blowups,ana:blowup-exceptional-image,ana:ideal-sheaves,ana:geometric-setting}
\lean{TranscendentalNakamaye.Pending.FactorizedWeightForm.WeightFactorization,TranscendentalNakamaye.Pending.FactorizedWeightForm.WeightFactorization.closedForm,TranscendentalNakamaye.Pending.BlanchardForm.exists_localFactor_of_flat,TranscendentalNakamaye.Pending.BlowupKahler.exists_blanchardFactorization_of_flat}
\leanok
Keep the notation of Lemma~\ref{ana:blanchard-weight}. Then:
\begin{enumerate}
\item every \(z\in Z\) has an open neighbourhood \(S\) carrying a holomorphic
      function \(h\) which vanishes identically on no nonempty open subset of
      \(S\) and a smooth function \(U'>0\) with \(U\circ g=|h|^2U'\) on \(S\);
\item the forms \(dd^c\log U'\) attached to any two such factorizations agree
      on the overlap of their domains, hence glue to a smooth closed real
      \((1,1)\)-form \(\Theta\) on \(Z\), equal to \(dd^c\log(U\circ g)\) on
      \(Z\setminus\operatorname{Exc}g\);
\item \(\Theta(v,\mathrm i v)>0\) at \(z\) for every \(z\in Z\) and every nonzero
      \(v\) with \(dg_z(v)=0\).
\end{enumerate}
\end{lemma}
\begin{proof}
\leanok
(1) If \(g(z)\notin C_2\), take \(S=Z\setminus g^{-1}(C_2)\), \(h=1\) and
\(U'=U\circ g\), which is positive there by
Lemma~\ref{ana:blanchard-weight}.

Let now \(g(z)\in C_2\). The supports of the \(\chi_a\) with \(a\notin A_0\)
are closed and disjoint from \(C_2\), so on a neighbourhood of \(g(z)\) one has
\(U=\sum_{a\in A_1}\chi_aN_a\), where \(A_1\subset A_0\) is the finite set of
indices with \(g(z)\in D_a\). Fix \(a_0\in A_1\) with \(\chi_{a_0}(g(z))>0\)
and write \(\zeta=\zeta^{(a_0)}\), \(k=k_{a_0}\). For \(a\in A_1\) the
functions \(\zeta^{(a)}_i\), \(i\ge k_a\), and \(\zeta_l\), \(l\ge k\), are two
systems of normal coordinates for the same submanifold \(C\) near \(g(z)\), so
\(n-k_a=n-k\) and, by Hadamard's lemma applied on a smaller neighbourhood,
there are holomorphic functions \(H^{(a)}_{il}\) with
\[
 \zeta^{(a)}_i=\sum_{l\ge k}H^{(a)}_{il}\zeta_l\quad(i\ge k),
 \qquad H^{(a_0)}_{il}=\delta_{il}.
\]
Because \(\mathcal I_C\mathcal O_Z\) is invertible and the blowup is unique up
to canonical isomorphism over its base, \(g\) is, over \(D_{a_0}\),
isomorphic to the blowup of the model \(\{\zeta\}\) along
\(\{\zeta_k=\cdots=\zeta_{n-1}=0\}\), which is covered by the monomial charts
\[
 \gamma_j(w):\quad \zeta_i=w_i\ (i<k),\qquad \zeta_j=w_j,\qquad
 \zeta_i=w_jw_i\ (i\ge k,\ i\ne j),
\]
for \(k\le j\le n-1\). Fix \(j\) with \(z\) in the image of \(\gamma_j\) and
put \(\hat w_j=1\) and \(\hat w_i=w_i\) for \(i\ne j\), so that
\(\zeta_i\circ\gamma_j=w_j\hat w_i\) for \(i\ge k\). Writing
\(c_a=\chi_a\circ\gamma_j\) we obtain
\[
 U(\gamma_j(w))=\sum_{a\in A_1}c_a(w)\sum_i
 \Bigl|\sum_{l\ge k}H^{(a)}_{il}(\gamma_j(w))\,w_j\hat w_l\Bigr|^2
 =|w_j|^2\,\tilde U(w),\qquad
 \tilde U(w)=\sum_{a\in A_1}c_a(w)\sum_i
 \Bigl|\sum_{l\ge k}H^{(a)}_{il}(\gamma_j(w))\,\hat w_l\Bigr|^2 .
\]
The function \(\tilde U\) is smooth, and it is positive near the point \(w_0\)
corresponding to \(z\): the block \(a=a_0\) contributes
\(c_{a_0}(w)\sum_{i\ge k}|\hat w_i|^2\ge c_{a_0}(w)>0\), because
\(\hat w_j=1\). This is a factorization as required, with \(h=w_j\), which
vanishes exactly on the exceptional divisor of the chart and hence on no
nonempty open set.

(2) Suppose \(|h_1|^2U_1'=|h_2|^2U_2'\) on an open set, with \(U_i'\) smooth
and positive and \(h_i\) holomorphic vanishing identically on no nonempty open
subset. Off the zero set of \(h_1h_2\), a dense open subset, one has
\(\log U_1'-\log U_2'=\log|h_2/h_1|^2\), which is pluriharmonic, so
\(dd^c\log U_1'=dd^c\log U_2'\) there; both sides are continuous, hence they
agree everywhere on the overlap. The local forms \(dd^c\log U'\) therefore glue
to a global form \(\Theta\), smooth, closed and real of type \((1,1)\) because
each of them is. Where \(U\circ g>0\) one may take \(h=1\) and \(U'=U\circ g\),
and \(g^{-1}(C_2)=\operatorname{Exc}g\) by
Lemma~\ref{ana:blowup-exceptional-image}, so \(\Theta=dd^c\log(U\circ g)\) on
\(Z\setminus\operatorname{Exc}g\).

(3) If \(dg_z(v)=0\) for some \(v\ne0\), then \(dg_z\) is not injective, so
\(z\in\operatorname{Exc}g\) and \(g(z)\in C_2\); keep the notation of (1) and
let \(w_0\) be the point of the monomial chart corresponding to \(z\). The
fibre of \(g\) through \(z\) is \(\{w:w_j=0,\ w_i=(w_0)_i\text{ for }i<k\}\),
because the normal coordinates \(\zeta_i\circ\gamma_j=w_j\hat w_i\) vanish
exactly when \(w_j=0\); hence \((w_0)_j=0\), and the kernel vectors \(v\) at
\(w_0\) are exactly those with \(v_j=0\) and \(v_i=0\) for \(i<k\). Along the
complex line \(w\mapsto w_0+tv\) the map \(\gamma_j\) is constant: the
coordinates \(\zeta_i\) with \(i<k\) are constant because \(v_i=0\), and those
with \(i\ge k\) vanish identically because \(w_j\) does. Therefore the
functions \(c_a\) and \(H^{(a)}_{il}\circ\gamma_j\) are constant along the
line, while \(\hat w_l(w_0+tv)=\hat w_l(w_0)+tv_l\) for every \(l\), so the
formula of (1) gives
\[
 \tilde U(w_0+tv)=\|X+tY\|^2,
\]
where \(X\) and \(Y\) are the vectors with components
\(\sqrt{c_a}\sum_lH^{(a)}_{il}\hat w_l(w_0)\) and
\(\sqrt{c_a}\sum_lH^{(a)}_{il}v_l\), indexed by the pairs \((a,i)\). The
\((a_0,j)\)-components are \(\sqrt{c_{a_0}}\ne0\) and \(\sqrt{c_{a_0}}v_j=0\),
and \(Y\ne0\) because \(v\ne0\) forces \(v_i\ne0\) for some \(i\ge k\),
\(i\ne j\), whose \((a_0,i)\)-component of \(Y\) is \(\sqrt{c_{a_0}}v_i\).
Hence \(X\) and \(Y\) are linearly independent over \(\mathbf C\): from
\(Y=\lambda X\) the \((a_0,j)\)-components give \(\lambda=0\) and then
\(Y=0\). Now
\[
 \partial_t\partial_{\bar t}\log\|X+tY\|^2\Big|_{t=0}
 =\frac{\|X\|^2\|Y\|^2-|\langle X,Y\rangle|^2}{\|X\|^4}>0
\]
by the Cauchy--Schwarz inequality, which is strict exactly because \(X\) and
\(Y\) are linearly independent. Since restricting \(\log\tilde U\) to the line
\(t\mapsto w_0+tv\) computes \(dd^c\log\tilde U(v,\mathrm i v)\) at \(w_0\) as a
positive multiple of this quantity, \(\Theta(v,\mathrm i v)>0\) at \(z\).
\end{proof}

\begin{lemma}[K\"ahler correction for a single blowup]
\label{ana:blowup-kahler-correction}
\uses{ana:blanchard-form,ana:blanchard-weight,ana:blowup-compact,ana:blowups,ana:geometric-setting}
\source{transcendental-nakamaye-copy4-20260911:lem:exceptional-metric}
\lean{TranscendentalNakamaye.Pending.BlowupKahler.exists_isKahler_pullback_add_of_flat}
\leanok
Let \(g:Z\to Y\) be the blowup of a compact complex manifold \(Y\) along a
nowhere dense closed submanifold \(C\), let \(\beta\) be a K\"ahler form on
\(Y\), and let \(\Theta\) be the form of Lemma~\ref{ana:blanchard-form}. Then
\(Z\) is compact and there is \(\varepsilon>0\) such that
\(g^*\beta+\varepsilon\Theta\) is a K\"ahler form on \(Z\).
\end{lemma}
\begin{proof}
\leanok
\(Z\) is compact by Lemma~\ref{ana:blowup-compact}. The form \(A:=g^*\beta\) is
smooth, closed and real of type \((1,1)\), and it is semipositive:
\(A(v,\mathrm i v)=\beta(dg_z(v),\mathrm i\,dg_z(v))\ge0\), with equality if and only if
\(dg_z(v)=0\), because \(\beta\) is strictly positive.

Fix \(z\in Z\) and read \(A\) and \(\Theta\) in a chart at \(z\) as alternating
two-forms \(a\) and \(b\) on \(\mathbf C^n\). Let \(S\) be the unit sphere and
\(K=\{v\in S:b(v,\mathrm i v)\le0\}\), a compact set. If \(v\in K\) then
\(b(v,\mathrm i v)\le0\) forces \(dg_z(v)\ne0\) by Lemma~\ref{ana:blanchard-form}(3),
so \(a(v,\mathrm i v)>0\); by compactness \(a(v,\mathrm i v)\ge m\) on \(K\) for some \(m>0\),
and \(|b(v,\mathrm i v)|\le B\) on \(S\) for some \(B\ge0\). Choose \(\varepsilon_0\)
with \(0<\varepsilon_0<m/(B+1)\). For \(v\in S\cap K\) we get
\((a+\varepsilon_0b)(v,\mathrm i v)\ge m-\varepsilon_0B>0\), and for \(v\in S\setminus
K\) both summands are nonnegative with the second positive; by homogeneity of
\(v\mapsto(a+\varepsilon_0b)(v,\mathrm i v)\), the form \(A+\varepsilon_0\Theta\) is
strictly positive at \(z\). Moreover, for \(0<\varepsilon\le\varepsilon_0\),
\[
 A+\varepsilon\Theta
 =\Bigl(1-\frac{\varepsilon}{\varepsilon_0}\Bigr)A
 +\frac{\varepsilon}{\varepsilon_0}\,(A+\varepsilon_0\Theta)
\]
is a sum of a semipositive and a strictly positive form, hence strictly
positive.

Strict positivity at a point is an open condition, because the coefficients of
the forms are continuous; so each \(z\) has a neighbourhood \(V_z\) on which
\(A+\varepsilon_0(z)\Theta\) is strictly positive. Cover the compact manifold
\(Z\) by finitely many \(V_{z_1},\dots,V_{z_N}\) and put
\(\varepsilon=\min_i\varepsilon_0(z_i)>0\). At a point of \(V_{z_i}\) the form
\(A\) is semipositive and \(A+\varepsilon_0(z_i)\Theta\) is strictly positive,
so the displayed convexity identity makes \(A+\varepsilon\Theta\) strictly
positive there. Thus \(g^*\beta+\varepsilon\Theta\) is a closed real
\((1,1)\)-form which is strictly positive everywhere, i.e. a K\"ahler form.
\end{proof}

\begin{lemma}[Exceptional correction of a pulled-back K\"ahler form]
\label{ana:exceptional-metric}
\uses{ana:blowups,ana:geometric-setting,ana:ideal-sheaves}
\source{transcendental-nakamaye-copy4-20260911:lem:exceptional-metric}
\lean{TranscendentalNakamaye.Pending.ExceptionalMetric.exists_exceptionalMetric}
\leanok
Let \(\pi:Y\to X\) be a modification of a compact K\"ahler manifold
\((X,\omega)\). There is an effective \(\mathbf Q\)-divisor \(E\) on \(Y\),
with a smooth Hermitian metric on \(\mathcal O_Y(E)\) and support contained in
\(\operatorname{Exc}\pi\), such that \(\omega_Y:=\pi^*\omega-\Theta_E\) is a
K\"ahler form; moreover \(E=0\) whenever \(\pi\) is an isomorphism over \(X\).
\end{lemma}
\begin{proof}
\leanok
\uses{ana:blanchard-form,ana:blanchard-weight,ana:blowup-compact,ana:blowup-exceptional-image,ana:blowup-kahler-correction,ana:strict-transform-composition}
We show by induction on the number of blowups in the sequence that for every
modification \(\pi:Y\to X\) the manifold \(Y\) is compact and carries an
effective \(\mathbf Q\)-divisor \(E\) with a smooth metric such that
\(\operatorname{supp}E\subset\operatorname{Exc}\pi\) and
\(\pi^*\omega-\Theta_E\) is a K\"ahler form.

If \(\pi=\operatorname{id}_X\), take \(E=0\) with the trivial metric: then
\(\Theta_E=0\), \(\pi^*\omega=\omega\) is a K\"ahler form, and
\(\operatorname{supp}E=\emptyset\).

For the inductive step write \(\pi=f\circ g\), where \(f:Y'\to X\) is a
modification for which the claim holds, with divisor \(E'\), and where
\(g:Z\to Y'\) is the blowup along a nowhere dense closed submanifold
\(C\subset Y'\). Thus \(Y'\) is compact and \(\beta:=f^*\omega-\Theta_{E'}\) is
a K\"ahler form on \(Y'\), and \(Z\) is compact by
Lemma~\ref{ana:blowup-compact}. Let \(U\) be the Blanchard weight of the
exceptional centre \(C_2=g(\operatorname{Exc}g)\)
(Lemma~\ref{ana:blanchard-weight}) and \(\Theta\) the associated form on \(Z\)
(Lemma~\ref{ana:blanchard-form}). By
Lemma~\ref{ana:blowup-kahler-correction} there is \(\varepsilon>0\) with
\(g^*\beta+\varepsilon\Theta\) K\"ahler; fix a rational \(q\) with
\(0<q\le\varepsilon\).

\emph{The divisor \(q\operatorname{div}(U\circ g)\).} By
Lemma~\ref{ana:blanchard-form} every point of \(Z\) has a neighbourhood on
which \(U\circ g=|h|^2U'\) with \(h\) holomorphic vanishing identically on no
nonempty open set and \(U'\) smooth and positive. Hence
\(\psi:=q\log(U\circ g)=q\log|h|^2+q\log U'\) is the weight of an effective
\(\mathbf Q\)-divisor with a smooth metric in the sense of
Definition~\ref{ana:ideal-sheaves}: one local equation \(h\), coefficient
\(q\in\mathbf Q_{>0}\) and smooth remainder \(q\log U'\); its curvature form is
\(-dd^c(q\log U')=-q\Theta\). Its support is the polar set of \(\psi\), that is
the zero set of \(U\circ g\), which equals \(g^{-1}(C_2)=\operatorname{Exc}g\)
by Lemmas~\ref{ana:blanchard-weight} and~\ref{ana:blowup-exceptional-image}.

\emph{The pullback \(g^*E'\).} Its weight is \(\psi_{E'}\circ g\) and its
curvature is \(g^*\Theta_{E'}\): pulling back a local presentation
\(\psi_{E'}=\sum_ia_i\log|g_i|^2+h'\) gives a presentation of the same shape,
with local equations \(g_i\circ g\) and smooth remainder \(h'\circ g\). These
pullbacks again vanish identically on no nonempty open set: if \(W\subset Z\)
is open and nonempty, then \(W\not\subset\operatorname{Exc}g\), because
\(\operatorname{Exc}g\subset g^{-1}(C)\) has empty interior
(Lemma~\ref{ana:blowup-exceptional-image}), and \(g\) maps a neighbourhood of a
point of \(W\setminus\operatorname{Exc}g\) onto a neighbourhood of its image
(Lemma~\ref{ana:strict-transform-composition}), so \(g(W)\) has nonempty
interior, on which some \(g_i\) is not identically zero. The support of
\(g^*E'\) is \(g^{-1}(\operatorname{supp}E')\subset
g^{-1}(\operatorname{Exc}f)\).

Put \(E:=g^*E'+q\operatorname{div}(U\circ g)\), with the product metric: its
weight is \(\psi_{E'}\circ g+q\log(U\circ g)\) and its curvature is
\(\Theta_E=g^*\Theta_{E'}-q\Theta\). Then
\[
 \pi^*\omega-\Theta_E=g^*\bigl(f^*\omega-\Theta_{E'}\bigr)+q\Theta
 =g^*\beta+q\Theta .
\]
The form \(A:=g^*\beta\) is closed, real of type \((1,1)\) and semipositive,
and \(A+\varepsilon\Theta\) is K\"ahler, so
\(A+q\Theta=(1-q/\varepsilon)A+(q/\varepsilon)(A+\varepsilon\Theta)\) is
strictly positive; thus \(\pi^*\omega-\Theta_E\) is a K\"ahler form. Finally
\[
 \operatorname{supp}E\subset
 g^{-1}(\operatorname{Exc}f)\cup\operatorname{Exc}g
 =\operatorname{Exc}(f\circ g)=\operatorname{Exc}\pi,
\]
since \(d(f\circ g)_z=df_{g(z)}\circ dg_z\) is a composition of linear maps
between complex vector spaces of the same finite dimension, and such a
composition is bijective if and only if both factors are. This completes the
induction.

If \(\pi\) is an isomorphism over \(X\), then \(d\pi\) is bijective at every
point of \(\pi^{-1}(X)=Y\), so \(\operatorname{Exc}\pi=\emptyset\) and
\(\operatorname{supp}E=\emptyset\), that is \(E=0\).
\end{proof}

\begin{lemma}[The global polar branch on a sequence of blowups]
\label{ana:polar-branch-step-three}
\uses{ana:exceptional-metric,ana:blowups,num:ambient-margin,hgeom:global-theta-pullback}
\source{transcendental-nakamaye-copy4-20260911:prop:principalization,transcendental-nakamaye-copy4-20260911:eq:QM-positivity}
\lean{TranscendentalNakamaye.Pending.PolarBranchStepThree.isKahlerCurrentPotential_polarBranch,TranscendentalNakamaye.Pending.PolarBranchStepThree.exists_exceptionalMetric_isKahlerCurrentPotential_polarBranch,TranscendentalNakamaye.Pending.PolarBranchStepThree.locallyNondegenerate_comp_add_smul}
\leanok
Let \((X,\omega)\) be compact K\"ahler, \(\theta\) a smooth closed real \((1,1)\)-form and
\(\pi:Y\to X\) a finite sequence of blowups with smooth centres. Let \(c_0\ge0\),
\(0<L<M\) and \(\delta_R\) be real numbers with \(c_0<L\delta_R\), and let
\(U\in\operatorname{PSH}(X,\theta+\tfrac{c_0}{M}\omega)\) and
\(\Phi\in\operatorname{PSH}(X,\theta-\delta_R\omega)\). Let \(E\) be an effective
\(\mathbf Q\)-divisor on \(Y\) with a smooth metric, weight \(\psi_E\) and curvature
\(\Theta_E\), supported in \(\operatorname{Exc}\pi\), such that
\(\omega_Y:=\pi^*\omega-\Theta_E\) is K\"ahler (Lemma~\ref{ana:exceptional-metric} provides
one). Put \(\epsilon_0=(L\delta_R-c_0)/M\) and \(Q=(1-\tfrac LM)U+\tfrac LM\Phi\). Then for
every \(0<\tau\le\epsilon_0\),
\[
 \varphi_G:=Q\circ\pi+\tau\psi_E
 \quad\text{satisfies}\quad
 \pi^*\theta+dd^c\varphi_G\ge\tau\omega_Y ,
\]
and \(\varphi_G\) is not identically \(-\infty\) near any point, so
\(\pi^*\theta+dd^c\varphi_G\) is a K\"ahler current.
\end{lemma}
\begin{proof}
\uses{ana:exceptional-metric,num:ambient-margin,hgeom:global-theta-pullback}
\leanok
By convexity, \(\theta+dd^cQ\ge\bigl((1-\tfrac LM)(-\tfrac{c_0}M)+\tfrac LM\delta_R\bigr)\omega
=\epsilon\omega\) with \(\epsilon\ge\epsilon_0>0\) by Lemma~\ref{num:ambient-margin}. Pulling
back (Proposition~\ref{hgeom:global-theta-pullback}),
\(\pi^*\theta+dd^c(Q\circ\pi)\ge\epsilon\,\pi^*\omega\). The weight satisfies
\(dd^c\psi_E=[E]-\Theta_E\ge-\Theta_E\). Adding,
\[
 \pi^*\theta+dd^c\varphi_G\ge\epsilon\,\pi^*\omega-\tau\Theta_E
 =\tau\omega_Y+(\epsilon-\tau)\pi^*\omega\ge\tau\omega_Y ,
\]
because \(\tau\le\epsilon_0\le\epsilon\) and \(\pi^*\omega\ge0\). For nondegeneracy: off
\(\operatorname{Exc}\pi\) the map \(\pi\) is a local biholomorphism and \(\psi_E\) is finite,
so \(\varphi_G\) is not identically \(-\infty\) near such a point because \(Q\) is a genuine
potential; and \(\operatorname{Exc}\pi\) has empty interior, so every neighbourhood of an
exceptional point meets the complement.
\end{proof}
This is Step~3 of Copy-4's proof of the principalization proposition, in a form generic
in the blowup sequence: nothing about the three-stage modification of Step~1 is used.
Copy-4 takes \(\tau<\min\{\epsilon_0,\delta_0/2\}\); only \(\tau\le\epsilon_0\) is needed here,
the bound by \(\delta_0/2\) being used in Step~4. The vanishing \(\varphi_G=-\infty\) on the
strict transform, which Copy-4 reads off from \(\varphi_{Q_M}=-\infty\) on \(V\), is not part
of this lemma.

\begin{lemma}[A quasi-plurisubharmonic weight for a coherent ideal]
\label{ana:ideal-function}
\uses{ana:ideal-sheaves,as:rational-type,ana:geometric-setting}
\source{transcendental-nakamaye-copy4-20260911:lem:ideal-function}
\lean{TranscendentalNakamaye.Pending.IdealFunctionMaxGluing.exists_idealFunction,TranscendentalNakamaye.Pending.IdealFunctionMaxGluing.exists_isThetaPshPotential_idealFunction}
\leanok
Let \((Z,\omega)\) be a compact K\"ahler manifold, \(\mathfrak a\) a coherent ideal sheaf
and \(x_0\in Z\). There are \(c>0\) and a function \(\Psi\) with
\(dd^c\Psi\ge-c\,\omega\), with analytic singularities of type \((1,\mathfrak a)\), that is
\(\Psi=\log\sum_i|g_i|^2+O(1)\) near every point for local generators \(g_i\) of
\(\mathfrak a\), and, if \(x_0\) is not in the zero locus of \(\mathfrak a\), smooth near
\(x_0\) and finite there. If \(Z\) is connected and \(x_0\notin Z(\mathfrak a)\), \(\Psi\) is a
genuine \(c\,\omega\)-plurisubharmonic potential.
\end{lemma}
\begin{proof}
\uses{ana:ideal-sheaves}
\leanok
Cover \(Z\) by finitely many open sets \(U_j\Subset V_j\), each \(\overline V_j\) inside a
chart on which \(\mathfrak a\) is generated by \(g_{j1},\dots,g_{jN_j}\), so that \(x_0\) has a
neighbourhood \(O\subseteq U_{j_0}\) disjoint from every other \(V_j\). Put
\(\psi_j=\log\sum_i|g_{ji}|^2\); since the two families generate the same ideal,
\(\psi_j\le\psi_l+C_B\) on \(V_j\cap V_l\). Take smooth cut-offs \(0\le\chi_j\le1\), equal to
\(1\) on \(U_j\) and supported in \(V_j\), put \(A=C_B+1\), and let \(\Psi\) be the maximum of
\(\psi_j+A\chi_j\) over the \(j\) with \(z\in V_j\). Near a point of \(\partial V_j\) the
cut-off \(\chi_j\) vanishes while some \(\chi_l\) equals \(1\), so the branch \(j\) is strictly
below the branch \(l\) and can be dropped: locally \(\Psi\) is a maximum of finitely many
functions each defined on a whole neighbourhood and each satisfying
\(dd^c(\psi_j+A\chi_j)\ge A\,dd^c\chi_j\), hence \(dd^c\Psi\ge-c\,\omega\) for a \(c\) with
\(A\,dd^c\chi_j\ge-c\,\omega\) for all \(j\). The comparison \(|\psi_j-\psi_l|\le C_B\) gives
\(\Psi=\psi_{j}+O(1)\) near each point. Near \(x_0\) only the branch \(j_0\) is present, which
is smooth off \(Z(\mathfrak a)\). On a connected \(Z\), finiteness at \(x_0\) excludes
\(\Psi\equiv-\infty\) on any open set by the identity principle.
\end{proof}
Copy-4 glues with a \emph{regularized} maximum and obtains smoothness on all of
\(Z\setminus Z(\mathfrak a)\). The plain maximum used here gives smoothness only near the
prescribed \(x_0\); this weaker statement is what Step~5 of the principalization argument
uses, the point being chosen before the divisor is known.

\begin{lemma}[The correction term of Step~5]
\label{ana:step-five-correction}
\uses{ana:ideal-function,ana:divisor-weight-smooth-potential,ana:ideal-sheaves}
\source{transcendental-nakamaye-copy4-20260911:prop:principalization}
\lean{TranscendentalNakamaye.Pending.StepFiveCorrectionTerm.thetaPlurisubharmonic_correction,TranscendentalNakamaye.Pending.StepFiveCorrectionTerm.isThetaPshPotentialOn_correction,TranscendentalNakamaye.Pending.StepFiveCorrectionTerm.isSmoothNear_correction,TranscendentalNakamaye.Pending.StepFiveCorrectionTerm.isSmoothNear_weight_of_notMem_support,TranscendentalNakamaye.Pending.StepFiveCorrectionTerm.moreSingularOn_correction_weight}
\leanok
Let \(\pi:Y\to X\) be a holomorphic map, \(\omega\) a K\"ahler form on \(X\), \(D\) an
effective \(\mathbf Q\)-divisor on \(Y\) with a smooth metric (weight \(\psi_D\), curvature
\(\Theta_D\)), and \(\Psi\) a function on \(X\) with \(dd^c\Psi\ge-c_\Psi\omega\). Let
\(\kappa>0\) and \(\delta\) satisfy \(\kappa c_\Psi\le\delta\), and put
\(\rho=\psi_D+\kappa\,\Psi\circ\pi\).
\begin{enumerate}
\item \(dd^c\rho\ge-(\Theta_D+\delta\,\pi^*\omega)\); if \(Y\) is connected and \(\rho\) is
finite at one point, \(\rho\) is a genuine \((\Theta_D+\delta\,\pi^*\omega)\)-plurisubharmonic
potential on every open set.
\item \(\rho\) is smooth near every \(y\notin\operatorname{Supp}D\) such that \(\Psi\) is smooth
near \(\pi(y)\).
\item If \(X\) is compact, \(\Psi\), \(\psi_{\mathfrak j}\) and \(\psi_{\mathfrak b}\) have analytic
singularities of types \((1,\mathfrak j\mathfrak b)\), \((1,\mathfrak j)\) and
\((1,\mathfrak b)\), and \(\psi_{\mathfrak b}\circ\pi\le\psi_{H}+O(1)\) on \(Y\) for a function
\(\psi_H\), then \(\kappa\,\Psi\circ\pi\le\kappa(\psi_{\mathfrak j}\circ\pi+\psi_H)+O(1)\).
\end{enumerate}
\end{lemma}
\begin{proof}
\uses{ana:divisor-weight-smooth-potential}
\leanok
(i) \(dd^c\psi_D=[D]-\Theta_D\ge-\Theta_D\), and
\(dd^c(\kappa\Psi\circ\pi)\ge-\kappa c_\Psi\pi^*\omega\ge-\delta\,\pi^*\omega\) since
\(\pi^*\omega\ge0\); add. A function satisfying the criterion and finite at one point of a
connected manifold is not identically \(-\infty\) on any open set (identity principle).
(ii) Off \(\operatorname{Supp}D\) the weight \(\psi_D\) is smooth
(Lemma~\ref{ana:divisor-weight-smooth-potential}); compositions, constant multiples and sums
of smooth functions are smooth. (iii) Local generators of \(\mathfrak j\mathfrak b\) are the
products of those of \(\mathfrak j\) and \(\mathfrak b\), so \(\Psi=\psi_{\mathfrak j}+
\psi_{\mathfrak b}+O(1)\) on the compact \(X\); pull back by \(\pi\), use the hypothesis on
\(\psi_{\mathfrak b}\circ\pi\), and multiply by \(\kappa\ge0\).
\end{proof}
In Copy-4's Step~5 this is applied with \(\Psi=\Psi_m\) of Lemma~\ref{ana:ideal-function} for
\(\mathfrak j\mathfrak B_m\), \(\kappa=(C-1)/M\), and \(H=H^{\mathrm{fix}}\); the relation
\(\psi_{\mathfrak B_m}\circ\pi\le\psi_{H^{\mathrm{fix}}}+O(1)\) comes from Step~1 and is a
hypothesis here.

\begin{lemma}[Step~6: a potential dominating a K\"ahler form near the strict transform]
\label{ana:step-six-local-geometry}
\uses{ana:geometric-setting}
\source{transcendental-nakamaye-copy4-20260911:prop:principalization}
\lean{TranscendentalNakamaye.Pending.StepSixLocalGeometry.exists_localGeometry_of_positiveOnOneFlattening}
\leanok
Let \((X,\omega)\) be compact K\"ahler, \(W\subseteq X\) closed and a submanifold of dimension
\(k\) at each of its points, \(\beta\) a smooth closed real \((1,1)\)-form and \(f\) a smooth
function such that \(\beta+dd^cf\) is strictly positive on the tangent directions of \(W\)
(read through one flattening chart at each point). Then there are an open \(U\supseteq W\) and
a smooth \(\sigma\) such that \(\beta+dd^c\sigma\) dominates a K\"ahler form on \(U\).
\end{lemma}
\begin{proof}
\leanok
Take \(\sigma=f+B\,r_W\) with \(r_W\ge0\) a smooth function vanishing to second order along
\(W\) whose complex Hessian is positive definite in the directions normal to \(W\) (built
from the flattening charts and a partition of unity). Along \(W\),
\(\beta+dd^c\sigma=\beta+dd^cf+B\,dd^cr_W\) is positive on tangent directions by hypothesis
and on normal directions once \(B\) is large, and the mixed terms are absorbed by
Cauchy--Schwarz; by compactness of \(W\) and continuity the positivity, with a uniform
K\"ahler lower bound, persists on a neighbourhood \(U\).
\end{proof}

\begin{lemma}[Positivity along \(T_W\) from the restriction to \(W\)]
\label{ana:restriction-positivity-along-w}
\uses{ana:geometric-setting,hgeom:form-pullback,hgeom:ddc-naturality}
\source{transcendental-nakamaye-copy4-20260911:prop:principalization}
\lean{TranscendentalNakamaye.Pending.SubmanifoldRestrictionPositivity.range_fderiv_eq_comap,TranscendentalNakamaye.Pending.SubmanifoldRestrictionPositivity.positiveOnOneFlattening_of_isHermitian_pullback,TranscendentalNakamaye.Pending.SubmanifoldRestrictionPositivity.pullback_ddcClosedForm_coeff,TranscendentalNakamaye.Pending.SubmanifoldRestrictionPositivity.positiveOnOneFlattening_add_ddc_of_isKahler}
\leanok
Let \(X\) be a complex manifold, \(W\subseteq X\) a set every point of which is a regular point
of dimension \(k\), and \(\iota:V\to X\) a holomorphic immersion of a complex \(k\)-manifold
with \(\iota(V)=W\). Let \(\beta\) be a smooth closed real \((1,1)\)-form on \(X\).
\begin{enumerate}
\item At every point of \(W\) and for every nondegenerate holomorphic flattening \(F\) of
\(W\) there, \((dF)^{-1}(\mathbf C^k\times0)\) is the image of the differential of \(\iota\),
read in charts.
\item If \(\iota^*\beta\) is a Hermitian form on \(V\), then \(\beta\) is strictly positive on
the tangent directions of \(W\), read through one flattening at each point of \(W\).
\item For a smooth \(f\) on \(X\), \(\iota^*(dd^cf)=dd^c(f\circ\iota)\).
\item If \(h\) is smooth on \(V\), \(\iota^*\beta+dd^ch\) is K\"ahler on \(V\), and \(\tilde f\) is a
smooth function on \(X\) with \(\tilde f\circ\iota=h\), then \(\beta+dd^c\tilde f\) is strictly
positive on the tangent directions of \(W\) in the sense of~(ii).
\end{enumerate}
\end{lemma}
\begin{proof}
\uses{hgeom:ddc-naturality}
\leanok
(i) \(F\circ\iota\) takes values in \(\mathbf C^k\times0\), so the image of \(d\iota\) lies in
\((dF)^{-1}(\mathbf C^k\times0)\). The differential of \(\iota\) is injective, so its image has
real dimension \(2k\); for a nondegenerate flattening, \((dF)^{-1}(\mathbf C^k\times0)\) has real
dimension at most that of \(\mathbf C^k\times0\), which is \(2k\). The two spaces are therefore
equal. (ii) A nonzero vector of \((dF)^{-1}(\mathbf C^k\times0)\) at \(\iota(v)\) is
\(d\iota(\xi)\) for some nonzero \(\xi\) by~(i), and
\(\beta(d\iota\,\xi,J\,d\iota\,\xi)=(\iota^*\beta)(\xi,J\xi)>0\), since \(d\iota\) is complex
linear. (iii) is Lemma~\ref{hgeom:ddc-naturality} in charts. (iv) By~(iii),
\(\iota^*(\beta+dd^c\tilde f)=\iota^*\beta+dd^c(\tilde f\circ\iota)=\iota^*\beta+dd^ch\), which is
K\"ahler, in particular Hermitian; apply~(ii).
\end{proof}

\begin{lemma}[Step~4: a K\"ahler class on the strict transform]
\label{ana:step-four-kahler-class-on-w}
\uses{ana:residual-nef,ana:restriction-positivity-along-w,ana:step-six-local-geometry}
\source{transcendental-nakamaye-copy4-20260911:prop:principalization}
\lean{TranscendentalNakamaye.Pending.StepFourKahlerClassOnW.exists_positiveOnOneFlattening_of_logSingularities,TranscendentalNakamaye.Pending.StepFourKahlerClassOnW.exists_localGeometry_of_logSingularities}
\leanok
Let \(X\), \(W\), \(\iota:V\to X\) and \(\beta\) be as in
Lemma~\ref{ana:restriction-positivity-along-w}, with \(V\) compact, connected and nonempty,
and assume that every smooth function on \(V\) is \(f\circ\iota\) for some smooth \(f\) on \(X\).
Let \(\theta'\) be a smooth closed real \((1,1)\)-form on \(V\) and \(\varphi'\) a potential with
\(\theta'+dd^c\varphi'\) a K\"ahler current and with logarithmic singularities along an
effective \(\mathbf Q\)-divisor \(D_W\) on \(V\) carrying a smooth metric of curvature
\(\Theta_{D_W}\), and suppose \(\theta'-\Theta_{D_W}=\iota^*\beta\) as forms. Then there is a
smooth \(\tilde f\) on \(X\) with \(\beta+dd^c\tilde f\) strictly positive on the tangent
directions of \(W\). If moreover \(X\) is compact K\"ahler and \(W\) is closed, there are an
open \(U\supseteq W\) and a smooth \(\sigma\) on \(X\) with \(\beta+dd^c\sigma\) dominating a
K\"ahler form on \(U\).
\end{lemma}
\begin{proof}
\uses{ana:residual-nef,ana:restriction-positivity-along-w,ana:step-six-local-geometry}
\leanok
The K\"ahler half of Lemma~\ref{ana:residual-nef}, applied on \(V\), gives a smooth \(h\) on
\(V\) with \(\theta'-\Theta_{D_W}+dd^ch=\iota^*\beta+dd^ch\) K\"ahler. Extend \(h\) to a smooth
\(\tilde f\) on \(X\); part~(iv) of Lemma~\ref{ana:restriction-positivity-along-w} gives the
first assertion, and Lemma~\ref{ana:step-six-local-geometry} applied to \(\beta\) and
\(\tilde f\) gives the second.
\end{proof}
Copy-4 applies Lemma~\ref{ana:residual-nef} ``on \(W\)'', to the restricted class
\(\{\theta_Y-\Theta_D-\tfrac12\delta_0\pi^*\omega\}|_W\). Here \(W\) carries its complex
structure through the immersion \(\iota\); the restriction of the divisor and its metric enters
only through the identity \(\theta'-\Theta_{D_W}=\iota^*\beta\) of forms on \(V\), and the
extension of smooth functions along \(\iota\) is a hypothesis. The assumption that every smooth
function extends forces \(\iota\) to be injective, so \(W\) is an embedded submanifold.


\begin{lemma}[A closed complex submanifold is a complex manifold, and its inclusion an immersion]
\label{ana:submanifold-complex-structure}
\uses{ana:geometric-setting,ana:blowups}
\lean{TranscendentalNakamaye.Pending.SubmanifoldChartedSpace.chartedSpace,TranscendentalNakamaye.Pending.SubmanifoldChartedSpace.isComplexManifold_chartedSpace,TranscendentalNakamaye.Pending.SubmanifoldChartedSpace.contMDiff_val,TranscendentalNakamaye.Pending.StepFourOnStrictTransform.fderiv_coordinateMap_val_injective}
\leanok
Let \(W\) be a subset of a complex manifold \(X\) of dimension \(n\) which is regular of
dimension \(k\le n\) at each of its points. With the subspace topology and the atlas of
flattening charts (at each point, a holomorphic chart of \(X\) followed by a biholomorphism
straightening \(W\) to \(\mathbf C^k\times0\), restricted to \(W\)), \(W\) is a complex manifold
of dimension \(k\), and the inclusion \(W\hookrightarrow X\) is holomorphic, with injective
differential in every pair of charts.
\end{lemma}
\begin{proof}
\leanok
The transition between two flattening charts is the composite of holomorphic maps restricted
to \(\mathbf C^k\times0\), hence holomorphic; Hausdorffness and second countability pass to
the subspace. In a flattening chart \(F\) and the ambient chart it is built from, the inclusion
reads \(\zeta\mapsto\Phi_F^{-1}(\zeta,0)\), which is holomorphic. In a flattening chart \(F\) and an
ambient chart \(d\), the inclusion read in coordinates has the holomorphic left inverse \(w\mapsto\operatorname{pr}_k(F(d^{-1}(w)))\)
near the image point, so its differential is injective.
\end{proof}

\begin{lemma}[Holomorphic maps into a submanifold]
\label{ana:submanifold-corestriction}
\uses{ana:submanifold-complex-structure}
\source{transcendental-nakamaye-copy4-20260911:prop:principalization}
\lean{TranscendentalNakamaye.Pending.SubmanifoldCodRestrict.contMDiff_codRestrict,TranscendentalNakamaye.Pending.SubmanifoldCodRestrict.contMDiff_restrictPreimage}
\leanok
Let \(W\subseteq X\) be as in Lemma~\ref{ana:submanifold-complex-structure}, with its complex
structure. If \(f:M\to X\) is a holomorphic map of complex manifolds with \(f(M)\subseteq W\), then
\(f\), regarded as a map \(M\to W\), is holomorphic. In particular, if \(\pi:Y\to X\) is holomorphic
and \(W'\subseteq Y\) is regular of some dimension \(k'\) at each of its points, with
\(\pi(W')\subseteq W\), then the restriction \(\pi|_{W'}:W'\to W\) is holomorphic.
\end{lemma}
\begin{proof}
\leanok
Let \(y\in M\) and let \(F\) be the flattening chart of \(W\) at \(f(y)\): a chart \(e_F\) of \(X\)
followed by a biholomorphism \(\Phi_F\) straightening \(W\) to \(\mathbf C^k\times0\). The map
\(f\) is continuous into \(W\) for the subspace topology, and in the chart \(F\) it reads
\(z\mapsto\operatorname{pr}_k(\Phi_F(e_F(f(z))))\) near \(y\), a composite of holomorphic maps. For the
second assertion apply the first to \(\pi\circ\iota_{W'}\), where the inclusion
\(\iota_{W'}:W'\to Y\) is holomorphic by Lemma~\ref{ana:submanifold-complex-structure}.
\end{proof}

\begin{lemma}[Smooth extension from a closed submanifold]
\label{ana:smooth-extension}
\uses{ana:submanifold-complex-structure}
\lean{TranscendentalNakamaye.Pending.StepFourOnStrictTransform.exists_local_extension,TranscendentalNakamaye.Pending.StepFourOnStrictTransform.exists_chartSmooth_extension}
\leanok
Let \(X\) be a compact complex manifold and \(W\subseteq X\) closed and regular of dimension
\(k\) at each of its points, with the complex structure of
Lemma~\ref{ana:submanifold-complex-structure}. Every smooth function on \(W\) is the
restriction of a smooth function on \(X\).
\end{lemma}
\begin{proof}
\leanok
Near a point of \(W\), on the domain \(S\) of a flattening chart \(F\), the function
\(h\circ F|_W^{-1}\circ\operatorname{pr}_k\circ F\) is smooth and agrees with \(h\) on \(W\cap S\)
(it is constant in the normal coordinates). Cover the compact \(W\) by finitely many such
domains, add the open set \(X\setminus W\) with the function \(0\), and glue with a smooth
partition of unity subordinate to this cover; on \(W\) only the first kind of term is present,
and each equals \(h\) there.
\end{proof}
This is the smooth extension lemma (Lee, \emph{Introduction to Smooth Manifolds}, Lemma~5.34),
in the chart-smooth setting used here.

\begin{definition}[Restriction of a metered divisor]
\label{ana:divisor-restriction}
\uses{ana:geometric-setting}
\source{transcendental-nakamaye-copy4-20260911:prop:principalization}
\lean{TranscendentalNakamaye.Pending.MeteredDivisorRestriction.restrictAlong}
\leanok
Let \(\iota:V\to X\) be a holomorphic map and \(D\) an effective \(\mathbf Q\)-divisor on \(X\)
with a smooth metric, of weight \(\psi_D\) and curvature \(\Theta_D\), such that
\(\psi_D\circ\iota\) is \(-\infty\) on no nonempty open subset of \(V\) (no component of \(D\)
contains \(\iota(V)\)). The \emph{restriction} \(\iota^*D\) is the effective
\(\mathbf Q\)-divisor on \(V\) with weight \(\psi_D\circ\iota\) and curvature
\(\iota^*\Theta_D\): the local presentations of \(D\) pull back along \(\iota\).
\end{definition}

\begin{lemma}[Copy-4 Step~4 on the strict transform]
\label{ana:step-four-on-w}
\uses{ana:residual-nef,ana:submanifold-complex-structure,ana:smooth-extension,ana:divisor-restriction,ana:restriction-positivity-along-w,ana:step-six-local-geometry}
\source{transcendental-nakamaye-copy4-20260911:prop:principalization}
\lean{TranscendentalNakamaye.Pending.MeteredDivisorRestriction.exists_positiveOnOneFlattening_of_restrictedDivisor,TranscendentalNakamaye.Pending.StepFourOnStrictTransform.exists_positiveOnOneFlattening_of_logSingularities_strictTransform,TranscendentalNakamaye.Pending.StepFourOnStrictTransform.exists_localGeometry_of_logSingularities_strictTransform}
\leanok
Let \(X\) be a compact complex manifold, \(W\subseteq X\) a closed connected submanifold of
dimension \(k\) with the complex structure of Lemma~\ref{ana:submanifold-complex-structure},
\(\theta_Y,\omega_Y\) smooth closed real \((1,1)\)-forms, \(\delta\in\mathbf R\), and \(D\) an
effective \(\mathbf Q\)-divisor with a smooth metric, no component of which contains \(W\). If
\(\varphi'\) is a potential on \(W\) of a K\"ahler current in the class of
\(\theta'=(\theta_Y-\delta\omega_Y)|_W\) with logarithmic singularities along \(D|_W\) (with the
restricted metric), then there is a smooth \(\tilde f\) on \(X\) such that
\(\theta_Y-\Theta_D-\delta\omega_Y+dd^c\tilde f\) is strictly positive on the tangent directions
of \(W\). The same holds with \(\theta'-\Theta_{D_W}=\beta|_W\) for any smooth closed real
\((1,1)\)-form \(\beta\) on \(X\) and divisor \(D_W\) on \(W\), with \(\beta+dd^c\tilde f\) in the
conclusion; and if \(X\) is K\"ahler, Lemma~\ref{ana:step-six-local-geometry} then gives an open
\(U\supseteq W\) and a smooth \(\sigma\) with \(\beta+dd^c\sigma\) dominating a K\"ahler form on
\(U\).
\end{lemma}
\begin{proof}
\uses{ana:step-four-kahler-class-on-w,ana:smooth-extension,ana:divisor-restriction}
\leanok
Apply Lemma~\ref{ana:step-four-kahler-class-on-w} to the inclusion \(\iota:W\hookrightarrow X\):
it is an immersion with image \(W\) (Lemma~\ref{ana:submanifold-complex-structure}), and smooth
functions on \(W\) extend (Lemma~\ref{ana:smooth-extension}). For the first form, the curvature
of \(D|_W\) is \(\iota^*\Theta_D\) (Definition~\ref{ana:divisor-restriction}), so
\(\theta'-\Theta_{D|_W}=\iota^*(\theta_Y-\Theta_D-\delta\omega_Y)\).
\end{proof}
This is Copy-4's Step~4 in its own terms: there \(\varphi'=(\pi|_W)^*\varphi_{K,\ell}+\tau
\psi_{E_{\mathrm{rel}}}|_W\), built from the positive current of Copy-4's section on a positive
current on the submanifold; here \(\varphi'\) is a hypothesis. Lemma~\ref{ana:step-four-current}
shows that this \(\varphi'\) is a potential of a K\"ahler current in the class of \(\theta'\) for every
potential \(\varphi_K\) on \(V\) with \(\theta|_V-\delta_0\omega|_V+dd^c\varphi_K\ge0\); its logarithmic
singularities along \(D|_W\) follow from the singularity type of \(\varphi_{K,\ell}\)
(Lemma~\ref{ana:step-four-log-singularities}).

\begin{lemma}[Copy-4 Step~4: the current on the strict transform]
\label{ana:step-four-current}
\uses{ana:submanifold-complex-structure,ana:submanifold-corestriction,hgeom:form-pullback-functorial,hgeom:global-theta-pullback,hgeom:pullback-semipositive,hgeom:pullback-kahler,ana:strict-transform-composition,ana:ideal-sheaves,ana:blowups}
\source{transcendental-nakamaye-copy4-20260911:prop:principalization}
\lean{TranscendentalNakamaye.Pending.StepFourPullbackCurrent.thetaPlurisubharmonic_stepFourCurrent,TranscendentalNakamaye.Pending.StepFourPullbackCurrent.locallyNondegenerate_stepFourCurrent,TranscendentalNakamaye.Pending.StepFourPullbackCurrent.isKahlerCurrentPotential_stepFourCurrent}
\leanok
Let \(X\) and \(Y\) be complex manifolds of dimension \(n\), \(\omega\) a K\"ahler form and \(\theta\) a
smooth closed real \((1,1)\)-form on \(X\), and \(\pi:Y\to X\) holomorphic. Let \(V\subseteq X\) and
\(W\subseteq Y\) be regular of dimension \(k\) at each of their points, with the complex structures
of Lemma~\ref{ana:submanifold-complex-structure} and inclusions \(\iota_V\), \(\iota_W\), and with
\(\pi(W)\subseteq V\). Let \(\delta_0,\delta,\tau\) be reals with \(0<\tau\) and \(\delta+\tau\le\delta_0\),
\(\varphi_K\) a function on \(V\) with \(\iota_V^*(\theta-\delta_0\omega)+dd^c\varphi_K\ge0\), and \(E\) an
effective \(\mathbf Q\)-divisor on \(Y\) with a smooth metric, of weight \(\psi_E\) and curvature
\(\Theta_E\). Put \(\theta'=\iota_W^*(\pi^*\theta-\delta\,\pi^*\omega)\) and
\(\varphi'=\varphi_K\circ\pi|_W+\tau\,\psi_E|_W\).
\begin{enumerate}
\item \(\theta'+dd^c\varphi'\ge\tau\,\iota_W^*(\pi^*\omega-\Theta_E)\) on \(W\), in the sense that
\(\varphi'\) is \((\theta'-\tau\,\iota_W^*(\pi^*\omega-\Theta_E))\)-plurisubharmonic.
\item Let \(W=\widetilde V\) be the strict transform of \(V\) under \(\pi\), and \(\psi\) a function on \(Y\)
which is finite off \(\operatorname{Exc}\pi\). If \(\varphi_K\) is not identically \(-\infty\) near any
point of \(V\), neither is \(\varphi_K\circ\pi|_W+\tau\psi|_W\) near any point of \(W\).
\item Let \(W=\widetilde V\), let \(\varphi_K\) be moreover not identically \(-\infty\) near any point of
\(V\), and let \(\operatorname{supp}E\subseteq\operatorname{Exc}\pi\) with \(\omega_Y=\pi^*\omega-\Theta_E\)
K\"ahler. Then \(\theta'+dd^c\varphi'\) is a K\"ahler current: \(\varphi'\) is a potential of
\(\theta'\) with \(\theta'+dd^c\varphi'\ge\tau\,\omega_Y|_W\), and \(\tau\,\omega_Y|_W\) is a K\"ahler form
on \(W\).
\end{enumerate}
\end{lemma}
\begin{proof}
\leanok
(i) The restriction \(\pi|_W:W\to V\) is holomorphic (Lemma~\ref{ana:submanifold-corestriction}), so
\(\varphi_K\circ\pi|_W\) is \((\pi|_W)^*\iota_V^*(\theta-\delta_0\omega)\)-plurisubharmonic
(Proposition~\ref{hgeom:global-theta-pullback}). Since \(\iota_V\circ\pi|_W=\pi\circ\iota_W\),
Proposition~\ref{hgeom:form-pullback-functorial} identifies this form with
\(\iota_W^*(\pi^*\theta-\delta_0\pi^*\omega)\). The weight satisfies \(dd^c\psi_E=[E]-\Theta_E\ge-\Theta_E\),
that is, \(\psi_E\) is \(\Theta_E\)-plurisubharmonic (locally \(\psi_E=\sum a_i\log|g_i|^2+h\) with
\(a_i\ge0\) and \(\Theta_E=-dd^ch\)), so \(\psi_E\circ\iota_W\) is \(\iota_W^*\Theta_E\)-plurisubharmonic, and
\(\varphi'\) is \((\iota_W^*(\pi^*\theta-\delta_0\pi^*\omega)+\tau\iota_W^*\Theta_E)\)-plurisubharmonic. Finally
\[
 \theta'-\tau\,\iota_W^*(\pi^*\omega-\Theta_E)
 =\iota_W^*(\pi^*\theta-\delta_0\pi^*\omega)+\tau\,\iota_W^*\Theta_E+(\delta_0-\delta-\tau)\,\iota_W^*\pi^*\omega ,
\]
and the last term is semipositive (\(\delta_0-\delta-\tau\ge0\), Proposition~\ref{hgeom:pullback-semipositive}),
so \(\varphi'\) is plurisubharmonic for this form as well.

(ii) Let \(w\in W\) and \(O\) an open neighbourhood of \(w\) in \(W\), \(O=O'\cap W\) with \(O'\) open in
\(Y\). As \(W\) is the closure of \(\pi^{-1}(V)\setminus\operatorname{Exc}\pi\), \(O'\) contains a point
\(z_0\in\pi^{-1}(V)\setminus\operatorname{Exc}\pi\). The set \(N=O'\setminus\operatorname{Exc}\pi\) is open
(Lemma~\ref{ana:strict-transform-composition}) and \(\pi(N)\) is a neighbourhood of \(\pi(z_0)\) (same
lemma). Since \(\varphi_K\) is not identically \(-\infty\) near \(\pi(z_0)\in V\), there is
\(z\in N\) with \(\pi(z)\in V\) and \(\varphi_K(\pi(z))>-\infty\). Then
\(z\in\pi^{-1}(V)\setminus\operatorname{Exc}\pi\subseteq W\), so \(z\in O\), and \(\psi(z)>-\infty\) because
\(z\notin\operatorname{Exc}\pi\); hence the function is finite at \(z\).

(iii) The weight \(\psi_E\) is finite off \(\operatorname{supp}E\), hence off \(\operatorname{Exc}\pi\), so
(ii) applies with \(\psi=\psi_E\). The inclusion \(\iota_W\) has injective differential
(Lemma~\ref{ana:submanifold-complex-structure}), so \(\omega_Y|_W=\iota_W^*\omega_Y\) is K\"ahler
(Proposition~\ref{hgeom:pullback-kahler}), and so is \(\tau\,\omega_Y|_W\). With (i), \(\varphi'\) is a
\((\theta'-\tau\,\omega_Y|_W)\)-plurisubharmonic potential.
\end{proof}
In Copy-4, \(\delta=\tfrac12\delta_0\), \(\tau<\tfrac12\delta_0\), \(E=E_{\mathrm{rel}}\) is the exceptional
divisor of Lemma~\ref{ana:exceptional-metric}, and \(\varphi_K=\varphi_{K,\ell}\); then \(\varphi'\) is the
potential of the field (vii) of Definition~\ref{ana:copy4-step-one-data}; its logarithmic singularities along
\(D|_W\) are Lemma~\ref{ana:step-four-log-singularities}.

\begin{lemma}[A uniform power in the analytic Nullstellensatz]
\label{ana:uniform-nullstellensatz-power}
\uses{ana:ideal-sheaves,lag:prime-nullstellensatz,lag:germ-noetherian}
\source{demailly-cadg:II.4.22}
\lean{TranscendentalNakamaye.Pending.UniformNullstellensatzPower.demaillyII422_exists_idealPow_le_of_zeroLocus_subset,TranscendentalNakamaye.Pending.RueckertNullstellensatz.exists_idealPow_le_of_zeroLocus_subset,TranscendentalNakamaye.Pending.RueckertNullstellensatz.exists_nhds_idealPow_le,TranscendentalNakamaye.Pending.RueckertNullstellensatz.exists_powFam_eq_sum,TranscendentalNakamaye.Pending.RueckertNullstellensatz.isInSpanAt_powFam_of_mem_idealPow,TranscendentalNakamaye.Pending.RueckertNullstellensatz.idealPow_le_of_le}
\leanok
Let \(X\) be a compact complex manifold and \(\mathfrak a,\mathfrak r\) coherent ideal sheaves
with \(Z(\mathfrak r)\subseteq Z(\mathfrak a)\). There is \(N>0\) with
\(\mathfrak a^N\subseteq\mathfrak r\).
\end{lemma}
\begin{proof}
\leanok
Fix \(x\in X\) and, in a chart centred at \(x\), generators \(a_1,\dots,a_K\) of \(\mathfrak a\) and
\(r_1,\dots,r_R\) of \(\mathfrak r\) near \(x\) (coherence). The germs of the \(a_i\) vanish on
\(\{r=0\}=Z(\mathfrak r)\) near \(x\), so by Theorem~\ref{lag:prime-nullstellensatz} they lie in
\(\sqrt{(r)}\). The ideal \((a)\) of \(\mathcal O_n\) is finitely generated (Theorem~\ref{lag:germ-noetherian}),
so \((a)^{N_x}\subseteq(r)\) for some \(N_x\): every product of \(N_x\) of the \(a_i\) is a holomorphic
combination of the \(r_j\), and these finitely many identities hold on a neighbourhood \(U_x\) of \(x\). A
member of \(\mathfrak a^{N_x}\) is locally a holomorphic combination of such products, so
\(\mathfrak a^{N_x}\subseteq\mathfrak r\) on \(U_x\). Finitely many \(U_x\) cover the compact \(X\); with
\(N\) larger than the corresponding \(N_x\), \(\mathfrak a^N\subseteq\mathfrak a^{N_x}\subseteq\mathfrak r\) on
each of them.
\end{proof}
It supplies the power containment \((\mathcal I_V\mathfrak j)^p\subseteq\mathfrak a_R\) behind the
lower bound for \(\Phi_R\) in the ambient data of Definition~\ref{ana:step-five-data}.

\begin{definition}[Step-5 data of a polar branch]
\label{ana:step-five-data}
\uses{ana:geometric-setting,ana:blowups}
\source{transcendental-nakamaye-copy4-20260911:prop:principalization,transcendental-nakamaye-copy4-20260911:sec:ambient}
\lean{TranscendentalNakamaye.Pending.PrincipalizationCopy4StepFiveCoreReduction.Copy4StepFiveExceptionalWeightData}
\leanok
Let \(Y\) be a complex manifold, \(\theta_Y\) a smooth closed real \((1,1)\)-form,
\(W\subseteq Y\) and \(\varphi_G\) a function on \(Y\). \emph{Step-5 data} for \(\varphi_G\)
consist of:
\begin{enumerate}
\item a point \(y\in W\), an open \(U\supseteq W\), a smooth closed real \((1,1)\)-form \(\beta\)
and a smooth \(\sigma\) with \(\beta+dd^c\sigma\) dominating a K\"ahler form on \(U\), and a
function \(\rho\) which is \((\theta_Y-\beta)\)-plurisubharmonic on \(U\) and smooth near \(y\);
\item reals \(\tau\), \(0<L<M\), \(A\ge0\) and \(C\) with \(LA\le C-1\), and a splitting
\[
 \rho=\psi_D+\chi,\qquad \psi_D=\psi_{D'}+\tau\psi_E,\qquad
 \varphi_G=\varphi_Q\circ\pi+\tau\psi_E,\qquad
 \varphi_Q=(1-\tfrac LM)\,U_M+\tfrac LM\,\Phi_R,
\]
for a holomorphic map \(\pi:Y\to X\) to a compact complex manifold and functions
\(U_M,\Phi_R\) on \(X\);
\item functions \(\psi_{\widetilde G}\), \(\psi_{H^{\mathrm{fix}}}\), \(\psi_{P^{\mathrm{mov}}}\),
\(\psi_{P^F}\) on \(Y\) and \(\psi_{\mathfrak j}\) on \(X\), all bounded above, such that, up to
additive constants on all of \(Y\),
\[
 \psi_{D'}\le\tfrac1M\psi_{\widetilde G}+\tfrac{C-1}M\psi_{P^{\mathrm{mov}}}+\tfrac CM\psi_{P^F},\quad
 \chi\le\tfrac{C-1}M\bigl(\psi_{\mathfrak j}\circ\pi+\psi_{H^{\mathrm{fix}}}\bigr),\quad
 \psi_{H^{\mathrm{fix}}}+\psi_{P^{\mathrm{mov}}}+\psi_{P^F}\le\psi_H ;
\]
\item the ambient comparisons: \(U_M\ge\tfrac1M\psi_{\mathfrak B_M}-O(1)\) on \(X\); functions
\(\Phi_R,\psi_{\mathcal I_V},\psi_{\mathfrak j}\) with analytic singularities of types
\((c_Rc,\mathfrak a_R)\), \((c,\mathcal I_V)\), \((c,\mathfrak j)\), where
\((\mathcal I_V\mathfrak j)^p\subseteq\mathfrak a_R\) and \(A=c_Rp\); \(\psi_{\mathfrak B_M}\)
and \(\psi_{\widetilde G}\) with analytic singularities of types \((c',\mathfrak B_M)\) and
\((c',\mathfrak B_M\mathcal O_Y)\); and \(\psi_{\mathcal I_V}\), \(\psi_H\), \(\psi_{\mathcal I_W}\)
with analytic singularities of types \((c'',\mathcal I_V)\), \((c'',\mathfrak h)\),
\((c'',\mathfrak i_W)\) where \(\mathcal I_V\mathcal O_Y=\mathfrak h\cdot\mathfrak i_W\) and
\(\psi_{\mathcal I_W}\) has analytic singularities along an ideal with zero locus \(W\).
\end{enumerate}
\end{definition}
These are the objects of Steps~3--5 of Copy-4's proof of the principalization proposition:
\(\varphi_Q\) is the ambient convex current, \(\psi_E\) the weight of the exceptional metric,
\(D=D'+\tau E\) the divisor of Copy-4's Step~4, \(\chi\) the correction term and
\(\psi_H\) the weight of the divisor \(H\) of Step~2.

\begin{lemma}[The pole of a polar branch along the strict transform]
\label{ana:polar-branch-pole}
\uses{ana:modification-big-nonkahler,ana:nonkahler-locus}
\source{transcendental-nakamaye-copy4-20260911:prop:principalization}
\lean{TranscendentalNakamaye.Pending.PrincipalizationPolarBranchStepFive.eq_bot_on_strictTransform_of_isKahlerCurrentPotential}
\leanok
Let \(X\) be compact K\"ahler, \(\alpha=\{\theta\}\) big, \(V\) a closed irreducible analytic
subvariety of positive dimension which is an irreducible component of
\(E_{\mathrm{nK}}(\alpha)\), and \(\pi:Y\to X\) a sequence of blowups adapted to \(V\), with
strict transform \(W\). Every potential \(\varphi_G\) of a K\"ahler current in
\(\pi^*\alpha\) equals \(-\infty\) on \(W\).
\end{lemma}
\begin{proof}
\uses{ana:modification-big-nonkahler}
\leanok
\(W\subseteq\pi^{-1}(V)\subseteq\pi^{-1}(E_{\mathrm{nK}}(\alpha))\subseteq
E_{\mathrm{nK}}(\pi^*\alpha)\) by Lemma~\ref{ana:modification-big-nonkahler}, so every K\"ahler
current of \(\pi^*\alpha\) has positive Lelong number at each point of \(W\), and a potential
with positive Lelong number at a point is \(-\infty\) there.
\end{proof}

\begin{lemma}[Step-5 data give the local branch and the collar comparison]
\label{ana:step-five-local-branch}
\uses{ana:step-five-data,ana:gluing-conditions}
\source{transcendental-nakamaye-copy4-20260911:prop:principalization}
\lean{TranscendentalNakamaye.Pending.PrincipalizationCopy4StepFiveCoreReduction.nonempty_anchoredData_of_nonempty_copy4StepFiveExceptionalWeightData,TranscendentalNakamaye.Pending.PrincipalizationSourceDataReduction.PrincipalizationStepSixAnchoredData.exists_localBranch}
\leanok
Let \(Y\) be compact, \(\theta_Y+dd^c\varphi_G\) a K\"ahler current with \(\varphi_G=-\infty\)
on \(W\), and suppose \(\varphi_G\) has Step-5 data. Then there are \(y\in W\), an open
\(U\supseteq W\) and \(\varphi_U\) such that \(\theta_Y+dd^c\varphi_U\) dominates a K\"ahler
form on \(U\), \(\varphi_U\) is smooth near \(y\), and for every compact
\(\Lambda\subseteq U\setminus W\) there is \(C_\Lambda\) with
\(\varphi_U\le\varphi_G+C_\Lambda\) on \(\Lambda\).
\end{lemma}
\begin{proof}
\uses{ana:step-five-data}
\leanok
Take \(y\), \(U\) from the data and \(\varphi_U=\sigma+\rho\). Then
\(\theta_Y+dd^c\varphi_U=(\beta+dd^c\sigma)+(\theta_Y-\beta+dd^c\rho)\) dominates the K\"ahler
form dominated by \(\beta+dd^c\sigma\), and \(\varphi_U\) is smooth near \(y\). For the collar,
\(\sigma\) is bounded on \(\Lambda\), and the comparison \(\psi_D+\chi\le\varphi_G+C_\Lambda\) is
Copy-4's collar computation in Step~5: the ambient comparisons give
\(\varphi_Q\circ\pi\ge\tfrac{1-L/M}M\psi_{\widetilde G}+\tfrac{LA}M(\psi_H+\psi_{\mathfrak j}
\circ\pi)-C_\Lambda\) on \(\Lambda\), because \(\psi_{\mathcal I_W}\) is bounded on the compact
\(\Lambda\) away from \(W\); the weight inequalities and \(C-1\ge LA\), \(1\ge1-L/M\), with all
the weights bounded above, give
\(\psi_{D'}+\chi\le\tfrac{1-L/M}M\psi_{\widetilde G}+\tfrac{LA}M(\psi_H+\psi_{\mathfrak j}\circ
\pi)+O(1)\); add \(\tau\psi_E\) to both sides.
\end{proof}

\begin{proposition}[Copy-4 Steps~3--6 on one choice of data]
\label{ana:step-five-assembly}
\uses{ana:step-five-data,ana:exceptional-metric,ana:polar-branch-step-three,ana:step-five-correction,ana:step-six-local-geometry,ana:ideal-function}
\source{transcendental-nakamaye-copy4-20260911:prop:principalization,transcendental-nakamaye-copy4-20260911:sec:fixed-data,transcendental-nakamaye-copy4-20260911:eq:constants}
\lean{TranscendentalNakamaye.Pending.StepFiveAssemblyOfStepOneStepFour.exists_polarBranch_stepFiveData_of_stepOne_stepFour,TranscendentalNakamaye.Pending.StepFiveAssemblyOfStepOneStepFour.copy4Divisor}
\leanok
Let \((X,\omega)\) be compact K\"ahler, \(\theta\) a smooth closed real \((1,1)\)-form,
\(\pi:Y\to X\) a sequence of blowups with \(Y\) compact and connected, and
\(W\subseteq Y\) a closed connected submanifold of dimension \(k\). Suppose given:
\begin{enumerate}
\item an effective \(\mathbf Q\)-divisor \(E\) with a smooth metric, supported in
\(\operatorname{Exc}\pi\), with \(\pi^*\omega-\Theta_E\) K\"ahler (Step~3);
\item the fixed data: \(\Phi_R\in\operatorname{PSH}(X,\theta-\delta_R\omega)\), functions
\(\psi_{\mathcal I_V}\), \(\psi_{\mathfrak j}\) with analytic singularities of types
\((1,\mathcal I_V)\), \((1,\mathfrak j)\), and the data of Copy-4's lower bound for \(\Phi_R\) with constant
\(A\ge0\);
\item \(c_0\ge0\), a level \(M>0\), \(U_M\in\operatorname{PSH}(X,\theta+\tfrac{c_0}M\omega)\)
with \(U_M\ge\tfrac1M\psi_{\mathfrak B_M}-O(1)\) for a function \(\psi_{\mathfrak B_M}\) with
analytic singularities of type \((1,\mathfrak B_M)\);
\item \(\Psi\) with \(dd^c\Psi\ge-c_\Psi\omega\) and analytic singularities of type
\((1,\mathfrak j\mathfrak b)\), and \(\psi_{\mathfrak b}\) of type \((1,\mathfrak b)\);
\item constants \(0<L<M\), \(C>1\), \(\tau>0\), \(\delta\) with \(c_0<L\delta_R\),
\(LA\le C-1\), \(\tfrac{C-1}M c_\Psi\le\delta\), \(\tau\le(L\delta_R-c_0)/M\);
\item (Step~1) effective \(\mathbf Q\)-divisors \(\widetilde G_M\), \(H^{\mathrm{fix}}\),
\(P^{\mathrm{mov}}\), \(P^F\) on \(Y\) with smooth metrics, such that \(\psi_{\widetilde G_M}\)
has analytic singularities of type \((1,\mathfrak B_M\mathcal O_Y)\),
\(\psi_{H^{\mathrm{fix}}}+\psi_{P^{\mathrm{mov}}}+\psi_{P^F}\) has analytic singularities of
type \((1,\mathfrak h)\) for an ideal with \(\mathcal I_V\mathcal O_Y=\mathfrak h\cdot\mathcal I_W\),
and \(\psi_{\mathfrak b}\circ\pi\le\psi_{H^{\mathrm{fix}}}+O(1)\);
\item a point \(y\in W\) off the supports of \(\widetilde G_M\), \(P^{\mathrm{mov}}\), \(P^F\) and
off \(\operatorname{Exc}\pi\), with \(\Psi\) smooth near \(\pi(y)\);
\item (Step~4) a smooth \(\tilde f\) on \(Y\) with \(\beta+dd^c\tilde f\) strictly positive on the
tangent directions of \(W\), where \(\beta=\pi^*\theta-\Theta_D-\delta\,\pi^*\omega\) and
\[
 D=\tfrac1M\widetilde G_M+\tfrac{C-1}MP^{\mathrm{mov}}+\tfrac CMP^F+\tau E .
\]
\end{enumerate}
Then \(\varphi_G=\varphi_Q\circ\pi+\tau\psi_E\), with
\(\varphi_Q=(1-\tfrac LM)U_M+\tfrac LM\Phi_R\), is the potential of a K\"ahler current in
\(\pi^*\{\theta\}\), and \(\varphi_G\) has Step-5 data.
\end{proposition}
\begin{proof}
\uses{ana:polar-branch-step-three,ana:step-five-correction,ana:step-six-local-geometry}
\leanok
\(\varphi_G\) is a K\"ahler-current potential by Lemma~\ref{ana:polar-branch-step-three}.
Lemma~\ref{ana:step-six-local-geometry}, applied to the K\"ahler form \(\pi^*\omega-\Theta_E\)
and Step~4's positivity, gives \(U\supseteq W\) and \(\sigma\) with \(\beta+dd^c\sigma\)
dominating a K\"ahler form on \(U\). The point \(y\) is off \(\operatorname{Supp}D\), since
\(\operatorname{Supp}E\subseteq\operatorname{Exc}\pi\); so Lemma~\ref{ana:step-five-correction},
with \(\kappa=(C-1)/M\), makes \(\rho=\psi_D+\kappa\,\Psi\circ\pi\) a
\((\pi^*\theta-\beta)\)-plurisubharmonic function on \(U\), smooth near \(y\). The weight
inequalities of the Step-5 data are the exact identities of the weights of \(D\) and of the
Step-1 divisors, the bound of Lemma~\ref{ana:step-five-correction}~(iii), and the upper bounds of
weights of metered divisors and of \(\psi_{\mathfrak j}\) on the compact manifolds; the ambient
comparisons are the hypotheses (ii), (iii) and (vi), together with a weight of the residual
ideal of \(W\).
\end{proof}
The constants \(c_0\), \(A\) and \(c_\Psi\) are inputs here and the inequalities between them
are hypotheses, because Copy-4 chooses \(c_0\) and \(A\) before \(L\) and \(C\). With the
quantitative regularization lemma (Lemma~\ref{ana:regularization}) supplying \(c_0\) and
\(U_M\) first, together with the global lower bound of
Corollary~\ref{ana:regularization-global-lower-bound}, the order can be met.

\begin{lemma}[Step~1's divisors with smooth metrics]
\label{ana:step-one-divisor-weights}
\uses{ana:ideal-sheaves,ana:divisor-metric}
\source{transcendental-nakamaye-copy4-20260911:prop:principalization}
\lean{TranscendentalNakamaye.Pending.StepOneDivisorWeights.stepOneWeights_of_isDivisorMetric,TranscendentalNakamaye.Pending.StepOneDivisorWeights.exists_stepOneWeights_of_divisors}
\leanok
Let \(\pi:Y\to X\) be a holomorphic map of compact complex manifolds, \(\widetilde G_M\),
\(H^{\mathrm{fix}}\), \(P^{\mathrm{mov}}\), \(P^F\) effective divisors on \(Y\) with
\(\mathfrak B_M\mathcal O_Y=\mathcal O_Y(-\widetilde G_M)\), and \(\mathfrak b\) a coherent ideal
on \(X\), \(\psi_{\mathfrak b}\) of analytic singularities of type \((1,\mathfrak b)\), with
\(\mathfrak b\mathcal O_Y\subseteq\mathcal O_Y(-H^{\mathrm{fix}})\). For any smooth metrics on the
four divisors (Lemma~\ref{ana:divisor-metric}), and such metrics exist with the same supports,
\(\psi_{\widetilde G_M}\) has analytic singularities of type
\((1,\mathfrak B_M\mathcal O_Y)\),
\(\psi_{H^{\mathrm{fix}}}+\psi_{P^{\mathrm{mov}}}+\psi_{P^F}\) has analytic singularities of type
\((1,\mathcal O_Y(-H^{\mathrm{fix}}-P^{\mathrm{mov}}-P^F))\), and
\(\psi_{\mathfrak b}\circ\pi\le\psi_{H^{\mathrm{fix}}}+O(1)\).
\end{lemma}
\begin{proof}
\leanok
Smooth metrics exist by Lemma~\ref{ana:divisor-metric}; fix any. The weight of a smooth metric
on an effective divisor has analytic singularities of the type of its ideal
(Lemma~\ref{ana:divisor-metric}), which gives the first two presentations (the ideal of a
sum of divisors is the product of the ideals). For the last, \(\psi_{\mathfrak b}\circ\pi\) has
analytic singularities of type \((1,\mathfrak b\mathcal O_Y)\), and the inclusion
\(\mathfrak b\mathcal O_Y\subseteq\mathcal O_Y(-H^{\mathrm{fix}})\) of coherent ideals on the
compact \(Y\) gives the comparison of weights.
\end{proof}
These are the Step-1 hypotheses (vi) of Proposition~\ref{ana:step-five-assembly}, with
\(\mathfrak h=\mathcal O_Y(-H^{\mathrm{fix}}-P^{\mathrm{mov}}-P^F)\); the factorization
\(\mathcal I_V\mathcal O_Y=\mathfrak h\cdot\mathcal I_W\) is Copy-4's identity for the divisor
\(H\) of Step~2.

\begin{proposition}[Copy-4 Steps~3--6 from the Step-1 data and the Step-4 current]
\label{ana:step-five-from-step-one-current}
\uses{ana:step-five-assembly,ana:step-four-on-w,ana:divisor-restriction,ana:step-one-divisor-weights}
\source{transcendental-nakamaye-copy4-20260911:prop:principalization}
\lean{TranscendentalNakamaye.Pending.StepFiveAssemblyOfStepOneCurrent.exists_polarBranch_stepFiveData_of_stepOne_stepFourCurrent}
\leanok
Under the hypotheses (i)--(vii) of Proposition~\ref{ana:step-five-assembly}, suppose further
that no component of the divisor \(D\) of that proposition contains \(W\), and that \(\varphi'\)
is a potential on \(W\) of a K\"ahler current in the class of
\((\pi^*\theta-\delta\,\pi^*\omega)|_W\) with logarithmic singularities along \(D|_W\). Then
\(\varphi_G=\varphi_Q\circ\pi+\tau\psi_E\) is the potential of a K\"ahler current in
\(\pi^*\{\theta\}\) and has Step-5 data.
\end{proposition}
\begin{proof}
\uses{ana:step-five-assembly,ana:step-four-on-w}
\leanok
Lemma~\ref{ana:step-four-on-w}, with \(\theta_Y=\pi^*\theta\) and \(\omega_Y=\pi^*\omega\),
supplies the hypothesis (viii) of Proposition~\ref{ana:step-five-assembly}, which then applies.
\end{proof}

\begin{lemma}[Supports of the Step-4 divisor]
\label{ana:copy4-divisor-support}
\uses{ana:ideal-sheaves,ana:step-five-assembly}
\source{transcendental-nakamaye-copy4-20260911:prop:principalization}
\lean{TranscendentalNakamaye.Pending.StepFiveAssemblyOfStepOneStepFour.copy4Divisor_support_subset,TranscendentalNakamaye.Pending.StepSixGenericPoint.notMem_support_of_smul_weight_ne_bot,TranscendentalNakamaye.Pending.StepSixGenericPoint.notMem_support_of_copy4Divisor_weight_ne_bot}
\leanok
Let \(\widetilde G_M\), \(P^{\mathrm{mov}}\), \(P^F\), \(E\) be effective \(\mathbf Q\)-divisors
with smooth metrics, \(M>0\), \(C\ge1\), \(\tau\ge0\), and
\(D=\tfrac1M\widetilde G_M+\tfrac{C-1}MP^{\mathrm{mov}}+\tfrac CMP^F+\tau E\) with the induced
metric. Then \(\operatorname{Supp}D\) is contained in the union of the four supports. If
\(C>1\) and \(\psi_D(y)>-\infty\), then \(y\) lies off \(\operatorname{Supp}\widetilde G_M\),
\(\operatorname{Supp}P^{\mathrm{mov}}\) and \(\operatorname{Supp}P^F\).
\end{lemma}
\begin{proof}
\leanok
\(\psi_D=\tfrac1M\psi_{\widetilde G_M}+\tfrac{C-1}M\psi_{P^{\mathrm{mov}}}+\tfrac CM\psi_{P^F}
+\tau\psi_E\), a sum of nonnegative multiples of functions with values in
\([-\infty,\infty)\). It is \(-\infty\) only where some term is, so only on one of the four
supports. If \(C>1\), the first three coefficients are positive, so a point of one of the first
three supports makes the corresponding term, and hence \(\psi_D\), equal to \(-\infty\).
\end{proof}

\begin{lemma}[No component of the Step-4 divisor contains \(W\)]
\label{ana:copy4-divisor-no-component}
\uses{ana:copy4-divisor-support,ana:strict-transform-dense-over-complement,ana:zero-locus-coherent,ana:divisor-metric}
\source{transcendental-nakamaye-copy4-20260911:prop:principalization,transcendental-nakamaye-copy4-20260911:sec:positive-current}
\lean{TranscendentalNakamaye.Pending.StepFourNoComponentContainsW.copy4Divisor_noComponent,TranscendentalNakamaye.Pending.StepFourNoComponentContainsW.copy4Divisor_noComponent_of_effectiveDivisors}
\leanok
Let \(X\) be a compact complex manifold, \(\pi:Y\to X\) a finite composition of blowups with
smooth centres, \(V\subset X\) a closed irreducible analytic subvariety with strict transform
\(W\), and \(D\) as in Lemma~\ref{ana:copy4-divisor-support}. Suppose that
\(\operatorname{Supp}P^{\mathrm{mov}}\), \(\operatorname{Supp}P^F\) and
\(\operatorname{Supp}E\) lie in \(\operatorname{Exc}\pi\), and that
\(\operatorname{Supp}\widetilde G_M\subset\pi^{-1}(Z)\) for an analytic subset \(Z\) with
\(V\not\subset Z\). Then \(\psi_D\) is identically \(-\infty\) on no nonempty relatively open
subset of \(W\). This holds in particular when \(\widetilde G_M\) is a smooth metric on the
divisor of \(\mathfrak B_M\mathcal O_Y\) for a coherent ideal \(\mathfrak B_M\) on \(X\) with
\(V\not\subset Z(\mathfrak B_M)\), and \(P^{\mathrm{mov}}\), \(P^F\) are smooth metrics on
effective divisors supported in \(\operatorname{Exc}\pi\).
\end{lemma}
\begin{proof}
\leanok
A nonempty relatively open subset of \(W\) is \(O\cap W\) with \(O\subset Y\) open and meeting
\(W\). Lemma~\ref{ana:strict-transform-dense-over-complement} gives \(p\in O\) with
\(p\notin\operatorname{Exc}\pi\) and \(\pi(p)\in V\setminus Z\); then \(p\in W\), as a point of
\(\pi^{-1}(V)\setminus\operatorname{Exc}\pi\). The point \(p\) is off
\(\operatorname{Supp}P^{\mathrm{mov}}\), \(\operatorname{Supp}P^F\), \(\operatorname{Supp}E\)
(all in \(\operatorname{Exc}\pi\)) and off \(\operatorname{Supp}\widetilde G_M\subset\pi^{-1}(Z)\),
so \(\psi_D(p)>-\infty\) by Lemma~\ref{ana:copy4-divisor-support}.

For the special case: \(\operatorname{Supp}\widetilde G_M\) is the support of the divisor of
\(\mathfrak B_M\mathcal O_Y\) (Lemma~\ref{ana:divisor-metric}), that is
\(Z(\mathfrak B_M\mathcal O_Y)=\pi^{-1}Z(\mathfrak B_M)\), and \(Z(\mathfrak B_M)\) is analytic
(Lemma~\ref{ana:zero-locus-coherent}); the supports of \(P^{\mathrm{mov}}\), \(P^F\) are those
of the underlying divisors (Lemma~\ref{ana:divisor-metric}).
\end{proof}
Since \(W\) is irreducible and \(\operatorname{Supp}D\cap W\) is analytic in \(W\), the
conclusion says exactly that no component of \(D\) contains \(W\), the form in which Copy-4
states it in Step~4. The hypotheses are Copy-4's facts that \(P^{\mathrm{mov}}\), \(P^F\) and
\(E_{\mathrm{rel}}\) are exceptional and that \(\mathfrak B_M\) is the unit ideal near a point of
\(V\) (its section on a positive current on the submanifold).

\begin{proposition}[Copy-4 Steps~3--6 with the point \(y\) derived]
\label{ana:step-five-generic-point}
\uses{ana:step-five-from-step-one-current,ana:strict-transform-generic-point,ana:copy4-divisor-support}
\source{transcendental-nakamaye-copy4-20260911:prop:principalization}
\lean{TranscendentalNakamaye.Pending.StepSixGenericPoint.exists_polarBranch_stepFiveData_of_stepOne_stepFourCurrent_strictTransform}
\leanok
Let \(V\subset X\) be a closed irreducible analytic subvariety and let \(W\) be its strict
transform under \(\pi\), assumed to be a closed connected submanifold of dimension \(k\). Under
the hypotheses (i)--(vi) of Proposition~\ref{ana:step-five-assembly} and the two further
hypotheses of Proposition~\ref{ana:step-five-from-step-one-current} (no component of \(D\)
contains \(W\); the potential \(\varphi'\) on \(W\)), with the hypothesis (vii) on the point
\(y\) replaced by: \emph{there is \(x_0\in V\) near which \(\Psi\) is smooth}, the conclusion
of Proposition~\ref{ana:step-five-from-step-one-current} holds.
\end{proposition}
\begin{proof}
\uses{ana:step-five-from-step-one-current,ana:strict-transform-generic-point,ana:copy4-divisor-support}
\leanok
Let \(U\ni x_0\) be open with \(\Psi\) smooth on \(U\). By
Lemma~\ref{ana:strict-transform-generic-point}~(ii) there is \(y_1\in W\setminus\operatorname{Exc}\pi\)
with \(\pi(y_1)\in U\), so \(W\cap(\pi^{-1}(U)\setminus\operatorname{Exc}\pi)\) is a nonempty
relatively open subset of \(W\). As no component of \(D\) contains \(W\), it contains a point
\(y\) with \(\psi_D(y)>-\infty\), which lies off \(\operatorname{Supp}\widetilde G_M\),
\(\operatorname{Supp}P^{\mathrm{mov}}\), \(\operatorname{Supp}P^F\)
(Lemma~\ref{ana:copy4-divisor-support}), off \(\operatorname{Exc}\pi\), and has \(\Psi\) smooth
near \(\pi(y)\). This is hypothesis (vii), and
Proposition~\ref{ana:step-five-from-step-one-current} applies.
\end{proof}
This is Copy-4's Step~6 choice: ``we may choose \(y\in W\) outside \(\operatorname{Supp}D\) near
which \(\chi\) is smooth''. A point \(x_0\in V\) near which \(\Psi\) is smooth is what
Lemma~\ref{ana:ideal-function} provides, for the ideal \(\mathfrak j\mathfrak b\), at a point of \(V\) off
\(Z(\mathfrak j\mathfrak b)\).

\begin{proposition}[Copy-4 Steps~3--6 from Step~1's effective divisors]
\label{ana:step-five-effective-divisors}
\uses{ana:step-five-generic-point,ana:step-one-divisor-weights,ana:divisor-metric}
\source{transcendental-nakamaye-copy4-20260911:prop:principalization}
\lean{TranscendentalNakamaye.Pending.StepFiveAssemblyOfEffectiveDivisors.exists_polarBranch_stepFiveData_of_effectiveDivisors}
\leanok
Proposition~\ref{ana:step-five-generic-point} holds with its hypothesis (vi) in Copy-4's
Step-1 form: effective divisors \(\widetilde G_M\), \(H^{\mathrm{fix}}\), \(P^{\mathrm{mov}}\),
\(P^F\) on \(Y\) with
\[
 \mathfrak B_M\mathcal O_Y=\mathcal O_Y(-\widetilde G_M),\qquad
 \mathfrak b\,\mathcal O_Y\subseteq\mathcal O_Y(-H^{\mathrm{fix}}),\qquad
 \mathcal I_V\mathcal O_Y=\mathcal O_Y(-H^{\mathrm{fix}}-P^{\mathrm{mov}}-P^F)\cdot\mathcal I_W,
\]
together with smooth metrics on them (Lemma~\ref{ana:divisor-metric}); the divisor \(D\) and
the logarithmic singularities of \(\varphi'\) along \(D|_W\) refer to these metrics.
\end{proposition}
\begin{proof}
\uses{ana:step-five-generic-point,ana:step-one-divisor-weights}
\leanok
For the given metrics, Lemma~\ref{ana:step-one-divisor-weights} supplies the three weight
relations of hypothesis (vi), with
\(\mathfrak h=\mathcal O_Y(-H^{\mathrm{fix}}-P^{\mathrm{mov}}-P^F)\), and the third ideal
identity is the factorization \(\mathcal I_V\mathcal O_Y=\mathfrak h\cdot\mathcal I_W\). Apply
Proposition~\ref{ana:step-five-generic-point}.
\end{proof}

\begin{proposition}[Copy-4 Steps~3--6 from Step~1, with \(W\not\subset\operatorname{Supp}D\) derived]
\label{ana:step-five-exceptional-divisors}
\uses{ana:step-five-effective-divisors,ana:copy4-divisor-no-component}
\source{transcendental-nakamaye-copy4-20260911:prop:principalization,transcendental-nakamaye-copy4-20260911:sec:positive-current}
\lean{TranscendentalNakamaye.Pending.StepFourNoComponentContainsW.exists_polarBranch_stepFiveData_of_effectiveDivisors_of_exceptional}
\leanok
Proposition~\ref{ana:step-five-effective-divisors} holds with the hypothesis that no component
of \(D\) contains \(W\) replaced by: \(P^{\mathrm{mov}}\) and \(P^F\) are supported in
\(\operatorname{Exc}\pi\), and \(V\not\subset Z(\mathfrak B_M)\).
\end{proposition}
\begin{proof}
\uses{ana:step-five-effective-divisors,ana:copy4-divisor-no-component}
\leanok
\(\operatorname{Supp}E\subset\operatorname{Exc}\pi\) by hypothesis (i), and \(\mathfrak B_M\) is
coherent since \(\psi_{\mathfrak B_M}\) has analytic singularities of its type. So
Lemma~\ref{ana:copy4-divisor-no-component} shows that no component of \(D\) contains \(W\), and
Proposition~\ref{ana:step-five-effective-divisors} applies.
\end{proof}

\begin{definition}[Copy-4 Steps~1--4 on a modification]
\label{ana:copy4-step-one-data}
\uses{ana:step-five-assembly,ana:divisor-metric,ana:copy4-divisor-no-component}
\source{transcendental-nakamaye-copy4-20260911:prop:principalization,transcendental-nakamaye-copy4-20260911:sec:fixed-data,transcendental-nakamaye-copy4-20260911:sec:positive-current}
\lean{TranscendentalNakamaye.Pending.PrincipalizationStepOneInterface.Copy4StepOneData}
\leanok
Let \((X,\omega)\) be compact K\"ahler, \(\theta\) a smooth closed real \((1,1)\)-form,
\(V\subset X\) a closed irreducible analytic subvariety, and \(\pi:Y\to X\) a sequence of
blowups whose strict transform \(W\) of \(V\) is a closed connected submanifold of dimension
\(k\). \emph{Step-1 data} on \(\pi\) consist of:
\begin{enumerate}
\item an effective \(\mathbf Q\)-divisor \(E\) with a smooth metric, supported in
\(\operatorname{Exc}\pi\), with \(\pi^*\omega-\Theta_E\) K\"ahler (Step~3);
\item the fixed data of hypothesis (ii) of Proposition~\ref{ana:step-five-assembly}
(\(\Phi_R\), \(\psi_{\mathcal I_V}\), \(\psi_{\mathfrak j}\), the constant \(A\ge0\));
\item \(c_0\ge0\), a rational level \(M>0\), \(U_M\in\operatorname{PSH}(X,\theta+\tfrac{c_0}M\omega)\)
with \(U_M\ge\tfrac1M\psi_{\mathfrak B_M}-O(1)\) for \(\psi_{\mathfrak B_M}\) of type
\((1,\mathfrak B_M)\), and \(V\not\subset Z(\mathfrak B_M)\);
\item the ideal function \(\Psi\) and \(\psi_{\mathfrak b}\) of hypothesis (iv), and the
constants \(L\), \(C\), \(\tau\), \(\delta\) of hypothesis (v), of
Proposition~\ref{ana:step-five-assembly};
\item effective divisors \(\widetilde G_M\), \(H^{\mathrm{fix}}\), \(P^{\mathrm{mov}}\), \(P^F\) on
\(Y\) with \(\mathfrak B_M\mathcal O_Y=\mathcal O_Y(-\widetilde G_M)\),
\(\mathfrak b\mathcal O_Y\subseteq\mathcal O_Y(-H^{\mathrm{fix}})\),
\(\mathcal I_V\mathcal O_Y=\mathcal O_Y(-H^{\mathrm{fix}}-P^{\mathrm{mov}}-P^F)\cdot\mathcal I_W\),
\(P^{\mathrm{mov}}\) and \(P^F\) supported in \(\operatorname{Exc}\pi\), and smooth metrics on
the four divisors (Lemma~\ref{ana:divisor-metric});
\item a point \(x_0\in V\) near which \(\Psi\) is smooth;
\item a potential \(\varphi'\) on \(W\) of a K\"ahler current in the class of
\((\pi^*\theta-\delta\,\pi^*\omega)|_W\) with logarithmic singularities along \(D|_W\), where
\(D=\tfrac1M\widetilde G_M+\tfrac{C-1}MP^{\mathrm{mov}}+\tfrac CMP^F+\tau E\) (no component of
\(D\) contains \(W\), by Lemma~\ref{ana:copy4-divisor-no-component}).
\end{enumerate}
These are exactly the hypotheses of Proposition~\ref{ana:step-five-exceptional-divisors}.
\end{definition}

\begin{proposition}[Copy-4 Steps~3--6 from the Step-1 data]
\label{ana:copy4-steps-three-to-six}
\uses{ana:copy4-step-one-data,ana:step-five-exceptional-divisors}
\source{transcendental-nakamaye-copy4-20260911:prop:principalization}
\lean{TranscendentalNakamaye.Pending.PrincipalizationStepOneInterface.Copy4StepOneData.exists_polarBranch_stepFiveData}
\leanok
If, in the setting of Definition~\ref{ana:copy4-step-one-data}, \(Y\) is compact and connected
and \(\pi\) carries Step-1 data, then
\(\varphi_G=\varphi_{Q_M}\circ\pi+\tau\psi_E\) is the potential of a K\"ahler current in
\(\pi^*\{\theta\}\) and has Step-5 data (Definition~\ref{ana:step-five-data}) for
\(\theta_Y=\pi^*\theta\) and \(W\).
\end{proposition}
\begin{proof}
\uses{ana:step-five-exceptional-divisors}
\leanok
The Step-1 data are the hypotheses of Proposition~\ref{ana:step-five-exceptional-divisors},
which gives the conclusion.
\end{proof}

\begin{proposition}[Copy-4 Step~1 and the data of Steps~1--4]
\label{ana:copy4-steps-one-to-five}
\uses{ana:copy4-step-one-data,ana:cartan-coherence,ana:copy4-positive-current-input,ana:copy4-trace-input-positive-current,ana:copy4-other-components-ideal,ana:copy4-inputs,ana:copy4-step-one-from-inputs,ana:copy4-ambient-input,ana:copy4-trace-input-at-point,ana:copy4-fixed-data,ana:copy4-step-one-from-fixed-data,ana:copy4-fixed-data-fields,ana:relative-principalization,ana:regularization,ana:ideal-function,ana:step-four-kahler-class-on-w,trace:qe-existence}
\source{transcendental-nakamaye-copy4-20260911:prop:principalization,transcendental-nakamaye-copy4-20260911:sec:positive-current,transcendental-nakamaye-copy4-20260911:sec:fixed-data}
\lean{TranscendentalNakamaye.Pending.PrincipalizationStepOneInterface.ct22Copy4_exists_stepOneData}
\leanok
Let \(X\) be a connected compact K\"ahler manifold of dimension \(n\), \(\alpha=\{\theta\}\) big,
and \(V\subset X\) a closed connected complex submanifold of dimension \(0<k<n\) which is an
irreducible component of \(E_{\mathrm{nK}}(\alpha)\), is not contained in
\(E_{\mathrm{nn}}(\alpha)\), and has \(\operatorname{RV}_{X\mid V}(\alpha)>0\). Assume relative
principalization (Lemma~\ref{ana:relative-principalization}) on compact connected complex
manifolds of dimension \(n\). Then some sequence of blowups \(\pi:Y\to X\) adapted to \(V\),
whose strict transform \(W\) is a closed connected submanifold of dimension \(k\), carries
Step-1 data (Definition~\ref{ana:copy4-step-one-data}).
\end{proposition}
\begin{proof}
\uses{ana:copy4-other-components-ideal,ana:cartan-coherence,ana:copy4-positive-current-input,ana:closure-difference-analytic}
\leanok
The set \(Z'=\overline{E_{\mathrm{nK}}(\alpha)\setminus V}\) is analytic (Theorem~\ref{ana:nonkahler-locus-analytic},
Proposition~\ref{ana:closure-difference-analytic}), so \(\mathcal I_{Z'}\) is coherent by Theorem~\ref{ana:cartan-coherence};
the data of (I3) for \(\varphi=V_\theta\) exist by Proposition~\ref{ana:copy4-positive-current-input}. The last assertion of
Proposition~\ref{ana:copy4-other-components-ideal} gives the conclusion.
\end{proof}
The proposition rests on Copy-4's positive current on \(V\)
(Proposition~\ref{ana:copy4-positive-current-input}), on Cartan's coherence theorem (Theorem~\ref{ana:cartan-coherence}), and on
relative principalization, a hypothesis. Everything Copy-4 builds on the modification is proved: given relative
principalization, the data of Definition~\ref{ana:copy4-fixed-data} produce \(\pi\) together with its Step-1 data
(Proposition~\ref{ana:copy4-step-one-from-fixed-data}) --- the three relative principalizations of Step~1 with the excess ideal
(Proposition~\ref{ana:copy4-step-one-excess}), the metrics of Step~3, and the potential \(\varphi'\) of Step~4 with its
logarithmic singularities along \(D|_W\). Those data are in turn produced, with Copy-4's constants chosen in its order and
\(U_M\), \(\Psi\) and \(x_0\) constructed (Lemma~\ref{ana:copy4-fixed-data-fields}), from three inputs on \(X\) and \(V\)
(Proposition~\ref{ana:copy4-step-one-from-inputs}): (I1)--(I3) of Definition~\ref{ana:copy4-inputs}, which
Propositions~\ref{ana:copy4-trace-input-positive-current} and~\ref{ana:copy4-other-components-ideal} reduce to the coherence of
\(\mathcal I_{Z'}\) for \(Z'=\overline{E_{\mathrm{nK}}(\alpha)\setminus V}\) (Cartan's theorem) and to (I3) for \(\varphi=V_\theta\)
(Proposition~\ref{ana:copy4-positive-current-input}); the other inputs are proved: the fixed data \(\Phi_R\), \(\mathfrak j\), \(A\) (proved from bigness, given a coherent ideal \(\mathfrak j\) cutting out the other
components of \(E_{\mathrm{nK}}(\alpha)\), in Proposition~\ref{ana:copy4-ambient-input}); Skoda's theorem (proved,
Proposition~\ref{ana:copy4-trace-input-at-point}: one point of \(V\) where \(\nu(V_\theta,\cdot)=0\) suffices) and subadditivity for the
multiplier ideals \(\mathfrak B_q=\mathcal I(qV_\theta)\) (proved, Theorem~\ref{dem:del00-subadditivity} and
Proposition~\ref{ana:copy4-trace-input-no-sub}); and, from Copy-4's section on a positive current on the submanifold (the
quasi-equisingular approximation of \(V_\theta\)), the potentials \(\varphi_{K,\ell}\) on \(V\) with ambient coherent ideals
\(\widetilde{\mathfrak B}\) restricting to \(\mathfrak B=\mathcal I(\ell\varphi_K)\) on \(V\) and the comparison of
Lemma~\ref{ana:ideal-comparison}. Relative principalization is a
hypothesis of the statement and is supplied on the route by Lemma~\ref{ana:relative-principalization}, so the
resolution-of-singularities content stays in that lemma. As in Copy-4, \(\pi\) is existential and \(\varphi_G\) is produced
together with its data (Proposition~\ref{ana:copy4-steps-three-to-six}).

\begin{lemma}[One stage of Copy-4's Step~1]
\label{ana:principalization-stage}
\uses{ana:blowups,ana:relative-principalization,ana:strict-transform-composition,ana:ideal-sheaves}
\source{transcendental-nakamaye-copy4-20260911:prop:principalization,transcendental-nakamaye-copy4-20260911:lem:relative-principalization}
\lean{TranscendentalNakamaye.Pending.PrincipalizationStepOneStages.isAdaptedBlowupSequence_comp,TranscendentalNakamaye.Pending.PrincipalizationStepOneStages.exists_principalizationStage,TranscendentalNakamaye.Pending.PrincipalizationStepOneStages.exists_vanishingIdeal_stage}
\leanok
Let \(X\) be a complex manifold, \(V\subset X\), and \(\pi:Y\to X\) a sequence of blowups adapted
to \(V\), with strict transform \(W=\widetilde V\).
\begin{enumerate}
\item If \(f:Y'\to Y\) is a sequence of blowups adapted to \(W\), then \(\pi\circ f\) is adapted
to \(V\), and the strict transform of \(V\) under \(\pi\circ f\) is that of \(W\) under \(f\).
\item Assume relative principalization (Lemma~\ref{ana:relative-principalization}) for compact
connected manifolds of dimension \(n=\dim X\). If \(Y\) is compact and connected, \(W\ne Y\) is a
closed connected submanifold of dimension \(k\), and \(\mathfrak a\) is a coherent ideal on \(Y\)
not contained in \(\mathcal I_W\), there is a sequence of blowups \(f:Y'\to Y\) adapted to \(W\)
with \(Y'\) compact and connected, such that \(\pi\circ f\) is adapted to \(V\), the strict
transform \(W'\) of \(V\) is a closed connected submanifold of dimension \(k\), and the
conclusions of Lemma~\ref{ana:relative-principalization} hold for \(f\), \(W\), \(\mathfrak a\):
effective divisors \(H\) (exceptional for \(f\)) and \(H'\) (disjoint from \(W'\)) with
\(\mathcal I_W\mathcal O_{Y'}=\mathcal O_{Y'}(-H)\mathcal I_{W'}\),
\(\mathfrak a\mathcal O_{Y'}=\mathcal O_{Y'}(-H-H')\), and the restriction identity on \(W'\).
\item If \(V\), \(W\) and \(W'\) are closed connected submanifolds,
\(\mathcal I_V\mathcal O_Y=\mathcal O_Y(-H_0)\mathcal I_W\) with \(H_0\) exceptional for \(\pi\), and
\(\mathcal I_W\mathcal O_{Y'}=\mathcal O_{Y'}(-H)\mathcal I_{W'}\) with \(H\) exceptional for \(f\),
then \(H_{\mathrm{tot}}=f^*H_0+H\) is exceptional for \(\pi\circ f\) and
\(\mathcal I_V\mathcal O_{Y'}=\mathcal O_{Y'}(-H_{\mathrm{tot}})\mathcal I_{W'}\).
\end{enumerate}
\end{lemma}
\begin{proof}
\uses{ana:relative-principalization,ana:strict-transform-composition}
\leanok
(i) Induction on the number of blowups of \(f\): the strict transform along \(\pi\circ g\) is the
strict transform along \(g\) of \(W\) (Lemma~\ref{ana:strict-transform-composition}), so each
centre of \(f\), adapted to the current strict transform of \(W\), is adapted to the current strict
transform of \(V\). (ii) Apply Lemma~\ref{ana:relative-principalization} on \(Y\) to \(W\) and
\(\mathfrak a\) (its restriction to \(W\) is nonzero because \(\mathfrak a\not\subset\mathcal I_W\)),
and use (i). (iii) Pull back the first identity along \(f\):
\(\mathcal I_V\mathcal O_{Y'}=(\mathcal O_Y(-H_0)\mathcal I_W)\mathcal O_{Y'}
=\mathcal O_{Y'}(-f^*H_0)\,\mathcal O_{Y'}(-H)\mathcal I_{W'}\). The support of \(f^*H_0\) is
\(f^{-1}(\operatorname{Supp}H_0)\subset f^{-1}(\operatorname{Exc}\pi)\), and
\(\operatorname{Exc}(\pi\circ f)=\operatorname{Exc}f\cup f^{-1}(\operatorname{Exc}\pi)\).
\end{proof}
These are the stages \(Y\to Y_1\to Y_{\mathrm{fix}}\to X\) of Copy-4's Step~1, independently of
which ideal is principalized at each stage, and the identity \(\mathcal I_V\mathcal O_Y=
\mathcal O_Y(-H)\mathcal I_W\) of its Step~2 across them. The ideals themselves come from the
quasi-equisingular approximation, so Proposition~\ref{ana:copy4-steps-one-to-five} stays a named
input.

\begin{lemma}[Dividing a coherent ideal by an invertible one]
\label{ana:invertible-colon}
\uses{ana:ideal-sheaves}
\source{transcendental-nakamaye-copy4-20260911:prop:principalization}
\lean{TranscendentalNakamaye.Pending.InvertibleIdealQuotient.exists_isCoherent_eq_mul,TranscendentalNakamaye.Pending.InvertibleIdealQuotient.exists_isCoherent_eq_mul_divisor}
\leanok
If \(\mathfrak a\) is a coherent ideal contained in an invertible ideal \(\mathfrak l\), then
\(\mathfrak a=\mathfrak l\cdot\mathfrak c\) with \(\mathfrak c=(\mathfrak a:\mathfrak l)\) coherent.
\end{lemma}
\begin{proof}
\leanok
Near a point \(x\) let \(s_1,\dots,s_N\) generate \(\mathfrak a\) and \(h\) generate \(\mathfrak l\).
Since \(\mathfrak a\subset\mathfrak l\), \(s_i=c_ih\) for germs \(c_i\) at \(x\), holomorphic on a common
neighbourhood \(U\). The \(c_i\) generate \((\mathfrak a:\mathfrak l)\) on \(U\): if \(gh\in\mathfrak a\)
then \(gh=\sum a_is_i=(\sum a_ic_i)h\), and \(h\) vanishes identically on no open set, so
\(g=\sum a_ic_i\) on the dense set \(\{h\ne0\}\), hence everywhere by continuity. Then
\(\mathfrak l\cdot(\mathfrak a:\mathfrak l)\subset\mathfrak a\) is clear, and \(s_i=h\,c_i\) gives the
reverse inclusion.
\end{proof}

\begin{lemma}[Copy-4's Step~1: the three-stage modification]
\label{ana:copy4-step-one-stages}
\uses{ana:principalization-stage,ana:relative-principalization,ana:invertible-colon,ana:strict-transform-dense-over-complement,ana:zero-locus-coherent,ana:blowups}
\source{transcendental-nakamaye-copy4-20260911:prop:principalization,transcendental-nakamaye-copy4-20260911:sec:fixed-principalization}
\lean{TranscendentalNakamaye.Pending.PrincipalizationStepOneAssembly.exists_firstTwoStages,TranscendentalNakamaye.Pending.PrincipalizationStepOneAssembly.exists_thirdStage,TranscendentalNakamaye.Pending.PrincipalizationStepOneAssembly.stepOneFields}
\leanok
Assume relative principalization (Lemma~\ref{ana:relative-principalization}) for compact connected
manifolds of dimension \(n\). Let \(X\) be a compact connected complex manifold of dimension \(n\),
\(V\subset X\) a closed irreducible analytic subvariety which is a closed connected submanifold of
dimension \(k<n\), and \(\mathfrak B_m,\mathfrak B_M\) coherent ideals with
\(\mathfrak B_M\subset\mathfrak B_m^q\), \(V\not\subset Z(\mathfrak B_m)\) and
\(V\not\subset Z(\mathfrak B_M)\).
\begin{enumerate}
\item (Stages 1 and 2.) There is a sequence of blowups \(\pi:Y\to X\) adapted to \(V\), with \(Y\)
compact and connected and strict transform \(W\) a closed connected submanifold of dimension \(k\),
and effective divisors \(H^{\mathrm{fix}},P^{\mathrm{mov}}\) supported in \(\operatorname{Exc}\pi\) and
\(H_m,P_M\) disjoint from \(W\), such that, with \(\widetilde G_m=H^{\mathrm{fix}}+H_m\),
\[
 \mathcal I_V\mathcal O_Y=\mathcal O_Y(-H^{\mathrm{fix}}-P^{\mathrm{mov}})\,\mathcal I_W,\quad
 \mathfrak B_m\mathcal O_Y=\mathcal O_Y(-\widetilde G_m),\quad
 \mathfrak B_M\mathcal O_Y=\mathcal O_Y(-q\widetilde G_m)\,\mathcal O_Y(-P^{\mathrm{mov}}-P_M).
\]
\item (Stage 3 and Step 2.) Given such an output and a coherent ideal \(\mathfrak f\) on \(Y\) not
contained in \(\mathcal I_W\), there is \(g:Y'\to Y\) adapted to \(W\), with \(Y'\) compact and
connected and \(\pi\circ g\) adapted to \(V\), such that the three identities of (i) hold on \(Y'\) for
the total transforms of the four divisors, \(\mathfrak f\mathcal O_{Y'}=\mathcal O_{Y'}(-P^F-P^{F\prime})\)
with \(P^F\) supported in \(\operatorname{Exc}(\pi\circ g)\) and \(P^{F\prime}\) disjoint from the strict
transform \(W'\), and
\(\mathcal I_V\mathcal O_{Y'}=\mathcal O_{Y'}(-H^{\mathrm{fix}}-P^{\mathrm{mov}}-P^F)\,\mathcal I_{W'}\).
\item (The Step-1 fields.) On such an output, \(\widetilde G_M\) with
\(\mathcal O(-\widetilde G_M)=\mathfrak B_M\mathcal O\) is an effective divisor, and
\(\mathfrak B_m\mathcal O\subset\mathcal O(-H^{\mathrm{fix}})\).
\end{enumerate}
\end{lemma}
\begin{proof}
\uses{ana:principalization-stage,ana:invertible-colon,ana:strict-transform-dense-over-complement}
\leanok
(i) Stage~1: relative principalization of \(\mathfrak B_m\) along \(V\) (its restriction to \(V\) is
nonzero because \(V\not\subset Z(\mathfrak B_m)\)) gives \(f_1:Y_1\to X\) with
\(\mathcal I_V\mathcal O_{Y_1}=\mathcal O(-H_1)\mathcal I_{W_1}\) and
\(\mathfrak B_m\mathcal O_{Y_1}=\mathcal O(-H_1-H_1')\), \(H_1\) exceptional, \(H_1'\) disjoint from
\(W_1\). From \(\mathfrak B_M\subset\mathfrak B_m^q\),
\(\mathfrak B_M\mathcal O_{Y_1}\subset\mathcal O(-q(H_1+H_1'))\), an invertible ideal, so
\(\mathfrak B_M\mathcal O_{Y_1}=\mathcal O(-q(H_1+H_1'))\,\mathfrak c_M\) with \(\mathfrak c_M\) coherent
(Lemma~\ref{ana:invertible-colon}). \(\mathfrak c_M\not\subset\mathcal I_{W_1}\): otherwise
\(W_1\subset Z(\mathfrak B_M\mathcal O_{Y_1})=f_1^{-1}(Z(\mathfrak B_M))\), but
Lemma~\ref{ana:strict-transform-dense-over-complement} gives points of \(W_1\) off
\(\operatorname{Exc}f_1\) mapping into \(V\setminus Z(\mathfrak B_M)\). Stage~2: Lemma~\ref{ana:principalization-stage}~(ii)
for \(\mathfrak c_M\) along \(W_1\) gives \(f_2\) with
\(\mathcal I_{W_1}\mathcal O=\mathcal O(-P)\mathcal I_W\), \(\mathfrak c_M\mathcal O=\mathcal O(-P-P')\).
Put \(\pi=f_1\circ f_2\), \(H^{\mathrm{fix}}=f_2^*H_1\), \(H_m=f_2^*H_1'\), \(P^{\mathrm{mov}}=P\),
\(P_M=P'\). The first identity is Lemma~\ref{ana:principalization-stage}~(iii), the other two are the
pullbacks of the stage-1 identities along \(f_2\); \(f_2^*H_1'\) is disjoint from \(W\) because \(W\) lies
in \(f_2^{-1}(W_1)\), and \(f_2^*H_1\) and \(P\) are exceptional for \(\pi\) because
\(\operatorname{Exc}\pi=\operatorname{Exc}f_2\cup f_2^{-1}(\operatorname{Exc}f_1)\).
(ii) The same with Lemma~\ref{ana:principalization-stage}~(ii)--(iii) for \(\mathfrak f\) along \(W\),
transporting the identities of (i) by pullback. (iii) \(\mathfrak B_M\mathcal O\) is a product of
invertible ideals, hence invertible; and
\(\mathfrak B_m\mathcal O=\mathcal O(-H^{\mathrm{fix}})\mathcal O(-H_m)\subset\mathcal O(-H^{\mathrm{fix}})\).
\end{proof}
The remaining fields of Copy-4's Step-1 data (the ideals \(\mathfrak B_q=\mathcal I(qV_\theta)\) and
their coherence, the regularization \(U_M\), the potential \(\varphi_{K,\ell}\) of Step~4) come from the
positive current on \(V\) (Proposition~\ref{ana:copy4-positive-current-input}); the excess
ideal \(\widetilde{\mathfrak f}_B\) and the restriction identity of Step~4 are constructed from them below
(Proposition~\ref{ana:copy4-step-one-excess}), and all of Steps~1--4 from the data on \(X\) and \(V\) in
Proposition~\ref{ana:copy4-step-one-from-fixed-data}.

\begin{definition}[Restricted inclusions of ideal sheaves]
\label{ana:restricted-ideals}
\uses{ana:ideal-sheaves}
\source{transcendental-nakamaye-copy4-20260911:prop:principalization}
\lean{TranscendentalNakamaye.Pending.RestrictedIdealWeight.RestrictLe,TranscendentalNakamaye.Pending.RestrictedIdealExtension.restrictOn,TranscendentalNakamaye.Pending.RestrictedIdealQuotient.IsInvertibleAlong}
\leanok
Let \(W\subseteq M\) and let \(\mathfrak a,\mathfrak b\) be ideal sheaves on \(M\). We write
\(\operatorname{Res}_W\mathfrak a\subseteq\operatorname{Res}_W\mathfrak b\) if at every \(x\in W\) every germ of
\(\mathfrak a_x\) agrees on \(W\), near \(x\), with a combination \(\sum_ic_is_i\) of germs \(s_i\) of
\(\mathfrak b_x\) with holomorphic \(c_i\); and \(\operatorname{Res}_W\mathfrak a=\operatorname{Res}_W\mathfrak b\) if both
inclusions hold. The ideal \(\widetilde{\mathfrak a}^W\) consists of the holomorphic germs that agree on \(W\),
near the point, with such a combination of germs of \(\mathfrak a\) (the preimage of the restriction of \(\mathfrak a\)
to \(W\)). An ideal \(\mathfrak l\) is \emph{invertible along \(W\)} if near every point of \(W\) it is generated by
one holomorphic function whose zero set contains no nonempty relatively open subset of \(W\).
\end{definition}

\begin{lemma}[Calculus of restricted inclusions]
\label{ana:restricted-calculus}
\uses{ana:restricted-ideals,ana:ideal-sheaves}
\lean{TranscendentalNakamaye.Pending.RestrictLeCalculus.RestrictLe.trans,TranscendentalNakamaye.Pending.RestrictLeCalculus.restrictLe_of_le,TranscendentalNakamaye.Pending.RestrictLeCalculus.restrictLe_refl,TranscendentalNakamaye.Pending.RestrictLeCalculus.restrictLe_pullback,TranscendentalNakamaye.Pending.RestrictLeCalculus.restrictLe_of_restrictLe_mul_mul,TranscendentalNakamaye.Pending.RestrictLeCalculus.isInvertibleAlong_of_isInvertible,TranscendentalNakamaye.Pending.RestrictLeCalculus.idealPow_add,TranscendentalNakamaye.Pending.RestrictedIdealExtension.restrictOn_eq_sum_vanishingIdeal,TranscendentalNakamaye.Pending.RestrictedIdealExtension.isCoherent_restrictOn_of_isClosedSubmanifold}
\leanok
Let \(\mathfrak a,\mathfrak b,\mathfrak c,\mathfrak l\) be ideal sheaves on \(M\) and \(W\subseteq M\).
\begin{enumerate}
\item Restricted inclusions are transitive; if \(\mathfrak b\) is coherent, \(\mathfrak a\subseteq\mathfrak b\) implies
\(\operatorname{Res}_W\mathfrak a\subseteq\operatorname{Res}_W\mathfrak b\), in particular
\(\operatorname{Res}_W\mathfrak b\subseteq\operatorname{Res}_W\mathfrak b\).
\item If \(f:Y\to M\) is holomorphic, \(W'\subseteq Y\) with \(f(W')\subseteq W\), \(\mathfrak b\) is coherent and
\(\operatorname{Res}_W\mathfrak a\subseteq\operatorname{Res}_W\mathfrak b\), then
\(\operatorname{Res}_{W'}(\mathfrak a\mathcal O_Y)\subseteq\operatorname{Res}_{W'}(\mathfrak b\mathcal O_Y)\).
\item If \(\mathfrak a,\mathfrak b\) are coherent, \(\mathfrak l\) is invertible along \(W\) and
\(\operatorname{Res}_W(\mathfrak a\mathfrak l)\subseteq\operatorname{Res}_W(\mathfrak b\mathfrak l)\), then
\(\operatorname{Res}_W\mathfrak a\subseteq\operatorname{Res}_W\mathfrak b\).
\item An invertible ideal whose zero locus contains no nonempty relatively open subset of \(W\) is invertible
along \(W\); and \(\mathfrak a^{r+p}=\mathfrak a^r\mathfrak a^p\) for coherent \(\mathfrak a\).
\item If \(\mathfrak a\) is coherent, \(\widetilde{\mathfrak a}^W=\mathfrak a+\mathcal I_W\); it is coherent when moreover
\(W\) is a closed submanifold.
\end{enumerate}
\end{lemma}
\begin{proof}
\leanok
(i) Substitute the combinations; a germ of a coherent \(\mathfrak b\) is the combination of itself. (ii) Compose a
relation \(a=\sum_id_ib_i\), valid on \(W\) near \(f(y)\), with \(f\): it holds on \(W'\) near \(y\) since
\(f(W')\subseteq W\), and the \(b_i\circ f\) are germs of \(\mathfrak b\mathcal O_Y\) (combinations of the
\(\mathfrak b\)-generators composed with \(f\), by coherence). (iii) Let \(h\) generate \(\mathfrak l\) near \(x\in W\) and
\(g\in\mathfrak a_x\): \(gh\) agrees on \(W\) with \(\sum_ic_is_ih\), \(s_i\in\mathfrak b_x\), so \(h(g-\sum_ic_is_i)=0\) on
\(W\) near \(x\), and \(g=\sum_ic_is_i\) on \(W\) near \(x\) because \(h\) vanishes on no nonempty relatively open subset
of \(W\) and both sides are continuous. (iv) The generator of an invertible ideal is the required function; the
power identity is induction on \(p\), with associativity and the unit ideal for coherent ideals. (v) If \(g\) agrees on
\(W\) with \(\sum_ic_is_i\), then \(g-\sum_ic_is_i\) is holomorphic and vanishes on \(W\) near the point, a germ of
\(\mathcal I_W\); the converse inclusion is clear. A sum of coherent ideals is coherent (concatenate the generators),
and \(\mathcal I_W\) is coherent for a closed submanifold (Hadamard's lemma in flattening coordinates).
\end{proof}

\begin{lemma}[Dividing along \(W\) by an ideal invertible along \(W\)]
\label{ana:restricted-colon}
\uses{ana:restricted-ideals,ana:restricted-calculus,ana:invertible-colon}
\source{transcendental-nakamaye-copy4-20260911:prop:principalization}
\lean{TranscendentalNakamaye.Pending.RestrictedIdealQuotient.colonAlong,TranscendentalNakamaye.Pending.RestrictedIdealQuotient.restrictLe_mul_colon,TranscendentalNakamaye.Pending.RestrictedIdealQuotient.restrictLe_mul_colon',TranscendentalNakamaye.Pending.RestrictedIdealQuotient.isCoherent_colonAlong,TranscendentalNakamaye.Pending.RestrictedIdealQuotient.exists_isCoherent_restrictLe_mul,TranscendentalNakamaye.Pending.RestrictLeCalculus.not_le_vanishingIdeal_of_restrictLe_mul}
\leanok
Let \(W\subseteq M\) be closed with \(\mathcal I_W\) coherent, \(\mathfrak a\) coherent, and \(\mathfrak l\) coherent and
invertible along \(W\) with \(\operatorname{Res}_W\mathfrak a\subseteq\operatorname{Res}_W\mathfrak l\). Then
\(\mathfrak c=(\widetilde{\mathfrak a}^W:\mathfrak l)\) is coherent and
\(\operatorname{Res}_W\mathfrak a=\operatorname{Res}_W(\mathfrak l\mathfrak c)\). If moreover \(\mathfrak a\not\subseteq\mathcal I_W\),
then \(\mathfrak c\not\subseteq\mathcal I_W\).
\end{lemma}
\begin{proof}
\leanok
Near \(x\in W\) let \(h\) generate \(\mathfrak l\). A germ \(g\in\mathfrak a_x\) agrees on \(W\) with \(\sum_ic_il_i\), \(l_i\in\mathfrak l_x\);
writing \(l_i=e_ih\), \(g\) agrees on \(W\) with \(h\sum_ic_ie_i\), and \(h\sum_ic_ie_i\in\widetilde{\mathfrak a}^W\), so
\(\sum_ic_ie_i\in\mathfrak c_x\) and \(\operatorname{Res}_W\mathfrak a\subseteq\operatorname{Res}_W(\mathfrak l\mathfrak c)\); no division is
performed. Conversely \(\mathfrak l\mathfrak c\subseteq\widetilde{\mathfrak a}^W=\mathfrak a+\mathcal I_W\)
(Lemma~\ref{ana:restricted-calculus}~(v)), which restricts into \(\mathfrak a\). Coherence: off the closed \(W\),
\(\widetilde{\mathfrak a}^W\) is the whole structure sheaf and \(\mathfrak c\) is generated by \(1\); near \(x\in W\),
\(\mathfrak c\) is generated by the quotients by \(h\) of the generators of \(\widetilde{\mathfrak a}^W\) restricted as above,
which exist by the first part, together with the generators of \(\mathcal I_W\), as in
Lemma~\ref{ana:invertible-colon}, the cancellation of \(h\) along \(W\) replacing that of \(h\) on \(M\). Finally, if
\(\mathfrak c\subseteq\mathcal I_W\), every germ of \(\mathfrak a\) restricts to \(0\) on \(W\), that is
\(\mathfrak a\subseteq\mathcal I_W\).
\end{proof}

\begin{lemma}[Copy-4's excess ideal]
\label{ana:excess-ideal}
\uses{ana:restricted-calculus,ana:restricted-colon,ana:principalization-stage,ana:strict-transform-dense-over-complement,ana:blowups}
\source{transcendental-nakamaye-copy4-20260911:prop:principalization,transcendental-nakamaye-copy4-20260911:lem:ideal-comparison}
\lean{TranscendentalNakamaye.Pending.ExcessIdealPrincipalization.restrictLe_pullback_divided,TranscendentalNakamaye.Pending.ExcessIdealPrincipalization.excessIdeal_spec,TranscendentalNakamaye.Pending.ExcessIdealPrincipalization.restrictLe_pullback_excess}
\leanok
Let \(X\) be a compact complex manifold, \(V\subset X\) a closed irreducible analytic subvariety, and
\(\mathfrak B_m,\mathfrak B_M,\widetilde{\mathfrak B}\) coherent ideals on \(X\) with
\(\operatorname{Res}_V(\widetilde{\mathfrak B}\,\mathfrak B_m^p)\subseteq\operatorname{Res}_V\mathfrak B_M\)
(the comparison of Lemma~\ref{ana:ideal-comparison}, \(\mathfrak B=\widetilde{\mathfrak B}|_V\)) and
\(V\not\subseteq Z(\mathfrak B_m)\), \(V\not\subseteq Z(\mathfrak B_M)\). Let \(\pi:Y_1\to X\) be a sequence of blowups with strict
transform \(W_1\), and \(\mathfrak h,\mathfrak P\) invertible ideals on \(Y_1\) with \(\mathfrak B_m\mathcal O_{Y_1}=\mathfrak h\) and
\(\mathfrak B_M\mathcal O_{Y_1}=\mathfrak h^{r+p}\mathfrak P\). Put \(\mathfrak l=\mathfrak h^r\mathfrak P\).
\begin{enumerate}
\item \(\operatorname{Res}_{W_1}(\widetilde{\mathfrak B}\mathcal O_{Y_1})\subseteq\operatorname{Res}_{W_1}\mathfrak l\).
\item If \(W_1\) is a closed connected submanifold and \(V\not\subseteq Z(\widetilde{\mathfrak B})\), the excess ideal
\(\widetilde{\mathfrak f}_B=(\widetilde{(\widetilde{\mathfrak B}\mathcal O_{Y_1})}^{W_1}:\mathfrak l)\) is coherent, not contained in
\(\mathcal I_{W_1}\), and \(\operatorname{Res}_{W_1}(\widetilde{\mathfrak B}\mathcal O_{Y_1})=\operatorname{Res}_{W_1}(\mathfrak l\,\widetilde{\mathfrak f}_B)\).
\item If \(\operatorname{Res}_{W_1}(\widetilde{\mathfrak B}\mathcal O_{Y_1})=\operatorname{Res}_{W_1}(\mathfrak l\mathfrak f)\) for coherent
\(\mathfrak l,\mathfrak f\), and \(g:Y\to Y_1\) is holomorphic with \(g(W)\subseteq W_1\) and
\(\mathfrak f\mathcal O_Y=\mathcal O_Y(-P^F-P^{F\prime})\), then
\(\operatorname{Res}_W(\widetilde{\mathfrak B}\mathcal O_Y)=\operatorname{Res}_W\bigl((\mathfrak l\mathcal O_Y)\,\mathcal O_Y(-P^F-P^{F\prime})\bigr)\).
\end{enumerate}
\end{lemma}
\begin{proof}
\leanok
(i) The strict transform maps into \(V\), so the comparison pulls back
(Lemma~\ref{ana:restricted-calculus}~(ii)): \(\operatorname{Res}_{W_1}(\widetilde{\mathfrak B}\mathcal O\cdot\mathfrak h^p)\subseteq
\operatorname{Res}_{W_1}(\mathfrak h^{r+p}\mathfrak P)=\operatorname{Res}_{W_1}(\mathfrak l\,\mathfrak h^p)\). The ideal \(\mathfrak h^p\) is
invertible, and its zero locus \(\pi^{-1}Z(\mathfrak B_m)\) contains no nonempty relatively open subset of \(W_1\):
every such subset contains points off \(\operatorname{Exc}\pi\) over \(V\setminus(Z(\mathfrak B_m)\cup Z(\mathfrak B_M))\)
(Lemma~\ref{ana:strict-transform-dense-over-complement}, applied twice). So \(\mathfrak h^p\) is invertible along \(W_1\), and
cancelling it (Lemma~\ref{ana:restricted-calculus}~(iii), (iv)) gives (i). (ii) \(\mathfrak l\) is invertible along \(W_1\) by
the same density argument, \(\mathcal I_{W_1}\) is coherent, and Lemma~\ref{ana:restricted-colon} applies with
\(\mathfrak a=\widetilde{\mathfrak B}\mathcal O_{Y_1}\); \(\widetilde{\mathfrak B}\mathcal O_{Y_1}\not\subseteq\mathcal I_{W_1}\) because a germ of
\(\widetilde{\mathfrak B}\) not vanishing on \(V\) pulls back to one not vanishing at points of \(W_1\) over \(V\setminus
Z(\widetilde{\mathfrak B})\). (iii) Pull both restricted inclusions back along \(g\)
(Lemma~\ref{ana:restricted-calculus}~(ii)), and use \((\mathfrak l\mathfrak f)\mathcal O_Y=(\mathfrak l\mathcal O_Y)(\mathfrak f\mathcal O_Y)\).
\end{proof}
On \(W_1\) the ideal \(\mathfrak l\) is Copy-4's \(\mathcal O_{W_1}(-(G_M-pG_m))\), up to the divisors \(H'\), \(P'\) disjoint from
\(W_1\) which Lemma~\ref{ana:copy4-step-one-stages} keeps ambiently, and (ii) is Copy-4's identity
\(\mathfrak B\,\mathcal O_{W_1}=\mathcal O_{W_1}(-(G_M-pG_m))\cdot\mathfrak f_B\); (iii) is its \texttt{eq:PF-restriction} read through
\(\mathfrak B\).

\begin{lemma}[The restriction identity of Step~4, weight form]
\label{ana:dl-restriction}
\uses{ana:excess-ideal,ana:divisor-metric,ana:restricted-calculus,ana:ideal-sheaves}
\source{transcendental-nakamaye-copy4-20260911:prop:principalization}
\lean{TranscendentalNakamaye.Pending.StepFourLogSingularities.gb_decomposition_ereal,TranscendentalNakamaye.Pending.StepFourLogSingularities.sameSingularityTypeOn_stepFourWeights}
\leanok
Let \(\pi:Y\to X\) be a holomorphic map of compact complex manifolds, \(W\subseteq Y\) compact,
\(\widetilde{\mathfrak B}\) a coherent ideal on \(X\) with a weight \(\psi_B\) of type \((1,\widetilde{\mathfrak B})\), and
\(H^{\mathrm{fix}},H_m,P^{\mathrm{mov}},P_M,P^F,P^{F\prime},\widetilde G_M\) effective divisors on \(Y\), with \(H_m,P_M,P^{F\prime}\)
disjoint from \(W\), \(\mathcal O(-\widetilde G_M)=\mathfrak h^q\mathfrak P\) and
\(\operatorname{Res}_W(\widetilde{\mathfrak B}\mathcal O_Y)=\operatorname{Res}_W(\mathfrak h^r\mathfrak P\,\mathcal O(-P^F-P^{F\prime}))\), where
\(\mathfrak h=\mathcal O(-H^{\mathrm{fix}}-H_m)\), \(\mathfrak P=\mathcal O(-P^{\mathrm{mov}}-P_M)\). Let \(M>0\), \(C\ge1\), \(r\ge1\) and
\(q=Cr\). For all smooth metrics on \(\widetilde G_M,P^{\mathrm{mov}},P^F\),
\[
 \tfrac CM\,\psi_B\circ\pi=\tfrac1M\psi_{\widetilde G_M}+\tfrac{C-1}M\psi_{P^{\mathrm{mov}}}+\tfrac CM\psi_{P^F}+O(1)
\]
on \(W\) (both sides \(-\infty\) at the same points).
\end{lemma}
\begin{proof}
\leanok
Weights of ideals with the same restriction to the compact \(W\) agree on \(W\) up to \(O(1)\), and weights of a pullback,
a product and a power are the composed weight, the sum and the multiple (Lemma~\ref{ana:divisor-metric} for the divisor
metrics). So \(\psi_B\circ\pi=r\,\psi_{\mathfrak h}+\psi_{\mathfrak P}+\psi_{P^F}+O(1)\) on \(W\), and
\(\psi_{\widetilde G_M}=q\,\psi_{\mathfrak h}+\psi_{\mathfrak P}+O(1)\); on \(W\), \(\psi_{\mathfrak h}=\psi_{H^{\mathrm{fix}}}+O(1)\) and
\(\psi_{\mathfrak P}=\psi_{P^{\mathrm{mov}}}+O(1)\), the divisors \(H_m,P_M,P^{F\prime}\) being disjoint from the compact \(W\). The
identity \(\tfrac CM(ra+b+c)=\tfrac1M(qa+b)+\tfrac{C-1}Mb+\tfrac CMc\) for \(q=Cr\), valid in \([-\infty,\infty)\), is
Copy-4's \texttt{eq:GB-decomposition}.
\end{proof}

\begin{lemma}[Copy-4 Step~4: the logarithmic singularities of \(\varphi'\)]
\label{ana:step-four-log-singularities}
\uses{ana:dl-restriction,ana:step-four-current,ana:divisor-restriction,ana:submanifold-complex-structure}
\source{transcendental-nakamaye-copy4-20260911:prop:principalization}
\lean{TranscendentalNakamaye.Pending.StepFourLogSingularities.logSingularitiesOfSameSingularityType,TranscendentalNakamaye.Pending.StepFourLogSingularities.stepFourLogSingularities}
\leanok
Let \(\pi:Y\to X\) be holomorphic with \(Y\) compact, \(W\subseteq Y\) a closed connected submanifold with
\(\pi(W)\subseteq V\), \(M,C,\tau\in\mathbf Q\) with \(M>0\), \(C\ge1\), \(\tau\ge0\), and
\(D=\tfrac1M\widetilde G_M+\tfrac{C-1}MP^{\mathrm{mov}}+\tfrac CMP^F+\tau E\) for metered effective divisors on \(Y\)
(the divisor of Definition~\ref{ana:copy4-step-one-data}), with \(\psi_D|_W\) not identically \(-\infty\) near any point
of \(W\). If \(\varphi_K\) has the singularity type of \(\tfrac CM\psi_B\) on \(V\), and the identity of
Lemma~\ref{ana:dl-restriction} holds on \(W\), then \(\varphi'=\varphi_K\circ\pi|_W+\tau\psi_E|_W\) has logarithmic
singularities along \(D|_W\) with the restricted metric (Definition~\ref{ana:divisor-restriction}).
\end{lemma}
\begin{proof}
\leanok
On \(W\), \(\varphi_K\circ\pi|_W=\tfrac CM\psi_B\circ\pi+O(1)\) and, by Lemma~\ref{ana:dl-restriction},
\(\tfrac CM\psi_B\circ\pi=\tfrac1M\psi_{\widetilde G_M}+\tfrac{C-1}M\psi_{P^{\mathrm{mov}}}+\tfrac CM\psi_{P^F}+O(1)\);
adding \(\tau\psi_E\), \(\varphi'=\psi_D|_W+O(1)\), both sides being \(-\infty\) at the same points. A function of the same
singularity type as the weight \(\psi_{D|_W}\), which is never \(+\infty\), has logarithmic singularities along \(D|_W\): its
residual \(\varphi'-\psi_{D|_W}\) (put \(0\) where both are \(-\infty\)) is bounded.
\end{proof}
In Copy-4, \(M=C\ell\) and \(\varphi_K=\varphi_{K,\ell}\) has analytic singularities of type \((\ell^{-1},\mathfrak B)=(C/M,\mathfrak B)\),
so it has the singularity type of \(\tfrac CM\psi_B\); then \(\varphi'\) is the potential of the field (vii) of
Definition~\ref{ana:copy4-step-one-data}, a K\"ahler-current potential by Lemma~\ref{ana:step-four-current}.

\begin{proposition}[Copy-4's Step~1 with the excess ideal]
\label{ana:copy4-step-one-excess}
\uses{ana:copy4-step-one-stages,ana:excess-ideal,ana:dl-restriction,ana:principalization-stage,ana:relative-principalization}
\source{transcendental-nakamaye-copy4-20260911:prop:principalization,transcendental-nakamaye-copy4-20260911:sec:fixed-principalization}
\lean{TranscendentalNakamaye.Pending.StepOneExcessAssembly.exists_stepOneModel}
\leanok
Assume relative principalization (Lemma~\ref{ana:relative-principalization}) for compact connected manifolds of dimension
\(n\). Let \(X\) be a compact connected complex manifold of dimension \(n\), \(V\subset X\) a closed irreducible analytic
subvariety which is a closed connected submanifold of dimension \(k<n\), and \(\mathfrak B_m,\mathfrak B_M,\widetilde{\mathfrak B}\)
coherent ideals with \(\mathfrak B_M\subseteq\mathfrak B_m^{r+p}\),
\(\operatorname{Res}_V(\widetilde{\mathfrak B}\mathfrak B_m^p)\subseteq\operatorname{Res}_V\mathfrak B_M\), and \(V\) contained in none of
\(Z(\mathfrak B_m)\), \(Z(\mathfrak B_M)\), \(Z(\widetilde{\mathfrak B})\). Then there is a sequence of blowups \(\pi:Y\to X\) adapted to
\(V\), with \(Y\) compact and connected and strict transform \(W\) a closed connected submanifold of dimension \(k\), and
effective divisors \(\widetilde G_M,H^{\mathrm{fix}},P^{\mathrm{mov}},P^F\) on \(Y\) with
\(\mathcal O(-\widetilde G_M)=\mathfrak B_M\mathcal O_Y\), \(\mathfrak B_m\mathcal O_Y\subseteq\mathcal O(-H^{\mathrm{fix}})\),
\(\mathcal I_V\mathcal O_Y=\mathcal O(-H^{\mathrm{fix}}-P^{\mathrm{mov}}-P^F)\,\mathcal I_W\), \(P^{\mathrm{mov}}\) and \(P^F\) supported in
\(\operatorname{Exc}\pi\), and such that for every weight \(\psi_B\) of type \((1,\widetilde{\mathfrak B})\), all smooth metrics on
\(\widetilde G_M,P^{\mathrm{mov}},P^F\), \(M>0\), \(C\ge1\) and \(r\ge1\) with \(r+p=Cr\), the identity of
Lemma~\ref{ana:dl-restriction} holds on \(W\).
\end{proposition}
\begin{proof}
\leanok
Lemma~\ref{ana:copy4-step-one-stages}~(i) gives the first two stages \(\pi_1:Y_1\to X\) with \(\mathfrak B_m\mathcal O_{Y_1}=\mathfrak h\)
and \(\mathfrak B_M\mathcal O_{Y_1}=\mathfrak h^{r+p}\mathfrak P\), \(\mathfrak h=\mathcal O(-H^{\mathrm{fix}}-H_m)\),
\(\mathfrak P=\mathcal O(-P^{\mathrm{mov}}-P_M)\) invertible. Lemma~\ref{ana:excess-ideal}~(ii) gives the excess ideal
\(\widetilde{\mathfrak f}_B\), coherent and not contained in \(\mathcal I_{W_1}\); Lemma~\ref{ana:principalization-stage}~(ii) principalizes
it along \(W_1\) by \(g:Y\to Y_1\), \(\widetilde{\mathfrak f}_B\mathcal O_Y=\mathcal O(-P^F-P^{F\prime})\) with \(P^F\) exceptional and
\(P^{F\prime}\) disjoint from \(W\), and Lemma~\ref{ana:principalization-stage}~(iii) carries \(\mathcal I_V\mathcal O\) across the third
stage. Lemma~\ref{ana:excess-ideal}~(iii) gives the restriction identity on \(W\) and Lemma~\ref{ana:dl-restriction} its weight form;
the remaining fields are pullbacks of the stage-1--2 identities, as in Lemma~\ref{ana:copy4-step-one-stages}~(ii)--(iii).
\end{proof}

\begin{definition}[The data of Copy-4 on \(X\) and \(V\)]
\label{ana:copy4-fixed-data}
\uses{ana:copy4-step-one-data,ana:restricted-ideals,ana:ideal-sheaves,ana:submanifold-complex-structure}
\source{transcendental-nakamaye-copy4-20260911:sec:fixed-data,transcendental-nakamaye-copy4-20260911:sec:positive-current,transcendental-nakamaye-copy4-20260911:prop:principalization}
\lean{TranscendentalNakamaye.Pending.Copy4FixedDataAssembly.Copy4FixedData}
\leanok
Let \((X,\omega)\) be compact K\"ahler, \(\theta\) a smooth closed real \((1,1)\)-form, and \(V\subset X\) a closed
irreducible analytic subvariety which is a closed connected submanifold, with the complex structure of
Lemma~\ref{ana:submanifold-complex-structure}. The \emph{data of Copy-4 on \(X\) and \(V\)} consist of:
\begin{enumerate}
\item the fields (ii), (iii), (iv) and (vi) of Definition~\ref{ana:copy4-step-one-data}, with \(\mathcal I_V\) the ideal of
\(V\): \(\Phi_R\), \(\psi_{\mathcal I_V}\), \(\psi_{\mathfrak j}\), \(A\ge0\); \(c_0\ge0\), the rational level \(M>0\), \(U_M\) and
\(\psi_{\mathfrak B_M}\) with \(V\not\subset Z(\mathfrak B_M)\); the ideal function \(\Psi\) of \(\mathfrak j\mathfrak b\), \(\psi_{\mathfrak b}\) of
type \((1,\mathfrak b)\) (Copy-4's \(\mathfrak B_m\)), and the constants \(L\), \(C\), \(\tau\), \(\delta\); a point \(x_0\in V\) near which
\(\Psi\) is smooth --- that is, every field of the Step-1 data that does not mention the modification;
\item integers \(r\ge1\), \(p\ge0\) with \(r+p=Cr\), and \(\mathfrak B_M\subseteq\mathfrak b^{\,r+p}\);
\item a coherent ideal \(\widetilde{\mathfrak B}\) on \(X\) with \(V\not\subset Z(\widetilde{\mathfrak B})\),
\(\operatorname{Res}_V(\widetilde{\mathfrak B}\,\mathfrak b^{\,p})\subseteq\operatorname{Res}_V\mathfrak B_M\)
(Definition~\ref{ana:restricted-ideals}; the comparison of Lemma~\ref{ana:ideal-comparison}), and a weight \(\psi_B\) of
type \((1,\widetilde{\mathfrak B})\);
\item a real \(\delta_0\) with \(\delta+\tau\le\delta_0\) and a \((\theta-\delta_0\omega)|_V\)-plurisubharmonic potential \(\varphi_K\) on
\(V\) (so \((\theta-\delta_0\omega)|_V+dd^c\varphi_K\ge0\), and \(\varphi_K\) is not identically \(-\infty\) near any point) with the
singularity type of \(\tfrac CM\psi_B\) on \(V\).
\end{enumerate}
\end{definition}
In Copy-4, \(r+p=q=M/m\), \(\mathfrak B_M\subseteq\mathfrak B_m^{\,q}\) is subadditivity of multiplier ideals, \(\widetilde{\mathfrak B}\) restricts
to \(\mathfrak B=\mathcal I(\ell\varphi_K)\) on \(V\), and \(\varphi_K=\varphi_{K,\ell}\) with \(\ell^{-1}=C/M\) comes from the quasi-equisingular
approximation of a positive current on \(V\).

\begin{proposition}[Copy-4 Steps~1--4 from the data on \(X\) and \(V\)]
\label{ana:copy4-step-one-from-fixed-data}
\uses{ana:copy4-fixed-data,ana:copy4-step-one-data,ana:copy4-step-one-excess,ana:dl-restriction,ana:exceptional-metric,ana:divisor-metric,ana:step-four-current,ana:step-four-log-singularities,ana:relative-principalization,ana:zero-locus-coherent,ana:ideal-sheaves}
\source{transcendental-nakamaye-copy4-20260911:prop:principalization}
\lean{TranscendentalNakamaye.Pending.Copy4FixedDataAssembly.Copy4FixedData.not_subset_zeroLocus_idealb,TranscendentalNakamaye.Pending.Copy4FixedDataAssembly.exists_copy4StepOneData_of_fixedData}
\leanok
Assume relative principalization (Lemma~\ref{ana:relative-principalization}) for compact connected manifolds of dimension
\(n\). Let \((X,\omega)\) be a compact connected K\"ahler manifold of dimension \(n\), and let \(\theta\) and \(V\), of dimension
\(k<n\), carry the data of Definition~\ref{ana:copy4-fixed-data}. Then \(V\not\subset Z(\mathfrak b)\), and some sequence of
blowups \(\pi:Y\to X\) adapted to \(V\), whose strict transform \(W\) is a closed connected submanifold of dimension \(k\), carries
Step-1 data (Definition~\ref{ana:copy4-step-one-data}).
\end{proposition}
\begin{proof}
\leanok
\(\Psi\) is smooth, hence finite, near \(x_0\in V\), and its pole set is \(Z(\mathfrak j\mathfrak b)\supseteq Z(\mathfrak b)\), since \(\Psi\) has
analytic singularities of type \((1,\mathfrak j\mathfrak b)\); so \(x_0\notin Z(\mathfrak b)\). Proposition~\ref{ana:copy4-step-one-excess} for
\(\mathfrak B_m=\mathfrak b\), \(\mathfrak B_M\) and \(\widetilde{\mathfrak B}\) gives \(\pi\), \(W\) and the divisors
\(\widetilde G_M,H^{\mathrm{fix}},P^{\mathrm{mov}},P^F\) of field (v), with the identity of Lemma~\ref{ana:dl-restriction} on \(W\) for
\(\psi_B\) and every choice of metrics. Lemma~\ref{ana:exceptional-metric} gives \(E\) (field (i)) and Lemma~\ref{ana:divisor-metric}
smooth metrics on the four divisors. Since \(\pi(W)\subseteq V\) and \(\delta+\tau\le\delta_0\),
\(\varphi'=\varphi_K\circ\pi|_W+\tau\psi_E|_W\) is the potential of a K\"ahler current in the class of
\((\pi^*\theta-\delta\,\pi^*\omega)|_W\) (Lemma~\ref{ana:step-four-current}~(iii)) with logarithmic singularities along \(D|_W\)
(Lemma~\ref{ana:step-four-log-singularities}, with \(r+p=Cr\)): field (vii). The fields (ii), (iii), (iv) and (vi) are those of the
data on \(X\) and \(V\).
\end{proof}

\begin{lemma}[Producing \(U_M\), \(\Psi\) and \(x_0\)]
\label{ana:copy4-fixed-data-fields}
\uses{ana:copy4-fixed-data,ana:regularization-global-lower-bound,dem:multiplier-coherent,ana:ideal-function,ana:irreducible-open-density,ana:zero-locus-coherent,ana:ideal-sheaves}
\source{transcendental-nakamaye-copy4-20260911:sec:ambient,transcendental-nakamaye-copy4-20260911:sec:fixed-principalization,transcendental-nakamaye-copy4-20260911:prop:principalization}
\lean{TranscendentalNakamaye.Pending.Copy4FixedDataProducers.exists_regularization_UM,TranscendentalNakamaye.Pending.Copy4FixedDataProducers.exists_mem_notMem_zeroLocus₂,TranscendentalNakamaye.Pending.Copy4FixedDataProducers.exists_idealFunction_of_not_subset}
\leanok
Let \(X\) be a complex manifold of dimension \(n\).
\begin{enumerate}
\item If \((X,\omega)\) is a connected compact K\"ahler manifold and \(n\ge1\), there is \(c_0>0\) such that for every smooth closed real
\((1,1)\)-form \(\theta\), every \(\theta\)-plurisubharmonic potential \(\varphi\) and every integer \(M\ge1\) there are a weight
\(\psi_{\mathfrak B_M}\) of type \((1,\mathcal I(M\varphi))\) and a \((\theta+\tfrac{c_0}M\omega)\)-plurisubharmonic potential \(U_M\) with
\(\tfrac1M\psi_{\mathfrak B_M}\le U_M+O(1)\) on \(X\).
\item If \(V\subset X\) is a closed irreducible analytic subvariety and \(\mathfrak a,\mathfrak b\) are coherent ideals with
\(V\not\subset Z(\mathfrak a)\) and \(V\not\subset Z(\mathfrak b)\), some point of \(V\) lies in neither \(Z(\mathfrak a)\) nor \(Z(\mathfrak b)\).
\item If \((X,\omega)\) is compact K\"ahler and \(V\), \(\mathfrak j\), \(\mathfrak b\) are as in (ii), there are \(x_0\in V\), \(c_\Psi>0\) and \(\Psi\)
with \(c_\Psi\omega+dd^c\Psi\ge0\), of analytic singularities of type \((1,\mathfrak j\mathfrak b)\), and smooth near \(x_0\).
\end{enumerate}
\end{lemma}
\begin{proof}
\leanok
(i) Corollary~\ref{ana:regularization-global-lower-bound} gives \(c_0\) first, and for \(\theta\), \(\varphi\), \(M\) and every weight of type
\((1,\mathcal I(M\varphi))\) a regularization \(U_M\) dominating \(\tfrac1M\) times it up to a constant on \(X\); such a weight exists because
\(\mathcal I(M\varphi)\) is coherent (Theorem~\ref{dem:multiplier-coherent}, \(n\ge1\)). (ii) \(Z(\mathfrak a)\) and \(Z(\mathfrak b)\) are analytic
subsets (Lemma~\ref{ana:zero-locus-coherent}). Pick \(x\in V\setminus Z(\mathfrak a)\); if \(V\setminus Z(\mathfrak a)\subseteq Z(\mathfrak b)\), then
\(V\cap U\subseteq Z(\mathfrak b)\) for the open \(U=X\setminus Z(\mathfrak a)\ni x\), so \(V\subseteq Z(\mathfrak b)\)
(Lemma~\ref{ana:irreducible-open-density}), a contradiction. (iii) Take \(x_0\) from (ii) and \(\Psi\) from Lemma~\ref{ana:ideal-function} for the
coherent ideal \(\mathfrak j\mathfrak b\) and \(x_0\notin Z(\mathfrak j)\cup Z(\mathfrak b)=Z(\mathfrak j\mathfrak b)\).
\end{proof}
This produces the fields \(c_0\), \(\psi_{\mathfrak B_M}\) and \(U_M\) (with \(\mathfrak B_M=\mathcal I(MV_\theta)\)), \(c_\Psi\), \(\Psi\) and
\(x_0\) of Definition~\ref{ana:copy4-fixed-data}. The constants are chosen in Copy-4's order; what remains is \(\Phi_R\), \(\mathfrak j\) and
\(A\), Skoda's \(V\not\subset Z(\mathfrak B_q)\), subadditivity \(\mathfrak B_M\subseteq\mathfrak B_m^{\,q}\), and the data of the positive current
on \(V\), collected in Definition~\ref{ana:copy4-inputs}.

\begin{definition}[The inputs of Copy-4's Steps~1--4 on \(X\) and \(V\)]
\label{ana:copy4-inputs}
\uses{ana:copy4-fixed-data,ana:step-five-assembly,ana:restricted-ideals,ana:ideal-sheaves,ana:submanifold-complex-structure}
\source{transcendental-nakamaye-copy4-20260911:sec:setup,transcendental-nakamaye-copy4-20260911:sec:fixed-data,transcendental-nakamaye-copy4-20260911:sec:positive-current,transcendental-nakamaye-copy4-20260911:eq:subadditivity}
\lean{TranscendentalNakamaye.Pending.Copy4FixedDataFromInputs.AmbientFixedInput,TranscendentalNakamaye.Pending.Copy4FixedDataFromInputs.TraceInput,TranscendentalNakamaye.Pending.Copy4FixedDataFromInputs.IsLevelPotential}
\leanok
Let \((X,\omega)\) be compact K\"ahler, \(\theta\) a smooth closed real \((1,1)\)-form, and \(V\subset X\) a closed irreducible
analytic subvariety which is a closed connected submanifold, with the complex structure of
Lemma~\ref{ana:submanifold-complex-structure}.
\begin{enumerate}
\item[(I1)] (\emph{Ambient fixed data.}) There are \(\delta_R>0\) and a \((\theta-\delta_R\omega)\)-plurisubharmonic potential \(\Phi_R\), a
coherent ideal \(\mathfrak j\) with \(V\not\subset Z(\mathfrak j)\), functions \(\psi_{\mathcal I_V}\), \(\psi_{\mathfrak j}\) with analytic singularities of
types \((1,\mathcal I_V)\), \((1,\mathfrak j)\), and \(A\ge0\) with the data of Copy-4's lower bound for \(\Phi_R\) (hypothesis (ii) of
Proposition~\ref{ana:step-five-assembly}).
\item[(I2)] (\emph{Multiplier ideals of the trace potential.}) There is a \(\theta\)-plurisubharmonic potential \(\varphi\) with
\(V\not\subset Z(\mathcal I(q\varphi))\) for every integer \(q\ge1\), and \(\mathcal I(qk\varphi)\subseteq\mathcal I(q\varphi)^k\) for all integers
\(q\ge1\), \(k\ge0\).
\item[(I3)] (\emph{A positive current on \(V\).}) For the \(\varphi\) of (I2) and every integer \(C\ge2\) there are an integer \(m\ge1\) and
\(\delta_0>0\) such that for arbitrarily large integers \(r\ge1\) there are a coherent ideal \(\widetilde{\mathfrak B}\) with
\(V\not\subset Z(\widetilde{\mathfrak B})\), a weight \(\psi_B\) of type \((1,\widetilde{\mathfrak B})\) with
\[
 \operatorname{Res}_V\bigl(\widetilde{\mathfrak B}\,\mathcal I(m\varphi)^{(C-1)r}\bigr)\subseteq\operatorname{Res}_V\mathcal I(Cmr\,\varphi)
\]
(Definition~\ref{ana:restricted-ideals}), and a \((\theta-\delta_0\omega)|_V\)-plurisubharmonic potential \(\varphi_K\) on \(V\) with the
singularity type of \(\tfrac1{mr}\psi_B\) on \(V\).
\end{enumerate}
\end{definition}
In Copy-4, (I1) is the fixed data of its section on fixed data (\(\mathfrak j=\mathcal I_{Z'}\) for the union \(Z'\) of the other components
of \(E_{\mathrm{nK}}(\alpha)\)); in (I2), \(\varphi=V_\theta\), the first property is Skoda's theorem (Theorem~\ref{ana:skoda-integrability}; see
Proposition~\ref{ana:copy4-trace-input-at-point}) at a point \(x_*\in V\) with \(\nu(V_\theta,x_*)=0\) and the second is the subadditivity theorem of Demailly--Ein--Lazarsfeld; (I3) is the section on a positive
current on the submanifold, with \(\ell=mr\), \(\varphi_K=\varphi_{K,\ell}\) and \(\widetilde{\mathfrak B}|_V=\mathfrak B\), and the comparison of
Lemma~\ref{ana:ideal-comparison}. None of the three mentions a modification or a constant of Copy-4. (I1) is produced from the bigness
of \(\alpha\) in Proposition~\ref{ana:copy4-ambient-input}, given a coherent ideal of the other components of \(E_{\mathrm{nK}}(\alpha)\).

\begin{proposition}[Copy-4 Steps~1--4 from the inputs on \(X\) and \(V\)]
\label{ana:copy4-step-one-from-inputs}
\uses{ana:copy4-inputs,ana:copy4-fixed-data,ana:copy4-fixed-data-fields,ana:copy4-step-one-from-fixed-data,dem:multiplier-coherent,ana:ideal-sheaves}
\source{transcendental-nakamaye-copy4-20260911:sec:fixed-data,transcendental-nakamaye-copy4-20260911:eq:constants,transcendental-nakamaye-copy4-20260911:sec:ambient,transcendental-nakamaye-copy4-20260911:prop:principalization}
\lean{TranscendentalNakamaye.Pending.Copy4FixedDataFromInputs.nonempty_copy4FixedData,TranscendentalNakamaye.Pending.Copy4FixedDataFromInputs.exists_copy4StepOneData_of_inputs}
\leanok
Let \((X,\omega)\) be a connected compact K\"ahler manifold of dimension \(n\), and let \(\theta\) and \(V\) satisfy (I1)--(I3) of
Definition~\ref{ana:copy4-inputs}.
\begin{enumerate}
\item If \(n\ge1\), \(\theta\) and \(V\) carry the data of Definition~\ref{ana:copy4-fixed-data}.
\item If \(0<\dim V<n\) and relative principalization (Lemma~\ref{ana:relative-principalization}) holds for compact connected manifolds of
dimension \(n\), some sequence of blowups \(\pi:Y\to X\) adapted to \(V\), whose strict transform \(W\) is a closed connected
submanifold of dimension \(\dim V\), carries Step-1 data (Definition~\ref{ana:copy4-step-one-data}).
\end{enumerate}
\end{proposition}
\begin{proof}
\leanok
(i) The constants are chosen in Copy-4's order. Lemma~\ref{ana:copy4-fixed-data-fields}~(i) gives \(c_0>0\), depending only on
\((X,\omega)\). Put \(L=c_0/\delta_R+1\), so \(c_0<L\delta_R\), and \(C=\lceil LA\rceil+2\), so \(C\ge2\) and \(LA\le C-1\). (I3) for \(C\)
gives \(m\) and \(\delta_0\); \(\mathfrak b=\mathcal I(m\varphi)\) is coherent (Theorem~\ref{dem:multiplier-coherent}) and has a weight
\(\psi_{\mathfrak b}\) of type \((1,\mathfrak b)\); since \(V\not\subset Z(\mathfrak j)\) and \(V\not\subset Z(\mathfrak b)\) by (I1), (I2),
Lemma~\ref{ana:copy4-fixed-data-fields}~(iii) gives \(x_0\in V\), \(c_\Psi>0\) and \(\Psi\) of type \((1,\mathfrak j\mathfrak b)\), smooth near
\(x_0\). (I3) then gives \(r\ge1\) so large that \(M=Cmr\) satisfies \(M>L\) and \(M\ge2(C-1)c_\Psi/\delta_0\), with
\(\widetilde{\mathfrak B}\), \(\psi_B\) and \(\varphi_K\). Lemma~\ref{ana:copy4-fixed-data-fields}~(i) at the level \(M\) gives \(U_M\) and a weight
\(\psi_{\mathfrak B_M}\) of \(\mathfrak B_M=\mathcal I(M\varphi)\), with \(V\not\subset Z(\mathfrak B_M)\) by (I2). Choose a rational \(\tau\) with
\(0<\tau<\min\bigl((L\delta_R-c_0)/M,\ \delta_0/2\bigr)\) and put \(\delta=\delta_0/2\), so \(\delta+\tau\le\delta_0\) and
\(\tfrac{C-1}Mc_\Psi\le\delta\). With \(p=(C-1)r\), \(r+p=Cr\); \(\mathfrak B_M=\mathcal I(m\cdot Cr\,\varphi)\subseteq\mathfrak b^{\,Cr}\) by (I2);
the comparison of (I3) is \(\operatorname{Res}_V(\widetilde{\mathfrak B}\mathfrak b^{\,p})\subseteq\operatorname{Res}_V\mathfrak B_M\); and
\(\tfrac CM=\tfrac1{mr}\), so \(\varphi_K\) has the singularity type of \(\tfrac CM\psi_B\). These are the data of
Definition~\ref{ana:copy4-fixed-data}. (ii) follows from (i) (\(n>\dim V>0\)) and Proposition~\ref{ana:copy4-step-one-from-fixed-data}.
\end{proof}

\begin{lemma}[Copy-4's lower bound for \(\Phi_R\) from the Nullstellensatz]
\label{ana:phir-lower-source}
\uses{ana:uniform-nullstellensatz-power,ana:zero-locus-coherent,ana:ideal-sheaves}
\source{transcendental-nakamaye-copy4-20260911:sec:fixed-data}
\lean{TranscendentalNakamaye.Pending.AmbientPhiRSourceOfAnalyticSingularities.exists_ambientPhiRLowerSource_of_type}
\leanok
Let \(X\) be a compact complex manifold, \(c,c_\Phi>0\) rational, \(\psi_{\mathcal I}\), \(\psi_{\mathfrak j}\) functions with analytic
singularities of types \((c,\mathcal I)\), \((c,\mathfrak j)\), and \(\Phi_R\) one of type \((c_\Phi,\mathfrak a_R)\), with
\(Z(\mathfrak a_R)\subseteq Z(\mathcal I\mathfrak j)\). Then there are \(A>0\) and the data of Copy-4's lower bound for \(\Phi_R\) with
constant \(A\) (hypothesis (ii) of Proposition~\ref{ana:step-five-assembly}): an integer \(N>0\) with
\((\mathcal I\mathfrak j)^N\subseteq\mathfrak a_R\), so that locally \(\Phi_R\ge A(\psi_{\mathcal I}+\psi_{\mathfrak j})-O(1)\).
\end{lemma}
\begin{proof}
The three ideals are coherent. Lemma~\ref{ana:uniform-nullstellensatz-power} for \(Z(\mathfrak a_R)\subseteq Z(\mathcal I\mathfrak j)\) gives
\(N>0\) with \((\mathcal I\mathfrak j)^N\subseteq\mathfrak a_R\). Locally, with generators \(g\) of \(\mathcal I\mathfrak j\) and \(f\) of
\(\mathfrak a_R\), each \(g^\beta\), \(|\beta|=N\), is a holomorphic combination of the \(f\), so
\(N\log\sum|g|^2\le\log\sum|f|^2+O(1)\); with \(\psi_{\mathcal I}+\psi_{\mathfrak j}=c\log\sum|g|^2+O(1)\) and
\(\Phi_R=c_\Phi\log\sum|f|^2+O(1)\) this is the bound with \(A=c_\Phi N/c\).
\end{proof}

\begin{proposition}[Copy-4's ambient fixed data from bigness]
\label{ana:copy4-ambient-input}
\uses{ana:copy4-inputs,ana:copy4-step-one-from-inputs,ana:nonkahler-locus-analytic,gnpp:kahler-domination,ana:phir-lower-source,ana:zero-locus-coherent,ana:ideal-sheaves}
\source{transcendental-nakamaye-copy4-20260911:sec:fixed-data}
\lean{TranscendentalNakamaye.Pending.AmbientFixedInputProducer.hasAnalyticSingularitiesOfType_sub_continuous,TranscendentalNakamaye.Pending.AmbientFixedInputProducer.subset_zeroLocus_vanishingIdeal,TranscendentalNakamaye.Pending.AmbientFixedInputProducer.ambientFixedInput_of_cover,TranscendentalNakamaye.Pending.AmbientFixedInputProducer.exists_copy4StepOneData_of_cover}
\leanok
Let \((X,\omega)\) be a compact K\"ahler manifold of dimension \(n\ge1\), \(\alpha=\{\theta\}\) big, \(V\subset X\) a closed irreducible analytic
subvariety which is a closed connected submanifold, and \(\mathfrak j\) a coherent ideal with
\(E_{\mathrm{nK}}(\alpha)\subseteq V\cup Z(\mathfrak j)\) and \(V\not\subset Z(\mathfrak j)\). Then (I1) of Definition~\ref{ana:copy4-inputs} holds.
If moreover \(X\) is connected, \(0<\dim V<n\), relative principalization (Lemma~\ref{ana:relative-principalization}) holds for compact
connected manifolds of dimension \(n\), and (I2)--(I3) hold, some sequence of blowups \(\pi:Y\to X\) adapted to \(V\), whose strict
transform \(W\) is a closed connected submanifold of dimension \(\dim V\), carries Step-1 data
(Definition~\ref{ana:copy4-step-one-data}).
\end{proposition}
\begin{proof}
Theorem~\ref{ana:nonkahler-locus-analytic} gives a K\"ahler current \(T=\theta_T+dd^c\varphi_T\ge\omega_T\) in \(\alpha\), \(\omega_T\) a K\"ahler
form, with analytic singularities of type \((c,\mathfrak a_R)\) and pole set \(E_{\mathrm{nK}}(\alpha)\). Write \(\theta_T=\theta+dd^ch\) with \(h\)
smooth and put \(\Phi_R=\varphi_T+h\): then \(\theta+dd^c\Phi_R=T\ge\omega_T\), and \(\Phi_R\) has the same analytic singularity type as
\(\varphi_T\), adding a continuous function changing only the bounded remainder. Lemma~\ref{gnpp:kahler-domination} gives \(A>0\) with
\(A\omega_T\ge\omega\), so \(\Phi_R\) is \((\theta-\delta_R\omega)\)-plurisubharmonic for \(\delta_R=A^{-1}\). The ideal \(\mathcal I_V\) of the
submanifold \(V\) is coherent, and the coherent ideals \(\mathcal I_V\), \(\mathfrak j\) have weights \(\psi_{\mathcal I_V}\), \(\psi_{\mathfrak j}\) of
types \((1,\mathcal I_V)\), \((1,\mathfrak j)\). By Lemma~\ref{ana:zero-locus-coherent},
\(Z(\mathfrak a_R)=\{\Phi_R=-\infty\}=E_{\mathrm{nK}}(\alpha)\subseteq V\cup Z(\mathfrak j)\subseteq Z(\mathcal I_V\mathfrak j)\), and
Lemma~\ref{ana:phir-lower-source} gives \(A_R>0\) with the lower bound. The second statement is
Proposition~\ref{ana:copy4-step-one-from-inputs}~(ii).
\end{proof}
In Copy-4, \(\mathfrak j=\mathcal I_{Z'}\) for the union \(Z'\) of the irreducible components of \(E_{\mathrm{nK}}(\alpha)\) other than \(V\),
coherent by Cartan's theorem; the two properties of \(\mathfrak j\) used here are proved for \(Z'=\overline{E_{\mathrm{nK}}(\alpha)\setminus V}\) in
Proposition~\ref{ana:copy4-other-components-ideal}.

\begin{proposition}[Input (I2) from one point of zero Lelong number]
\label{ana:copy4-trace-input-at-point}
\uses{ana:copy4-inputs,ana:skoda-integrability,ana:copy4-step-one-from-inputs,ana:nonkahler-locus}
\source{transcendental-nakamaye-copy4-20260911:sec:setup}
\lean{TranscendentalNakamaye.Pending.TraceInputSkoda.TraceInputAtPoint,TranscendentalNakamaye.Pending.TraceInputSkoda.traceInput_of_atPoint,TranscendentalNakamaye.Pending.TraceInputSkoda.exists_copy4StepOneData_of_inputs_atPoint}
\leanok
Let \((X,\omega)\), \(\theta\) and \(V\) be as in Definition~\ref{ana:copy4-inputs}, and let \((\mathrm{I}2')\) be \((\mathrm{I}2)+(\mathrm{I}3)\) with the
clause ``\(V\not\subset Z(\mathcal I(q\varphi))\) for every integer \(q\ge1\)'' replaced by ``\(\nu(\varphi,x_*)=0\) for some \(x_*\in V\)''. Then
\((\mathrm{I}2')\) implies \((\mathrm{I}2)+(\mathrm{I}3)\). Hence, if \(X\) is connected of dimension \(n\), \(0<\dim V<n\), relative principalization
(Lemma~\ref{ana:relative-principalization}) holds for compact connected manifolds of dimension \(n\), and (I1) and \((\mathrm{I}2')\) hold, some
sequence of blowups \(\pi:Y\to X\) adapted to \(V\), whose strict transform \(W\) is a closed connected submanifold of dimension \(\dim V\),
carries Step-1 data (Definition~\ref{ana:copy4-step-one-data}).
\end{proposition}
\begin{proof}
\leanok
By Theorem~\ref{ana:skoda-integrability}, \(\nu(\varphi,x_*)=0\) with \(x_*\in V\) gives \(V\not\subset Z(\mathcal I(q\varphi))\) for every \(q\ge1\); the other
clauses are unchanged. The second statement is then Proposition~\ref{ana:copy4-step-one-from-inputs}~(ii).
\end{proof}

\begin{proposition}[Input (I2) without the subadditivity clause]
\label{ana:copy4-trace-input-no-sub}
\uses{ana:copy4-trace-input-at-point,dem:del00-subadditivity,ana:copy4-inputs,ana:copy4-step-one-from-inputs}
\source{transcendental-nakamaye-copy4-20260911:sec:setup,transcendental-nakamaye-copy4-20260911:eq:subadditivity}
\lean{TranscendentalNakamaye.Pending.MultiplierSubadditivity.TraceInputAtPointNoSub,TranscendentalNakamaye.Pending.MultiplierSubadditivity.traceInputAtPoint_of_noSub,TranscendentalNakamaye.Pending.MultiplierSubadditivity.exists_copy4StepOneData_of_inputs_noSub}
\leanok
In the setting of Proposition~\ref{ana:copy4-trace-input-at-point}, the subadditivity clause
\(\mathcal I(qk\varphi)\subseteq\mathcal I(q\varphi)^k\) of its input may be dropped: a \(\theta\)-PSH potential \(\varphi\) with
\(\nu(\varphi,x_*)=0\) at one point \(x_*\in V\), together with the data of (I3), gives (I2)--(I3); hence, given relative
principalization and (I1), the conclusion of Proposition~\ref{ana:copy4-step-one-from-inputs}~(ii) holds.
\end{proposition}
\begin{proof}
\leanok
The dropped clause holds for every \(\theta\)-plurisubharmonic \(\varphi\) when \(n\ge1\) (Theorem~\ref{dem:del00-subadditivity}),
and \(n>\dim V>0\). The rest is Proposition~\ref{ana:copy4-trace-input-at-point}.
\end{proof}
In Copy-4, \(x_*\in V\setminus E_{\mathrm{nn}}(\alpha)\) (\(V\not\subset E_{\mathrm{nn}}(\alpha)\) by the setup) and \(\varphi=V_\theta\), whose Lelong number
vanishes at \(x_*\).

\begin{proposition}[The minimal-singularity potential has zero Lelong number off the non-nef locus]
\label{ana:minimal-singularity-lelong-vanishing}
\uses{ana:nonkahler-locus,fpenv:envelopes,fpenv:v-universal,gnpp:kahler-domination,ana:lelong-additivity,bc:quotient}
\source{transcendental-nakamaye-copy4-20260911:sec:nK,transcendental-nakamaye-copy4-20260911:sec:setup}
\lean{TranscendentalNakamaye.Pending.MinimalSingularityLelongVanishing.isThetaPshPotential_thetaPshSupremum,TranscendentalNakamaye.Pending.MinimalSingularityLelongVanishing.lelongNumber_thetaPshSupremum_le,TranscendentalNakamaye.Pending.MinimalSingularityLelongVanishing.lelongNumber_thetaPshSupremum_eq_zero}
\leanok
Let \(Y\) be a compact complex manifold of dimension \(n\) with a K\"ahler form \(\omega\), and let
\(\alpha=\{\theta\}\) be big. Then \(V_\theta\) is a \(\theta\)-PSH potential. If \(T_0\) is a K\"ahler current in
\(\alpha\) dominating a K\"ahler form \(\omega_0\), \(A>0\) with \(A\omega_0-\omega\ge0\), \(\varepsilon>0\) and \(T\) a
K\"ahler current in \(\alpha+\varepsilon\{\omega\}\), then with \(s=1/(1+\varepsilon A)\), at every point \(x\),
\[
 \nu(V_\theta,x)\le s\,\nu(T,x)+(1-s)\,\nu(T_0,x).
\]
Consequently, if \(n\ge1\) and \(x\notin E_{\mathrm{nn}}(\alpha)\), then \(\nu(V_\theta,x)=0\).
\end{proposition}
\begin{proof}
\leanok
A K\"ahler current \(T_0=\theta_0+dd^cu_0\) in \(\alpha\) gives \(\theta_0=\theta+dd^cf\) with \(f\) smooth, and \(u_0+f\) is a
\(\theta\)-PSH potential below a constant on the compact \(Y\), hence a competitor; so \(V_\theta\) is a \(\theta\)-PSH
potential. For the estimate write \(T=\theta_T+dd^cu\) and \(w=s\,u+(1-s)\,u_0\). Since \(T\ge\omega_T\) and
\(T_0\ge\omega_0\), the current \(\theta_w=s\,\theta_T+(1-s)(\theta_0-\omega_0)+dd^cw\) is positive, and
\(s\varepsilon(A\omega_0-\omega)\ge0\) with \(s\varepsilon A=1-s\) shows that
\(\theta'=s\,\theta_T+(1-s)\theta_0-s\varepsilon\omega\) satisfies \(\theta'+dd^cw\ge0\). The class of \(\theta'\) is
\(s(\alpha+\varepsilon\{\omega\})+(1-s)\alpha-s\varepsilon\{\omega\}=\alpha\), so \(\theta'=\theta+dd^cg\) with \(g\) smooth and
\(w+g\) is \(\theta\)-PSH, bounded above; minimality of \(V_\theta\) gives \(w+g\le V_\theta+C\). Lelong numbers decrease
under \(\le\) up to constants, are unchanged by smooth terms, are subadditive
(Lemma~\ref{ana:lelong-additivity}) and satisfy \(\nu(cu,x)\le c\,\nu(u,x)\) for \(c\ge0\), so \(\nu(V_\theta,x)\le\nu(w,x)\le s\,\nu(u,x)+(1-s)\,\nu(u_0,x)\).

If \(x\notin E_{\mathrm{nn}}(\alpha)\), then for every \(\varepsilon>0\), \(x\notin E_{\mathrm{nK}}(\alpha+\varepsilon\{\omega\})\): some
K\"ahler current \(T\) in \(\alpha+\varepsilon\{\omega\}\) has \(\nu(T,x)=0\). Fix \(T_0\) (bigness) and \(A\)
(Lemma~\ref{gnpp:kahler-domination}). Then \(\nu(V_\theta,x)\le(1-s)\nu(T_0,x)\le\varepsilon A\,\nu(T_0,x)\), and
\(\nu(T_0,x)<\infty\) as \(n\ge1\) (Lemma~\ref{ana:lelong-additivity}); let \(\varepsilon\to0\).
\end{proof}
This is the inclusion \(\{\nu(V_\theta,\cdot)>0\}\subseteq E_{\mathrm{nn}}(\alpha)\) of Copy-4's description of the non-nef
locus, the direction its setup uses to choose \(x_*\in V\setminus E_{\mathrm{nn}}(\alpha)\) with \(\nu(V_\theta,x_*)=0\).

\begin{proposition}[Input (I2) reduced to the positive current on \(V\)]
\label{ana:copy4-trace-input-positive-current}
\uses{ana:minimal-singularity-lelong-vanishing,ana:copy4-trace-input-no-sub,ana:copy4-trace-input-at-point,ana:copy4-inputs,ana:copy4-step-one-from-inputs,ana:copy4-ambient-input}
\source{transcendental-nakamaye-copy4-20260911:sec:setup,transcendental-nakamaye-copy4-20260911:sec:positive-current}
\lean{TranscendentalNakamaye.Pending.Copy4TraceInputReduction.PositiveCurrentInput,TranscendentalNakamaye.Pending.Copy4TraceInputReduction.traceInputAtPointNoSub_of_positiveCurrentInput,TranscendentalNakamaye.Pending.Copy4TraceInputReduction.traceInput_of_positiveCurrentInput,TranscendentalNakamaye.Pending.Copy4TraceInputReduction.exists_copy4StepOneData_of_positiveCurrentInput,TranscendentalNakamaye.Pending.Copy4TraceInputReduction.exists_copy4StepOneData_of_cover_positiveCurrentInput}
\leanok
Let \((X,\omega)\), \(\theta\) and \(V\) be as in Definition~\ref{ana:copy4-inputs}, with \(\dim X\ge1\), \(\{\theta\}\) big and
\(V\not\subset E_{\mathrm{nn}}(\{\theta\})\). If the data of (I3) exist for \(\varphi=V_\theta\), then (I2) and (I3) hold with
\(\varphi=V_\theta\). Hence, if \(X\) is connected, \(0<\dim V<\dim X\) and relative principalization
(Lemma~\ref{ana:relative-principalization}) holds for compact connected manifolds of dimension \(\dim X\), the conclusion of
Proposition~\ref{ana:copy4-step-one-from-inputs}~(ii) follows from (I1) and (I3) for \(V_\theta\), and also from a coherent ideal
\(\mathfrak j\) with \(E_{\mathrm{nK}}(\{\theta\})\subseteq V\cup Z(\mathfrak j)\), \(V\not\subset Z(\mathfrak j)\) and (I3) for \(V_\theta\).
\end{proposition}
\begin{proof}
\leanok
\(V_\theta\) is a \(\theta\)-PSH potential, and at a point \(x_*\in V\setminus E_{\mathrm{nn}}(\{\theta\})\) its Lelong number vanishes
(Proposition~\ref{ana:minimal-singularity-lelong-vanishing}). Proposition~\ref{ana:copy4-trace-input-no-sub} turns this point and the
data of (I3) into (I2)--(I3). The last two assertions are Proposition~\ref{ana:copy4-step-one-from-inputs}~(ii) and
Proposition~\ref{ana:copy4-ambient-input}.
\end{proof}

\begin{proposition}[The ideal of the other components]
\label{ana:copy4-other-components-ideal}
\uses{ana:component-local-filling,ana:closure-difference-analytic,ana:nonkahler-locus-analytic,ana:ideal-sheaves,ana:copy4-trace-input-positive-current}
\source{transcendental-nakamaye-copy4-20260911:sec:fixed-data}
\lean{TranscendentalNakamaye.Pending.CartanComponentIdeal.CartanCoherence,TranscendentalNakamaye.Pending.CartanComponentIdeal.subset_zeroLocus_vanishingIdeal,TranscendentalNakamaye.Pending.CartanComponentIdeal.subset_union_zeroLocus_vanishingIdeal_closure_diff,TranscendentalNakamaye.Pending.CartanComponentIdeal.not_subset_zeroLocus_vanishingIdeal_closure_diff,TranscendentalNakamaye.Pending.CartanComponentIdeal.exists_copy4StepOneData_of_cartan_positiveCurrentInput,TranscendentalNakamaye.Pending.CartanComponentIdeal.exists_copy4StepOneData_of_cartanCoherence}
\leanok
Let \(A\) be a subset of a complex manifold, \(V\subseteq A\), \(Z'=\overline{A\setminus V}\) and \(\mathcal I_{Z'}\) the ideal of germs of
holomorphic functions vanishing on \(Z'\). Then \(A\subseteq V\cup Z(\mathcal I_{Z'})\). If \(A\) is analytic and \(V\) is an irreducible
component of \(A\), then \(V\not\subset Z(\mathcal I_{Z'})\). Consequently, in the setting of
Proposition~\ref{ana:copy4-steps-one-to-five}, its conclusion holds as soon as \(\mathcal I_{Z'}\) is coherent for
\(Z'=\overline{E_{\mathrm{nK}}(\alpha)\setminus V}\) and the data of (I3) exist for \(V_\theta\); in particular it follows from Cartan's
theorem that the ideal sheaf of every analytic subset is coherent.
\end{proposition}
\begin{proof}
\leanok
A point of \(A\setminus V\) lies in \(Z'\subseteq Z(\mathcal I_{Z'})\). If \(V\) is an irreducible component of the analytic set \(A\),
Proposition~\ref{ana:component-local-filling} gives an open \(U\) meeting \(V\) at a point \(v\) with \(A\cap U\subseteq V\); then
\(U\cap Z'=\emptyset\) (a point of \(Z'\) in \(U\) would be a limit of points of \((A\setminus V)\cap U=\emptyset\)), so the constant
\(1\) is a germ of \(\mathcal I_{Z'}\) at \(v\), and \(v\notin Z(\mathcal I_{Z'})\). In the setting of
Proposition~\ref{ana:copy4-steps-one-to-five}, \(E_{\mathrm{nK}}(\alpha)\) is analytic (Theorem~\ref{ana:nonkahler-locus-analytic}), so
\(\mathfrak j=\mathcal I_{Z'}\) satisfies the hypotheses of the last assertion of
Proposition~\ref{ana:copy4-trace-input-positive-current}. For Cartan's theorem, \(Z'\) is analytic by
Proposition~\ref{ana:closure-difference-analytic}.
\end{proof}
Copy-4 takes \(\mathfrak j=\mathcal I_{Z'}\) for the union \(Z'\) of the irreducible components of \(E_{\mathrm{nK}}(\alpha)\) other than
\(V\); here \(Z'=\overline{E_{\mathrm{nK}}(\alpha)\setminus V}\), and only the two properties of the proposition are used. Coherence of
\(\mathcal I_{Z'}\) is Cartan's theorem (H.~Cartan 1950), Theorem~\ref{ana:cartan-coherence}.

\begin{proposition}[An ideal of the other components without Cartan's theorem]
\label{ana:copy4-colon-component-ideal}
\uses{oka:coherent-colon,ana:uniform-nullstellensatz-power,ana:component-local-filling,ana:nonkahler-locus-analytic,ana:ideal-sheaves,ana:copy4-trace-input-positive-current,ana:copy4-other-components-ideal}
\source{transcendental-nakamaye-copy4-20260911:sec:fixed-data}
\lean{TranscendentalNakamaye.Pending.Copy4ColonComponentIdeal.exists_nhds_idealPow_le_of_eventually,TranscendentalNakamaye.Pending.Copy4ColonComponentIdeal.one_mem_idealPow,TranscendentalNakamaye.Pending.Copy4ColonComponentIdeal.one_mem_vanishingIdeal,TranscendentalNakamaye.Pending.Copy4ColonComponentIdeal.exists_isCoherent_cover_not_subset,TranscendentalNakamaye.Pending.Copy4ColonComponentIdeal.exists_copy4StepOneData_of_positiveCurrentInput}
\leanok
Let \(X\) be a compact complex manifold of dimension \(n>0\), \(\alpha=\{\theta\}\) big, and \(V\subset X\) a closed connected
submanifold which is an irreducible component of \(E_{\mathrm{nK}}(\alpha)\). Let \(\mathfrak a_R\) be the ideal of the analytic
singularities of a K\"ahler current in \(\alpha\) whose polar set is \(E_{\mathrm{nK}}(\alpha)\)
(Theorem~\ref{ana:nonkahler-locus-analytic}) and \(\mathcal I_V\) the ideal of \(V\). For a suitable \(N>0\) the colon ideal
\(\mathfrak j=(\mathfrak a_R:\mathcal I_V^N)\) is coherent, \(E_{\mathrm{nK}}(\alpha)\subseteq V\cup Z(\mathfrak j)\) and \(V\not\subset Z(\mathfrak j)\).
Consequently, in the setting of Proposition~\ref{ana:copy4-steps-one-to-five}, its conclusion holds as soon as the data of (I3)
exist for \(\varphi=V_\theta\); Cartan's theorem is not needed.
\end{proposition}
\begin{proof}
\leanok
The ideal \(\mathfrak a_R\) is coherent with \(Z(\mathfrak a_R)=E_{\mathrm{nK}}(\alpha)\), and \(\mathcal I_V\), locally generated by
coordinate functions of the submanifold \(V\), is coherent, hence so is \(\mathcal I_V^N\). By
Proposition~\ref{ana:component-local-filling} there are \(v\in V\) and an open \(W\ni v\) with \(E_{\mathrm{nK}}(\alpha)\cap W\subseteq V\);
so near \(v\), \(Z(\mathfrak a_R)\subseteq Z(\mathcal I_V)\), and the Nullstellensatz at \(v\)
(Lemma~\ref{ana:uniform-nullstellensatz-power}, whose proof is local) gives \(N>0\) and a neighbourhood of \(v\) on which
\(\mathcal I_V^N\subseteq\mathfrak a_R\). With this \(N\), \(\mathfrak j\) is coherent by Lemma~\ref{oka:coherent-colon}. At a point
\(x\notin V\), \(1\in\mathcal I_{V,x}\), so \(1\in(\mathcal I_V^N)_x\) and \(\mathfrak j_x\subseteq\mathfrak a_{R,x}\); hence every point of
\(E_{\mathrm{nK}}(\alpha)\setminus V\subseteq Z(\mathfrak a_R)\) lies in \(Z(\mathfrak j)\). Near \(v\), \(\mathcal I_V^N\subseteq\mathfrak a_R\) gives
\(1\in\mathfrak j_v\), so \(v\in V\setminus Z(\mathfrak j)\). For the last assertion, \(\mathfrak j\) satisfies the hypotheses of the last
assertion of Proposition~\ref{ana:copy4-trace-input-positive-current}, exactly as \(\mathcal I_{Z'}\) does in
Proposition~\ref{ana:copy4-other-components-ideal}.
\end{proof}
Copy-4 takes \(\mathfrak j=\mathcal I_{Z'}\); only the three properties above are used, and the colon ideal has them without
Cartan's theorem, its coherence resting on Oka's theorem alone.

\begin{theorem}[Cartan's coherence theorem]
\label{ana:cartan-coherence}
\uses{ana:ideal-sheaves,fpgeom:complex-analytic-subvariety,oka:cartan-chart-reduction,oka:cartan-local,oka:cartan-irreducible-reduction,oka:cartan-irreducible-fg}
\source{demailly-cadg:II.4.29}
\lean{TranscendentalNakamaye.Pending.CartanCoherenceInterface.cartanCoherence,TranscendentalNakamaye.Pending.CartanIrreducible.cartanCoherence}
\leanok
For every analytic subset \(A\) of a complex manifold \(M\), the ideal sheaf \(\mathcal I_A\) of germs of holomorphic functions
vanishing on \(A\) is coherent: near every point it is generated, at every point of a neighbourhood, by finitely many
holomorphic functions.
\end{theorem}
\begin{proof}
\leanok
Coherence is local: by Lemma~\ref{oka:cartan-chart-reduction} it suffices to prove the local form in \(\mathbb C^n\),
Theorem~\ref{oka:cartan-local}, which Lemma~\ref{oka:cartan-irreducible-reduction} reduces to irreducible germs at the origin,
Corollary~\ref{oka:cartan-irreducible-fg}.
\end{proof}
H.~Cartan, Bull.\ SMF \textbf{78} (1950); the proof is Demailly's (Theorem~II.4.29), with the choice of primitive elements of
Proposition~\ref{oka:cartan-irreducible}. Dimension \(0\) is included.

\begin{proposition}[Copy-4's positive current on \(V\) for \(V_\theta\)]
\label{ana:copy4-positive-current-input}
\uses{ana:copy4-inputs,ana:copy4-trace-input-positive-current,ana:restricted-volume}
\source{transcendental-nakamaye-copy4-20260911:sec:positive-current}
\lean{TranscendentalNakamaye.Pending.PositiveCurrentInputInterface.ct22Copy4_positiveCurrentInput,TranscendentalNakamaye.Pending.RestrictedApproximantMassAssembly.positiveCurrentInput_of_numericalRestrictedVolume_pos}
\leanok
Let \(X\) be a connected compact K\"ahler manifold of dimension \(n\), \(\alpha=\{\theta\}\) big, and \(V\subset X\) a closed
connected submanifold of dimension \(0<k<n\) which is an irreducible component of \(E_{\mathrm{nK}}(\alpha)\), is not contained in
\(E_{\mathrm{nn}}(\alpha)\), and has \(\operatorname{RV}_{X\mid V}(\alpha)>0\). Then the data of (I3) of
Definition~\ref{ana:copy4-inputs} exist for \(\varphi=V_\theta\): for every integer \(C\ge2\) there are an integer \(m\ge1\) and
\(\delta_0>0\) such that for arbitrarily large integers \(r\ge1\) there are a coherent ideal \(\widetilde{\mathfrak B}\) with
\(V\not\subset Z(\widetilde{\mathfrak B})\), a weight \(\psi_B\) of type \((1,\widetilde{\mathfrak B})\) with
\(\operatorname{Res}_V\bigl(\widetilde{\mathfrak B}\,\mathcal I(mV_\theta)^{(C-1)r}\bigr)\subseteq\operatorname{Res}_V\mathcal I(Cmr\,V_\theta)\),
and a \((\theta-\delta_0\omega)|_V\)-plurisubharmonic potential \(\varphi_K\) on \(V\) with the singularity type of
\(\tfrac1{mr}\psi_B\).
\end{proposition}
\begin{proof}
\uses{pc:positive-current-from-u-part,pc:f-u-part,pc:f-restricted-mass}
\leanok
Let \((\Phi_j)\) be a quasi-equisingular approximation of \(V_\theta\) of global type with positive errors, as in Proposition~\ref{pc:f-u-part}.
By Proposition~\ref{pc:f-restricted-mass} the masses \(b_j\) of its restrictions to \(V\) are bounded below by a positive constant (its proof gives
\(\operatorname{RV}_{X\mid V}(\alpha)\)), so Proposition~\ref{pc:positive-current-from-u-part} gives the data (I3) for \(V_\theta\). The hypothesis that \(V\) is
an irreducible component of \(E_{\mathrm{nK}}(\alpha)\) is not used.
\end{proof}
This is Copy-4's section on a positive current on the submanifold, built from a quasi-equisingular approximation of \(V_\theta\) and the
comparison of Lemma~\ref{ana:ideal-comparison}. Copy-4 obtains the potential \(u\) on \(V\) as the trace \(\operatorname{Tr}^\theta_V(V_\theta)\);
the proof above is the trace-free variant (Proposition~\ref{pc:f-u-part} and the lemmas before it), and through Proposition~\ref{pc:bigness-from-mass}
it rests on Boucksom's criterion (Theorem~\ref{ana:ambient-volume-uniform-big}). The proposition does not follow from the conclusion of
Proposition~\ref{ana:copy4-steps-one-to-five}, which lives on a modification and says nothing about the multiplier ideals of \(V_\theta\) on
\(X\).

\begin{proposition}[Gluing a global and a local positive branch]
\label{ana:gluing}
\uses{ana:gluing-conditions,ana:geometric-setting,psh:maximum,top:collar,top:upper-semicontinuity,top:local-branch}
\source{transcendental-nakamaye-copy4-20260911:prop:gluing}
\lean{TranscendentalNakamaye.Pending.AnalyticGluing.exists_kahlerCurrent_isSmoothNear_of_gluingHypotheses}
\leanok
Let \(Y\) be a connected compact K\"ahler manifold, let \(\theta_Y\) be a smooth
closed real \((1,1)\)-form, let \(U_1\Subset U_2\) be open subsets of
\(Y\), and let \(p\in U_1\). Suppose that \(\varphi_G\) is a
quasi-plurisubharmonic function on \(Y\) and \(\varphi_U\) is a
quasi-plurisubharmonic function on a neighborhood of \(\overline U_2\),
such that:
\begin{enumerate}
\item \(\theta_Y+d d^c\varphi_G\) is a K\"ahler current and
      \(\varphi_G(p)=-\infty\);
\item \(\theta_Y+d d^c\varphi_U\) dominates a K\"ahler form near
      \(\overline U_2\), and \(\varphi_U\) is smooth near \(p\);
\item \(\varphi_U\le\varphi_G+C\) on
      \(\overline U_2\setminus U_1\) for some \(C\in\mathbb R\).
\end{enumerate}
Then \(\{\theta_Y\}\) contains a K\"ahler current which is smooth near
\(p\). In particular, \(p\notin E_{\mathrm{nK}}(\{\theta_Y\})\).
\end{proposition}
\begin{proof}
On \(U_2\) take the maximum of \(\varphi_U-C\) and \(\varphi_G\), and
use \(\varphi_G\) on \(Y\setminus\overline U_1\). The collar inequality
makes these definitions agree on the overlap. Compactness supplies one
K\"ahler form dominated by both branch currents on the regions where
they are used; the maximum of the corresponding plurisubharmonic local
potentials preserves that lower bound. Thus the glued potential defines a
K\"ahler current in \(\{\theta_Y\}\). Near \(p\), upper semicontinuity
and \(\varphi_G(p)=-\infty\) force the smooth local branch to dominate,
so the glued current is smooth there.
\end{proof}

\begin{lemma}[Compact collars from principalization gluing data]
\label{ana:local-gluing-data}
\uses{ana:gluing-conditions,ana:principalization-data,ana:geometric-setting}
\source{transcendental-nakamaye-copy4-20260911:prop:principalization,transcendental-nakamaye-copy4-20260911:sec:gluing}
\lean{TranscendentalNakamaye.Pending.LocalGluingData.exists_collars_gluingHypotheses_of_principalizationGluingConclusions}
\leanok
Assume the data and conclusions of Proposition~\ref{ana:principalization-data}.
Put \(\theta_Y=\pi^*\theta\). There are open neighborhoods
\(W\Subset U_1\Subset U_2\Subset U\), so that \(y\in U_1\), such that the supplied point
\(y\), global function \(\varphi_G\), and local function \(\varphi_U\)
satisfy all three hypotheses of Proposition~\ref{ana:gluing} for
\(Y,\theta_Y,U_1,U_2,y\).
\end{lemma}
\begin{proof}
The compact submanifold \(W\) lies in the open set \(U\), so choose
\(W\Subset U_1\Subset U_2\Subset U\). The first two hypotheses of
Proposition~\ref{ana:gluing} are exactly the first two conclusions of
Proposition~\ref{ana:principalization-data}. Its third conclusion,
applied to the compact set \(\overline U_2\setminus U_1\), supplies
the required collar constant.
\end{proof}

\begin{lemma}[Embedded resolution of a subvariety]
\label{ana:embedded-resolution}
\uses{ana:blowups}
\source{transcendental-nakamaye-copy4-20260911:lem:resolution}
\lean{TranscendentalNakamaye.Pending.EmbeddedResolution.exists_embeddedResolution,TranscendentalNakamaye.Pending.EmbeddedResolution.exists_embeddedResolution_smoothData}
\leanok
Let \(X\) be a compact complex manifold and let \(V\subset X\) be a closed
irreducible analytic subvariety. There is a finite composition of blowups with
smooth centres \(\mu:\widetilde X\to X\) which is an isomorphism over
\(X\setminus V_{\mathrm{sing}}\) and whose strict transform of \(V\) is a closed
connected complex submanifold of \(\widetilde X\) of dimension \(\dim V\).
\end{lemma}
\begin{proof}
\leanok
\uses{ana:blowups,mc:hironaka-from-seq-invariant,hir:canonical-centres,ana:cartan-coherence}
The input taken from resolution of singularities is a
finite composition of blowups with smooth centres \(\mu\) whose exceptional image lies in
\(V_{\mathrm{sing}}\) and whose strict transform of \(V\) is regular of dimension \(\dim V\) at
each of its points. It is Theorem~\ref{mc:hironaka-from-seq-invariant}(3), applied to the resolution invariant in dimension
\(\le n\) of Theorem~\ref{hir:canonical-centres}, the ideal sheaf \(\mathcal I_V\) being coherent by Cartan's theorem
(Theorem~\ref{ana:cartan-coherence}). Then \(\mu\) is an isomorphism over \(X\setminus V_{\mathrm{sing}}\), since
every centre lies over \(V_{\mathrm{sing}}\); the strict transform is closed, as a closure;
and it is connected: since \(V_{\mathrm{reg}}\ne\emptyset\) is not in the exceptional image, the
strict transform is an irreducible analytic subvariety of \(\widetilde X\), and an irreducible
subvariety is the closure of one connected component of its regular locus.
\end{proof}

\begin{lemma}[Proper images of subvarieties are subvarieties]
\label{ana:proper-image}
\uses{ana:geometric-setting}
\lean{TranscendentalNakamaye.Pending.ProperImageAnalytic.exists_image_subvariety}
\leanok
Let \(Z\) be a compact complex manifold and \(M\) a complex manifold of the same
dimension \(n\), and let \(f:Z\to M\) be holomorphic. The image under \(f\) of a closed
irreducible analytic subvariety of \(Z\) is a closed irreducible analytic subvariety
of \(M\).
\end{lemma}
\begin{proof}
\uses{ana:proper-image-analytic,ana:analytic-preimage,lag:regular-locus-dense,fpgeom:complex-analytic-subvariety}
\leanok
Write \(A=f(W)\), which is analytic by
Lemma~\ref{ana:proper-image-analytic}. Beyond that analyticity the statement has no
further classical content.

\(A\) is closed: \(W\) is closed in the compact \(Z\), hence compact, so \(A\) is the
continuous image of a compact set, hence compact and therefore closed in \(M\).

\(A\) is nonempty because \(W\) is.

\(A\) is irreducible. Let \(A=B\cup C\) with \(B,C\) analytic subsets of \(M\)
contained in \(A\). Preimages of analytic sets are analytic
(Lemma~\ref{ana:analytic-preimage}), so \(W\cap f^{-1}(B)\) and \(W\cap f^{-1}(C)\) are
analytic subsets of \(Z\) contained in \(W\), and every point of \(W\) is sent into
\(B\cup C\), so \(W=(W\cap f^{-1}(B))\cup(W\cap f^{-1}(C))\). Irreducibility of \(W\)
makes one of the two pieces all of \(W\); say \(W=W\cap f^{-1}(B)\). Then
\(A=f(W)\subseteq B\), and \(B\subseteq A\) gives \(A=B\).

Finally \(A\) carries a certified dimension: by density of the regular locus
(Theorem~\ref{lag:regular-locus-dense}) any point of \(A\) has a regular point of \(A\)
in every neighbourhood, and the dimension of such a point serves.
\end{proof}

\begin{lemma}[Remmert: the image of an analytic set with compact source is analytic]
\label{ana:proper-image-analytic}
\uses{ana:geometric-setting,fpgeom:complex-analytic-subvariety}
\source{demailly-cadg:II.8.8,remmert57-holomorphe-abbildungen:page-0030}
\lean{TranscendentalNakamaye.Pending.ProperImageAnalytic.remmert1956_image_isComplexAnalyticSubset}
\leanok
Let \(Z\) be a compact complex manifold and \(M\) a complex manifold of the same
dimension \(n\), let \(f:Z\to M\) be holomorphic, and let \(W\subseteq Z\) be a closed
irreducible analytic subvariety. Then \(f(W)\) is an analytic subset of \(M\).
\end{lemma}
\begin{proof}
\uses{rem:proper-mapping}
\leanok
This is Remmert's proper mapping theorem (Remmert, Math.\ Ann.\ \textbf{130} (1956), 410--441, and
Math.\ Ann.\ \textbf{133} (1957), 328--370, \S\,4, Satz~23 and its compact-source form Satz~24, p.~356), proved in Theorem~\ref{rem:proper-mapping}: the image of an analytic set under a
holomorphic map whose restriction to it is proper is analytic. Compactness of \(Z\) makes this hypothesis hold: \(W\) is closed in the
compact \(Z\), so \(W\cap f^{-1}(K)\) is compact for every compact \(K\subset M\).

Only analyticity is asserted. Closedness of \(f(W)\) is a separate and elementary
consequence of compactness, and is not a substitute for analyticity: a compact set need
not be analytic.
\end{proof}

\begin{lemma}[Reduction to a smooth component]
\label{ana:resolution}
\uses{ana:blowups,ana:restricted-volume,ana:locus-properties}
\source{transcendental-nakamaye-copy4-20260911:lem:resolution}
\lean{TranscendentalNakamaye.Pending.SmoothComponentReduction.exists_modification_smooth_component}
\leanok
Let \(\alpha\) be big and let \(V\) be a positive-dimensional
irreducible component of \(E_{\mathrm{nK}}(\alpha)\), with
\(V\not\subset E_{\mathrm{nn}}(\alpha)\). There is a modification
\(\mu:\widetilde X\to X\), an isomorphism over
\(X\setminus V_{\mathrm{sing}}\), whose strict transform
\(\widetilde V\) is a closed connected smooth submanifold. It is an
irreducible component of \(E_{\mathrm{nK}}(\mu^*\alpha)\), is not
contained in \(E_{\mathrm{nn}}(\mu^*\alpha)\), has dimension
\(\dim V\), and satisfies
\[
 \operatorname{RV}_{\widetilde X\mid\widetilde V}(\mu^*\alpha)
 =\operatorname{RV}_{X\mid V}(\alpha).
\]
\end{lemma}
\begin{proof}
\uses{ana:embedded-resolution,ana:component-preimage-maximal,ana:nonkahler-locus-analytic,ana:locus-properties,ana:irreducible-components}
\leanok
Let \(\mu:\widetilde X\to X\) be an embedded resolution of \(V\) as in
Lemma~\ref{ana:embedded-resolution}, and let \(\widetilde V\) be the strict
transform of \(V\); it is a closed connected submanifold of dimension
\(\dim V\). Blowups with smooth centres of a compact K\"ahler manifold are
compact K\"ahler and the pullback of a big class under a modification is big, so
\(\widetilde X\) and \(\mu^*\alpha\) satisfy the standing assumptions.

Since \(\mu\) is an isomorphism over \(X\setminus V_{\mathrm{sing}}\), its
exceptional image satisfies \(\mu(\operatorname{Exc}\mu)\subset
V_{\mathrm{sing}}\). The \(\dim V\)-regular locus of \(V\) is nonempty, so
\(V\not\subset\mu(\operatorname{Exc}\mu)\); Proposition~\ref{ana:locus-properties}
therefore applies to \(V\) and gives both the birational invariance
\(\operatorname{RV}_{\widetilde X\mid\widetilde V}(\mu^*\alpha)
=\operatorname{RV}_{X\mid V}(\alpha)\) and the two transformation rules
\(E_{\mathrm{nn}}(\mu^*\alpha)=\mu^{-1}(E_{\mathrm{nn}}(\alpha))\) and
\(E_{\mathrm{nK}}(\mu^*\alpha)=\mu^{-1}(E_{\mathrm{nK}}(\alpha))\cup
\operatorname{Exc}\mu\). From the first rule,
\(\widetilde V\subset E_{\mathrm{nn}}(\mu^*\alpha)\) would give
\(V=\mu(\widetilde V)\subset E_{\mathrm{nn}}(\alpha)\), which is excluded; so
\(\widetilde V\not\subset E_{\mathrm{nn}}(\mu^*\alpha)\). From the second,
\(\widetilde V\subset\mu^{-1}(V)\subset E_{\mathrm{nK}}(\mu^*\alpha)\).

It remains to see that \(\widetilde V\) is a \emph{maximal} irreducible
subvariety of \(E_{\mathrm{nK}}(\mu^*\alpha)\). Apply
Lemma~\ref{ana:component-preimage-maximal} with \(A=E_{\mathrm{nK}}(\alpha)\), which is analytic
(Theorem~\ref{ana:nonkahler-locus-analytic}), the closed set \(E=\operatorname{Exc}\mu\) and
\(S=E_{\mathrm{nK}}(\mu^*\alpha)=\mu^{-1}(E_{\mathrm{nK}}(\alpha))\cup\operatorname{Exc}\mu\). We have
\(\widetilde V\subset S\) (above), and
\(\mu^{-1}(V)\setminus\operatorname{Exc}\mu\subset\widetilde V\) because \(\widetilde V\) is the closure
of that set. A point \(x\) at which \(V\) is regular of dimension \(\dim V\) lies outside
\(V_{\mathrm{sing}}\supset\mu(\operatorname{Exc}\mu)\); as \(\mu\) is onto, \(x=\mu(y)\) with
\(y\notin\operatorname{Exc}\mu\), and then \(y\in\mu^{-1}(V)\setminus\operatorname{Exc}\mu\subset\widetilde V\).
The lemma shows that \(\widetilde V\) is an irreducible component of \(S\).
\end{proof}
Copy-4 argues the maximality through the image: \(\mu(Z)\) of an irreducible \(Z\) with
\(\widetilde V\subset Z\subset E_{\mathrm{nK}}(\mu^*\alpha)\) is an irreducible analytic subset of
\(E_{\mathrm{nK}}(\alpha)\) containing \(V\), hence equal to \(V\), which needs Remmert's proper mapping
theorem (Lemma~\ref{ana:proper-image-analytic}). The route above never forms \(\mu(Z)\): it uses only
that an irreducible component fills \(E_{\mathrm{nK}}(\alpha)\) near one of its regular points
(Proposition~\ref{ana:component-local-filling}).

\begin{theorem}[Vanishing on components of the non-K\"ahler locus]
\label{ana:component-vanishing}
\uses{ana:nonkahler-locus,ana:restricted-volume,ana:locus-properties,ana:irreducible-components,ana:resolution,ana:principalization-data,ana:local-gluing-data,ana:gluing}
\source{transcendental-nakamaye-copy4-20260911:thm:main,transcendental-nakamaye-copy4-20260911:sec:gluing}
\lean{TranscendentalNakamaye.Pending.ComponentVanishing.numericalRestrictedVolume_eq_zero_of_isIrreducibleComponent}
\leanok
Let \(X\) be a nonempty connected compact K\"ahler manifold of
positive complex dimension, let \(\alpha\in H^{1,1}(X,\mathbb R)\)
be pseudoeffective, and let \(V\) be a positive-dimensional
irreducible component of \(E_{\mathrm{nK}}(\alpha)\). Then
\[
 \operatorname{RV}_{X\mid V}(\alpha)=0.
\]
\end{theorem}
\begin{proof}
\leanok
If \(\alpha\) is not big, then \(E_{\mathrm{nK}}(\alpha)=X\).
A connected complex manifold is irreducible, so \(V=X\). Proposition~
\ref{ana:locus-properties} gives
\(\operatorname{RV}_{X\mid X}(\alpha)=\operatorname{vol}(\alpha)=0\).
Hence assume \(\alpha\) is big. If
\(V\subset E_{\mathrm{nn}}(\alpha)\), the result is the defining zero
case of restricted volume. Otherwise, restricted volume is
nonnegative; if it were nonzero, it would have a value \(h>0\).
Apply Lemma~\ref{ana:resolution}. It preserves \(h\), bigness, the
component condition, and the condition of not lying in the non-nef
locus. We may therefore replace \((X,\alpha,V)\) by that resolved
triple and assume \(V\) is a closed connected smooth submanifold.
Since the non-K\"ahler locus of a big class is proper,
\(0<\dim V<\dim X\).

Use Proposition~\ref{ana:principalization-data} to obtain
\(\pi:Y\to X\), \(W\), \(y\in W\), \(U\), and the global and local
potentials \(\varphi_G,\varphi_U\). Lemma~\ref{ana:local-gluing-data}
chooses compactly nested neighborhoods \(U_1\Subset U_2\Subset U\)
on which these potentials satisfy Proposition~\ref{ana:gluing} in
the class \(\pi^*\alpha\). The proposition therefore gives a K\"ahler
current in \(\pi^*\alpha\) which is smooth near \(y\), so
\(y\notin E_{\mathrm{nK}}(\pi^*\alpha)\).

On the other hand \(W\subset\pi^{-1}(V)\), so Proposition~
\ref{ana:locus-properties} gives
\[
 W\subset\pi^{-1}(V)
 \subset\pi^{-1}(E_{\mathrm{nK}}(\alpha))
 \subset E_{\mathrm{nK}}(\pi^*\alpha).
\]
This contradicts the choice of \(y\). Thus positive restricted
volume is impossible, and its nonnegativity implies the asserted
vanishing.
\end{proof}

\begin{theorem}[Equality of the non-K\"ahler locus and null locus]
\label{ana:locus-equality}
\uses{ana:nonkahler-locus,ana:restricted-volume,ana:locus-properties,ana:irreducible-components,ana:no-isolated-points,ana:component-vanishing}
\source{transcendental-nakamaye-copy4-20260911:thm:main,transcendental-nakamaye-copy4-20260911:sec:gluing}
\lean{TranscendentalNakamaye.Pending.NakamayeMainTheorem.nonKahlerLocus_eq_nullLocus}
\leanok
Let \(X\) be a nonempty connected compact K\"ahler manifold of
positive complex dimension and let \(\alpha\) be a pseudoeffective
real \((1,1)\)-class. Then
\[
 E_{\mathrm{nK}}(\alpha)=\operatorname{Null}(\alpha).
\]
\end{theorem}
\begin{proof}
\leanok
Suppose first that \(\alpha\) is not big. Then
\(E_{\mathrm{nK}}(\alpha)=X\), and
\(\operatorname{RV}_{X\mid X}(\alpha)=0\). The connected manifold
\(X\) is irreducible and has positive dimension, so it is one of the
subvarieties in the defining union for \(\operatorname{Null}(\alpha)\).
That union is contained in \(X\), hence equals \(X\).

Now suppose that \(\alpha\) is big. Every point of its non-K\"ahler
locus lies on an irreducible component of this analytic subset.
By Lemma~\ref{ana:no-isolated-points} the subset has no isolated points, so every such component has
positive dimension. Theorem~\ref{ana:component-vanishing} gives zero
restricted volume on each component; therefore each component occurs
in the defining union for \(\operatorname{Null}(\alpha)\). This proves
\(E_{\mathrm{nK}}(\alpha)\subset\operatorname{Null}(\alpha)\).

For the reverse inclusion, let \(V\) be any subvariety occurring in
that union. If \(V\not\subset E_{\mathrm{nK}}(\alpha)\),
Proposition~\ref{ana:locus-properties} would imply
\(\operatorname{RV}_{X\mid V}(\alpha)>0\), contrary to its defining
zero volume. Thus every such \(V\) is contained in
\(E_{\mathrm{nK}}(\alpha)\). Taking their union gives the reverse
inclusion and proves the equality.
\end{proof}
